\PassOptionsToPackage{hyperfootnotes=false}{hyperref}
\input{preamble_ptbr}
\ifdefined\pdfinfoomitdate\pdfinfoomitdate=1\fi
\ifdefined\pdftrailerid\pdftrailerid{}\fi
\ifdefined\pdfsuppressptexinfo\pdfsuppressptexinfo=15\fi
\setlength{\emergencystretch}{3em}

\hypersetup{
  pdftitle={Elementos de geometria algébrica III, primeira parte},
  pdfauthor={Alexander Grothendieck e Jean Dieudonné; tradução brasileira},
  pdfsubject={Estudo cohomológico dos feixes coerentes}
}

\begin{document}

\title{Elementos de geometria algébrica\\
III. Estudo cohomológico dos feixes coerentes\\
Primeira parte}
\author{Alexander Grothendieck e Jean Dieudonné}
\date{}
\maketitle

\begin{center}
\large Tradução brasileira
\end{center}

\phantomsection
\label{section:ega0}
\pdfbookmark[0]{Capítulo 0 (continuação)}{chapter.0.continued}
\addcontentsline{toc}{part}{Capítulo 0 (continuação)}
\section*{CAPÍTULO 0 (continuação)}
\setcounter{section}{7}
\setcounter{subsection}{0}
\setcounter{equation}{0}

\section{Funtores representáveis}
\label{section:0.8}

\subsection{Funtores representáveis}
\label{subsection:0.8.1}

\oldpage[0\textsubscript{III}]{5}
\begin{env}[8.1.1]
\label{0.8.1.1}
Denotamos por $\Set$ a categoria de conjuntos.
Seja $\C$ uma categoria; para dois objetos $X$, $Y$ de $\C$, pomos $h_X(Y)=\Hom(Y,X)$; para cada morfismo $u:Y\to Y'$ de $\C$, denotamos por $h_X(u)$ a aplicação $v\mapsto vu$ de $\Hom(Y',X)$ em $\Hom(Y,X)$.
É imediato que, com estas definições, $h_X:\C\to\Set$ é um \emph{funtor contravariante}, isto é, um objeto da categoria $\CHom(\C\op,\Set)$ dos funtores covariantes da categoria $\C\op$ (a categoria dual de $\C$) na categoria $\Set$ (T, 1.7, (d) e \cite{III-29}).
\end{env}

\begin{env}[8.1.2]
\label{0.8.1.2}
Seja agora $w:X\to X'$ um morfismo de $\C$; para cada $Y\in\C$ e cada $v\in\Hom(Y,X)=h_X(Y)$,
temos $wv\in\Hom(Y,X')=h_{X'}(Y)$; denotamos por $h_w(Y)$ a aplicação $v\mapsto wv$ de
$h_X(Y)$ em $h_{X'}(Y)$. É imediato que, para cada morfismo $u:Y\to Y'$ de $\C$, o
diagrama
\[
  \xymatrix{
    h_X(Y')\ar[r]^{h_X(u)}\ar[d]_{h_w(Y')} &
    h_X(Y)\ar[d]^{h_w(Y)}\\
    h_{X'}(Y')\ar[r]_{h_{X'}(u)} &
    h_{X'}(Y)
  }
\]
é comutativo; em outros termos, $h_w$ é uma \emph{transformação natural (ou morfismo funtorial)
$h_X\to h_{X'}$} (T, 1.2), e também um morfismo da categoria $\CHom(\C\op,\Set)$ (T, 1.7, (d)). As
definições de $h_X$ e de $h_w$ constituem, portanto, a definição de um \emph{funtor covariante
canônico}
\[
  h:\C\to\CHom(\C\op,\Set),\quad X\mapsto h_X.
  \tag{8.1.2.1}
\]
\end{env}

\begin{env}[8.1.3]
\label{0.8.1.3}
Seja $X$ um objeto de $\C$, e seja $F$ um funtor contravariante de $\C$ em $\Set$
(um objeto de $\CHom(\C\op,\Set)$). Seja $g:h_X\to F$ uma \emph{transformação natural}: para todo $Y\in\C$,
\oldpage[0\textsubscript{III}]{6}
$g(Y)$ é, pois, uma aplicação $h_X(Y)\to F(Y)$ tal que, para cada morfismo $u:Y\to Y'$ de $\C$,
o diagrama
\[
  \xymatrix{
    h_X(Y')\ar[r]^{h_X(u)}\ar[d]_{g(Y')} &
    h_X(Y)\ar[d]^{g(Y)}\\
    F(Y')\ar[r]_{F(u)} &
    F(Y)
  }
  \tag{8.1.3.1}\label{0.8.1.3.1}
\]
é comutativo. Em particular, temos uma aplicação $g(X):h_X(X)=\Hom(X,X)\to F(X)$ e, por conseguinte, um elemento
\[
  \alpha(g)=(g(X))(1_X)\in F(X)
  \tag{8.1.3.2}
\]
e, consequentemente, uma aplicação canônica
\[
  \alpha:\Hom(h_X,F)\to F(X).
  \tag{8.1.3.3}
\]

Reciprocamente, consideremos um elemento $\xi\in F(X)$; para cada morfismo $v:Y\to X$ de $\C$, $F(v)$ é uma
aplicação $F(X)\to F(Y)$; consideremos a aplicação
\[
  v\mapsto(F(v))(\xi)
  \tag{8.1.3.4}
\]
de $h_X(Y)$ em $F(Y)$; se denotarmos esta aplicação por $(\beta(\xi))(Y)$,
\[
  \beta(\xi):h_X\to F
  \tag{8.1.3.5}
\]
é uma \emph{transformação natural}, pois para cada morfismo $u:Y\to Y'$ de $\C$ temos
$(F(vu))(\xi)=(F(u)\circ F(v))(\xi)$, o que torna (\hyperref[0.8.1.3.1]{8.1.3.1}) comutativo para $g=\beta(\xi)$.
Temos definido assim uma aplicação canônica
\[
  \beta:F(X)\to\Hom(h_X,F).
  \tag{8.1.3.6}
\]
\end{env}

\begin{proposition}[8.1.4]
\label{0.8.1.4}
As aplicações $\alpha$ e $\beta$ são bijeções inversas uma da outra.
\end{proposition}

\begin{proof}
Calculemos $\alpha(\beta(\xi))$ para $\xi\in F(X)$; para cada $Y\in\C$, $(\beta(\xi))(Y)$ é a aplicação
$g_1(Y):v\mapsto(F(v))(\xi)$ de $h_X(Y)$ em $F(Y)$. Temos, pois,
\[
  \alpha(\beta(\xi))=(g_1(X))(1_X)=(F(1_X))(\xi)=1_{F(X)}(\xi)=\xi.
\]
Calculemos agora $\beta(\alpha(g))$ para $g\in\Hom(h_X,F)$; para cada $Y\in\C$, $(\beta(\alpha(g)))(Y)$
é a aplicação $v\mapsto(F(v))((g(X))(1_X))$; pela comutatividade de (\hyperref[0.8.1.3.1]{8.1.3.1}), esta aplicação
não é senão $v\mapsto(g(Y))((h_X(v))(1_X))=(g(Y))(v)$, por definição de $h_X(v)$; em outros termos,
é igual a $g(Y)$, o que termina a demonstração.
\end{proof}

\begin{env}[8.1.5]
\label{0.8.1.5}
Recordemos que uma \emph{subcategoria} $\C'$ de uma categoria $\C$ se define pela condição de que seus objetos
sejam objetos de $\C$, e de que, se $X'$, $Y'$ são dois objetos de $\C'$, o conjunto
$\Hom_{\C'}(X',Y')$ de morfismos $X'\to Y'$ \emph{em $\C'$} seja um subconjunto do conjunto $\Hom_\C(X',Y')$ de
morfismos $X'\to Y'$ \emph{em $\C$}, sendo a aplicação canônica de ``composição de morfismos''
\[
  \Hom_{\C'}(X',Y')\times\Hom_{\C'}(Y',Z')\to\Hom_{\C'}(X',Z')
\]
\oldpage[0\textsubscript{III}]{7}
a restrição da aplicação canônica
\[
  \Hom_\C(X',Y')\times\Hom_\C(Y',Z')\to\Hom_\C(X',Z').
\]

Dizemos que $\C'$ é uma subcategoria \emph{plena} de $\C$ se $\Hom_{\C'}(X',Y')=\Hom_\C(X',Y')$ para todo
par de objetos de $\C'$. A subcategoria $\C''$ de $\C$ formada pelos objetos de $\C$ isomorfos a
objetos de $\C'$ é então também uma subcategoria plena de $\C$, \emph{equivalente} (T, 1.2) a $\C'$, como se
verifica facilmente.

Um funtor covariante $F:\C_1\to\C_2$ chama-se \emph{plenamente fiel} se, para todo par de objetos
$X_1$, $Y_1$ de $\C_1$, a aplicação $u\mapsto F(u)$ de $\Hom(X_1,Y_1)$ em $\Hom(F(X_1),F(Y_1))$ é
\emph{bijetiva}; isto implica que a subcategoria $F(\C_1)$ de $\C_2$ é \emph{plena}. Além disso,
se dois objetos $X_1$, $X_1'$ têm a mesma imagem $X_2$, existe um único isomorfismo
$u:X_1\to X_1'$ tal que $F(u)=1_{X_2}$. Para cada objeto $X_2$ de $F(\C_1)$, seja $G(X_2)$ um dos
os objetos $X_1$ de $\C_1$ tais que $F(X_1)=X_2$ ($G$ se define mediante o axioma da escolha); para
cada morfismo $v:X_2\to Y_2$ de $F(\C_1)$, $G(v)$ será o único morfismo $u:G(X_2)\to G(Y_2)$ tal
que $F(u)=v$; $G$ é então um \emph{funtor} de $F(\C_1)$ em $\C_1$; $FG$ é o funtor identidade de
$F(\C_1)$, e o anterior mostra que existe um isomorfismo de funtores $\vphi:1_{\C_1}\to GF$
tal que $F$, $G$, $\vphi$ e a identidade $1_{F(\C_1)}\to FG$ definem uma \emph{equivalência} entre
a categoria $\C_1$ e a subcategoria plena $F(\C_1)$ de $\C_2$ (T, 1.2).
\end{env}

\begin{env}[8.1.6]
\label{0.8.1.6}
Aplicamos a Proposição \sref{0.8.1.4} ao caso em que $F$ é $h_{X'}$, sendo $X'$ um
objeto qualquer de $\C$; a aplicação $\beta:\Hom(X,X')\to\Hom(h_X,h_{X'})$ não é senão a aplicação
$w\mapsto h_w$ definida em \sref{0.8.1.2}; uma vez que esta aplicação é \emph{bijetiva}, vemos,
com a terminologia de \sref{0.8.1.5}, que:
\end{env}

\begin{proposition}[8.1.7]
\label{0.8.1.7}
O funtor canônico $h:\C\to\CHom(\C\op,\Set)$ é plenamente fiel.
\end{proposition}

\begin{env}[8.1.8]
\label{0.8.1.8}
Seja $F$ um funtor contravariante de $\C$ em $\Set$; dizemos que $F$ é \emph{representável} se
existe um objeto $X\in\C$ tal que $F$ é \emph{isomorfo} a $h_X$; da
Proposição \sref{0.8.1.7} segue-se que os dados de um $X\in\C$ e de um isomorfismo de funtores
$g:h_X\to F$ determinam $X$ salvo isomorfismo único. A Proposição \sref{0.8.1.7} implica então
que $h$ define uma \emph{equivalência} entre $\C$ e a subcategoria plena de
$\CHom(\C\op,\Set)$ formada pelos \emph{funtores contravariantes representáveis}. Da
Proposição \sref{0.8.1.4} segue-se que o dado de uma transformação natural $g:h_X\to F$ é
equivalente ao de um elemento $\xi\in F(X)$; dizer que $g$ é um \emph{isomorfismo} equivale
à seguinte condição sobre $\xi$: \emph{para todo objeto $Y$ de $\C$, a aplicação
$v\mapsto(F(v))(\xi)$ de $\Hom(Y,X)$ em $F(Y)$ é bijetiva}. Quando $\xi$ satisfaz esta condição,
dizemos que o par $(X,\xi)$ \emph{representa} o funtor representável $F$. Por abuso de linguagem,
dizemos também que o objeto $X\in\C$ representa $F$ se existe um $\xi\in F(X)$ tal que
$(X,\xi)$ representa $F$, isto é, se $h_X$ é isomorfo a $F$.

Sejam $F$, $F'$ dois funtores contravariantes representáveis de $\C$ em $\Set$, e sejam $h_X\to F$ e
$h_{X'}\to F'$ dois isomorfismos de funtores. Segue-se então de \sref{0.8.1.6} que
existe uma correspondência bijetiva canônica entre $\Hom(X,X')$ e o conjunto $\Hom(F,F')$ de
transformações naturais $F\to F'$.
\end{env}

\begin{env}[8.1.9]
\label{0.8.1.9}
\textit{Exemplo I. Limites projetivos}.
A noção de funtor contravariante representável compreende, em particular, a noção ``dual'' da noção usual de ``solução de um problema universal''.
\oldpage[0\textsubscript{III}]{8}
Mais geralmente, veremos que a noção de \emph{limite projetivo} é um caso particular da noção de funtor representável.
Recordemos que, em uma categoria $\C$, define-se um \emph{sistema projetivo} pelos dados de um conjunto pré-ordenado $I$, de uma família $(A_\alpha)_{\alpha\in I}$ de objetos de $\C$ e, para cada par de índices $(\alpha,\beta)$ tal que $\alpha\leq\beta$, de um morfismo $u_{\alpha\beta}:A_\beta\to A_\alpha$.
Um \emph{limite projetivo} deste sistema em $\C$ consiste em um objeto $B$ de $\C$ (denotado $\varprojlim A_\alpha$) e, para cada $\alpha\in I$, um morfismo $u_\alpha:B\to A_\alpha$ tais que: 1\textsuperscript{o}. $u_\alpha=u_{\alpha\beta}u_\beta$ para $\alpha\leq\beta$; 2\textsuperscript{o}. para todo objeto $X$ de $\C$ e toda família $(v_\alpha)_{\alpha\in I}$ de morfismos $v_\alpha:X\to A_\alpha$ tal que $v_\alpha=u_{\alpha\beta}v_\beta$ para $\alpha\leq\beta$, existe um único morfismo $v:X\to B$ (denotado $\varprojlim v_\alpha$) tal que $v_\alpha=u_\alpha v$ para todo $\alpha\in I$ (T, 1.8).
Isto pode ser interpretado da seguinte maneira: os $u_{\alpha\beta}$ definem canonicamente aplicações
\[
  \overline{u}_{\alpha\beta}:\Hom(X,A_\beta)\to\Hom(X,A_\alpha)
\]
que definem um \emph{sistema projetivo} de conjuntos $(\Hom(X,A_\alpha),\overline{u}_{\alpha\beta})$, e $(v_\alpha)$ é, por definição, um elemento do conjunto $\varprojlim\Hom(X,A_\alpha)$; é claro que $X\mapsto\varprojlim\Hom(X,A_\alpha)$ é um \emph{funtor contravariante} de $\C$ em $\Set$, e a existência do limite projetivo $B$ equivale a dizer que $(v_\alpha)\mapsto\varprojlim v_\alpha$ é um \emph{isomorfismo} de funtores em $X$
\[
  \varprojlim\Hom(X,A_\alpha)\isoto\Hom(X,B),
  \tag{8.1.9.1}
\]
isto é, que o funtor $X\mapsto\varprojlim\Hom(X,A_\alpha)$ é \emph{representável}.
\end{env}

\begin{env}[8.1.10]
\label{0.8.1.10}
\textit{Exemplo II. Objetos finais}.
Seja $\C$ uma categoria, e seja $\{a\}$ um conjunto unitário.
Consideremos o funtor contravariante $F:\C\to\Set$ que envia cada objeto $X$ de $\C$ ao conjunto $\{a\}$, e cada morfismo $X\to X'$ de $\C$ à única aplicação $\{a\}\to\{a\}$.
Dizer que este funtor é \emph{representável} significa que existe um objeto $e\in\C$ tal que, para todo $Y\in\C$, $\Hom(Y,e)=h_e(Y)$ é um \emph{conjunto unitário}; dizemos que $e$ é um \emph{objeto final} de $\C$, e é claro que dois objetos finais de $\C$ são isomorfos (o que permite definir, em geral mediante o axioma da escolha, \emph{um} objeto final de $\C$, que denotamos então por $e_\C$).
Por exemplo, na categoria $\Set$, os objetos finais são os conjuntos unitários; na categoria de \emph{álgebras aumentadas} sobre um corpo $K$ (onde os morfismos são os homomorfismos de álgebras compatíveis com o aumento), $K$ é um objeto final; na categoria de \emph{$S$-esquemas} \sref[I]{I.2.5.1}, $S$ é um objeto final.
\end{env}

\begin{env}[8.1.11]
\label{0.8.1.11}
Para dois objetos $X$ e $Y$ de uma categoria $\C$, ponhamos $h_X'(Y)=\Hom(X,Y)$ e, para todo morfismo $u:Y\to Y'$, seja $h_X'(u)$ a aplicação $v\mapsto uv$ de $\Hom(X,Y)$ em $\Hom(X,Y')$; $h_X'$ é então um \emph{funtor covariante $\C\to\Set$}, de modo que deduzimos, como em \sref{0.8.1.2}, a definição de um funtor covariante canônico $h':\C\op\to\CHom(\C,\Set)$; um funtor \emph{covariante} $F$ de $\C$ em $\Set$, isto é, um objeto de $\CHom(\C,\Set)$, é então \emph{representável} se existe um objeto $X\in\C$ (necessariamente único salvo isomorfismo único) tal que $F$ é \emph{isomorfo a $h_X'$}; deixamos ao leitor o desenvolvimento das noções ``duais'' das anteriores, que desta vez compreendem a noção de \emph{limite indutivo}, e em particular a noção usual de ``solução de um problema universal''.
\end{env}

\subsection{Estruturas algébricas em categorias}
\label{subsection:0.8.2}

\begin{env}[8.2.1]
\label{0.8.2.1}
\oldpage[0\textsubscript{III}]{9}
Dados dois funtores contravariantes $F$ e $F'$ de $\C$ em $\Set$, recordemos que, para todo objeto $Y\in\C$, pomos $(F\times F')(Y)=F(Y)\times F'(Y)$, e para todo morfismo $u:Y\to Y'$ de $\C$, pomos $(F\times F')(u)=F(u)\times F'(u)$, que é a aplicação $(t,t')\mapsto(F(u)(t),F'(u)(t'))$ de $F(Y')\times F'(Y')$ em $F(Y)\times F'(Y)$; $F\times F':\C\to\Set$ é assim um \emph{funtor contravariante} (que não é senão o \emph{produto} dos objetos $F$ e $F'$ na categoria $\CHom(\C\op,\Set)$).
Dado um objeto $X\in\C$, chamamos \emph{lei de composição interna} sobre $X$ a uma \emph{transformação natural}
\[
\label{0.8.2.1.1}
  \gamma_X:h_X\times h_X\to h_X.
  \tag{8.2.1.1}
\]

Em outros termos (T, 1.2), para todo objeto $Y\in\C$, $\gamma_X(Y)$ é uma aplicação $h_X(Y)\times h_X(Y)\to h_X(Y)$ (portanto, por definição, uma \emph{lei de composição interna} sobre o conjunto $h_X(Y)$) com a condição de que, para todo morfismo $u:Y\to Y'$ de $\C$, o diagrama
\[
  \xymatrix{
    h_X(Y')\times h_X(Y')\ar[rr]^{h_X(u)\times h_X(u)}\ar[dd]_{\gamma_X(Y')} & &
    h_X(Y)\times h_X(Y)\ar[dd]^{\gamma_X(Y)}\\\\
    h_X(Y')\ar[rr]_{h_X(u)} & &
    h_X(Y)
  }
\]
seja comutativo; isto implica que, para as leis de composição $\gamma_X(Y)$ e $\gamma_X(Y')$, $h_X(u)$ é um \emph{homomorfismo} de $h_X(Y')$ em $h_X(Y)$.

De maneira análoga, dados dois objetos $Z$ e $X$ de $\C$, chamamos \emph{lei de composição externa sobre $X$, com $Z$ como domínio de operadores} a uma transformação natural
\[
\label{0.8.2.1.2}
  \omega_{X,Z}:h_Z\times h_X\to h_X.
  \tag{8.2.1.2}
\]

Vemos, como antes, que, para todo $Y\in\C$, $\omega_{X,Z}(Y)$ é uma lei de composição externa sobre $h_X(Y)$, com $h_Z(Y)$ como domínio de operadores, e tal que, para todo morfismo $u:Y\to Y'$, $h_X(u)$ e $h_Z(u)$ formam um \emph{di-homomorfismo} de $(h_Z(Y'),h_X(Y'))$ em $(h_Z(Y),h_X(Y))$.
\end{env}

\begin{env}[8.2.2]
\label{0.8.2.2}
Seja $X'$ um segundo objeto de $\C$, e suponhamos dada uma lei de composição interna $\gamma_{X'}$ sobre $X'$; dizemos que um morfismo $w:X\to X'$ de $\C$ é um \emph{homomorfismo} para as leis de composição se, para todo $Y\in\C$, $h_w(Y):h_X(Y)\to h_{X'}(Y)$ é um \emph{homomorfismo} para as leis de composição $\gamma_X(Y)$ e $\gamma_{X'}(Y)$.
Se $X''$ é um terceiro
\oldpage[0\textsubscript{III}]{10}
objeto de $\C$ munido de uma lei de composição interna $\gamma_{X''}$, e $w':X'\to X''$ é um morfismo de $\C$ que é um homomorfismo para $\gamma_{X'}$ e $\gamma_{X''}$, é claro que o morfismo $w'w:X\to X''$ é um homomorfismo para as leis de composição $\gamma_X$ e $\gamma_{X''}$.
Um isomorfismo $w:X\isoto X'$ de $\C$ chama-se \emph{isomorfismo para as leis de composição $\gamma_X$ e $\gamma_{X'}$} se $w$ é um homomorfismo para estas leis de composição, e se seu morfismo inverso $w^{-1}$ é um homomorfismo para as leis de composição $\gamma_{X'}$ e $\gamma_X$.

Definimos de maneira análoga os \emph{di-homomorfismos} para pares de objetos de $\C$ munidos de leis de composição externa.
\end{env}

\begin{env}[8.2.3]
\label{0.8.2.3}
Quando uma lei de composição interna $\gamma_X$ sobre um objeto $X\in\C$ é tal que $\gamma_X(Y)$ é uma \emph{lei de grupo} sobre $h_X(Y)$ para \emph{todo} $Y\in\C$, dizemos que $X$, munido desta lei, é um \emph{$\C$-grupo} ou um \emph{objeto grupo de $\C$}.
Definimos analogamente os \emph{$\C$-anéis}, \emph{$\C$-módulos}, etc.
\end{env}

\begin{env}[8.2.4]
\label{0.8.2.4}
Suponhamos que o \emph{produto} $X\times X$ de um objeto $X\in\C$ por si mesmo existe em $\C$; por definição, temos então $h_{X\times X}=h_X\times h_X$ salvo isomorfismo canônico, pois é um caso particular do limite projetivo \sref{0.8.1.9}; uma lei de composição interna sobre $X$ pode, portanto, ser considerada como um morfismo funtorial $\gamma_X:h_{X\times X}\to h_X$, e determina assim canonicamente \sref{0.8.1.6} um elemento $c_X\in\Hom(X\times X,X)$ tal que $h_{c_X}=\gamma_X$; neste caso, o dado de uma lei de composição interna sobre $X$ equivale ao dado de um morfismo $X\times X\to X$; quando $\C$ é a categoria $\Set$, recuperamos a noção clássica de lei de composição interna sobre um conjunto.
Temos um resultado análogo para uma lei de composição externa quando o produto $Z\times X$ existe em $\C$.
\end{env}

\begin{env}[8.2.5]
\label{0.8.2.5}
Com a notação anterior, suponhamos além disso que $X\times X\times X$ existe em $\C$; a caracterização do produto como objeto que representa um funtor \sref{0.8.1.9} implica a existência de isomorfismos canônicos
\[
  (X\times X)\times X\isoto X\times X\times X\isoto X\times(X\times X);
\]
se identificarmos canonicamente $X\times X\times X$ com $(X\times X)\times X$, a aplicação $\gamma_X(Y)\times 1_{h_X(Y)}$ identifica-se com $h_{c_X\times 1_X}(Y)$ para todo $Y\in\C$.
Consequentemente, é equivalente dizer que, para todo $Y\in\C$, a lei interna $\gamma_X(Y)$ é associativa, que o diagrama de aplicações
\[
  \xymatrix{
    h_X(Y)\times h_X(Y)\times h_X(Y)\ar[rr]^{\gamma_X(Y)\times 1}\ar[dd]_{1\times\gamma_X(Y)} & &
    h_X(Y)\times h_X(Y)\ar[dd]^{\gamma_X(Y)}\\\\
    h_X(Y)\times h_X(Y)\ar[rr]_{\gamma_X(Y)} & &
    h_X(Y)
  }
\]
\oldpage[0\textsubscript{III}]{11}
é comutativo, e que o diagrama de morfismos
\[
  \xymatrix{
    X\times X\times X\ar[rr]^{c_X\times 1_X}\ar[dd]_{1_X\times c_X} & &
    X\times X\ar[dd]^{c_X}\\\\
    X\times X\ar[rr]_{c_X} & &
    X
  }
\]
é comutativo.
\end{env}

\begin{env}[8.2.6]
\label{0.8.2.6}
Sob as hipóteses de \sref{0.8.2.5}, se quisermos expressar, para todo $Y\in\C$, que a lei interna $\gamma_X(Y)$ é uma \emph{lei de grupo}, é necessário, em primeiro lugar, que seja associativa e, em segundo lugar, que exista uma aplicação $\alpha_X(Y):h_X(Y)\to h_X(Y)$ com as propriedades da operação de \emph{inverso} de um grupo; como vimos que, para todo morfismo $u:Y\to Y'$ de $\C$, $h_X(u)$ deve ser um homomorfismo de grupos $h_X(Y')\to h_X(Y)$, vemos primeiro que $\alpha_X:h_X\to h_X$ deve ser uma \emph{transformação natural}.
Por outro lado, podem-se expressar as propriedades características do inverso $s\mapsto s^{-1}$ de um grupo $G$ sem fazer intervir o elemento identidade: basta verificar que as duas aplicações compostas
\[
  (s,t)\mapsto(s,s^{-1},t)\mapsto(s,s^{-1}t)\mapsto s(s^{-1}t),
\]
\[
  (s,t)\mapsto(s,s^{-1},t)\mapsto(s,ts^{-1})\mapsto(ts^{-1})s
\]
são iguais à segunda projeção $(s,t)\mapsto t$ de $G\times G$ em $G$.
Por \sref{0.8.1.3}, temos $\alpha_X=h_{a_X}$, onde $a_X\in\Hom(X,X)$; a primeira condição anterior expressa então que o morfismo composto
\[
  X\times X\xrightarrow{(1_X,a_X)\times 1_X}X\times X\times X\xrightarrow{1_X\times c_X}X\times X\xrightarrow{c_X}X
\]
é a segunda projeção $X\times X\to X$ em $\C$, e a segunda condição é análoga.
\end{env}

\begin{env}[8.2.7]
\label{0.8.2.7}
Suponhamos agora que existe um \emph{objeto final} $e$ \sref{0.8.1.10} em $\C$.
Suponhamos sempre que $\gamma_X(Y)$ é uma lei de grupo sobre $h_X(Y)$ para todo $Y\in\C$, e denotemos por $\eta_X(Y)$ o elemento identidade de $\gamma_X(Y)$.
Como, para todo morfismo $u:Y\to Y'$ de $\C$, $h_X(u)$ é um homomorfismo de grupos, temos $\eta_X(Y)=(h_X(u))(\eta_X(Y'))$; tomando em particular $Y'=e$, caso no qual $u$ é o único elemento $\varepsilon$ de $\Hom(Y,e)$, vemos que o elemento $\eta_X(e)$ determina por completo $\eta_X(Y)$ para todo $Y\in\C$.
Ponhamos $e_X=\eta_X(X)$, o elemento identidade do grupo $h_X(X)=\Hom(X,X)$; a comutatividade do diagrama
\[
  \xymatrix{
    h_X(e)\ar[r]^{h_X(\varepsilon)}\ar[d]_{h_{e_X}(e)} &
    h_X(Y)\ar[d]^{h_{e_X}(Y)}\\
    h_X(e)\ar[r]_{h_X(\varepsilon)} &
    h_X(Y)
  }
\]
(cf. \sref{0.8.1.2}) mostra que, sobre o conjunto $h_X(Y)$, a aplicação $h_{e_X}(Y)$ não é senão
\oldpage[0\textsubscript{III}]{12}
$s\mapsto\eta_X(Y)$, que envia todo elemento ao elemento identidade.
Comprovamos então que o fato de $\eta_X(Y)$ ser o elemento identidade de $\gamma_X(Y)$ para todo $Y\in\C$ equivale a dizer que o morfismo composto
\[
  X\xrightarrow{(1_X,1_X)}X\times X\xrightarrow{1_X\times e_X}X\times X\xrightarrow{c_X}X,
\]
e seu análogo no qual se intercambiam $1_X$ e $e_X$, são ambos \emph{iguais} a $1_X$.
\end{env}

\begin{env}[8.2.8]
\label{0.8.2.8}
Naturalmente, poder-se-iam ampliar com facilidade os exemplos de estruturas algébricas em categorias.
O exemplo dos grupos foi tratado com suficiente detalhe, mas daqui em diante deixaremos, em geral, ao leitor o desenvolvimento das noções análogas para os exemplos de estruturas algébricas que encontrarmos.
\end{env}

\section{Conjuntos construtíveis}
\label{section:0.9}

\subsection{Conjuntos construtíveis}
\label{subsection:0.9.1}

\begin{definition}[9.1.1]
\label{0.9.1.1}
Dizemos que uma aplicação contínua $f:X\to Y$ é \emph{quase-compacta} se, para todo subconjunto aberto quase-compacto $U$ de $Y$, $f^{-1}(U)$ é quase-compacto.
Dizemos que um subconjunto $Z$ de um espaço topológico $X$ é \emph{retrocompacto} em $X$ se a injeção canônica $Z\to X$ é quase-compacta; em outros termos, se para todo subconjunto aberto quase-compacto $U$ de $X$, $U\cap Z$ é quase-compacto.
\end{definition}

Um subconjunto \emph{fechado} de $X$ é retrocompacto em $X$, mas um subconjunto quase-compacto de $X$ não é necessariamente retrocompacto em $X$.
Se $X$ é quase-compacto, todo subconjunto aberto retrocompacto de $X$ é quase-compacto.
É claro que toda união \emph{finita} de conjuntos retrocompactos em $X$ é retrocompacta em $X$, pois toda união finita de conjuntos quase-compactos é quase-compacta.
Toda interseção finita de abertos retrocompactos de $X$ é um aberto retrocompacto de $X$.
Em um espaço \emph{localmente noetheriano} $X$, todo conjunto quase-compacto é um subespaço noetheriano e, consequentemente, \emph{todo subconjunto} de $X$ é retrocompacto em $X$.

\begin{definition}[9.1.2]
\label{0.9.1.2}
Dado um espaço topológico $X$, dizemos que um subconjunto de $X$ é \emph{construtível} se pertence ao menor conjunto de subconjuntos $\mathfrak{F}$ de $X$ que contém todos os subconjuntos abertos retrocompactos de $X$ e é estável por interseções finitas e complementares (o que implica que $\mathfrak{F}$ também é estável por uniões finitas).
\end{definition}

\begin{proposition}[9.1.3]
\label{0.9.1.3}
Para que um subconjunto de $X$ seja construtível, é necessário e suficiente que seja uma união finita de conjuntos da forma $U\cap\complement{V}$, onde $U$ e $V$ são abertos retrocompactos de $X$.
\end{proposition}

\begin{proof}
É claro que a condição é suficiente.
Para ver que é necessária, consideremos o conjunto $\mathfrak{G}$ das uniões finitas de conjuntos da forma $U\cap\complement{V}$, onde $U$ e $V$ são abertos retrocompactos de $X$; basta ver que o complementar de todo conjunto de $\mathfrak{G}$ pertence a $\mathfrak{G}$.
Seja $Z=\bigcup_{i\in I}(U_i\cap\complement{V_i})$, onde $I$ é finito e os $U_i$ e $V_i$ são abertos retrocompactos de $X$; temos
$\complement{Z}=\bigcap_{i\in I}(V_i\cup\complement{U_i})$, de modo que $\complement{Z}$ é uma união finita de conjuntos que são interseções de certo número dos $V_i$ e de certo número dos $\complement{U_i}$, portanto da forma
\oldpage[0\textsubscript{III}]{13}
$V\cap\complement{U}$, onde $U$ é a união de certo número dos $U_i$ e $V$ é a interseção de certo número dos $V_i$; mas observamos anteriormente que as uniões e interseções finitas de abertos retrocompactos de $X$ são abertos retrocompactos de $X$, de onde se segue a conclusão.
\end{proof}

\begin{corollary}[9.1.4]
\label{0.9.1.4}
Todo subconjunto construtível de $X$ é retrocompacto em $X$.
\end{corollary}

\begin{proof}
Basta demonstrar que, se $U$ e $V$ são abertos retrocompactos de $X$, então $U\cap\complement{V}$ é retrocompacto em $X$; se $W$ é um aberto quase-compacto de $X$, então $W\cap U\cap\complement{V}$ é fechado no espaço quase-compacto $W\cap U$ e, portanto, é quase-compacto.
\end{proof}

Em particular:
\begin{corollary}[9.1.5]
\label{0.9.1.5}
Para que um subconjunto aberto $U$ de $X$ seja construtível, é necessário e suficiente que seja retrocompacto em $X$.
Para que um subconjunto fechado $F$ de $X$ seja construtível, é necessário e suficiente que o aberto $\complement{F}$ seja retrocompacto.
\end{corollary}

\begin{env}[9.1.6]
\label{0.9.1.6}
Um caso importante é aquele em que todo subconjunto aberto quase-compacto de $X$ é retrocompacto; em outros termos, aquele em que a interseção de dois subconjuntos abertos quase-compactos de $X$ é quase-compacta (cf.~\sref[I]{I.5.5.6}).
Quando $X$ é, além disso, quase-compacto, isto implica que os subconjuntos abertos retrocompactos de $X$ coincidem com os subconjuntos abertos quase-compactos de $X$, e que os subconjuntos construtíveis de $X$ são as uniões finitas de conjuntos da forma $U\cap\complement{V}$, onde $U$ e $V$ são abertos quase-compactos.
\end{env}

\begin{corollary}[9.1.7]
\label{0.9.1.7}
Para que um subconjunto de um espaço noetheriano $X$ seja construtível, é necessário e suficiente que seja uma união finita de subconjuntos localmente fechados de $X$.
\end{corollary}

\begin{proposition}[9.1.8]
\label{0.9.1.8}
Seja $X$ um espaço topológico e $U$ um subconjunto aberto de $X$.
\begin{enumerate}
  \item[{\rm(i)}] Se $T$ é um subconjunto construtível de $X$, então $T\cap U$ é um subconjunto construtível de $U$.
  \item[{\rm(ii)}] Suponhamos, além disso, que $U$ é retrocompacto em $X$. Para que um subconjunto $Z$ de $U$ seja construtível em $X$, é necessário e suficiente que seja construtível em $U$.
\end{enumerate}
\end{proposition}

\begin{proof}
\medskip\noindent
\begin{enumerate}
  \item[(i)] Mediante a Proposição \sref{0.9.1.3}, reduzimo-nos a demonstrar que, se $T$ é um aberto retrocompacto de $X$, então $T\cap U$ é um aberto retrocompacto de $U$; em outros termos, para todo aberto quase-compacto $W\subset U$, $T\cap U\cap W=T\cap W$ é quase-compacto, o que decorre imediatamente da hipótese.
  \item[(ii)] A condição é necessária por (i), de modo que resta demonstrar que é suficiente.
    Pela Proposição \sref{0.9.1.3}, basta considerar o caso em que $Z$ é um aberto retrocompacto \emph{em $U$}, pois então se seguirá que $U\setmin Z$ é construtível em $X$, e se $Z$ e $Z'$ são dois abertos retrocompactos \emph{em $U$}, então $Z\cap(U\setmin Z')$ será construtível em $X$.
    Se $W$ é um aberto quase-compacto de $X$ e $Z$ um aberto retrocompacto de $U$, temos $Z\cap W=Z\cap(W\cap U)$, e por hipótese $W\cap U$ é um aberto quase-compacto de $U$; portanto, $W\cap Z$ é quase-compacto e, consequentemente, $Z$ é um aberto retrocompacto de $X$, e \emph{a fortiori} construtível em $X$.
\end{enumerate}
\end{proof}

\begin{corollary}[9.1.9]
\label{0.9.1.9}
Seja $X$ um espaço topológico e $(U_i)_{i\in I}$ um recobrimento finito de $X$ por abertos retrocompactos de $X$.
Para que um subconjunto $Z$ de $X$ seja construtível em $X$, é necessário e suficiente que $Z\cap U_i$ seja construtível em $U_i$ para todo $i\in I$.
\end{corollary}

\begin{env}[9.1.10]
\label{0.9.1.10}
Suponhamos, em particular, que $X$ é quase-compacto e que todo ponto
\oldpage[0\textsubscript{III}]{14}
de $X$ admite um sistema fundamental de vizinhanças abertas retrocompactas em $X$ (e \emph{a fortiori} quase-compactas); então a condição para que um subconjunto $Z$ de $X$ seja construtível em $X$ é de natureza \emph{local}; em outros termos, é necessário e suficiente que, para todo $x\in X$, exista uma vizinhança aberta $V$ de $x$ tal que $V\cap Z$ seja construtível em $V$.
Com efeito, se esta condição é satisfeita, então para todo $x\in X$ existe uma vizinhança aberta $V$ de $x$ que é \emph{retrocompacta em $X$} e tal que $V\cap Z$ é construtível em $V$, pelas hipóteses sobre $X$ e pela Proposição \sref{0.9.1.8}[i]; basta então recobrir $X$ por um número finito dessas vizinhanças e aplicar o Corolário \sref{0.9.1.9}.
\end{env}

\begin{definition}[9.1.11]
\label{0.9.1.11}
Seja $X$ um espaço topológico.
Dizemos que um subconjunto $T$ de $X$ é \emph{localmente construtível} em $X$ se, para todo $x\in X$, existe uma vizinhança aberta $V$ de $x$ tal que $T\cap V$ é construtível em $V$.
\end{definition}

Da Proposição \sref{0.9.1.8}[i] segue-se que, se $V$ é tal que $V\cap T$ é construtível em $V$, então para todo aberto $W\subset V$, $W\cap T$ é construtível em $W$.
Se $T$ é localmente construtível em $X$, então, para todo aberto $U$ de $X$, $T\cap U$ é localmente construtível em $U$, em virtude da observação anterior.
A mesma observação mostra que o conjunto de subconjuntos localmente construtíveis de $X$ é estável por uniões finitas e interseções finitas; por outro lado, é claro que também é estável pela passagem ao complementar.

\begin{proposition}[9.1.12]
\label{0.9.1.12}
Seja $X$ um espaço topológico.
Todo conjunto construtível em $X$ é localmente construtível em $X$.
A recíproca é verdadeira se $X$ é quase-compacto e se sua topologia admite uma base formada pelos conjuntos retrocompactos de $X$.
\end{proposition}

\begin{proof}
A primeira afirmação decorre da Definição \sref{0.9.1.11}, e a segunda de \sref{0.9.1.10}.
\end{proof}

\begin{corollary}[9.1.13]
\label{0.9.1.13}
Seja $X$ um espaço topológico cuja topologia admite uma base formada pelos conjuntos retrocompactos de $X$.
Então todo subconjunto $T$ localmente construtível de $X$ é retrocompacto em $X$.
\end{corollary}

\begin{proof}
Seja $U$ um aberto quase-compacto de $X$; $T\cap U$ é localmente construtível em $U$, logo construtível em $U$ pela Proposição~\sref{0.9.1.12} e, consequentemente, quase-compacto pelo Corolário \sref{0.9.1.4}.
\end{proof}

\subsection{Subconjuntos construtíveis de espaços noetherianos}
\label{subsection:0.9.2}

\begin{env}[9.2.1]
\label{0.9.2.1}
Vimos em \sref{0.9.1.7} que, em um espaço noetheriano $X$, os subconjuntos construtíveis de $X$ são as \emph{uniões finitas de subconjuntos localmente fechados de $X$}.

A imagem inversa de um conjunto construtível de $X$ por uma aplicação contínua de um espaço noetheriano $X'$ em $X$ é construtível em $X'$.
Se $Y$ é um subconjunto construtível de um espaço noetheriano $X$, então os subconjuntos de $Y$ que são construtíveis como subespaços de $Y$ coincidem com os que são construtíveis como subespaços de $X$.
\end{env}

\begin{proposition}[9.2.2]
\label{0.9.2.2}
Seja $X$ um espaço noetheriano irredutível, e seja $E$ um subconjunto construtível de $X$.
Para que $E$ seja denso em toda parte em $X$, é necessário e suficiente que $E$ contenha um subconjunto aberto não vazio de $X$.
\end{proposition}

\begin{proof}
\oldpage[0\textsubscript{III}]{15}
A condição é evidentemente suficiente, pois todo aberto não vazio é denso em $X$.
Reciprocamente, seja $E=\bigcup_{i=1}^n(U_i\cap F_i)$ um subconjunto construtível de $X$, sendo os $U_i$ abertos não vazios e os $F_i$ fechados em $X$; temos então $\overline{E}\subset\bigcup_i F_i$.
Consequentemente, se $\overline{E}=X$, então $X$ é igual a um dos $F_i$, logo $E\supset U_i$, o que termina a demonstração.
\end{proof}

Quando $X$ admite um ponto genérico $x$ \sref[0]{0.2.1.2}, a condição da Proposição \sref{0.9.2.2} equivale à relação $x\in E$.

\begin{proposition}[9.2.3]
\label{0.9.2.3}
Seja $X$ um espaço noetheriano.
Para que um subconjunto $E$ de $X$ seja construtível, é necessário e suficiente que, para todo subconjunto fechado irredutível $Y$ de $X$, $E\cap Y$ seja raro em $Y$ ou contenha um subconjunto aberto não vazio de $Y$.
\end{proposition}

\begin{proof}
A necessidade da condição decorre de que $E\cap Y$ deve ser um subconjunto construtível de $Y$ e da Proposição \sref{0.9.2.2}, pois um subconjunto não denso de $Y$ é necessariamente raro no espaço irredutível $Y$ \sref[0]{0.2.1.1}.
Para demonstrar que a condição é suficiente, apliquemos o princípio de indução noetheriana \sref[0]{0.2.2.2} ao conjunto $\mathfrak{F}$ de subconjuntos fechados $Y$ de $X$ tais que $Y\cap E$ é construtível (com respeito a $Y$ ou com respeito a $X$, o que é equivalente): podemos supor assim que, para todo subconjunto fechado $Y\neq X$ de $X$, $E\cap Y$ é construtível.
Suponhamos primeiro que $X$ não é irredutível, e sejam $X_i$ ($1\leq i\leq m$) seus componentes irredutíveis, necessariamente em número finito \sref[0]{0.2.2.5}; por hipótese, os $E\cap X_i$ são construtíveis, logo também o é sua união $E$.
Suponhamos agora que $X$ é irredutível; então, por hipótese, se $E$ é raro, tem-se $\overline{E}\neq X$ e $E=E\cap\overline{E}$ é construtível; se $E$ contém um subconjunto aberto não vazio $U$ de $X$, é a união de $U$ e $E\cap(X\setmin U)$; mas $X\setmin U$ é um conjunto fechado distinto de $X$, pelo que $E\cap(X\setmin U)$ é construtível; consequentemente, o próprio $E$ é construtível, o que termina a demonstração.
\end{proof}

\begin{corollary}[9.2.4]
\label{0.9.2.4}
Seja $X$ um espaço noetheriano e $(E_\alpha)$ uma família filtrante crescente de subconjuntos construtíveis de $X$, tal que
\begin{enumerate}
  \item[\emph{(1.o)}] $X$ é a união da família $(E_\alpha)$.
  \item[\emph{(2.o)}] Todo subconjunto fechado irredutível de $X$ está contido no fecho de um dos $E_\alpha$.
\end{enumerate}

Então existe um índice $\alpha$ tal que $X=E_\alpha$.

Quando todo subconjunto fechado irredutível de $X$ admite um ponto genérico, pode-se omitir a hipótese \emph{(2.o)}.
\end{corollary}

\begin{proof}
Apliquemos o princípio de indução noetheriana \sref[0]{0.2.2.2} ao conjunto $\mathfrak{M}$ de subconjuntos fechados de $X$ contidos em ao menos um dos $E_\alpha$; podemos supor assim que todo subconjunto fechado $Y\neq X$ de $X$ está contido em um dos $E_\alpha$.
A proposição é evidente se $X$ não é irredutível, porque cada um dos componentes irredutíveis $X_i$ de $X$ ($1\leq i\leq m$) está contido em um $E_{\alpha_i}$, e existe um $E_\alpha$ que contém todos os $E_{\alpha_i}$.
Suponhamos agora que $X$ é irredutível.
Por hipótese, existe um $\beta$ tal que $X=\overline{E_\beta}$, de modo que \sref{0.9.2.2} $E_\beta$ contém um subconjunto aberto não vazio $U$ de $X$.
Mas então o conjunto fechado $X\setmin U$ está contido em um $E_\gamma$, e basta tomar um $E_\alpha$ que contenha $E_\beta$ e $E_\gamma$.
Quando todo subconjunto fechado irredutível $Y$ de $X$
\oldpage[0\textsubscript{III}]{16}
admite um ponto genérico $y$, existe $\alpha$ tal que $y\in E_\alpha$, de modo que $Y=\overline{\{y\}}\subset\overline{E_\alpha}$, e a condição (2.o) é, portanto, consequência de (1.o).
\end{proof}

\begin{proposition}[9.2.5]
\label{0.9.2.5}
Seja $X$ um espaço noetheriano, seja $x$ um ponto de $X$ e seja $E$ um subconjunto construtível de $X$.
Para que $E$ seja uma vizinhança de $x$, é necessário e suficiente que, para todo subconjunto fechado irredutível $Y$ de $X$ que contenha $x$, $E\cap Y$ seja denso em $Y$ (se existe um ponto genérico $y$ de $Y$, isto significa também \sref{0.9.2.2} que $y\in E$).
\end{proposition}

\begin{proof}
A condição é evidentemente necessária; demonstraremos que é suficiente.
Aplicando o princípio de indução noetheriana ao conjunto $\mathfrak{M}$ de subconjuntos fechados $Y$ de $X$ que contêm $x$ e tais que $E\cap Y$ é uma vizinhança de $x$ em $Y$, podemos supor que todo subconjunto fechado $Y\neq X$ de $X$ que contém $x$ pertence a $\mathfrak{M}$.
Se $X$ não é irredutível, cada um dos componentes irredutíveis $X_i$ de $X$ que contêm $x$ é distinto de $X$, logo $E\cap X_i$ é uma vizinhança de $x$ com respeito a $X_i$; consequentemente, $E$ é uma vizinhança de $x$ na união dos componentes irredutíveis de $X$ que contêm $x$, e como esta união é uma vizinhança de $x$ em $X$, também o é $E$.
Se $X$ é irredutível, $E$ é denso em $X$ por hipótese, pelo que contém um subconjunto aberto não vazio $U$ de $X$ \sref{0.9.2.2}; a proposição é então evidente se $x\in U$; em caso contrário, por hipótese, $x$ é um ponto interior de $E\cap(X\setmin U)$ relativo a $X\setmin U$, de modo que o fecho de $X\setmin E$ em $X$ não contém $x$, e o complementar desse fecho é uma vizinhança de $x$ contida em $E$, o que termina a demonstração.
\end{proof}

\begin{corollary}[9.2.6]
\label{0.9.2.6}
Seja $X$ um espaço noetheriano e seja $E$ um subconjunto de $X$.
Para que $E$ seja um aberto de $X$, é necessário e suficiente que, para todo subconjunto fechado irredutível $Y$ de $X$ que interseca $E$, $E\cap Y$ contenha um subconjunto aberto não vazio de $Y$.
\end{corollary}

\begin{proof}
A condição é evidentemente necessária; reciprocamente, se for satisfeita, implica que $E$ é construtível pela Proposição \sref{0.9.2.3}.
Além disso, a Proposição \sref{0.9.2.5} mostra que $E$ é então uma vizinhança de cada um de seus pontos, de onde se segue a conclusão.
\end{proof}

\subsection{Funções construtíveis}
\label{subsection:0.9.3}

\begin{definition}[9.3.1]
\label{0.9.3.1}
Seja $h$ uma aplicação de um espaço topológico $X$ em um conjunto $T$.
Dizemos que $h$ é \emph{construtível} se $h^{-1}(t)$ é construtível para todo $t\in T$, e vazio salvo para um número finito de valores de $t$; então, para todo subconjunto $S$ de $T$, $h^{-1}(S)$ é construtível.
Dizemos que $h$ é \emph{localmente construtível} se todo $x\in X$ possui uma vizinhança aberta $V$ tal que $h|V$ é construtível.
\end{definition}

Toda função construtível é localmente construtível; a recíproca é verdadeira quando $X$ é quase-compacto e admite uma base formada pelos abertos retrocompactos de $X$ (em particular, quando $X$ é noetheriano).

\begin{proposition}[9.3.2]
\label{0.9.3.2}
Seja $h$ uma aplicação de um espaço noetheriano $X$ em um conjunto $T$.
Para que $h$ seja construtível, é necessário e suficiente que, para todo subconjunto fechado irredutível $Y$ de $X$, exista um subconjunto não vazio $U$ de $Y$, aberto com respeito a $Y$, sobre o qual $h$ seja constante.
\end{proposition}

\begin{proof}
A condição é necessária: com efeito, por hipótese, $h$ somente toma um número finito de valores $t_i$ sobre $Y$, e cada um dos conjuntos $h^{-1}(t_i)\cap Y$ é construtível em $Y$ \sref{0.9.2.1}; como não podem ser todos subconjuntos raros do espaço $Y$, ao menos um deles contém um aberto não vazio \sref{0.9.2.3}.

\oldpage[0\textsubscript{III}]{17}
Para ver que a condição é suficiente, aplicamos o princípio de indução noetheriana ao conjunto $\mathfrak{M}$ de subconjuntos fechados $Y$ de $X$ tais que a restrição $h|Y$ é construtível; podemos supor assim que, para todo subconjunto fechado $Y\neq X$ de $X$,
$h|Y$ é construtível.
Se $X$ não é irredutível, a restrição de $h$ a cada um dos componentes irredutíveis (em número finito) $X_i$ de $X$ é construtível, e segue-se então imediatamente da Definição \sref{0.9.3.1} que $h$ é construtível.
Se $X$ é irredutível, existe por hipótese um subconjunto aberto não vazio $U$ de $X$ sobre o qual $h$ é constante; por outro lado, a restrição de $h$ a $X\setmin U$ é construtível por hipótese, e segue-se imediatamente que $h$ é construtível.
\end{proof}

\begin{corollary}[9.3.3]
\label{0.9.3.3}
Seja $X$ um espaço noetheriano no qual todo subconjunto fechado irredutível admite um ponto genérico.
Se $h$ é uma aplicação de $X$ em um conjunto $T$ tal que, para todo $t\in T$, $h^{-1}(t)$ é construtível, então $h$ é construtível.
\end{corollary}

\begin{proof}
Se $Y$ é um subconjunto fechado irredutível de $X$ e $y$ seu ponto genérico, então $Y\cap h^{-1}(h(y))$ é construtível e contém $y$; portanto, \sref{0.9.2.2} este conjunto contém um subconjunto aberto não vazio de $Y$, e basta aplicar a Proposição \sref{0.9.3.2}.
\end{proof}

\begin{proposition}[9.3.4]
\label{0.9.3.4}
Seja $X$ um espaço noetheriano no qual todo subconjunto fechado irredutível admite um ponto genérico, e seja $h$ uma aplicação construtível de $X$ em um conjunto ordenado.
Para que $h$ seja semicontínua superiormente sobre $X$, é necessário e suficiente que, para todo $x\in X$ e toda generalização \sref[0]{0.2.1.2} $x'$ de $x$, se tenha $h(x')\leq h(x)$.
\end{proposition}

\begin{proof}
A função $h$ somente toma um número finito de valores; portanto, dizer que é semicontínua superiormente significa que, para todo $x\in X$, o conjunto $E$ dos $y\in X$ tais que $h(y)\leq h(x)$ é uma vizinhança de $x$.
Por hipótese, $E$ é um subconjunto construtível de $X$; por outro lado, dizer que um subconjunto fechado irredutível $Y$ de $X$ contém $x$ significa que seu ponto genérico $y$ é uma generalização de $x$; a conclusão decorre então da Proposição \sref{0.9.2.5}.
\end{proof}

\section{Complemento sobre os módulos planos}
\label{section:0.10}

Para as demonstrações que faltam em \hyperref[subsection:0.10.1]{(10.1)} e \hyperref[subsection:0.10.2]{(10.2)}, remetemos o leitor a Bourbaki,~\emph{Alg.~comm.}, cap.~II e~III.

\subsection{Relações entre módulos planos e módulos livres}
\label{subsection:0.10.1}

\begin{env}[10.1.1]
\label{0.10.1.1}
Seja $A$ um anel, $\mathfrak{J}$ um ideal de $A$ e $M$ um $A$-módulo;
para todo inteiro $p\geq0$, temos um homomorfismo canônico de $(A/\mathfrak{J})$-módulos
\[
\label{0.10.1.1.1}
  \vphi_p:(M/\mathfrak{J}M)\otimes_{A/\mathfrak{J}}(\mathfrak{J}^p/\mathfrak{J}^{p+1})\to\mathfrak{J}^pM/\mathfrak{J}^{p+1}M,
  \tag{10.1.1.1}
\]
que é evidentemente \emph{sobrejetivo}.
Denotamos por $\gr(A)=\oplus_{p\geq0}\mathfrak{J}^p/\mathfrak{J}^{p+1}$ o anel graduado associado a $A$ filtrado pelos $\mathfrak{J}^p$, e por $\gr(M)=\oplus_{p\geq0}\mathfrak{J}^pM/\mathfrak{J}^{p+1}M$ o $\gr(A)$-módulo graduado associado a $M$ filtrado pelos $\mathfrak{J}^pM$;
temos então $\gr_p(A)=\mathfrak{J}^p/\mathfrak{J}^{p+1}$ e $\gr_p(M)=\mathfrak{J}^pM/\mathfrak{J}^{p+1}M$;
os $\vphi$ definem um homomorfismo \emph{sobrejetivo} de $\gr(A)$-módulos graduados
\[
\label{0.10.1.1.2}
  \vphi:\gr_0(M)\otimes_{\gr_0(A)}\gr(A)\to\gr(M).
  \tag{10.1.1.2}
\]
\end{env}

\oldpage[0\textsubscript{III-1}]{18}
\begin{env}[10.1.2]
\label{0.10.1.2}
Suponhamos que se verifica \emph{uma} das hipóteses seguintes:
\begin{enumerate}
  \item[(i)] $\mathfrak{J}$ é nilpotente;
  \item[(ii)] $A$ é noetheriano, $\mathfrak{J}$ está contido no radical de $A$ e $M$ é de tipo finito.
\end{enumerate}
Então as propriedades seguintes são equivalentes.
\begin{enumerate}
  \item[(a)] $M$ é um $A$-módulo livre.
  \item[(b)] $M/\mathfrak{J}M=M\otimes_A(A/\mathfrak{J})$ é um $(A/\mathfrak{J})$-módulo livre e $\Tor_1^A(M,A/\mathfrak{J})=0$.
  \item[(c)] $M/\mathfrak{J}M$ é um $(A/\mathfrak{J})$-módulo livre e o homomorfismo canônico \sref{0.10.1.1.2} é injetivo (e, portanto, bijetivo).
\end{enumerate}
\end{env}

\begin{env}[10.1.3]
\label{0.10.1.3}
Suponhamos que $A/\mathfrak{J}$ é um \emph{corpo} (em outros termos, que $\mathfrak{J}$ é maximal) e que se verifica uma das hipóteses (i) e (ii) de \sref{0.10.1.2}.
Então as propriedades seguintes são equivalentes.
\begin{enumerate}
  \item[(a)] $M$ é um $A$-módulo livre.
  \item[(b)] $M$ é um $A$-módulo projetivo.
  \item[(c)] $M$ é um $A$-módulo plano.
  \item[(d)] $\Tor_1^A(M,A/\mathfrak{J})=0$.
  \item[(e)] O homomorfismo canônico \sref{0.10.1.1.2} é bijetivo.
\end{enumerate}

Este resultado pode ser aplicado, em particular, aos dois casos seguintes:
\begin{enumerate}
  \item[(i)] $M$ é um módulo \emph{arbitrário} sobre um anel local $A$ cujo ideal maximal $\mathfrak{J}$ é \emph{nilpotente} (por exemplo, um anel artiniano local);
  \item[(ii)] $M$ é um módulo \emph{de tipo finito} sobre um anel \emph{local noetheriano}.
\end{enumerate}
\end{env}

\subsection{Critérios locais de planicidade}
\label{subsection:0.10.2}

\begin{env}[10.2.1]
\label{0.10.2.1}
Com as hipóteses e a notação de \sref{0.10.1.1}, consideremos as condições seguintes.
\begin{enumerate}
  \item[(a)] $M$ é um $A$-módulo plano.
  \item[(b)] $M/\mathfrak{J}M$ é um $(A/\mathfrak{J})$-módulo plano e $\Tor_1^A(M,A/\mathfrak{J})=0$.
  \item[(c)] $M/\mathfrak{J}M$ é um $(A/\mathfrak{J})$-módulo plano e o homomorfismo canônico \sref{0.10.1.1.2} é bijetivo.
  \item[(d)] Para todo $n\geq1$, $M/\mathfrak{J}^nM$ é um $(A/\mathfrak{J}^n)$-módulo plano.
\end{enumerate}

Então temos as implicações
\[
  \text{(a)}\implies\text{(b)}\implies\text{(c)}\implies\text{(d)},
\]
e, se $\mathfrak{J}$ é \emph{nilpotente}, as quatro condições são \emph{equivalentes}.
Ocorre também assim se $A$ é noetheriano e $M$ é \emph{idealmente separado}, isto é, se para todo ideal $\mathfrak{a}$ de $A$, o $A$-módulo $\mathfrak{a}\otimes_A M$ está \emph{separado} para a topologia $\mathfrak{J}$-pré-ádica.
\end{env}

\begin{env}[10.2.2]
\label{0.10.2.2}
Seja $A$ um anel noetheriano, $B$ uma $A$-álgebra comutativa noetheriana, $\mathfrak{J}$ um ideal de $A$ tal que $\mathfrak{J}B$ está contido no radical de $B$, e $M$ um $B$-\emph{módulo} de tipo finito.
Então, quando $M$ é considerado como $A$-\emph{módulo}, as quatro condições de \sref{0.10.2.1} são \emph{equivalentes}.
\oldpage[0\textsubscript{III-1}]{19}
Este resultado se aplica, em primeiro lugar, ao caso em que $A$ e $B$ são anéis noetherianos \emph{locais} e o homomorfismo $A\to B$ é um homomorfismo \emph{local}.
Mais precisamente, se $\mathfrak{J}$ é então o ideal \emph{maximal} de $A$, podemos suprimir, nas condições~\emph{(b)} e \emph{(c)}, a hipótese de que $M/\mathfrak{J}M$ seja plano, pois ela se verifica automaticamente, e a condição~\emph{(d)} implica que os módulos $M/\mathfrak{J}^nM$ são \emph{livres} sobre os $A/\mathfrak{J}^n$.
\end{env}

\begin{env}[10.2.3]
\label{0.10.2.3}
Com as hipóteses sobre $A$, $B$, $\mathfrak{J}$ e $M$ do início de \sref{0.10.2.2}, seja $\widehat{A}$ o completamento separado de $A$ para a topologia $\mathfrak{J}$-pré-ádica, e $\widehat{M}$ o completamento separado de $M$ para a topologia $\mathfrak{J}B$-pré-ádica.
Para que $M$ seja um $A$-módulo plano, é necessário e suficiente que $\widehat{M}$ seja um $\widehat{A}$-módulo plano.
\end{env}

\begin{env}[10.2.4]
\label{0.10.2.4}
Seja $\rho:A\to B$ um homomorfismo local de anéis noetherianos locais, seja $k$ o corpo residual de $A$, e sejam $M$ e $N$ dois $B$-módulos de tipo finito, supondo $N$ \emph{$A$-plano}.
Seja $u:M\to N$ um $B$-homomorfismo.
Então as condições seguintes são equivalentes.
\begin{enumerate}
  \item[(a)] $u$ é injetivo e $\Coker(u)$ é um $A$-módulo plano.
  \item[(b)] $u\otimes1:M\otimes_A k\to N\otimes_A k$ é injetivo.
\end{enumerate}
\end{env}

\begin{env}[10.2.5]
\label{0.10.2.5}
Sejam $\rho:A\to B$ e $\sigma:B\to C$ homomorfismos locais de anéis noetherianos locais, seja $k$ o corpo residual de $A$, e seja $M$ um $C$-módulo de tipo finito.
Suponhamos que $B$ é um $A$-módulo \emph{plano}.
Então as condições seguintes são equivalentes.
\begin{enumerate}
  \item[(a)] $M$ é um $B$-módulo plano.
  \item[(b)] $M$ é um $A$-módulo plano e $M\otimes_A k$ é um $(B\otimes_A k)$-módulo plano.
\end{enumerate}
\end{env}

\begin{proposition}[10.2.6]
\label{0.10.2.6}
Sejam $A$ e $B$ anéis noetherianos locais, seja $\rho:A\to B$ um homomorfismo local, seja $\mathfrak{J}$ um ideal de $B$ contido no ideal maximal, e seja $M$ um $B$-módulo de tipo finito.
Suponhamos que, para todo $n\geq0$, $M_n=M/\mathfrak{J}^{n+1}M$ é um $A$-módulo plano.
Então $M$ é um $A$-módulo plano.
\end{proposition}

\begin{proof}
Devemos demonstrar que, para todo homomorfismo injetivo $u:N'\to N$ de $A$-módulos de tipo finito, $v=1\otimes u:M\otimes_A N'\to M\otimes_A N$ é injetivo.
Mas $M\otimes_A N'$ e $M\otimes_A N$ são $B$-módulos de tipo finito e, portanto, estão separados para a topologia $\mathfrak{J}$-pré-ádica \sref[0\textsubscript{I}]{0.7.3.5};
basta, pois, demonstrar que o homomorfismo $\widehat{v}:\widehat{M\otimes_A N'}\to\widehat{M\otimes_A N}$ dos completamentos separados é injetivo.
Mas $\widehat{v}=\varprojlim v_n$, onde $v_n$ é o homomorfismo $1\otimes u:M_n\otimes_A N'\to M_n\otimes_A N$;
como, por hipótese, $M_n$ é $A$-plano, $v_n$ é injetivo para todo $n$, e também o é $\widehat{v}$, pois o funtor $\varprojlim$ é exato à esquerda.
\end{proof}

\begin{corollary}[10.2.7]
\label{0.10.2.7}
Seja $A$ um anel noetheriano, seja $B$ um anel noetheriano local, seja $\rho:A\to B$ um homomorfismo, seja $f$ um elemento do ideal maximal de $B$, e seja $M$ um $B$-módulo de tipo finito.
Suponhamos que a homotecia $f_M:x\to fx$ sobre $M$ é injetiva e que $M/fM$ é um $A$-módulo plano.
Então $M$ é um $A$-módulo plano.
\end{corollary}

\begin{proof}
Seja $M_i=f^iM$ para $i\geq0$;
como $f_M$ é injetiva, $M_i/M_{i+1}$ é isomorfo a $M/fM$ e, portanto, $A$-plano para todo $i\geq0$;
a sucessão exata
\[
  0\to M_i/M_{i+1}\to M/M_{i+1}\to M/M_i\to 0
\]
nos dá, por indução sobre $i$, que $M/M_i$ é \emph{$A$-plano} para todo $i\geq0$ \sref[0\textsubscript{I}]{0.6.1.2};
podemos, pois, aplicar \sref{0.10.2.6}.
Também podemos raciocinar diretamente como segue:
para todo $A$-módulo $N$ de tipo finito, $M\otimes_A N$ é um $B$-módulo de tipo finito;
como $f$ pertence ao radical $\mathfrak{n}$ de $B$, a topologia $(f)$-ádica sobre $M\otimes_A N$ é mais fina que a topologia
\oldpage[0\textsubscript{III-1}]{20}
$\mathfrak{n}$-ádica, e sabemos que esta última está \emph{separada} \sref[0\textsubscript{I}]{0.7.3.5}.
Ora, como $M/M_i$ é $A$-plano, temos
\[
  f^i(M\otimes_A N)=\Im(M_i\otimes_A N\to M\otimes_A N)=\Ker(M\otimes_A N\to(M/M_i)\otimes_A N)
\]
por \sref[0\textsubscript{I}]{0.6.1.2}.
Seja, pois, $N$ um $A$-módulo de tipo finito, e seja $N'$ um submódulo de $N$, com injeção canônica $j:N'\to N$; no diagrama comutativo
\[
  \xymatrix{
    M\otimes_A N'\ar[r]\ar[d]_{1_M\otimes j} &
    (M/M_i)\otimes_A N'\ar[d]^{1_{M/M_i}\otimes j}\\
    M\otimes_A N\ar[r] &
    (M/M_i)\otimes_A N
  }
\]
$1_{M/M_i}\otimes j$ é injetivo, porque $M/M_i$ é $A$-plano; concluímos, pois, que
\[
  \Ker(M\otimes_A N'\to M\otimes_A N)\subset\Ker(M\otimes_A N'\to(M/M_i)\otimes_A N')
\]
para todo $i$; como a interseção destes últimos núcleos é $0$, como vimos anteriormente, também o é o primeiro núcleo e, portanto, $M$ é $A$-plano.
\end{proof}

\begin{proposition}[10.2.8]
\label{0.10.2.8}
Seja $A$ um anel noetheriano reduzido e seja $M$ um $A$-módulo de tipo finito.
Suponhamos que, para toda $A$-álgebra $B$ que seja um anel de valoração discreta, $M\otimes_A B$ é um $B$-módulo plano (e, portanto, livre \sref{0.10.1.3}).
Então $M$ é um $A$-módulo plano.
\end{proposition}

\begin{proof}
Sabemos que, para que $M$ seja plano, é necessário e suficiente que $M_\mathfrak{m}$ seja um $A_\mathfrak{m}$-módulo plano para todo ideal maximal $\mathfrak{m}$ de $A$ \sref[0\textsubscript{I}]{0.6.3.3};
podemos, pois, restringir-nos ao caso em que $A$ é \emph{local} \sref[0\textsubscript{I}]{0.1.2.8}.
Seja $\mathfrak{m}$ o ideal maximal de $A$, sejam $\mathfrak{p}_i$ ($1\leq i\leq r$) os ideais primos mínimos de $A$, e seja $k$ o corpo residual $A/\mathfrak{m}$.
Sabemos \sref[II]{II.7.1.7} que existe, para cada $i$, um anel de valoração discreta $B_i$ que tem o mesmo corpo de frações $K_i$ que o domínio integral $A/\mathfrak{p}_i$ e que, além disso, domina $A/\mathfrak{p}_i$.
Seja $M_i=M\otimes_A B_i$.
Por hipótese, $M_i$ é livre sobre $B_i$ e, denotando por $k_i$ o corpo residual de $B_i$, temos
\[
\label{0.10.2.8.1}
  \rg_{k_i}(M_i\otimes_{B_i}k_i)=\rg_{K_i}(M_i\otimes_{B_i}K_i).
  \tag{10.2.8.1}
\]

Mas está claro que o homomorfismo composto $A\to A/\mathfrak{p}_i\to B_i$ é local, de modo que $k_i$ é uma extensão de $k$, e que temos $M_i\otimes_{B_i}k_i = M\otimes_A k_i = (M\otimes_A k)\otimes_k k_i$, assim como $M_i\otimes_{B_i}K_i = M\otimes_A K_i$.
A equação~\sref{0.10.2.8.1} implica, portanto, que
\[
  \rg_k(M\otimes_A k)=\rg_{K_i}(M\otimes_A K_i)\quad\mbox{para $1\leq i\leq r$}
\]
e, como $A$ é reduzido, sabemos que esta condição implica que $M$ é um $A$-módulo \emph{livre} (Bourbaki, \emph{Alg. comm.}, cap.~II, \textsection~3, n\textsuperscript{o}~2, prop.~7).
\end{proof}

\subsection{Existência de extensões planas de anéis locais}
\label{subsection:0.10.3}

\begin{proposition}[10.3.1]
\label{0.10.3.1}
Seja $A$ um anel noetheriano local, com ideal maximal $\mathfrak{J}$ e corpo residual $k=A/\mathfrak{J}$.
Seja $K$ uma extensão do corpo $k$.
Então existe um homomorfismo local de $A$ em um anel noetheriano local $B$, tal que $B/\mathfrak{J}B$ é $k$-isomorfo a $K$ e tal que $B$ é um $A$-módulo plano.
\end{proposition}

O restante desta seção é dedicado a demonstrar esta proposição passo a passo.
\begin{env}[10.3.1.1]
\label{0.10.3.1.1}
Suponhamos primeiro que $K=k(T)$, onde $T$ é uma indeterminada.
No anel de polinômios $A'=A[T]$, consideremos o ideal primo $\mathfrak{J}'=\mathfrak{J}A'$, formado pelos
\oldpage[0\textsubscript{III-1}]{21}
polinômios cujos coeficientes pertencem ao ideal $\mathfrak{J}$;
está claro que $A'/\mathfrak{J}'$ é canonicamente isomorfo a $k[T]$.
Mostraremos que o anel de frações $B=A'_{\mathfrak{J}'}$ é o que buscamos (isto é, um anel que satisfaz as condições da conclusão da proposição);
é claramente um anel noetheriano local, com ideal maximal $\mathfrak{L}=\mathfrak{J}B$.
Além disso, $B/\mathfrak{L}=(A'/\mathfrak{J}')_{\mathfrak{J}'}=(k[T])_{\mathfrak{J}'}$ é precisamente o corpo de frações $K$ de $k[T]$.
Finalmente, $B$ é um $A'$-módulo plano e $A'$ é um $A$-módulo livre, logo $B$ é um $A$-módulo plano \sref[0\textsubscript{I}]{0.6.2.1}.
\end{env}

\begin{env}[10.3.1.2]
\label{0.10.3.1.2}
Suponhamos agora que $K=k(t)=k[t]$, onde $t$ é algébrico sobre $k$;
seja $f\in k[T]$ o polinômio mínimo de $t$;
existe um polinômio mônico $F\in A[T]$ cuja imagem canônica em $k[T]$ é $f$.
Seja, pois, $A'=A[T]$, e seja $\mathfrak{J}'$ o ideal $\mathfrak{J}A'+(F)$ de $A'$.
Veremos que o anel quociente $B=A'/(F)$ é o que buscamos.
Em primeiro lugar, está claro que $B$ é um $A$-módulo \emph{livre} e, portanto, plano.
O anel $A'/\mathfrak{J}'$ é isomorfo a
\[
  (A'/\mathfrak{J}A')/\big((\mathfrak{J}A'+(F))/\mathfrak{J}A'\big)=k[T]/(f)=K;
\]
a imagem $\mathfrak{L}$ de $\mathfrak{J}'$ em $B$ é, portanto, maximal, e temos evidentemente $\mathfrak{L}=\mathfrak{J}B$.
Finalmente, $B$ é um anel semilocal, porque é um $A$-módulo de tipo finito (Bourbaki, \emph{Alg. comm.}, cap.~IV, \textsection~2, n\textsuperscript{o}~5, cor.~3 da prop.~9), e seus ideais maximais estão em correspondência bijetiva com os de $B/\mathfrak{J}B$ (\cite[vol.~I, p.~259]{I-13});
os argumentos anteriores demonstram então que $B$ é um anel local.
\end{env}

\begin{lemma}[10.3.1.3]
\label{0.10.3.1.3}
Seja $(A_\lambda,f_{\mu\lambda})$ um sistema indutivo filtrante de anéis locais, tal que os $f_{\mu\lambda}$ são homomorfismos locais;
seja $\mathfrak{m}_\lambda$ o ideal maximal de $A_\lambda$, e seja $K_\lambda=A_\lambda/\mathfrak{m}_\lambda$.
Então $A'=\varinjlim A_\lambda$ é um anel local, com ideal maximal $\mathfrak{m}'=\varinjlim\mathfrak{m}_\lambda$ e corpo residual $K=\varinjlim K_\lambda$.
Além disso, se $\mathfrak{m}_\mu=\mathfrak{m}_\lambda A_\mu$ para $\lambda<\mu$, temos $\mathfrak{m}'=\mathfrak{m}_\lambda A'$ para todo $\lambda$.
Se, além disso, para $\lambda<\mu$, $A_\mu$ é um $A_\lambda$-módulo plano e todos os $A_\lambda$ são noetherianos, então $A'$ é noetheriano e é um $A_\lambda$-módulo plano para todo $\lambda$.
\end{lemma}

\begin{proof}
Como, por hipótese, $f_{\mu\lambda}(\mathfrak{m}_\lambda)\subset\mathfrak{m}_\mu$ para $\lambda<\mu$, os $\mathfrak{m}_\lambda$ formam um sistema indutivo e seu limite $\mathfrak{m}'$ é evidentemente um ideal de $A'$.
Além disso, se $x'\not\in\mathfrak{m}'$, existe um $\lambda$ tal que $x'=f_\lambda(x_\lambda)$ para algum $x_\lambda\in A_\lambda$ (onde $f_\lambda:A_\lambda\to A'$ denota o homomorfismo canônico);
como $x'\not\in\mathfrak{m}'$, necessariamente $x_\lambda\not\in\mathfrak{m}_\lambda$, de modo que $x_\lambda$ admite um inverso $y_\lambda$ em $A_\lambda$, e $y'=f_\lambda(y_\lambda)$ é o inverso de $x'$ em $A'$, o que demonstra que $A'$ é um anel local com ideal maximal $\mathfrak{m}'$;
a afirmação sobre $K$ decorre imediatamente do fato de que $\varinjlim$ é um funtor exato.
A hipótese $\mathfrak{m}_\mu=\mathfrak{m}_\lambda A_\mu$ implica que a aplicação canônica $\mathfrak{m}_\lambda\otimes_{A_\lambda}A_\mu\to\mathfrak{m}_\mu$ é sobrejetiva;
a igualdade $\mathfrak{m}'=\mathfrak{m}_\lambda A'$ decorre então, novamente, do fato de que o funtor $\varinjlim$ é exato e comuta com o produto tensorial.

Suponhamos agora que, para $\lambda<\mu$, temos $\mathfrak{m}_\mu=\mathfrak{m}_\lambda A_\mu$ e que $A_\mu$ é um $A_\lambda$-módulo plano.
Então $A'$ é um $A_\lambda$-módulo plano para todo $\lambda$, por \sref[0\textsubscript{I}]{0.6.2.3};
como $A'$ e $A_\lambda$ são anéis locais e $\mathfrak{m}'=\mathfrak{m}_\lambda A'$, $A'$ é até mesmo um módulo \emph{fielmente plano} sobre $A_\lambda$ \sref[0\textsubscript{I}]{0.6.6.2}.
Suponhamos finalmente, além disso, que os $A_\lambda$ são \emph{noetherianos};
as topologias $\mathfrak{m}_\lambda$-ádicas estão então separadas \sref[0\textsubscript{I}]{0.7.3.5};
mostraremos agora que disso decorre que a topologia $\mathfrak{m}'$-ádica sobre $A'$ está \emph{separada}.
Com efeito, se $x'\in A'$ pertence a todos os $\mathfrak{m}^{'n}$ ($n>0$), então é a imagem de algum $x_\mu\in A_\mu$ para certo índice $\mu$, e como a imagem inversa em $A_\mu$ de $\mathfrak{m}^{'n}=\mathfrak{m}_\mu^n A'$ é $\mathfrak{m}_\mu^n$ \sref[0\textsubscript{I}]{0.6.6.1}, $x_\mu$ pertence a todos os $\mathfrak{m}_\mu^n$, logo $x_\mu=0$ por hipótese e, portanto, $x'=0$.
Denotemos por $\widehat{A}'$ o completamento de $A'$ para a topologia $\mathfrak{m}'$-ádica;
o anterior mostra que temos $A'\subset\widehat{A}'$.
Mostraremos agora que $\widehat{A}'$ é \emph{noetheriano} e \emph{$A_\lambda$-plano} para todo $\lambda$;
disso,
\oldpage[0\textsubscript{III-1}]{22}
seguirá que $\widehat{A}'$ é $A'$-plano \sref[0\textsubscript{I}]{0.6.2.3} e, como $\mathfrak{m}'\widehat{A}'\neq\widehat{A}'$, que $\widehat{A}'$ é um $A'$-módulo fielmente plano \sref[0\textsubscript{I}]{0.6.6.2}; daí se obtém finalmente que $A'$ é \emph{noetheriano} \sref[0\textsubscript{I}]{0.6.5.2}, o que terminará a demonstração do lema.

Temos $\widehat{A}'=\varprojlim_n A'/\mathfrak{m}^{\prime n}$;
pelo fato de que $A'$ é $A_\lambda$-plano, temos
\[
  \mathfrak{m}^{\prime n}/\mathfrak{m}^{\prime n+1}=(\mathfrak{m}_\lambda^{n}/\mathfrak{m}_\lambda^{n+1})\otimes_{A_\lambda}A'=(\mathfrak{m}_\lambda^{n}/\mathfrak{m}_\lambda^{n+1})\otimes_{K_\lambda}(K_\lambda\otimes_{A_\lambda}A')=(\mathfrak{m}_\lambda^{n}/\mathfrak{m}_\lambda^{n+1})\otimes_{K_\lambda}K;
\]
como $\mathfrak{m}_\lambda^{n}/\mathfrak{m}_\lambda^{n+1}$ é um $K_\lambda$-espaço vetorial de dimensão finita, $\mathfrak{m}^{\prime n}/\mathfrak{m}^{\prime n+1}$ é um $K$-espaço vetorial de dimensão finita para todo $n\geq 0$.
Segue-se então de \sref[0\textsubscript{I}]{0.7.2.12} e \sref[0\textsubscript{I}]{0.7.2.8} que $\widehat{A}'$ é \emph{noetheriano}.
Sabemos além disso que o ideal maximal de $\widehat{A}'$ é $\mathfrak{m}'\widehat{A}'$, e que $\widehat{A}'/\mathfrak{m}^{\prime n}\widehat{A}'$ é isomorfo a $A'/\mathfrak{m}^{\prime n}$;
como $A'/\mathfrak{m}^{\prime n}=(A_\lambda/\mathfrak{m}_\lambda^n)\otimes_{A_\lambda}A'$, vemos que $A'/\mathfrak{m}^{\prime n}$ é um $(A_\lambda/\mathfrak{m}_\lambda^n)$-módulo plano \sref[0\textsubscript{I}]{0.6.2.1};
o critério~\sref{0.10.2.2} é, portanto, aplicável à $A_\lambda$-álgebra noetheriana $\widehat{A}'$, e mostra que $\widehat{A}'$ é \emph{$A_\lambda$-plano}.
\end{proof}

\begin{env}[10.3.1.4]
\label{0.10.3.1.4}
Tratemos agora o caso geral.
Existe um ordinal $\gamma$ e, para todo ordinal $\lambda\leq\gamma$, um subcorpo $k_\lambda$ de $K$ que contém $k$, tais que (i) para todo $\lambda<\gamma$, $k_{\lambda+1}$ é uma extensão de $k_\lambda$ gerada por um único elemento; (ii) para todo ordinal limite $\mu$, $k_\mu=\bigcup_{\lambda<\mu}k_\lambda$; e (iii) $K=k_\gamma$.
Com efeito, basta considerar uma bijeção $\xi\mapsto t_\xi$ do conjunto de ordinais $\xi\leq\beta$ (para algum $\beta$ apropriado) em $K$, e definir $k_\lambda$ por indução transfinita (para $\lambda\leq\beta$) como a união dos $k_\mu$ para $\mu<\lambda$ se $\lambda$ é um ordinal limite, e como $k_\nu(t_\xi)$ se $\lambda=\nu+1$, onde $\xi$ é o menor ordinal tal que $t_\xi\not\in k_\nu$;
$\gamma$ é então, por definição, o menor ordinal $\leq\beta$ tal que $k_\gamma=K$.

Tendo isso em mente, definiremos, por indução transfinita, uma família de anéis noetherianos locais $A_\lambda$ para $\lambda\leq\gamma$, e homomorfismos locais $f_{\mu\lambda}:A_\lambda\to A_\mu$ para $\lambda\leq\mu$, que satisfaçam as condições seguintes:
\begin{enumerate}
  \item[(i)] $(A_\lambda,f_{\mu\lambda})$ é um sistema indutivo e $A_0=A$;
  \item[(ii)] para todo $\lambda$, temos um $k$-isomorfismo $A_\lambda/\mathfrak{J}A_\lambda\isoto k_\lambda$;
  \item[(iii)] para $\lambda\leq\mu$, $A_\mu$ é um $A_\lambda$-módulo plano.
\end{enumerate}

Suponhamos, pois, definidos os $A_\lambda$ e os $f_{\mu\lambda}$ para $\lambda<\mu<\xi$, e suponhamos primeiro que $\xi=\zeta+1$, de modo que $k_\xi=k_\zeta(t)$.
Se $t$ é transcendente sobre $k_\zeta$, definimos $A_\xi$, seguindo o procedimento de \sref{0.10.3.1.1}, como $(A_\zeta[t])_{\mathfrak{J}A_\zeta[t]}$;
a aplicação canônica é $f_{\xi\zeta}$ e, para $\lambda<\zeta$, tomamos $f_{\xi\lambda}=f_{\xi\zeta}\circ f_{\zeta\lambda}$;
a verificação das condições~(i) a (iii) é então imediata, em virtude do que foi demonstrado em \sref{0.10.3.1.1}.
Suponhamos agora que $t$ é algébrico, e seja $h$ seu polinômio mínimo em $k_\zeta[T]$, e seja $H$ um polinômio mônico de $A_\zeta[T]$ cuja imagem em $k_\zeta[T]$ é $h$;
tomamos então $A_\xi$ igual a $A_\zeta[T]/(H)$, definindo os $f_{\xi\lambda}$ como antes;
a verificação das condições~(i) a (iii) decorre então do que foi demonstrado em \sref{0.10.3.1.2}.

Suponhamos agora que $\xi$ não tem predecessor;
tomamos então $A_\xi$ como limite indutivo do sistema indutivo de anéis locais $(A_\lambda,f_{\mu\lambda})$ para $\lambda<\xi$;
definimos $f_{\xi\lambda}$ como a aplicação canônica para $\lambda<\xi$.
O fato de que $A_\xi$ seja local e noetheriano, de que os $f_{\xi\lambda}$ sejam homomorfismos locais e de que as condições~(i) a (iii) sejam satisfeitas para $\lambda\leq\xi$
\oldpage[0\textsubscript{III-1}]{23}
decorre então da hipótese de indução e do Lema~\sref{0.10.3.1.3}.
Com esta construção, está claro que o anel $B=A_\gamma$ satisfaz as condições de \sref{0.10.3.1}.\qed
\end{env}

Observemos que, por \sref{0.10.2.1}[c], temos um isomorfismo canônico
\[
\label{0.10.3.1.5}
  \gr(A)\otimes_k K\simto\gr(B).
  \tag{10.3.1.5}
\]

Também podemos substituir $B$ por seu completamento $\mathfrak{J}B$-ádico $\widehat{B}$ sem modificar as conclusões de \sref{0.10.3.1}, porque $\widehat{B}$ é um $B$-módulo plano \sref[0\textsubscript{I}]{0.7.3.3} e, portanto, um $A$-módulo plano \sref[0\textsubscript{I}]{0.6.2.1}.

Demonstramos também o seguinte:
\begin{corollary}[10.3.2]
\label{0.10.3.2}
Se $K$ é uma extensão de grau finito, podemos supor que $B$ é uma $A$-álgebra finita.
\end{corollary}

\section{Suplemento sobre álgebra homológica}
\label{section:0.11}

\subsection{Revisão das sequências espectrais}
\label{subsection:0.11.1}

\begin{env}[11.1.1]
\label{0.11.1.1}
No que segue usamos uma noção de sequência espectral mais geral do que a definida em (T, 2.4); conservando a notação de (T, 2.4), chamamos \emph{sequência espectral} em uma categoria abeliana $\cat{C}$ a um sistema $E$ constituído pelas seguintes partes:
\begin{enumerate}
  \item[(a) ] Uma família $(E_r^{pq})$ de objetos de $\cat{C}$ definida para $p,q\in\bb{Z}$ e $r\geq 2$.
  \item[(b) ] Uma família de morfismos $d_r^{pq}:E_r^{pq}\to E_r^{p+r,q-r+1}$ tais que $d_r^{p+r,q-r+1}d_r^{pq}=0$.
    Colocamos $Z_{r+1}(E_r^{pq})=\Ker(d_r^{pq})$ e $B_{r+1}(E_r^{pq})=\Im(d_r^{p-r,q+r-1})$, para que
    \[
      B_{r+1}(E_r^{pq})\subset Z_{r+1}(E_r^{pq})\subset E_r^{pq}.
    \]
  \item[(c) ] Uma família de isomorfismos $\alpha_r^{pq}:Z_{r+1}(E_r^{pq})/B_{r+1}(E_r^{pq})\isoto E_{r+1}^{pq}$.

  Definimos então, para $k\geq r+1$, por indução sobre $k$, os subobjetos $B_k(E_r^{pq})$ e $Z_k(E_r^{pq})$ como as imagens inversas, pelo morfismo canônico $E_r^{pq}\to E_r^{pq}/B_{r+1}(E_r^{pq})$, dos subobjetos desse quociente identificados mediante $\alpha_r^{pq}$ com os subobjetos $B_k(E_{r+1}^{pq})$ e $Z_k(E_{r+1}^{pq})$, respectivamente.
É claro que temos então, exceto isomorfismo,
  \[
  \label{0.11.1.1.1}
    Z_k(E_r^{pq})/B_k(E_r^{pq})=E_k^{pq}\text{ para }k\geq r+1,
    \tag{11.1.1.1}
  \]
  e, se colocarmos $B_r(E_r^{pq})=0$ e $Z_r(E_r^{pq})=E_r^{pq}$, temos as relações de inclusão
  \[
  \label{0.11.1.1.2}
    0=B_r(E_r^{pq})\subset B_{r+1}(E_r^{pq})\subset B_{r+2}(E_r^{pq})\subset\cdots\subset Z_{r+2}(E_r^{pq})\subset Z_{r+1}(E_r^{pq})\subset Z_r(E_r^{pq})=E_r^{pq}.
    \tag{11.1.1.2}
  \]
  As outras partes dos dados de $E$ são então:
  \item[(d) ] Dois subobjetos $B_\infty(E_2^{pq})$ e $Z_\infty(E_2^{pq})$ de $E_2^{pq}$ tais que $B_\infty(E_2^{pq})\subset Z_\infty(E_2^{pq})$ e, para todo $k\geq 2$,
    \[
      B_k(E_2^{pq})\subset B_\infty(E_2^{pq})\text{ e }Z_\infty(E_2^{pq})\subset Z_k(E_2^{pq}).
    \]
    Definimos
    \[
    \label{0.11.1.1.3}
      E_\infty^{pq}=Z_\infty(E_2^{pq})/B_\infty(E_2^{pq}).
      \tag{11.1.1.1.3}
    \]
  \item[(e)]
\oldpage[0\textsubscript{III}]{24}
    Uma família $(E^n)$ de objetos de $\cat{C}$, cada um provido de uma \emph{filtração decrescente $(F^p(E^n))_{p\in\bb{Z}}$}.
    Como de costume, denotamos por $\gr(E^n)$ o objeto graduado associado ao objeto filtrado $E^n$, soma direta dos $\gr_p(E^n)=F^p(E^n)/F^{p+1}(E^n)$.
  \item[(f) ] Para todo par $(p,q)\in\bb{Z}\times\bb{Z}$, um isomorfismo $\beta^{pq}:E_\infty^{pq}\isoto\gr_p(E^{p+q})$.
\end{enumerate}

A família $(E^n)$, sem as filtrações, chama-se o \emph{termo limite} (ou \emph{limite}) da sequência espectral $E$.

Suponhamos que a categoria $\cat{C}$ admita somas diretas infinitas, ou que, para todo $r\geq 2$ e todo $n\in\bb{Z}$, exista apenas um número finito de pares $(p,q)$ tais que $p+q=n$ e $E_r^{pq}\neq 0$ (basta que isso se verifique para $r=2$).
Então são definidos os $E_r^{(n)}=\sum_{p+q=n}E_r^{pq}$ e, se denotarmos por $d_r^{(n)}$ o morfismo $E_r^{(n)}\to E_r^{(n+1)}$ cuja restrição a $E_r^{pq}$ é $d_r^{pq}$ para todo par $(p,q)$ tal que $p+q=n$, então $d_r^{(n+1)}\circ d_r^{(n)}=0$; em outras palavras, $(E_r^{(n)})_{n\in\bb{Z}}$ é um \emph{complexo $E_r^{(\bullet)}$} em $\cat{C}$, com diferencial de grau $+1$, e de (c) deduz-se que $\HH^n(E_r^{(\bullet)})$ é \emph{isomorfo a $E_{r+1}^{(n)}$} para todo $r\geq 2$.
\end{env}

\begin{env}[11.1.2]
\label{0.11.1.2}
Um \emph{morfismo $u:E\to E'$} de uma sequência espectral $E$ em uma sequência espectral $E'=(E_r^{\prime pq},E^{\prime n})$ consiste em sistemas de morfismos $u_r^{pq}:E_r^{pq}\to E_r^{\prime pq}$ e $u^n:E^n\to E^{\prime n}$, sendo os $u^n$ compatíveis com as filtrações de $E^n$ e $E^{\prime n}$, e tais que os diagramas
\[
  \xymatrix{
    E_r^{pq}\ar[r]^-{d_r^{pq}}\ar[d]_{u_r^{pq}} &
    E_r^{p+r,q-r+1}\ar[d]^{u_r^{p+r,q-r+1}}\\
    E_r^{\prime pq}\ar[r]_{d_r^{\prime pq}} &
    E_r^{\prime p+r,q-r+1}
  }
\]
sejam comutativos; além disso, passando aos quocientes, $u_r^{pq}$ dá um morfismo $\overline{u}_r^{pq}:Z_{r+1}(E_r^{pq})/B_{r+1}(E_r^{pq})\to Z_{r+1}(E_r^{\prime pq})/B_{r+1}(E_r^{\prime pq})$ e deve-se ter $\alpha_r^{\prime pq}\circ\overline{u}_r^{pq}=u_{r+1}^{pq}\circ\alpha_r^{pq}$; finalmente, deve-se ter $u_2^{pq}(B_\infty(E_2^{pq}))\subset B_\infty(E_2^{\prime pq})$ e $u_2^{pq}(Z_\infty(E_2^{pq}))\subset Z_\infty(E_2^{\prime pq})$; passando aos quocientes, $u_2^{pq}$ dá então um morfismo $u_\infty^{pq}:E_\infty^{pq}\to E_\infty^{\prime pq}$, e o diagrama
\[
  \xymatrix{
    E_\infty^{pq}\ar[r]^-{u_\infty^{pq}}\ar[d]_{\beta^{pq}} &
    E_\infty^{\prime pq}\ar[d]^{\beta^{\prime pq}}\\
    \gr_p(E^{p+q})\ar[r]_{\gr_p(u^{p+q})} &
    \gr_p(E^{\prime p+q})
  }
\]
deve ser comutativo.

As definições anteriores mostram, por indução sobre $r$, que se os $u_2^{pq}$ são \emph{isomorfismos}, também são os $u_r^{pq}$ para $r\geq 2$; \emph{se também sabemos que $u_2^{pq}(B_\infty(E_2^{pq}))=B_\infty(E_2^{\prime pq})$ e $u_2^{pq}(Z_\infty(E_2^{pq}))=Z_\infty(E_2^{\prime pq})$ e que os $u^n$ são isomorfismos}, então podemos concluir que $u$ é um \emph{isomorfismo}.
\end{env}

\begin{env}[11.1.3]
\label{0.11.1.3}
\oldpage[0\textsubscript{III}]{25}
Lembremos que, se $(F^p(X))_{p\in\bb{Z}}$ é uma \emph{filtração} (decrescente) de um objeto $X\in\cat{C}$, diz-se que essa filtração é \emph{separada} se $\inf(F^p(X))=0$, \emph{discreta} se existe um $p$ tal que $F^p(X)=0$, \emph{exaustiva} (ou \emph{cosseparada}) se $\sup(F^p(X))=X$, e \emph{codiscreta} se existe um $p$ tal que $F^p(X)=X$.

Diz-se que uma sequência espectral $E=(E_r^{pq},E^n)$ é \emph{fracamente convergente} se tiver $B_\infty(E_2^{pq})=\sup_k(B_k(E_2^{pq}))$ e $Z_\infty(E_2^{pq})=\inf_k(Z_k(E_2^{pq}))$ (em outras palavras, os objetos $B_\infty(E_2^{pq})$ e $Z_\infty(E_2^{pq})$ são determinados pelos dados (a)--(c) da sequência espectral $E$).
Diz-se que a sequência espectral $E$ é \emph{regular} se for fracamente convergente e se, além disso:
\begin{enumerate}
  \item[(1st) ] Para todo par $(p,q)$, a sequência decrescente $(Z_k(E_2^{pq}))_{k\geq 2}$ é \emph{estável}; a hipótese de que $E$ é fracamente convergente implica então que $Z_\infty(E_2^{pq})=Z_k(E_2^{pq})$ para $k$ suficientemente grande (dependendo de $p$ e $q$).
  \item[(2nd) ] Para todo $n$, a filtração $(F^p(E^n))_{p\in\bb{Z}}$ de $E^n$ é \emph{discreta} e \emph{exaustiva}.
\end{enumerate}

Diz-se que a sequência espectral $E$ é \emph{corregular} se for fracamente convergente e se, além disso:
\begin{enumerate}
  \item[(3rd) ] Para todo par $(p,q)$, a sequência crescente $(B_k(E_2^{pq}))_{k\geq 2}$ é \emph{estável}, o que implica que $B_\infty(E_2^{pq})=B_k(E_2^{pq})$ e, consequentemente, $E_\infty^{pq}=\inf E_k^{pq}$.
  \item[(4th) ] Para todo $n$, a filtração de $E^n$ é \emph{codiscreta}.
\end{enumerate}

Finalmente, diz-se que $E$ é \emph{birregular} se é ao mesmo tempo regular e corregular, ou seja, se as seguintes condições são cumpridas:
\begin{enumerate}
  \item[(a) ] Para todo par $(p,q)$, as sequências $(B_k(E_2^{pq}))_{k\geq 2}$ e $(Z_k(E_2^{pq}))_{k\geq 2}$ são \emph{estáveis} e tem-se $B_\infty(E_2^{pq})=B_k(E_2^{pq})$ e $Z_\infty(E_2^{pq})=Z_k(E_2^{pq})$ para $k$ suficientemente grande (o que implica que $E_\infty^{pq}=E_k^{pq}$).
  \item[(b)] Para todo $n$, a filtração $(F^p(E^n))_{p\in\bb{Z}}$ é \emph{discreta} e \emph{codiscreta} (também chamada \emph{finita}).
\end{enumerate}

As sequências espectrais definidas em (T, 2.4) são, portanto, sequências espectrais birregulares.
\end{env}

\begin{env}[11.1.4]
\label{0.11.1.4}
Suponhamos que na categoria $\cat{C}$ existam os limites indutivos filtrantes e que o funtor $\varinjlim$ seja \emph{exato} (o que equivale a dizer que se satisfaz o axioma (AB~5) de (T, 1.5) (cf.~T, 1.8)).
A condição de que a filtração $(F^p(X))_{p\in\bb{Z}}$ de um objeto $X\in\cat{C}$ seja exaustiva é então expressa por $\varinjlim_{p\to-\infty}F^p(X)=X$.
Se uma sequência espectral $E$ é fracamente convergente, temos $B_\infty(E_2^{pq})=\varinjlim_{k\to\infty}B_k(E_2^{pq})$; se além disso $u:E\to E'$ é um morfismo de $E$ em uma sequência espectral fracamente convergente $E'$ em $\cat{C}$, então temos $u_2^{pq}(B_\infty(E_2^{pq}))=B_\infty(E_2^{\prime pq})$, pela exatidão de $\varinjlim$.
Além disso:
\end{env}

\begin{proposition}[11.1.5]
\label{0.11.1.5}
Seja $\cat{C}$ uma categoria abeliana em que os limites indutivos filtrantes são exatos, sejam $E$ e $E'$ duas sequências espectrais regulares em $\cat{C}$, e seja $u:E\to E'$ um morfismo de sequências espectrais.
Se os $u_2^{pq}$ são isomorfismos, então $u$ também é.
\end{proposition}

\begin{proof}
Já sabemos por \sref{0.11.1.2} que os $u_r^{pq}$ são isomorfismos e que
\[
  u_2^{pq}(B_\infty(E_2^{pq}))=B_\infty(E_2^{\prime pq});
\]
\oldpage[0\textsubscript{III}]{26}
a hipótese de que $E$ e $E'$ são regulares também implica que $u_2^{pq}(Z_\infty(E_2^{pq}))=Z_\infty(E_2^{\prime pq})$ e, como $u_2^{pq}$ é um isomorfismo, também é $u_\infty^{pq}$; concluímos assim que $\gr_p(u^{p+q})$ é também um isomorfismo.
Mas, como as filtrações dos $E^n$ e os $E^{\prime n}$ são discretas e exaustivas, isso implica que os $u^n$ são também isomorfismos (Bourbaki, \emph{Alg. comm.}, cap.~III, \textsection2, n\textsuperscript{o}8, th.~1).
\end{proof}

\begin{env}[11.1.6]
\label{0.11.1.6}
De (\hyperref[0.11.1.1.2]{11.1.1.2}) e da definição (\hyperref[0.11.1.1.3]{11.1.1.3}) deduz-se que se, para uma sequência espectral $E$, temos $E_r^{pq}=0$, então $E_k^{pq}=0$ para $k\geq r$ e $E_\infty^{pq}=0$.
Diz-se que uma sequência espectral \emph{degenera} se existe um inteiro $r\geq 2$ e, para todo inteiro $n\in\bb{Z}$, um inteiro $q(n)$ tal que \emph{$E_r^{n-q,q}=0$ para todo $q\neq q(n)$}.
A partir da observação anterior deduzimos primeiro que também temos $E_k^{n-q,q}=0$ para $k\geq r$ (incluindo $k=\infty$) e $q\neq q(n)$.
Além disso, a definição de $E_{r+1}^{pq}$ mostra que temos $E_{r+1}^{n-q(n),q(n)}=E_r^{n-q(n),q(n)}$; se $E$ é \emph{fracamente convergente}, temos também $E_\infty^{n-q(n),q(n)}=E_r^{n-q(n),q(n)}$; em outras palavras, para todo $n\in\bb{Z}$, $\gr_p(E^n)=0$ para $p\neq q(n)$ e $\gr_{q(n)}(E^n)=E_r^{n-q(n),q(n)}$.
Se, além disso, a filtração de $E^n$ é \emph{discreta} e \emph{exaustiva}, então a sequência espectral $E$ é \emph{regular}, e tem-se $E^n=E_r^{n-q(n),q(n)}$ salvo isomorfismo.
\end{env}

\begin{env}[11.1.7]
\label{0.11.1.7}
Suponhamos que na categoria $\cat{C}$ existam os limites indutivos filtrantes e sejam exatos, e seja $(E_\lambda,u_{\mu\lambda})$ um sistema indutivo (sobre um conjunto filtrante de índices) de sequências espectrais em $\cat{C}$.
Então o \emph{limite indutivo} deste sistema indutivo existe na categoria aditiva das sequências espectrais de objetos de $\cat{C}$: para vê-lo basta definir $E_r^{pq}$, $d_r^{pq}$, $\alpha_r^{pq}$, $B_\infty(E_2^{pq})$, $Z_\infty(E_2^{pq})$, $E^n$, $F^p(E^n)$ e $\beta^{pq}$ como os respectivos limites indutivos dos $E_{r,\lambda}^{pq}$, $d_{r,\lambda}^{pq}$, $\alpha_{r,\lambda}^{pq}$, $B_\infty(E_{2,\lambda}^{pq})$, $Z_\infty(E_{2,\lambda}^{pq})$, $E_\lambda^n$, $F^p(E_\lambda^n)$ e $\beta_\lambda^{pq}$; a verificação das condições de \sref{0.11.1.1} deduz-se da exatidão do funtor $\varinjlim$ sobre $\cat{C}$.
\end{env}

\begin{remark}[11.1.8]
\label{0.11.1.8}
Suponhamos que a categoria $\cat{C}$ é a categoria dos $A$-módulos sobre um anel $A$ \emph{noetheriano} (resp. um anel $A$).
As definições de \sref{0.11.1.1} mostram então que se, para um dado $r$, os $E_r^{pq}$ são $A$-módulos \emph{de tipo finito} (resp. \emph{de comprimento finito}), também o são todos os módulos $E_s^{pq}$ para $s\geq r$ e, portanto, também $E_\infty^{pq}$.
Se, além disso, a filtração do termo limite $(E^n)$ é \emph{discreta} e \emph{codiscreta} para todo $n$, concluímos que cada $E^n$ é também um $A$-módulo \emph{de tipo finito} (resp. \emph{de comprimento finito}).
\end{remark}

\begin{env}[11.1.9]
\label{0.11.1.9}
Teremos de considerar condições que garantam que uma sequência espectral $E$ seja birregular de forma ``uniforme'' em $p+q=n$.
Usaremos então o seguinte lema:
\end{env}

\begin{lemma}[11.1.10]
\label{0.11.1.10}
Seja $(E_r^{pq})$ uma família de objetos de $\cat{C}$ relacionada pelos dados \emph{(a) }, \emph{(b) } e \emph{(c) } de \sref{0.11.1.1}.
Para um inteiro fixo $n$, as seguintes propriedades são equivalentes:
\begin{enumerate}
  \item[{\rm(a) }] Existe um inteiro $r(n)$ tal que, para $r\geq r(n)$, $p+q=n$ ou $p+q=n-1$, os morfismos $d_r^{pq}$ são nulos.
  \item[{\rm(b) }] Existe um inteiro $r(n)$ tal que, para $p+q=n$ ou $p+q=n+1$, há $B_r(E_2^{pq})=B_s(E_2^{pq})$ para $s\geq r\geq r(n)$.
  \item[{\rm(c) }] Existe um inteiro $r(n)$ tal que, para $p+q=n$ ou $p+q=n-1$, temos $Z_r(E_2^{pq})=Z_s(E_2^{pq})$ para $s\geq r\geq r(n)$.
  \item[{\rm(d) }] Existe um inteiro $r(n)$ tal que, para $p+q=n$, temos $B_r(E_2^{pq})=B_s(E_2^{pq})$ e $Z_r(E_2^{pq})=Z_s(E_2^{pq})$ para $s\geq r\geq r(n)$.
\end{enumerate}
\end{lemma}

\begin{proof}
\oldpage[0\textsubscript{III}]{27}
De acordo com as condições (a), (b) e (c) de \sref{0.11.1.1}, dizer que $Z_{r+1}(E_2^{pq})=Z_r(E_2^{pq})$ equivale a dizer que $d_r^{pq}=0$, e dizer que $B_r(E_2^{p+r,q-r+1})=B_{r+1}(E_2^{p+r,q-r+1})$ equivale a dizer que $d_r^{pq}=0$; o lema deduz-se imediatamente desta observação.
\end{proof}

\subsection{A sequência espectral de um complexo filtrado}
\label{subsection:0.11.2}

\begin{env}[11.2.1]
\label{0.11.2.1}
Seja $\cat{C}$ uma categoria abeliana; concordamos em denotar por notações como $K^\bullet$ os \emph{complexos $(K^i)_{i\in\bb{Z}}$} de objetos de $\cat{C}$ cujo diferencial é de grau $+1$, e por notações como $K_\bullet$ os complexos $(K_i)_{i\in\bb{Z}}$ de objetos de $\cat{C}$ cujo diferencial é de grau $-1$.
A todo complexo $K^\bullet=(K^i)$ cujo diferencial $d$ é de grau $+1$ podemos associar um complexo $K_\bullet'=(K_i')$ colocando $K_i'=K^{-i}$, sendo o diferencial $K_i'\to K_{i-1}'$ o operador $d:K^{-i}\to K^{-i+1}$; e \emph{vice-versa}, o que permitirá, conforme as circunstâncias, considerar um ou outro dos tipos de complexos e traduzir qualquer resultado de um dos tipos em resultados para o outro.
Denotamos de forma análoga por notações como $K^{\bullet\bullet}=(K^{ij})$ (resp. $K_{\bullet\bullet}=(K_{ij})$) os \emph{bicomplexos} (ou \emph{complexos duplos}) de objetos de $\cat{C}$ em que os \emph{dois} diferenciais são de grau $+1$ (resp. $-1$); ainda se pode passar de um tipo para outro alterando os sinais dos índices, e temos notações e observações análogas para multicomplexos quaisquer.
As notações $K^\bullet$ e $K_\bullet$ serão também usadas para \emph{objetos $\bb{Z}$-graduados} de $\cat{C}$, que não são necessariamente complexos (podem ser considerados como tais para os diferenciais \emph{nulos}); por exemplo, escrevemos $\HH^\bullet(K^\bullet)=(\HH^i(K^\bullet))_{i\in\bb{Z}}$ para a \emph{cohomologia} de um complexo $K^\bullet$ cujo diferencial é de grau $+1$, e $\HH_\bullet(K_\bullet)=(\HH_i(K_\bullet))_{i\in\bb{Z}}$ para a \emph{homologia} de um complexo $K_\bullet$ cujo diferencial é de grau $-1$; quando se passa de $K^\bullet$ para $K_\bullet'$ pelo procedimento descrito acima, tem-se $\HH_i(K_\bullet')=\HH^{-i}(K^\bullet)$.

Lembre-se neste caso que, para um complexo $K^\bullet$ (resp. $K_\bullet$), vamos escrever em geral $Z^i(K^\bullet)=\Ker(K^i\to K^{i+1})$ (``objeto de cociclos'') e $B^i(K^\bullet)=\Im(K^{i-1}\to K^i)$ (``objeto de cobordes'') (resp. $Z_i(K_\bullet)=\Ker(K_i\to K_{i-1})$ (``objeto de ciclos'') e $B_i(K_\bullet)=\Im(K_{i+1}\to K_i)$ (``objeto de bordas'')), de modo que $\HH^i(K^\bullet)=Z^i(K^\bullet)/B^i(K^\bullet)$ (resp. $\HH_i(K_\bullet)=Z_i(K_\bullet)/B_i(K_\bullet)$).

Se $K^\bullet=(K^i)$ (resp. $K_\bullet=(K_i)$) é um complexo em $\cat{C}$ e $T:\cat{C}\to\cat{C}'$ um funtor de $\cat{C}$ em uma categoria abeliana $\cat{C}'$, denotamos por $T(K^\bullet)$ (resp. $T(K_\bullet)$) o complexo $(T(K^i))$ (resp. $(T(K_i))$) em $\cat{C}'$.

Não vamos rever a definição dos \emph{$\partial$-funtores} (T, 2.1), exceto para observar que se diz \emph{também} $\partial$-funtor em vez de $\partial^*$-funtor quando o morfismo $\partial$ diminui o grau em uma unidade; se puder haver confusão, o contexto esclarecerá este ponto.

Finalmente, diz-se que um \emph{objeto graduado $(A_i)_{i\in\bb{Z}}$} de $\cat{C}$ é \emph{limitado inferiormente} (resp. \emph{superiormente}) se existe um $i_0$ tal que $A_i=0$ para $i<i_0$ (resp. $i>i_0$).
\end{env}

\begin{env}[11.2.2]
\label{0.11.2.2}
Seja $K^\bullet$ um complexo em $\cat{C}$ cujo diferencial $d$ é de grau $+1$, e suponhamos que está munido de uma \emph{filtração $F(K^\bullet)=(F^p(K^\bullet))_{p\in\bb{Z}}$} formada por subobjetos \emph{graduados}
\oldpage[0\textsubscript{III}]{28}
de $K^\bullet$, isto é, $F^p(K^\bullet)=(K^i\cap F^p(K^\bullet))_{i\in\bb{Z}}$; supomos ainda que $d(F^p(K^\bullet))\subset F^p(K^\bullet)$ para todo $p\in\bb{Z}$.
Lembremos rapidamente como se define \emph{funtorialmente} uma sequência espectral $E(K^\bullet)$ a partir de $K^\bullet$ (M,~XV,~4 e G,~I,~4.3).
Para $r\geq 2$, o morfismo canônico $F^p(K^\bullet)/F^{p+r}(K^\bullet)\to F^p(K^\bullet)/F^{p+1}(K^\bullet)$ define um morfismo em cohomologia
\[
  \HH^{p+q}(F^p(K^\bullet)/F^{p+r}(K^\bullet))\to\HH^{p+q}(F^p(K^\bullet)/F^{p+1}(K^\bullet)).
\]

Denotamos por $Z_r^{pq}(K^\bullet)$ a imagem deste morfismo.
Analogamente, da sequência exata
\[
  0\to F^p(K^\bullet)/F^{p+1}(K^\bullet)\to F^{p-r+1}(K^\bullet)/F^{p+1}(K^\bullet)\to F^{p-r+1}(K^\bullet)/F^p(K^\bullet)\to 0,
\]
deduzimos da sequência exata da cohomologia um morfismo
\[
  \HH^{p+q-1}(F^{p-r+1}(K^\bullet)/F^p(K^\bullet))\to\HH^{p+q}(F^p(K^\bullet)/F^{p+1}(K^\bullet)),
\]
e denotamos por $B_r^{pq}(K^\bullet)$ a imagem deste morfismo; demonstra-se que $B_r^{pq}(K^\bullet)\subset Z_r^{pq}(K^\bullet)$ e toma-se $E_r^{pq}(K^\bullet)=Z_r^{pq}(K^\bullet)/B_r^{pq}(K^\bullet)$; não precisaremos da definição dos $d_r^{pq}$ nem dos $\alpha_r^{pq}$.

Observemos aqui que todos os $Z_r^{pq}(K^\bullet)$ e $B_r^{pq}(K^\bullet)$, para $p$ e $q$ fixos, são subobjetos do mesmo objeto $\HH^{p+q}(F^p(K^\bullet)/F^{p+1}(K^\bullet))$, que denotamos por $Z_1^{pq}(K^\bullet)$; colocamos $B_1^{pq}(K^\bullet)=0$, de modo que as definições anteriores de $Z_r^{pq}(K^\bullet)$ e $B_r^{pq}(K^\bullet)$ também são válidas para $r=1$; colocamos também $E_1^{pq}(K^\bullet)=Z_1^{pq}(K^\bullet)$.
Definimos $d_1^{pq}$ e $\alpha_1^{pq}$ de modo que as condições de \sref{0.11.1.1} para $r=1$ sejam satisfeitas.
Por outro lado, definimos os subobjetos $Z_\infty^{pq}(K^\bullet)$ como a imagem do morfismo
\[
  \HH^{p+q}(F^p(K^\bullet))\to\HH^{p+q}(F^p(K^\bullet)/F^{p+1}(K^\bullet))=E_1^{pq}(K^\bullet),
\]
e $B_\infty^{pq}(K^\bullet)$ como a imagem do morfismo
\[
  \HH^{p+q-1}(K^\bullet/F^p(K^\bullet))\to\HH^{p+q}(F^p(K^\bullet)/F^{p+1}(K^\bullet))=E_1^{pq}(K^\bullet),
\]
induzido como antes pela sequência exata de cohomologia.
Definimos $Z_\infty(E_2^{pq}(K^\bullet))$ e $B_\infty(E_2^{pq}(K^\bullet))$ como as imagens canônicas em $E_2^{pq}(K^\bullet)$ de $Z_\infty^{pq}(K^\bullet)$ e $B_\infty^{pq}(K^\bullet)$.

Finalmente, denotamos por $F^p(\HH^n(K^\bullet))$ a imagem em $\HH^n(K^\bullet)$ do morfismo $\HH^n(F^p(K^\bullet))\to\HH^n(K^\bullet)$ induzido pela injeção canônica $F^p(K^\bullet)\to K^\bullet$; pela sequência exata de cohomologia, este também é o núcleo do morfismo $\HH^n(K^\bullet)\to\HH^n(K^\bullet/F^p(K^\bullet))$.
Isto define uma filtração sobre $E^n(K^\bullet)=\HH^n(K^\bullet)$; não vamos dar aqui a definição dos isomorfismos $\beta^{pq}$.
\end{env}

\begin{env}[11.2.3]
\label{0.11.2.3}
O caráter \emph{funtorial} de $E(K^\bullet)$ entende-se do seguinte modo: dados dois complexos \emph{filtrados} $K^\bullet$ e $K^{\prime\bullet}$ em $\cat{C}$ e um morfismo de complexos $u:K^\bullet\to K^{\prime\bullet}$ \emph{compatível com as filtrações}, induzem-se de maneira evidente os morfismos $u_r^{pq}$ (para $r\geq 1$) e $u^n$, e demonstra-se que esses morfismos são compatíveis com $d_r^{pq}$, $\alpha_r^{pq}$ e $\beta^{pq}$ no sentido de \sref{0.11.1.2}, e definem assim um morfismo bem definido $E(u):E(K^\bullet)\to E(K^{\prime\bullet})$ de sequências espectrais.
Além disso, é demonstrado que se $u$ e $v$ são morfismos $K^\bullet\to K^{\prime\bullet}$ do tipo anterior, \emph{homotópicos em grau $\leq k$}, então $u_r^{pq}=v_r^{pq}$ para $r>k$ e $u^n=v^n$ para todo $n$ (M,~XV,~3.1).
\end{env}

\begin{env}[11.2.4]
\label{0.11.2.4}
\oldpage[0\textsubscript{III}]{29}
Suponhamos que os limites indutivos filtrantes em $\cat{C}$ são exatos.
Então, se a filtração $(F^p(K^\bullet))$ de $K^\bullet$ é \emph{exaustiva}, também o é a filtração $(F^p(\HH^n(K^\bullet)))$ para todo $n$, pois por hipótese temos $K^\bullet=\varinjlim_{p\to-\infty}F^p(K^\bullet)$ e a hipótese sobre $\cat{C}$ implica que a cohomologia comuta com os limites indutivos.
Além disso, pela mesma razão, temos $B_\infty(E_2^{pq}(K^\bullet))=\sup_k B_k(E_2^{pq}(K^\bullet))$.
Diz-se que a filtração $(F^p(K^\bullet))$ de $K^\bullet$ é \emph{regular} se para todo $n$ existe um inteiro $u(n)$ tal que $\HH^n(F^p(K^\bullet))=0$ para $p>u(n)$.
Este é, em particular, o caso quando a filtração de $K^\bullet$ é \emph{discreta}.
Quando a filtração de $K^\bullet$ é regular e exaustiva, e os limites indutivos filtrantes são exatos em $\cat{C}$, obtém-se (M,~XV,~4) que a sequência espectral $E(K^\bullet)$ é \emph{regular}.
\end{env}

\subsection{As sequências espectrais de um bicomplexo}
\label{subsection:0.11.3}

\begin{env}[11.3.1]
\label{0.11.3.1}
Quanto às convenções para os bicomplexos, seguimos as de ~(T,~2.4) em vez das de ~(M), supondo assim que os dois diferenciais $d'$ e $d''$ (de grau~$+1$) de um bicomplexo $K^{\bullet\bullet}=(K^{ij})$ são \emph{permutáveis}.
Suponhamos que \emph{uma} das duas condições seguintes seja satisfeita:
\begin{enumerate}
   \item Existem \emph{somas diretas infinitas} em $\cat{C}$;
   \item Para todo $n\in\bb{Z}$, há apenas um número \emph{finito} de pares $(p,q)$ tais como $p+q=n$ e $K^{pq}\neq 0$.
 \end{enumerate}
Então o bicomplexo $K^{\bullet\bullet}$ define um \emph{complexo} (simples) $(K^{\prime n})_{n\in\bb{Z}}$ com $K^{\prime n}=\sum_{i+j=n}K^{ij}$, cujo diferencial $d$ (de grau~$+1$) é dado por $dx=d'x+(-1)^i d''x$ para $x\in K^{ij}$.
Quando mais adiante falarmos do \emph{complexo (simples) definido por um bicomplexo $K^{\bullet\bullet}$}, ficará sempre subentendido que uma das condições anteriores é satisfeita.
Adotamos convenções análogas para os multicomplexos.

Denotamos por $K^{i,\bullet}$ (resp.~$K^{\bullet,j}$) o complexo simples $(K^{ij})_{j\in\bb{Z}}$ (resp.~$(K^{ij})_{i\in\bb{Z}}$), e por $Z_\mathrm{II}^p(K^{i,\bullet})$, $B_\mathrm{II}^p(K^{i,\bullet})$, $\HH_\mathrm{II}^p(K^{i,\bullet})$ (resp.~$Z_\mathrm{I}^p(K^{\bullet,j})$, $B_\mathrm{I}^p(K^{\bullet,j})$, $\HH_\mathrm{I}^p(K^{\bullet,j})$) seus objetos de cociclos, de cobordes e de cohomologia de grau $p$, respectivamente;
o diferencial $d':K^{i,\bullet}\to K^{i+1,\bullet}$ é um morfismo de complexos, que dá assim um operador sobre os cociclos, os cobordes e a cohomologia,
\begin{align*}
  d'&:Z_\mathrm{II}^p(K^{i,\bullet})\to Z_\mathrm{II}^p(K^{i+1,\bullet})
\\d'&:B_\mathrm{II}^p(K^{i,\bullet})\to B_\mathrm{II}^p(K^{i+1,\bullet})
\\d'&:\HH_\mathrm{II}^p(K^{i,\bullet})\to\HH_\mathrm{II}^p(K^{i+1,\bullet})
\end{align*}
e é claro que, para estes operadores, $(Z_\mathrm{II}^p(K^{i,\bullet}))_{i\in\bb{Z}}$, $(B_\mathrm{II}^p(K^{i,\bullet}))_{i\in\bb{Z}}$ e $(\HH_\mathrm{II}^p(K^{i,\bullet}))_{i\in\bb{Z}}$ são complexos;
denotamos o complexo $(\HH_\mathrm{II}^p(K^{i,\bullet}))_{i\in\bb{Z}}$ por $\HH_\mathrm{II}^p(K^{\bullet\bullet})$, e seus objetos de cociclos, cobordes e cohomologia de grau $q$ por $Z_\mathrm{I}^q(\HH_\mathrm{II}^p(K^{\bullet\bullet}))$, $B_\mathrm{I}^q(\HH_\mathrm{II}^p(K^{\bullet\bullet}))$ e $\HH_\mathrm{I}^q(\HH_\mathrm{II}^p(K^{\bullet\bullet}))$.
Definem-se de forma análoga os complexos $\HH_\mathrm{I}^p(K^{\bullet\bullet})$ e os seus objetos de cohomologia $\HH_\mathrm{II}^q(\HH_\mathrm{I}^p(K^{\bullet\bullet}))$.
Lembre-se, no entanto, que $\HH^n(K^{\bullet\bullet})$ denota o objeto de grau $n$ da cohomologia do complexo \emph{(simples)} definido por $K^{\bullet\bullet}$.
\end{env}

\begin{env}[11.3.2]
\label{0.11.3.2}
Sobre o complexo definido por um bicomplexo $K^{\bullet\bullet}$ podem ser consideradas duas filtrações canônicas, $(F_\mathrm{I}^p(K^{\bullet\bullet}))$ e $(F_\mathrm{II}^p(K^{\bullet\bullet}))$, dadas por
\[
\label{0.11.3.2.1}
  F_\mathrm{I}^p(K^{\bullet\bullet})
  = \left(\sum_{i+j=n,i\geq p}K^{ij}\right)_{n\in\bb{Z}}
  \quad\text{e}\quad
  F_\mathrm{II}^p(K^{\bullet\bullet})
  = \left(\sum_{i+j=n,j\geq p}K^{ij}\right)_{n\in\bb{Z}}
  \tag{11.3.2.1}
\]
\oldpage[0\textsubscript{III}]{30}
que, por definição, são subobjetos graduados do complexo (simples) definido por $K^{\bullet\bullet}$ e, assim, tornam este complexo um complexo filtrado;
além disso, é claro que estas filtrações são \emph{exaustivas} e \emph{separadas}.

Cada uma dessas filtrações corresponde a uma sequência espectral~\sref{0.11.2.2};
designamos por $'E(K^{\bullet\bullet})$ e $''E(K^{\bullet\bullet})$ as sequências espectrais correspondentes a $(F_\mathrm{I}^p(K^{\bullet\bullet}))$ e $(F_\mathrm{II}^p(K^{\bullet\bullet}))$, respectivamente, chamadas \emph{sequências espectrais do bicomplexo $K^{\bullet\bullet}$}, e que têm ambas por termo limite a cohomologia $(\HH^n(K^{\bullet\bullet}))$.
Também é demonstrado~(M,~XV,~6) que se tem
\[
\label{0.11.3.2.2}
  'E_2^{pq}(K^{\bullet\bullet})
  = \HH_\mathrm{I}^p(\HH_\mathrm{II}^q(K^{\bullet\bullet}))
  \quad\text{e}\quad
  ''E_2^{pq}(K^{\bullet\bullet})
  = \HH_\mathrm{II}^q(\HH_\mathrm{I}^p(K^{\bullet\bullet})).
  \tag{11.3.2.2}
\]

Todo morfismo $u:K^{\bullet\bullet}\to K^{\prime\bullet\bullet}$ de bicomplexos é \emph{ipso facto} compatível com as filtrações do mesmo tipo de $K^{\bullet\bullet}$ e $K^{\prime\bullet\bullet}$, e define assim um morfismo para cada uma das duas sequências espectrais;
além disso, dois morfismos \emph{homotópicos} definem uma homotopia \emph{de ordem $\leq 1$} dos complexos (simples) filtrados correspondentes, e, portanto, o \emph{mesmo} morfismo para cada uma das duas sequências espectrais~(M,~XV,~6.1).
\end{env}

\begin{proposition}[11.3.3]
\label{0.11.3.3}
Seja $K^{\bullet\bullet}=(K^{ij})$ um bicomplexo em uma categoria abeliana $\cat{C}$.
\begin{enumerate}
  \item[{\rm(i) }] Se existem $i_0$ e $j_0$ tais que $K^{ij}=0$ para $i<i_0$ ou $j<j_0$ (resp.~$i>i_0$ ou $j>j_0$), então as duas sequências espectrais $'E(K^{\bullet\bullet})$ e $''E(K^{\bullet\bullet})$ são bi-regulares.
  \item[{\rm(ii) }] Se existem $i_0$ e $i_1$ tais como $K^{ij}=0$ para $i<i_0$ ou $i>i_1$ (resp.~se existem $j_0$ e $j_1$ tais como $K^{ij}=0$ para $j<j_0$ ou $j>j_1$), então as duas sequências espectrais $'E(K^{\bullet\bullet})$ e $''E(K^{\bullet\bullet})$ são birregulares.
\end{enumerate}
Suponhamos ainda que os limites indutivos filtrantes existam e sejam exatos em $\cat{C}$. Então:
\begin{enumerate}
  \item[{\rm(iii) }] Se existe $i_0$ tal que $K^{ij}=0$ para $i>i_0$ (resp.~se existe $j_0$ tal que $K^{ij}=0$ para $j<j_0$), então a sequência espectral $'E(K^{\bullet\bullet})$ é regular.
  \item[{\rm(iv) }] Se existe $i_0$ tal que $K^{ij}=0$ para $i<i_0$ (resp.~se existe $j_0$ tal que $K^{ij}=0$ para $j>j_0$), então a sequência espectral $''E(K^{\bullet\bullet})$ é regular.
\end{enumerate}
\end{proposition}

\begin{proof}
A proposição deduz-se imediatamente das definições ~\sref{0.11.1.3} e de ~\sref{0.11.2.4}, bem como das seguintes observações relativas à filtração $F_\mathrm{I}$ (e das observações análogas que são deduzidas para $F_\mathrm{II}$ trocando os papéis dos dois índices em $K^{\bullet\bullet}$):
\begin{enumerate}
  \item[{$1^{\circ}$}] Se existe $i_0$ tal que $K^{ij}=0$ para $i>i_0$, então a filtração $F_\mathrm{I}(K^{\bullet\bullet})$ é \emph{discreta}.
    \item[{$2^{\circ}$}] Se existe $i_0$ tal que $K^{ij}=0$ para $i<i_0$, então a filtração $F_\mathrm{I}(K^{\bullet\bullet})$ é \emph{codiscreta}.
    Deduz-se imediatamente que o mesmo ocorre para a filtração correspondente $F_\mathrm{I}(\HH^n(K^{\bullet\bullet}))$, para todo $n$.
    Além disso, a definição dos $B_{r}^{pq}$ correspondentes à filtração $F_\mathrm{I}(K^{\bullet\bullet})$ \sref{0.11.2.2} mostra que, para todo par $(p,q)$, a sequência $(B_{r}^{pq})_{r\geq 2}$ é estável.
  \item[{$3^{\circ}$}] Se há $j_0$ tal que $K^{ij}=0$ para $j<j_0$, então temos
    \[
      F_\mathrm{I}^{p+r}(K^{\bullet\bullet})
      \cap (\sum_{i+j=n}K^{ij})
      = 0
    \]
    quando $p+r+j_{0}>n$, logo $Z_{r}^{pq} = Z_\infty(E_2^{pq})$ para $r>q-j_{0}+1$.
    Por outro lado, $\HH^{n}(F_\mathrm{I}^{p}(K^{\bullet\bullet})) = 0$ para $p>n-j_{0}+1$.
  \item[{$4^{\circ}$}] Se há $j_0$ tal que $K^{ij}=0$ para $j>j_0$, então temos
    \[
      F_\mathrm{I}^{p-r+1}(K^{\bullet\bullet})
      \cap (\sum_{i+j=n}K^{ij})
      = \sum_{i+j=n}K^{ij}
    \]
    \oldpage[0\textsubscript{III}]{31}
    quando $p-r+1+j_{0}<n$, logo $B_{r}^{pq} = B_\infty(E_2^{pq})$ para $r>j_{0}-q+1$.
    Por outro lado, $\HH^{n}(F_\mathrm{I}^{p}(K^{\bullet\bullet})) = \HH^{n}(K^{\bullet\bullet})$ para $p+j_{0}<n-1$.
\end{enumerate}
\end{proof}

\begin{env}[11.3.4]
\label{0.11.3.4}
Suponhamos que o bicomplexo $K^{\bullet\bullet}=(K^{ij})$ é tal que $K^{ij}=0$ para $i<0$ ou $j<0$.
Sabemos que pode ser definido, para todo $p\in\bb{Z}$, um homomorfismo canônico de borda.
\[
\label{0.11.3.4.1}
  'E_2^{p0}(K^{\bullet\bullet})
  \rightarrow \HH^{p}(K^{\bullet\bullet})
\tag{11.3.4.1}
\]
(M,~XV,~6).
Lembremos que isso se deve, por um lado, à sequência espectral
$'E(K^{\bullet\bullet})$, pois $Z_r^{p0}=Z_{\mathrm{I}}^{p}(Z_{\mathrm{II}}^{0}(K^{\bullet\bullet}))$ para $2\leq r \leq +\infty$, e, por outro lado, ao fato de que $\HH^{p}(F_\mathrm{I}^{p+1}(K^{\bullet\bullet})) = 0$, de modo que o isomorfismo $\beta^{p0}:\ 'E_\infty^{p0} \simto \HH^{p}(F_\mathrm{I}^p)/\HH^{p}(F_\mathrm{I}^{p+1})$ dá um homomorfismo $'E_\infty^{p0} \rightarrow \HH^{p}(F_\mathrm{I}^p(K^{\bullet\bullet})) \rightarrow \HH^{p}(K^{\bullet\bullet})$.
A igualdade de todos os $Z_r^{p0}$ permite definir homomorfismos canônicos $'E_r^{p0} \rightarrow 'E_s^{p0}$ para $r \leq s$ e, em particular, um homomorfismo $'E_2^{p0} \rightarrow 'E_\infty^{p0}$, de onde, por composição, o homomorfismo de borda $'E_2^{p0}(K^{\bullet\bullet}) \rightarrow \HH^{p}(K^{\bullet\bullet})$.
Além disso, verifica-se imediatamente que o homomorfismo de borda assim definido envia a classe módulo $B_2^{p0}$ de um elemento $z\in Z_\mathrm{II}^0(K^{\bullet\bullet})\subset K^{p0}$ tal que $d'z=0$ na classe de $z$ módulo $B_\infty^{p0}$ em $'E_\infty^{p0}$, e depois na classe de cohomologia de $z$ em $\HH^p(K^{\bullet\bullet})$.
Vemos, portanto, finalmente que o \emph{homomorfismo de borda} \sref{0.11.3.4.1} \emph{provém, ao passar à cohomologia, da injeção canônica} $Z_\mathrm{II}^0(K^{\bullet\bullet}) \rightarrow K^{\bullet\bullet}$ (onde $K^{\bullet\bullet}$ é considerado como complexo simples).
Interpretamos naturalmente, da mesma maneira, o homomorfismo de borda
\[
\label{0.11.3.4.2}
  ''E_2^{p0}(K^{\bullet\bullet})
  \rightarrow \HH^{p}(K^{\bullet\bullet})
\tag{11.3.4.2}
\]
como proveniente da injeção canônica $Z_\mathrm{I}^0(K^{\bullet\bullet}) \rightarrow K^{\bullet\bullet}$.
\end{env}

\begin{env}[11.3.5]
\label{0.11.3.5}
Seja agora $K_{\bullet\bullet} = (K_{ij})$ um bicomplexo em $\cat{C}$ cujos dois operadores diferenciais são de grau~$-1$.
Escrevemos então $K_{i,\bullet}$ e $K_{\bullet,j}$ para designar, respectivamente, os complexos simples $(K_{ij})_{j\in\bb{Z}}$ e $(K_{ij})_{i\in\bb{Z}}$; $\HH_{p}^\mathrm{II}(K_{i, \bullet})$ e $\HH_{p}^\mathrm{I}(K_{\bullet, j})$ para designar os respectivos objetos de homologia em grau $p$; $\HH_{p}^\mathrm{II}(K_{\bullet\bullet})$ e $\HH_{p}^\mathrm{I}(K_{\bullet\bullet})$ para designar, respectivamente, os complexos $(\HH_{p}^\mathrm{II}(K_{i, \bullet}))_{i \in\bb{Z}}$ e $(\HH_{p}^\mathrm{I}(K_{\bullet, j}))_{j \in\bb{Z}}$; e $\HH_q^\mathrm{I}(\HH_{p}^\mathrm{II}(K_{\bullet\bullet}))$ e $\HH_q^\mathrm{II}(\HH_{p}^\mathrm{I}(K_{\bullet\bullet}))$ para designar os respectivos objetos de homologia em grau $q$;
utilizamos notações análogas para os objetos de ciclos e os objetos de bordas;
finalmente, denotaremos por $\HH_n(K_{\bullet\bullet})$ o objeto de homologia de grau $n$ (quando existir) do complexo simples (com um operador diferencial de grau~$-1$) definido por $K_{\bullet\bullet}$.
Seja $K'^{\bullet\bullet}=(K'^{ij})$, com $K'^{ij}=K_{-i,-j}$, o bicomplexo com operadores diferenciais de grau~$+1$ associado a $K_{\bullet\bullet}$.
Por definição, as \emph{sequências espectrais} de $K_{\bullet\bullet}$ são as de $K'^{\bullet\bullet}$, que denotamos por $'E(K_{\bullet\bullet})$ e $''E(K_{\bullet\bullet})$, mudando, no entanto, a notação ao colocar
\[
  'E_{pq}^{r}(K_{\bullet\bullet})
  = 'E_{r}^{-p,-q}(K'^{\bullet\bullet})
  \quad\text{e}\quad
  ''E_{pq}^{r}(K_{\bullet\bullet})
  = ''E_{r}^{-p,-q}(K'^{\bullet\bullet}), 
\]
para $2\leq r \leq \infty$.
Com esta notação, temos
\[
  'E_{pq}^{2}(K_{\bullet\bullet})
  = \HH_{p}^\mathrm{I}(\HH_{q}^\mathrm{II}(K_{\bullet\bullet}))
  \quad\text{e}\quad
  ''E_{pq}^{2}(K_{\bullet\bullet})
  = \HH_{q}^\mathrm{II}(\HH_{p}^\mathrm{I}(K_{\bullet\bullet})).
\]
Para evitar erros de sinal, será preferível, em geral, voltar ao complexo $K'^{\bullet\bullet}$ para as relações entre essas sequências espectrais e suas convergências.
Assinalemos, no entanto, os critérios correspondentes a \sref{0.11.3.3}:
\end{env}

\begin{env}[11.3.6]
\label{0.11.3.6}
As sequências espectrais $'E(K_{\bullet\bullet})$ e $''E(K_{\bullet\bullet})$ são \emph{birregulares} nos seguintes casos:
\begin{enumerate}
  \item[(a) ] Existem $i_0$ e $j_0$ tais que $K_{ij}=0$ para $i>i_0$ ou $j>j_0$ (resp. para $i<i_0$
    \oldpage[0\textsubscript{III}]{32}
    ou $j<j_0$);
  \item[(b) ] Existem $i_0$ e $i_1$ tais que $K_{ij}=0$ para $i<i_0$ e $i>i_1$;
  \item[(c) ] Existem $j_0$ e $j_1$ tais que $K_{ij}=0$ para $j<j_0$ e $j>j_1$.
\end{enumerate}

A sequência $'E(K_{\bullet\bullet})$ é \emph{regular} se existe $i_0$ tal que $K_{ij}=0$ para $i<i_0$, ou se existe $j_0$ tal que $K_{ij}=0$ para $j>j_0$.

A sequência $''E(K_{\bullet\bullet})$ é \emph{regular} se existe $i_0$ tal que $K_{ij}=0$ para $i>i_0$, ou se existe $j_0$ tal que $K_{ij}=0$ para $j<j_0$.
\end{env}


\subsection{Hipercohomologia de um funtor em relação a um complexo $K^{\bullet}$}
\label{subsection:0.11.4}

\begin{env}[11.4.1]
\label{0.11.4.1}
Seja $\cat{C}$ uma categoria abeliana.
Lembre-se de que se definiu uma \emph{resolução pela direita} (ou \emph{resolução cohomológica}) de um objeto $A$ de $\cat{C}$ como um complexo de objetos de $\cat{C}$, cujo operador diferencial é de grau $+1$,
\[
  0 \rightarrow L^0 \rightarrow L^1 \rightarrow L^2 \rightarrow \ldots
\]
munido de um morfismo $\varepsilon: A \rightarrow L^0$, chamado de \emph{aumentação} da resolução, que pode ser considerado como um morfismo de complexos
\[
  \xymatrix{
    0 \ar[r]
    & A \ar[r] \ar[d]_{\varepsilon}
    & 0 \ar[r] \ar[d]
    & 0 \ar[r] \ar[d]
    & \ldots
  \\0 \ar[r]
    & L^0 \ar[r]
    & L^1 \ar[r]
    & L^2 \ar[r]
    & \ldots
  }
\]
tal que a sequência
\[
  0 \rightarrow A \xrightarrow{\varepsilon} L^0 \rightarrow L^1 \ldots
\]
seja exata.
Analogamente, uma \emph{resolução pela esquerda} (ou \emph{resolução homológica}) de $A$ é um complexo $0 \leftarrow L_0 \leftarrow L_1 \leftarrow L_2 \ldots$ de objetos de $\cat{C}$, cujo operador diferencial é de grau~$-1$, munido de uma \emph{aumentação} $L_0 \xrightarrow{\varepsilon} A$ tal que a sequência
\[
  0 \leftarrow A \xleftarrow{\varepsilon} L_0 \leftarrow L_1 \ldots
\]
seja exata.

Quando a resolução à direita $(L^i)_{i\geq 0}$ de um objeto $A$ é tal que $L^i=0$ para $i\geq n+1$, diz-se que esta resolução é \emph{de comprimento $\leq n$}.
Define-se, analogamente, uma resolução pela esquerda de comprimento $\leq n$.
Uma resolução que é de comprimento $\leq n$ para algum inteiro $n$ chama-se \emph{finita}.

Uma resolução de $A$ chama-se \emph{projetiva} (resp. \emph{injetiva}) se os objetos de $\cat{C}$ diferentes de $A$ que compõem a mesma são \emph{projetivos} (resp. \emph{injetivos}).
Quando $\cat{C}$ é a categoria dos módulos (à esquerda, por exemplo) sobre um anel, diz-se que uma resolução de $A$ é \emph{plana} (resp. \emph{livre}) quando os módulos diferentes de $A$ que a compõem são \emph{planos} (resp. \emph{livres}).
\end{env}

\begin{env}[11.4.2]
\label{0.11.4.2}
Seja $K^\bullet = (K^i)_{i\in\bb{Z}}$ um complexo de objetos de $\cat{C}$, cujo operador diferencial é de grau~$+1$.
Define-se uma \emph{resolução de Cartan--Eilenberg à direita} de $K^\bullet$ como um par formado por um bicomplexo $L^{\bullet\bullet} = (L^{ij})$ com operadores diferenciais de grau~$+1$ e tal que $L^{ij}=0$ para $j<0$, e um morfismo de complexos simples $\varepsilon: K^\bullet \rightarrow L^{\bullet, 0}$, de modo que se cumpram as seguintes condições:
\oldpage[0\textsubscript{III}]{33}
\begin{enumerate}
  \item[(i) ] Para cada índice $i$, as sequências
    \[
      \begin{aligned}
        0 \rightarrow
        & K^i
        \xrightarrow{\varepsilon} L^{i0}
        \rightarrow L^{i1}
        \rightarrow \ldots
      \\0 \rightarrow
        & B^i(K^\bullet)
        \xrightarrow{\varepsilon} B_\mathrm{I}^i(L^{\bullet,0})
        \rightarrow B_\mathrm{I}^i(L^{\bullet,1})
        \rightarrow \ldots
      \\0 \rightarrow
        & Z^i(K^\bullet)
        \xrightarrow{\varepsilon} Z_\mathrm{I}^i(L^{\bullet,0})
        \rightarrow Z_\mathrm{I}^i(L^{\bullet,1})
        \rightarrow \ldots
      \\0 \rightarrow
        & \HH^i(K^\bullet)
        \xrightarrow{\varepsilon} \HH_\mathrm{I}^i(L^{\bullet,0})
        \rightarrow \HH_\mathrm{I}^i(L^{\bullet,1})
        \rightarrow \ldots
      \end{aligned}
    \]
    são exatas ou, em outras palavras, $(L^{i,\bullet}), (B_\mathrm{I}^i(L^{\bullet\bullet})), (Z_\mathrm{I}^i(L^{\bullet\bullet}))$ e $(\HH_\mathrm{I}^i(L^{\bullet\bullet}))$ são \emph{resoluções} de $K^i, B^i(K^\bullet), Z^i(K^\bullet)$ e $\HH^i(K^\bullet)$, respectivamente.
  \item[(ii) ] Para cada $j$, o complexo simples $L^{\bullet, j}$ é \emph{cindido} ou, em outras palavras, as sequências exatas
    \[
      \label{0.11.4.2.1}
      0 \rightarrow B_\mathrm{I}^i(L^{\bullet,j}) \rightarrow Z_\mathrm{I}^i(L^{\bullet,j}) \rightarrow \HH_\mathrm{I}^i(L^{\bullet,j}) \rightarrow 0
      \tag{11.4.2.1}
    \]
    \[
      \label{0.11.4.2.2}
      0 \rightarrow Z_\mathrm{I}^i(L^{\bullet,j}) \rightarrow L^{ij} \rightarrow B_\mathrm{I}^{i+1}(L^{\bullet,j}) \rightarrow 0
      \tag{11.4.2.2}
    \]
    são \emph{cindidas}.
\end{enumerate}

É demonstrado (M,~XVII,~1.2) que, se todo objeto de $\cat{C}$ é subobjeto de um objeto injetivo, então todo complexo $K^\bullet$ de $\cat{C}$ admite uma resolução de Cartan--Eilenberg \emph{injetiva}, ou seja, formada por objetos injetivos $L^{ij}$ (a condição (ii) anterior implica então que os $(B_\mathrm{I}^i(L^{\bullet,j})), (Z_\mathrm{I}^i(L^{\bullet,j}))$ e $(\HH_\mathrm{I}^i(L^{\bullet,j}))$ também são formados por objetos injetivos).
Além disso, para todo morfismo $f: K^\bullet \rightarrow K'^\bullet$ de complexos de $\cat{C}$,
para toda resolução de Cartan--Eilenberg $L^{\bullet\bullet}$ de $K^\bullet$ e toda resolução de Cartan--Eilenberg \emph{injetiva} $L'^{\bullet\bullet}$ de $K'^\bullet$, existe um morfismo de bicomplexos $F: L^{\bullet\bullet} \rightarrow L'^{\bullet\bullet}$ compatível com $f$ e com as aumentações; e, se $f$ e $g$ são morfismos homotópicos de $K^\bullet$ em $K'^\bullet$, então os morfismos correspondentes de $L^{\bullet\bullet}$ em $L'^{\bullet\bullet}$ também são homotópicos \emph{(loc. cit.)}.

Quando $K^{\bullet}$ está \emph{limitado inferiormente } (resp. \emph{superiormente}), pode ser tomado $L^{\bullet\bullet}$ tal que $L^{ij}=0$ para $i<i_0$ (resp. $i>i_0$) se $K^i=0$ para $i<i_0$ (resp. $i>i_0$) (M,~XVII,~1.3).

Suponhamos, por outro lado, que existe um inteiro $n$ tal que todo objeto de $\cat{C}$ admita uma \emph{resolução injetiva de comprimento $\leq n$};
então pode-se supor que $L^{ij}=0$ para $j>n$ (M,~XVII,~1.4).
\end{env}


\begin{env}[11.4.3]
\label{0.11.4.3}
Seja agora $T$ um \emph{funtor covariante aditivo} de $\cat{C}$ em uma categoria abeliana $\cat{C'}$.
Dado um complexo $K^\bullet$ de $\cat{C}$ e uma resolução de Cartan--Eilenberg \emph{injetiva} $L^{\bullet\bullet}$ de $K^{\bullet}$, suponhamos que exista o complexo (simples) definido pelo bicomplexo $T(L^{\bullet\bullet})$ \sref{0.11.3.1};
então as duas sequências espectrais $'E(T(L^{\bullet\bullet}))$ e $''E(T(L^{\bullet\bullet}))$ deste bicomplexo são chamadas de \emph{sequências espectrais de hipercohomologia} de $T$ em relação a $K^\bullet$;
por \sref{0.11.4.2} e \sref{0.11.3.2}, dependem apenas do complexo $K^\bullet$ e não da escolha da resolução injetiva de Cartan--Eilenberg $L^{\bullet\bullet}$;
além disso, dependem \emph{funtorialmente} de $K^\bullet$.
Elas têm o mesmo termo limite $\HH^\bullet(T(L^{\bullet\bullet}))$, também chamado de \emph{hipercohomologia} de $T$ em relação a $K^\bullet$, e denotado por $\RR^\bullet T(K^\bullet)$.
Demonstra-se que os termos $E_2$ das duas sequências espectrais anteriores são dados por
\[
\label{0.11.4.3.1}
  'E_2^{pq}
  = \HH^p(\RR^q T(K^\bullet))
\tag{11.4.3.1}
\]
\[
\label{0.11.4.3.2}
  ''E_2^{pq}
  = \RR^p T(\HH^q(K^\bullet))
\tag{11.4.3.2}
\]
\oldpage[0\textsubscript{III}]{34}
onde $\RR^p T$ denota, como de costume, o \emph{funtor derivado} de grau $p$ de $T$ para $p\in\bb{Z}$;
$\RR^q T(K^\bullet)$ denota o complexo $(\RR^q T(K^i))_{i\in\bb{Z}}$.
\phantomsection\label{0.11.4.4}
Salvo menção explícita em contrário, vamos assumir que todo objeto de $\cat{C}$ é subobjeto de um objeto injetivo de $\cat{C}$, de modo que existem resoluções de Cartan--Eilenberg injetivas para todo complexo de $\cat{C}$.
Como $L^{ij}=0$ para $j<0$, os critérios de \sref{0.11.3.3} mostram que as \emph{duas} sequências espectrais de hipercohomologia de $T$ em relação a $K^\bullet$ existem e são \emph{birregulares} em cada um dos dois casos seguintes:
\begin{enumerate}
  \item $K^\bullet$ é limitado inferiormente;
  \item Todos os objetos de $\cat{C}$ admitem uma resolução injetiva de comprimento máximo $n$, para algum inteiro $n$ independente do objeto considerado.
\end{enumerate}
Com efeito, no primeiro caso, por \sref{0.11.4.2}, pode-se supor que exista $i_0$ tal que $L^{ij}=0$ para $i<i_0$ e, no segundo caso, que exista $j_1$ tal que $L^{ij}=0$ para $j>j_1$;
em cada um destes dois casos, é claro que, para cada $n$ dado, só existe um número finito de pares $(i,j)$ tais como $L^{ij}\neq 0$ e $i+j=n$, o que demonstra nossas afirmações.

Se assumirmos que os limites indutivos filtrantes existem em $\cat{C}'$ e são exatos (o que implica, em particular, a existência de somas diretas infinitas em $\cat{C}'$), então existe o complexo definido pelo bicomplexo $T(L^{\bullet\bullet})$, e o critério \sref{0.11.3.3} mostra que a sequência $'E(T(L^{\bullet\bullet}))$ é sempre \emph{regular}.

Se $K^\bullet$ é um complexo tal que todos os $K^i$ são nulos exceto um único $K^{i_0}$, então
$\RR^n T(K^\bullet)$ é isomorfo a $\RR^{n-i_0} T(K^{i_0})$, como se deduz imediatamente das definições tomando uma resolução de Cartan--Eilenberg $L^{\bullet\bullet}$ tal que $L^{ij}=0$ para $i\neq i_0$.

Se $K^\bullet$ e $K'^\bullet$ são dois complexos de $\cat{C}$, e $f$ e $g$ são morfismos homotópicos de $K^\bullet$ em $K'^\bullet$, então os morfismos $\RR^\bullet T(K^\bullet) \rightarrow \RR^\bullet T(K'^\bullet)$ induzidos por $f$ e $g$ são idênticos, e o mesmo acontece com os morfismos das sequências espectrais de hipercohomologia.
\end{env}

\begin{proposition}[11.4.5]
\label{0.11.4.5}
Suponhamos que os limites indutivos filtrantes existem em $\cat{C}'$ e são exatos.
Se $\RR^n T(K^i) = 0$ para todo $n>0$ e todo $i\in\bb{Z}$, então temos isomorfismos funtoriais
\[
\label{0.11.4.5.1}
  \RR^i T(K^\bullet) \simto \HH^i(T(K^\bullet))
\tag{11.4.5.1}
\]
para todo $i\in\bb{Z}$.
\end{proposition}

\begin{proof}
Os termos $E_2$ não nulos da primeira sequência espectral \sref{0.11.4.3.1} são então $'E_2^{p0} = \HH^p(T(K^\bullet))$;
Em outras palavras, esta sequência é \emph{degenerada};
Como é regular \sref{0.11.4.4}, a conclusão deduz-se de \sref{0.11.1.6}.
\end{proof}

\begin{env}[11.4.6]
\label{0.11.4.6}
Consideremos agora, por exemplo, um bifuntor covariante $(M, N) \rightarrow T(M, N)$ de $\cat{C} \times \cat{C}'$ em $\cat{C}''$, onde $\cat{C}$, $\cat{C}'$ e $\cat{C}''$ são três categorias abelianas;
Supomos, por simplicidade, que $T$ é aditivo em cada um dos seus argumentos e, além disso, que todo objeto de $\cat{C}$ e todo objeto de $\cat{C}'$ são subobjetos de um objeto injetivo,
e que os limites indutivos filtrantes existem em $\cat{C}''$ e são exatos.
Define-se então a \emph{hipercohomologia} de $T$ em relação aos complexos $K^\bullet$ e $K'^\bullet$ de $\cat{C}$ e $\cat{C}'$, respectivamente, com operadores diferenciais de grau~$+1$, substituindo $K^\bullet$ (resp. $K'^\bullet$) por uma resolução injetiva de Cartan--Eilenberg $L^{\bullet\bullet}$ (resp. $L'^{\bullet\bullet}$);
Então $T(L^{\bullet\bullet}, L'^{\bullet\bullet})$ é um quadricomplexo de $\cat{C}''$, que consideramos como um \emph{bicomplexo} de $\cat{C}''$, sendo os graus de $T(L^{ij}, L'^{hk})$ os inteiros $i+h$ e $j+k$.
A \emph{hipercohomologia} de $T$ em relação a $K^\bullet$ e $K'^\bullet$ é, por definição, a cohomologia $\HH^\bullet(T(L^{\bullet\bullet}, L'^{\bullet\bullet}))$ deste bicomplexo (em outras palavras, a do complexo simples associado) e denota-se por $\RR^\bullet T(K^\bullet, K'^\bullet)$;
\oldpage[0\textsubscript{III}]{35}
é o limite de duas sequências espectrais cujos termos $E_2$ são dados por
\[
\label{0.11.4.6.1}
  'E_2^{pq}
  = \HH^p(\RR^q T(K^\bullet, K'^\bullet))
\tag{11.4.6.1}
\]
\[
\label{0.11.4.6.2}
  ''E_2^{pq}
  = \sum_{q'+q''=q}\RR^p T(\HH^{q'}(K^\bullet),\HH^{q''}(K'^\bullet))
\tag{11.4.6.2}
\]
(cf.~M,~XVII,~2).

Aqui $\RR^q T(K^\bullet, K'^\bullet)$ é o bicomplexo $(\RR^q T(K^i, K'^j))_{(i,j)\in \bb{Z}\times \bb{Z}}$, e o segundo membro de \sref{0.11.4.6.1} é sua cohomologia quando ele é considerado como complexo simples.

Além disso, a primeira sequência espectral é sempre regular, e as duas sequências espectrais são birregulares quando existe $n$ tal que todo objeto de $\cat{C}$ e todo objeto de $\cat{C}'$ admitam uma resolução injetiva de comprimento $\leq n$, ou quando $K^\bullet$ e $K'^\bullet$ estão limitados inferiormente;
neste último caso, pode-se ainda omitir a hipótese de que existam os limites indutivos em $\cat{C}$ e $\cat{C}'$.

Se $K_1^\bullet$ e ${K'_1}^\bullet$ são outros dois complexos de $\cat{C}$ e $\cat{C}'$, respectivamente, $f$ e $g$ são morfismos \emph{homotópicos } de $K^\bullet$ em $K_1^\bullet$,
e $f'$ e $g'$ são morfismos \emph{homotópicos} de $K'^\bullet$ em ${K'_1}^\bullet$, então os morfismos $\RR^\bullet  T(K^\bullet, K'^\bullet) \rightarrow \RR^\bullet  T(K_1^\bullet, {K'_1}^\bullet)$ induzidos por $f$ e $f'$, por um lado, e por $g$ e $g'$, por outro, são idênticos, e o mesmo acontece com os morfismos das sequências espectrais de hipercohomologia.

Isso se generaliza facilmente a qualquer multifuntor covariante aditivo.
\end{env}

\begin{proposition}[11.4.7]
\label{0.11.4.7}
Suponhamos que, para todo objeto injetivo $\cat{I}$ de $\cat{C}$ (resp. $\cat{I'}$ de $\cat{C}'$), $\cat{A'}\mapsto T(\cat{I}, \cat{A'})$ (resp. $\cat{A}\mapsto T(\cat{A}, \cat{I'})$) é um funtor exato.
Então, com a notação de \sref{0.11.4.6}, temos um isomorfismo canônico
\[
  \RR^\bullet T(K^\bullet, K'^\bullet) \simto \HH^\bullet(T(L^{\bullet\bullet}, K'^\bullet)) \simto \HH^\bullet(T(K^\bullet, L'^{\bullet\bullet})) 
\]
onde os dois últimos termos são as cohomologias dos complexos simples definidos pelos tricomplexos $T(L^{\bullet\bullet}, K'^\bullet)$ e $T(K^\bullet, L'^{\bullet\bullet})$, respectivamente.
\end{proposition}

\begin{proof}
Definimos, por exemplo, o primeiro destes isomorfismos.
O quadricomplexo $T(L^{\bullet\bullet}, L'^{\bullet\bullet})$ pode ser considerado como um bicomplexo, onde os graus de $T(L^{ij}, L'^{hk})$ são os inteiros $i+j+h$ e $k$.
Como, para cada $h$, ${L'}^{h,\bullet}$ é uma \emph{resolução} de ${K'}^h$, tem-se, para este bicomplexo, sob as hipóteses sobre $T$, $\HH_\mathrm{II}^q(T(L^{\bullet\bullet}, L'^{\bullet\bullet}))=0$ para $q\neq 0$, e $\HH_\mathrm{II}^0(T(L^{\bullet\bullet}, L'^{\bullet\bullet})) = T(L^{\bullet\bullet}, K'^\bullet)$;
A primeira sequência espectral deste bicomplexo é, portanto, \emph{degenerada};
Como $L'^{hk}=0$ para $k<0$, esta sequência é também regular \sref{0.11.3.3}, e a conclusão deduz-se então de \sref{0.11.1.6}.
\end{proof}

Os resultados são análogos para um multifuntor covariante de qualquer número $n$ de argumentos: no cálculo da hipercohomologia não é necessário substituir \emph{todos} os complexos
por uma resolução de Cartan--Eilenberg, mas apenas $n-1$ deles, desde que, ao fixar $n-1$ argumentos arbitrários tomando-os como objetos \emph{injetivos}, o funtor covariante no argumento restante seja \emph{exato}.


\subsection{Passo para limite indutivo na hipercohomologia}
\label{subsection:0.11.5}

\begin{lemma}[11.5.1]
\label{0.11.5.1}
Seja $K^\bullet=(K^i)_{i\in\bb{Z}}$ um complexo de $\cat{C}$ e, para todo inteiro $r\in\bb{Z}$, seja $K^\bullet_{(r)}$ o complexo tal que $K^i_{(r)}=0$ para $i<r$ e $K^i_{(r)}=K^i$ para $i\geq r$.
Seja $T$ um funtor covariante aditivo
\oldpage[0\textsubscript{III}]{36}
de $\cat{C}$ em $\cat{C}'$ que comuta com os limites indutivos; supomos que os limites indutivos filtrantes existem e são exatos em $\cat{C}$ e $\cat{C}'$.
Então $\RR^\bullet T(K^\bullet)$ é isomorfo ao limite indutivo
\[
  \varinjlim_{r\to-\infty}\RR^\bullet T(K^\bullet_{(r)}).
\]
\end{lemma}

Uma resolução injetiva de Cartan--Eilenberg de $K^\bullet$ é construída escolhendo arbitrariamente, para cada $i$, uma resolução injetiva $(X_{\mathrm{B}}^{ij})_{j\geq 0}$ de $\mathrm{B}^i(K^\bullet)$ e uma resolução injetiva $(X_{\mathrm{H}}^{ij})_{j\geq 0}$ de $\HH^i(K^\bullet)$.
Feito isso, a construção mostra que a resolução injetiva $(L^{ij})_{j\geq 0}$ de $K^i$ e os operadores diferenciais $L^{i,j}\to L^{i+1,j}$ dependem apenas das resoluções $(X_{\mathrm{B}}^{i,\bullet})$, $(X_{\mathrm{H}}^{i,\bullet})$ e $(X_{\mathrm{B}}^{i+1,\bullet})$ (M,~XVII,~1.2).

Está claro agora que $\mathrm{B}^i(K^\bullet_{(r)})=\mathrm{B}^i(K^\bullet)$ e $\HH^i(K^\bullet_{(r)})=\HH^i(K^\bullet)$ para $i\geq r+1$.
Por outro lado, para todo $r$ há uma injeção canônica
\[
  \varphi^\bullet_{r-1,r}:K^\bullet_{(r)}\longrightarrow K^\bullet_{(r-1)},
\]
onde $\varphi^i_{r-1,r}$ é a identidade para $i\neq r-1$.
A observação anterior mostra que, se $L^{\bullet\bullet}=(L^{ij})$ é uma resolução injetiva de Cartan--Eilenberg de $K^\bullet$, então para todo $r$ pode-se definir uma resolução injetiva de Cartan--Eilenberg $L^{\bullet\bullet}_{(r)}=(L^{ij}_{(r)})$ de $K^\bullet_{(r)}$ tal que $L^{ij}_{(r)}=0$ para $i<r$ e $L^{ij}_{(r)}=L^{ij}$ para $i\geq r+1$.
Também pode ser definido um morfismo de bicomplexos
\[
  \Phi^{\bullet\bullet}_{r-1,r}:L^{\bullet\bullet}_{(r)}
  \longrightarrow L^{\bullet\bullet}_{(r-1)}
\]
correspondente a $\varphi^\bullet_{r-1,r}$; a construção deste morfismo (\emph{loc.\ cit.}) mostra novamente que pode ser escolhido de modo que $\Phi^{ij}_{r-1,r}$ seja a identidade para $i\neq r$ e $i\neq r-1$.

Assim, definimos um sistema indutivo $(L^{\bullet\bullet}_{(r)})$ de bicomplexos de $\cat{C}$ cujo limite indutivo, quando $r$ tende a $-\infty$, é evidentemente $L^{\bullet\bullet}$.
Como as somas diretas comutam com os limites indutivos, o complexo simples associado a $L^{\bullet\bullet}$ é igualmente o limite indutivo dos complexos simples associados a $L^{\bullet\bullet}_{(r)}$.
Por hipótese, $T$ comuta com os limites indutivos, assim como a cohomologia, porque o funtor $\varinjlim$ é exato.
Consequentemente, exceto isomorfismo,
\[
  \HH^\bullet(T(L^{\bullet\bullet}))
  =\varinjlim_{r\to-\infty}\HH^\bullet(T(L^{\bullet\bullet}_{(r)})).
\]

O Lema~\sref{0.11.5.1} permite, passando ao limite indutivo, estender a complexos arbitrários $K^\bullet$ propriedades que valem para complexos \emph{limitados inferiormente}.
Como primeiro exemplo, demonstramos o seguinte.

\begin{proposition}[11.5.2]
\label{0.11.5.2}
Sob as hipóteses de \sref{0.11.5.1} relativas a $\cat{C}$, $\cat{C}'$ e $T$, o funtor $\RR^\bullet T(K^\bullet)$ é um funtor cohomológico sobre a categoria abeliana dos complexos de $\cat{C}$.
\end{proposition}

\begin{proof}
Primeiro, mostremos que basta tratar os complexos limitados inferiormente.
Dada uma sequência exata
\[
  0\longrightarrow K'^\bullet\longrightarrow K^\bullet
  \longrightarrow K''^\bullet\longrightarrow 0
\]
de complexos, para todo $r$ induz uma sequência exata
\[
  0\longrightarrow K'^\bullet_{(r)}\longrightarrow K^\bullet_{(r)}
  \longrightarrow K''^\bullet_{(r)}\longrightarrow 0.
\]
Por hipótese, a última dá uma sequência exata
\[
  \cdots\longrightarrow\RR^nT(K'^\bullet_{(r)})
  \longrightarrow\RR^nT(K^\bullet_{(r)})
  \longrightarrow\RR^nT(K''^\bullet_{(r)})
  \xrightarrow{\partial}\RR^{n+1}T(K'^\bullet_{(r)})
  \longrightarrow\cdots .
\]
Estas sequências exatas formam um sistema indutivo.
O Lema~\sref{0.11.5.1} e a exatidão de $\varinjlim$ dão, portanto, uma sequência exata
\[
  \cdots\longrightarrow\RR^nT(K'^\bullet)
  \longrightarrow\RR^nT(K^\bullet)
  \longrightarrow\RR^nT(K''^\bullet)
  \xrightarrow{\partial}\RR^{n+1}T(K'^\bullet)
  \longrightarrow\cdots .
\]
\end{proof}

Para tratar os complexos limitados abaixo, basta considerar aqueles para os quais $K^i=0$ quando $i<0$; estes formam manifestamente uma categoria abeliana, que denotamos por $\cat{K}$.

\begin{lemma}[11.5.2.1]
\label{0.11.5.2.1}
Em $\cat{K}$, seja $\mathfrak{J}$ a classe dos complexos $Q^\bullet=(Q^i)_{i\geq0}$ que possuem as
\oldpage[0\textsubscript{III}]{37}
propriedades seguintes:
\begin{enumerate}
  \item[{\rm(1) }] todo $Q^i$ é um objeto injetivo de $\cat{C}$;
  \item[{\rm(2) }] para todo $i\geq0$, $\mathrm{Z}^i(Q^\bullet)=\mathrm{B}^i(Q^\bullet)$, e $\mathrm{Z}^i(Q^\bullet)$ é um somando direto de $Q^i$.
\end{enumerate}
Então:
\begin{enumerate}
  \item[{\rm(i) }] todo $Q^\bullet\in\mathfrak{J}$ é um objeto injetivo de $\cat{K}$;
  \item[{\rm(ii) }] todo objeto de $\cat{K}$ é isomorfo a um subcomplexo de um complexo pertencente a $\mathfrak{J}$.
\end{enumerate}
\end{lemma}

\begin{proof}
{\rm(i)}
Seja $A^\bullet=(A^i)$ um objeto de $\cat{K}$, seja $A'^\bullet=(A'^i)$ um subobjeto de $A^\bullet$, seja $Q^\bullet=(Q^i)$ pertencente a $\mathfrak{J}$, e suponhamos dado um morfismo
\[
  f=(f^i):A'^\bullet\longrightarrow Q^\bullet
\]
.
Temos que estendê-lo a um morfismo $g=(g^i):A^\bullet\to Q^\bullet$.
Por simplicidade, usamos a linguagem da categoria de módulos (cf.~[27]).

Identificemos $Q^i$ com $\mathrm{B}^i(Q^\bullet)\oplus\mathrm{B}^{i+1}(Q^\bullet)$.
Procedamos por indução sobre $i$.
Suponhamos que $g^j$ foi definido para $j<i$, compatível com os diferenciais $d^j:A^j\to A^{j+1}$ e $d'^j:Q^j\to Q^{j+1}$ para $j<i-1$, e que além disso:
\begin{enumerate}
  \item[{\rm(1)}] $g^{i-1}(\mathrm{Z}^{i-1}(A^\bullet))\subset\mathrm{Z}^{i-1}(Q^\bullet)$;
  \item[{\rm(2) }] se $C^j=(d^j)^{-1}(A'^{j+1})$ para todo $j$, então $d'^{i-1}\circ g^{i-1}$ coincide com $f^i\circ d^{i-1}$ sobre $C^{i-1}$.
\end{enumerate}
Compondo $f^i:A'^i\to Q^i$ com as duas projeções, obtemos morfismos
\[
  f'^i:A'^i\longrightarrow\mathrm{B}^i(Q^\bullet),
  \qquad
  f''^i:A'^i\longrightarrow\mathrm{B}^{i+1}(Q^\bullet).
\]
Como $d'^{i-1}\circ g^{i-1}$ envia $A^{i-1}$ em $\mathrm{B}^i(Q^\bullet)$ e se anula sobre $\mathrm{Z}^{i-1}(A^\bullet)$, define um morfismo
\[
  h^i:\mathrm{B}^i(A^\bullet)\longrightarrow\mathrm{B}^i(Q^\bullet).
\]
A condição ~{\rm(2)} diz que $h^i$ coincide com $f'^i$ sobre $\mathrm{B}^i(A^\bullet)\cap A'^i$.
Como $\mathrm{B}^i(Q^\bullet)$ é um somando direto de $Q^i$, é injetivo; existe, portanto, um morfismo
\[
  g'^i:A^i\longrightarrow\mathrm{B}^i(Q^\bullet)
\]
que coincide com $h^i$ sobre $\mathrm{B}^i(A^\bullet)$ e com $f'^i$ sobre $A'^i$.

Por outro lado, consideremos o morfismo
\[
  f''^{i+1}\circ d^i:C^i\longrightarrow\mathrm{B}^{i+1}(Q^\bullet),
\]
que se anula sobre $\mathrm{Z}^i(A^\bullet)$.
Como $\mathrm{B}^{i+1}(Q^\bullet)$ é injetivo, existe um morfismo
\[
  g''^i:A^i\longrightarrow\mathrm{B}^{i+1}(Q^\bullet)
\]
que corresponde a $f''^{i+1}\circ d^i$ sobre $C^i$ e a zero sobre $\mathrm{Z}^i(A^\bullet)$.
Tomando $g^i=g'^i+g''^i$, pode-se continuar a indução.

\smallskip
\noindent{\rm(ii)}
Para mergulhar $A^\bullet=(A^i)$ em um complexo pertencente a $\mathfrak{J}$, escolhemos para todo $i\geq1$ um objeto injetivo $Q'^i$ de $\cat{C}$ e uma injeção $f'^i:A^i\to Q'^i$.
Ponhamos
\[
  Q^i=0\quad(i<0),\qquad Q^0=Q'^1,\qquad
  Q^i=Q'^i\oplus Q'^{i+1}\quad(i\geq1),
\]
com o diferencial evidente.
Para todo $i$, vamos colocar
\[
  f^i=f'^i+(f'^{i+1}\circ d^i),
\]
onde $f'^i=0$ para $i\leq0$.
É imediato que todo $f^i$ é injetivo e que os $f^i$ comutam com os diferenciais.
\end{proof}

\begin{corollary}[11.5.2.2]
\label{0.11.5.2.2}
Todo objeto $K^\bullet$ de $\cat{K}$ admite uma resolução à direita formada por objetos de $\mathfrak{J}$.
Se $L^{\bullet\bullet}$ é tal resolução, então para toda resolução $L'^{\bullet\bullet}$ de $K^\bullet$ formada por objetos de $\cat{K}$ existe um morfismo de bicomplexos
\[
  F:L'^{\bullet\bullet}\longrightarrow L^{\bullet\bullet}
\]
compatível com as aumentações, e quaisquer dois morfismos $F,F'$ deste tipo são homotópicos.
\end{corollary}

\begin{proof}
Isto é precisamente (M, V, 1.1~a)) aplicado à categoria abeliana $\cat{K}$.
\end{proof}

\begin{env}[11.5.2.3]
\label{0.11.5.2.3}
Estabelecidos estes preliminares, seja $L'^{\bullet\bullet}$ uma resolução de Cartan--Eilenberg injetiva de $K^\bullet$ e seja $L^{\bullet\bullet}$ uma resolução de $K^\bullet$ formada por objetos de $\mathfrak{J}$.
Mostremos que existe um isomorfismo
\[
  \HH^\bullet(T(L'^{\bullet\bullet}))
  \simto
  \HH^\bullet(T(L^{\bullet\bullet})).
\]
O Corolário~\sref{0.11.5.2.2} dá um morfismo de bicomplexos
\[
  T(L'^{\bullet\bullet})\longrightarrow T(L^{\bullet\bullet}),
\]
e, portanto, um morfismo
\[
  'E(T(L'^{\bullet\bullet}))\longrightarrow{}'E(T(L^{\bullet\bullet}))
\]
entre suas primeiras sequências espectrais.
Por \sref{0.11.3.3}, estas sequências espectrais são regulares, então por \sref{0.11.1.5} basta verificar que este morfismo é um isomorfismo sobre os termos $E_2$.
Equivalentemente, é preciso demonstrar que
\[
  \HH_\mathrm{II}^q(T(L^{i,\bullet}))=\RR^qT(K^i).
\]
Como $L^{i,\bullet}$ é uma resolução à direita de $K^i$, isso reduz a questão ao seguinte lema.
\end{env}

\oldpage[0\textsubscript{III}]{38}
\begin{lemma}[11.5.2.4]
\label{0.11.5.2.4}
Seja $L^\bullet=(L^i)_{i\in\bb{Z}}$ uma resolução à direita de um objeto $A$ de $\cat{C}$ tal que $\RR^nT(L^i)=0$ para todo $i\in\bb{Z}$ e todo $n>0$.
Então
\[
  \HH^\bullet(T(L^\bullet))=\RR^\bullet T(A).
\]
\end{lemma}

\begin{proof}
Este é um caso particular de (T, 2.5.2).
\end{proof}

\begin{env}[11.5.2.5]
\label{0.11.5.2.5}
A demonstração de \sref{0.11.5.2} deduz-se agora imediatamente, já que $L'^{\bullet\bullet}$ é uma resolução injetiva de $K^\bullet$ na categoria abeliana $\cat{K}$.
Em outras palavras, o funtor
\[
  K^\bullet\longmapsto\HH^\bullet(T(L'^{\bullet\bullet}))
\]
é precisamente o funtor cohomológico derivado pela direita de $T$ na categoria $\cat{K}$ (T, 2.3).
\end{env}

\begin{proposition}[11.5.3]
\label{0.11.5.3}
Sob as hipóteses de \sref{0.11.5.1} relativas a $\cat{C}$, $\cat{C}'$ e $T$, seja $L^{\bullet\bullet}=(L^{ij})$ um bicomplexo tal que $L^{ij}=0$ para $j<0$ e tal que, para todo $i$, $L^{i,\bullet}$ é uma resolução de $K^i$.
Suponhamos ainda que $\RR^nT(L^{ij})=0$ para todo par $(i,j)$ e todo $n>0$.
Então há um isomorfismo funtorial
\[
\label{0.11.5.3.1}
\tag{11.5.3.1}
  \RR^\bullet T(K^\bullet)
  \simto
  \HH^\bullet(T(L^{\bullet\bullet})).
\]
\end{proposition}

\begin{proof}
Seja $L^{\bullet\bullet}_{(r)}=(L^{ij}_{(r)})$ o bicomplexo definido por
\[
  L^{ij}_{(r)}=0\quad(i<r),\qquad
  L^{ij}_{(r)}=L^{ij}\quad(i\geq r).
\]
É imediato que $L^{\bullet\bullet}$ é o limite indutivo dos $L^{\bullet\bullet}_{(r)}$ quando $r$ tende a $-\infty$.
Pela hipótese sobre $T$ e por \sref{0.11.5.1}, basta, portanto, demonstrar a proposição quando $K^\bullet$ está limitado inferiormente, por exemplo quando $K^i=0$ para $i<0$, com $L^{ij}=0$ para $i<0$.
Seja $L'^{\bullet\bullet}=(L'^{ij})$ uma resolução à direita de $K^\bullet$ formada por objetos de $\mathfrak{J}$, como em \sref{0.11.5.2.2}.
Existe um morfismo de bicomplexos
\[
  L^{\bullet\bullet}\longrightarrow L'^{\bullet\bullet}
\]
compatível com as aumentações e, portanto, um morfismo
\[
  'E(T(L^{\bullet\bullet}))\longrightarrow{}'E(T(L'^{\bullet\bullet}))
\]
entre as primeiras sequências espectrais.
O Lema~\sref{0.11.5.2.4} mostra, exatamente como em \sref{0.11.5.2.3}, que este morfismo é um isomorfismo, o que demonstra a afirmação.
\end{proof}

\begin{remark}[11.5.4]
\label{0.11.5.4}
Os argumentos anteriores mostram que as conclusões de \sref{0.11.5.2} e \sref{0.11.5.3} se verificam na categoria dos complexos $K^\bullet$ limitados inferiormente sem supor que $T$ comuta com os limites indutivos filtrantes.
Além disso, se considerarmos apenas a categoria $\cat{K}$ dos complexos $K^\bullet$ tais que $K^i=0$ para $i<0$, a caracterização de $\RR^\bullet T(K^\bullet)$ como o sistema de funtores derivados pela direita de $T$ em $\cat{K}$ mostra que esse funtor cohomológico é \emph{universal} (T, 2.3).

Outro caso em que \sref{0.11.5.2} é cumprido sem qualquer hipótese adicional sobre $T$ é aquele em que existe um inteiro $m>0$ tal que todo objeto de $\cat{C}$ admite uma resolução injetiva de comprimento máximo $m$.
Na demonstração de \sref{0.11.5.1}, todas as resoluções injetivas de objetos de $\cat{C}$ que aparecem podem ser escolhidas de comprimento no máximo $m$.
Deduz-se imediatamente que, para um grau total fixo $n$ e para um $r$ suficientemente grande, os termos de grau total $n$ em $T(L^{\bullet\bullet}_{(r)})$ são iguais aos de $T(L^{\bullet\bullet})$ e são finitos em número.
Consequentemente,
\[
  \HH^n(T(L^{\bullet\bullet}))
  =\HH^n(T(L^{\bullet\bullet}_{(r)}))
\]
para $r$ suficientemente grande.
Com a notação de \sref{0.11.5.2}, temos analogamente
\[
  \RR^nT(K^\bullet_{(r)})=\RR^nT(K^\bullet)
\]
para $r$ suficientemente grande (dependendo de $n$), e de forma análoga para $K'^\bullet$ e $K''^\bullet$, o que demonstra a conclusão.
Do mesmo modo, \sref{0.11.5.3} vale sem condição adicional sobre $T$ quando $\cat{C}$ satisfaz a hipótese anterior e se tomam as resoluções $L^{i,\bullet}$ de comprimento no máximo $m$.
\end{remark}

\begin{env}[11.5.5]
\label{0.11.5.5}
O resultado de \sref{0.11.5.2} estende-se aos multifuntores covariantes.
Consideremos, por exemplo, a situação de \sref{0.11.4.6}, em que se supõe que os limites indutivos filtrantes existem e são exatos em $\cat{C}$, $\cat{C}'$ e $\cat{C}''$,
\oldpage[0\textsubscript{III}]{39}
e supõe-se que $T$ comuta com os limites indutivos.
Então $\RR^\bullet T(K^\bullet,K'^\bullet)$ é um bifuntor cohomológico em cada um dos complexos $K^\bullet,K'^\bullet$.
Para vê-lo, reduz-se, como em \sref{0.11.5.2}, ao caso em que $K^\bullet$ e $K'^\bullet$ estão limitados inferiormente.
Escolhendo resoluções injetivas de $K^\bullet$ e $K'^\bullet$ do tipo descrito em \sref{0.11.5.2.2}, ficamos reduzidos à propriedade geral demonstrada em (M, V, 4.1).
\end{env}

\begin{env}[11.5.6]
\label{0.11.5.6}
De modo análogo, os resultados de \sref{0.11.4.7} e \sref{0.11.5.3} estendem-se da seguinte forma.
Sob as hipóteses de \sref{0.11.5.5}, suponhamos dados dois bicomplexos
\[
  L^{\bullet\bullet}=(L^{ij}),\qquad
  L'^{\bullet\bullet}=(L'^{ij})
\]
tais que $L^{ij}=L'^{ij}=0$ para $j<0$, tais que, para todo $i$, $L^{i,\bullet}$ é uma resolução de $K^i$ e $L'^{i,\bullet}$ é uma resolução de $K'^i$, e tais que
\[
  \RR^nT(L^{ij},L'^{hk})=0
\]
Para $n>0$ e todo sistema de índices $(i,j,h,k)$.
Então há um isomorfismo, funtorial em $K^\bullet$ e $K'^\bullet$,
\[
\label{0.11.5.6.1}
\tag{11.5.6.1}
  \RR^\bullet T(K^\bullet,K'^\bullet)
  \simto
  \HH^\bullet(T(L^{\bullet\bullet},L'^{\bullet\bullet})).
\]
Isto é demonstrado como em \sref{0.11.5.3}, reduzindo-se ao caso em que $K^\bullet$ e $K'^\bullet$ estão limitados inferiormente.

Suponhamos ainda que, para todo par $(i,j)$ e todo par $(h,k)$, os funtores
\[
  A\longmapsto T(A,L'^{hk}),
  \qquad
  A'\longmapsto T(L^{ij},A')
\]
são exatos em $\cat{C}$ e $\cat{C}'$, respectivamente.
Então existem também isomorfismos funtoriais
\[
\label{0.11.5.6.2}
\tag{11.5.6.2}
  \RR^\bullet T(K^\bullet,K'^\bullet)
  \simto
  \HH^\bullet(T(L^{\bullet\bullet},K'^\bullet))
  \simto
  \HH^\bullet(T(K^\bullet,L'^{\bullet\bullet})).
\]
A demonstração é análoga à de \sref{0.11.4.7}.
\end{env}

\begin{env}[11.5.7]
\label{0.11.5.7}
Os resultados de \sref{0.11.5.5} e \sref{0.11.5.6} também valem sem supor que $T$ comute com os limites indutivos, contanto que nos limitemos a complexos $K^\bullet,K'^\bullet$ limitados inferiormente, ou que se suponha que todo objeto de $\cat{C}$ (respectivamente, de $\cat{C}'$) admita uma resolução injetiva de comprimento limitado e, em \sref{0.11.5.6}, que os segundos graus de $L^{\bullet\bullet}$ e $L'^{\bullet\bullet}$ sejam limitados superiormente.
\end{env}

\subsection{Hiperhomologia de um funtor em relação a um complexo $K_{\bullet}$}
\label{subsection:0.11.6}

\begin{env}[11.6.1]
\label{0.11.6.1}
Seja $\cat{C}$ uma categoria abeliana e seja $K_\bullet=(K_i)_{i\in\bb{Z}}$ um complexo de objetos de $\cat{C}$ cujo diferencial é de grau $-1$.
Uma \emph{resolução de Cartan--Eilenberg pela esquerda} de $K_\bullet$ consiste em um bicomplexo $L_{\bullet\bullet}=(L_{ij})$, cujos diferenciais são de grau $-1$ e para o qual $L_{ij}=0$ quando $j<0$, juntamente com um morfismo de complexos simples
\[
  \varepsilon:L_{\bullet,0}\longrightarrow K_\bullet,
\]
de modo que sejam satisfeitas as condições obtidas das de \sref{0.11.4.2} ``invertendo as setas''.
Se todo objeto de $\cat{C}$ é quociente de um objeto projetivo, então todo complexo $K_\bullet$ de $\cat{C}$ admite uma \emph{resolução de Cartan--Eilenberg projetiva}, isto é, formada por objetos projetivos $L_{ij}$; ela possui propriedades funtoriais análogas às de \sref{0.11.4.2}.
Além disso, se $K_\bullet$ for limitado inferiormente (respectivamente, superiormente), pode ser exigido $L_{ij}=0$ para $i<i_0$ (respectivamente, $i>i_0$) sempre que $K_i=0$ para $i<i_0$ (respectivamente, $i>i_0$).
Se todo objeto de $\cat{C}$ admite uma resolução projetiva de comprimento no máximo $n$, pode-se exigir $L_{ij}=0$ para $j>n$.
\end{env}

\begin{env}[11.6.2]
\label{0.11.6.2}
\oldpage[0\textsubscript{III}]{40}
Suponhamos que $T$ seja um funtor covariante aditivo de $\cat{C}$ em uma categoria abeliana $\cat{C}'$.
A \emph{hiperhomologia} $\LL_\bullet T(K_\bullet)$ e as \emph{sequências espectrais de hiperhomologia} de $T$ em relação a um complexo $K_\bullet$ de $\cat{C}$ (quando existem) são definidas a partir de \sref{0.11.4.3}, novamente ``invertendo as flechas''.
Os termos $E^2$ das duas sequências espectrais assim obtidos são
\[
\label{0.11.6.2.1}
\tag{11.6.2.1}
  {}'E^2_{pq}
  =\HH_p(\LL_qT(K_\bullet))
\]
e
\[
\label{0.11.6.2.2}
\tag{11.6.2.2}
  {}''E^2_{pq}
  =\LL_pT(\HH_q(K_\bullet)),
\]
onde $\LL_pT$ denota o funtor derivado pela esquerda do grau $p$ de $T$ para $p\geq0$, e zero para $p<0$, enquanto $\LL_qT(K_\bullet)$ denota o complexo $(\LL_qT(K_i))_{i\in\bb{Z}}$.

As propriedades da hiperhomologia não se deduzem todas das da hipercohomologia apenas ``invertendo as flechas'' (a menos que se imponham a $\cat{C}'$ hipóteses adicionais de tipo (AB~5*) de (T, 1.5)).
A razão são as condições de regularidade das duas sequências espectrais anteriores, às quais agora devem ser aplicados os critérios de \sref{0.11.3.4}.
Estes critérios mostram que, se os limites indutivos filtrantes existem e são exatos em $\cat{C}'$, então existe o complexo definido pelo bicomplexo $T(L_{\bullet\bullet})$ e a \emph{segunda} sequência espectral ${}''E(T(L_{\bullet\bullet}))$ é regular.
Se $K_\bullet$ estiver limitado inferiormente ou se existir um inteiro $n$ tal que todo objeto de $\cat{C}$ admita uma resolução projetiva de comprimento no máximo $n$, então as \emph{duas} sequências espectrais de hiperhomologia existem (sem hipótese alguma sobre $\cat{C}'$) e são birregulares.
\end{env}

\begin{proposition}[11.6.3]
\label{0.11.6.3}
Sejam $\cat{C}$ e $\cat{C}'$ duas categorias abelianas, e seja $T$ um funtor covariante aditivo de $\cat{C}$ em $\cat{C}'$.
Então:
\begin{enumerate}
  \item[{\rm(i) }] \emph{A hiperhomologia $\LL_\bullet T(K_\bullet)$ é um funtor homológico sobre a categoria abeliana dos complexos de $\cat{C}$ limitados inferiormente.}

  \item[{\rm(ii) }] \emph{Seja $K_\bullet$ um complexo de $\cat{C}$ limitado abaixo.
  Se $\LL_nT(K_i)=0$ para todo $n>0$ e todo $i\in\bb{Z}$, existem isomorfismos funtoriais}
  \[
  \label{0.11.6.3.1}
  \tag{11.6.3.1}
    \LL_iT(K_\bullet)
    \simto
    \HH_i(T(K_\bullet))
    \qquad(i\in\bb{Z}).
  \]

  \item[{\rm(iii) }] \emph{Seja $K_\bullet$ um complexo de $\cat{C}$ limitado inferiormente.
  Seja $L_{\bullet\bullet}=(L_{ij})$ um bicomplexo tal que $L_{ij}=0$ para $j<0$ e, para todo $i$, $L_{i,\bullet}$ é uma resolução de $K_i$.
  Suponhamos ainda que $\LL_nT(L_{ij})=0$ para todo par $(i,j)$ e todo $n>0$.
  Então há um isomorfismo funtorial}
  \[
  \label{0.11.6.3.2}
  \tag{11.6.3.2}
    \LL_\bullet T(K_\bullet)
    \simto
    \HH_\bullet(T(L_{\bullet\bullet})).
  \]
\end{enumerate}

As demonstrações procedem como as de \sref{0.11.5.2}, \sref{0.11.4.5} e \sref{0.11.5.3} no caso de complexos limitados inferiormente.
Deixamos ao leitor os detalhes destes raciocínios.
\end{proposition}

\begin{env}[11.6.4]
\label{0.11.6.4}
Os resultados são inteiramente análogos para os multifuntores covariantes aditivos em cada argumento.
Por exemplo, para um bifuntor $T$, as duas sequências espectrais de hiperhomologia têm termos $E^2$
\oldpage[0\textsubscript{III}]{41}
\[
\label{0.11.6.4.1}
\tag{11.6.4.1}
  {}'E^2_{pq}
  =\HH_p(\LL_qT(K_\bullet,K'_\bullet))
\]
e
\[
\label{0.11.6.4.2}
\tag{11.6.4.2}
  {}''E^2_{pq}
  =
  \sum_{q'+q''=q}
  \LL_pT(\HH_{q'}(K_\bullet),\HH_{q''}(K'_\bullet)).
\]
Também aqui é a \emph{segunda} sequência espectral que é regular.
As duas sequências são birregulares quando $K_\bullet$ e $K'_\bullet$ estão limitados inferiormente, ou quando os objetos das categorias abelianas consideradas admitem resoluções projetivas de comprimento limitado.

Além disso:
\end{env}

\begin{proposition}[11.6.5]
\label{0.11.6.5}
Sejam $\cat{C}$, $\cat{C}'$ e $\cat{C}''$ três categorias abelianas, e seja $T$ um bifuntor covariante biaditivo de $\cat{C}\times\cat{C}'$ em $\cat{C}''$.
\begin{enumerate}
  \item[{\rm(i) }] \emph{$\LL_\bullet T(K_\bullet,K'_\bullet)$ é um bifuntor homólogo em cada um dos complexos limitados inferiormente $K_\bullet,K'_\bullet$, formados respectivamente por objetos de $\cat{C}$ e de $\cat{C}'$.}

  \item[{\rm(ii) }] \emph{Suponha que $K_\bullet$ e $K'_\bullet$ estão limitados abaixo.
  Sejam $L_{\bullet\bullet}=(L_{ij})$ e $L'_{\bullet\bullet}=(L'_{ij})$ dois bicomplexos tais que $L_{ij}=L'_{ij}=0$ para $j<0$ e tais que, para todo $i$, $L_{i,\bullet}$ é uma resolução de $K_i$ e $L'_{i,\bullet}$ é uma resolução de $K'_i$.
  Vamos finalmente supor que $\LL_nT(L_{ij},L'_{hk})=0$ para todo $n>0$ e todo sistema de índices $(i,j,h,k)$.
  Então há um isomorfismo funtorial}
  \[
  \label{0.11.6.5.1}
  \tag{11.6.5.1}
    \LL_\bullet T(K_\bullet,K'_\bullet)
    \simto
    \HH_\bullet(T(L_{\bullet\bullet},L'_{\bullet\bullet})).
  \]

  \item[{\rm(iii) }] \emph{Suponhamos ainda que, para todo par $(i,j)$ e todo par $(h,k)$, os funtores
  \[
    A\longmapsto T(A,L'_{hk}),
    \qquad
    A'\longmapsto T(L_{ij},A')
  \]
  são exatos em $\cat{C}$ e $\cat{C}'$, respectivamente.
  Então existem isomorfismos funtoriais}
  \[
  \label{0.11.6.5.2}
  \tag{11.6.5.2}
    \LL_\bullet T(K_\bullet,K'_\bullet)
    \simto
    \HH_\bullet(T(L_{\bullet\bullet},K'_\bullet))
    \simto
    \HH_\bullet(T(K_\bullet,L'_{\bullet\bullet})).
  \]
\end{enumerate}

As demonstrações são análogas às de \sref{0.11.5.5} e \sref{0.11.5.6}.
\end{proposition}

\subsection{Hiperhomologia de um funtor em relação a um bicomplexo $K_{\bullet\bullet}$}
\label{subsection:0.11.7}

\begin{env}[11.7.1]
\label{0.11.7.1}
Seja $\cat{C}$ uma categoria abeliana em que todo objeto é quociente de um
objeto projetivo.
Consideremos um bicomplexo $K_{\bullet\bullet}=(K_{ij})$ de objetos de $\cat{C}$ cujos
ambos os graus estão limitados inferiormente.
Podemos sempre nos limitar ao caso em que $K_{ij}=0$ para $i<0$ ou
$j<0$, e assim faremos no futuro.
Podemos considerar $K_{\bullet\bullet}$ como um complexo ordinário formado por
objetos
\[
  K_{i,\bullet}=(K_{ij})_{j\geq 0}
\]
da categoria abeliana $\cat{K}$ dos complexos não negativos de objetos de
$\cat{C}$.
Do lema~\sref{0.11.5.2.1} (ou, mais precisamente, do lema
``dual'' obtido ``invertendo as flechas'') e de (M,~V,~2.2) deduz-se
que $K_{\bullet\bullet}$ admite uma \emph{resolução de Cartan--Eilenberg
projetiva} na categoria $\cat{K}$.
Tal resolução é um tricomplexo
$M_{\bullet\bullet\bullet}=(M_{ijk})$ em $\cat{C}$, cujos graus são todos
não negativos e cujos objetos são projetivos, tal que, para todo $i$,
\[
  M_{i,\bullet\bullet},\qquad
  \mathrm{B}^{\mathrm{I}}_i(M_{\bullet\bullet\bullet}),\qquad
  \mathrm{Z}^{\mathrm{I}}_i(M_{\bullet\bullet\bullet}),\qquad
  \HH^{\mathrm{I}}_i(M_{\bullet\bullet\bullet})
\]
são resoluções projetivas em $\cat{K}$, respectivamente, de
\[
  K_{i,\bullet},\qquad
  \mathrm{B}^{\mathrm{I}}_i(K_{\bullet\bullet}),\qquad
  \mathrm{Z}^{\mathrm{I}}_i(K_{\bullet\bullet}),\qquad
  \HH^{\mathrm{I}}_i(K_{\bullet\bullet}).
\]
Em particular, para todo par $(i,j)$, $M_{ij,\bullet}$ é uma resolução
projetiva de $K_{ij}$ em $\cat{C}$.
\end{env}

\begin{proposition}[11.7.2]
\label{0.11.7.2}
Seja $T$ um funtor covariante aditivo de $\cat{C}$ em uma categoria
abeliana $\cat{C}'$.
Com a notação de \sref{0.11.7.1}, a homologia
\[
  \HH_\bullet\bigl(T(M_{\bullet\bullet\bullet})\bigr)
\]
do complexo simples associado ao tricomplexo
$T(M_{\bullet\bullet\bullet})$ é canonicamente isomorfa à
hiperhomologia $\LL_\bullet T(K_{\bullet\bullet})$ do
\oldpage[0\textsubscript{III}]{42}
complexo simples associado a $K_{\bullet\bullet}$
\sref{0.11.6.2}.
É também o termo limite comum de seis sequências espectrais birregulares denotadas
${}^{(t)}\mathrm{E}$, onde $t=a,b,a',b',c$ ou $d$, cujos
termos $\mathrm{E}^2$ são
\[
\begin{aligned}
  {}^{(a)}\mathrm{E}^2_{pq}
    &= \LL_pT\bigl(\HH_q(K_{\bullet\bullet})\bigr),\\
  {}^{(b)}\mathrm{E}^2_{pq}
    &= \HH_p\bigl(\LL^{\mathrm{II}}_qT(K_{\bullet\bullet})\bigr),\\
  {}^{(a')}\mathrm{E}^2_{pq}
    &= \LL_pT\bigl(\HH^{\mathrm{I}}_q(K_{\bullet\bullet})\bigr),\\
  {}^{(b')}\mathrm{E}^2_{pq}
    &= \HH_p\bigl(\LL_qT(K_{\bullet\bullet})\bigr),\\
  {}^{(c)}\mathrm{E}^2_{pq}
    &= \LL_pT\bigl(\HH^{\mathrm{II}}_q(K_{\bullet\bullet})\bigr),\\
  {}^{(d)}\mathrm{E}^2_{pq}
    &= \HH_p\bigl(\LL^{\mathrm{I}}_qT(K_{\bullet\bullet})\bigr).
\end{aligned}
\]
\end{proposition}

Lembre-se que, para um complexo $A_\bullet=(A_i)$, usamos a notação
$F(A_\bullet)$ para o complexo cujos objetos são os $F(A_i)$.
Por exemplo, $\LL^{\mathrm{II}}_qT(K_{\bullet\bullet})$ denota o complexo
\[
  \bigl(\LL^{\mathrm{II}}_qT(K_{i,\bullet})\bigr)_{i\geq 0},
\]
onde $\LL^{\mathrm{II}}_qT(K_{i,\bullet})$ é a hiperhomologia de índice
$q$ da hiperhomologia do funtor $T$ relativa ao complexo simples $K_{i,\bullet}$.

\begin{proof}
Seja $L_\bullet$ o complexo simples associado a
$K_{\bullet\bullet}$, de modo que
\[
  L_i=\bigoplus_{r+s=i}K_{rs},
\]
e vamos colocar
\[
  N_{ij}=\bigoplus_{r+s=i}M_{rsj}.
\]
É claro que, para todo $i$, $N_{i,\bullet}$ é uma resolução projetiva
de $L_i$ em $\cat{C}$.
De \sref{0.11.6.3} e \sref{0.11.6.4} deduz-se que existe um
isomorfismo funtorial
\[
  \LL_\bullet T(L_\bullet)
  \isoto
  \HH_\bullet\bigl(T(N_{\bullet\bullet})\bigr).
\]
Como o complexo simples associado ao bicomplexo
$T(N_{\bullet\bullet})$ é também o complexo simples associado ao
tricomplexo $T(M_{\bullet\bullet\bullet})$, isto demonstra a primeira afirmação.

Além disso, $\LL_\bullet T(L_\bullet)$ é o termo limite das duas
sequências espectrais de hiperhomologia \sref{0.11.6.2} de $T$ relativas ao
complexo simples $L_\bullet$.
Estas são precisamente ${}^{(b')}\mathrm{E}$ e ${}^{(a)}\mathrm{E}$,
respectivamente.

Consideremos agora $M_{\bullet\bullet\bullet}$ como um bicomplexo
$U_{\bullet\bullet}$, onde
\[
  U_{ij}=\bigoplus_{r+s=j}M_{irs}.
\]
A homologia $\HH_\bullet(T(M_{\bullet\bullet\bullet}))$ é também o
termo limite das duas sequências espectrais do bicomplexo
$T(U_{\bullet\bullet})$.
Para todo $i$, o bicomplexo $M_{i,\bullet\bullet}$ satisfaz as
hipóteses de \sref{0.11.6.3} relativas ao complexo simples
$K_{i,\bullet}$.
Consequentemente,
\[
  \HH^{\mathrm{II}}_q\bigl(T(U_{i,\bullet})\bigr)
  =
  \LL^{\mathrm{II}}_qT(K_{i,\bullet}),
\]
e a primeira sequência espectral de $T(U_{\bullet\bullet})$ é precisamente
${}^{(b)}\mathrm{E}$.

Por outro lado, para todo $r$, $M_{\bullet,r,\bullet}$ é uma
resolução de Cartan--Eilenberg do complexo simples $K_{\bullet,r}$.
O cálculo realizado em (M,~XV,~2) mostra que
\[
  \HH^{\mathrm{I}}_q\bigl(T(M_{\bullet,rs})\bigr)
  =
  T\bigl(\HH^{\mathrm{I}}_q(M_{\bullet,rs})\bigr),
\]
e, portanto,
\[
  \HH^{\mathrm{I}}_q\bigl(T(U_{\bullet,j})\bigr)
  =
  \bigoplus_{r+s=j}
  T\bigl(\HH^{\mathrm{I}}_q(M_{\bullet,rs})\bigr).
\]
Em outras palavras, o complexo simples
$\HH^{\mathrm{I}}_q(T(U_{\bullet\bullet}))$ é o complexo simples
associado ao bicomplexo
$T(\HH^{\mathrm{I}}_q(M_{\bullet\bullet\bullet}))$.
Para todo $q$, $\HH^{\mathrm{I}}_q(M_{\bullet\bullet\bullet})$ é uma
resolução projetiva em $\cat{K}$ do complexo simples
$\HH^{\mathrm{I}}_q(K_{\bullet\bullet})$.
Aplicando \sref{0.11.6.3}, obtemos
\[
  \HH^{\mathrm{II}}_p
  \bigl(T(\HH^{\mathrm{I}}_q(M_{\bullet\bullet\bullet}))\bigr)
  =
  \LL_pT\bigl(\HH^{\mathrm{I}}_q(K_{\bullet\bullet})\bigr),
\]
e, portanto, a sequência espectral ${}^{(a')}\mathrm{E}$.
Finalmente, as sequências espectrais ${}^{(c)}\mathrm{E}$ e
${}^{(d)}\mathrm{E}$ obtêm-se permutando os papéis dos
índices $i$ e $j$ na definição do tricomplexo
$M_{\bullet\bullet\bullet}$ e aplicando os raciocínios anteriores ao
tricomplexo resultante.
\end{proof}

Chamamos $\LL_\bullet T(K_{\bullet\bullet})$ a \emph{hiperhomologia} de
$T$ relativa ao bicomplexo $K_{\bullet\bullet}$.

\begin{env}[11.7.3]
\label{0.11.7.3}
\emph{Observações}.
\begin{enumerate}
  \item[{\rm(i)}]
  De \sref{0.11.6.3} deduz-se que
  $\LL_\bullet T(K_{\bullet\bullet})$ É um funtor homólogo sobre a
  categoria dos bicomplexos $K_{\bullet\bullet}$ em $\cat{C}$ que estão
  limitados inferiormente em cada grau.

\oldpage[0\textsubscript{III}]{43}
  \item[{\rm(ii)}]
  Seja $M_{\bullet\bullet\bullet}$ um tricomplexo em $\cat{C}$ tal que,
  Para cada par $(r,s)$, $M_{rs,\bullet}$ seja uma resolução de $K_{rs}$,
  e tal que
  \[
    \LL_nT(M_{ijk})=0
  \]
  para toda terna $(i,j,k)$ e todo $n>0$.
  Então existe um isomorfismo
  \[
    \LL_\bullet T(K_{\bullet\bullet})
    \isoto
    \HH_\bullet\bigl(T(M_{\bullet\bullet\bullet})\bigr).
  \]
  A notação utilizada na demonstração de \sref{0.11.7.2}
  $N_{i,\bullet}$ é uma resolução de $L_i$ tal que
  $\LL_nT(N_{ij})=0$ para todo $n>0$ e tudo em par $(i,j)$, por isso
  Basta aplicar \sref{0.11.6.3},~(iii).

  \item[{\rm(iii)}]
  Os resultados de \sref{0.11.7.2} são imediatamente generalizados para os
  multifuntores covariantes.
  Por exemplo, seja $\cat{C}'$ uma segunda categoria abeliana em que todo
  objeto seja quociente de um objeto projetivo, seja
  $K'_{\bullet\bullet}$ um bicomplexo em $\cat{C}'$ cujos dois graus são
  limitados inferiormente, e seja $T$ um bifuntor covariante aditivo de
  $\cat{C}\times\cat{C}'$ em uma categoria abeliana $\cat{C}''$.
  Se $L_\bullet$ e $L'_\bullet$ são os complexos simples associados a
  $K_{\bullet\bullet}$ e $K'_{\bullet\bullet}$, respectivamente, define-se
  A hiperhomologia de $T$ relativa aos dois bicomplexos
  $K_{\bullet\bullet}$ e $K'_{\bullet\bullet}$ como
  \[
    \LL_\bullet T(L_\bullet,L'_\bullet).
  \]
  Aplicando \sref{0.11.6.4} e \sref{0.11.6.5}, obtêm-se, como em
  \sref{0.11.7.2}, seis sequências espectrais convergentes a esta hiperhomologia;
  Deixamos ao leitor a sua formulação explícita.
\end{enumerate}
\end{env}

\subsection{Suplemento sobre a cohomologia dos complexos simpliciais}
\label{subsection:0.11.8}

\begin{env}[11.8.1]
\label{0.11.8.1}
Seja $A$ um conjunto finito, e seja $\Sigma(A)$ o conjunto das sequências
finitas
\[
  \sigma=(\alpha_0,\ldots,\alpha_h)
\]
de elementos de $A$ (os ``símplices'' de $A$); vamos colocar
$|\sigma|=\{\alpha_0,\ldots,\alpha_h\}$.
Lembre-se que o complexo de cadeias $\mathrm{C}_\bullet(A)$ é o grupo abeliano
livre graduado gerado pelos elementos de $\Sigma(A)$, com
$(\alpha_0,\ldots,\alpha_h)$ em grau $h$, e com diferencial
\[
  d(\alpha_0,\ldots,\alpha_h)
  =
  \sum_{j=0}^{h}(-1)^j
  (\alpha_0,\ldots,\widehat{\alpha_j},\ldots,\alpha_h).
\]
O subgrupo $\mathrm{D}_\bullet(A)$ de $\mathrm{C}_\bullet(A)$ gerado
por cadeias $\sigma=(\alpha_0,\ldots,\alpha_h)$ para as quais dois dos
$\alpha_i$ são iguais, e pelas cadeias
\[
  \pi(\sigma)-\varepsilon_\pi\sigma,
\]
onde, para cada permutação $\pi$,
\[
  \pi(\sigma)
  =
  (\alpha_{\pi(0)},\ldots,\alpha_{\pi(h)})
\]
e $\varepsilon_\pi$ é o signo de $\pi$, é um subcomplexo de
$\mathrm{C}_\bullet(A)$.
Seus elementos são chamados de \emph{cadeias degeneradas}.
Ponhamos
\[
  \mathrm{L}_\bullet(A)
  =
  \mathrm{C}_\bullet(A)/\mathrm{D}_\bullet(A);
\]
Existe um homomorfismo natural de complexos
\[
  p:\mathrm{C}_\bullet(A)\longrightarrow\mathrm{L}_\bullet(A).
\]

Reciprocamente, definamos um homomorfismo de complexos
\[
  j:\mathrm{L}_\bullet(A)\longrightarrow\mathrm{C}_\bullet(A)
\]
como segue.
Escolhemos uma ordem total sobre $A$.
A classe módulo $\mathrm{D}_\bullet(A)$ de um símplice
$\sigma=(\alpha_0,\ldots,\alpha_h)$, associamos $0$ se dois dos
$\alpha_i$ são iguais e, caso contrário, associamos a sequência
$(\beta_0,\ldots,\beta_h)$ obtida ordenando os $\alpha_i$ em
ordem crescente.
É claro que $p\circ j$ é a identidade de
$\mathrm{L}_\bullet(A)$.
\end{env}

\begin{env}[11.8.2]
\label{0.11.8.2}
Seja $B$ um segundo conjunto finito.
Se $d'$ e $d''$ são os diferenciais de
$\mathrm{C}_\bullet(A)$ e $\mathrm{C}_\bullet(B)$, o complexo produto
tensorial
\[
  \mathrm{C}_\bullet(A)\otimes\mathrm{C}_\bullet(B)
\]
pode ser considerado como o grupo abeliano livre gerado pelos elementos de
$\Sigma(A)\times\Sigma(B)$, com diferencial
\[
  d(\sigma,\tau)
  =
  (d'\sigma,\tau)+(-1)^{h+1}(\sigma,d''\tau)
  \qquad\text{se }\operatorname{Card}|\sigma|=h+1.
\]

Os homomorfismos naturais
$\mathrm{C}_\bullet(A)\to\mathrm{L}_\bullet(A)$ e
$\mathrm{C}_\bullet(B)\to\mathrm{L}_\bullet(B)$ definem um homomorfismo
\[
  p:
  \mathrm{C}_\bullet(A)\otimes\mathrm{C}_\bullet(B)
  \longrightarrow
  \mathrm{L}_\bullet(A)\otimes\mathrm{L}_\bullet(B),
\]
sendo este último produto tensorial isomórfico a
\[
  \frac{\mathrm{C}_\bullet(A)\otimes\mathrm{C}_\bullet(B)}
  {\mathrm{C}_\bullet(A)\otimes\mathrm{D}_\bullet(B)
   +\mathrm{D}_\bullet(A)\otimes\mathrm{C}_\bullet(B)}.
\]
Analogamente, utilizando os homomorfismos
$\mathrm{L}_\bullet(A)\to\mathrm{C}_\bullet(A)$ e
$\mathrm{L}_\bullet(B)\to\mathrm{C}_\bullet(B)$ definidos em
\sref{0.11.8.1}, com ordens totais sobre $A$ e $B$, obtém-se um
homomorfismo
\[
  j:
  \mathrm{L}_\bullet(A)\otimes\mathrm{L}_\bullet(B)
  \longrightarrow
  \mathrm{C}_\bullet(A)\otimes\mathrm{C}_\bullet(B)
\]
tal que $p\circ j$ é a identidade.
\end{env}

\oldpage[0\textsubscript{III}]{44}
Com esta notação:
\begin{proposition}[11.8.3]
\label{0.11.8.3}
Existe uma homotopia
\[
  h:
  \mathrm{C}_\bullet(A)\otimes\mathrm{C}_\bullet(B)
  \longrightarrow
  \mathrm{C}_\bullet(A)\otimes\mathrm{C}_\bullet(B)
\]
tal que $h(\sigma,\tau)$ é uma combinação linear de pares de símplices
$(\sigma_i,\tau_i)$ que satisfazem
$|\sigma_i|\subset|\sigma|$ e $|\tau_i|\subset|\tau|$, e tal que,
se $f=j\circ p$, então
\[
  f-1=h\circ d+d\circ h.
  \tag{11.8.3.1}\label{0.11.8.3.1}
\]
\end{proposition}

\begin{proof}
Basta definir $h$ sobre cada par $(\sigma,\tau)$ de símplices, por
indução sobre a soma dos graus de $\sigma$ e $\tau$; quando essa soma
é $0$, pode ser tomado $h=0$.
Ponhamos
\[
  \omega
  =
  f(\sigma,\tau)-(\sigma,\tau)-h\bigl(d(\sigma,\tau)\bigr).
\]
Pela hipótese de indução e pela definição de $d$, temos
\[
  \omega
  \in
  \mathrm{C}_\bullet(|\sigma|)
  \otimes
  \mathrm{C}_\bullet(|\tau|).
\]
Além disso,
\[
  d\omega
  =
  f\bigl(d(\sigma,\tau)\bigr)
  -d(\sigma,\tau)
  -d\bigl(h(d(\sigma,\tau))\bigr)
  =
  h\bigl(d(d(\sigma,\tau))\bigr)
  =
  0
\]
por \eqref{0.11.8.3.1} e a hipótese de indução.
Ora,
\[
  \HH_q\bigl(\mathrm{C}_\bullet(A)\bigr)=0
  \qquad(q>0)
\]
(G, I, 3.7.4), e, consequentemente,
\[
  \HH_q\bigl(
    \mathrm{C}_\bullet(A)\otimes\mathrm{C}_\bullet(B)
  \bigr)=0
  \qquad(q>0)
\]
pela fórmula de K\"unneth (G, I, 5.5.2).
Aplicando esta observação com $A$ substituído por $|\sigma|$ e $B$ por
$|\tau|$, encontra-se
\[
  \omega'
  \in
  \mathrm{C}_\bullet(|\sigma|)
  \otimes
  \mathrm{C}_\bullet(|\tau|)
\]
tal que $\omega=d\omega'$.
Tomar $h(\sigma,\tau)=\omega'$ verifica \eqref{0.11.8.3.1} para o par
$(\sigma,\tau)$ e permite continuar a indução.
\end{proof}

\begin{env}[11.8.4]
\label{0.11.8.4}
Com a notação de \sref{0.11.8.2}, nós escrevemos
\[
  (\sigma,\tau)\leq(\sigma',\tau')
\]
se $|\sigma|\subset|\sigma'|$ e $|\tau|\subset|\tau'|$.
Um \emph{sistema de coeficientes} $\sh{S}$ sobre
$\Sigma(A)\times\Sigma(B)$ consiste em uma família
$(\Gamma_{\sigma,\tau})$ de grupos abelianos, onde
$\Gamma_{\sigma,\tau}$ depende apenas dos conjuntos $|\sigma|$ e
$|\tau|$, juntamente com homomorfismos
\[
  \Gamma_{\sigma,\tau}
  \longrightarrow
  \Gamma_{\sigma',\tau'}
  \qquad
  \bigl((\sigma,\tau)\leq(\sigma',\tau')\bigr)
\]
que formam um sistema indutivo para este preordenamento.

Definimos um complexo de cocadeias $\mathrm{C}^\bullet(A,B;\sh{S})$ cujos elementos
são as famílias
\[
  \lambda=\bigl(\lambda(\sigma,\tau)\bigr),
  \qquad
  \lambda(\sigma,\tau)\in\Gamma_{\sigma,\tau},
\]
quando $(\sigma,\tau)$ percorre $\Sigma(A)\times\Sigma(B)$.
Seu diferencial define-se da seguinte forma.
Se
\[
  d(\sigma,\tau)=\sum_i\pm(\sigma_i,\tau_i),
\]
então $|\sigma_i|\subset|\sigma|$ e $|\tau_i|\subset|\tau|$ para todo
$i$, e põe-se
\[
  d\lambda(\sigma,\tau)
  =
  \sum_i\pm\lambda_i(\sigma_i,\tau_i),
\]
onde $\lambda_i(\sigma_i,\tau_i)$ denota a imagem canônica de
$\lambda(\sigma_i,\tau_i)$ em $\Gamma_{\sigma,\tau}$.

Uma cocadeia $\lambda\in\mathrm{C}^\bullet(A,B;\sh{S})$ chama-se
\emph{bialternada} se $\lambda(\sigma,\tau)=0$ sempre que $\sigma$
ou $\tau$ tenham dois termos iguais, e se
\[
  \lambda(\pi(\sigma),\tau)
  =
  \varepsilon_\pi\lambda(\sigma,\tau),
  \qquad
  \lambda(\sigma,\pi'(\tau))
  =
  \varepsilon_{\pi'}\lambda(\sigma,\tau)
\]
para permutações arbitrárias $\pi$ e $\pi'$ dos índices.
Estas cocadeias geram claramente um subcomplexo
$\mathrm{L}^\bullet(A,B;\sh{S})$ de
$\mathrm{C}^\bullet(A,B;\sh{S})$.
\end{env}

\begin{proposition}[11.8.5]
\label{0.11.8.5}
A injeção canônica
\[
  \mathrm{L}^\bullet(A,B;\sh{S})
  \longrightarrow
  \mathrm{C}^\bullet(A,B;\sh{S})
\]
induz um isomorfismo em cohomologia.
\end{proposition}

\begin{proof}
Se $p$ e $j$ tiverem os significados atribuídos em \sref{0.11.8.2}, então
as aplicações
\[
  {}'\!p:\lambda\longmapsto\lambda\circ p,
  \qquad
  {}'\!j:\lambda\longmapsto\lambda\circ j
\]
são definidas em $\mathrm{L}^\bullet(A,B;\sh{S})$ e
$\mathrm{C}^\bullet(A,B;\sh{S})$, respectivamente; a primeira é precisamente
a injeção canônica.
Como ${}'\!j\circ{}'\!p$ é a identidade, só resta provar
que ${}'\!p\circ{}'\!j$ é homotópica à identidade.
Por \sref{0.11.8.3}, a aplicação
\[
  {}'\!h:\lambda\longmapsto\lambda\circ h
\]
é definida sobre $\mathrm{C}^\bullet(A,B;\sh{S})$, de modo que, ao transpor
\eqref{0.11.8.3.1}, obtém-se o resultado desejado.
\end{proof}

\begin{env}[11.8.6]
\label{0.11.8.6}
\oldpage[0\textsubscript{III}]{45}
A Proposição~\sref{0.11.8.5} reduz o cálculo da cohomologia de
$\mathrm{L}^\bullet(A,B;\sh{S})$ ao de
$\mathrm{C}^\bullet(A,B;\sh{S})$.
Por outro lado, pelo teorema de Eilenberg-Zilber
(G, I, 3.10.2), esta última cohomologia é canonicamente isomorfa à
cohomologia obtida a partir do complexo de cadeias
$\mathrm{P}_\bullet(A,B)$ de combinações lineares de pares
$(\sigma,\tau)\in\Sigma(A)\times\Sigma(B)$ para os quais $\sigma$ e
$\tau$ têm o mesmo grau.
O diferencial deste complexo é
\[
  d(\sigma,\tau)
  =
  \sum_{j,k}(-1)^{j+k}(\sigma_j,\tau_k)
\]
quando
\[
  d\sigma=\sum_j(-1)^j\sigma_j,
  \qquad
  d\tau=\sum_k(-1)^k\tau_k.
\]
Existem dois homomorfismos canônicos de complexos
\[
  f:\mathrm{P}_\bullet(A,B)
  \longrightarrow
  \mathrm{C}_\bullet(A)\otimes\mathrm{C}_\bullet(B),
  \qquad
  g:\mathrm{C}_\bullet(A)\otimes\mathrm{C}_\bullet(B)
  \longrightarrow
  \mathrm{P}_\bullet(A,B),
\]
e demonstra-se que existem homotópias $h$ e $h'$ tais que
\[
  f\circ g-1=d\circ h+h\circ d,
  \qquad
  g\circ f-1=d\circ h'+h'\circ d.
\]
Além disso,
\[
  f(\sigma,\tau)
  \in
  \mathrm{C}_\bullet(|\sigma|)
  \otimes
  \mathrm{C}_\bullet(|\tau|),
  \qquad
  g(\sigma,\tau)
  \in
  \mathrm{P}_\bullet(|\sigma|,|\tau|),
\]
e as homotópias podem ser escolhidas de modo a
\[
  h(\sigma,\tau)
  \in
  \mathrm{C}_\bullet(|\sigma|)
  \otimes
  \mathrm{C}_\bullet(|\tau|),
  \qquad
  h'(\sigma,\tau)
  \in
  \mathrm{P}_\bullet(|\sigma|,|\tau|).
\]
Com efeito, $h(\sigma,\tau)$ e $h'(\sigma,\tau)$ podem ser definidos por
indução sobre a soma dos graus de $\sigma$ e $\tau$, enquanto
\[
  \HH_q\bigl(
    \mathrm{C}_\bullet(|\sigma|)
    \otimes
    \mathrm{C}_\bullet(|\tau|)
  \bigr)
  =
  \HH_q\bigl(\mathrm{P}_\bullet(|\sigma|,|\tau|)\bigr)
  =
  0
  \qquad(q>0)
\]
(\emph{loc.\ cit.}); o raciocínio é o mesmo que em
\sref{0.11.8.3}.

Agora definimos $\mathrm{P}^\bullet(A,B;\sh{S})$ como o conjunto de famílias
\[
  \lambda=\bigl(\lambda(\sigma,\tau)\bigr),
  \qquad
  \lambda(\sigma,\tau)\in\Gamma_{\sigma,\tau},
\]
onde $(\sigma,\tau)$ percorre os pares de símplices do mesmo grau.
Quando
\[
  d\sigma=\sum_j(-1)^j\sigma_j,
  \qquad
  d\tau=\sum_k(-1)^k\tau_k,
\]
a fórmula
\[
  d\lambda(\sigma,\tau)
  =
  \sum_{j,k}(-1)^{j+k}\lambda(\sigma_j,\tau_k)
\]
está definida e fornece o diferencial de
$\mathrm{P}^\bullet(A,B;\sh{S})$.
As aplicações
\[
  {}'\!f:\lambda\longmapsto\lambda\circ f,\quad
  {}'\!g:\lambda\longmapsto\lambda\circ g,\quad
  {}'\!h:\lambda\longmapsto\lambda\circ h,\quad
  {}'\!h':\lambda\longmapsto\lambda\circ h'
\]
são todas definidas pelas propriedades de suporte anteriores.
Assim, ${}'\!f\circ{}'\!g$ e ${}'\!g\circ{}'\!f$ são homotópicas à
identidade, dando o isomorfismo requerido entre a cohomologia de
$\mathrm{C}^\bullet(A,B;\sh{S})$ e a de
$\mathrm{P}^\bullet(A,B;\sh{S})$.
\end{env}

\begin{env}[11.8.7]
\label{0.11.8.7}
\emph{Observação}.
O raciocínio de \sref{0.11.8.3}, aplicado a
$\mathrm{C}_\bullet(A)$ e $\mathrm{L}_\bullet(A)$, mostra que se
$j$ e $p$ são definidos como em \sref{0.11.8.1}, então
$f=j\circ p$ satisfaz novamente \eqref{0.11.8.3.1}, com
$|h(\sigma)|\subset|\sigma|$.
Como em \sref{0.11.8.5}, obtém-se um isomorfismo da cohomologia de
$\mathrm{L}^\bullet(A;\sh{S})$ na de
$\mathrm{C}^\bullet(A;\sh{S})$, definindo os dois complexos de forma a
evidente.
Este é o resultado cuja demonstração é esboçada em (G, I, 3.8.1).
\end{env}

\begin{env}[11.8.8]
\label{0.11.8.8}
Voltemos agora à notação e hipóteses de \sref{0.11.8.2}, e seja
\[
  \sh{S}^\bullet=(\sh{S}^k)
\]
um complexo de sistemas de coeficientes sobre
$\Sigma(A)\times\Sigma(B)$.
Para cada $(\sigma,\tau)$, os grupos $\Gamma^k_{\sigma,\tau}$
formam, portanto, um complexo de grupos abelianos $(k\in\bb{Z})$, e para
$(\sigma,\tau)\leq(\sigma',\tau')$ existem diagramas comutativos
\[
  \xymatrix{
    \Gamma^k_{\sigma,\tau}\ar[r]\ar[d] &
    \Gamma^{k+1}_{\sigma,\tau}\ar[d]\\
    \Gamma^k_{\sigma',\tau'}\ar[r] &
    \Gamma^{k+1}_{\sigma',\tau'}.
  }
\]
\oldpage[0\textsubscript{III}]{46}
Deduz-se imediatamente que
\[
  \mathrm{C}^\bullet(A,B;\sh{S}^\bullet)
  =
  \bigl(
    \mathrm{C}^h(A,B;\sh{S}^k)
  \bigr)_{(h,k)\in\bb{Z}\times\bb{Z}}
\]
é um bicomplexo de grupos abelianos, e que
\[
  \mathrm{L}^\bullet(A,B;\sh{S}^\bullet)
  =
  \bigl(
    \mathrm{L}^h(A,B;\sh{S}^k)
  \bigr)
\]
é um sub-bicomplexo.
\end{env}

\begin{proposition}[11.8.9]
\label{0.11.8.9}
A injeção canônica
\[
  \mathrm{L}^\bullet(A,B;\sh{S}^\bullet)
  \longrightarrow
  \mathrm{C}^\bullet(A,B;\sh{S}^\bullet)
\]
induz um isomorfismo na cohomologia destes dois bicomplexos.
\end{proposition}

\begin{proof}
Por brevidade, ponhamos
\[
  \mathrm{C}^{\bullet\bullet}
  =
  \mathrm{C}^\bullet(A,B;\sh{S}^\bullet),
  \qquad
  \mathrm{L}^{\bullet\bullet}
  =
  \mathrm{L}^\bullet(A,B;\sh{S}^\bullet).
\]
Como $\mathrm{C}^{hk}=\mathrm{L}^{hk}=0$ para $h<0$, as segundas
sequências espectrais destes bicomplexos são regulares
\sref{0.11.3.3}.
O homomorfismo
$\mathrm{L}^{\bullet\bullet}\to\mathrm{C}^{\bullet\bullet}$ dá, portanto,
um morfismo de sequências espectrais
\[
  {}''\mathrm{E}(\mathrm{L}^{\bullet\bullet})
  \longrightarrow
  {}''\mathrm{E}(\mathrm{C}^{\bullet\bullet})
\]
que, sobre os termos $\mathrm{E}_2$, é
\[
  \HH^p_{\mathrm{II}}
  \bigl(\HH^q_{\mathrm{I}}(\mathrm{L}^{\bullet\bullet})\bigr)
  \longrightarrow
  \HH^p_{\mathrm{II}}
  \bigl(\HH^q_{\mathrm{I}}(\mathrm{C}^{\bullet\bullet})\bigr).
  \tag{11.8.9.1}\label{0.11.8.9.1}
\]
No entanto, para todo $k\in\bb{Z}$, \sref{0.11.8.5} mostra que
\[
  \HH^q_{\mathrm{I}}(\mathrm{L}^{\bullet,k})
  \longrightarrow
  \HH^q_{\mathrm{I}}(\mathrm{C}^{\bullet,k})
\]
é bijetivo.
A conclusão deduz-se de \sref{0.11.1.5}.
\end{proof}

\begin{env}[11.8.10]
\label{0.11.8.10}
Assim, com a notação de \sref{0.11.8.6}, existe um
homomorfismo canônico de bicomplexos
\[
  \mathrm{C}^\bullet(A,B;\sh{S}^\bullet)
  \longrightarrow
  \mathrm{P}^\bullet(A,B;\sh{S}^\bullet),
\]
com notação evidente.
O raciocínio de \sref{0.11.8.9}, usando agora \sref{0.11.8.6}, mostra que
este homomorfismo também induz um isomorfismo em cohomologia.
\end{env}

\subsection{Um lema sobre os complexos de tipo finito}
\label{subsection:0.11.9}

\begin{proposition}[11.9.1]
\label{0.11.9.1}
Seja $\cat{C}$ uma categoria abeliana, e sejam $K'$ e $K''$ subconjuntos
do conjunto de objetos de $\cat{C}$, com $K''\subset K'$, que satisfaçam
as seguintes condições:
\begin{enumerate}
  \item[{\rm(i)}]
  \emph{Para todo $A'\in K'$ e todo epimorfismo $u:A\to A'$ em
  $\cat{C}$, existe um objeto $B\in K''$ e um morfismo $v:B\to A$
  tais que $u\circ v$ é um epimorfismo.}

  \item[{\rm(ii)}]
  \emph{Para todo par $A,B\in K'$ e todo epimorfismo $u:A\to B$,
  o objeto $\Ker(u)$ pertence a $K'$.}

  \item[{\rm(iii)}]
  \emph{O produto de dois objetos de $K''$ pertence a $K''$.}
\end{enumerate}

\emph{Seja $P_\bullet=(P_i)_{i\in\bb{Z}}$ um complexo em $\cat{C}$ tal
que $\HH_i(P_\bullet)\in K'$ para todo $i$, e vamos supor que existe um
inteiro $d$ para o qual $\HH_i(P_\bullet)=0$ quando $i<d$.
Então há um complexo $Q_\bullet=(Q_i)_{i\in\bb{Z}}$ em $\cat{C}$
tal que $Q_i\in K''$ para todo $i$ e $Q_i=0$ para $i<d$, e um
homomorfismo de complexos
\[
  u:Q_\bullet\longrightarrow P_\bullet
\]
cujo morfismo induzido
\[
  \HH_\bullet(Q_\bullet)\longrightarrow\HH_\bullet(P_\bullet)
\]
é um isomorfismo.}
\end{proposition}

\begin{proof}
Primeiro, demonstremos a seguinte consequência da condição ~(i).

\medskip\noindent
\emph{(i bis)}
\emph{Seja $u:C\to B$ um epimorfismo em $\cat{C}$, seja $A$ um
objeto de $K'$, e seja $v:A\to B$ um morfismo em $\cat{C}$.
Então há um objeto $D\in K''$, um epimorfismo $u':D\to A$ e um
morfismo $v':D\to C$ tais como o diagrama}
\[
  \xymatrix{
    D\ar[r]^{u'}\ar[d]_{v'} & A\ar[d]^{v}\\
    C\ar[r]_{u} & B
  }
  \tag{11.9.1.1}\label{0.11.9.1.1}
\]
\emph{é comutativo.}

Vamos formar o produto fibrado $C\times_B A$ em $\cat{C}$ e sejam
\[
  p:C\times_B A\longrightarrow C,
  \qquad
  q:C\times_B A\longrightarrow A
\]
as projeções canônicas, de modo que
\oldpage[0\textsubscript{III}]{47}
\[
  \xymatrix{
    C\times_B A\ar[r]^{q}\ar[d]_{p} & A\ar[d]^{v}\\
    C\ar[r]_{u} & B
  }
  \tag{11.9.1.2}\label{0.11.9.1.2}
\]
é comutativo.
O conúcleo de $q$ é o quociente de $A$ por $v^{-1}(u(C))$
([27, p.~1-12]).
Como $u$ é um epimorfismo, $u(C)=B$ e
$v^{-1}(u(C))=A$; portanto, $q$ é um epimorfismo.
Aplicando ~(i) a $q:C\times_B A\to A$, obtém-se um objeto
$D\in K''$ e um morfismo $w:D\to C\times_B A$ tais que $q\circ w$ é
um epimorfismo.
Basta tomar $u'=q\circ w$ e $v'=p\circ w$.

\medskip
Vamos agora demonstrar a proposição por indução.
Suponhamos, para algum $i\geq d-1$, que para todo $j\leq i$ temos
construído objetos $Q_j$ e diferenciais
\[
  d^Q_j:Q_j\longrightarrow Q_{j-1}
\]
com $Q_j=0$ para $j<d$ e
$d^Q_{j-1}\circ d^Q_j=0$, e morfismos
\[
  u_j:Q_j\longrightarrow P_j
\]
tais que
\[
  d^P_j\circ u_j=u_{j-1}\circ d^Q_j
  \qquad(j\leq i).
\]
Suponhamos, além disso, que as seguintes condições sejam cumpridas:
\begin{enumerate}
  \item[$(\mathrm{I}_i)$]
  $Q_j\in K''$ para $j\leq i$, e
  $\mathrm{B}_j(Q_\bullet)\in K'$ para $j<i$.

  \item[$(\mathrm{II}_i)$]
  Para $j<i$, o homomorfismo
  \[
    \HH_j(Q_\bullet)\longrightarrow\HH_j(P_\bullet)
  \]
  induzido por $(u_k)_{k\leq i}$ é um isomorfismo.

  \item[$(\mathrm{III}_i)$]
  O composto
  \[
    v_i:
    \mathrm{Z}_i(Q_\bullet)
    \longrightarrow
    \mathrm{Z}_i(P_\bullet)
    \longrightarrow
    \HH_i(P_\bullet),
  \]
  onde a primeira flecha é a restrição de $u_i$ e a segunda é
  canônica, é um epimorfismo.
\end{enumerate}

Por~(ii), $\mathrm{Z}_i(Q_\bullet)$ pertence a $K'$: é o núcleo do
epimorfismo
\[
  Q_i\longrightarrow\mathrm{B}_{i-1}(Q_\bullet),
\]
e aqui é usado $(\mathrm{I}_i)$.
A condição ~(ii) e $(\mathrm{III}_i)$ também mostram que
\[
  N_i=\Ker(v_i)
\]
pertence a $K'$.
Por (i bis), há um objeto $Q'_{i+1}\in K''$, um epimorfismo
\[
  d'_{i+1}:Q'_{i+1}\longrightarrow N_i,
\]
e um morfismo $u'_{i+1}:Q'_{i+1}\to P_{i+1}$ tais que
\[
  \xymatrix{
    Q'_{i+1}\ar[r]^{d'_{i+1}}\ar[d]_{u'_{i+1}} &
    N_i\ar[d]^{u_i}\\
    P_{i+1}\ar[r]_{d^P_{i+1}} &
    \mathrm{B}_i(P_\bullet)
  }
  \tag{11.9.1.3}\label{0.11.9.1.3}
\]
é comutativo.

O morfismo canônico
\[
  \mathrm{Z}_{i+1}(P_\bullet)
  \longrightarrow
  \HH_{i+1}(P_\bullet)
\]
é um epimorfismo, e $\HH_{i+1}(P_\bullet)\in K'$.
A condição ~(i) fornece, portanto, um objeto $Q''_{i+1}\in K''$ e um
morfismo
\[
  u''_{i+1}:Q''_{i+1}\longrightarrow
  \mathrm{Z}_{i+1}(P_\bullet)
\]
tal que o composto
\[
  Q''_{i+1}
  \longrightarrow
  \mathrm{Z}_{i+1}(P_\bullet)
  \longrightarrow
  \HH_{i+1}(P_\bullet)
\]
é um epimorfismo.
Se $d''_{i+1}:Q''_{i+1}\to N_i$ é o morfismo nulo, então
\[
  \xymatrix{
    Q''_{i+1}\ar[r]^{d''_{i+1}}\ar[d]_{u''_{i+1}} &
    N_i\ar[d]^{u_i}\\
    \mathrm{Z}_{i+1}(P_\bullet)\ar[r]_{d^P_{i+1}} &
    P_i
  }
  \tag{11.9.1.4}\label{0.11.9.1.4}
\]
é comutativo; a seta horizontal inferior é nula.

\oldpage[0\textsubscript{III}]{48}
Ponhamos
\[
  Q_{i+1}=Q'_{i+1}\times Q''_{i+1},
  \qquad
  d^Q_{i+1}=d'_{i+1}+d''_{i+1},
  \qquad
  u_{i+1}=u'_{i+1}+u''_{i+1}.
\]
Por ~(iii), $Q_{i+1}$ pertence a $K''$.
Como
\[
  d^Q_{i+1}(Q_{i+1})
  =
  d'_{i+1}(Q'_{i+1})
  =
  N_i
  \subset
  \mathrm{Z}_i(Q_\bullet),
\]
tem-se
$d^Q_i\circ d^Q_{i+1}=0$ e
$\mathrm{B}_i(Q_\bullet)=N_i$; isso demonstra
$(\mathrm{I}_{i+1})$.
A comutabilidade de \eqref{0.11.9.1.3} e
\eqref{0.11.9.1.4} da
\[
  d^P_{i+1}\circ u_{i+1}
  =
  u_i\circ d^Q_{i+1}.
\]
Por definição de $N_i$, o morfismo
\[
  \HH_i(Q_\bullet)
  =
  \mathrm{Z}_i(Q_\bullet)/N_i
  \longrightarrow
  \HH_i(P_\bullet)
\]
induzido pelos $u_k$ para $k\leq i+1$ é um isomorfismo, já que $v_i$ é
um epimorfismo; portanto, cumpre-se $(\mathrm{II}_{i+1})$.
Finalmente, o fator $Q''_{i+1}$ está contido em
$\mathrm{Z}_{i+1}(Q_\bullet)$.
A escolha de $u''_{i+1}$ mostra que
\[
  v_{i+1}:
  \mathrm{Z}_{i+1}(Q_\bullet)
  \longrightarrow
  \mathrm{Z}_{i+1}(P_\bullet)
  \longrightarrow
  \HH_{i+1}(P_\bullet)
\]
é um epimorfismo, porque sua restrição a $Q''_{i+1}$ já o
é.
Cumpre-se assim $(\mathrm{III}_{i+1})$; a indução continua e deduz-se a
proposição.
\end{proof}

\begin{corollary}[11.9.2]
\label{0.11.9.2}
Seja $A$ um anel noetheriano, não necessariamente comutativo, e seja
$P_\bullet=(P_i)_{i\in\bb{Z}}$ um complexo de $A$-módulos à direita.
Suponhamos que todo $\HH_i(P_\bullet)$ é um $A$-módulo finitamente
gerado e que existe um inteiro $d$ tal que
$\HH_i(P_\bullet)=0$ para $i<d$.
Então existe um complexo $Q_\bullet=(Q_i)_{i\in\bb{Z}}$ de
$A$-módulos livres de posto finito, com $Q_i=0$ para $i<d$, e um homomorfismo
de complexos
\[
  u:Q_\bullet\longrightarrow P_\bullet
\]
cujo morfismo induzido
\[
  \HH_\bullet(Q_\bullet)\longrightarrow\HH_\bullet(P_\bullet)
\]
é um isomorfismo.
\end{corollary}

\begin{proof}
Aplicamos \sref{0.11.9.1} tomando como $\cat{C}$ a categoria dos
$A$-módulos à direita, como $K'$ o conjunto dos $A$-módulos finitamente gerados e
como $K''$ o conjunto de $A$-módulos livres de posto finito.
As condições ~(i), (ii) e ~(iii) de \sref{0.11.9.1} são deduzidas imediatamente
da hipótese de que $A$ é noetheriano.
\end{proof}

\begin{env}[11.9.3]
\label{0.11.9.3}
\emph{Observações}.
\begin{enumerate}
  \item[{\rm(i)}]
  Sob a hipótese de \sref{0.11.9.2}, supomos ainda que os
  $P_i$ são $A$-módulos planos à direita.
  Então, para todo $A$-módulo à esquerda $M$, o homomorfismo de complexos
  \[
    u\otimes 1:
    Q_\bullet\otimes_A M
    \longrightarrow
    P_\bullet\otimes_A M
  \]
  induz novamente um isomorfismo
  \[
    \HH_\bullet(Q_\bullet\otimes_A M)
    \isoto
    \HH_\bullet(P_\bullet\otimes_A M),
  \]
  Como será visto no Capítulo~III.

  \item[{\rm(ii)}]
  A conclusão de \sref{0.11.9.2} não tem por que ser cumprida quando $A$ não é
  noetheriano.
  Com efeito, aplicá-la a um complexo que tivesse apenas um termo não nulo
  implicaria que todo $A$-módulo à direita finitamente gerado admitiria uma resolução por
  módulos livres finitamente gerados, o que é falso em geral
(Bourbaki, \emph{Alg\`ebre commutative}, cap.~I, \S2, exerc.~6).

  No entanto, é possível substituir a hipótese de que $A$ é
  noetheriano pela suposição de que os $\HH_i(P_\bullet)$ admitam
  $\infty$-apresentações finitas (cf.\ Capítulo~IV).
\end{enumerate}
\end{env}

\subsection{Característica de Euler--Poincar\'e de um complexo de módulos de comprimento finito}
\label{subsection:0.11.10}

\begin{env}[11.10.1]
\label{0.11.10.1}
Seja $A$ um anel, não necessariamente comutativo, e seja
\[
  M^\bullet:
  0\longrightarrow M^0\longrightarrow M^1
  \longrightarrow\cdots\longrightarrow M^n\longrightarrow 0
  \tag{11.10.1.1}\label{0.11.10.1.1}
\]
um complexo de $A$-módulos à esquerda de comprimento finito.
A \emph{característica de Euler--Poincar\'e} deste complexo é o
inteiro
\oldpage[0\textsubscript{III}]{49}
\[
  \chi(M^\bullet)
  =
  \sum_{i=0}^{n}(-1)^i\operatorname{length}(M^i).
  \tag{11.10.1.2}\label{0.11.10.1.2}
\]
\end{env}

\begin{proposition}[11.10.2]
\label{0.11.10.2}
Para todo complexo finito $M^\bullet$ de $A$-módulos à esquerda de comprimento
finito,
\[
  \chi(M^\bullet)=\chi\bigl(\HH^\bullet(M^\bullet)\bigr),
\]
onde $\HH^\bullet(M^\bullet)$ é considerado como um complexo com diferencial
nulo.
Em particular, se a sequência \eqref{0.11.10.1.1} for exata, então
$\chi(M^\bullet)=0$.
\end{proposition}

\begin{proof}
Por brevidade, escrevamos
\[
  \mathrm{B}^i=\mathrm{B}^i(M^\bullet),
  \qquad
  \mathrm{Z}^i=\mathrm{Z}^i(M^\bullet),
  \qquad
  \HH^i=\HH^i(M^\bullet)=\mathrm{Z}^i/\mathrm{B}^i.
\]
Os módulos $\mathrm{B}^i$, $\mathrm{Z}^i$ e $\HH^i$ têm todos
comprimento finito.
As sequências exatas
\[
  0\longrightarrow\mathrm{B}^i
  \longrightarrow\mathrm{Z}^i
  \longrightarrow\HH^i
  \longrightarrow0
\]
e
\[
  0\longrightarrow\mathrm{Z}^i
  \longrightarrow M^i
  \longrightarrow\mathrm{B}^{i+1}
  \longrightarrow0
\]
dão
\[
  \operatorname{length}(\mathrm{Z}^i)
  =
  \operatorname{length}(\HH^i)
  +
  \operatorname{length}(\mathrm{B}^i)
\]
e
\[
  \operatorname{length}(M^i)
  =
  \operatorname{length}(\mathrm{Z}^i)
  +
  \operatorname{length}(\mathrm{B}^{i+1}).
\]
Consequentemente,
\[
  \operatorname{length}(M^i)-\operatorname{length}(\HH^i)
  =
  \operatorname{length}(\mathrm{B}^{i+1})
  +
  \operatorname{length}(\mathrm{B}^i).
\]
Multipliquemos esta igualdade por $(-1)^i$ e somemos para $0\leq i\leq n$.
Como $\mathrm{B}^0=\mathrm{B}^{n+1}=0$, deduz-se a igualdade desejada.
\end{proof}

\begin{corollary}[11.10.3]
\label{0.11.10.3}
Seja $\mathrm{E}=(\mathrm{E}_r^{pq})$ uma sequência espectral na
categoria de módulos sobre um anel $A$.
Suponhamos que os $\mathrm{E}_2^{pq}$ são $A$-módulos de comprimento finito e
que $\mathrm{E}_2^{pq}\neq0$ apenas para um número finito de pares $(p,q)$.
Então as características de Euler--Poincar\'e
\[
  \chi\bigl(\mathrm{E}_r^{(\bullet)}\bigr)
\]
de todos os complexos
\[
  \mathrm{E}_r^{(\bullet)}
  =
  \bigl(\mathrm{E}_r^{(n)}\bigr)_{n\in\bb{Z}}
\]
de \sref{0.11.1.1} são iguais.
Se, além disso, $\mathrm{E}$ é fracamente convergente e se define
\[
  \mathrm{E}_\infty^{(n)}
  =
  \bigoplus_{p+q=n}\mathrm{E}_\infty^{pq}
  \qquad(n\in\bb{Z}),
\]
então
\[
  \chi\bigl(\mathrm{E}_\infty^{(\bullet)}\bigr)
  =
  \chi\bigl(\mathrm{E}_2^{(\bullet)}\bigr),
\]
onde
\[
  \mathrm{E}_\infty^{(\bullet)}
  =
  \bigl(\mathrm{E}_\infty^{(n)}\bigr)_{n\in\bb{Z}}
\]
é considerado como um complexo com diferencial nulo.
\end{corollary}

\begin{proof}
Vamos observar primeiro que se $\mathrm{E}_2^{pq}=0$, então
$\mathrm{E}_r^{pq}=0$ para $2\leq r\leq+\infty$.
Então, todo complexo $\mathrm{E}_r^{(\bullet)}$ é finito e é formado por
$A$-módulos de comprimento finito.
Para todo $r$ finito, a igualdade
\[
  \chi\bigl(\mathrm{E}_r^{(\bullet)}\bigr)
  =
  \chi\bigl(\mathrm{E}_{r+1}^{(\bullet)}\bigr)
\]
deduz-se, por conseguinte, de \sref{0.11.10.2} e do isomorfismo
\[
  \HH^\bullet\bigl(\mathrm{E}_r^{(\bullet)}\bigr)
  \simeq
  \mathrm{E}_{r+1}^{(\bullet)}
\]
de complexos com diferencial nulo.

Como os $\mathrm{E}_r^{pq}$ têm comprimento finito, para cada par
$(p,q)$, as sequências
\[
  \bigl(\mathrm{B}_k(\mathrm{E}_2^{pq})\bigr)_{k\geq2},
  \qquad
  \bigl(\mathrm{Z}_k(\mathrm{E}_2^{pq})\bigr)_{k\geq2}
\]
são estacionárias.
A fraca convergência, juntamente com o fato de que
$\mathrm{E}_2^{pq}=0$ exceto para um número finito de pares $(p,q)$, implica, portanto,
que algum inteiro $r\geq2$ satisfaça
\[
  \mathrm{E}_\infty^{pq}=\mathrm{E}_r^{pq}
\]
para todo par $(p,q)$.
Deduz-se a afirmação relativa a
$\chi(\mathrm{E}_\infty^{(\bullet)})$.
\end{proof}

\section{Suplemento sobre a cohomologia de feixes}
\label{section:0.12}

\subsection{Cohomologia de feixes de módulos sobre espaços anelados}
\label{subsection:0.12.1}

\begin{env}[12.1.1]
\label{0.12.1.1}
\oldpage[0\textsubscript{III}]{49}
Seja $(X,\sh{O}_X)$ um espaço anelado.
Lembremos que, para todo $\sh{O}_X$-módulo $\sh{F}$, define-se a
cohomologia $\HH^\bullet(X,\sh{F})$, que é um funtor cohomológico universal
(T, 2.2) da categoria $\cat{C}(X)$ de
$\sh{O}_X$-módulos na categoria de grupos abelianos; é o funtor derivado
do funtor exato à esquerda
\[
  \sh{F}\longmapsto\Gamma(X,\sh{F}).
\]
O funtor $\HH^\bullet(X,\sh{F})$ é isomorfo
\oldpage[0\textsubscript{III}]{50}
à restrição à categoria $\cat{C}(X)$ do funtor cohomológico
definido do mesmo modo sobre a categoria de feixes de grupos abelianos
em $X$ (G, II, 7.2.1).
\end{env}

\begin{env}[12.1.2]
\label{0.12.1.2}
Ponhamos $A=\Gamma(X,\sh{O}_X)$.
Como todo elemento de $A$ define um endomorfismo do grupo abeliano
$\Gamma(X,\sh{F})$, define por funtorialidade um endomorfismo do
$\delta$-funtor $\HH^\bullet(X,\sh{F})$.
Estes endomorfismos dotam cada $\HH^p(X,\sh{F})$ de uma
estrutura de $A$-módulo, e o operador de conexão $\partial$ é
$A$-linear.
Além disso, para quaisquer dois inteiros não negativos $p,q$ e quaisquer dois
$\sh{O}_X$-módulos $\sh{F},\sh{G}$, existe um homomorfismo de
$A$-módulos, chamado \emph{produto cup},
\[
\label{0.12.1.2.1}
\tag{12.1.2.1}
  \HH^p(X,\sh{F})\otimes_A\HH^q(X,\sh{G})
  \longrightarrow
  \HH^{p+q}\bigl(X,\sh{F}\otimes_{\sh{O}_X}\sh{G}\bigr)
\]
(G, II, 6.6).
Estes homomorfismos fazem da soma direta $S$ dos
$\HH^p(X,\sh{O}_X)$, para $p\geq0$, uma $A$-álgebra graduada
anticomutativa, e da soma direta dos $\HH^p(X,\sh{F})$ um
$S$-módulo graduado.

Para todo recobrimento aberto $\mathfrak{U}$ de $X$, denotaremos sempre
por $\check{\mathrm{C}}^\bullet(\mathfrak{U},\sh{F})$ (contrariamente a
(G, II, 5.1)) o complexo de cocadeias alternadas do nervo de
$\mathfrak{U}$ com valores no sistema de coeficientes
$\Gamma(U_\alpha,\sh{F})$.
Está claro que
$\check{\mathrm{C}}^\bullet(\mathfrak{U},\sh{F})$ é um
$A$-módulo graduado, de modo que os grupos de cohomologia
$\check{\HH}^p(\mathfrak{U},\sh{F})$ deste complexo estão dotados de
uma estrutura de $A$-módulo.
Além disso, as aplicações canônicas
\[
  \check{\HH}^p(\mathfrak{U},\sh{F})
  \longrightarrow
  \HH^p(X,\sh{F})
\]
(G, II, 5.4) são $A$-lineares.
\end{env}

\begin{env}[12.1.3]
\label{0.12.1.3}
Seja $(X',\sh{O}_{X'})$ um segundo espaço anelado, e seja
$f=(\psi,\theta)$ um morfismo de $X'$ em $X$.
Ponhamos $A'=\Gamma(X',\sh{O}_{X'})$; o $\psi$-morfismo $\theta$
define canonicamente um homomorfismo de anéis $A\to A'$.
Seja $\sh{F}$ um $\sh{O}_X$-módulo e $\sh{F}'$ um
$\sh{O}_{X'}$-módulo.
Para todo $f$-morfismo $u:\sh{F}\to\sh{F}'$
\sref[0\textsubscript{I}]{0.4.4.1}, veremos que se pode definir,
para todo $p\geq0$, um dihomomorfismo
\[
\label{0.12.1.3.1}
\tag{12.1.3.1}
  u_p:\HH^p(X,\sh{F})\longrightarrow\HH^p(X',\sh{F}').
\]

Com efeito, como $\psi^*$ é exato na categoria de feixes de grupos
abelianos sobre $X$, o funtor
\[
  \sh{F}\longmapsto\HH^\bullet(X',\psi^*(\sh{F}))
\]
é um $\delta$-funtor sobre dita categoria, e sabe-se que existe um
homomorfismo canônico de $\delta$-funtores
\[
\label{0.12.1.3.2}
\tag{12.1.3.2}
  \HH^\bullet(X,\sh{F})
  \longrightarrow
  \HH^\bullet(X',\psi^*(\sh{F}))
\]
determinado de maneira única pela condição de que, em grau $0$, reduz-se
ao homomorfismo canônico
\[
  \Gamma(X,\sh{F})
  \longrightarrow
  \Gamma(X',\psi^*(\sh{F}))
\]
(T, 3.2.2).
Além disso, todo elemento de $A$ determina um endomorfismo $\mu$ de
$\Gamma(X,\sh{F})$ e um endomorfismo $\mu'$ de
$\Gamma(X',\psi^*(\sh{F}))$ tal que o diagrama
\[
\label{0.12.1.3.3}
\tag{12.1.3.3}
  \xymatrix{
    \Gamma(X,\sh{F})\ar[r]\ar[d]_\mu
      & \Gamma(X',\psi^*(\sh{F}))\ar[d]^{\mu'}\\
    \Gamma(X,\sh{F})\ar[r]
      & \Gamma(X',\psi^*(\sh{F}))
  }
\]
é comutativo.
\oldpage[0\textsubscript{III}]{51}
Pela propriedade de unicidade para prolongar morfismos de funtores
cohomológicos universais (T, 2.2), obtêm-se prolongações únicas de
$\mu$ e $\mu'$ à cohomologia que fazem comutativos os diagramas
análogos a \sref{0.12.1.3.3}.
Assim, \sref{0.12.1.3.2} é um homomorfismo de $A$-módulos.

Observemos agora que
\[
  f^*(\sh{F})
  =
  \psi^*(\sh{F})
  \otimes_{\psi^*(\sh{O}_X)}
  \sh{O}_{X'},
\]
e, portanto, que existe um dihomomorfismo canônico
\[
  \psi^*(\sh{F})\longrightarrow f^*(\sh{F})
\]
do $\psi^*(\sh{O}_X)$-módulo $\psi^*(\sh{F})$ no
$\sh{O}_{X'}$-módulo $f^*(\sh{F})$.
Por funtorialidade, obtém-se, pois, um dihomomorfismo funtorial
\[
\label{0.12.1.3.4}
\tag{12.1.3.4}
  \HH^p(X',\psi^*(\sh{F}))
  \longrightarrow
  \HH^p(X',f^*(\sh{F})),
\]
sendo $A$ e $A'$ os anéis correspondentes.
Compondo este dihomomorfismo com \sref{0.12.1.3.2}, obtém-se um
dihomomorfismo canônico, funtorial em $\sh{F}$,
\[
\label{0.12.1.3.5}
\tag{12.1.3.5}
  \theta_p:\HH^p(X,\sh{F})
  \longrightarrow
  \HH^p(X',f^*(\sh{F})).
\]

Finalmente, por funtorialidade, o homomorfismo
$u^\sharp:f^*(\sh{F})\to\sh{F}'$ dá um homomorfismo de
$A'$-módulos
\[
  \HH^p(X',f^*(\sh{F}))
  \longrightarrow
  \HH^p(X',\sh{F}'),
\]
que, composto com \sref{0.12.1.3.5}, dá
\sref{0.12.1.3.1}.

Seja $f'=(\psi',\theta'):X''\to X'$ um segundo morfismo de espaços
anelados, e seja $f''=f\circ f'$ o morfismo composto.
Levando em conta a compatibilidade do funtor $\psi'^*$ com os
produtos tensoriais \sref[0\textsubscript{I}]{0.4.3.3}, verifica-se
imediatamente que o composto do dihomomorfismo
\[
  \HH^p(X',f^*(\sh{F}))
  \longrightarrow
  \HH^p\bigl(X'',f'^*(f^*(\sh{F}))\bigr)
\]
com \sref{0.12.1.3.5} é o dihomomorfismo correspondente
\[
  \HH^p(X,\sh{F})
  \longrightarrow
  \HH^p(X'',f''^*(\sh{F})).
\]
\end{env}

\begin{env}[12.1.4]
\label{0.12.1.4}
Pode-se dar uma definição direta do homomorfismo
\sref{0.12.1.3.2} como segue.
Tomemos uma resolução injetiva
$\sh{L}^\bullet=(\sh{L}^i)$ de $\sh{F}$ constituída por feixes de
grupos abelianos sobre $X$.
Como o funtor $\psi^*$ é exato,
$\psi^*(\sh{L}^\bullet)$ é uma resolução de $\psi^*(\sh{F})$
constituída por feixes sobre $X'$.
Se $\sh{L}'^\bullet=(\sh{L}'^i)$ é uma resolução injetiva de
$\psi^*(\sh{F})$ na categoria de feixes de grupos abelianos sobre
$X'$, existe um morfismo
\[
  \psi^*(\sh{L}^\bullet)\longrightarrow\sh{L}'^\bullet
\]
de complexos de feixes de grupos abelianos, compatível com os
aumentos (M, V, 1.1.a), e determinado salvo homotopia.
Obtêm-se assim homomorfismos
\[
  \Gamma(X,\sh{L}^\bullet)
  \longrightarrow
  \Gamma(X',\psi^*(\sh{L}^\bullet))
  \longrightarrow
  \Gamma(X',\sh{L}'^\bullet)
\]
de complexos de grupos abelianos.
Seu composto, ao passar à cohomologia, dá um morfismo de
$\delta$-funtores
\[
  \HH^\bullet(X,\sh{F})
  \longrightarrow
  \HH^\bullet(X',\psi^*(\sh{F}));
\]
como coincide com \sref{0.12.1.3.2} em grau $0$, é idêntico a ele
(T, 2.2).

Consideremos agora um recobrimento aberto
$\mathfrak{U}=(U_\alpha)$ de $X$, e seja
$\mathfrak{U}'=(f^{-1}(U_\alpha))$ o recobrimento aberto por imagens inversas
de $X'$.
Para todo subconjunto aberto $V$ de $X$, os homomorfismos canônicos
\[
  \Gamma(V,\sh{F})
  \longrightarrow
  \Gamma(f^{-1}(V),f^*(\sh{F}))
\]
definem imediatamente (cf.\ G, II, 5.1) um homomorfismo de complexos
\[
  \check{\mathrm{C}}^\bullet(\mathfrak{U},\sh{F})
  \longrightarrow
  \check{\mathrm{C}}^\bullet(\mathfrak{U}',f^*(\sh{F})),
\]
e, portanto, homomorfismos canônicos
\[
\label{0.12.1.4.1}
\tag{12.1.4.1}
  \theta_p:
  \check{\HH}^p(\mathfrak{U},\sh{F})
  \longrightarrow
  \check{\HH}^p(\mathfrak{U}',f^*(\sh{F})).
\]

\oldpage[0\textsubscript{III}]{52}
Além disso, existem diagramas comutativos
\[
\label{0.12.1.4.2}
\tag{12.1.4.2}
  \xymatrix{
    \check{\HH}^p(\mathfrak{U},\sh{F})
      \ar[r]^-{\theta_p}\ar[d]
      & \check{\HH}^p(\mathfrak{U}',f^*(\sh{F}))\ar[d]\\
    \HH^p(X,\sh{F})
      \ar[r]_-{\theta_p}
      & \HH^p(X',f^*(\sh{F})),
  }
\]
onde as setas verticais são os homomorfismos canônicos de
(G, II, 5.2).
Para demonstrar a comutatividade de \sref{0.12.1.4.2}, consideremos o
complexo de feixes de cocadeias alternadas de $\sh{F}$ relativo a
$\mathfrak{U}$,
$\sh{C}^\bullet(\mathfrak{U},\sh{F})$, para o qual
\[
  \Gamma\bigl(X,\sh{C}^\bullet(\mathfrak{U},\sh{F})\bigr)
  =
  \check{\mathrm{C}}^\bullet(\mathfrak{U},\sh{F})
\]
(G, II, 5.2).
Os homomorfismos canônicos
\[
  \Gamma(V,\sh{F})
  \longrightarrow
  \Gamma(f^{-1}(V),\psi^*(\sh{F}))
\]
definem então um $\psi$-morfismo
\[
  \sh{C}^\bullet(\mathfrak{U},\sh{F})
  \longrightarrow
  \sh{C}^\bullet(\mathfrak{U}',\psi^*(\sh{F})),
\]
e, com a notação precedente, existe um diagrama comutativo
\[
  \xymatrix{
    \Gamma\bigl(X,\sh{C}^\bullet(\mathfrak{U},\sh{F})\bigr)
      \ar[r]\ar[d]
      & \Gamma\bigl(X',
          \sh{C}^\bullet(\mathfrak{U}',\psi^*(\sh{F}))\bigr)\ar[d]\\
    \Gamma(X,\sh{L}^\bullet)\ar[r]
      & \Gamma(X',\psi^*(\sh{L}^\bullet))\ar[r]
      & \Gamma(X',\sh{L}'^\bullet).
  }
\]
Ao passar à cohomologia, obtêm-se diagramas comutativos
\[
  \xymatrix{
    \check{\HH}^p(\mathfrak{U},\sh{F})\ar[r]\ar[d]
      & \check{\HH}^p(\mathfrak{U}',\psi^*(\sh{F}))\ar[d]\\
    \HH^p(X,\sh{F})\ar[r]
      & \HH^p(X',\psi^*(\sh{F})),
  }
\]
onde as setas verticais são os homomorfismos canônicos de
(G, II, 5.2).
Basta combinar estes diagramas com os diagramas comutativos
\[
  \xymatrix{
    \check{\HH}^p(\mathfrak{U}',\psi^*(\sh{F}))\ar[r]\ar[d]
      & \check{\HH}^p(\mathfrak{U}',f^*(\sh{F}))\ar[d]\\
    \HH^p(X',\psi^*(\sh{F}))\ar[r]
      & \HH^p(X',f^*(\sh{F})),
  }
\]
\oldpage[0\textsubscript{III}]{53}
que provêm do homomorfismo
$\psi^*(\sh{F})\to f^*(\sh{F})$ e da funtorialidade dos
homomorfismos canônicos de (G, II, 5.2), para obter a comutatividade
de \sref{0.12.1.4.2}.

Observemos que se
$A=\Gamma(X,\sh{O}_X)$ e $A'=\Gamma(X',\sh{O}_{X'})$, então
\sref{0.12.1.4.1} é um dihomomorfismo de módulos correspondente aos
anéis $A$ e $A'$.
Existe para \sref{0.12.1.4.1} uma propriedade de transitividade com respeito
a um composto de dois morfismos, análoga à de
\sref{0.12.1.3.5}.
Finalmente, nas definições precedentes, em lugar de uma resolução
injetiva $\sh{L}^\bullet$ de $\sh{F}$, poder-se-ia ter partido igualmente
de uma resolução tal que
\[
  \HH^p(X,\sh{L}^i)=0
  \qquad\text{para todo $i$ e todo $p>0$}
\]
(G, II, 4.7.1).
\end{env}

\begin{env}[12.1.5]
\label{0.12.1.5}
Sejam $\sh{F},\sh{G},\sh{H}$ três $\sh{O}_X$-módulos, e consideremos
um $\sh{O}_X$-homomorfismo
\[
  u:\sh{F}\otimes_{\sh{O}_X}\sh{G}\longrightarrow\sh{H},
\]
que dá, em cohomologia, os homomorfismos
\[
\label{0.12.1.5.1}
\tag{12.1.5.1}
  \HH^p(X,\sh{F})\otimes_A\HH^q(X,\sh{G})
  \longrightarrow
  \HH^{p+q}(X,\sh{H})
\]
deduzidos do produto cup \sref{0.12.1.2.1}.
Demonstraremos que, sob as hipóteses e com a notação de
\sref{0.12.1.3}, os diagramas
\[
\label{0.12.1.5.2}
\tag{12.1.5.2}
  \xymatrix@C=10pt{
    \HH^p(X,\sh{F})\otimes_A\HH^q(X,\sh{G})
      \ar[r]\ar[d]
      & \HH^{p+q}(X,\sh{H})\ar[d]\\
    \HH^p(X',f^*(\sh{F}))\otimes_{A'}
      \HH^q(X',f^*(\sh{G}))
      \ar[r]
      & \HH^{p+q}(X',f^*(\sh{H}))
  }
\]
são comutativos, onde as setas verticais provêm dos
homomorfismos canônicos \sref{0.12.1.3.5}.

Lembremos para isso que \sref{0.12.1.5.1} pode ser obtido a partir das
resoluções canônicas (G, II, 4.3)
$\sh{L}^\bullet,\sh{M}^\bullet,\sh{N}^\bullet$ de
$\sh{F},\sh{G},\sh{H}$ respectivamente, que estão constituídas por
$\sh{O}_X$-módulos, e da aplicação linear
\[
  \sh{L}^\bullet\otimes_{\sh{O}_X}\sh{M}^\bullet
  \longrightarrow
  \sh{N}^\bullet
\]
de complexos de $\sh{O}_X$-módulos correspondente a $u$.
Esta aplicação proporciona um homomorfismo de complexos de $A$-módulos
\[
  \Gamma(X,\sh{L}^\bullet)\otimes_A\Gamma(X,\sh{M}^\bullet)
  \longrightarrow
  \Gamma(X,\sh{N}^\bullet),
\]
e, portanto, ao passar à cohomologia, homomorfismos
\[
  \HH^p\bigl(\Gamma(X,\sh{L}^\bullet)\bigr)
  \otimes_A
  \HH^q\bigl(\Gamma(X,\sh{M}^\bullet)\bigr)
  \longrightarrow
  \HH^{p+q}\bigl(\Gamma(X,\sh{N}^\bullet)\bigr)
\]
(G, II, 6.6).
Existe claramente um diagrama comutativo
\[
\label{0.12.1.5.3}
\tag{12.1.5.3}
  \xymatrix@C=8pt{
    \Gamma(X,\sh{L}^\bullet)\otimes_A
      \Gamma(X,\sh{M}^\bullet)
      \ar[r]\ar[d]
      & \Gamma(X,\sh{N}^\bullet)\ar[d]\\
    \Gamma(X',\psi^*(\sh{L}^\bullet))
      \otimes_{\Gamma(X',\psi^*(\sh{O}_X))}
      \Gamma(X',\psi^*(\sh{M}^\bullet))
      \ar[r]
      & \Gamma(X',\psi^*(\sh{N}^\bullet)).
  }
\]
\oldpage[0\textsubscript{III}]{54}
Ao passar à cohomologia, obtêm-se diagramas comutativos
\[
\label{0.12.1.5.4}
\tag{12.1.5.4}
  \xymatrix@C=7pt{
    \HH^p(X,\sh{F})\otimes_A\HH^q(X,\sh{G})
      \ar[r]\ar[d]
      & \HH^{p+q}(X,\sh{H})\ar[d]\\
    \HH^p\bigl(\Gamma(X',\psi^*(\sh{L}^\bullet))\bigr)
      \otimes_{\Gamma(X',\psi^*(\sh{O}_X))}
      \HH^q\bigl(\Gamma(X',\psi^*(\sh{M}^\bullet))\bigr)
      \ar[r]
      & \HH^{p+q}\bigl(\Gamma(X',\psi^*(\sh{N}^\bullet))\bigr).
  }
\]
Mas como
$\psi^*(\sh{L}^\bullet)$,
$\psi^*(\sh{M}^\bullet)$, e
$\psi^*(\sh{N}^\bullet)$
são resoluções de
$\psi^*(\sh{F})$,
$\psi^*(\sh{G})$, e
$\psi^*(\sh{H})$ respectivamente, existe um diagrama comutativo
(G, II, 6.6.1)
\[
\label{0.12.1.5.5}
\tag{12.1.5.5}
  \xymatrix@C=7pt{
    \HH^p\bigl(\Gamma(X',\psi^*(\sh{L}^\bullet))\bigr)
      \otimes_{\Gamma(X',\psi^*(\sh{O}_X))}
      \HH^q\bigl(\Gamma(X',\psi^*(\sh{M}^\bullet))\bigr)
      \ar[r]\ar[d]
      & \HH^{p+q}\bigl(\Gamma(X',\psi^*(\sh{N}^\bullet))\bigr)
        \ar[d]\\
    \HH^p(X',\psi^*(\sh{F}))
      \otimes_{\Gamma(X',\psi^*(\sh{O}_X))}
      \HH^q(X',\psi^*(\sh{G}))
      \ar[r]
      & \HH^{p+q}(X',\psi^*(\sh{H})).
  }
\]
Finalmente, por funtorialidade, existe um diagrama comutativo
\[
\label{0.12.1.5.6}
\tag{12.1.5.6}
  \xymatrix@C=7pt{
    \HH^p(X',\psi^*(\sh{F}))
      \otimes_{\Gamma(X',\psi^*(\sh{O}_X))}
      \HH^q(X',\psi^*(\sh{G}))
      \ar[r]\ar[d]
      & \HH^{p+q}(X',\psi^*(\sh{H}))\ar[d]\\
    \HH^p(X',f^*(\sh{F}))\otimes_{A'}
      \HH^q(X',f^*(\sh{G}))
      \ar[r]
      & \HH^{p+q}(X',f^*(\sh{H})).
  }
\]
Combinando os três diagramas \sref{0.12.1.5.4},
\sref{0.12.1.5.5} e \sref{0.12.1.5.6}, obtém-se o diagrama
comutativo desejado \sref{0.12.1.5.2}.
\end{env}

\begin{remark}[12.1.6]
\label{0.12.1.6}
\oldpage[0\textsubscript{III}]{55}
Com a notação de \sref{0.12.1.3}, suponhamos que existe um
diagrama comutativo
\[
\label{0.12.1.6.1}
\tag{12.1.6.1}
  \xymatrix{
    0\ar[r] & \sh{F}\ar[r]^r\ar[d]_u
      & \sh{G}\ar[r]^s\ar[d]_v
      & \sh{H}\ar[r]\ar[d]_w & 0\\
    0\ar[r] & \sh{F}'\ar[r]_{r'}
      & \sh{G}'\ar[r]_{s'}
      & \sh{H}'\ar[r] & 0,
  }
\]
onde $r,s$ são homomorfismos de $\sh{O}_X$-módulos,
$r',s'$ são homomorfismos de $\sh{O}_{X'}$-módulos,
$u,v,w$ são $f$-morfismos, e as linhas são exatas.
Segue-se que o diagrama
\[
\label{0.12.1.6.2}
\tag{12.1.6.2}
  \xymatrix@C=6pt{
    \cdots\ar[r]
      & \HH^p(X,\sh{F})\ar[r]\ar[d]_{u_p}
      & \HH^p(X,\sh{G})\ar[r]\ar[d]_{v_p}
      & \HH^p(X,\sh{H})\ar[r]^-\partial\ar[d]_{w_p}
      & \HH^{p+1}(X,\sh{F})\ar[r]\ar[d]_{u_{p+1}}
      & \cdots\\
    \cdots\ar[r]
      & \HH^p(X',\sh{F}')\ar[r]
      & \HH^p(X',\sh{G}')\ar[r]
      & \HH^p(X',\sh{H}')\ar[r]_-\partial
      & \HH^{p+1}(X',\sh{F}')\ar[r]
      & \cdots
  }
\]
é comutativo.
Com efeito, \sref{0.12.1.6.1} se fatora como
\[
  \xymatrix@R=14pt{
    0\ar[r] & \sh{F}\ar[r]\ar[d]
      & \sh{G}\ar[r]\ar[d]
      & \sh{H}\ar[r]\ar[d] & 0\\
    0\ar[r] & \psi^*(\sh{F})\ar[r]\ar[d]
      & \psi^*(\sh{G})\ar[r]\ar[d]
      & \psi^*(\sh{H})\ar[r]\ar[d] & 0\\
    0\ar[r] & \sh{F}'\ar[r]
      & \sh{G}'\ar[r]
      & \sh{H}'\ar[r] & 0,
  }
\]
onde a linha intermediária é exata
\sref[0\textsubscript{I}]{0.3.7.2}.
Basta então utilizar o fato de que \sref{0.12.1.3.2} é um
homomorfismo de $\delta$-funtores e que os
$\HH^p(X',\sh{F}')$ formam um $\delta$-funtor em $\sh{F}'$.
\end{remark}

\begin{env}[12.1.7]
\label{0.12.1.7}
Com as hipóteses e a notação de \sref{0.12.1.3}, consideremos agora o
caso
\[
  \sh{F}=f_*(\sh{F}')=\psi_*(\sh{F}').
\]
Veremos que o dihomomorfismo definido em \sref{0.12.1.3},
\[
  \HH^p(X,f_*(\sh{F}'))\longrightarrow\HH^p(X',\sh{F}'),
  \tag{12.1.7.1}
  \label{0.12.1.7.1}
\]
pode ser obtido (salvo um automorfismo de $\HH^p(X',\sh{F}')$) como um
homomorfismo de bordo de uma sucessão espectral para o funtor composto
\[
  \sh{F}'\leadsto\Gamma(X,\psi_*(\sh{F}'))
\]
(T, 2.4).
A descrição de \sref{0.12.1.7.1} dada em \sref{0.12.1.4} mostra
aqui que este homomorfismo pode ser obtido como segue.
Tomemos resoluções injetivas $\sh{L}^\bullet$ e
${\sh{L}'}^\bullet$ de $\psi_*(\sh{F}')$ e $\sh{F}'$,
respectivamente, e depois um homomorfismo de complexos
\[
  v:\psi^*(\sh{L}^\bullet)\longrightarrow{\sh{L}'}^\bullet
\]
``acima'' do homomorfismo canônico
$\psi^*(\psi_*(\sh{F}'))\to\sh{F}'$.
\oldpage[0\textsubscript{III}]{56}
Observemos em seguida que
\[
  \Gamma(X',{\sh{L}'}^\bullet)
  =\Gamma(X,\psi_*({\sh{L}'}^\bullet))
\]
e que o homomorfismo composto
\[
  \Gamma(X,\sh{L}^\bullet)
  \longrightarrow\Gamma(X',\psi^*(\sh{L}^\bullet))
  \xrightarrow{\Gamma(v)}\Gamma(X',{\sh{L}'}^\bullet)
\]
não é senão
\[
  \Gamma(v^\flat):
  \Gamma(X,\sh{L}^\bullet)
  \longrightarrow\Gamma(X,\psi_*({\sh{L}'}^\bullet))
  \tag{12.1.7.2}
  \label{0.12.1.7.2}
\]
\sref[0\textsubscript{I}]{0.3.7.1}; e \sref{0.12.1.7.1} se
obtém passando à cohomologia em \sref{0.12.1.7.2}.

Por outro lado, as sucessões espectrais para o funtor composto
\[
  \sh{F}'\leadsto\Gamma(X,\psi_*(\sh{F}'))
\]
são obtidas tomando uma resolução injetiva de Cartan--Eilenberg
\[
  \sh{M}^{\bullet\bullet}=(\sh{M}^{ij})
\]
do complexo $\psi_*({\sh{L}'}^\bullet)$ na categoria de
feixes de grupos abelianos sobre $X$.
As sucessões espectrais em questão são as do bicomplexo
$\Gamma(X,\sh{M}^{\bullet\bullet})$; são birregulares, pois
$\sh{M}^{ij}=0$ para $i<0$ ou $j<0$.
Ora, a primeira sucessão espectral deste bicomplexo degenera, pois
os feixes $\psi_*({\sh{L}'}^i)$ são flácidos (G, II, 3.1.1), e
portanto
\[
  \HH_\mathrm{II}^q
    \bigl(\Gamma(X,\sh{M}^{i,\bullet})\bigr)
  =\HH^q\bigl(\psi_*({\sh{L}'}^i)\bigr)=0
  \qquad(q>0)
\]
(G, II, 4.4.3).
Por conseguinte, existem homomorfismos de bordo bijetivos
\sref{0.11.1.6}
\[
  {}'E_2^{i0}
  =\HH^i\!\left(
      \HH_\mathrm{II}^0
        \bigl(\Gamma(X,\sh{M}^{\bullet\bullet})\bigr)
    \right)
  \longrightarrow
  \HH^i\bigl(\Gamma(X,\sh{M}^{\bullet\bullet})\bigr),
  \tag{12.1.7.3}
  \label{0.12.1.7.3}
\]
e sabemos por \sref{0.11.3.4} que este homomorfismo provém, ao passar
à cohomologia, do aumento
\[
  \Gamma(X,\psi_*({\sh{L}'}^\bullet))
  \longrightarrow\Gamma(X,\sh{M}^{\bullet\bullet}),
  \tag{12.1.7.4}
  \label{0.12.1.7.4}
\]
que, por sua vez, provém do aumento
\[
  \eta:\psi_*({\sh{L}'}^\bullet)
  \longrightarrow\sh{M}^{\bullet\bullet}.
\]

Por outro lado, para a segunda sucessão espectral, existem homomorfismos
de bordo
\[
  {}''E_2^{i0}
  =\HH^i\!\left(
      \HH_\mathrm{I}^0
        \bigl(\Gamma(X,\sh{M}^{\bullet\bullet})\bigr)
    \right)
  \longrightarrow
  \HH^i\bigl(\Gamma(X,\sh{M}^{\bullet\bullet})\bigr),
  \tag{12.1.7.5}
  \label{0.12.1.7.5}
\]
que provêm, por \sref{0.11.3.4}, ao passar à cohomologia, do
homomorfismo de complexos
\[
  Z_\mathrm{I}^0
    \bigl(\Gamma(X,\sh{M}^{\bullet\bullet})\bigr)
  \longrightarrow\Gamma(X,\sh{M}^{\bullet\bullet}).
\]
Como $\psi_*$ é exato à esquerda, a sucessão
\[
  0\longrightarrow\psi_*(\sh{F}')
  \longrightarrow\psi_*({\sh{L}'}^0)
  \longrightarrow\psi_*({\sh{L}'}^1)
\]
é exata.
Pela definição de uma resolução de Cartan--Eilenberg
\sref{0.11.4.2}, podemos tomar, portanto,
\[
  B_\mathrm{I}^0(\sh{M}^{\bullet\bullet})=0,
  \qquad
  Z_\mathrm{I}^0(\sh{M}^{\bullet\bullet})=\sh{L}^\bullet.
\]
Como o diagrama
\[
  \xymatrix@C=50pt@R=24pt{
    \sh{L}^0\ar[r]^{i^0}
      & \sh{M}^{00}\\
    \psi_*(\sh{F}')\ar[u]^\varepsilon
      \ar[r]_{\varepsilon'}\ar[ur]^{\varepsilon''}
      & \psi_*({\sh{L}'}^0)\ar[u]_{\eta^0}
  }
\]
é comutativo, a injeção de complexos
$i:\sh{L}^\bullet\to\sh{M}^{\bullet\bullet}$ é compatível com os
aumentos $\varepsilon$ e $\varepsilon''$.
Existem, pois, dois homomorfismos de complexos de $\sh{L}^\bullet$
em $\sh{M}^{\bullet\bullet}$,
\[
  \xymatrix@C=52pt@R=22pt{
    \sh{L}^\bullet\ar[r]^i\ar[dr]_{v^\flat}
      & \sh{M}^{\bullet\bullet}\\
    & \psi_*({\sh{L}'}^\bullet)\ar[u]_\eta
  }
\]
\oldpage[0\textsubscript{III}]{57}
compatíveis com os aumentos $\varepsilon$ e
$\varepsilon''$.
Como $\sh{L}^\bullet$ é uma resolução injetiva e
$\sh{M}^{\bullet\bullet}$ está constituído por feixes injetivos, segue-se
de (M, V, 1.1(a)) que estes dois homomorfismos são homotópicos.
O mesmo vale, portanto, para os dois homomorfismos correspondentes
\[
  \Gamma(X,\sh{L}^\bullet)
  \longrightarrow\Gamma(X,\sh{M}^{\bullet\bullet}),
\]
e, ao passar à cohomologia, dão por conseguinte o mesmo
homomorfismo.
Em outras palavras, temos demonstrado que o homomorfismo de bordo
\sref{0.12.1.7.5}, escrito
\[
  \HH^p(X,\psi_*(\sh{F}'))
  \longrightarrow\HH^p\bigl(\Gamma(X,\sh{M}^{\bullet\bullet})\bigr),
\]
é o composto de \sref{0.12.1.7.1} e
\sref{0.12.1.7.3}, estando este último escrito
\[
  \HH^p(X',\sh{F}')
  \longrightarrow\HH^p\bigl(\Gamma(X,\sh{M}^{\bullet\bullet})\bigr)
\]
e tendo-se visto que é um isomorfismo.
Isto demonstra a asserção.
\end{env}

\subsection{Imagens diretas superiores}
\label{subsection:0.12.2}

\begin{env}[12.2.1]
\label{0.12.2.1}
Sejam $(X,\sh{O}_X)$ e $(Y,\sh{O}_Y)$ dois espaços anelados, e seja
$f=(\psi,\theta)$ um morfismo de $X$ em $Y$.  Esse morfismo define o funtor imagem
direta
\[
  f_*:\cat{C}(X)\longrightarrow\cat{C}(Y),
\]
que é além disso idêntico à restrição a $\cat{C}(X)$ do
funtor $\psi_*$ definido sobre a categoria de feixes de grupos abelianos
em $X$.
Este último funtor é aditivo e exato à esquerda; e como todo feixe de
grupos abelianos sobre $X$ é isomorfo a um subfeixe de um feixe injetivo
de grupos abelianos, definem-se os funtores derivados à direita
\[
  \sh{F}\longmapsto\RR^p\psi_*(\sh{F})
\]
de $\psi_*$.  Os $\RR^p\psi_*(\sh{F})$ são feixes de grupos abelianos
sobre $Y$, e os $\RR^p\psi_*$ formam um funtor cohomológico universal
(T, 2.3).

Além disso, $\RR^p\psi_*(\sh{F})$ é o feixe associado ao
pré-feixe
\[
  V\longmapsto\HH^p(f^{-1}(V),\sh{F})
\]
(T, 3.7.2).
Suponhamos agora que $\sh{F}$ é um $\sh{O}_X$-módulo.
Então $\HH^p(f^{-1}(V),\sh{F})$ está dotado naturalmente da estrutura
de $\Gamma(f^{-1}(V),\sh{O}_X)$-módulo, e portanto da de
$\Gamma(V,\psi_*(\sh{O}_X))$-módulo; o homomorfismo
\[
  \theta:\sh{O}_Y\longrightarrow\psi_*(\sh{O}_X)
\]
dota-o então da estrutura de
$\Gamma(V,\sh{O}_Y)$-módulo.
Para as estruturas assim definidas, a restrição de um subconjunto aberto
$V$ a um subconjunto aberto $V'\subset V$ define um dihomomorfismo.  Isto
define, portanto, sobre todo $\RR^p\psi_*(\sh{F})$ a estrutura de um
$\sh{O}_Y$-módulo.  Denotamos este $\sh{O}_Y$-módulo por
$\RR^p f_*(\sh{F})$; assim, $\RR^p f_*$ fica definido como um funtor
aditivo de $\cat{C}(X)$ em $\cat{C}(Y)$.

Além disso, os $\RR^p f_*$ formam um $\delta$-funtor.  Com efeito, se
\[
  0\longrightarrow\sh{F}'\longrightarrow\sh{F}
  \longrightarrow\sh{F}''\longrightarrow0
\]
é uma sucessão exata de $\sh{O}_X$-módulos, a descrição precedente
dos $\RR^p\psi_*$ e da estrutura de $\sh{O}_Y$-módulo sobre
$\RR^p\psi_*(\sh{F})$ mostra imediatamente que o homomorfismo de conexão
\[
  \partial:\RR^p\psi_*(\sh{F}'')
  \longrightarrow\RR^{p+1}\psi_*(\sh{F}')
\]
é um homomorfismo de $\sh{O}_Y$-módulos.
Finalmente, os $\RR^p f_*$ identificam-se com os funtores derivados à
direita de $f_*$.  Todo $\sh{O}_X$-módulo admite uma resolução injetiva por
$\sh{O}_X$-módulos, e tal resolução está constituída por feixes flácidos
de grupos abelianos (G, II, 7.1); pode, portanto, ser utilizada para calcular
os $\RR^p\psi_*(\sh{F})$, pois $\RR^n\psi_*(\sh{G})=0$ para $n\geq1$
para todo feixe flácido $\sh{G}$ (T, 2.4.1, Observação~3, e o
 corolário da Proposição~\EGAUnavailableHyperref{0.3.3.2}{3.3.2}).
Segue-se que os $\RR^p f_*$ formam um funtor cohomológico universal
de $\cat{C}(X)$ em $\cat{C}(Y)$ (T, 2.3).
\end{env}

\begin{env}[12.2.2]
\label{0.12.2.2}
Sejam $\sh{F}$ e $\sh{G}$ dois $\sh{O}_X$-módulos.
Com a notação de \sref{0.12.2.1}, para todo subconjunto aberto $V$ de
$Y$ existe o homomorfismo de produto cup \sref{0.12.1.2.1}
\[
  \begin{aligned}
  &\HH^p(f^{-1}(V),\sh{F})
    \otimes_{\Gamma(f^{-1}(V),\sh{O}_X)}
    \HH^q(f^{-1}(V),\sh{G})\\
  &\hspace{55mm}\longrightarrow
    \HH^{p+q}\bigl(f^{-1}(V),
      \sh{F}\otimes_{\sh{O}_X}\sh{G}\bigr).
  \end{aligned}
\]
\oldpage[0\textsubscript{III}]{58}
Segue-se imediatamente da definição do produto cup
(G, II, 6.6) que estes homomorfismos comutam com a restrição de
$V$ a um subespaço aberto $V'$ de $V$.
Por outro lado, $\theta$ dá um homomorfismo de anéis
\[
  \Gamma(V,\sh{O}_Y)
  \longrightarrow\Gamma(V,\psi_*(\sh{O}_X))
  =\Gamma(f^{-1}(V),\sh{O}_X),
\]
e, portanto, um homomorfismo canônico de produtos tensoriais
\[
  \begin{aligned}
  &\HH^p(f^{-1}(V),\sh{F})
    \otimes_{\Gamma(V,\sh{O}_Y)}
    \HH^q(f^{-1}(V),\sh{G})\\
  &\quad\longrightarrow
    \HH^p(f^{-1}(V),\sh{F})
    \otimes_{\Gamma(f^{-1}(V),\sh{O}_X)}
    \HH^q(f^{-1}(V),\sh{G}),
  \end{aligned}
\]
que é igualmente compatível com a restrição de $V$ a $V'$.
Por composição, obtém-se assim um homomorfismo de
$\Gamma(V,\sh{O}_Y)$-módulos.  Para os feixes associados aos
pré-feixes em questão, define um homomorfismo canônico,
funtorial em $\sh{F}$ e $\sh{G}$,
\[
  \RR^p f_*(\sh{F})\otimes_{\sh{O}_Y}\RR^q f_*(\sh{G})
  \longrightarrow
  \RR^{p+q}f_*
    \bigl(\sh{F}\otimes_{\sh{O}_X}\sh{G}\bigr).
  \tag{12.2.2.1}
  \label{0.12.2.2.1}
\]
Para $p=q=0$, este homomorfismo reduz-se a
\sref[0\textsubscript{I}]{0.4.2.2.1}.
\end{env}

\begin{proposition}[12.2.3]
\label{0.12.2.3}
Para todo $\sh{O}_X$-módulo $\sh{F}$ e todo $\sh{O}_Y$-módulo
$\sh{L}$ localmente livre de posto finito, existem isomorfismos canônicos
funtoriais
\[
  \RR^p f_*(\sh{F})\otimes_{\sh{O}_Y}\sh{L}
  \simto
  \RR^p f_*
    \bigl(\sh{F}\otimes_{\sh{O}_X}f^*(\sh{L})\bigr).
  \tag{12.2.3.1}
  \label{0.12.2.3.1}
\]
\end{proposition}

\begin{proof}
O homomorfismo \sref{0.12.2.3.1} é obtido compondo o
seguinte caso particular de \sref{0.12.2.2.1},
\[
  \RR^p f_*(\sh{F})
    \otimes_{\sh{O}_Y}f_*\bigl(f^*(\sh{L})\bigr)
  \longrightarrow
  \RR^p f_*
    \bigl(\sh{F}\otimes_{\sh{O}_X}f^*(\sh{L})\bigr),
  \tag{12.2.3.2}
  \label{0.12.2.3.2}
\]
com o homomorfismo do primeiro membro de
\sref{0.12.2.3.1} ao de \sref{0.12.2.3.2} induzido pelo
homomorfismo canônico \sref[0\textsubscript{I}]{0.4.4.3.2}.
Para verificar que \sref{0.12.2.3.1} é um isomorfismo quando $\sh{L}$ é
localmente livre, pode reduzir-se imediatamente ao caso
$\sh{L}=\sh{O}_Y$, pois a questão é local sobre $Y$ e os
funtores em questão são aditivos em $\sh{L}$.
Nesse caso, em vista da definição de \sref{0.12.2.2.1}, a
proposição reduz-se a verificar que o homomorfismo correspondente
de pré-feixes é bijetivo; isto segue-se imediatamente de
$f^*(\sh{O}_Y)=\sh{O}_X$.
\end{proof}

\begin{env}[12.2.4]
\label{0.12.2.4}
Seja $(Z,\sh{O}_Z)$ um terceiro espaço anelado, e seja
$g:Y\to Z$ um morfismo de espaços anelados.
Sabemos (G, II, 7.1 e 3.1.1) que, para todo
$\sh{O}_X$-módulo injetivo $\sh{G}$, $f_*(\sh{G})$ é um feixe flácido de
grupos abelianos.  Por conseguinte, por \sref{0.12.2.1},
\[
  \RR^p g_*\bigl(f_*(\sh{G})\bigr)=0
  \qquad(p>0).
\]
Segue-se (T, 2.4.1) que a sucessão espectral de Leray para funtores
compostos se aplica ao composto $g_*\circ f_*$.  Existe uma
sucessão espectral birregular cujo termo limite é o funtor
$\RR^\bullet h_*$, onde $h=g\circ f$, e cujo termo $E_2$ é
\[
  E_2^{pq}=\RR^p g_*\bigl(\RR^q f_*(\sh{F})\bigr).
  \tag{12.2.4.1}
  \label{0.12.2.4.1}
\]
\end{env}

\begin{env}[12.2.5]
\label{0.12.2.5}
Sob as condições de \sref{0.12.2.4}, definiremos diretamente
homomorfismos canônicos de $\sh{O}_Z$-módulos
\[
  \RR^n g_*\bigl(f_*(\sh{F})\bigr)
  \longrightarrow\RR^n h_*(\sh{F})
  \tag{12.2.5.1}
  \label{0.12.2.5.1}
\]
e
\[
  \RR^n h_*(\sh{F})
  \longrightarrow g_*\bigl(\RR^n f_*(\sh{F})\bigr),
  \tag{12.2.5.2}
  \label{0.12.2.5.2}
\]
\oldpage[0\textsubscript{III}]{59}
que podem identificar-se com os ``homomorfismos de bordo'' da sucessão
espectral de Leray (cf.\ \sref{0.12.1.7}).
Basta trabalhar com os pré-feixes aos quais estão associados os feixes de
imagens diretas superiores \sref{0.12.2.1}.
Para isso, seja $W$ um subconjunto aberto qualquer de $Z$, e consideremos
sua imagem inversa $g^{-1}(W)$ em $Y$.  Existe um
dihomomorfismo canônico
\[
  \begin{aligned}
  \HH^n\bigl(g^{-1}(W),f_*(\sh{F})\bigr)
  &\longrightarrow
  \HH^n\bigl(f^{-1}(g^{-1}(W)),
    f^*(f_*(\sh{F}))\bigr),
  \end{aligned}
  \tag{12.2.5.3}
  \label{0.12.2.5.3}
\]
sendo os anéis correspondentes
$\Gamma(g^{-1}(W),\sh{O}_Y)$ e
$\Gamma(h^{-1}(W),\sh{O}_X)$.
Por outro lado, o homomorfismo canônico
\sref[0\textsubscript{I}]{0.4.4.3.3} dá, por funtorialidade,
homomorfismos canônicos
\[
  \HH^n\bigl(h^{-1}(W),f^*(f_*(\sh{F}))\bigr)
  \longrightarrow
  \HH^n\bigl(h^{-1}(W),\sh{F}\bigr),
  \tag{12.2.5.4}
  \label{0.12.2.5.4}
\]
que são homomorfismos de
$\Gamma(h^{-1}(W),\sh{O}_X)$-módulos.
Levando em conta o homomorfismo de anéis
\[
  \Gamma(W,\sh{O}_Z)
  \longrightarrow\Gamma(g^{-1}(W),\sh{O}_Y),
\]
vemos que, compondo \sref{0.12.2.5.3} com
\sref{0.12.2.5.4}, se obtém um homomorfismo de pré-feixes e, portanto, o
homomorfismo de feixes \sref{0.12.2.5.1}.

A definição de \sref{0.12.2.5.2} é ainda mais simples.
Por definição, $\RR^n h_*(\sh{F})$ está associado ao pré-feixe
\[
  W\longmapsto
  \HH^n\bigl(f^{-1}(g^{-1}(W)),\sh{F}\bigr),
\]
enquanto que $\RR^n f_*(\sh{F})$ está associado ao pré-feixe
\[
  V\longmapsto\HH^n(f^{-1}(V),\sh{F}).
\]
Existe, portanto, um homomorfismo canônico
\[
  \HH^n\bigl(f^{-1}(g^{-1}(W)),\sh{F}\bigr)
  \longrightarrow
  \Gamma\bigl(g^{-1}(W),\RR^n f_*(\sh{F})\bigr),
\]
e é imediato que estes homomorfismos definem um homomorfismo de
pré-feixes, que, por sua vez, define \sref{0.12.2.5.2}.
\end{env}

\begin{env}[12.2.6]
\label{0.12.2.6}
Sob as hipóteses de \sref{0.12.2.4}, sejam $\sh{F}$, $\sh{G}$
e $\sh{K}$ três $\sh{O}_X$-módulos, e seja
\[
  u:\sh{F}\otimes_{\sh{O}_X}\sh{G}\longrightarrow\sh{K}
\]
um $\sh{O}_X$-homomorfismo.
Então os diagramas seguintes são comutativos:
\[
  \xymatrix@C=10pt@R=24pt{
    \RR^p g_*(f_*(\sh{F}))
      \otimes_{\sh{O}_Z}\RR^q g_*(f_*(\sh{G}))
      \ar[r]\ar[d]
      & \RR^{p+q}g_*(f_*(\sh{K}))\ar[d]\\
    \RR^p h_*(\sh{F})
      \otimes_{\sh{O}_Z}\RR^q h_*(\sh{G})
      \ar[r]
      & \RR^{p+q}h_*(\sh{K}),
  }
  \tag{12.2.6.1}
  \label{0.12.2.6.1}
\]
e
\[
  \xymatrix@C=10pt@R=24pt{
    \RR^p h_*(\sh{F})
      \otimes_{\sh{O}_Z}\RR^q h_*(\sh{G})
      \ar[r]\ar[d]
      & \RR^{p+q}h_*(\sh{K})\ar[d]\\
    g_*(\RR^p f_*(\sh{F}))
      \otimes_{\sh{O}_Z}g_*(\RR^q f_*(\sh{G}))
      \ar[r]
      & g_*(\RR^{p+q}f_*(\sh{K})).
  }
  \tag{12.2.6.2}
  \label{0.12.2.6.2}
\]
\oldpage[0\textsubscript{III}]{60}
Aqui as setas horizontais provêm de \sref{0.12.2.2.1} (a última
combinada com \sref[0\textsubscript{I}]{0.4.2.2.1}), e as setas
verticais de \sref{0.12.2.5.1} e \sref{0.12.2.5.2}, respectivamente.

Com efeito, basta verificar a asserção para os homomorfismos de pré-feixes
correspondentes.  Voltando às definições de
\sref{0.12.2.2} e \sref{0.12.2.5}, reduz-se imediatamente, para
\sref{0.12.2.6.1}, aos diagramas comutativos
\sref{0.12.1.5.2}; a verificação é ainda mais simples para
\sref{0.12.2.6.2}.
\end{env}

\subsection{Suplementos sobre os funtores Ext de feixes}
\label{subsection:0.12.3}

\begin{env}[12.3.1]
\label{0.12.3.1}
Seja $(X,\sh{O}_X)$ um espaço anelado.
Não recordaremos a definição nem as propriedades principais dos
bifuntores
\[
  \Ext^p_{\sh{O}_X}(X;\sh{F},\sh{G}),
  \qquad
  \shExt^p_{\sh{O}_X}(\sh{F},\sh{G}),
\]
o primeiro com valores na categoria de
$\Gamma(X,\sh{O}_X)$-módulos e o segundo na categoria de
$\sh{O}_X$-módulos, nem a sucessão espectral birregular
$E(\sh{F},\sh{G})$ que os relaciona (T, 4.2 e G, II, 7.3).
\end{env}

\begin{env}[12.3.2]
\label{0.12.3.2}
Assim como em (M, XIV, 1), define-se a noção de uma
\emph{extensão} de um $\sh{O}_X$-módulo $\sh{F}$ por um
$\sh{O}_X$-módulo $\sh{G}$ e a lei de composição sobre as classes de
extensões equivalentes: os argumentos utilizados para os módulos
transportam-se imediatamente para uma categoria abeliana arbitrária.
A segunda demonstração de (M, XIV, 1.1), que somente utiliza a existência de
imersões em objetos injetivos, continua sendo, portanto, válida na
categoria de $\sh{O}_X$-módulos.  Mostra que
\[
  \Ext^1_{\sh{O}_X}(X;\sh{F},\sh{G})
\]
identifica-se canonicamente com o grupo abeliano de classes de
extensões de $\sh{F}$ por $\sh{G}$.
\end{env}

\begin{proposition}[12.3.3]
\label{0.12.3.3}
Seja $(X,\sh{O}_X)$ um espaço anelado cujo feixe de anéis
$\sh{O}_X$ é coerente.
Então, para todo par de $\sh{O}_X$-módulos coerentes $\sh{F}$ e
$\sh{G}$ e todo $p\geq0$,
\[
  \shExt^p_{\sh{O}_X}(\sh{F},\sh{G})
\]
é um $\sh{O}_X$-módulo coerente.
\end{proposition}

\begin{proof}
Os $\shExt^p_{\sh{O}_X}(\sh{F},\sh{G})$ formam um funtor cohomológico
contravariante em $\sh{F}$.
Como $\sh{F}$ é coerente, para todo $p$ e todo ponto $x\in X$
existem uma vizinhança aberta $U$ de $x$ em $X$ e uma sucessão exata
de $(\sh{O}_X|U)$-módulos
\[
  0\longrightarrow\sh{R}\longrightarrow\sh{L}_{p-1}
  \longrightarrow\cdots\longrightarrow\sh{L}_0
  \longrightarrow\sh{F}|U\longrightarrow0,
\]
onde cada $\sh{L}_i$ ($0\leq i\leq p-1$) é isomorfo a
$\sh{O}_X^{n_i}|U$, e $\sh{R}$ é coerente.
Isto decorre por indução sobre $p$ de
\sref[0\textsubscript{I}]{0.5.3.2} e
\sref[0\textsubscript{I}]{0.5.3.4}, pois $\sh{O}_X$ é coerente.

Além disso, para $p\geq1$,
\[
  \shExt^p_{\sh{O}_X|U}(\sh{L}|U,\sh{G}|U)=0
\]
para todo $\sh{O}_X$-módulo $\sh{L}$ tal que $\sh{L}|U$ seja
isomorfo a $\sh{O}_X^n|U$ (T, 4.2.3).
O argumento de (M, V, 7.2) se aplica, portanto, ao funtor cohomológico
contravariante
\[
  \sh{F}\longmapsto
  \shHom_{\sh{O}_X|U}(\sh{F}|U,\sh{G}|U)
\]
e dá uma sucessão exata
\[
  \begin{aligned}
  \shHom_{\sh{O}_X|U}(\sh{L}_{p-1},\sh{G}|U)
  &\longrightarrow
  \shHom_{\sh{O}_X|U}(\sh{R},\sh{G}|U)\\
  &\longrightarrow
  \shExt^p_{\sh{O}_X|U}(\sh{F}|U,\sh{G}|U)
  \longrightarrow0.
  \end{aligned}
\]
Os dois primeiros termos são $(\sh{O}_X|U)$-módulos coerentes
\sref[0\textsubscript{I}]{0.5.3.5}; portanto, também o é o terceiro
\sref[0\textsubscript{I}]{0.5.3.4}.
\end{proof}

\oldpage[0\textsubscript{III}]{61}
\begin{proposition}[12.3.4]
\label{0.12.3.4}
Seja $f:X\to Y$ um morfismo plano de espaços anelados, e sejam
$\sh{F}$ e $\sh{G}$ dois $\sh{O}_Y$-módulos.
\begin{enumerate}
  \item[\rm(i)] Existe um homomorfismo de bifuntores cohomológicos
  \[
    f^*\bigl(\shExt^\bullet_{\sh{O}_Y}(\sh{F},\sh{G})\bigr)
    \longrightarrow
    \shExt^\bullet_{\sh{O}_X}
      \bigl(f^*(\sh{F}),f^*(\sh{G})\bigr),
    \tag{12.3.4.1}
    \label{0.12.3.4.1}
  \]
  que, em grau $0$, se reduz ao homomorfismo canônico
  \sref[0\textsubscript{I}]{0.4.4.6}.

  \item[\rm(ii)] Existe um morfismo canônico de sucessões espectrais
  \[
    E(\sh{F},\sh{G})
    \longrightarrow E\bigl(f^*(\sh{F}),f^*(\sh{G})\bigr),
    \tag{12.3.4.2}
    \label{0.12.3.4.2}
  \]
  cuja aplicação sobre os termos $E_2$ é dada pelos homomorfismos
  \[
    \begin{aligned}
    \HH^p\bigl(Y,\shExt^q_{\sh{O}_Y}(\sh{F},\sh{G})\bigr)
    \longrightarrow
    \HH^p\bigl(X,\shExt^q_{\sh{O}_X}
      (f^*(\sh{F}),f^*(\sh{G}))\bigr),
    \end{aligned}
    \tag{12.3.4.3}
    \label{0.12.3.4.3}
  \]
  deduzidos de \sref{0.12.3.4.1} e \sref{0.12.1.3.1}.
\end{enumerate}
\end{proposition}

\begin{proof}
\begin{enumerate}
  \item[\rm(i)]
  Como $f^*$ é exato sobre a categoria de $\sh{O}_Y$-módulos
  \sref[0\textsubscript{I}]{0.6.7.2}, os funtores
  \[
    \sh{G}\longmapsto
      f^*\bigl(\shHom_{\sh{O}_Y}(\sh{F},\sh{G})\bigr),
    \qquad
    \sh{G}\longmapsto
      \shHom_{\sh{O}_X}\bigl(f^*(\sh{F}),f^*(\sh{G})\bigr)
  \]
  são exatos à esquerda.  O homomorfismo canônico
  \sref[0\textsubscript{I}]{0.4.4.6} induz, portanto, canonicamente um
  homomorfismo entre seus funtores derivados.

  Para calcular estes funtores derivados, tomemos uma resolução injetiva
  $\sh{L}^\bullet=(\sh{L}^i)$ de $\sh{G}$.  Têm-se então
  homomorfismos
  \[
    \begin{aligned}
    \sh{H}^p\!\left(
      f^*\bigl(\shHom_{\sh{O}_Y}
        (\sh{F},\sh{L}^\bullet)\bigr)\right)
    \longrightarrow
    \sh{H}^p\!\left(
      \shHom_{\sh{O}_X}
        (f^*(\sh{F}),f^*(\sh{L}^\bullet))\right)
    \end{aligned}
  \]
  entre os feixes de cohomologia de complexos.
  Pela exatidão de $f^*$, o primeiro membro é
  \[
    f^*\bigl(\shExt^p_{\sh{O}_Y}(\sh{F},\sh{G})\bigr).
  \]
  A exatidão de $f^*$ implica também que
  $f^*(\sh{L}^\bullet)$ é uma resolução de $f^*(\sh{G})$.
  Se ${\sh{L}'}^\bullet=({\sh{L}'}^i)$ é uma resolução injetiva
  de $f^*(\sh{G})$ na categoria de $\sh{O}_X$-módulos, existe um
  homomorfismo de complexos
  \[
    f^*(\sh{L}^\bullet)\longrightarrow{\sh{L}'}^\bullet,
  \]
  determinado salvo homotopia, e portanto um homomorfismo bem definido em
  cohomologia.  Seu composto com o homomorfismo precedente é
  \sref{0.12.3.4.1}.

  \item[\rm(ii)]
  Com a notação precedente, existe um homomorfismo de complexos de
  $\sh{O}_X$-módulos
  \[
    f^*\bigl(\shHom_{\sh{O}_Y}(\sh{F},\sh{L}^\bullet)\bigr)
    \longrightarrow
    \shHom_{\sh{O}_X}
      \bigl(f^*(\sh{F}),{\sh{L}'}^\bullet\bigr).
  \]
  Seja $\sh{M}^{\bullet\bullet}$ uma resolução injetiva de
  Cartan--Eilenberg de
  $\shHom_{\sh{O}_Y}(\sh{F},\sh{L}^\bullet)$ na categoria de
  $\sh{O}_Y$-módulos.  Como $f^*$ é exato,
  $f^*(\sh{M}^{\bullet\bullet})$ é uma resolução de Cartan--Eilenberg de
  $f^*(\shHom_{\sh{O}_Y}(\sh{F},\sh{L}^\bullet))$.
  Se ${\sh{M}'}^{\bullet\bullet}$ é uma resolução injetiva de
  Cartan--Eilenberg de
  $\shHom_{\sh{O}_X}(f^*(\sh{F}),{\sh{L}'}^\bullet)$, existe
  \sref{0.11.4.2} um homomorfismo, determinado salvo homotopia,
  \[
    f^*(\sh{M}^{\bullet\bullet})
    \longrightarrow{\sh{M}'}^{\bullet\bullet},
  \]
  compatível com o homomorfismo precedente; em outras palavras, existe
  um $f$-morfismo
  $\sh{M}^{\bullet\bullet}\to{\sh{M}'}^{\bullet\bullet}$ de
  bicomplexos de feixes.
  Induz um dihomomorfismo
  \[
    \Gamma(Y,\sh{M}^{\bullet\bullet})
    \longrightarrow\Gamma(X,{\sh{M}'}^{\bullet\bullet})
  \]
  de bicomplexos de módulos, determinado salvo homotopia, e portanto
  um morfismo bem definido de sucessões espectrais
  \sref{0.11.3.2}.  Este é precisamente o morfismo desejado
  \sref{0.12.3.4.2}, e a descrição
  \sref{0.12.3.4.3} decorre imediatamente das definições.
\end{enumerate}
\end{proof}

\begin{proposition}[12.3.5]
\label{0.12.3.5}
Sob as hipóteses de \sref{0.12.3.4}, suponhamos além disso que o
feixe de anéis $\sh{O}_Y$ é coerente.
Então, para todo $\sh{O}_Y$-módulo coerente $\sh{F}$, os
homomorfismos canônicos \sref{0.12.3.4.1} são bijetivos.
\end{proposition}

\begin{proof}
\oldpage[0\textsubscript{III}]{62}
Como a questão é local sobre $Y$, podemos supor que existe uma
sucessão exata
\[
  0\longrightarrow\sh{R}\longrightarrow\sh{O}_Y^n
  \longrightarrow\sh{F}\longrightarrow0.
\]
Então $\sh{R}$ é também um $\sh{O}_Y$-módulo coerente
\sref[0\textsubscript{I}]{0.5.3.4}.
Para demonstrar que os homomorfismos
\[
  f^*\bigl(\shExt^p_{\sh{O}_Y}(\sh{F},\sh{G})\bigr)
  \longrightarrow
  \shExt^p_{\sh{O}_X}
    \bigl(f^*(\sh{F}),f^*(\sh{G})\bigr)
\]
são bijetivos, raciocinamos por indução sobre $p$.
Para $p=0$, a asserção é
\sref[0\textsubscript{I}]{0.6.7.6.1}.
Existe um diagrama comutativo com linhas exatas
\[
{\tiny
  \xymatrix@C=2pt@R=24pt{
  f^*\!\left(\shExt^{p-1}_{\sh{O}_Y}(\sh{O}_Y^n,\sh{G})\right)
    \ar[r]\ar[d]
  & f^*\!\left(\shExt^{p-1}_{\sh{O}_Y}(\sh{R},\sh{G})\right)
    \ar[r]^\partial\ar[d]
  & f^*\!\left(\shExt^p_{\sh{O}_Y}(\sh{F},\sh{G})\right)
    \ar[r]\ar[d]
  & f^*\!\left(\shExt^p_{\sh{O}_Y}(\sh{O}_Y^n,\sh{G})\right)
    \ar[d]\\
  \shExt^{p-1}_{\sh{O}_X}(\sh{O}_X^n,f^*(\sh{G}))
    \ar[r]
  & \shExt^{p-1}_{\sh{O}_X}(f^*(\sh{R}),f^*(\sh{G}))
    \ar[r]^\partial
  & \shExt^p_{\sh{O}_X}(f^*(\sh{F}),f^*(\sh{G}))
    \ar[r]
  & \shExt^p_{\sh{O}_X}(\sh{O}_X^n,f^*(\sh{G})).
  }}
\]
Aqui temos utilizado $f^*(\sh{O}_Y)=\sh{O}_X$; como $f^*$ é exato,
ambas as linhas são exatas.
Além disso,
\[
  \shExt^p_{\sh{O}_Y}(\sh{O}_Y^n,\sh{G})=0,
  \qquad
  \shExt^p_{\sh{O}_X}(\sh{O}_X^n,f^*(\sh{G}))=0
  \qquad(p>0)
\]
(T, 4.2.3).
Pela hipótese de indução, as duas primeiras setas verticais do
diagrama precedente são isomorfismos, enquanto que os termos da
direita são nulos.  Segue-se que
\[
  f^*\bigl(\shExt^p_{\sh{O}_Y}(\sh{F},\sh{G})\bigr)
  \longrightarrow
  \shExt^p_{\sh{O}_X}
    \bigl(f^*(\sh{F}),f^*(\sh{G})\bigr)
\]
é um isomorfismo.
\end{proof}

\subsection{Hipercohomologia do funtor imagem direta}
\label{subsection:0.12.4}

\begin{env}[12.4.1]
\label{0.12.4.1}
Sejam $(X,\sh{O}_X)$ e $(Y,\sh{O}_Y)$ dois espaços anelados, e seja
$f:X\longrightarrow Y$ um morfismo de espaços anelados.
Pode-se formar a hipercohomologia de $f_*$ com respeito a um complexo
arbitrário de $\sh{O}_X$-módulos
$\sh{K}^\bullet=(\sh{K}^i)_{i\in\bb{Z}}$
\sref{0.11.4.4}, pois os limites indutivos filtrantes existem e são exatos
na categoria abeliana de $\sh{O}_Y$-módulos (T, 3.1.1).
Os $\sh{O}_Y$-módulos de hipercohomologia
$\RR^p f_*(\sh{K}^\bullet)$ serão também denotados por
$\sh{H}^p(f,\sh{K}^\bullet)$, ou simplesmente
$\sh{H}^p_f(\sh{K}^\bullet)$.
Lembremos que $\sh{H}^\bullet(f,\sh{K}^\bullet)$ é a cohomologia do
bicomplexo de $\sh{O}_Y$-módulos $f_*(\sh{L}^{\bullet\bullet})$,
onde $\sh{L}^{\bullet\bullet}$ é uma resolução injetiva de
Cartan--Eilenberg de $\sh{K}^\bullet$ na categoria de
$\sh{O}_X$-módulos.
É o termo limite de duas sucessões espectrais cujos termos $E_2$
são dados por
\begin{equation}
\label{0.12.4.1.1}
\tag{12.4.1.1}
  {}'E_2^{pq}
  =\sh{H}^p\bigl(\sh{H}^q(f,\sh{K}^\bullet)\bigr)
\end{equation}
e
\begin{equation}
\label{0.12.4.1.2}
\tag{12.4.1.2}
  {}''E_2^{pq}
  =\sh{H}^p\bigl(f,\sh{H}^q(\sh{K}^\bullet)\bigr)
  \quad
  \bigl(=\RR^p f_*(\sh{H}^q(\sh{K}^\bullet))\bigr).
\end{equation}
Nestas fórmulas adotamos a notação geral $T(A^\bullet)$
para o transformado de um complexo por um funtor \sref{0.11.2.1}, e, para
um $\sh{O}_X$-módulo $\sh{F}$, escrevemos
$\sh{H}^p(f,\sh{F})$ em lugar de $\RR^p f_*(\sh{F})$.
Lembremos também que a primeira sucessão espectral é sempre regular.
Ambas as sucessões espectrais são birregulares quando $\sh{K}^\bullet$ está
limitado inferiormente,
\oldpage[0\textsubscript{III}]{63}
ou quando existe um inteiro $m$ tal que todo $\sh{O}_X$-módulo
admite uma resolução flácida de comprimento $\leq m$ \sref{0.11.4.4}.
\end{env}

\begin{env}[12.4.2]
\label{0.12.4.2}
Denotaremos analogamente por $\bb{H}^\bullet(X,\sh{K}^\bullet)$ a
hipercohomologia do funtor $\Gamma$ com respeito a um complexo
$\sh{K}^\bullet$ de $\sh{O}_X$-módulos.
Assim, os $\bb{H}^p(X,\sh{K}^\bullet)$ são
$\Gamma(X,\sh{O}_X)$-módulos.
Além disso, $\bb{H}^\bullet(X,\sh{K}^\bullet)$ pode ser considerado um
caso particular de $\sh{H}^\bullet(f,\sh{K}^\bullet)$, onde $f$ é um
morfismo de $(X,\sh{O}_X)$ a um espaço anelado constituído por um ponto
dotado do anel $\Gamma(X,\sh{O}_X)$.

Para todo subconjunto aberto $V$ de $X$, escreveremos
$\bb{H}^\bullet(V,\sh{K}^\bullet)$ em lugar de
$\bb{H}^\bullet(V,\sh{K}^\bullet|V)$.
\end{env}

\begin{proposition}[12.4.3]
\label{0.12.4.3}
Para todo inteiro $p\in\bb{Z}$, o $\sh{O}_Y$-módulo
$\sh{H}^p(f,\sh{K}^\bullet)$ é canonicamente isomorfo ao feixe
associado ao pré-feixe
\[
  U\longmapsto
  \bb{H}^p\bigl(f^{-1}(U),\sh{K}^\bullet\bigr)
\]
sobre $Y$.
\end{proposition}

\begin{proof}
Com a notação de \sref{0.12.4.1}, o feixe de cohomologia
$\sh{H}^p(f_*(\sh{L}^{\bullet\bullet}))$ está associado ao
pré-feixe
\[
  U\longmapsto
  \HH^p\bigl(\Gamma(U,f_*(\sh{L}^{\bullet\bullet}))\bigr)
  =
  \HH^p\bigl(\Gamma(f^{-1}(U),\sh{L}^{\bullet\bullet})\bigr).
\]
Mas está claro que
$\sh{L}^{\bullet\bullet}|f^{-1}(U)$ é uma resolução injetiva de
Cartan--Eilenberg de $\sh{K}^\bullet|f^{-1}(U)$ (T, 3.1.3).
Por conseguinte,
\[
  \HH^p\bigl(\Gamma(f^{-1}(U),\sh{L}^{\bullet\bullet})\bigr)
  =
  \bb{H}^p\bigl(f^{-1}(U),\sh{K}^\bullet\bigr)
\]
por definição.
\end{proof}

\begin{proposition}[12.4.4]
\label{0.12.4.4}
A hipercohomologia $\sh{H}^\bullet(f,\sh{K}^\bullet)$ é um
funtor cohomológico de $\sh{K}^\bullet$ em cada um dos casos
seguintes:
\begin{enumerate}
  \item[\rm{a)}] $\sh{K}^\bullet$ varia na categoria de complexos
    limitados inferiormente;
  \item[\rm{b)}] existe um inteiro $m$ tal que todo
    $\sh{O}_X$-módulo admite uma resolução flácida de comprimento $\leq m$;
  \item[\rm{c)}] $X$ é um espaço noetheriano.
\end{enumerate}
\end{proposition}

\begin{proof}
Os casos~a) e~b) são casos particulares de \sref{0.11.5.4}.
Por outro lado, o caso~c) segue-se de \sref{0.11.5.2}, pois, neste
caso, o funtor $f_*$ comuta com os limites indutivos
(G, II, 3.10.1).
\end{proof}

\begin{env}[12.4.5]
\label{0.12.4.5}
Seja agora $\mathfrak{U}=(U_\alpha)$ um recobrimento aberto de $X$.
Para todo complexo de \emph{pré-feixes}
$\sh{K}^\bullet=(\sh{K}^j)$ sobre $X$, consideremos o bicomplexo
$\check{\mathrm{C}}^\bullet(\mathfrak{U},\sh{K}^\bullet)$ cuja
componente de bigrado $(i,j)$ é
$\check{\mathrm{C}}^i(\mathfrak{U},\sh{K}^j)$, o grupo de
$i$-cocadeias alternadas do nervo de $\mathfrak{U}$ com valores
em $\sh{K}^j$ (G, II, 5.1).
Chamamos a cohomologia deste bicomplexo de
\emph{hipercohomologia do recobrimento $\mathfrak{U}$ com
coeficientes em $\sh{K}^\bullet$}, e a denotamos por
\[
  \bb{H}^\bullet(\mathfrak{U},\sh{K}^\bullet)
  =
  \HH^\bullet
  \bigl(\check{\mathrm{C}}^\bullet
    (\mathfrak{U},\sh{K}^\bullet)\bigr).
\]
A sucessão espectral de Leray de um recobrimento
(T, 3.8.1 e G, II, 5.9.1) se generaliza à hipercohomologia como segue.
\end{env}

\begin{proposition}[12.4.6]
\label{0.12.4.6}
Seja $\sh{K}^\bullet=(\sh{K}^i)$ um complexo de
$\sh{O}_X$-módulos.
Existe um funtor espectral regular de $\sh{K}^\bullet$ com
termo limite $\bb{H}^\bullet(X,\sh{K}^\bullet)$ e termo $E_2$
\begin{equation}
\label{0.12.4.6.1}
\tag{12.4.6.1}
  E_2^{pq}
  =
  \CHH^p\bigl(\mathfrak{U},h^q(\sh{K}^\bullet)\bigr),
\end{equation}
onde $h^q(\sh{K}^\bullet)$ denota o complexo de pré-feixes
\[
  V\longmapsto\HH^q(V,\sh{K}^\bullet)
\]
sobre $X$.
Esta sucessão espectral é birregular se $\sh{K}^\bullet$ está limitado
inferiormente.
\end{proposition}

\begin{proof}
Seja $\sh{L}^{\bullet\bullet}$ uma resolução injetiva de
Cartan--Eilenberg de $\sh{K}^\bullet$, e consideremos o tricomplexo
\[
  \check{\mathrm{C}}^\bullet
  (\mathfrak{U},\sh{L}^{\bullet\bullet})
  =
  \bigl(\check{\mathrm{C}}^i
    (\mathfrak{U},\sh{L}^{jk})\bigr).
\]
Consideremos primeiro este tricomplexo como um bicomplexo nos graus $i$ e
$j+k$.
Como $i$ somente toma valores não negativos, a segunda sucessão espectral
deste bicomplexo é regular \sref{0.11.3.3} e degenerada, porque
\[
  \CHH^q(\mathfrak{U},\sh{L}^{jk})=0
  \qquad(q>0),
\]
sendo os $\sh{O}_X$-módulos $\sh{L}^{jk}$ feixes flácidos
\oldpage[0\textsubscript{III}]{64}
(G, II, 5.2.3).
Segue-se que existe um isomorfismo canônico
\[
  \HH^n\bigl(
    \check{\mathrm{C}}^\bullet
      (\mathfrak{U},\sh{L}^{\bullet\bullet})
  \bigr)
  \isoto
  \HH^n\bigl(\Gamma(X,\sh{L}^{\bullet\bullet})\bigr)
\]
\sref{0.11.1.6}, em virtude de (G, II, 5.2.2), e portanto, pela
definição \sref{0.12.4.2}, um isomorfismo
\[
  \HH^n\bigl(
    \check{\mathrm{C}}^\bullet
      (\mathfrak{U},\sh{L}^{\bullet\bullet})
  \bigr)
  \isoto
  \bb{H}^n(X,\sh{K}^\bullet).
\]

Consideremos agora
$\check{\mathrm{C}}^\bullet
(\mathfrak{U},\sh{L}^{\bullet\bullet})$
como um bicomplexo nos graus $i+j$ e $k$.
Como $k$ somente toma valores não negativos, a primeira sucessão espectral
deste bicomplexo é sempre regular.
É birregular se $\sh{L}^{jk}=0$ para $j<j_0$, isto é, quando
$\sh{K}^\bullet$ está limitado inferiormente \sref{0.11.3.3}.
Esta é a sucessão espectral desejada.
Com efeito, para todo $j$, $\sh{L}^{j,\bullet}$ é uma resolução
injetiva de $\sh{K}^j$; por conseguinte, quando $i$ varia,
\[
  \HH^q\bigl(
    \check{\mathrm{C}}^i
      (\mathfrak{U},\sh{L}^{j,\bullet})
  \bigr)
\]
é precisamente o complexo de cocadeias
$\check{\mathrm{C}}^i(\mathfrak{U},h^q(\sh{K}^j))$.
Isto demonstra a proposição.
\end{proof}

\begin{corollary}[12.4.7]
\label{0.12.4.7}
Se, para todo simplexo $\sigma$ do nervo de $\mathfrak{U}$ e
todo inteiro $i$, tem-se
\[
  \HH^q(U_\sigma,\sh{K}^i)=0
  \qquad(q>0),
\]
então existe um isomorfismo canônico
\begin{equation}
\label{0.12.4.7.1}
\tag{12.4.7.1}
  \bb{H}^\bullet(\mathfrak{U},\sh{K}^\bullet)
  \isoto
  \bb{H}^\bullet(X,\sh{K}^\bullet).
\end{equation}
\end{corollary}

\begin{proof}
A hipótese implica que
$\check{\mathrm{C}}^\bullet
(\mathfrak{U},h^q(\sh{K}^j))=0$ para $q>0$, e portanto que
$E_2^{pq}=0$ para $q>0$.
Como a sucessão espectral \sref{0.12.4.6.1} é degenerada e
regular, a conclusão segue-se da definição
\sref{0.12.4.5} de
$\bb{H}^\bullet(\mathfrak{U},\sh{K}^\bullet)$ e de
\sref{0.11.1.6}.
\end{proof}

\begin{env}[12.4.8]
\label{0.12.4.8}
Seja $(X',\sh{O}_{X'})$ um segundo espaço anelado, e seja
$f=(\psi,\theta)$ um morfismo de $X'$ em $X$.
Pelo mesmo método que em \sref{0.12.1.3} e \sref{0.12.1.4}, se
define um dihomomorfismo para a hipercohomologia de um complexo
$\sh{K}^\bullet$ de $\sh{O}_X$-módulos:
\begin{equation}
\label{0.12.4.8.1}
\tag{12.4.8.1}
  \bb{H}^p(X,\sh{K}^\bullet)
  \longrightarrow
  \bb{H}^p(X',f^*(\sh{K}^\bullet)).
\end{equation}
Partindo de uma resolução injetiva de Cartan--Eilenberg
$\sh{L}^{\bullet\bullet}$ de $\sh{K}^\bullet$, observemos que, como
$\psi^*$ é exato, $\psi^*(\sh{L}^{\bullet\bullet})$ é uma
resolução de Cartan--Eilenberg de $\psi^*(\sh{K}^\bullet)$ na
categoria de $\psi^*(\sh{O}_X)$-módulos.
Existe, portanto, um morfismo
\[
  \psi^*(\sh{L}^{\bullet\bullet})
  \longrightarrow
  {\sh{L}'}^{\bullet\bullet},
\]
onde ${\sh{L}'}^{\bullet\bullet}$ é uma resolução injetiva de
Cartan--Eilenberg de $\psi^*(\sh{K}^\bullet)$.
Induz um morfismo em cohomologia
\[
  \bb{H}^\bullet(X,\sh{K}^\bullet)
  \longrightarrow
  \bb{H}^\bullet
  (X',\psi^*(\sh{K}^\bullet)).
\]
Compondo-o com o morfismo induzido por funtorialidade a partir de
$\psi^*(\sh{K}^\bullet)\longrightarrow f^*(\sh{K}^\bullet)$ obtém-se
o morfismo desejado \sref{0.12.4.8.1}.

Partindo de \sref{0.12.4.8.1} e \sref{0.12.4.3}, e raciocinando como
em \sref{0.12.2.5}, podem-se então definir, para dois morfismos de
espaços anelados
\[
  f:X\longrightarrow Y,
  \qquad
  g:Y\longrightarrow Z,
\]
com $h=g\circ f$, homomorfismos para a hipercohomologia de um complexo
$\sh{K}^\bullet$ de $\sh{O}_X$-módulos:
\begin{equation}
\label{0.12.4.8.2}
\tag{12.4.8.2}
  \sh{H}^n(g,f_*(\sh{K}^\bullet))
  \longrightarrow
  \sh{H}^n(h,\sh{K}^\bullet),
\end{equation}
e
\begin{equation}
\label{0.12.4.8.3}
\tag{12.4.8.3}
  \sh{H}^n(h,\sh{K}^\bullet)
  \longrightarrow
  g_*\bigl(\sh{H}^n(f,\sh{K}^\bullet)\bigr).
\end{equation}
Deixamos ao leitor os detalhes das definições.
\end{env}

\section{Limites projetivos em álgebra homológica}
\label{section:0.13}

\subsection{A condição de Mittag-Leffler}
\label{subsection:0.13.1}

\begin{env}[13.1.1]
\label{0.13.1.1}
Seja $\cat{C}$ uma categoria abeliana na qual existam os produtos
infinitos (axioma AB~3* de T, 1.5).
Então existe o ínfimo de uma família de subobjetos de um objeto de
$\cat{C}$, e todo sistema projetivo de objetos de $\cat{C}$ possui um
limite projetivo; este último é um funtor exato à esquerda do
sistema projetivo em questão (T, 1.8).
\oldpage[0\textsubscript{III}]{65}
Seja $(A_\alpha,f_{\alpha\beta})$ um sistema projetivo de objetos de
$\cat{C}$ cujo conjunto de índices $\mathrm{I}$ está filtrado à direita.
Ponhamos
\[
  A=\varprojlim A_\alpha,
\]
e, para cada $\alpha\in\mathrm{I}$, seja
$f_\alpha:A\longrightarrow A_\alpha$ o morfismo canônico.
Para todo $\alpha\in\mathrm{I}$, os subobjetos
$f_{\alpha\beta}(A_\beta)$, para $\alpha\leq\beta$, formam uma família
filtrada decrescente de subobjetos de $A_\alpha$.
O subobjeto
\[
  A'_\alpha
  =
  \inf_{\beta\geq\alpha}f_{\alpha\beta}(A_\beta)
\]
chama-se o subobjeto das \emph{imagens universais} em $A_\alpha$.
Está claro que
\[
  f_\alpha(A)\subset A'_\alpha,
  \qquad
  f_{\alpha\beta}(A'_\beta)\subset A'_\alpha
  \quad(\alpha\leq\beta).
\]
Assim,
$(A'_\alpha,f_{\alpha\beta}|A'_\beta)$ é um sistema projetivo e
\[
  A=\varprojlim A'_\alpha.
\]
\end{env}

\begin{env}[13.1.2]
\label{0.13.1.2}
Para um sistema projetivo $(A_\alpha,f_{\alpha\beta})$ em $\cat{C}$,
chama-se \emph{condição de Mittag-Leffler} à seguinte:
\begin{itemize}
  \item[\textbf{(ML)}]
  \emph{Para todo índice $\alpha$, existe um $\beta\geq\alpha$ tal
  que, para todo $\gamma\geq\beta$,}
  \[
    f_{\alpha\gamma}(A_\gamma)
    =
    f_{\alpha\beta}(A_\beta).
  \]
\end{itemize}

Está claro que a condição (ML) é satisfeita se todos os
$f_{\alpha\beta}$ são epimorfismos.
Reciprocamente, suponhamos que (ML) seja satisfeita, e para cada
$\alpha\in\mathrm{I}$ seja $A'_\alpha$ o subobjeto das imagens
universais em $A_\alpha$.
Então, para $\alpha\leq\beta$, a restrição de
$f_{\alpha\beta}$ a $A'_\beta$ é um \emph{epimorfismo}
\[
  A'_\beta\longrightarrow A'_\alpha.
\]
Com efeito, escolhamos $\gamma\geq\beta$ tal que
\[
  f_{\beta\delta}(A_\delta)
  =
  f_{\beta\gamma}(A_\gamma)
  \qquad(\delta\geq\gamma).
\]
Então $A'_\beta=f_{\beta\gamma}(A_\gamma)$, e portanto
\[
  f_{\alpha\beta}(A'_\beta)
  =
  f_{\alpha\gamma}(A_\gamma),
  \qquad
  A'_\alpha
  =
  f_{\alpha\gamma}(A_\gamma)
  =
  f_{\alpha\beta}(A'_\beta).
\]

A condição (ML) também é satisfeita quando os objetos $A_\alpha$ são
\emph{artinianos} em $\cat{C}$, isto é, quando toda família de
subobjetos de $A_\alpha$ possui um elemento minimal.
Com efeito, um elemento minimal da família filtrada decrescente
$\bigl(f_{\alpha\beta}(A_\beta)\bigr)$ de subobjetos de $A_\alpha$
é necessariamente o menor desses subobjetos.
\end{env}

\begin{env}[13.1.3]
\label{0.13.1.3}
\emph{Observação} (13.1.3).
A condição (ML) também pode ser formulada quando $\cat{C}$ é, por
exemplo, a categoria de conjuntos.
Pode-se definir novamente o subconjunto das imagens universais em
$A_\alpha$, e as observações feitas em \sref{0.13.1.1} e
\sref{0.13.1.2} continuam válidas.
\end{env}

\subsection{A condição de Mittag-Leffler para grupos abelianos}
\label{subsection:0.13.2}

\begin{proposition}[13.2.1]
\label{0.13.2.1}
Seja
\[
  0\longrightarrow A_\alpha
  \xrightarrow{u_\alpha}B_\alpha
  \xrightarrow{v_\alpha}C_\alpha
  \longrightarrow0
\]
uma sucessão exata de sistemas projetivos de grupos abelianos, relativos
ao mesmo conjunto de índices filtrado $\mathrm{I}$.
\begin{enumerate}
  \item[{\rm(i)}]
  Se $(B_\alpha)$ satisfaz {\rm(ML)}, então também $(C_\alpha)$.
  \item[{\rm(ii)}]
  Se $(A_\alpha)$ e $(C_\alpha)$ satisfazem {\rm(ML)}, então também
  $(B_\alpha)$.
\end{enumerate}
\end{proposition}

\begin{proof}
Sejam $(f_{\alpha\beta})$, $(g_{\alpha\beta})$ e
$(h_{\alpha\beta})$ os sistemas de homomorfismos que definem
$(A_\alpha)$, $(B_\alpha)$ e $(C_\alpha)$, respectivamente.

(i) Suponhamos que
\[
  g_{\alpha\beta}(B_\beta)
  =
  g_{\alpha\lambda}(B_\lambda)
  \qquad(\lambda\geq\beta).
\]
Como $v_\beta$ e $v_\lambda$ são sobrejetivos,
\[
  h_{\alpha\beta}(C_\beta)
  =
  v_\alpha\bigl(g_{\alpha\beta}(B_\beta)\bigr)
  =
  v_\alpha\bigl(g_{\alpha\lambda}(B_\lambda)\bigr)
  =
  h_{\alpha\lambda}(C_\lambda)
\]
para $\lambda\geq\beta$.

(ii) Seja $\alpha\in\mathrm{I}$.
Escolhamos $\beta\geq\alpha$ tal que
\[
  f_{\alpha\beta}(A_\beta)
  =
  f_{\alpha\lambda}(A_\lambda)
  \qquad(\lambda\geq\beta),
\]
e escolhamos depois $\gamma\geq\beta$ tal que
\[
  h_{\beta\gamma}(C_\gamma)
  =
  h_{\beta\lambda}(C_\lambda)
  \qquad(\lambda\geq\gamma).
\]
Seja $y_\alpha\in g_{\alpha\gamma}(B_\gamma)$, escrevamos
$y_\alpha=g_{\alpha\gamma}(y_\gamma)$ com $y_\gamma\in B_\gamma$,
e ponhamos $y_\beta=g_{\beta\gamma}(y_\gamma)$.
Então
\[
  v_\beta(y_\beta)
  =
  h_{\beta\gamma}\bigl(v_\gamma(y_\gamma)\bigr).
\]
Para todo $\lambda\geq\gamma$, existe, por hipótese, um
$y_\lambda\in B_\lambda$ tal que
\[
  h_{\beta\gamma}\bigl(v_\gamma(y_\gamma)\bigr)
  =
  h_{\beta\lambda}\bigl(v_\lambda(y_\lambda)\bigr).
\]
Daí,
$v_\beta(y_\beta-g_{\beta\lambda}(y_\lambda))=0$, e portanto
existe um $x_\beta\in A_\beta$
\oldpage[0\textsubscript{III}]{66}
tal que
\[
  y_\beta
  =
  g_{\beta\lambda}(y_\lambda)+u_\beta(x_\beta).
\]
Segue-se que
\[
  y_\alpha
  =
  g_{\alpha\lambda}(y_\lambda)
  +
  u_\alpha\bigl(f_{\alpha\beta}(x_\beta)\bigr).
\]
Como $\lambda\geq\beta$, existe um $x_\lambda\in A_\lambda$ tal
que
\[
  f_{\alpha\beta}(x_\beta)=f_{\alpha\lambda}(x_\lambda).
\]
Por conseguinte,
\[
  y_\alpha
  =
  g_{\alpha\lambda}
  \bigl(y_\lambda+u_\lambda(x_\lambda)\bigr)
  \in g_{\alpha\lambda}(B_\lambda),
\]
o que demonstra a asserção.
\end{proof}

\begin{proposition}[13.2.2]
\label{0.13.2.2}
Seja $\mathrm{I}$ um conjunto ordenado filtrado que possui um subconjunto
cofinal enumerável, e seja
\[
  0\longrightarrow A_\alpha
  \xrightarrow{u_\alpha}B_\alpha
  \xrightarrow{v_\alpha}C_\alpha
  \longrightarrow0
\]
uma sucessão exata de sistemas projetivos de grupos abelianos indexados
por $\mathrm{I}$.
Se $(A_\alpha)$ satisfaz {\rm(ML)}, então a sucessão
\[
  0\longrightarrow\varprojlim A_\alpha
  \longrightarrow\varprojlim B_\alpha
  \longrightarrow\varprojlim C_\alpha
  \longrightarrow0
\]
é exata.
\end{proposition}

\begin{proof}
Resta apenas demonstrar que o homomorfismo
\[
  v=\varprojlim v_\alpha:
  \varprojlim B_\alpha\longrightarrow\varprojlim C_\alpha
\]
é sobrejetivo.
Seja $z=(z_\alpha)\in\varprojlim C_\alpha$, e ponhamos
\[
  E_\alpha=v_\alpha^{-1}(z_\alpha).
\]
Os $E_\alpha$ formam manifestamente um sistema projetivo de conjuntos não
vazios mediante as restrições dos homomorfismos
$g_{\alpha\beta}:B_\beta\longrightarrow B_\alpha$.

Demonstremos que este sistema projetivo satisfaz (ML).
Identificando $A_\alpha$ com um subgrupo de $B_\alpha$ por meio de
$u_\alpha$, escolhamos, para todo $\alpha\in\mathrm{I}$, um índice
$\beta\geq\alpha$ tal que
\[
  g_{\alpha\beta}(A_\beta)
  =
  g_{\alpha\lambda}(A_\lambda)
  \qquad(\lambda\geq\beta).
\]
Afirmamos que também
\[
  g_{\alpha\beta}(E_\beta)
  =
  g_{\alpha\lambda}(E_\lambda)
  \qquad(\lambda\geq\beta).
\]
Com efeito, tomemos $y_\lambda\in E_\lambda$ e ponhamos
\[
  y_\beta=g_{\beta\lambda}(y_\lambda),
  \qquad
  y_\alpha=g_{\alpha\lambda}(y_\lambda).
\]
Seja $y'_\alpha=g_{\alpha\beta}(y'_\beta)$ para algum
$y'_\beta\in E_\beta$.
Então $y'_\beta-y_\beta=x_\beta\in A_\beta$, e por hipótese existe
um $x_\lambda\in A_\lambda$ tal que
\[
  g_{\alpha\beta}(x_\beta)
  =
  g_{\alpha\lambda}(x_\lambda).
\]
Assim,
\[
  y'_\alpha
  =
  g_{\alpha\beta}(y_\beta)+g_{\alpha\beta}(x_\beta)
  =
  g_{\alpha\lambda}(y_\lambda+x_\lambda)
  \in g_{\alpha\lambda}(E_\lambda),
\]
o que demonstra a afirmação.

Sabe-se (Bourbaki, \emph{Top.\ g\'{e}n.}, cap.~II,
\textsection3, teorema~1) que, sob as hipóteses impostas a
$\mathrm{I}$, um sistema projetivo de conjuntos não vazios que satisfaz (ML)
possui um limite projetivo não vazio.
Existe, portanto, um ponto
$y=(y_\alpha)\in\varprojlim E_\alpha$.
Como $v_\alpha(y_\alpha)=z_\alpha$ para todo $\alpha$, tem-se
$z=v(y)$.
\end{proof}

\begin{proposition}[13.2.3]
\label{0.13.2.3}
Sob as hipóteses sobre $\mathrm{I}$ de \sref{0.13.2.2}, seja
$(\mathrm{K}^\bullet_\alpha)_{\alpha\in\mathrm{I}}$ um sistema projetivo
de complexos de grupos abelianos
\[
  \mathrm{K}^\bullet_\alpha
  =
  (\mathrm{K}^n_\alpha)_{n\in\bb{Z}},
\]
cujas diferenciais têm grau $+1$.
Para todo $n$, existe um homomorfismo canônico funtorial
\begin{equation}
\label{0.13.2.3.1}
\tag{13.2.3.1}
  h_n:
  \HH^n\bigl(\varprojlim\mathrm{K}^\bullet_\alpha\bigr)
  \longrightarrow
  \varprojlim\HH^n(\mathrm{K}^\bullet_\alpha).
\end{equation}
Se, para todo grau $n$, o sistema projetivo de grupos abelianos
$(\mathrm{K}^n_\alpha)_{\alpha\in\mathrm{I}}$ satisfaz {\rm(ML)},
então todo $h_n$ é sobrejetivo.
Se, além disso, o sistema projetivo
$(\HH^{n-1}(\mathrm{K}^\bullet_\alpha))_{\alpha\in\mathrm{I}}$
satisfaz {\rm(ML)}, então $h_n$ é bijetivo.
\end{proposition}

\begin{proof}
Ponhamos
\[
  \mathrm{K}^n=\varprojlim\mathrm{K}^n_\alpha
  \qquad(n\in\bb{Z}).
\]
Os homomorfismos $h_n$ estão definidos pelos diagramas comutativos
\[
  \xymatrix{
    \cdots\ar[r]
      & \mathrm{K}^{n-1}\ar[r]\ar[d]
      & \mathrm{K}^n\ar[r]\ar[d]
      & \mathrm{K}^{n+1}\ar[r]\ar[d]
      & \cdots\\
    \cdots\ar[r]
      & \mathrm{K}^{n-1}_\alpha\ar[r]
      & \mathrm{K}^n_\alpha\ar[r]
      & \mathrm{K}^{n+1}_\alpha\ar[r]
      & \cdots,
  }
\]
\oldpage[0\textsubscript{III}]{67}
sendo as diferenciais em $\mathrm{K}^\bullet$ os limites projetivos
das diferenciais correspondentes nos
$\mathrm{K}^\bullet_\alpha$.
Consideremos as sucessões exatas
\[
  (*_n)\qquad
  0\longrightarrow
  \mathrm{B}^n(\mathrm{K}^\bullet_\alpha)
  \longrightarrow
  \mathrm{Z}^n(\mathrm{K}^\bullet_\alpha)
  \longrightarrow
  \HH^n(\mathrm{K}^\bullet_\alpha)
  \longrightarrow0
\]
e
\[
  (**_n)\qquad
  0\longrightarrow
  \mathrm{Z}^{n-1}(\mathrm{K}^\bullet_\alpha)
  \longrightarrow
  \mathrm{K}^{n-1}_\alpha
  \longrightarrow
  \mathrm{B}^n(\mathrm{K}^\bullet_\alpha)
  \longrightarrow0.
\]
A hipótese e a Proposição~\sref{0.13.2.1}~(i) mostram que o
sistema projetivo
$(\mathrm{B}^n(\mathrm{K}^\bullet_\alpha))_{\alpha\in\mathrm{I}}$
satisfaz (ML) para todo $n$.
Segue-se de \sref{0.13.2.2} que
\[
  (***_n)\qquad
  0\longrightarrow
  \varprojlim_\alpha\mathrm{B}^n(\mathrm{K}^\bullet_\alpha)
  \longrightarrow
  \varprojlim_\alpha\mathrm{Z}^n(\mathrm{K}^\bullet_\alpha)
  \longrightarrow
  \varprojlim_\alpha\HH^n(\mathrm{K}^\bullet_\alpha)
  \longrightarrow0
\]
é exata.
Ora,
$\varprojlim_\alpha\mathrm{B}^n(\mathrm{K}^\bullet_\alpha)$
identifica-se com um subgrupo de $\mathrm{K}^{n}$ que contém
$\mathrm{B}^n(\mathrm{K}^\bullet)$, enquanto que
$\varprojlim_\alpha\mathrm{Z}^n(\mathrm{K}^\bullet_\alpha)$
identifica-se com um subgrupo de $\mathrm{Z}^n(\mathrm{K}^\bullet)$.
Por conseguinte, $h_n$ é sobrejetivo.

Suponhamos agora, além disso, que
$(\HH^{n-1}(\mathrm{K}^\bullet_\alpha))_{\alpha\in\mathrm{I}}$
satisfaz (ML).
As sucessões exatas $(*_{n-1})$ e a
Proposição~\sref{0.13.2.1}~(ii) mostram que
$(\mathrm{Z}^{n-1}(\mathrm{K}^\bullet_\alpha))_{\alpha\in\mathrm{I}}$
satisfaz (ML).
Aplicando \sref{0.13.2.2} às sucessões exatas $(**_n)$ obtém-se uma
sucessão exata
\[
  0\longrightarrow
  \varprojlim_\alpha
    \mathrm{Z}^{n-1}(\mathrm{K}^\bullet_\alpha)
  \longrightarrow
  \mathrm{K}^{n-1}
  \xrightarrow{u}
  \varprojlim_\alpha
    \mathrm{B}^n(\mathrm{K}^\bullet_\alpha)
  \longrightarrow0.
\]
Como
$\varprojlim_\alpha\mathrm{B}^n(\mathrm{K}^\bullet_\alpha)
\supset\mathrm{B}^n(\mathrm{K}^\bullet)$
e o composto de $u$ com a injeção
\[
  \varprojlim_\alpha\mathrm{B}^n(\mathrm{K}^\bullet_\alpha)
  \longrightarrow\mathrm{K}^n
\]
é a diferencial
$\mathrm{K}^{n-1}\longrightarrow\mathrm{K}^n$, a sobrejetividade de
$u$ implica
\[
  \varprojlim_\alpha\mathrm{B}^n(\mathrm{K}^\bullet_\alpha)
  =
  \mathrm{B}^n(\mathrm{K}^\bullet).
\]
Assim, $h_n$ é injetivo.
\end{proof}

\begin{env}[13.2.4]
\label{0.13.2.4}
\emph{Observações} (13.2.4).
\begin{enumerate}
  \item[{\rm(i)}]
  O argumento de \sref{0.13.2.2}
  (cf.\ Bourbaki, \emph{loc.\ cit.}) mostra que a conclusão de
  essa proposição continua válida se supusermos apenas que os
  $A_\alpha$ podem ser dotados de estruturas de espaços metrizáveis completos
  nos quais as translações são homeomorfismos, que as aplicações
  $f_{\alpha\beta}:A_\beta\longrightarrow A_\alpha$ que definem o
  sistema projetivo são uniformemente contínuas para as métricas em
  questão, e que $(A_\alpha)$ satisfaz a condição
  \begin{itemize}
    \item[\textbf{(ML')}]
    \emph{Para todo índice $\alpha$, existe um $\beta\geq\alpha$ tal
    que, para todo $\gamma\geq\beta$,
    $f_{\alpha\gamma}(A_\gamma)$ é denso em
    $f_{\alpha\beta}(A_\beta)$.}
  \end{itemize}

  Isto dá um suplemento análogo a \sref{0.13.2.3}.
  Suponhamos que $\mathrm{K}^n_\alpha=0$ para $n<0$ e todo $\alpha$;
  que
  $(\mathrm{K}^n_\alpha)_{\alpha\in\mathrm{I}}$ satisfaz (ML) para
  $n\geq0$; e que
  $A_\alpha=\HH^0(\mathrm{K}^\bullet_\alpha)$ pode ser dotado de
  estruturas de espaço métrico que tenham as propriedades recém-enunciadas.
  Então as conclusões de \sref{0.13.2.3} não mudam para
  $n\geq2$, e $h_1$ é além disso \emph{bijetivo}.
  Com efeito, o argumento de \sref{0.13.2.2} mostra também que
  $(\mathrm{B}^1(\mathrm{K}^\bullet_\alpha))_{\alpha\in\mathrm{I}}$
  satisfaz (ML), que a sucessão $(***_1)$ é exata e, finalmente,
  que
  \[
    \varprojlim_\alpha
    \mathrm{B}^1(\mathrm{K}^\bullet_\alpha)
    =
    \mathrm{B}^1(\mathrm{K}^\bullet).
  \]
  Isto demonstra, entre outras coisas, as asserções de (T, 3.10.2).

  \item[{\rm(ii)}]
  Podem-se introduzir os funtores derivados à direita
  $\varprojlim^{(n)}$ do funtor $\varprojlim$, e obter enunciados mais
  completos que os precedentes~\cite{III-28}.
\end{enumerate}
\end{env}

\subsection{Aplicação: cohomologia de um limite projetivo de feixes}
\label{subsection:0.13.3}

\begin{proposition}[13.3.1]
\label{0.13.3.1}
\oldpage[0\textsubscript{III}]{68}
Seja $X$ um espaço topológico, seja
$(\sh{F}_k)_{k\in\bb{N}}$ um sistema projetivo de feixes de
grupos abelianos sobre $X$, e ponhamos
$\sh{F}=\varprojlim_k\sh{F}_k$.
Suponhamos que sejam satisfeitas as condições seguintes:
\begin{enumerate}
  \item[{\rm(i)}]
  Existe uma base $\mathfrak{B}$ da topologia de $X$ tal que,
  para todo $U\in\mathfrak{B}$ e todo $i\geq0$, o sistema
  projetivo
  $(\HH^i(U,\sh{F}_k))_{k\in\bb{N}}$ satisfaz (ML).

  \item[{\rm(ii)}]
  Para todo $x\in X$ e todo $i>0$, tem-se
  \[
    \varinjlim_U\bigl(\varprojlim_k
    \HH^i(U,\sh{F}_k)\bigr)=0,
  \]
  onde $U$ percorre as vizinhanças de $x$ pertencentes a
  $\mathfrak{B}$.

  \item[{\rm(iii)}]
  Os homomorfismos
  $u_{hk}:\sh{F}_k\longrightarrow\sh{F}_h$ $(h\leq k)$ que definem o
  sistema projetivo $(\sh{F}_k)$ são sobrejetivos.
\end{enumerate}
Sob estas condições, para todo $i>0$, o homomorfismo canônico
\[
  h_i:\HH^i(X,\sh{F})
  \longrightarrow
  \varprojlim_k\HH^i(X,\sh{F}_k)
\]
é sobrejetivo. Se, além disso, para um valor dado de $i$ o sistema
projetivo
$(\HH^{i-1}(X,\sh{F}_k))_{k\in\bb{N}}$ satisfaz (ML), então $h_i$ é
bijetivo.
\end{proposition}

\begin{proof}
\emph{a)}
Suponhamos primeiro que os $\sh{F}_k$, assim como os núcleos
$\sh{N}_{hk}$ dos $u_{hk}$, são flácidos; mostraremos então que
a condição~{\rm(iii)} do enunciado implica
$\HH^i(X,\sh{F})=0$ para $i>0$.
Basta demonstrar que, para todo subconjunto aberto $U$ de $X$ e todo
recobrimento $\mathfrak{U}$ de $U$ por subconjuntos abertos de $U$, tem-se
\[
  \CHH^i(\mathfrak{U},\sh{F})=0
  \qquad (i>0).
\]
Com efeito, segue-se então primeiro que, para a cohomologia de \v{C}ech,
$\check{\HH}^{\,i}(U,\sh{F})=0$ para todo $i>0$, e portanto, por
(G, II, 5.9.2) aplicado ao conjunto de todos os subconjuntos abertos de
$X$, que $\HH^i(X,\sh{F})=0$ para todo $i>0$.

Como os $\sh{F}_k$ são flácidos, tem-se
$\CHH^i(\mathfrak{U},\sh{F}_k)=0$ para todo $i>0$
(G, II, 5.2.3).
Para cada $k$, consideremos o complexo
$\check{\mathrm{C}}^\bullet(\mathfrak{U},\sh{F}_k)$ de cocadeias
alternadas sobre o nervo do recobrimento $\mathfrak{U}$ (G, II, 5.1);
estes complexos formam manifestamente um sistema projetivo de complexos de
grupos abelianos.
Demonstremos que todas as aplicações
\[
  \check{\mathrm{C}}^\bullet(\mathfrak{U},\sh{F}_k)
  \longrightarrow
  \check{\mathrm{C}}^\bullet(\mathfrak{U},\sh{F}_h)
  \qquad(h\leq k)
\]
são sobrejetivas.
Por definição, basta mostrar que, para todo subconjunto aberto $V$
de $X$, a aplicação
$\Gamma(V,\sh{F}_k)\longrightarrow\Gamma(V,\sh{F}_h)$ é
  sobrejetiva. Mas a sucessão
\[
  0\longrightarrow\sh{N}_{hk}\longrightarrow\sh{F}_k
  \longrightarrow\sh{F}_h\longrightarrow0
\]
é exata por hipótese e dá, portanto, a sucessão exata de
cohomologia
\[
  \Gamma(V,\sh{F}_k)\longrightarrow\Gamma(V,\sh{F}_h)
  \longrightarrow\HH^1(V,\sh{N}_{hk})=0,
\]
pois $\sh{N}_{hk}$ é flácido.

O sistema projetivo
$(\check{\mathrm{C}}^\bullet(\mathfrak{U},\sh{F}_k))_{k\in\bb{N}}$
satisfaz, portanto, (ML), e o mesmo vale para
$(\CHH^i(\mathfrak{U},\sh{F}_k))_{k\in\bb{N}}$ para todo $i\geq0$:
para $i>0$ isto é trivial, enquanto que para $i=0$ segue-se do
precedente, pois
$\CHH^0(\mathfrak{U},\sh{F}_k)=\Gamma(U,\sh{F}_k)$
(G, II, 5.2.2).
Podemos, portanto, aplicar \sref{0.13.2.3}, que dá
\[
  \CHH^i(\mathfrak{U},\sh{F})
  =
  \varprojlim_k\CHH^i(\mathfrak{U},\sh{F}_k)
  =0
  \qquad(i>0).
\]

\emph{b)}
Passemos agora ao caso geral. Para cada $k\in\bb{N}$, consideremos a
\emph{resolução canônica}
\[
  \sh{C}^\bullet(X,\sh{F}_k)
  =
  \bigl(\sh{C}^i(X,\sh{F}_k)\bigr)_{i\geq0}
\]
de $\sh{F}_k$ por feixes flácidos (G, II, 4.3).
Para todo $i\geq0$, está claro que
$(\sh{C}^i(X,\sh{F}_k))_{k\in\bb{N}}$ é um sistema projetivo de
feixes flácidos; demonstremos que satisfaz as condições da
parte~\emph{a)}.
\oldpage[0\textsubscript{III}]{69}
Com efeito, se $\sh{N}_{hk}$ denota o núcleo de $u_{hk}$ para $h\leq k$,
então a sucessão
\[
  0\longrightarrow\sh{N}_{hk}\longrightarrow\sh{F}_k
  \longrightarrow\sh{F}_h\longrightarrow0
\]
é exata por~{\rm(iii)}, e nossa asserção segue-se da exatidão
do funtor
$\sh{A}\longmapsto\sh{C}^i(X,\sh{A})$ (G, II, 4.3).

Ponhamos
\[
  \sh{G}^i=\varprojlim_k\sh{C}^i(X,\sh{F}_k).
\]
Pela parte~\emph{a)}, tem-se
$\HH^j(X,\sh{G}^i)=0$ para $j>0$ e $i\geq0$.
Demonstraremos que
$\sh{G}^\bullet=(\sh{G}^i)_{i\geq0}$ é uma resolução do feixe
$\sh{F}$. Como $\HH^j(X,\sh{G}^i)=0$ para $j>0$, a cohomologia
$\HH^\bullet(X,\sh{F})$ será então igual a
$\HH^\bullet(\Gamma(X,\sh{G}^\bullet))$ (G, II, 4.7.1).

Tomando o limite projetivo das sucessões exatas
\[
  0\longrightarrow\sh{F}_k
  \longrightarrow\sh{C}^0(X,\sh{F}_k)
  \longrightarrow\sh{C}^1(X,\sh{F}_k)
  \longrightarrow\cdots
\]
obtém-se manifestamente um complexo de feixes de grupos abelianos
\[
  0\longrightarrow\sh{F}
  \longrightarrow\sh{G}^0
  \longrightarrow\sh{G}^1
  \longrightarrow\cdots.
\]
Para demonstrar nossa asserção, resta estabelecer que
$\sh{H}^i(\sh{G}^\bullet)=0$ para $i>0$.
Este feixe é gerado pelo pré-feixe
\[
  U\longmapsto\HH^i\bigl(\Gamma(U,\sh{G}^\bullet)\bigr)
\]
(G, II, 4.1).
Ora, o complexo $\Gamma(U,\sh{G}^\bullet)$ é o limite projetivo
do sistema projetivo de complexos de grupos abelianos
\[
  \bigl(\Gamma(U,\sh{C}^\bullet(X,\sh{F}_k))\bigr)_{k\in\bb{N}}
\]
(0\textsubscript{I}, \EGAUnavailableHyperref{0.3.2.6}{3.2.6}).
Vimos na parte~\emph{a)} que, para todo $i\geq0$, as aplicações
\[
  \Gamma(U,\sh{C}^i(X,\sh{F}_k))
  \longrightarrow
  \Gamma(U,\sh{C}^i(X,\sh{F}_h))
  \qquad(h\leq k)
\]
são sobrejetivas. Por outro lado,
\[
  \HH^i(U,\sh{F}_k)
  =
  \HH^i\bigl(\Gamma(U,\sh{C}^\bullet(X,\sh{F}_k))\bigr),
\]
pois a resolução canônica $\sh{C}^\bullet(U,\sh{F}_k|U)$ é
induzida sobre $U$ por $\sh{C}^\bullet(X,\sh{F}_k)$.
Pela hipótese~{\rm(i)}, para todo $U\in\mathfrak{B}$ podemos aplicar
\sref{0.13.2.3} a este sistema projetivo de complexos, e portanto
\[
  \HH^i\bigl(\Gamma(U,\sh{G}^\bullet)\bigr)
  =
  \varprojlim_k\HH^i(U,\sh{F}_k)
  \qquad(i\geq0).
\]
A hipótese~{\rm(ii)} mostra então, por definição, que os feixes
$\sh{H}^i(\sh{G}^\bullet)$ se anulam para $i>0$.

Assim, para todo $i\geq0$,
\[
  \HH^i(X,\sh{F})
  =
  \HH^i\bigl(\Gamma(X,\sh{G}^\bullet)\bigr),
  \qquad
  \Gamma(X,\sh{G}^\bullet)
  =
  \varprojlim_k\Gamma(X,\sh{C}^\bullet(X,\sh{F}_k)).
\]
Acabamos de observar que todas as aplicações
\[
  \Gamma(X,\sh{C}^i(X,\sh{F}_k))
  \longrightarrow
  \Gamma(X,\sh{C}^i(X,\sh{F}_h))
  \qquad(h\leq k)
\]
são sobrejetivas; a conclusão segue-se, portanto, uma vez mais de
\sref{0.13.2.3}.
\end{proof}

\begin{env}[13.3.2]
\label{0.13.3.2}
\emph{Observações} (13.3.2).
\begin{enumerate}
  \item[{\rm(i)}]
  O enunciado de \sref{0.13.3.1} somente tem interesse para $i>0$,
  pois, para $i=0$, $h_i$ é sempre um isomorfismo sem hipótese
  alguma (0\textsubscript{I}, \EGAUnavailableHyperref{0.3.2.6}{3.2.6}).

  \item[{\rm(ii)}]
  As condições~{\rm(i)} e~{\rm(ii)} de \sref{0.13.3.1} são satisfeitas,
  em particular, se
  $\HH^i(U,\sh{F}_k)=0$ para todo $k$, todo $i>0$ e todo
  $U\in\mathfrak{B}$, e se, para $U\in\mathfrak{B}$, as aplicações
  \[
    \Gamma(U,\sh{F}_k)\longrightarrow\Gamma(U,\sh{F}_h)
  \]
  são sobrejetivas. Este é o caso mais frequente no qual se aplica
  \sref{0.13.3.1}.
\end{enumerate}
\end{env}

\subsection{Condição de Mittag-Leffler e objetos graduados associados a sistemas projetivos}
\label{subsection:0.13.4}
\label{0.13.4}

\begin{env}[13.4.1]
\label{0.13.4.1}
Seja
$\bb{A}=(A_k,u_{kh})_{k\in\bb{Z}}$ um sistema projetivo em uma
categoria abeliana $\cat{C}$.
Dizemos que está \emph{limitado inferiormente} se existe um $k_0$ tal que
$A_k=0$ para $k<k_0$.
Sobre cada $A_k$ definimos uma \emph{filtração}
$(\mathrm{F}^p(A_k))_{p\in\bb{Z}}$ por
\begin{equation}
\label{0.13.4.1.1}
\tag{13.4.1.1}
  \mathrm{F}^p(A_k)=
  \begin{cases}
    \Ker(A_k\longrightarrow A_{p-1}) & \text{para }p\leq k+1,\\
    0 & \text{para }p\geq k+1.
  \end{cases}
\end{equation}

\oldpage[0\textsubscript{III}]{70}
Por hipótese,
$\mathrm{F}^{k_0}(A_k)=A_k$ e
$\mathrm{F}^{k+1}(A_k)=0$; em outras palavras, a filtração considerada
é finita (\sref{0.11.1.3}).
Os objetos graduados associados a esta filtração são, portanto,
\[
  \mathrm{gr}^p(A_k)
  =
  \Ker(A_k\longrightarrow A_{p-1})
  \big/
  \Ker(A_k\longrightarrow A_p).
\]
Por conseguinte, $\mathrm{gr}^p(A_k)$ é isomorfo à imagem por
$A_k\longrightarrow A_p$ de
$\Ker(A_k\longrightarrow A_{p-1})$.
Pela transitividade dos morfismos que definem um sistema projetivo,
tem-se, portanto,
\begin{equation}
\label{0.13.4.1.2}
\tag{13.4.1.2}
  \mathrm{gr}^p(A_k)
  =
  \Ker(A_p\longrightarrow A_{p-1})
  \cap
  \Im(A_k\longrightarrow A_p).
\end{equation}
Mas \sref{0.13.4.1.1} dá
$\Ker(A_p\longrightarrow A_{p-1})=\mathrm{gr}^p(A_p)$, e portanto
\begin{equation}
\label{0.13.4.1.3}
\tag{13.4.1.3}
  \mathrm{gr}^p(A_k)
  =
  \mathrm{gr}^p(A_p)
  \cap
  \Im(A_k\longrightarrow A_p).
\end{equation}

As definições precedentes mostram também que, para $k\leq h$,
\[
  u_{kh}\bigl(\mathrm{F}^p(A_h)\bigr)
  \subset
  \mathrm{F}^p(A_k),
\]
e, por conseguinte, para todo $p\in\bb{Z}$, os
$\mathrm{gr}^p(u_{kh})$ definem um sistema projetivo
$(\mathrm{gr}^p(A_k))_{k\in\bb{Z}}$.
\end{env}

\begin{env}[13.4.2]
\label{0.13.4.2}
Dizemos que o sistema projetivo $\bb{A}$ é
\emph{essencialmente constante} se os morfismos
$A_{k+1}\longrightarrow A_k$ são isomorfismos para todo $k$
suficientemente grande.
Dizemos que $\bb{A}$ é \emph{estrito} se os morfismos
$A_i\longrightarrow A_j$ $(j\leq i)$ são epimorfismos.

Quando $\bb{A}$ é estrito, segue-se de \sref{0.13.4.1.3} que, para
$p\leq k\leq h$, o morfismo canônico
\[
  \mathrm{gr}^p(A_h)\longrightarrow\mathrm{gr}^p(A_k)
\]
é um isomorfismo; em outras palavras, o sistema projetivo
$(\mathrm{gr}^p(A_k))_{k\in\bb{Z}}$ é essencialmente constante.
A sucessão de objetos $\mathrm{gr}^p(A_p)$, identificada para todo
$p$ com $\varprojlim_k\mathrm{gr}^p(A_k)$, denota-se então
$\mathrm{gr}^\bullet(\bb{A})$ e chama-se o
\emph{objeto graduado associado ao sistema projetivo estrito}
$\bb{A}=(A_k)$.

Suponhamos agora que o sistema projetivo $\bb{A}$, limitado
inferiormente, satisfaz (ML).
Por \sref{0.13.1.2}, o sistema projetivo
$\bb{A}'=(A'_k)$ de objetos de ``imagem universal'' é estrito; também está
limitado inferiormente.
O objeto graduado $\mathrm{gr}^\bullet(\bb{A}')$ associado a
$\bb{A}'$ chama-se igualmente o objeto graduado
\emph{associado a} $\bb{A}$ e é denotado
$\mathrm{gr}^\bullet(\bb{A})$.
\end{env}

\begin{proposition}[13.4.3]
\label{0.13.4.3}
Seja $\bb{A}=(A_k)_{k\in\bb{Z}}$ um sistema projetivo limitado
inferiormente em uma categoria abeliana $\cat{C}$.
As condições seguintes são equivalentes:
\begin{enumerate}
  \item[{\rm a)}] $\bb{A}$ satisfaz (ML).
  \item[{\rm b)}] Para todo $p\in\bb{Z}$, o sistema projetivo
  $(\mathrm{gr}^p(A_k))_{k\in\bb{Z}}$ é essencialmente constante.
\end{enumerate}
Além disso, quando essas condições são satisfeitas, existe, para todo
$p\in\bb{Z}$, um isomorfismo canônico
\begin{equation}
\label{0.13.4.3.1}
\tag{13.4.3.1}
  \mathrm{gr}^p(\bb{A})
  \isoto
  \varprojlim_k\mathrm{gr}^p(A_k).
\end{equation}
\end{proposition}

\begin{proof}
Segue-se imediatamente de \sref{0.13.4.1.2} que \emph{a)} implica
\emph{b)}. A mesma fórmula, aplicada ao sistema projetivo
$\bb{A}'$ com a notação de \sref{0.13.4.2}, da
\sref{0.13.4.3.1} por definição.

Para $k\leq h$, ponhamos
\[
  A_{kh}=\Im(A_h\longrightarrow A_k).
\]
Se $k\leq h\leq j$, então
$A_{kj}\subset A_{kh}\subset A_k$.
Dotemos $A_{kh}$ da filtração induzida por
$(\mathrm{F}^p(A_k))$.
A transitividade dos morfismos que definem $\bb{A}$ mostra
imediatamente que esta é também a filtração quociente de
$(\mathrm{F}^p(A_h))$; por conseguinte,
\begin{equation}
\label{0.13.4.3.2}
\tag{13.4.3.2}
  \mathrm{gr}^p(A_{kh})
  =
  \Im\bigl(
    \mathrm{gr}^p(A_h)\longrightarrow\mathrm{gr}^p(A_k)
  \bigr).
\end{equation}

\oldpage[0\textsubscript{III}]{71}
Suponhamos agora que \emph{b)} seja satisfeita.
Para todo $p\in\bb{Z}$ e todo $k\geq p$, existe um inteiro
$\mathrm{L}(p,k)$ tal que o membro direito de
\sref{0.13.4.3.2} é constante para $h\geq\mathrm{L}(p,k)$.
Como $\mathrm{gr}^p(A_k)=0$ para $p<k_0$, para $k$ fixo somente
um número finito de inteiros $\mathrm{L}(p,k)$ pode ser não nulo quando
$p$ percorre os inteiros $\leq k$.
Seja $\mathrm{L}(k)=m$ o maior destes inteiros.
Para todo $h\geq m$, tem-se
$A_{kh}\subset A_{km}$, e, pela definição de $m$, a injeção
canônica
$A_{kh}\longrightarrow A_{km}$ induz um isomorfismo
\[
  \mathrm{gr}^\bullet(A_{kh})
  \simto
  \mathrm{gr}^\bullet(A_{km}).
\]
Como as filtrações são finitas, esta injeção é ela própria bijetiva
(Bourbaki, \emph{Alg.\ comm.}, cap.~III, \textsection2,
n\textsuperscript{o}~8, teorema~1).
Isto demonstra que $\bb{A}$ satisfaz (ML).
\end{proof}

\begin{env}[13.4.4]
\label{0.13.4.4}
Suponhamos que o limite projetivo
$A=\varprojlim_k A_k$ existe em $\cat{C}$.
Nas definições de \sref{0.13.4.1}, se pode então substituir $A_k$
por $A$, e a filtração assim definida sobre $A$ continua satisfazendo
\begin{equation}
\label{0.13.4.4.1}
\tag{13.4.4.1}
  \mathrm{gr}^p(A)
  =
  \mathrm{gr}^p(A_p)
  \cap
  \Im(A\longrightarrow A_p).
\end{equation}
\end{env}

\begin{corollary}[13.4.5]
\label{0.13.4.5}
Suponhamos que $\cat{C}$ é a categoria de grupos abelianos.
Se o sistema projetivo $\bb{A}$ satisfaz (ML) e
$A=\varprojlim_k A_k$, então para todo $p\in\bb{Z}$ existe um
isomorfismo canônico
\begin{equation}
\label{0.13.4.5.1}
\tag{13.4.5.1}
  \mathrm{gr}^p(A)
  \isoto
  \varprojlim_k\mathrm{gr}^p(A_k).
\end{equation}
\end{corollary}

\begin{proof}
Com efeito,
\[
  \Im(A_k\longrightarrow A_p)
  =
  \Im(A\longrightarrow A_p)
\]
para todo $k$ suficientemente grande
(Bourbaki, \emph{Top.\ g\'{e}n.}, cap.~II,
3\textsuperscript{e}~\'{e}d., \textsection3,
n\textsuperscript{o}~5, teorema~1), e a conclusão decorre de
\sref{0.13.4.1.3} e \sref{0.13.4.4.1}.
\end{proof}

\subsection{Limites projetivos de sucessões espectrais de complexos filtrados}
\label{subsection:0.13.5}

\begin{env}[13.5.1]
\label{0.13.5.1}
Seja $\cat{C}$ uma categoria abeliana e seja $X^\bullet$ um complexo
de objetos de $\cat{C}$ dotado de uma filtração
$(\mathrm{F}^p(X^\bullet))_{p\in\bb{Z}}$ tal que
$\mathrm{F}^{p_0}(X^\bullet)=X^\bullet$ para algum índice $p_0$.
Para todo $k\in\bb{Z}$, consideremos o complexo
\[
  X^\bullet_k=X^\bullet/\mathrm{F}^{k+1}(X^\bullet);
\]
está dotado canonicamente da filtração constituída por
\[
  \mathrm{F}^p(X^\bullet_k)
  =
  \mathrm{F}^p(X^\bullet)/\mathrm{F}^{k+1}(X^\bullet)
  \quad\text{para }p\leq k,
\]
e $\mathrm{F}^p(X^\bullet_k)=0$ para $p\geq k+1$.
Além disso, existem morfismos canônicos
$X^\bullet_{k+1}\longrightarrow X^\bullet_k$, que fazem de
$\bb{X}^\bullet=(X^\bullet_k)_{k\in\bb{Z}}$ um \emph{sistema projetivo
de complexos filtrados de objetos de $\cat{C}$}.
Observemos que este sistema projetivo é \emph{estrito} e que
$X^\bullet_k=0$ para $k<p_0$.
\end{env}

\begin{env}[13.5.2]
\label{0.13.5.2}
Mais geralmente, consideremos um sistema projetivo estrito limitado
inferiormente
$\bb{X}^\bullet=(X^\bullet_k)_{k\in\bb{Z}}$ de complexos de objetos de
$\cat{C}$, e dotemos cada $X^\bullet_k$ da filtração definida em
\sref{0.13.4.1} (aplicada na categoria abeliana de complexos limitados
inferiormente de objetos de $\cat{C}$).
Os morfismos
$X^\bullet_k\longrightarrow X^\bullet_p$ ($p\leq k$) tornam-se então
em morfismos de complexos \emph{filtrados} com filtrações finitas.
A funtorialidade das sucessões espectrais de complexos filtrados
(\sref{0.11.2.3}) mostra que os morfismos de transição de
$\bb{X}^\bullet$ induzem morfismos que fazem de
\[
  \mathrm{E}(\bb{X}^\bullet)
  =
  \bigl(\mathrm{E}(X^\bullet_k)\bigr)_{k\in\bb{Z}}
\]
um \emph{sistema projetivo de sucessões espectrais}.
\end{env}

\begin{lemma}[13.5.3]
\label{0.13.5.3}
Suponhamos que o sistema projetivo
$\bb{X}^\bullet=(X^\bullet_k)_{k\in\bb{Z}}$ de complexos filtrados se
obtém como em \sref{0.13.5.2}. Então:
\begin{enumerate}
  \item[{\rm a)}] Para $r\geq p-p_0$, tem-se
  \[
    \mathrm{B}_r^{pq}(X^\bullet_k)
    =
    \mathrm{B}_\infty^{pq}(X^\bullet_k)
  \]
  para todo $k\in\bb{Z}$.
  \item[{\rm b)}] Para $k+1\leq p+r$, tem-se
  \[
    \mathrm{Z}_r^{pq}(X^\bullet_k)
    =
    \mathrm{Z}_\infty^{pq}(X^\bullet_k).
  \]
  \item[{\rm c)}] Para $k+1\geq p+r$, os morfismos
  \[
    \mathrm{Z}_r^{pq}(X^\bullet_h)
    \longrightarrow
    \mathrm{Z}_r^{pq}(X^\bullet_k)
    \quad\text{e}\quad
    \mathrm{B}_r^{pq}(X^\bullet_h)
    \longrightarrow
    \mathrm{B}_r^{pq}(X^\bullet_k)
  \]
  são isomorfismos para todo $h\geq k$.
\end{enumerate}
\end{lemma}

\oldpage[0\textsubscript{III}]{72}
Estas três propriedades seguem-se imediatamente das definições de
\sref{0.11.2.2}, levando em conta que
$\mathrm{F}^{p-r+1}(X^\bullet_k)=X^\bullet_k$ quando $p-r<p_0$.

\begin{env}[13.5.4]
\label{0.13.5.4}
Suponhamos que sejam satisfeitas as hipóteses de \sref{0.13.5.3}.
Então, para $p$, $q$ e $r$ fixos (com $r$ finito), os sistemas
projetivos
\[
  \bigl(\mathrm{Z}_r^{pq}(X^\bullet_k)\bigr)_{k\in\bb{Z}},
  \qquad
  \bigl(\mathrm{B}_r^{pq}(X^\bullet_k)\bigr)_{k\in\bb{Z}},
  \qquad
  \bigl(\mathrm{E}_r^{pq}(X^\bullet_k)\bigr)_{k\in\bb{Z}}
\]
são \emph{essencialmente constantes}; seus respectivos limites projetivos
serão denotados por
$\mathrm{Z}_r^{pq}(\bb{X}^\bullet)$,
$\mathrm{B}_r^{pq}(\bb{X}^\bullet)$, e
\[
  \mathrm{E}_r^{pq}(\bb{X}^\bullet)
  =
  \mathrm{Z}_r^{pq}(\bb{X}^\bullet)/
  \mathrm{B}_r^{pq}(\bb{X}^\bullet).
\]
Os objetos $\mathrm{Z}_r^{pq}(\bb{X}^\bullet)$ e
$\mathrm{B}_r^{pq}(\bb{X}^\bullet)$ identificam-se canonicamente com
subobjetos de $\mathrm{E}_1^{pq}(\bb{X}^\bullet)$.
A definição dos $d_r^{pq}$ (M, XV, 1) mostra que estes
morfismos, tomados com respeito aos $X^\bullet_k$, são também essencialmente
constantes, e definem, portanto, morfismos
\begin{equation}
\label{0.13.5.4.1}
\tag{13.5.4.1}
  d_r^{pq}:
  \mathrm{E}_r^{pq}(\bb{X}^\bullet)
  \longrightarrow
  \mathrm{E}_r^{p+r,q-r+1}(\bb{X}^\bullet)
\end{equation}
tal que
$d_r^{p+r,q-r+1}\circ d_r^{pq}=0$.
Além disso, existem isomorfismos canônicos de
$\Ker(d_r^{pq})$ sobre
\[
  \mathrm{Z}_{r+1}^{pq}(\bb{X}^\bullet)/
  \mathrm{B}_r^{pq}(\bb{X}^\bullet)
\]
e de $\Im(d_r^{pq})$ sobre
\[
  \mathrm{B}_{r+1}^{p+r,q-r+1}(\bb{X}^\bullet)/
  \mathrm{B}_r^{p+r,q-r+1}(\bb{X}^\bullet).
\]
\end{env}

\begin{lemma}[13.5.5]
\label{0.13.5.5}
Sob as hipóteses de \sref{0.13.5.3}, para
$s\geq r>p-p_0$ existe um monomorfismo canônico
\begin{equation}
\label{0.13.5.5.1}
\tag{13.5.5.1}
  i:
  \mathrm{E}_s^{pq}(\bb{X}^\bullet)
  \longrightarrow
  \mathrm{E}_r^{pq}(\bb{X}^\bullet)
\end{equation}
e um isomorfismo canônico
\begin{equation}
\label{0.13.5.5.2}
\tag{13.5.5.2}
  j_r:
  \mathrm{E}_r^{pq}(\bb{X}^\bullet)
  \simto
  \mathrm{E}_\infty^{pq}(X^\bullet_{p+r-1})
\end{equation}
tal que o diagrama
\begin{equation}
\label{0.13.5.5.3}
\tag{13.5.5.3}
  \xymatrix{
    \mathrm{E}_s^{pq}(\bb{X}^\bullet)
      \ar[r]^{j_s}\ar[d]_i
    &
    \mathrm{E}_\infty^{pq}(X^\bullet_{p+s-1})
      \ar[d]
    \\
    \mathrm{E}_r^{pq}(\bb{X}^\bullet)
      \ar[r]^{j_r}
    &
    \mathrm{E}_\infty^{pq}(X^\bullet_{p+r-1})
  }
\end{equation}
é comutativo, onde a seta vertical direita está induzida pelo
morfismo
$X^\bullet_{p+s-1}\longrightarrow X^\bullet_{p+r-1}$.
\end{lemma}

\begin{proof}
A existência de $i$ segue-se de
\[
  \mathrm{B}_r^{pq}(X^\bullet_k)
  =
  \mathrm{B}_\infty^{pq}(X^\bullet_k)
  \qquad (r>p-p_0)
\]
(\sref{0.13.5.3},~a)).
Para $k+1\leq p+r$, tem-se
\[
  \mathrm{Z}_r^{pq}(X^\bullet_k)
  =
  \mathrm{Z}_\infty^{pq}(X^\bullet_k)
\]
(\sref{0.13.5.3},~b)); em particular,
\[
  \mathrm{Z}_r^{pq}(X^\bullet_{p+r-1})
  =
  \mathrm{Z}_\infty^{pq}(X^\bullet_{p+r-1}).
\]
Por outro lado,
$\mathrm{Z}_r^{pq}(X^\bullet_{p+r-1})$ e
$\mathrm{B}_r^{pq}(X^\bullet_{p+r-1})$ identificam-se canonicamente com
$\mathrm{Z}_r^{pq}(\bb{X}^\bullet)$ e
$\mathrm{B}_r^{pq}(\bb{X}^\bullet)$ por
\sref{0.13.5.3},~c).
Isto demonstra a existência de $j_r$ e a comutatividade de
\sref{0.13.5.5.3}.
\end{proof}

\begin{corollary}[13.5.6]
\label{0.13.5.6}
Sob as hipóteses de \sref{0.13.5.3}, se um dos limites
projetivos
\[
  \varprojlim_r\mathrm{E}_r^{pq}(\bb{X}^\bullet),
  \qquad
  \varprojlim_k\mathrm{E}_\infty^{pq}(X^\bullet_k)
\]
existe, então também existe o outro, e há um isomorfismo canônico
\begin{equation}
\label{0.13.5.6.1}
\tag{13.5.6.1}
  j_\infty:
  \varprojlim_r\mathrm{E}_r^{pq}(\bb{X}^\bullet)
  \simto
  \varprojlim_k\mathrm{E}_\infty^{pq}(X^\bullet_k).
\end{equation}
Além disso, o sistema projetivo
$(\mathrm{E}_r^{pq}(\bb{X}^\bullet))_{r\in\bb{Z}}$ é essencialmente
constante (\sref{0.13.4.2}) se e somente se o sistema projetivo
$(\mathrm{E}_\infty^{pq}(X^\bullet_k))_{k\in\bb{Z}}$ é essencialmente
constante.
\end{corollary}

\begin{env}[13.5.7]
\label{0.13.5.7}
Denotamos por
$\mathrm{B}_\infty^{pq}(\bb{X}^\bullet)$ e
$\mathrm{Z}_\infty^{pq}(\bb{X}^\bullet)$ os subobjetos de
$\mathrm{E}_1^{pq}(\bb{X}^\bullet)$ iguais respectivamente a
$\mathrm{B}_r^{pq}(\bb{X}^\bullet)$ para $r>p-p_0$ e a
$\inf_r\mathrm{Z}_r^{pq}(\bb{X}^\bullet)$, quando este último existe.
Assim,
\[
  \varprojlim_r\mathrm{E}_r^{pq}(\bb{X}^\bullet)
\]
\oldpage[0\textsubscript{III}]{73}
identifica-se canonicamente com
\[
  \mathrm{E}_\infty^{pq}(\bb{X}^\bullet)
  =
  \mathrm{Z}_\infty^{pq}(\bb{X}^\bullet)/
  \mathrm{B}_\infty^{pq}(\bb{X}^\bullet).
\]
Observemos que os objetos
$\mathrm{Z}_r^{pq}(\bb{X}^\bullet)$,
$\mathrm{B}_r^{pq}(\bb{X}^\bullet)$,
$\mathrm{E}_r^{pq}(\bb{X}^\bullet)$
($1\leq r\leq+\infty$), e os morfismos $d_r^{pq}$ dependem
\emph{funtorialmente} do sistema projetivo $\bb{X}^\bullet$, com as
restrições de \sref{0.13.5.5}, e os morfismos definidos em
\sref{0.13.5.5} e \sref{0.13.5.6} são funtoriais.
\end{env}

\subsection{Sucessão espectral de um funtor relativa a um objeto dotado de uma filtração finita}
\label{subsection:0.13.6}

\begin{env}[13.6.1]
\label{0.13.6.1}
Sejam $\cat{C}$ e $\cat{C}'$ duas categorias abelianas, e seja
$\mathrm{T}:\cat{C}\longrightarrow\cat{C}'$ um funtor covariante
aditivo.
Suponhamos que todo objeto de $\cat{C}$ é isomorfo a um subobjeto de
um objeto injetivo, de modo que existam os funtores derivados à direita
$\RR^p\mathrm{T}$ ($p\geq0$).
\end{env}

\begin{lemma}[13.6.2]
\label{0.13.6.2}
Seja $A$ um objeto de $\cat{C}$ dotado de uma filtração finita
$(\mathrm{F}^i(A))_{i\in\bb{Z}}$.
Existe uma resolução injetiva
$X^\bullet=(X^j)_{j\geq0}$ de $A$, dotada de uma filtração finita
$(\mathrm{F}^i(X^\bullet))_{i\in\bb{Z}}$, tal que
\[
  \mathrm{F}^i(A)=A
  \quad\text{(respectivamente, }\mathrm{F}^i(A)=0\text{)}
\]
implica
\[
  \mathrm{F}^i(X^\bullet)=X^\bullet
  \quad\text{(respectivamente, }\mathrm{F}^i(X^\bullet)=0\text{)},
\]
e tal que, para todo $i\in\bb{Z}$,
$\mathrm{F}^i(X^\bullet)$ é uma resolução injetiva de
$\mathrm{F}^i(A)$.
\end{lemma}

\begin{proof}
Seja $p$ (respectivamente, $q>p$) o maior índice tal que
$\mathrm{F}^p(A)=A$ (respectivamente, o menor índice tal que
$\mathrm{F}^q(A)=0$).
Raciocinamos por indução sobre $q-p$, sendo imediato o lema quando
$q-p=1$.
Suponhamos construída uma resolução injetiva ${X''}^\bullet$ de
$A/\mathrm{F}^{q-1}(A)$ que tenha as propriedades requeridas.
Consideremos a sucessão exata
\[
  0\longrightarrow\mathrm{F}^{q-1}(A)
  \longrightarrow A
  \longrightarrow A/\mathrm{F}^{q-1}(A)
  \longrightarrow0,
\]
tomemos uma resolução injetiva ${X'}^\bullet$ de
$\mathrm{F}^{q-1}(A)$, e escolhamos depois uma resolução injetiva
$X^\bullet$ de $A$ de modo que se obtenha uma sucessão exata
\[
  0\longrightarrow {X'}^\bullet
  \longrightarrow X^\bullet
  \longrightarrow {X''}^\bullet
  \longrightarrow0
\]
compatível com a sucessão precedente (M, V, 2.2).
Está claro que $X^\bullet$ tem as propriedades requeridas.
\end{proof}

\begin{corollary}[13.6.3]
\label{0.13.6.3}
Seja $B$ um segundo objeto de $\cat{C}$, dotado de uma filtração
finita $(\mathrm{F}^i(B))_{i\in\bb{Z}}$, seja $s$ um inteiro,
e seja $u:A\longrightarrow B$ um morfismo tal que
\[
  u(\mathrm{F}^i(A))\subset\mathrm{F}^{i+s}(B)
\]
para todo $i\in\bb{Z}$.
Se $Y^\bullet=(Y^j)_{j\geq0}$ é uma resolução injetiva de $B$
dotada de uma filtração
$(\mathrm{F}^i(Y^\bullet))_{i\in\bb{Z}}$ que tenha as propriedades
enunciadas em \sref{0.13.6.2}, então existe um morfismo
$v:X^\bullet\longrightarrow Y^\bullet$, compatível com $u$, tal
que
\[
  v(\mathrm{F}^i(X^\bullet))
  \subset
  \mathrm{F}^{i+s}(Y^\bullet)
\]
para todo $i\in\bb{Z}$.
Além disso, quaisquer dois desses morfismos $v$ e $v'$ são homotópicos.
\end{corollary}

Isto segue-se imediatamente, por indução sobre $q-p$, da construção
precedente e de (M, V, 2.3).

\begin{env}[13.6.4]
\label{0.13.6.4}
Sob as hipóteses de \sref{0.13.6.2}, consideremos o complexo
$\mathrm{T}(X^\bullet)$ em $\cat{C}'$.
Está filtrado pelos complexos
$\mathrm{T}(\mathrm{F}^i(X^\bullet))$, pois
$\mathrm{F}^i(X^\bullet)$ é um sumando direto de $X^\bullet$.
Segue-se de \sref{0.13.6.3} que a \emph{sucessão espectral} de
este complexo filtrado depende, salvo isomorfismo, unicamente do
objeto filtrado $A$.
Seu termo limite é a cohomologia $\RR^\bullet\mathrm{T}(A)$, dotada
da filtração
\begin{equation}
\label{0.13.6.4.1}
\tag{13.6.4.1}
  \begin{split}
  \mathrm{F}^p(\RR^n\mathrm{T}(A))
  &=
  \Im\bigl(
    \RR^n\mathrm{T}(\mathrm{F}^p(A))
    \longrightarrow
    \RR^n\mathrm{T}(A)
  \bigr)
  \\
  &=
  \Ker\bigl(
    \RR^n\mathrm{T}(A)
    \longrightarrow
    \RR^n\mathrm{T}(A/\mathrm{F}^p(A))
  \bigr),
  \end{split}
\end{equation}
por \sref{0.11.2.2}, e seu termo $\mathrm{E}_1$ é
\begin{equation}
\label{0.13.6.4.2}
\tag{13.6.4.2}
  \mathrm{E}_1^{pq}
  =
  \RR^{p+q}\mathrm{T}(\mathrm{gr}^p(A)),
\end{equation}
onde, como de costume,
$\mathrm{gr}^p(A)=\mathrm{F}^p(A)/\mathrm{F}^{p+1}(A)$.
Deduz-se de \sref{0.11.2.2} que a filtração do termo limite
é finita e que, para $p$ e $q$ fixos, as sucessões
\oldpage[0\textsubscript{III}]{74}
\[
  \mathrm{B}_r^{pq}(A)
  =
  \mathrm{B}_r^{pq}(\mathrm{T}(X^\bullet))
  \quad\text{e}\quad
  \mathrm{Z}_r^{pq}(A)
  =
  \mathrm{Z}_r^{pq}(\mathrm{T}(X^\bullet))
\]
são estacionárias.
A sucessão espectral precedente é, portanto, \emph{birregular}
(\sref{0.11.1.3}).
Denotamo-la por
$\mathrm{E}(A)=(\mathrm{E}_r^{pq}(A))$ e chamamo-la de
\emph{sucessão espectral do funtor $\mathrm{T}$ relativa ao
objeto filtrado $A$}.
\end{env}

\begin{env}[13.6.5]
\label{0.13.6.5}
Suponhamos agora que sejam satisfeitas as hipóteses de \sref{0.13.6.3}, e
conservemos sua notação.
Como $\mathrm{F}^i(X^\bullet)$ (respectivamente,
$\mathrm{F}^i(Y^\bullet)$) é um sumando direto de $X^\bullet$
(respectivamente, $Y^\bullet$), tem-se
\[
  (\mathrm{T}v)\bigl(\mathrm{T}(\mathrm{F}^i(X^\bullet))\bigr)
  \subset
  \mathrm{T}(\mathrm{F}^{i+s}(Y^\bullet))
\]
para todo $i\in\bb{Z}$.
As definições de \sref{0.11.2.2} mostram então que, para
$1\leq r\leq+\infty$, $\mathrm{T}v$ define morfismos
\[
  \mathrm{B}_r^{pq}(\mathrm{T}(X^\bullet))
  \longrightarrow
  \mathrm{B}_r^{p+s,q-s}(\mathrm{T}(Y^\bullet))
\]
e
\[
  \mathrm{Z}_r^{pq}(\mathrm{T}(X^\bullet))
  \longrightarrow
  \mathrm{Z}_r^{p+s,q-s}(\mathrm{T}(Y^\bullet)),
\]
e, portanto, um morfismo
\[
  w_r:
  \mathrm{E}_r^{pq}(A)
  \longrightarrow
  \mathrm{E}_r^{p+s,q-s}(B).
\]
Analogamente, sobre os termos limite têm-se morfismos
\[
  u_n:
  \RR^n\mathrm{T}(A)
  \longrightarrow
  \RR^n\mathrm{T}(B)
\]
tais que
\[
  u_n\bigl(\mathrm{F}^p(\RR^n\mathrm{T}(A))\bigr)
  \subset
  \mathrm{F}^{p+s}(\RR^n\mathrm{T}(B)).
\]
A definição dos $d_r^{pq}$ (M, XV, 1) mostra também que os
diagramas
\[
  \xymatrix{
    \mathrm{E}_r^{pq}(A)
      \ar[r]^{d_r^{pq}}\ar[d]_{w_r}
    &
    \mathrm{E}_r^{p+r,q-r+1}(A)
      \ar[d]^{w_r}
    \\
    \mathrm{E}_r^{p+s,q-s}(B)
      \ar[r]_{d_r^{p+s,q-s}}
    &
    \mathrm{E}_r^{p+r+s,q-r-s+1}(B)
  }
\]
são comutativos.
Segue-se que existe um diagrama comutativo análogo para os
isomorfismos $\alpha_r^{pq}$, cuja explicitação deixamos ao leitor.
Finalmente (\emph{loc.\ cit.}), têm-se também diagramas comutativos
sobre os termos limite:
\[
  \xymatrix{
    \mathrm{E}_\infty^{pq}(A)
      \ar[r]^{\beta^{pq}}\ar[d]_{w_\infty}
    &
    \mathrm{gr}^p(\RR^{p+q}\mathrm{T}(A))
      \ar[d]^{u_{p+q}}
    \\
    \mathrm{E}_\infty^{p+s,q-s}(B)
      \ar[r]_{\beta^{p+s,q-s}}
    &
    \mathrm{gr}^{p+s}(\RR^{p+q}\mathrm{T}(B)).
  }
\]
\end{env}

\begin{env}[13.6.6]
\label{0.13.6.6}
Suponhamos, em particular, que existe um anel $S$, dotado de uma
\emph{filtração} $(\mathrm{F}^i(S))_{i\in\bb{Z}}$, e um homomorfismo
de anéis
\begin{equation}
\label{0.13.6.6.1}
\tag{13.6.6.1}
  h:S\longrightarrow\Hom_{\cat{C}}(A,A)
\end{equation}
\oldpage[0\textsubscript{III}]{75}
tal que, para todo $t\in\mathrm{F}^i(S)$,
\[
  h_t(\mathrm{F}^j(A))
  \subset
  \mathrm{F}^{i+j}(A)
\]
para todo par $i,j$.
Para abreviar, dizemos então que $A$ está dotado da estrutura de um
\emph{$S$--$\cat{C}$-módulo filtrado} sobre o anel filtrado $S$.

Ao passar aos objetos graduados associados, todo $h_t$, para
$t\in\mathrm{F}^i(S)$, define um endomorfismo graduado
$\overline{h}_t$ de $\mathrm{gr}^\bullet(A)$, homogêneo de grau
$i$.
Além disso, este endomorfismo depende unicamente da classe de $t$ em
$\mathrm{gr}^i(S)$, e obtém-se assim um homomorfismo de
\emph{anéis graduados}
\[
  \overline{h}:
  \mathrm{gr}^\bullet(S)
  \longrightarrow
  \Hom_{\cat{C}}^\bullet
  \bigl(\mathrm{gr}^\bullet(A),\mathrm{gr}^\bullet(A)\bigr),
\]
onde o membro direito é o anel de endomorfismos graduados de
$\mathrm{gr}^\bullet(A)$.
Dizemos que $\mathrm{gr}^\bullet(A)$ está dotado da estrutura
de um \emph{$\mathrm{gr}^\bullet(S)$--$\cat{C}$-módulo graduado}.

Segue-se de \sref{0.13.6.5} que, para
$1\leq r\leq+\infty$, todo
$\overline{t}\in\mathrm{gr}^i(S)$ define canonicamente, sobre os
objetos bigraduados
\[
  \bigl(\mathrm{B}_r^{pq}(A)\bigr)_{(p,q)\in\bb{Z}\times\bb{Z}},
  \qquad
  \bigl(\mathrm{Z}_r^{pq}(A)\bigr)_{(p,q)\in\bb{Z}\times\bb{Z}},
\]
e
\[
  \mathrm{E}_r(A)
  =
  \bigl(\mathrm{E}_r^{pq}(A)\bigr)_{(p,q)\in\bb{Z}\times\bb{Z}},
\]
endomorfismos bigraduados de grau $(i,-i)$.
Sobre $\mathrm{E}_r(A)$, para $r$ finito, este endomorfismo comuta com
o endomorfismo bigraduado definido pelos $d_r^{pq}$.
Estes endomorfismos satisfazem as condições usuais de associatividade e
distributividade com respeito à adição em
$\mathrm{gr}^\bullet(S)$ e nos objetos bigraduados em questão.
Para abreviar, dizemos que estes últimos são
\emph{$\mathrm{gr}^\bullet(S)$--$\cat{C}'$-módulos bigraduados};
é imediato que os $\alpha_r^{pq}$ definem um isomorfismo para
este tipo de estrutura.

Para todo inteiro $n$, seja $\mathrm{B}_r^{(n)}(A)$
(respectivamente, $\mathrm{Z}_r^{(n)}(A)$ e
$\mathrm{E}_r^{(n)}(A)$) o subobjeto graduado de
$\mathrm{B}_r^{\bullet,\bullet}(A)$
(respectivamente, $\mathrm{Z}_r^{\bullet,\bullet}(A)$ e
$\mathrm{E}_r^{\bullet,\bullet}(A)$) constituído pelos termos para
os quais $p+q=n$ ($1\leq r\leq+\infty$).
É imediato que estes são
\emph{$\mathrm{gr}^\bullet(S)$--$\cat{C}'$-módulos graduados}.

Finalmente, todo $\overline{t}\in\mathrm{gr}^i(S)$ define, para todo
$n$, um endomorfismo graduado de grau $i$ sobre o objeto graduado
$\mathrm{gr}^\bullet(\RR^n\mathrm{T}(A))$, que adquire assim a
estrutura de um
$\mathrm{gr}^\bullet(S)$--$\cat{C}'$-módulo graduado.
Para $p+q=n$, os $\beta^{pq}$ definem um isomorfismo
\[
  \mathrm{E}^{(n)}(A)
  \isoto
  \mathrm{gr}^\bullet(\RR^n\mathrm{T}(A))
\]
para este tipo de estrutura.

Observemos que, quando $\cat{C}'$ é a categoria de grupos abelianos, as
estruturas de $S$--$\cat{C}'$-módulo (respectivamente, de
$\mathrm{gr}^\bullet(S)$--$\cat{C}'$-módulo graduado ou bigraduado) são
simplesmente as estruturas usuais de $S$-módulo (respectivamente, de
$\mathrm{gr}^\bullet(S)$-módulo graduado ou bigraduado).
\end{env}

\subsection{Funtores derivados de um limite projetivo de argumentos}
\label{subsection:0.13.7}
\label{0.13.7}

\begin{env}[13.7.1]
\label{0.13.7.1}
Sejam $\cat{C}$ e $\cat{C}'$ duas categorias abelianas, e suponhamos
que todo objeto de $\cat{C}$ é isomorfo a um subobjeto de um
objeto injetivo.
Seja
\[
  \mathrm{T}:\cat{C}\longrightarrow\cat{C}'
\]
um funtor covariante aditivo.
Consideremos um sistema projetivo estrito
$\bb{A}=(A_k)_{k\in\bb{Z}}$ em $\cat{C}$, limitado inferiormente; mais
precisamente, suponhamos que $A_k=0$ para $k<k_0$.
A este sistema associamos canonicamente, sobre cada $A_k$, a
filtração $(\mathrm{F}^p(A_k))_{p\in\bb{Z}}$ dada por
\sref{0.13.4.1.1}.
Como o sistema projetivo é estrito, os morfismos canônicos
\begin{equation}
\label{0.13.7.1.1}
\tag{13.7.1.1}
  \mathrm{F}^i(A_h)/\mathrm{F}^j(A_h)
  \longrightarrow
  \mathrm{F}^i(A_k)/\mathrm{F}^j(A_k)
  \qquad(h\geq k)
\end{equation}
são isomorfismos para $i\leq j\leq k+1$.
Lembremos também que
$\mathrm{F}^{k_0}(A_k)=A_k$ e
$\mathrm{F}^{k+1}(A_k)=0$ para todo $k$.
\end{env}

\begin{env}[13.7.2]
\label{0.13.7.2}
Para cada $k$, construamos uma resolução injetiva
$\mathrm{X}_k^\bullet=(\mathrm{X}_k^j)_{j\geq0}$ de $A_k$ que tenha as
propriedades de \sref{0.13.6.2}.
Os morfismos canônicos $A_{k+1}\longrightarrow A_k$ nos permitem, por
\sref{0.13.6.3}, definir para cada $k$ um morfismo de complexos
\[
  \mathrm{X}_{k+1}^\bullet\longrightarrow\mathrm{X}_k^\bullet
\]
compatível com as filtrações e que faz de
$\bb{X}^\bullet=(\mathrm{X}_k^\bullet)_{k\in\bb{Z}}$ um sistema
projetivo de complexos.
\oldpage[0\textsubscript{III}]{76}
Podemos supor além disso que este sistema projetivo é estrito.
Com efeito, por \sref{0.13.7.1.1},
\[
  A_k
  \simeq
  A_{k+1}/\mathrm{F}^{k+1}(A_{k+1}).
\]
Assim, ao construir $\mathrm{X}_{k+1}^\bullet$, podemos tomar a
resolução injetiva de
$A_{k+1}/\mathrm{F}^{k+1}(A_{k+1})$ \emph{igual} a
$\mathrm{X}_k^\bullet$.
Segue-se então de (M, V, 2.3) que o morfismo de complexos
$\mathrm{X}_{k+1}^\bullet\longrightarrow\mathrm{X}_k^\bullet$ pode
ser construído de modo que, ao passar aos quocientes, induza um
isomorfismo
\[
  \mathrm{X}_{k+1}^\bullet/
  \mathrm{F}^{k+1}(\mathrm{X}_{k+1}^\bullet)
  \simto
  \mathrm{X}_k^\bullet
\]
que respeita as filtrações; esta é precisamente a condição de
\sref{0.13.5.1}.
\end{env}

\begin{env}[13.7.3]
\label{0.13.7.3}
Por construção, o sistema projetivo
$\mathrm{T}(\bb{X}^\bullet)$ de complexos
$\mathrm{T}(\mathrm{X}_k^\bullet)$ satisfaz as hipóteses de
\sref{0.13.5.3}.
Os resultados de \sref{0.13.5.4}, \sref{0.13.5.5} e
\sref{0.13.5.6} aplicam-se, portanto, às sucessões espectrais
\[
  \mathrm{E}\bigl(\mathrm{T}(\mathrm{X}_k^\bullet)\bigr)
  =
  \mathrm{E}(A_k).
\]
Para $1\leq r\leq+\infty$, escrevemos
$\mathrm{E}_r^{pq}(\bb{A})$ em lugar de
$\mathrm{E}_r^{pq}(\mathrm{T}(\bb{X}^\bullet))$
(cf.\ \sref{0.13.5.7} quando $r=+\infty$), e utilizamos a mesma convenção
para a notação análoga.
Em particular,
\begin{equation}
\label{0.13.7.3.1}
\tag{13.7.3.1}
  \mathrm{E}_1^{pq}(\bb{A})
  =
  \RR^{p+q}\mathrm{T}\bigl(\mathrm{gr}^p(\bb{A})\bigr),
\end{equation}
por \sref{0.13.6.4.2} e o fato de que o sistema projetivo
$(\mathrm{gr}^p(A_k))_{k\in\bb{Z}}$ é essencialmente constante.
Estes resultados, junto com \sref{0.13.4.3}, dão a proposição
seguinte, demonstrada pela primeira vez por Shih Weishu mediante um método
diferente (inédito).
\end{env}

\begin{proposition}[13.7.4]
\label{0.13.7.4}
\emph{(Shih).}
Seja $n$ um inteiro.
As duas condições seguintes são equivalentes:
\begin{enumerate}
  \item[{\rm a)}] Para todo par $(p,q)$ tal que $p+q=n$, o
  sistema projetivo
  $(\mathrm{E}_r^{pq}(\bb{A}))_{r\geq2}$ é essencialmente constante.
  \item[{\rm b)}] O sistema projetivo
  \[
    \RR^n\mathrm{T}(\bb{A})
    =
    \bigl(\RR^n\mathrm{T}(A_k)\bigr)_{k\in\bb{Z}}
  \]
  satisfaz {\rm(ML)}.
\end{enumerate}
Além disso, quando essas condições são satisfeitas, existe um
isomorfismo canônico
\begin{equation}
\label{0.13.7.4.1}
\tag{13.7.4.1}
  \mathrm{gr}^p\bigl(\RR^n\mathrm{T}(\bb{A})\bigr)
  \isoto
  \mathrm{E}_\infty^{p,n-p}(\bb{A})
  \qquad(p\in\bb{Z}).
\end{equation}
\end{proposition}

\begin{proof}
Por \sref{0.13.5.6}, a condição {\rm a)} diz precisamente que, para
$p+q=n$, o sistema projetivo
$(\mathrm{E}_\infty^{pq}(A_k))_{k\in\bb{Z}}$ é essencialmente
constante.
Por outro lado,
$\mathrm{gr}^p(\RR^n\mathrm{T}(A_k))$ é canonicamente isomorfo a
$\mathrm{E}_\infty^{p,n-p}(A_k)$.
Segue-se de \sref{0.13.4.3} que {\rm a)} e {\rm b)} são
equivalentes, e \sref{0.13.7.4.1} não é senão
\sref{0.13.5.6.1} no caso considerado.
\end{proof}

\begin{corollary}[13.7.5]
\label{0.13.7.5}
Seja $\sh{F}=(\sh{F}_k)_{k\in\bb{N}}$ um sistema projetivo de
feixes de grupos abelianos que satisfaz as condições {\rm(i)}, {\rm(ii)}
e {\rm(iii)} de \sref{0.13.3.1}, e ponhamos
$\sh{F}=\varprojlim_k\sh{F}_k$.
Suponhamos que, para o funtor
$\sh{G}\longmapsto\Gamma(X,\sh{G})$, o sistema projetivo
$(\mathrm{E}_r^{pq}(\sh{F}))_r$ é essencialmente constante sempre que
$p+q=n$ ou $p+q=n+1$.
Sobre $\HH^{n+1}(X,\sh{F})$, consideremos a filtração
\[
  \mathrm{F}^p\bigl(\HH^{n+1}(X,\sh{F})\bigr)
  =
  \Ker\bigl(
    \HH^{n+1}(X,\sh{F})
    \longrightarrow
    \HH^{n+1}(X,\sh{F}_{p-1})
  \bigr).
\]
Então existe um isomorfismo canônico
\begin{equation}
\label{0.13.7.5.1}
\tag{13.7.5.1}
  \mathrm{gr}^p\bigl(\HH^{n+1}(X,\sh{F})\bigr)
  \isoto
  \mathrm{E}_\infty^{p,n-p+1}(\sh{F})
  \qquad(p\in\bb{Z}).
\end{equation}
\end{corollary}

\begin{proof}
\oldpage[0\textsubscript{III}]{77}
Apliquemos \sref{0.13.7.4} tomando como $\cat{C}$ a categoria de feixes de
grupos abelianos sobre $X$, como $\cat{C}'$ a categoria de grupos abelianos,
e $\mathrm{T}=\Gamma$.
Dá, para todo $p\in\bb{Z}$, um isomorfismo canônico
\[
  \mathrm{gr}^p\bigl(\RR^{n+1}\Gamma(\sh{F})\bigr)
  \isoto
  \mathrm{E}_\infty^{p,n-p+1}(\sh{F}).
\]
Além disso, por \sref{0.13.7.4}, o sistema projetivo
$(\HH^n(X,\sh{F}_k))_{k\in\bb{Z}}$ satisfaz {\rm(ML)}; portanto
\sref{0.13.3.1} dá um isomorfismo canônico
\[
  \HH^{n+1}(X,\sh{F})
  \isoto
  \varprojlim_k\HH^{n+1}(X,\sh{F}_k).
\]

Como o sistema projetivo
$\RR^{n+1}\Gamma(\sh{F})$ satisfaz {\rm(ML)} por
\sref{0.13.7.4}, existe um isomorfismo canônico
\[
  \mathrm{gr}^p\bigl(\RR^{n+1}\Gamma(\sh{F})\bigr)
  \isoto
  \varprojlim_k
  \mathrm{gr}^p\bigl(\HH^{n+1}(X,\sh{F}_k)\bigr)
\]
por \sref{0.13.4.3}, e um isomorfismo canônico
\[
  \varprojlim_k
  \mathrm{gr}^p\bigl(\HH^{n+1}(X,\sh{F}_k)\bigr)
  \isoto
  \mathrm{gr}^p\left(
    \varprojlim_k\HH^{n+1}(X,\sh{F}_k)
  \right)
\]
por \sref{0.13.4.5}.
Resta apenas verificar que o isomorfismo precedente é
compatível com as filtrações de seus dois lados.
Isto segue-se imediatamente das definições e, para todo $p$,
da comutatividade de
\[
  \xymatrix{
    \HH^{n+1}(X,\sh{F})
      \ar[r]^-{\sim}\ar[dr]
    &
    \varprojlim_k\HH^{n+1}(X,\sh{F}_k)
      \ar[d]
    \\
    &
    \HH^{n+1}(X,\sh{F}_{p-1}).
  }
\]
\end{proof}

\begin{env}[13.7.6]
\label{0.13.7.6}
Seja $S$ um anel dotado de uma filtração
$(\mathrm{F}^i(S))_{i\in\bb{Z}}$ tal que
$\mathrm{F}^0(S)=S$ (e portanto
$\mathrm{gr}^i(S)=0$ para $i<0$).
Suponhamos que cada $A_k$, com a filtração definida em
\sref{0.13.7.1}, é um
$S$--$\cat{C}$-módulo filtrado no sentido de \sref{0.13.6.6}, e que os
morfismos $A_h\longrightarrow A_k$ para $k\leq h$ são morfismos de
$S$--$\cat{C}$-módulos filtrados.
Para abreviar, dizemos então que $\bb{A}$ é um
\emph{sistema projetivo de $S$--$\cat{C}$-módulos filtrados}.

É imediato que, para $k\leq h$ e
$1\leq r\leq+\infty$, os morfismos
\[
  \mathrm{B}_r^{pq}(A_h)\longrightarrow\mathrm{B}_r^{pq}(A_k),
  \qquad
  \mathrm{Z}_r^{pq}(A_h)\longrightarrow\mathrm{Z}_r^{pq}(A_k)
\]
são morfismos de
$\mathrm{gr}^\bullet(S)$--$\cat{C}'$-módulos bigraduados
(\sref{0.13.6.5}), e que, para $r$ finito, as famílias
\[
  \bigl(\mathrm{Z}_r^{pq}(\bb{A})\bigr),\qquad
  \bigl(\mathrm{B}_r^{pq}(\bb{A})\bigr),\qquad
  \bigl(\mathrm{E}_r^{pq}(\bb{A})\bigr)
\]
são $\mathrm{gr}^\bullet(S)$--$\cat{C}'$-módulos
bigraduados, sendo os dois primeiros submódulos de
$(\mathrm{E}_1^{pq}(\bb{A}))$.
Escrevemos de novo
$\mathrm{Z}_r^{(n)}(\bb{A})$,
$\mathrm{B}_r^{(n)}(\bb{A})$ e
$\mathrm{E}_r^{(n)}(\bb{A})$ para os respectivos subobjetos graduados
obtidos conservando unicamente os termos com $p+q=n$; estes são
$\mathrm{gr}^\bullet(S)$--$\cat{C}'$-módulos graduados.

Quando o sistema projetivo
$(\mathrm{E}_r^{pq}(\bb{A}))_r$ é essencialmente constante,
$\mathrm{E}_\infty^{pq}(\bb{A})$ é igualmente um
$\mathrm{gr}^\bullet(S)$--$\cat{C}'$-módulo bigraduado, e cada
$\mathrm{E}_\infty^{(n)}(\bb{A})$ é um
$\mathrm{gr}^\bullet(S)$--$\cat{C}'$-módulo graduado.
Além disso, para cada $k$, os isomorfismos
\[
  \beta^{p,n-p}:
  \mathrm{E}_\infty^{p,n-p}(A_k)
  \isoto
  \mathrm{gr}^p\bigl(\RR^n\mathrm{T}(A_k)\bigr)
\]
formam um isomorfismo de
$\mathrm{gr}^\bullet(S)$--$\cat{C}'$-módulos graduados de
$\mathrm{E}_\infty^{(n)}(A_k)$ sobre
$\mathrm{gr}^\bullet(\RR^n\mathrm{T}(A_k))$.
Sob as hipóteses precedentes, o isomorfismo
\[
  \beta^{p,n-p}:
  \mathrm{E}_\infty^{(n)}(\bb{A})
  \isoto
  \varprojlim_k
  \mathrm{gr}^\bullet\bigl(\RR^n\mathrm{T}(A_k)\bigr)
\]
respeita, portanto, também estas estruturas.
O mesmo vale manifestamente para o isomorfismo canônico
\[
  \mathrm{gr}^\bullet\bigl(\RR^n\mathrm{T}(\bb{A})\bigr)
  \isoto
  \varprojlim_k
  \mathrm{gr}^\bullet\bigl(\RR^n\mathrm{T}(A_k)\bigr);
\]
por conseguinte, os isomorfismos \sref{0.13.7.4.1} formam um
isomorfismo de
$\mathrm{gr}^\bullet(S)$--$\cat{C}'$-módulos graduados.
\end{env}

\begin{proposition}[13.7.7]
\label{0.13.7.7}
Seja $S$ um anel noetheriano $\mathfrak{J}$-ádico.
Suponhamos que $\cat{C}$ é uma categoria abeliana na qual todo objeto
é isomorfo a um subobjeto de um objeto injetivo, e seja
$\mathrm{T}$ um funtor covariante aditivo de $\cat{C}$ na
categoria de grupos abelianos.
Seja $\bb{A}=(A_k)_{k\in\bb{Z}}$ um
\oldpage[0\textsubscript{III}]{78}
sistema projetivo estrito de $S$--$\cat{C}$-módulos filtrados, para a
filtração $\mathfrak{J}$-ádica sobre $S$, limitado inferiormente.
Para um inteiro fixo $n$, suponhamos que seja satisfeita a condição seguinte:
\begin{itemize}
  \item[$(\mathrm{F}_n)$]
  \emph{O $\mathrm{gr}^\bullet(S)$-módulo graduado}
  \[
    \mathrm{E}_1^{(m)}(\bb{A})
    =
    \bigl(
      \RR^m\mathrm{T}(\mathrm{gr}^p(\bb{A}))
    \bigr)_{p\in\bb{Z}}
  \]
  \emph{de \sref{0.13.7.3.1} é de tipo finito para
  $m=n$ e $m=n+1$.}
\end{itemize}
Então:
\begin{enumerate}
  \item[{\rm(i)}] Os sistemas projetivos
  $(\RR^n\mathrm{T}(A_k))_{k\in\bb{Z}}$ e
  $(\RR^{n+1}\mathrm{T}(A_k))_{k\in\bb{Z}}$ satisfazem {\rm(ML)}.
  \item[{\rm(ii)}] Se
  \[
    \RR'^{\,n}\mathrm{T}(\bb{A})
    =
    \varprojlim_k\RR^n\mathrm{T}(A_k),
  \]
  então $\RR'^{\,n}\mathrm{T}(\bb{A})$ é um $S$-módulo de tipo
  finito.
  \item[{\rm(iii)}] A filtração sobre
  $\RR'^{\,n}\mathrm{T}(\bb{A})$ definida por
  \[
    \mathrm{F}^p\bigl(\RR'^{\,n}\mathrm{T}(\bb{A})\bigr)
    =
    \Ker\bigl(
      \RR'^{\,n}\mathrm{T}(\bb{A})
      \longrightarrow
      \RR^n\mathrm{T}(A_{p-1})
    \bigr)
    \qquad(p\in\bb{Z})
  \]
  é $\mathfrak{J}$-boa; isto é,
  \[
    \mathfrak{J}\,
    \mathrm{F}^p\bigl(\RR'^{\,n}\mathrm{T}(\bb{A})\bigr)
    \subset
    \mathrm{F}^{p+1}\bigl(\RR'^{\,n}\mathrm{T}(\bb{A})\bigr)
  \]
  para todo $p$, com igualdade para todo $p$ suficientemente grande.
  Em particular, a topologia sobre
  $\RR'^{\,n}\mathrm{T}(\bb{A})$ definida por esta filtração é a
  topologia $\mathfrak{J}$-ádica.
  \item[{\rm(iv)}] O sistema projetivo
  $(\mathrm{E}_r^{pq}(\bb{A}))_r$ é essencialmente constante sempre que
  $p+q=n$ ou $p+q=n+1$.
  Assim, $\mathrm{E}_\infty^{pq}(\bb{A})$ está definido
  (\sref{0.13.5.7}), e os isomorfismos
  \begin{equation}
  \label{0.13.7.7.1}
  \tag{13.7.7.1}
    \mathrm{gr}^p\bigl(\RR'^{\,n}\mathrm{T}(\bb{A})\bigr)
    \isoto
    \mathrm{E}_\infty^{p,n-p}(\bb{A})
    \qquad(p\in\bb{Z})
  \end{equation}
  formam um isomorfismo de
  $\mathrm{gr}^\bullet(S)$-módulos graduados.
\end{enumerate}
\end{proposition}

\begin{proof}
Em vista de \sref{0.13.7.7.1}, e levando em conta os
isomorfismos \sref{0.13.7.4.1}, por abuso de notação escreveremos também
$\RR^n\mathrm{T}(\bb{A})$ para o limite projetivo
$\RR'^{\,n}\mathrm{T}(\bb{A})$ do sistema projetivo
$(\RR^n\mathrm{T}(A_k))$.

O anel graduado $\mathrm{gr}^\bullet(S)$ é noetheriano
(Bourbaki, \emph{Alg.\ comm.}, cap.~III, \textsection2,
n\textsuperscript{o}~9, cor.~5 do teorema~2).
Portanto, a sucessão crescente de
$\mathrm{gr}^\bullet(S)$-submódulos graduados
$\mathrm{B}_r^{(m)}(\bb{A})$ de
$\mathrm{E}_1^{(m)}(\bb{A})$ é estacionária para $m=n$ e
$m=n+1$.
A condição {\rm(b)} de \sref{0.11.1.10} fica, portanto, satisfeita.
Segue-se que a condição {\rm(a)} de \sref{0.13.7.4} é satisfeita para
$n$ e para $n+1$, o que demonstra {\rm(i)}.

Além disso, os isomorfismos \sref{0.13.7.4.1}, junto com as
observações de \sref{0.13.7.6}, mostram que
$\mathrm{gr}^\bullet(\RR^n\mathrm{T}(\bb{A}))$ é um
$\mathrm{gr}^\bullet(S)$-módulo graduado isomorfo a
\[
  \mathrm{E}_\infty^{(n)}(\bb{A})
  =
  \mathrm{Z}_\infty^{(n)}(\bb{A})
  /
  \mathrm{B}_\infty^{(n)}(\bb{A}).
\]
Como $\mathrm{Z}_\infty^{(n)}(\bb{A})$ é um submódulo de
$\mathrm{E}_1^{(n)}(\bb{A})$, é de tipo finito, e também o é
$\mathrm{E}_\infty^{(n)}(\bb{A})$.
Finalmente, para a filtração
$(\mathrm{F}^p(\RR'^{\,n}\mathrm{T}(\bb{A})))$, segue-se de
\sref{0.13.4.5} que
\[
  \mathrm{gr}^\bullet\bigl(\RR^n\mathrm{T}(\bb{A})\bigr)
  \quad\hbox{e}\quad
  \mathrm{gr}^\bullet\bigl(\RR'^{\,n}\mathrm{T}(\bb{A})\bigr)
\]
são $\mathrm{gr}^\bullet(S)$-módulos graduados isomorfos.
Isto demonstra {\rm(iv)}.
As asserções {\rm(ii)} e {\rm(iii)} se seguirão destes resultados
e do lema seguinte.
\end{proof}

\begin{lemma}[13.7.7.2]
\label{0.13.7.7.2}
Seja $S$ um anel noetheriano $\mathfrak{J}$-ádico, e seja $M$
um $S$-módulo dotado de uma filtração codiscreta
$(\mathrm{F}^p(M))_{p\in\bb{Z}}$ tal que
\[
  \mathfrak{J}\mathrm{F}^p(M)
  \subset
  \mathrm{F}^{p+1}(M).
\]
Assim, $M$ é um módulo filtrado sobre $S$ dotado de seu
filtração $\mathfrak{J}$-ádica.
Suponhamos além disso que $M$ é separado para a topologia definida por
$(\mathrm{F}^p(M))$.
Então as condições seguintes são equivalentes:
\begin{enumerate}
  \item[{\rm a)}] $M$ é um $S$-módulo de tipo finito e
  $(\mathrm{F}^p(M))$ é uma filtração $\mathfrak{J}$-boa.
  \item[{\rm b)}] $\mathrm{gr}^\bullet(M)$ é um
  $\mathrm{gr}^\bullet(S)$-módulo de tipo finito.
  \item[{\rm c)}] Os $\mathrm{gr}^p(M)$ são $S$-módulos de
  tipo finito e, para todo $p$ suficientemente grande, os homomorfismos
  canônicos
  \begin{equation}
  \label{0.13.7.7.3}
  \tag{13.7.7.3}
    \mathfrak{J}\otimes_S\mathrm{gr}^p(M)
    \longrightarrow
    \mathrm{gr}^{p+1}(M)
  \end{equation}
  \oldpage[0\textsubscript{III}]{79}
  são sobrejetivos.
  Estão induzidos por
  $\mathfrak{J}\otimes_S\mathrm{F}^p(M)
  \longrightarrow\mathrm{F}^{p+1}(M)$, levando em conta o fato
  de que a imagem do homomorfismo composto
  \[
    \mathfrak{J}\otimes_S\mathrm{F}^{p+1}(M)
    \longrightarrow
    \mathfrak{J}\otimes_S\mathrm{F}^p(M)
    \longrightarrow
    \mathrm{F}^{p+1}(M)
  \]
  é
  $\mathfrak{J}\mathrm{F}^{p+1}(M)
  \subset\mathrm{F}^{p+2}(M)$.
\end{enumerate}

Para a demonstração, veja-se Bourbaki, \emph{Alg.\ comm.}, cap.~III,
\textsection3, n\textsuperscript{o}~1, prop.~3.
\end{lemma}

\begin{env}[13.7.7.4]
\label{0.13.7.7.4}
Para aplicar o lema \sref{0.13.7.7.2}, resta observar que a
topologia sobre $\RR'^{\,n}\mathrm{T}(\bb{A})$ definida pela filtração
considerada o torna um $S$-módulo separado e completo:
esta é a topologia de limite projetivo dos grupos discretos
$\RR^n\mathrm{T}(A_k)$.
Por outro lado, se $A_k=0$ para $k<k_0$, então também
$\RR^n\mathrm{T}(A_k)=0$ para $k<k_0$; portanto
\[
  \mathrm{F}^{k_0}\bigl(\RR'^{\,n}\mathrm{T}(\bb{A})\bigr)
  =
  \RR'^{\,n}\mathrm{T}(\bb{A}).
\]
Assim, todas as hipóteses do lema ficam efetivamente satisfeitas.
\end{env}

\begin{corollary}[13.7.8]
\label{0.13.7.8}
Se a hipótese $(\mathrm{F}_n)$ é satisfeita, então para todo $f\in S$
existe um isomorfismo canônico
\begin{equation}
\label{0.13.7.8.1}
\tag{13.7.8.1}
  \varprojlim_k\bigl(\RR^n\mathrm{T}(A_k)\bigr)_f
  \isoto
  \RR^n\mathrm{T}(\bb{A})\otimes_S S_{\{f\}}.
\end{equation}
\end{corollary}

\begin{proof}
Com efeito, $\RR^n\mathrm{T}(\bb{A})$ é um $S$-módulo de tipo finito,
e $S_{\{f\}}$ é uma $S$-álgebra ádica noetheriana
(\sref[0\textsubscript{I}]{0.7.6.11}), a saber, o completamento separado
de $S_f$ para a topologia $\mathfrak{J}$-pré-ádica
(\sref[0\textsubscript{I}]{0.7.6.2}).
Segue-se de \sref[0\textsubscript{I}]{0.7.7.8} e
\sref[0\textsubscript{I}]{0.7.7.1} que
\[
  \RR^n\mathrm{T}(\bb{A})\otimes_S S_{\{f\}}
\]
é isomorfo ao completamento separado de
$\RR^n\mathrm{T}(\bb{A})\otimes_S S_f$ para a
topologia $\mathfrak{J}$-pré-ádica.
Um sistema fundamental de vizinhanças de zero para esta topologia
é
\[
  \bigl(
    \mathfrak{J}^p\RR^n\mathrm{T}(\bb{A})
  \bigr)\otimes_S S_f;
\]
os grupos
\[
  \mathrm{F}^p\bigl(\RR^n\mathrm{T}(\bb{A})\bigr)
  \otimes_S S_f
\]
formam, portanto, também um sistema tal.
Este último grupo é o núcleo da aplicação canônica
\[
  \bigl(\RR^n\mathrm{T}(\bb{A})\bigr)_f
  \longrightarrow
  \bigl(\RR^n\mathrm{T}(A_{p-1})\bigr)_f.
\]
Por conseguinte, o grupo separado associado a
$\RR^n\mathrm{T}(\bb{A})\otimes_S S_f$ identifica-se com um subgrupo
$G$ de
\[
  \varprojlim_k\bigl(\RR^n\mathrm{T}(A_k)\bigr)_f.
\]
Mas o sistema projetivo
$((\RR^n\mathrm{T}(A_k))_f)_k$ satisfaz manifestamente {\rm(ML)}, e a
imagem de $(\RR^n\mathrm{T}(\bb{A}))_f$ em cada
$(\RR^n\mathrm{T}(A_k))_f$ é igual à imagem comum dos
$(\RR^n\mathrm{T}(A_h))_f$ para todo $h\geq k$ suficientemente grande.
Segue-se imediatamente que $G$ é \emph{denso em toda parte} em
$\varprojlim_k(\RR^n\mathrm{T}(A_k))_f$.
Como este último grupo é separado e completo, fica demonstrado o
corolário.
\end{proof}


\clearpage
\phantomsection
\label{section:ega3}
\pdfbookmark[0]{Capítulo III}{chapter.III}
\addcontentsline{toc}{part}{Capítulo III: Estudo cohomológico dos feixes coerentes}
\section*{CAPÍTULO III}
\setcounter{section}{0}
\setcounter{subsection}{0}
\setcounter{equation}{0}
\begin{center}
\Large\bfseries ESTUDO COHOMOLÓGICO DOS FEIXES COERENTES
\end{center}
\oldpage[III]{81}

\section*{Sumário}

\begin{longtable}{@{}p{0.08\textwidth}p{0.86\textwidth}@{}}
  \hyperref[section:III.1]{\textsection1}. & Cohomologia dos esquemas afins.\\
  \hyperref[section:III.2]{\textsection2}. & Estudo cohomológico dos morfismos projetivos.\\
  \hyperref[section:III.3]{\textsection3}. & Teorema de finitude para os morfismos próprios.\\
  \hyperref[section:III.4]{\textsection4}. & Teorema fundamental dos morfismos próprios. Aplicações.\\
  \hyperref[section:III.5]{\textsection5}. & Um teorema de existência para os feixes algébricos coerentes.\\
  \textsection6. & Funtores $\Tor$ e $\Ext$ locais e globais; fórmula de Künneth.\\
  \textsection7. & Mudança de base para os funtores homológicos de feixes de módulos.\\
  \textsection8. & Teorema de dualidade para os fibrados projetivos.\\
  \textsection9. & Cohomologia relativa e cohomologia local; dualidade local.\\
  \textsection10. & Relações entre a cohomologia projetiva e a cohomologia local. Técnica de completamento formal ao longo de um divisor.\\
  \textsection11. & Grupos de Picard globais e locais.\footnotemark
\end{longtable}
\footnotetext{EGA IV não depende das \textsection\textsection8--11 e provavelmente será publicado antes destes capítulos. \emph{[Trad.] Estes quatro últimos capítulos nunca foram publicados.}}
\bigskip

Este capítulo apresenta os teoremas fundamentais relativos à cohomologia dos feixes algébricos coerentes, com exceção dos teoremas que explicam a teoria dos resíduos (teoremas de dualidade), que serão objeto de um capítulo posterior.
Entre todos os resultados reunidos aqui há essencialmente seis teoremas fundamentais, cada um dos quais constitui o objeto de uma das seis primeiras seções.
Esses resultados serão ferramentas essenciais em tudo o que segue, inclusive para questões cuja natureza não é propriamente cohomológica;
o leitor verá os primeiros exemplos a partir do \hyperref[section:III.4]{\textsection4}.
A \textsection7 contém resultados de caráter mais técnico, mas de uso constante nas aplicações.
Finalmente, nas \textsection\textsection8--11 desenvolveremos certos resultados, ligados à dualidade dos feixes coerentes, particularmente importantes para as aplicações e que podem ser expostos antes mesmo de introduzir a teoria geral completa dos resíduos.

\oldpage[III]{82}
O conteúdo das \hyperref[section:III.1]{\textsection1} e \hyperref[section:III.2]{\textsection2} deve-se a J.-P.~Serre, e o leitor observará que não fizemos senão seguir (FAC).
As \textsection8 e 9 também se inspiram em (FAC) (embora as modificações exigidas pelos diferentes contextos sejam menos evidentes).
Finalmente, como dissemos na Introdução, a \hyperref[section:III.4]{\textsection4} deve ser considerada como a formalização, em linguagem moderna, do ``teorema de invariância'' fundamental da ``teoria das funções holomorfas'' de Zariski.

Chamamos a atenção para o fato de que os resultados do n.\textsuperscript{o}\hyperref[subsection:III.3.4]{3.4} (e as proposições preliminares de \sref[0]{0.13.4} a \sref[0]{0.13.7}) não serão utilizados no restante deste \hyperref[section:ega3]{Capítulo~III}, podendo portanto ser omitidos numa primeira leitura.
\bigskip

\section{Cohomologia dos esquemas afins}
\label{section:III.1}

\subsection{Revisão do complexo da álgebra exterior}
\label{subsection:III.1.1}

\begin{env}[1.1.1]
\label{III.1.1.1}
Seja $A$ um anel e seja $\mathbf{f}=(f_i)_{1\leq i\leq r}$ um sistema de $r$ elementos de $A$.
O \emph{complexo da álgebra exterior $K_\bullet(\mathbf{f})$} correspondente a $\mathbf{f}$ é um complexo de cadeias (G, I, 2.2) definido do modo seguinte: o $A$-módulo graduado $K_\bullet(\mathbf{f})$ é igual à \emph{álgebra exterior $\wedge(A^r)$}, graduada da maneira usual, e a aplicação bordo é a \emph{multiplicação interior $i_\mathbf{f}$} por $\mathbf{f}$, considerado como elemento do dual $\dual{(A^r)}$; lembremos que $i_\mathbf{f}$ é uma \emph{antiderivação} de grau $-1$ de $\wedge(A^r)$, e que, se $(\mathbf{e}_i)_{1\leq i\leq r}$ é a base canônica de $A^r$, tem-se $i_\mathbf{f}(\mathbf{e}_i)=f_i$; a verificação da condição $i_\mathbf{f}\circ i_\mathbf{f}=0$ é imediata.

Uma definição equivalente é a seguinte: para cada $i$, considera-se um complexo de cadeias $K_\bullet(f_i)$ definido por $K_0(f_i)=K_1(f_i)=A$, $K_n(f_i)=0$ para $n\neq 0,1$; a aplicação bordo fica definida pela condição de que $d_1:A\to A$ seja a \emph{multiplicação por $f_i$}.
Toma-se então como $K_\bullet(\mathbf{f})$ o \emph{produto tensorial $K_\bullet(f_1)\otimes K_\bullet(f_2)\otimes\cdots\otimes K_\bullet(f_r)$} (G, I, 2.7), provido de seu grau total; a verificação do isomorfismo deste complexo com o definido anteriormente é imediata.
\end{env}

\begin{env}[1.1.2]
\label{III.1.1.2}
Para todo $A$-módulo $M$, define-se o \emph{complexo de cadeias}
\[
\label{III.1.1.2.1}
  K_\bullet(\mathbf{f},M)=K_\bullet(\mathbf{f})\otimes_A M
  \tag{1.1.2.1}
\]
e o \emph{complexo de cocadeias} (G, I, 2.2)
\[
\label{III.1.1.2.2}
  K^\bullet(\mathbf{f},M)=\Hom_A(K_\bullet(\mathbf{f}),M).
  \tag{1.1.2.2}
\]

Se $g$ é uma $k$-cocadeia deste último complexo e se põe
\[
  g(i_1,\dots,i_k)=g(\mathbf{e}_{i_1}\wedge\cdots\wedge\mathbf{e}_{i_k}),
\]
então $g$ identifica-se com uma aplicação \emph{alternada} de $[1,r]^k$ em $M$, e das definições anteriores deduz-se que
\[
\label{III.1.1.2.3}
  d^k g(i_1,i_2,\dots,i_{k+1})=\sum_{h=1}^{k+1}(-1)^{h-1}f_{i_h}g(i_1,\dots,\widehat{i_h},\dots,i_{k+1}).
  \tag{1.1.2.3}
\]
\end{env}

\begin{env}[1.1.3]
\label{III.1.1.3}
\oldpage[III]{83}
Dos complexos anteriores deduzem-se, como de costume, os \emph{$A$-módulos de homologia e de cohomologia} (G, I, 2.2)
\[
  \HH_\bullet(\mathbf{f},M)=\HH_\bullet(K_\bullet(\mathbf{f},M)),
  \tag{1.1.3.1}
\]
\[
  \HH^\bullet(\mathbf{f},M)=\HH^\bullet(K^\bullet(\mathbf{f},M)).
  \tag{1.1.3.2}
\]

Define-se um \emph{$A$-isomorfismo $K_i(\mathbf{f},M)\isoto K^{r-i}(\mathbf{f},M)$} que envia cada cadeia $z=\sum(\mathbf{e}_{i_1}\wedge\cdots\wedge\mathbf{e}_{i_k})\otimes z_{i_1,\dots,i_k}$ à cocadeia $g_z$ tal que $g_z(j_1,\dots,j_{r-k})=\varepsilon z_{i_1,\dots,i_k}$, onde $(j_h)_{1\leq h\leq r-k}$ é a sucessão estritamente crescente complementar à sucessão estritamente crescente $(i_h)_{1\leq h\leq k}$ em $[1,r]$ e $\varepsilon=(-1)^\nu$, sendo $\nu$ o número de inversões da permutação $i_1,\dots,i_k,j_1,\dots,j_{r-k}$ de $[1,r]$.
Verifica-se que $g_{dz}=d(g_z)$, o que dá um isomorfismo
\[
  \HH^i(\mathbf{f},M)\isoto\HH_{r-i}(\mathbf{f},M)\text{ para }0\leq i\leq r.
  \tag{1.1.3.3}\label{III.1.1.3.3}
\]

Neste capítulo se considerarão especialmente os módulos de cohomologia $\HH^\bullet(\mathbf{f},M)$.

Para um $\mathbf{f}$ dado, é imediato (G, I, 2.1) que $M\mapsto\HH^\bullet(\mathbf{f},M)$ é um \emph{funtor cohomológico} (T, II, 2.1) da categoria de $A$-módulos na categoria de $A$-módulos graduados, nulo nos graus $<0$ e $>r$.
Além disso, tem-se
\[
  \HH^0(\mathbf{f},M)=\Hom_A(A/(\mathbf{f}),M),
  \tag{1.1.3.4}
\]
denotando por $(\mathbf{f})$ o ideal de $A$ gerado por $f_1,\dots,f_r$; isto deduz-se imediatamente de (\hyperref[III.1.1.2.3]{1.1.2.3}), e é claro que $\HH^0(\mathbf{f},M)$ identifica-se com o submódulo de $M$ \emph{anulado por $(\mathbf{f})$}.
De maneira análoga, por (\hyperref[III.1.1.2.3]{1.1.2.3}) tem-se
\[
  \HH^r(\mathbf{f},M)=M/\bigg(\sum_{i=1}^r f_i M\bigg)=(A/(\mathbf{f}))\otimes_A M.
  \tag{1.1.3.5}\label{III.1.1.3.5}
\]

Utilizar-se-á o seguinte resultado conhecido, do qual se recordará uma demonstração para sermos completos:
\end{env}

\begin{proposition}[1.1.4]
\label{III.1.1.4}
Sejam $A$ um anel, $\mathbf{f}=(f_i)_{1\leq i\leq r}$ uma família finita de elementos de $A$ e $M$ um $A$-módulo.
Se, para $1\leq i\leq r$, a aplicação $z\mapsto f_i\cdot z$ sobre $M_{i-1}=M/(f_1 M+\cdots+f_{i-1}M)$ é injetiva, então $\HH^i(\mathbf{f},M)=0$ para $i\neq r$.
\end{proposition}

\begin{proof}
Basta provar que $\HH_i(\mathbf{f},M)=0$ para todo $i>0$, em virtude de (\hyperref[III.1.1.3.3]{1.1.3.3}).
Procede-se por indução sobre $r$, sendo trivial o caso $r=0$.
Coloque-se $\mathbf{f}'=(f_i)_{1\leq i\leq r-1}$; esta família satisfaz as condições do enunciado, de modo que, se pusermos $L_\bullet=K_\bullet(\mathbf{f}',M)$, por hipótese tem-se $\HH_i(L_\bullet)=0$ para $i>0$, e $\HH_0(L_\bullet)=M_{r-1}$ em virtude de (\hyperref[III.1.1.3.3]{1.1.3.3}) e (\hyperref[III.1.1.3.5]{1.1.3.5}).
Para abreviar, coloque-se $K_\bullet=K_\bullet(f_r)=K_0\oplus K_1$, com $K_0=K_1=A$ e $d_1:K_1\to K_0$ igual à multiplicação por $f_r$; por definição \sref{III.1.1.1}, tem-se $K_\bullet(\mathbf{f},M)=K_\bullet\otimes_A L_\bullet$.
Dispõe-se do seguinte lema:

\begin{lemma}[1.1.4.1]
\label{III.1.1.4.1}
Seja $K_\bullet$ um complexo de cadeias de $A$-módulos livres, nulo salvo nas dimensões $0$ e $1$.
Para todo complexo de cadeias $L_\bullet$ de $A$-módulos tem-se uma sucessão exata
\[
  0\to\HH_0(K_\bullet\otimes\HH_p(L_\bullet))\to\HH_p(K_\bullet\otimes L_\bullet)\to\HH_1(K_\bullet\otimes\HH_{p-1}(L_\bullet))\to 0
\]
para todo índice $p$.
\end{lemma}

\oldpage[III]{84}
Este é um caso particular de uma sucessão exata de termos de grau baixo da sucessão espectral de K\"unneth (M, XVII, 5.2 (a) e G, I, 5.5.2); pode demonstrar-se diretamente como segue.
Considerem-se $K_0$ e $K_1$ como complexos de cadeias (nulos nas dimensões $\neq 0$ e $\neq 1$, respectivamente); tem-se então uma sucessão exata de complexos
\[
  0\to K_0\otimes L_\bullet\to K_\bullet\otimes L_\bullet\to K_1\otimes L_\bullet\to 0,
\]
à qual pode aplicar-se a sucessão exata de homologia
\[
  \begin{aligned}
  \cdots&\to\HH_{p+1}(K_1\otimes L_\bullet)
  \xrightarrow{\partial}\HH_p(K_0\otimes L_\bullet)
  \to\HH_p(K_\bullet\otimes L_\bullet)\\
  &\to\HH_p(K_1\otimes L_\bullet)
  \xrightarrow{\partial}\HH_{p-1}(K_0\otimes L_\bullet)\to\cdots.
  \end{aligned}
\]
Mas é evidente que $\HH_p(K_0\otimes L_\bullet)=K_0\otimes\HH_p(L_\bullet)$ e $\HH_p(K_1\otimes L_\bullet)=K_1\otimes\HH_{p-1}(L_\bullet)$ para todo $p$; além disso, verifica-se imediatamente que o operador $\partial:K_1\otimes\HH_p(L_\bullet)\to K_0\otimes\HH_p(L_\bullet)$ não é senão $d_1\otimes 1$; o lema deduz-se assim da sucessão exata anterior e da definição de $\HH_0(K_\bullet\otimes\HH_p(L_\bullet))$ e $\HH_1(K_\bullet\otimes\HH_{p-1}(L_\bullet))$.

Estabelecido o lema, o final da demonstração da Proposição \sref{III.1.1.4} é imediato: a hipótese de indução e o Lema \sref{III.1.1.4.1} dão $\HH_p(K_\bullet\otimes L_\bullet)=0$ para $p\geq 2$; além disso, se provarmos que $\HH_1(K_\bullet,\HH_0(L_\bullet))=0$, deduz-se também do Lema \sref{III.1.1.4.1} que $\HH_1(K_\bullet\otimes L_\bullet)=0$; mas, por definição, $\HH_1(K_\bullet,\HH_0(L_\bullet))$ não é senão o núcleo da aplicação $z\mapsto f_r\cdot z$ sobre $M_{r-1}$, e, como por hipótese este núcleo é nulo, a demonstração está concluída.
\end{proof}

\begin{env}[1.1.5]
\label{III.1.1.5}
Seja $\mathbf{g}=(g_i)_{1\leq i\leq r}$ uma segunda sucessão de $r$ elementos de $A$, e coloque-se $\mathbf{f}\mathbf{g}=(f_i g_i)_{1\leq i\leq r}$.
Pode-se definir um homomorfismo canônico de complexos
\[
  \vphi_\mathbf{g}:K_\bullet(\mathbf{f}\mathbf{g})\to K_\bullet(\mathbf{f})
  \tag{1.1.5.1}
\]
como a extensão canônica à álgebra exterior $\wedge(A^r)$ da aplicação $A$-linear $(x_1,\dots,x_r)\mapsto(g_1 x_1,\dots,g_r x_r)$ de $A^r$ em si mesmo.
Para ver que se trata de um homomorfismo de complexos, basta observar, em geral, que, se $u:E\to F$ é uma aplicação $A$-linear, $\mathbf{x}\in\dual{F}$ e $\mathbf{y}={}^t u(\mathbf{x})\in\dual{E}$, tem-se a fórmula
\[
  (\wedge u)\circ i_\mathbf{y}=i_\mathbf{x}\circ(\wedge u);
  \tag{1.1.5.2}
\]
com efeito, ambos os membros são antiderivações de $\wedge F$, e basta verificar que coincidem sobre $F$, o que se deduz imediatamente das definições.

Quando se identifica $K_\bullet(\mathbf{f})$ com o produto tensorial dos $K_\bullet(f_i)$ \sref{III.1.1.1}, $\vphi_\mathbf{g}$ é o produto tensorial dos $\vphi_{g_i}$, onde $\vphi_{g_i}$ é a identidade em grau $0$ e a multiplicação por $g_i$ em grau $1$.
\end{env}

\begin{env}[1.1.6]
\label{III.1.1.6}
Em particular, para todo par de inteiros $m$ e $n$ tais que $0\leq n\leq m$, têm-se homomorfismos de complexos
\[
  \vphi_{\mathbf{f}^{m-n}}:K_\bullet(\mathbf{f}^m)\to K_\bullet(\mathbf{f}^n)
  \tag{1.1.6.1}
\]
e, por conseguinte, homomorfismos
\[
  \vphi_{\mathbf{f}^{m-n}}:K^\bullet(\mathbf{f}^n,M)\to K^\bullet(\mathbf{f}^m,M),
  \tag{1.1.6.2}
\]
\[
  \vphi_{\mathbf{f}^{m-n}}:\HH^\bullet(\mathbf{f}^n,M)\to\HH^\bullet(\mathbf{f}^m,M).
  \tag{1.1.6.3}
\]

\oldpage[III]{85}
Estes últimos homomorfismos satisfazem evidentemente a condição de transitividade $\vphi_{\mathbf{f}^{m-p}}=\vphi_{\mathbf{f}^{m-n}}\circ\vphi_{\mathbf{f}^{n-p}}$ para $p\leq n\leq m$; definem, portanto, dois \emph{sistemas indutivos} de $A$-módulos; põe-se
\[
  C^\bullet((\mathbf{f}),M)=\varinjlim_n K^\bullet(\mathbf{f}^n,M),
  \tag{1.1.6.4}
\]
\[
  \HH^\bullet((\mathbf{f}),M)=\HH^\bullet(C^\bullet((\mathbf{f}),M))=\varinjlim_n\HH^\bullet(\mathbf{f}^n,M),
  \tag{1.1.6.5}
\]
deduzindo-se a última igualdade do fato de que o passo ao limite indutivo comuta com o funtor $\HH^\bullet$ (G, I, 2.1).
Ver-se-á mais adiante \sref{III.1.4.3} que $\HH^\bullet((\mathbf{f}),M)$ não depende do \emph{ideal $(\mathbf{f})$} de $A$ (e, de maneira análoga, da topologia $(\mathbf{f})$-pré-ádica sobre $A$), o que justifica as notações.

É claro que $M\mapsto C^\bullet((\mathbf{f}),M)$ é um funtor exato $A$-linear e que $M\mapsto\HH^\bullet((\mathbf{f}),M)$ é um funtor cohomológico.
\end{env}

\begin{env}[1.1.7]
\label{III.1.1.7}
Coloquem-se $\mathbf{f}=(f_i)\in A^r$ e $\mathbf{g}=(g_i)\in A^r$; denote-se por $e_\mathbf{g}$ a multiplicação à esquerda pelo vetor $\mathbf{g}\in A^r$ sobre a álgebra exterior $\wedge(A^r)$; sabe-se que se tem a \emph{fórmula de homotopia}
\[
  i_\mathbf{f}e_\mathbf{g}+e_\mathbf{g}i_\mathbf{f}=\langle\mathbf{g},\mathbf{f}\rangle 1
  \tag{1.1.7.1}
\]
no $A$-módulo $\wedge(A^r)$ ($1$ denota o automorfismo identidade de $\wedge(A^r)$); esta relação implica também que \emph{no complexo $K_\bullet(\mathbf{f})$} tem-se
\[
  de_\mathbf{g}+e_\mathbf{g}d=\langle\mathbf{g},\mathbf{f}\rangle 1.
  \tag{1.1.7.2}
\]

Se o ideal $(\mathbf{f})$ é igual a $A$, existe $\mathbf{g}\in A^r$ tal que $\langle\mathbf{g},\mathbf{f}\rangle=\sum_{i=1}^r g_i f_i=1$.
Por conseguinte (G, I, 2.4):
\end{env}

\begin{proposition}[1.1.8]
\label{III.1.1.8}
Suponha-se que o ideal $(\mathbf{f})$ gerado pelos $f_i$ é igual a $A$.
Então o complexo $K_\bullet(\mathbf{f})$ é homotopicamente trivial, assim como os complexos $K_\bullet(\mathbf{f},M)$ e $K^\bullet(\mathbf{f},M)$ para todo $A$-módulo $M$.
\end{proposition}

\begin{corollary}[1.1.9]
\label{III.1.1.9}
Se $(\mathbf{f})=A$, então $\HH^\bullet(\mathbf{f},M)=0$ e $\HH^\bullet((\mathbf{f}),M)=0$ para todo $A$-módulo $M$.
\end{corollary}

\begin{proof}
Com efeito, tem-se então $(\mathbf{f}^n)=A$ para todo $n$.
\end{proof}

\begin{remark}[1.1.10]
\label{III.1.1.10}
Com as mesmas notações anteriores, coloque-se $X=\Spec(A)$ e seja $Y$ o subpré-esquema fechado de $X$ definido pelo ideal $(\mathbf{f})$.
Provar-se-á em \textsection9 que $\HH^\bullet((\mathbf{f}),M)$ é isomorfo à cohomologia $\HH_Y^\bullet(X,\widetilde{M})$ correspondente ao antifiltro $\Phi$ de subconjuntos fechados de $Y$ (T, 3.2).
Mostrar-se-á também que a Proposição \sref{III.1.2.3}, aplicada a $X$ e a $\sh{F}=\widetilde{M}$, é um caso particular de uma sucessão exata de cohomologia
\[
  \cdots\to\HH_Y^p(X,\sh{F})\to\HH^p(X,\sh{F})\to\HH^p(X\setmin Y,\sh{F})\to\HH_Y^{p+1}(X,\sh{F})\to\cdots.
\]
\end{remark}

\subsection{Cohomologia de \v Cech de um recobrimento aberto}
\label{subsection:III.1.2}

\begin{notation}[1.2.1]
\label{III.1.2.1}
Nesta seção, denota-se por:
\begin{enumerate}
  \item $X$ um pré-esquema;
  \item $\sh{F}$ um $\sh{O}_X$-módulo quase-coerente;
\oldpage[III]{86}
  \item $A=\Gamma(X,\sh{O}_X)$, $M=\Gamma(X,\sh{F})$;
  \item $\mathbf{f}=(f_i)_{1\leq i\leq r}$ um sistema finito de elementos de $A$;
  \item $U_i=X_{f_i}$, o aberto \sref[0]{0.5.5.2} dos $x\in X$ tais que $f_i(x)\neq 0$;
  \item $U=\bigcup_{i=1}^r U_i$;
  \item $\mathfrak{U}$ o recobrimento $(U_i)_{1\leq i\leq r}$ de $U$.
\end{enumerate}
\end{notation}

\begin{env}[1.2.2]
\label{III.1.2.2}
Suponha-se que $X$ é um pré-esquema cujo espaço subjacente é \emph{noetheriano}, ou um \emph{esquema} cujo espaço subjacente é \emph{quase-compacto}.
Sabe-se então \sref[I]{I.9.3.3} que $\Gamma(U_i,\sh{F})=M_{f_i}$.
Coloque-se
\[
  U_{i_0 i_1\cdots i_p}=\bigcap_{k=0}^p U_{i_k}=X_{f_{i_0}f_{i_1}\cdots f_{i_p}}
\]
\sref[0]{0.5.5.3}; tem-se, portanto, também
\[
  \Gamma(U_{i_0 i_1\cdots i_p},\sh{F})=M_{f_{i_0}f_{i_1}\cdots f_{i_p}}.
  \tag{1.2.2.1}
\]

Por \sref[0]{0.1.6.1}, $M_{f_{i_0}f_{i_1}\cdots f_{i_p}}$ identifica-se com o limite indutivo $\varinjlim_n M_{i_0 i_1\cdots i_p}^{(n)}$, onde o sistema indutivo é formado pelos $M_{i_0 i_1\cdots i_p}^{(n)}=M$, sendo os homomorfismos $\vphi_{nm}:M_{i_0 i_1\cdots i_p}^{(m)}\to M_{i_0 i_1\cdots i_p}^{(n)}$ a multiplicação por $(f_{i_0}f_{i_1}\cdots f_{i_p})^{n-m}$ para $m\leq n$.
Denota-se por $C_n^p(M)$ o conjunto de aplicações \emph{alternadas} de $[1,r]^{p+1}$ em $M$ (para todo $n$); estes $A$-módulos formam também um sistema indutivo com respeito aos $\vphi_{nm}$.
Se $C^p(\mathfrak{U},\sh{F})$ é o grupo das $p$-cocadeias \emph{alternadas} de \v Cech relativas ao recobrimento $\mathfrak{U}$, com coeficientes em $\sh{F}$ (G, II, 5.1), do que precede deduz-se que pode escrever-se
\[
  C^p(\mathfrak{U},\sh{F})=\varinjlim_n C_n^p(M).
  \tag{1.2.2.2}
\]

Com as notações de \sref{III.1.1.2}, $C_n^p(M)$ identifica-se com $K^{p+1}(\mathbf{f}^n,M)$, e a aplicação $\vphi_{nm}$ identifica-se com a aplicação $\vphi_{\mathbf{f}^{n-m}}$ definida em \sref{III.1.1.6}.
Tem-se assim, para todo $p\geq 0$, um isomorfismo canônico funtorial
\[
  C^p(\mathfrak{U},\sh{F})\isoto C^{p+1}((\mathbf{f}),M).
  \tag{1.2.2.3}\label{III.1.2.2.3}
\]

Além disso, a fórmula (\hyperref[III.1.1.2.3]{1.1.2.3}) e a definição da cohomologia de um recobrimento (G, II, 5.1) mostram que os isomorfismos (\hyperref[III.1.2.2.3]{1.2.2.3}) são compatíveis com as aplicações cobordo.
\end{env}

\begin{proposition}[1.2.3]
\label{III.1.2.3}
Se $X$ é um pré-esquema cujo espaço subjacente é noetheriano, ou um esquema cujo espaço subjacente é quase-compacto, existe um isomorfismo canônico funtorial em $\sh{F}$
\[
  \HH^p(\mathfrak{U},\sh{F})\isoto\HH^{p+1}((\mathbf{f}),M)\text{ para }p\geq 1.
  \tag{1.2.3.1}\label{III.1.2.3.1}
\]

Além disso, tem-se uma sucessão exata funtorial em $\sh{F}$
\[
  0\to\HH^0((\mathbf{f}),M)\to M\to\HH^0(\mathfrak{U},\sh{F})\to\HH^1((\mathbf{f}),M)\to 0.
  \tag{1.2.3.2}\label{III.1.2.3.2}
\]
\end{proposition}

\oldpage[III]{87}
\begin{proof}
Os isomorfismos (\hyperref[III.1.2.3.1]{1.2.3.1}) são consequências imediatas do que foi visto em \sref{III.1.2.2}.
Por outro lado, tem-se $C^0(\mathfrak{U},\sh{F})=C^1((\mathbf{f}),M)$; por conseguinte, $\HH^0(\mathfrak{U},\sh{F})$ identifica-se com o subgrupo dos $1$-cociclos de $C^1((\mathbf{f}),M)$; como $M=C^0((\mathbf{f}),M)$, a sucessão exata (\hyperref[III.1.2.3.2]{1.2.3.2}) não é senão a dada pela definição dos grupos de cohomologia $\HH^0((\mathbf{f}),M)$ e $\HH^1((\mathbf{f}),M)$.
\end{proof}

\begin{corollary}[1.2.4]
\label{III.1.2.4}
Suponha-se que os $X_{f_i}$ são quase-compactos e que existem $g_i\in\Gamma(U,\sh{O}_X)$ tais que $\sum_i g_i(f_i|U)=1|U$.
Então, para todo $(\sh{O}_X|U)$-módulo quase-coerente $\sh{G}$, tem-se $\HH^p(\mathfrak{U},\sh{G})=0$ para $p>0$; se além disso $U=X$, o homomorfismo canônico (\hyperref[III.1.2.3.2]{1.2.3.2}) $\Gamma(X,\sh{G})\to\HH^0(\mathfrak{U},\sh{G})$ é bijetivo.
\end{corollary}

\begin{proof}
Como, por hipótese, os $U_i=X_{f_i}$ são quase-compactos, também o é $U$, e pode reduzir-se ao caso $U=X$; pondo $N=\Gamma(X,\sh{G})$, a hipótese implica então que $\HH^p((\mathbf{f}),N)=0$ para todo $p\geq 0$ \sref{III.1.1.9}.
O corolário deduz-se imediatamente de (\hyperref[III.1.2.3.1]{1.2.3.1}) e (\hyperref[III.1.2.3.2]{1.2.3.2}).
\end{proof}

Observe-se que, dado que $\HH^0(\mathfrak{U},\sh{F})=\HH^0(U,\sh{F})$ (G, II, 5.2.2), demonstrou-se novamente \sref[I]{I.1.3.7} como caso particular.

\begin{remark}[1.2.5]
\label{III.1.2.5}
Suponha-se que $X$ é um \emph{esquema afim}; então os $U_i=X_{f_i}=D(f_i)$ são abertos afins, assim como os $U_{i_0 i_1\cdots i_p}$ (mas $U$ não é necessariamente afim).
Neste caso, os funtores $\Gamma(X,\sh{F})$ e $\Gamma(U_{i_0 i_1\cdots i_p},\sh{F})$ são exatos em $\sh{F}$ \sref[I]{I.1.3.11}.
Se houver uma sucessão exata $0\to\sh{F}'\to\sh{F}\to\sh{F}''\to 0$ de $\sh{O}_X$-módulos quase-coerentes, a sucessão de complexos
\[
  0\to C^\bullet(\mathfrak{U},\sh{F}')\to C^\bullet(\mathfrak{U},\sh{F})\to C^\bullet(\mathfrak{U},\sh{F}'')\to 0
\]
é exata e dá, portanto, uma sucessão exata de cohomologia
\[
  \cdots\to\HH^p(\mathfrak{U},\sh{F}')\to\HH^p(\mathfrak{U},\sh{F})\to\HH^p(\mathfrak{U},\sh{F}'')\xrightarrow{\partial}\HH^{p+1}(\mathfrak{U},\sh{F}')\to\cdots.
\]
Por outro lado, se pusermos $M'=\Gamma(X,\sh{F}')$ e $M''=\Gamma(X,\sh{F}'')$, a sucessão $0\to M'\to M\to M''\to 0$ é exata; como $C^\bullet((\mathbf{f}),M)$ é um funtor exato em $M$, tem-se também a sucessão exata de cohomologia
\[
  \cdots\to\HH^p((\mathbf{f}),M')\to\HH^p((\mathbf{f}),M)\to\HH^p((\mathbf{f}),M'')\xrightarrow{\partial}\HH^{p+1}((\mathbf{f}),M')\to\cdots.
\]

Assim, como o diagrama
\[
  \xymatrix{
    0\ar[r] &
    C^\bullet(\mathfrak{U},\sh{F}')\ar[r]\ar[d] &
    C^\bullet(\mathfrak{U},\sh{F})\ar[r]\ar[d] &
    C^\bullet(\mathfrak{U},\sh{F}'')\ar[r]\ar[d] &
    0\\
    0\ar[r] &
    C^\bullet((\mathbf{f}),M')\ar[r] &
    C^\bullet((\mathbf{f}),M)\ar[r] &
    C^\bullet((\mathbf{f}),M'')\ar[r] &
    0
  }
\]
\oldpage[III]{88}
é comutativo, conclui-se que os diagramas
\[
  \xymatrix{
    \HH^p(\mathfrak{U},\sh{F}'')\ar[r]^\partial\ar[d] &
    \HH^{p+1}(\mathfrak{U},\sh{F}')\ar[d]\\
    \HH^{p+1}((\mathbf{f}),M'')\ar[r]^\partial &
    \HH^{p+2}((\mathbf{f}),M')
  }
  \tag{1.2.5.1}
\]
são comutativos para todo $p$~(G, I, 2.1.1).
\end{remark}

\subsection{Cohomologia de um esquema afim}
\label{subsection:III.1.3}

\begin{theorem}[1.3.1]
\label{III.1.3.1}
Seja $X$ um esquema afim.
Para todo $\sh{O}_X$-módulo quase-coerente $\sh{F}$, tem-se $\HH^p(X,\sh{F})=0$ para todo $p>0$.
\end{theorem}

\begin{proof}
Seja $\mathfrak{U}$ um recobrimento finito de $X$ pelos abertos afins $X_{f_i}=D(f_i)$ ($1\leq i\leq r$); sabe-se então que o ideal de $A=\Gamma(X,\sh{O}_X)$ gerado pelos $f_i$ é igual a $A$.
Do Corolário~\sref{III.1.2.4} conclui-se, pois, que $\HH^p(\mathfrak{U},\sh{F})=0$ para $p>0$.
Como existem recobrimentos finitos de $X$ por abertos afins arbitrariamente finos \sref[I]{I.1.1.10}, a definição da cohomologia de \v Cech (G, II, 5.8) mostra que também $\CHH^p(X,\sh{F})=0$ para $p>0$.
Mas isto se aplica também a todo pré-esquema $X_f$, para $f\in A$ \sref[I]{I.1.3.6}; portanto, $\CHH^p(X_f,\sh{F})=0$ para $p>0$.
Como $X_f\cap X_g=X_{fg}$, deduz-se também que $\HH^p(X,\sh{F})=0$ para todo $p>0$, em virtude de (G, II, 5.9.2).
\end{proof}

\begin{corollary}[1.3.2]
\label{III.1.3.2}
Seja $Y$ um pré-esquema e seja $f:X\to Y$ um morfismo afim \sref[II]{II.1.6.1}.
Para todo $\sh{O}_X$-módulo quase-coerente $\sh{F}$, tem-se $\RR^q f_*(\sh{F})=0$ para $q>0$.
\end{corollary}

\begin{proof}
Por definição, $\RR^q f_*(\sh{F})$ é o $\sh{O}_Y$-módulo associado ao pré-feixe $U\mapsto\HH^q(f^{-1}(U),\sh{F})$, onde $U$ percorre os abertos de $Y$.
Mas os abertos afins formam uma base de $Y$, e, para um aberto tal $U$, $f^{-1}(U)$ é afim \sref[II]{II.1.3.2}; portanto, $\HH^q(f^{-1}(U),\sh{F})=0$ pelo Teorema~\sref{III.1.3.1}, o que prova o corolário.
\end{proof}

\begin{corollary}[1.3.3]
\label{III.1.3.3}
Seja $Y$ um pré-esquema e seja $f:X\to Y$ um morfismo afim.
Para todo $\sh{O}_X$-módulo quase-coerente $\sh{F}$, o homomorfismo canônico $\HH^p(Y,f_*(\sh{F}))\to\HH^p(X,\sh{F})$ (\textbf{0},~12.1.3.1) é bijetivo para todo $p$.
\end{corollary}

\begin{proof}
Basta (por \sref[0]{0.12.1.7}) mostrar que os homomorfismos de bordo $''E_2^{p0}=\HH^p(Y,f_*(\sh{F}))\to\HH^p(X,\sh{F})$ da segunda sucessão espectral do funtor composto $\Gamma f_*$ são bijetivos.
Mas o termo $E_2$ desta sucessão espectral é dado por $''E_2^{pq}=\HH^p(Y,\RR^q f_*(\sh{F}))$ (G, II, 4.17.1); pelo Corolário~\sref{III.1.3.2}, $''E_2^{pq}=0$ para $q>0$, e a sucessão espectral degenera; daí a afirmação \sref[0]{0.11.1.6}.
\end{proof}

\begin{corollary}[1.3.4]
\label{III.1.3.4}
Sejam $f:X\to Y$ um morfismo afim e $g:Y\to Z$ um morfismo.
\oldpage[III]{89}
Para todo $\sh{O}_X$-módulo quase-coerente $\sh{F}$, o homomorfismo canônico $\RR^p g_*(f_*(\sh{F}))\to\RR^p(g\circ f)_*(\sh{F})$ (\textbf{0},~12.2.5.1) é bijetivo para todo $p$.
\end{corollary}

\begin{proof}
Basta observar que, segundo o Corolário~\sref{III.1.3.3}, para todo aberto afim $W$ de $Z$, o homomorfismo canônico $\HH^p(g^{-1}(W),f_*(\sh{F}))\to\HH^p(f^{-1}(g^{-1}(W)),\sh{F})$ é bijetivo; isto prova que o homomorfismo de pré-feixes que define o homomorfismo canônico $\RR^p g_*(f_*(\sh{F}))\to\RR^p(g\circ f)_*(\sh{F})$ é bijetivo \sref[0]{0.12.2.5}.
\end{proof}

\subsection{Aplicação à cohomologia de pré-esquemas arbitrários}
\label{subsection:III.1.4}

\begin{proposition}[1.4.1]
\label{III.1.4.1}
Seja $X$ um esquema e seja $\mathfrak{U}=(U_\alpha)$ um recobrimento de $X$ por abertos afins.
Para todo $\sh{O}_X$-módulo quase-coerente $\sh{F}$, os módulos de cohomologia $\HH^\bullet(X,\sh{F})$ e $\HH^\bullet(\mathfrak{U},\sh{F})$ (sobre $\Gamma(X,\sh{O}_X)$) são canonicamente isomorfos.
\end{proposition}

\begin{proof}
Como $X$ é um esquema, toda interseção finita $V$ de abertos do recobrimento $\mathfrak{U}$ é afim \sref[I]{I.5.5.6}; portanto, $\HH^q(V,\sh{F})=0$ para $q\geq 1$ pelo Teorema~\sref{III.1.3.1}.
A proposição deduz-se então de um teorema de Leray (G, II, 5.4.1).
\end{proof}

\begin{remark}[1.4.2]
\label{III.1.4.2}
Observe-se que o resultado da Proposição~\sref{III.1.4.1} continua válido quando as interseções finitas dos conjuntos $U_\alpha$ são afins, mesmo sem supor necessariamente que $X$ seja um esquema.
\end{remark}

\begin{corollary}[1.4.3]
\label{III.1.4.3}
Seja $X$ um esquema de espaço subjacente quase-compacto, seja $A=\Gamma(X,\sh{O}_X)$ e seja $\mathbf{f}=(f_i)_{1\leq i\leq r}$ uma sucessão finita de elementos de $A$ tal que os $X_{f_i}$ (notação de \sref{III.1.2.1}) sejam afins.
Então (com as notações de \sref{III.1.2.1}), para todo $\sh{O}_X$-módulo quase-coerente $\sh{F}$, tem-se um isomorfismo canônico funtorial em $\sh{F}$
\[
  \HH^q(U,\sh{F})\isoto\HH^{q+1}((\mathbf{f}),M)\text{ para }q\geq 1,
  \tag{1.4.3.1}\label{III.1.4.3.1}
\]
e uma sucessão exata funtorial em $\sh{F}$
\[
  0\to\HH^0((\mathbf{f}),M)\to M\to\HH^0(U,\sh{F})\to\HH^1((\mathbf{f}),M)\to 0.
  \tag{1.4.3.2}
\]
\end{corollary}

\begin{proof}
Isto deduz-se imediatamente das Proposições~\sref{III.1.4.1} e \sref{III.1.2.3}.
\end{proof}

\begin{env}[1.4.4]
\label{III.1.4.4}
Se $X$ é um \emph{esquema afim}, da Observação~\sref{III.1.2.5} e da Proposição~\sref{III.1.4.1} deduz-se que, para todo $q\geq 0$, os diagramas
\[
  \xymatrix{
    \HH^q(U,\sh{F}'')\ar[r]^\partial\ar[d] &
    \HH^{q+1}(U,\sh{F}')\ar[d]\\
    \HH^{q+1}((\mathbf{f}),M'')\ar[r]^\partial &
    \HH^{q+2}((\mathbf{f}),M')
  }
  \tag{1.4.4.1}
\]
correspondentes a uma sucessão exata $0\to\sh{F}'\to\sh{F}\to\sh{F}''\to 0$ de $\sh{O}_X$-módulos quase-coerentes (com as notações da Observação~\sref{III.1.2.5}) são comutativos.
\end{env}

\begin{proposition}[1.4.5]
\label{III.1.4.5}
\oldpage[III]{90}
Seja $X$ um esquema quase-compacto, seja $\sh{L}$ um $\sh{O}_X$-módulo invertível, e considere-se o anel graduado $A_\bullet=\Gamma_\bullet(\sh{L})$ \sref[0]{0.5.4.6}; então $\HH^\bullet(\sh{F},\sh{L})=\bigoplus_{n\in\bb{Z}}\HH^\bullet(X,\sh{F}\otimes\sh{L}^{\otimes n})$ é um $A_\bullet$-módulo graduado e, para todo $f\in A_n$, tem-se um isomorfismo canônico
\[
  \HH^\bullet(X_f,\sh{F})\isoto(\HH^\bullet(\sh{F},\sh{L}))_{(f)}
  \tag{1.4.5.1}\label{III.1.4.5.1}
\]
de $(A_\bullet)_{(f)}$-módulos.
\end{proposition}

\begin{proof}
Como $X$ é um esquema quase-compacto, pode calcular-se a cohomologia de todos os $\sh{O}_X$-módulos $\sh{F}\otimes\sh{L}^{\otimes n}$ utilizando um mesmo recobrimento finito $\mathfrak{U}=(U_i)$, formado por abertos afins tais que a restrição $\sh{L}|U_i$ seja isomorfa a $\sh{O}_X|U_i$ para cada $i$ \sref{III.1.4.1}.
É então imediato que os $U_i\cap X_f$ são abertos afins \sref[I]{I.1.3.6}, e pode calcular-se assim a cohomologia $\HH^\bullet(X_f,\sh{F}\otimes\sh{L}^{\otimes n})$ utilizando o recobrimento $\mathfrak{U}|X_f=(U_i\cap X_f)$ \sref{III.1.4.1}.
É imediato que, para todo $f\in A_n$, a multiplicação por $f$ define um homomorfismo $C^\bullet(\mathfrak{U},\sh{F}\otimes\sh{L}^m)\to C^\bullet(\mathfrak{U},\sh{F}\otimes\sh{L}^{\otimes(m+n)})$ e, portanto, um homomorfismo $\HH^\bullet(\mathfrak{U},\sh{F}\otimes\sh{L}^{\otimes m})\to\HH^\bullet(\mathfrak{U},\sh{F}\otimes\sh{L}^{\otimes(m+n)})$, o que estabelece a primeira afirmação.
Por outro lado, para um $f\in A_n$ dado, de \sref[I]{I.9.3.2} deduz-se que se tem um isomorfismo de complexos de $(A_\bullet)_{(f)}$-módulos
\[
  C^\bullet(\mathfrak{U}|X_f,\sh{F})\isoto\bigg(\!\!C^\bullet\bigg(\mathfrak{U},\bigoplus_{n\in\bb{Z}}\sh{F}\otimes\sh{L}^{\otimes n}\bigg)\!\!\bigg)_{(f)},
\]
levando em conta \sref[I]{I.1.3.9}[ii].
Ao passar à cohomologia destes complexos obtém-se o isomorfismo (\hyperref[III.1.4.5.1]{1.4.5.1}), recordando que o funtor $M\mapsto M_{(f)}$ é exato sobre a categoria de $A_\bullet$-módulos graduados.
\end{proof}

\begin{corollary}[1.4.6]
\label{III.1.4.6}
Suponha-se que sejam satisfeitas as hipóteses da Proposição~\sref{III.1.4.5} e, além disso, que $\sh{L}=\sh{O}_X$.
Se pusermos $A=\Gamma(X,\sh{O}_X)$, então, para todo $f\in A$, tem-se um isomorfismo canônico $\HH^\bullet(X_f,\sh{F})\isoto(\HH^\bullet(X,\sh{F}))_f$ de $A_f$-módulos.
\end{corollary}

\begin{corollary}[1.4.7]
\label{III.1.4.7}
Seja $X$ um esquema quase-compacto e seja $f$ um elemento de $\Gamma(X,\sh{O}_X)$.
\begin{enumerate}
  \item[{\rm(i)}] Suponha-se que o aberto $X_f$ é afim.
    Então, para todo $\sh{O}_X$-módulo quase-coerente $\sh{F}$, todo $i>0$ e todo $\xi\in\HH^i(X,\sh{F})$, existe um inteiro $n>0$ tal que $f^n\xi=0$.
  \item[{\rm(ii)}] Reciprocamente, suponha-se que $X_f$ é quase-compacto e que, para todo feixe quase-coerente de ideais $\sh{J}$ de $\sh{O}_X$ e todo $\zeta\in\HH^1(X,\sh{J})$, existe um $n>0$ tal que $f^n\zeta=0$.
    Então $X_f$ é afim.
\end{enumerate}
\end{corollary}

\begin{proof}
\medskip\noindent
\begin{enumerate}
  \item[(i)] Se $X_f$ é afim, tem-se $\HH^i(X_f,\sh{F})=0$ para todo $i>0$ \sref{III.1.3.1}, de modo que a afirmação deduz-se diretamente do Corolário~\sref{III.1.4.6}.
  \item[(ii)] Em virtude do critério de Serre \sref[II]{II.5.2.1}, basta provar que, para todo feixe quase-coerente de ideais $\sh{K}$ de $\sh{O}_X|X_f$, tem-se $\HH^1(X_f,\sh{K})=0$.
    Como $X_f$ é um aberto quase-compacto de um esquema quase-compacto $X$, existe um feixe quase-coerente de ideais $\sh{J}$ de $\sh{O}_X$ tal que $\sh{K}=\sh{J}|X_f$ \sref[I]{I.9.4.2}.
    Segundo o Corolário~\sref{III.1.4.6}, tem-se $\HH^1(X_f,\sh{K})=(\HH^1(X,\sh{J}))_f$, e a hipótese implica que o membro da direita é nulo, o que dá a afirmação.
\end{enumerate}
\end{proof}

\begin{remark}[1.4.8]
\label{III.1.4.8}
Observe-se que o Corolário~\sref{III.1.4.7}[i] dá uma demonstração mais simples da relação (\textbf{II},~4.5.13.2).
\end{remark}

\begin{lemma}[1.4.9]
\label{III.1.4.9}
\oldpage[III]{91}
Seja $X$ um esquema quase-compacto, seja $\mathfrak{U}=(U_i)_{1\leq i\leq n}$ um recobrimento finito de $X$ por abertos afins e seja $\sh{F}$ um $\sh{O}_X$-módulo quase-coerente.
O complexo de feixes $\sh{C}^\bullet(\mathfrak{U},\sh{F})$ definido pelo recobrimento $\mathfrak{U}$ (G, II, 5.2) é então um $\sh{O}_X$-módulo quase-coerente.
\end{lemma}

\begin{proof}
Das definições (G, II, 5.2) deduz-se que $\sh{C}^p(\mathfrak{U},\sh{F})$ é a soma direta dos feixes imagem direta dos $\sh{F}|U_{i_0\cdots i_p}$ pela imersão canônica $U_{i_0\cdots i_p}\to X$.
A hipótese de que $X$ é um esquema implica que estas imersões são morfismos afins \sref[I]{I.5.5.6}; portanto, os $\sh{C}^p(\mathfrak{U},\sh{F})$ são quase-coerentes \sref[II]{II.1.2.6}.
\end{proof}

\begin{proposition}[1.4.10]
\label{III.1.4.10}
Seja $u:X\to Y$ um morfismo separado e quase-compacto.
Para todo $\sh{O}_X$-módulo quase-coerente $\sh{F}$, os $\RR^q u_*(\sh{F})$ são $\sh{O}_Y$-módulos quase-coerentes.
\end{proposition}

\begin{proof}
A questão é local sobre $Y$, de modo que pode supor-se que $Y$ é afim.
Então $X$ é uma união finita de abertos afins $U_i$ ($1\leq i\leq n$); seja $\mathfrak{U}$ o recobrimento $(U_i)$.
Além disso, como $Y$ é um esquema, de \sref[I]{I.5.5.10} deduz-se que, para todo aberto afim $V\subset Y$, a imersão canônica $u^{-1}(V)\to X$ é um morfismo afim; conclui-se (Proposição~\sref{III.1.4.1} e (G, II, 5.2)) que se tem um isomorfismo canônico
\[
  \HH^\bullet(u^{-1}(V),\sh{F})\isoto\HH^\bullet(\Gamma(V,\sh{K}^\bullet)),
  \tag{1.4.10.1}\label{III.1.4.10.1}
\]
onde se põe $\sh{K}^\bullet=u_*(\sh{C}^\bullet(\mathfrak{U},\sh{F}))$.
Segundo o Lema~\sref{III.1.4.9} e \sref[I]{I.9.2.2}, $\sh{K}^\bullet$ é um $\sh{O}_Y$-módulo quase-coerente; além disso, constitui um \emph{complexo de feixes}, dado que também o é $\sh{C}^\bullet(\mathfrak{U},\sh{F})$.
Da definição da cohomologia $\sh{H}^\bullet(\sh{K}^\bullet)$ (G, II, 4.1) deduz-se então que esta última é formada por $\sh{O}_Y$-módulos quase-coerentes \sref[I]{I.4.1.1}.
Como (para $V$ afim em $Y$) o funtor $\Gamma(V,\sh{G})$ é exato em $\sh{G}$ sobre a categoria de $\sh{O}_Y$-módulos quase-coerentes, tem-se (G, II, 4.1)
\[
  \HH^\bullet(\Gamma(V,\sh{K}^\bullet))=\Gamma(V,\sh{H}^\bullet(\sh{K}^\bullet)).
  \tag{1.4.10.2}\label{III.1.4.10.2}
\]

Finalmente, observe-se que da definição do homomorfismo canônico
\[
  \HH^\bullet(\mathfrak{U},\sh{F})\to\HH^\bullet(X,\sh{F}),
\]
dado em (G, II, 5.2), deduz-se que, se $V'\subset V$ é outro aberto afim de $Y$, o
diagrama
\[
  \xymatrix{
    \HH^\bullet(u^{-1}(V),\sh{F})\ar[r]^\sim\ar[d] &
    \HH^\bullet(\Gamma(V,\sh{K}^\bullet))\ar[d]\\
    \HH^\bullet(u^{-1}(V'),\sh{F})\ar[r]^\sim &
    \HH^\bullet(\Gamma(V',\sh{K}^\bullet))
  }
\]
é comutativo.
Do que precede conclui-se que os isomorfismos (\hyperref[III.1.4.10.1]{1.4.10.1}) definem um isomorfismo de $\sh{O}_Y$-módulos
\[
  \RR^\bullet u_*(\sh{F})\isoto\sh{H}^\bullet(\sh{K}^\bullet),
  \tag{1.4.10.3}\label{III.1.4.10.3}
\]
\oldpage[III]{92}
e, por conseguinte, $\RR^\bullet u_*(\sh{F})$ é quase-coerente.
\end{proof}

Além disso, de (\hyperref[III.1.4.10.3]{1.4.10.3}), (\hyperref[III.1.4.10.2]{1.4.10.2}) e (\hyperref[III.1.4.10.1]{1.4.10.1}) deduz-se que:
\begin{corollary}[1.4.11]
\label{III.1.4.11}
Sob as hipóteses da Proposição~\sref{III.1.4.10}, para todo aberto afim $V$ de $Y$, o homomorfismo canônico
\[
  \HH^q(u^{-1}(V),\sh{F})\to\Gamma(V,\RR^q u_*(\sh{F}))
  \tag{1.4.11.1}\label{III.1.4.11.1}
\]
é um isomorfismo para todo $q\geq 0$.
\end{corollary}

\begin{corollary}[1.4.12]
\label{III.1.4.12}
Suponha-se que sejam satisfeitas as hipóteses da Proposição~\sref{III.1.4.10} e, além disso, que $Y$ é quase-compacto.
Então existe um inteiro $r>0$ tal que, para todo $\sh{O}_X$-módulo quase-coerente $\sh{F}$ e todo inteiro $q>r$, tem-se $\RR^q u_*(\sh{F})=0$.
Se $Y$ é afim, pode tomar-se como $r$ um inteiro tal que exista um recobrimento de $X$ formado por $r$ abertos afins.
\end{corollary}

\begin{proof}
Como se pode recobrir $Y$ por um número finito de abertos afins, em virtude do Corolário~\sref{III.1.4.11} pode reduzir-se a demonstração à segunda afirmação.
Se $\mathfrak{U}$ é um recobrimento de $X$ por $r$ abertos afins, tem-se $\HH^q(\mathfrak{U},\sh{F})=0$ para $q>r$, dado que as cocadeias de $C^q(\mathfrak{U},\sh{F})$ são alternadas; a afirmação deduz-se então da Proposição~\sref{III.1.4.1}.
\end{proof}

\begin{corollary}[1.4.13]
\label{III.1.4.13}
Suponha-se que sejam satisfeitas as hipóteses da Proposição~\sref{III.1.4.10} e, além disso, que $Y=\Spec(A)$ é afim.
Então, para todo $\sh{O}_X$-módulo quase-coerente $\sh{F}$ e todo $f\in A$, tem-se
\[
  \Gamma(Y_f,\RR^q u_*(\sh{F}))=(\Gamma(Y,\RR^q u_*(\sh{F})))_f
\]
salvo isomorfismo canônico.
\end{corollary}

\begin{proof}
Isto deduz-se do fato de que $\RR^q u_*(\sh{F})$ é um $\sh{O}_Y$-módulo quase-coerente \sref[I]{I.1.3.7}.
\end{proof}

\begin{proposition}[1.4.14]
\label{III.1.4.14}
Sejam $f:X\to Y$ um morfismo separado e quase-compacto e $g:Y\to Z$ um morfismo afim.
Para todo $\sh{O}_X$-módulo quase-coerente $\sh{F}$, o homomorfismo canônico $\RR^p(g\circ f)_*(\sh{F})\to g_*(\RR^p f_*(\sh{F}))$ (\textbf{0},~12.2.5.2) é bijetivo para todo $p$.
\end{proposition}

\begin{proof}
Para todo aberto afim $W$ de $Z$, $g^{-1}(W)$ é um aberto afim de $Y$.
O homomorfismo de pré-feixes que define o homomorfismo canônico
\[
  \RR^p(g\circ f)_*(\sh{F})\to g_*(\RR^p f_*(\sh{F}))
\]
\sref[0]{0.12.2.5} é, portanto, bijetivo em virtude do Corolário~\sref{III.1.4.11}.
\end{proof}

\begin{proposition}[1.4.15]
\label{III.1.4.15}
Sejam $u:X\to Y$ um morfismo separado de tipo finito e $v:Y'\to Y$ um morfismo plano de pré-esquemas \sref[0]{0.6.7.1}; seja $u'=u_{(Y')}$, de modo que se tem o diagrama comutativo
\[
  \xymatrix{
    X\ar[d]_u &
    X'=X_{(Y')}\ar[l]_{v'}\ar[d]^{u'}\\
    Y &
    Y'\ar[l]_{v}.
  }
  \tag{1.4.15.1}
  \label{III.1.4.15.1}
\]
Então, para todo $\sh{O}_X$-módulo quase-coerente $\sh{F}$, $\RR^q u_*'(\sh{F}')$ é canonicamente isomorfo a $\RR^q u_*(\sh{F})\otimes_{\sh{O}_Y}\sh{O}_{Y'}=v^*(\RR^q u_*(\sh{F}))$ para todo $q\geq 0$, onde $\sh{F}'={v'}^*(\sh{F})=\sh{F}\otimes_{\sh{O}_Y}\sh{O}_{Y'}$.
\end{proposition}

\oldpage[III]{93}
\begin{proof}
O homomorfismo canônico $\rho:\sh{F}\to v_*'({v'}^*(\sh{F}))$ (\textbf{0},~4.4.3.2) define por funtorialidade um homomorfismo
\[
  \RR^q u_*(\sh{F})\to\RR^q u_*(v_*'(\sh{F}')).
  \tag{1.4.15.2}\label{III.1.4.15.2}
\]

Por outro lado, pondo $w=u\circ v'=v\circ u'$, têm-se os homomorfismos canônicos (\textbf{0}~,12.2.5.1 e 12.2.5.2)
\[
  \RR^q u_*(v_*'(\sh{F}'))\to\RR^q w_*(\sh{F}')\to v_*(\RR^q u_*'(\sh{F}')).
  \tag{1.4.15.3}\label{III.1.4.15.3}
\]

Compondo (\hyperref[III.1.4.15.3]{1.4.15.3}) e (\hyperref[III.1.4.15.2]{1.4.15.2}), obtém-se um homomorfismo
\[
  \psi:\RR^q u_*(\sh{F})\to v_*(\RR^q u_*'(\sh{F}')),
\]
e obtém-se finalmente um homomorfismo canônico (cuja definição não impõe \emph{hipótese alguma} sobre $v$)
\[
  \psi^\sharp:v^*(\RR^q u_*(\sh{F}))\to\RR^q u_*'(\sh{F}'),
  \tag{1.4.15.4}
\]
e é preciso provar que é um isomorfismo quando $v$ é \emph{plano}.
É claro que a questão é local sobre $Y$ e $Y'$; pode supor-se, portanto, que $Y=\Spec(A)$ e $Y'=\Spec(B)$; utilizar-se-á além disso o seguinte lema:
\begin{lemma}[1.4.15.5]
\label{III.1.4.15.5}
Sejam $\vphi:A\to B$ um homomorfismo de anéis, $Y=\Spec(A)$, $X=\Spec(B)$, $f:X\to Y$ o morfismo correspondente a $\vphi$, e $M$ um $B$-módulo.
Para que o $\sh{O}_X$-módulo $\widetilde{M}$ seja $f$-plano \sref[0]{0.6.7.1}, é necessário e suficiente que $M$ seja um $A$-módulo plano.
Em particular, para que o morfismo $f$ seja plano, é necessário e suficiente que $B$ seja um $A$-módulo plano.
\end{lemma}

Isto deduz-se da definição \sref[0]{0.6.7.1} e de \sref[0]{0.6.3.3}, levando em conta \sref[I]{I.1.3.4}.

Assim, de (\hyperref[III.1.4.11.1]{1.4.11.1}) e das definições dos homomorfismos (\hyperref[III.1.4.15.3]{1.4.15.3}) (cf.~\sref[0]{0.12.2.5}) deduz-se que $\psi$ corresponde então ao morfismo composto
\[
  \HH^q(X,\sh{F})\xrightarrow{\rho_q}\HH^q(X,v_*'({v'}^*(\sh{F})))\xrightarrow{\theta_q}\HH^q(X',{v'}^*(v_*'({v'}^*(\sh{F}))))\xrightarrow{\sigma_q}\HH^q(X',{v'}^*(\sh{F})),
\]
onde $\rho_q$ e $\sigma_q$ são os homomorfismos de cohomologia correspondentes aos morfismos canônicos $\rho$ e $\sigma:{v'}^*(v_*'(\sh{G}'))\to\sh{G}'$, e $\theta_q$ é o $\vphi$-morfismo (\textbf{0},~12.1.3.1) relativo ao $\sh{O}_X$-módulo $v_*'({v'}^*(\sh{F}))$.
Mas, pela funtorialidade de $\theta_q$, tem-se o diagrama comutativo
\[
  \xymatrix{
    \HH^q(X,\sh{F})\ar[rr]^{\rho_q}\ar[dd]_{\theta_q} & &
    \HH^q(X,v_*'({v'}^*(\sh{F})))\ar[dd]^{\theta_q}\\\\
    \HH^q(X',{v'}^*(\sh{F}))\ar[rr]^{{v'}^*(\rho_q)} & &
    \HH^q(X',{v'}^*(v_*'({v'}^*(\sh{F})))),
  }
\]
\oldpage[III]{94}
e, como por definição \sref[0]{0.4.4.3} ${v'}^*(\rho)$ é o inverso de $\sigma$, vê-se que o morfismo composto considerado anteriormente não é, afinal, senão $\theta_q$; por conseguinte, $\psi^\sharp$ é o $B$-homomorfismo associado $\HH^q(X,\sh{F})\otimes_A B\to\HH^q(X',\sh{F}')$.
Como $u$ é de tipo finito, $X$ é uma união finita de abertos afins $U_i$ ($1\leq i\leq r$); seja $\mathfrak{U}$ o recobrimento $(U_i)$.
Como $v$ é um morfismo afim, também o é $v'$ \sref[II]{II.1.6.2}[iii]; por conseguinte, os $U_i'={v'}^{-1}(U_i)$ formam um recobrimento aberto afim $\mathfrak{U}'$ de $X'$.
Sabe-se então \sref[0]{0.12.1.4.2} que o diagrama
\[
  \xymatrix{
    \HH^q(\mathfrak{U},\sh{F})\ar[r]^{\theta_q}\ar[d] &
    \HH^q(\mathfrak{U}',\sh{F}')\ar[d]\\
    \HH^q(X,\sh{F})\ar[r]^{\theta_q} &
    \HH^q(X',\sh{F}')
  }
\]
é comutativo, e as setas verticais são isomorfismos dado que $X$ e $X'$ são esquemas \sref{III.1.4.1}.
Por conseguinte, basta provar que o $\vphi$-morfismo canônico $\theta_q:\HH^q(\mathfrak{U},\sh{F})\to\HH^q(\mathfrak{U}',\sh{F}')$ é tal que o $B$-homomorfismo associado
\[
  \HH^q(\mathfrak{U},\sh{F})\otimes_A B\to\HH^q(\mathfrak{U}',\sh{F}')
\]
seja um isomorfismo.
Para toda sucessão $\mathbf{s}=(i_k)_{0\leq k\leq p}$ de $p+1$ índices de $[1,r]$, coloquem-se $U_\mathbf{s}=\bigcap_{k=0}^p U_{i_k}$, $U_\mathbf{s}'=\bigcap_{k=0}^p U_{i_k}'={v'}^{-1}(U_\mathbf{s})$, $M_\mathbf{s}=\Gamma(U_\mathbf{s},\sh{F})$ e $M_\mathbf{s}'=\Gamma(U_\mathbf{s}',\sh{F}')$.
A aplicação canônica $M_\mathbf{s}\otimes_A B\to M_\mathbf{s}'$ é um isomorfismo \sref[I]{I.1.6.5}; portanto, a aplicação canônica $C^p(\mathfrak{U},\sh{F})\otimes_A B\to C^p(\mathfrak{U}',\sh{F}')$ é um isomorfismo, mediante o qual $d\otimes 1$ identifica-se com a aplicação cobordo $C^p(\mathfrak{U}',\sh{F}')\to C^{p+1}(\mathfrak{U}',\sh{F}')$.
Como $B$ é um $A$-módulo \emph{plano}, da definição dos módulos de cohomologia deduz-se que a aplicação canônica $\HH^q(\mathfrak{U},\sh{F})\otimes_A B\to\HH^q(\mathfrak{U}',\sh{F}')$ é um isomorfismo \sref[0]{0.6.1.1}.
Este resultado se generalizará mais adiante em \EGAUnavailableHyperref{section:III.6}{\textsection6}.
\end{proof}

\begin{corollary}[1.4.16]
\label{III.1.4.16}
Sejam $A$ um anel, $X$ um $A$-esquema de tipo finito e $B$ uma $A$-álgebra fielmente plana sobre $A$.
Para que $X$ seja afim, é necessário e suficiente que $X\otimes_A B$ o seja.
\end{corollary}

\begin{proof}
A condição é evidentemente necessária \sref[I]{I.3.2.2}; provaremos que é suficiente.
Como $X$ é separado sobre $A$ e o morfismo $\Spec(B)\to\Spec(A)$ é plano, da Proposição~\sref{III.1.4.15} deduz-se que
\[
\label{III.1.4.16.1}
  \HH^i(X\otimes_A B,\sh{F}\otimes_A B)=\HH^i(X,\sh{F})\otimes_A B
  \tag{1.4.16.1}
\]
para todo $i\geq 0$ e todo $\sh{O}_X$-módulo quase-coerente $\sh{F}$.
Se $X\otimes_A B$ é afim, o membro esquerdo de (\hyperref[III.1.4.16.1]{1.4.16.1}) é nulo para $i=1$; portanto, também o é $\HH^1(X,\sh{F})$, dado que $B$ é um $A$-módulo fielmente plano.
Como $X$ é um esquema quase-compacto, a demonstração termina pelo critério de Serre \sref[II]{II.5.2.1}.
\end{proof}

\begin{proposition}[1.4.17]
\label{III.1.4.17}
Seja $X$ um pré-esquema e seja $0\to\sh{F}\xrightarrow{u}\sh{G}\xrightarrow{v}\sh{H}\to 0$ uma sucessão exata de $\sh{O}_X$-módulos.
Se $\sh{F}$ e $\sh{H}$ são quase-coerentes, também o é $\sh{G}$.
\end{proposition}

\oldpage[III]{95}
\begin{proof}
A questão é local sobre $X$, de modo que pode supor-se que $X=\Spec(A)$ é afim; basta então provar que $\sh{G}$ satisfaz as condições (d1) e (d2) de \sref[I]{I.1.4.1} (com $V=X$).
A verificação de (d2) é imediata, pois, se $t\in\Gamma(X,\sh{G})$ se anula ao restringi-lo a $D(f)$, o mesmo ocorre com sua imagem $v(t)\in\Gamma(X,\sh{H})$; portanto, existe um $m>0$ tal que $f^m v(t)=v(f^m t)=0$ \sref[I]{I.1.4.1}, e, como $\Gamma$ é exato à esquerda, $f^m t=u(s)$, onde $s\in\Gamma(X,\sh{F})$; como $u$ é injetivo, a restrição de $s$ a $D(f)$ é nula, logo por \sref[I]{I.1.4.1} existe um inteiro $n>0$ tal que $f^n s=0$; deduz-se finalmente que $f^{m+n}t=u(f^n s)=0$.

Verifique-se agora (d1); seja $t'\in\Gamma(D(f),\sh{G})$; como $\sh{H}$ é quase-coerente, existe um inteiro $m$ tal que $f^m v(t')=v(f^m t')$ se prolongue a uma seção $z\in\Gamma(X,\sh{H})$ \sref[I]{I.1.4.1}.
Mas, em virtude do Teorema~\sref{III.1.3.1} (o de (\textbf{I},~5.1.9.2)) aplicado ao $\sh{O}_X$-módulo quase-coerente $\sh{F}$, a sucessão $\Gamma(X,\sh{G})\to\Gamma(X,\sh{H})\to 0$ é exata; existe, pois, $t\in\Gamma(X,\sh{G})$ tal que $z=v(t)$; vê-se assim que $v(f^m t'-t'')=0$, denotando por $t''$ a restrição de $t$ a $D(f)$; portanto, $f^m t'-t''=u(s')$, onde $s'\in\Gamma(D(f),\sh{F})$.
Mas, como $\sh{F}$ é quase-coerente, existe um inteiro $n>0$ tal que $f^n s'$ se prolongue a uma seção $s\in\Gamma(X,\sh{F})$; como $f^{m+n}t'-f^n t''=u(f^n s')$, vê-se que $f^{m+n}t'$ é a restrição a $D(f)$ de uma seção $f^n t+u(s)\in\Gamma(X,\sh{G})$, o que termina a demonstração.
\end{proof}

\section{Estudo cohomológico dos morfismos projetivos}
\label{section:III.2}


\subsection{Cálculos explícitos de certos grupos de cohomologia}
\label{subsection:III.2.1}

\begin{env}[2.1.1]
\label{III.2.1.1}
\oldpage[III]{95}
Seja $X$ um pré-esquema e seja $\sh{L}$ um $\sh{O}_X$-módulo invertível; consideremos o anel graduado \sref[0]{0.5.4.6}
\[
\label{III.2.1.1.1}
  S = \Gamma_*(X,\sh{L}) = \bigoplus_{n\in\bb{Z}} \Gamma(X,\sh{L}^{\otimes n}).
  \tag{2.1.1.1}
\]

Seja $(f_i)_{1\leq i\leq r}$ uma família finita de elementos \emph{homogêneos} de $S$, digamos $f_i\in S_{d_i}$; ponhamos $U_i = X_{f_i}$, $U = \bigcup_i U_i$, e denotemos por $\mathfrak{U}$ o recobrimento $(U_i)$ de $U$.
Para todo $\sh{O}_X$-módulo quase-coerente $\sh{F}$, ponhamos
\[
\label{III.2.1.1.2}
  \HH^\bullet(U,\sh{F}(*)) = \bigoplus_{n\in\bb{Z}} \HH^\bullet(U,\sh{F}\otimes\sh{L}^{\otimes n})
  \tag{2.1.1.2}
\]
\[
\label{III.2.1.1.3}
  \HH^\bullet(\mathfrak{U},\sh{F}(*)) = \bigoplus_{n\in\bb{Z}} \HH^\bullet(\mathfrak{U},\sh{F}\otimes\sh{L}^{\otimes n}).
  \tag{2.1.1.3}
\]

Os grupos abelianos (\hyperref[III.2.1.1.2]{2.1.1.2}) e (\hyperref[III.2.1.1.3]{2.1.1.3}) são \emph{bigraduados}, tomando
\[
  (\HH^\bullet(U,\sh{F}(*)))_{mn} = \HH^m(U,\sh{F}\otimes\sh{L}^{\otimes n})
\]
e uma definição análoga para (\hyperref[III.2.1.1.3]{2.1.1.3}).
Para a segunda graduação, é claro que esses grupos são $S$-módulos graduados, como se deduz, por exemplo, do fato de que $\sh{F}\mapsto\HH^m(U,\sh{F})$ e $\sh{F}\mapsto\HH^m(\mathfrak{U},\sh{F})$ são funtores.
\end{env}

\begin{env}[2.1.2]
\label{III.2.1.2}
Consideremos agora o $S$-módulo graduado \sref[0]{0.5.4.6}
\[
\label{III.2.1.2.1}
  M = \Gamma_*(\sh{L},\sh{F}) = \HH^0(X,\sh{F}(*)) = \bigoplus_{n\in\bb{Z}} \Gamma(X,\sh{F}\otimes\sh{L}^{\otimes n}).
  \tag{2.1.2.1}
\]

\oldpage[III]{96}
Se $X$ é um pré-esquema cujo espaço subjacente é noetheriano, ou um esquema quase compacto, deduz-se de \sref[I]{I.9.3.1} que, colocando como de costume $U_{i_0 i_1\cdots i_p} = \bigcap_{k=0}^p U_{i_k}$, tem-se, salvo isomorfismo canônico,
\[
  \Gamma(U_{i_0 i_1\cdots i_p},\sh{F}(*)) = \HH^0(U_{i_0 i_1\cdots i_p},\sh{F}(*)) = M_{f_{i_0}f_{i_1}\cdots f_{i_p}}.
\]

Além disso, com a notação de (\hyperref[III.1.2.2]{1.2.2}), podemos identificar $M_{f_{i_0}f_{i_1}\cdots f_{i_p}}$ com $\varinjlim_n M_{i_0 i_1\cdots i_p}^{(n)}$.
Esta identificação é um isomorfismo de $S$-módulos \emph{graduados} se definirmos o grau de um elemento homogêneo $z\in\varinjlim_n M_{i_0 i_1\cdots i_p}^{(n)}$ da seguinte forma: $z$ é a imagem canônica de um elemento homogêneo $x\in M_{i_0 i_1\cdots i_p}^{(n)} = M$ de grau $m$, e tomamos como grau de $z$ o número $m-n(d_{i_0}+d_{i_1}+\cdots+d_{i_p})$.
Considerando a definição dos homomorfismos $\vphi_{kh}:M_{i_0 i_1\cdots i_p}^{(h)}\to M_{i_0 i_1\cdots i_p}^{(k)}$ (\hyperref[III.1.2.2]{1.2.2}), vê-se imediatamente que esta definição não depende do ``representante'' $x$ de $z$ escolhido.
Denotando, como em (\hyperref[III.1.2.2]{1.2.2}), por $C_n^p(M)$ o conjunto de aplicações alternadas de $[1,r]^{p+1}$ em $M$ (para todo $n$), definimos de forma análoga uma estrutura de $S$-módulo graduado sobre $\varinjlim_n C_n^p(M)$; como em (\hyperref[III.1.2.2]{1.2.2}), temos então
\[
\label{III.2.1.2.2}
  C^p(\mathfrak{U},\sh{F}(*)) = \varinjlim_n C_n^p(M)
  \tag{2.1.2.2}
\]
onde o isomorfismo entre ambos os membros \emph{respeita os graus}.
Tem-se também, como em (\hyperref[III.1.2.2]{1.2.2}),
\[
\label{III.2.1.2.3}
  C^p(\mathfrak{U},\sh{F}(*)) = C^{p+1}((\mathbf{f}),M) = \varinjlim_n K^{p+1}(\mathbf{f}^n,M)
  \tag{2.1.2.3}
\]
onde o isomorfismo \emph{conserva os graus}: o grau de um elemento de $\varinjlim_n K^{p+1}(\mathbf{f}^n,M)$ que seja a imagem canônica de uma cocadeia $g\in K^{p+1}(\mathbf{f}^n,M)$ cujos valores $g(i_0,\dots,i_p)$ pertençam a uma mesma componente homogênea $M_m$ de $M$ é $m-n(d_{i_0}+\cdots+d_{i_p})$, independentemente da escolha dessa cocadeia como representante do elemento em causa.

Como os isomorfismos anteriores são compatíveis com os operadores de coborde, conclui-se, como em (\hyperref[III.1.2.2]{1.2.2}), que:
\end{env}

\begin{proposition}[2.1.3]
\label{III.2.1.3}
Seja $X$ um pré-esquema cujo espaço subjacente é noetheriano, ou um esquema quase compacto.
Existe um isomorfismo canônico, funtorial em $\sh{F}$,
\[
\label{III.2.1.3.1}
  \HH^p(\mathfrak{U},\sh{F}(*)) \simto \HH^{p+1}((\mathbf{f}),M)
  \qquad\text{para todo } p\geq 1.
  \tag{2.1.3.1}
\]
Além disso, existe uma sequência exata, funtorial em $\sh{F}$,
\[
\label{III.2.1.3.2}
  0 \to \HH^0((\mathbf{f}),M) \to M \to \HH^0(\mathfrak{U},\sh{F}(*)) \to \HH^1((\mathbf{f}),M) \to 0.
  \tag{2.1.3.2}
\]
Todos os homomorfismos introduzidos são de grau $0$ para as estruturas de $S$-módulo graduado, onde $S$ é o anel (\hyperref[III.2.1.1.1]{2.1.1.1}).
\end{proposition}

\begin{corollary}[2.1.4]
\label{III.2.1.4}
\oldpage[III]{97}
Se $X$ é um esquema quase-compacto e os $U_i=X_{f_i}$ são afins, existe um isomorfismo canônico, funtorial em $\sh{F}$ e de grau $0$,
\[
\label{III.2.1.4.1}
  \HH^p(U,\sh{F}(*)) \simto \HH^{p+1}((\mathbf{f}),M)
  \qquad\text{para todo } p\geq 1,
  \tag{2.1.4.1}
\]
e uma sequência exata, funtorial em $\sh{F}$,
\[
\label{III.2.1.4.2}
  0 \to \HH^0((\mathbf{f}),M) \to M \to \HH^0(U,\sh{F}(*)) \to \HH^1((\mathbf{f}),M) \to 0,
  \tag{2.1.4.2}
\]
na qual todos os homomorfismos são de grau $0$.
\end{corollary}

\begin{proof}
Basta aplicar \sref{III.1.4.1} ao resultado de \sref{III.2.1.3}.
\end{proof}

A proposição ``local'' análoga a \sref{III.2.1.3} é a seguinte.

\begin{proposition}[2.1.5]
\label{III.2.1.5}
Seja $S$ um anel graduado positivamente, sejam os $f_i$ ($1\leq i\leq r$) elementos homogêneos de $S_+$ de graus $d_i$, e seja $M$ um $S$-módulo graduado.
Seja $X=\Proj(S)$ o espectro primo homogêneo de $S$, e ponhamos $U_i=\mathrm{D}_+(f_i)$, $U=\bigcup_i U_i$ e $\HH^\bullet(U,\widetilde{M}(*))=\bigoplus_{n\in\bb{Z}}\HH^\bullet(U,(M(n))^\sim)$.
Existem isomorfismos canônicos, funtoriais em $M$ e de grau $0$ para estruturas de $S$-módulo graduado,
\[
\label{III.2.1.5.1}
  \HH^p(U,\widetilde{M}(*)) \simto \HH^{p+1}((\mathbf{f}),M)
  \qquad\text{para todo } p\geq 1,
  \tag{2.1.5.1}
\]
e uma sequência exata, funtorial em $M$,
\[
\label{III.2.1.5.2}
  0 \to \HH^0((\mathbf{f}),M) \to M \to \HH^0(U,\widetilde{M}(*)) \to \HH^1((\mathbf{f}),M) \to 0,
  \tag{2.1.5.2}
\]
na qual todos os homomorfismos são de grau $0$.
\end{proposition}

\begin{proof}
De fato, por definição \sref[II]{II.2.5.2},
\[
  \Gamma(U_{i_0i_1\cdots i_p},(M(n))^\sim)
  =(M_{f_{i_0}\cdots f_{i_p}})_n,
\]
e, portanto, $\Gamma(U_{i_0i_1\cdots i_p},\widetilde{M}(*))=M_{f_{i_0}\cdots f_{i_p}}$.
O resto do raciocínio é o mesmo que na demonstração de \sref{III.2.1.4}, levando em conta que $X$ é um esquema.
\end{proof}

\begin{env}[2.1.6]
\label{III.2.1.6}
\emph{Observações.} ---
(i) Sob as hipóteses de \sref{III.2.1.5}, os funtores $\Gamma(U_{i_0i_1\cdots i_p},\widetilde{M}(*))$ são exatos em $M$, por \sref[0]{0.1.3.2}.
O mesmo raciocínio que em (\hyperref[III.1.2.5]{1.2.5}) mostra, portanto, que, se $0\to M'\to M\to M''\to0$ é uma sequência exata de $S$-módulos graduados cujos homomorfismos têm grau $0$, então para todo $p\geq0$ existem diagramas comutativos
\[
\label{III.2.1.6.1}
  \xymatrix{
    \HH^p(U,\widetilde{M''}(*)) \ar[r]^-\partial \ar[d] &
    \HH^{p+1}(U,\widetilde{M'}(*)) \ar[d] \\
    \HH^{p+1}((\mathbf{f}),M'') \ar[r]^-\partial &
    \HH^{p+2}((\mathbf{f}),M')
  }
  \tag{2.1.6.1}
\]

(ii) A Proposição~\sref{III.2.1.5} é especialmente útil quando $S$ é uma $A$-álgebra gerada por um número finito de elementos de grau $1$, com $A$ noetheriano; nesse caso, todo $\sh{O}_X$-módulo quase-coerente é da forma $\widetilde{M}$ \sref[II]{II.2.7.7}.
\end{env}

\begin{env}[2.1.7]
\label{III.2.1.7}
\oldpage[III]{98}
Aplicamos agora \sref{III.2.1.5} ao caso $S=A[\mathrm{T}_0,\dots,\mathrm{T}_r]$, onde $A$ é um anel arbitrário e os $\mathrm{T}_i$ são indeterminados, tomando $M=S$ e $f_i=\mathrm{T}_i$.
Ficamos assim essencialmente reduzidos a calcular $\HH^\bullet((\mathbf{T}),S)$, onde $\mathbf{T}=(\mathrm{T}_i)_{0\leq i\leq r}$.
\end{env}

\begin{lemma}[2.1.8]
\label{III.2.1.8}
Se $S=A[\mathrm{T}_0,\dots,\mathrm{T}_r]$, então, com $\mathbf{T}=(\mathrm{T}_i)_{0\leq i\leq r}$,
\[
\label{III.2.1.8.1}
  \HH^i(\mathbf{T}^n,S)=0
  \qquad\text{se }i\neq r+1,
  \tag{2.1.8.1}
\]
\[
\label{III.2.1.8.2}
  \HH^{r+1}(\mathbf{T}^n,S)=S/(\mathbf{T}^n).
  \tag{2.1.8.2}
\]
Assim, o $A$-módulo $\HH^{r+1}(\mathbf{T}^n,S)$ tem uma $A$-base formada pelas classes módulo $(\mathbf{T}^n)$ dos monômios
\[
  \mathrm{T}_{\mathbf{p}}=\mathrm{T}_0^{p_0}\cdots\mathrm{T}_r^{p_r},
  \qquad \mathbf{p}=(p_0,\dots,p_r),\quad 0\leq p_i<n\ \text{Para tudo }i.
\]
\end{lemma}

\begin{proof}
Isto deduz-se imediatamente de (\hyperref[III.1.1.3.5]{1.1.3.5}) e da Proposição~\sref{III.1.1.4}, cujas hipóteses são satisfeitas trivialmente.
\end{proof}

\begin{env}[2.1.9]
\label{III.2.1.9}
Passando ao limite indutivo em relação a $n$, as relações (\hyperref[III.2.1.8.1]{2.1.8.1}) dão $\HH^i((\mathbf{T}),S)=0$ para $i\neq r+1$.
Para $i=r+1$, o sistema indutivo consiste nos módulos $S/(\mathbf{T}^n)$, sendo o homomorfismo
\[
  \vphi_{nm}:S/(\mathbf{T}^n)\longrightarrow S/(\mathbf{T}^m)
  \qquad (0\leq n\leq m)
\]
a multiplicação por $(\mathrm{T}_0\cdots\mathrm{T}_r)^{m-n}$.
Para $n\geq\sup(p_i)_{0\leq i\leq r}$, denotemos por $\xi_{\mathbf{p}}^{(n)}=\xi_{p_0,\dots,p_r}^{(n)}$ a classe de $\mathrm{T}_0^{n-p_0}\cdots\mathrm{T}_r^{n-p_r}$ módulo $(\mathbf{T}^n)$.
Então $\vphi_{nm}(\xi_{\mathbf{p}}^{(n)})=\xi_{\mathbf{p}}^{(m)}$, de modo que estes elementos têm a mesma imagem canônica $\xi_{\mathbf{p}}=\xi_{p_0,\dots,p_r}$ no limite indutivo $\HH^{r+1}((\mathbf{T}),S)$.
Pela definição do grau em (\hyperref[III.2.1.2]{2.1.2}), o grau de $\xi_{\mathbf{p}}$ é $-|\mathbf{p}|=-(p_0+p_1+\cdots+p_r)$.
É claro que os $\xi_{\mathbf{p}}^{(n)}$ com $0<p_i\leq n$ para $0\leq i\leq r$ formam uma base de $S/(\mathbf{T}^n)$.
Deduz-se, então, imediatamente de \sref{III.2.1.8}:
\end{env}

\begin{corollary}[2.1.10]
\label{III.2.1.10}
Com a notação de \sref{III.2.1.8},
\[
\label{III.2.1.10.1}
  \HH^i((\mathbf{T}),S)=0
  \qquad\text{para }i\neq r+1,
  \tag{2.1.10.1}
\]
e $\HH^{r+1}((\mathbf{T}),S)$ é um $A$-módulo livre com base nos elementos $\xi_{p_0,\dots,p_r}$ tais que $p_i>0$ para $0\leq i\leq r$.
\end{corollary}

\begin{env}[2.1.11]
\label{III.2.1.11}
\emph{Observação.} ---
Seja $N$ um $A$-módulo qualquer e seja $M=S\otimes_A N$.
O raciocínio de \sref{III.2.1.8} mostra mais geralmente que
\[
\label{III.2.1.11.1}
  \HH^i(\mathbf{T}^n,M)=0
  \qquad\text{se }i\neq r+1,
  \tag{2.1.11.1}
\]
\[
\label{III.2.1.11.2}
  \HH^{r+1}(\mathbf{T}^n,M)=(S/(\mathbf{T}^n))\otimes_A N.
  \tag{2.1.11.2}
\]
Com efeito, esta última fórmula deduz-se diretamente de (\hyperref[III.1.1.3.5]{1.1.3.5}).
Além disso, é claro que
\[
  M/(\mathrm{T}_0^nM+\cdots+\mathrm{T}_{i-1}^nM)
  =\bigl(S/(\mathrm{T}_0^nS+\cdots+\mathrm{T}_{i-1}^nS)\bigr)\otimes_A N,
\]
pois o ideal $\mathrm{T}_0^nS+\cdots+\mathrm{T}_{i-1}^nS$ é um somando direto do $A$-módulo $S$.
Isso permite aplicar \sref{III.1.1.4} a $M$ e dá (\hyperref[III.2.1.11.1]{2.1.11.1}).
\end{env}

Combinando \sref{III.2.1.10} e \sref{III.2.1.5}, obtém-se:

\begin{proposition}[2.1.12]
\label{III.2.1.12}
Seja $A$ um anel, seja $r$ um inteiro $>0$ e seja $X=\bb{P}_A^r$ \sref[II]{II.4.1.1}.
Então:
\begin{enumerate}
  \item[{\rm(i) }] $\HH^i(X,\sh{O}_X(*))=0$ para $i\neq0,r$.
  \item[{\rm(ii) }] O homomorfismo canônico $\alpha:S\to\HH^0(X,\sh{O}_X(*))$ \sref[II]{II.2.6.2} é bijectivo.
\oldpage[III]{99}
  \item[{\rm(iii) }] $\HH^r(X,\sh{O}_X(*))$ é um $A$-módulo livre com base nos elementos $\xi_{p_0,\dots,p_r}$, onde $p_i>0$ para $0\leq i\leq r$; o elemento $\xi_{p_0,\dots,p_r}$ tem grau $-|\mathbf{p}|=-(p_0+\cdots+p_r)$, e
  \[
    \mathrm{T}_i\xi_{p_0,\dots,p_r}
    =\xi_{p_0,\dots,p_{i-1},p_i-1,p_{i+1},\dots,p_r}.
  \]
\end{enumerate}
\end{proposition}

\begin{proof}
Na sequência exata (\hyperref[III.2.1.5.2]{2.1.5.2}) aplicada a $M=S=A[\mathrm{T}_0,\dots,\mathrm{T}_r]$, temos $\HH^0((\mathbf{T}),S)=0$ e $\HH^1((\mathbf{T}),S)=0$ por \sref{III.2.1.10.1}; além disso, a Proposição~\sref{III.2.1.5} é aplicada com $U=X$, já que $X$ é a união dos $\mathrm{D}_+(\mathrm{T}_i)$ \sref[II]{II.2.3.14}.
Resta identificar a aplicação $S\to\HH^0(X,\sh{O}_X(*))$ de (\hyperref[III.2.1.5.2]{2.1.5.2}) com a aplicação canônica $\alpha$.
Isso se deduz da identificação canônica de $\HH^0(U,\sh{O}_X(*))$ e $\HH^0(\mathfrak{U},\sh{O}_X(*))$.
\end{proof}

\begin{corollary}[2.1.13]
\label{III.2.1.13}
Os únicos pares $(i,n)$ para os quais pode ocorrer que $\HH^i(X,\sh{O}_X(n))\neq0$ são $i=0$, $n\geq0$, e $i=r$, $n\leq-(r+1)$.
\end{corollary}

Se $A\neq0$, então temos efetivamente $\HH^i(X,\sh{O}_X(n))\neq0$ para todos os pares listados em \sref{III.2.1.13}.
Isto deduz-se de \sref{III.2.1.12}, já que $S_n\neq0$ para todo grau $n\geq0$.

Nas aplicações deste capítulo, utilizaremos principalmente o seguinte resultado menos preciso.

\begin{corollary}[2.1.14]
\label{III.2.1.14}
Os $A$-módulos $\HH^i(X,\sh{O}_X(n))$ são livres de tipo finito; se $i>0$, anulam-se para $n>0$.
\end{corollary}

\begin{proposition}[2.1.15]
\label{III.2.1.15}
Seja $Y$ um pré-esquema, seja $\sh{E}$ um $\sh{O}_Y$-módulo localmente livre de posto $r+1$, seja $X=\bb{P}(\sh{E})$ o fibrado projetivo definido por $\sh{E}$, e seja $f:X\to Y$ o morfismo estrutural.
Os únicos valores de $i$ e $n$ para os quais pode ocorrer que $\RR^if_*(\sh{O}_X(n))\neq0$ são $i=0$, $n\geq0$, e $i=r$, $n\leq-(r+1)$.
Além disso, o homomorfismo canônico \sref[II]{II.3.3.2}
\[
  \alpha:\bb{S}_{\sh{O}_Y}(\sh{E})
  \longrightarrow \Gamma_*(\sh{O}_X)
  =\RR^0f_*(\sh{O}_X(*))
  =\bigoplus_{n\in\bb{Z}}f_*(\sh{O}_X(n))
\]
é um isomorfismo.
\end{proposition}

\begin{proof}
A questão é local sobre $Y$, portanto podemos supor que $Y$ é afim, de anel $A$, e que $\sh{E}=\widetilde{E}$, onde $E=A^{r+1}$.
Então, ficamos imediatamente reduzidos a \sref{III.2.1.12}, tendo em conta \sref{III.1.4.11}.
\end{proof}

\begin{env}[2.1.16]
\label{III.2.1.16}
\emph{Observação.} ---
Mais adiante, completaremos os resultados de \sref{III.2.1.15} demonstrando as seguintes proposições.
Ponhamos
\[
  \boldsymbol{\omega}=f^*\bigl(\bigwedge^{r+1}\sh{E}\bigr)(-r-1),
\]
que é um $\sh{O}_X$-módulo invertível.
Então:
\begin{enumerate}
  \item[{\rm(i) }] Existe um isomorfismo canônico
  \[
  \label{III.2.1.16.1}
    \rho:\sh{O}_Y \simto \RR^rf_*(\boldsymbol{\omega}).
    \tag{2.1.16.1}
  \]
  \item[{\rm(ii) }] A acoplamento dado pelo produto cup \sref[0]{0.12.2.2}
  \[
  \label{III.2.1.16.2}
    \RR^rf_*(\sh{O}_X(n))
    \times \RR^0f_*(\boldsymbol{\omega}(-n))
    \longrightarrow \RR^rf_*(\boldsymbol{\omega}),
    \tag{2.1.16.2}
  \]
  composto por $\rho^{-1}$, define um \emph{isomorfismo} de $\RR^rf_*(\sh{O}_X(n))$ sobre o \emph{dual} do $\sh{O}_Y$-módulo localmente livre
  \[
    \RR^0f_*(\boldsymbol{\omega}(-n))
    =\bigl(\bigwedge^{r+1}\sh{E}\bigr)
      \otimes_{\sh{O}_Y}
      \bigl(\bb{S}_{\sh{O}_Y}(\sh{E})\bigr)_{-n}.
  \]
\end{enumerate}
\end{env}

\subsection{O teorema fundamental dos morfismos projetivos}
\label{subsection:III.2.2}
\oldpage[III]{100}

\begin{theorem}[2.2.1]
\label{III.2.2.1}
(Serre).
Seja $Y$ um pré-esquema noetheriano, seja $f:X\to Y$ um morfismo próprio e seja $\sh{L}$ um $\sh{O}_X$-módulo invertível amplo em relação a $f$.
Para todo $\sh{O}_X$-módulo $\sh{F}$, ponhamos
\[
  \sh{F}(n)=\sh{F}\otimes_{\sh{O}_X}\sh{L}^{\otimes n}
  \qquad(n\in\bb{Z}).
\]
Então, para todo $\sh{O}_X$-módulo coerente $\sh{F}$:
\begin{enumerate}
  \item[{\rm(i) }] os $\sh{O}_Y$-módulos $\RR^qf_*(\sh{F})$ são coerentes;
  \item[{\rm(ii) }] existe um inteiro $N$ tal que, para $n\geq N$,
  \[
    \RR^qf_*(\sh{F}(n))=0
    \qquad\text{Para tudo }q>0;
  \]
  \item[{\rm(iii) }] existe um inteiro $N$ tal que, para $n\geq N$, o homomorfismo canônico
  \[
    f^*\bigl(f_*(\sh{F}(n))\bigr)\longrightarrow\sh{F}(n)
  \]
  é sobrejetivo.
\end{enumerate}
\end{theorem}

\begin{proof}
Observemos primeiro que, se o teorema for verdadeiro depois de substituir $\sh{L}$ por $\sh{L}^{\otimes d}$, onde $d>0$, então será verdadeiro na sua forma original.
Com efeito, escrevendo
\[
  \sh{F}(n)
  =\bigl(\sh{F}\otimes\sh{L}^{\otimes r}\bigr)
   \otimes\sh{L}^{\otimes hd},
   \qquad h>0,\quad 0\leq r<d,
\]
A hipótese dá, para cada $r$, um inteiro $N_r$ tal que (ii) e (iii) valem para $\sh{F}\otimes\sh{L}^{\otimes r}$ sempre que $h\geq N_r$.
Tomando para $N$ o maior dos inteiros $dN_r$, as afirmações (ii) e (iii) valem, portanto, para $n\geq N$.

Podemos, portanto, supor que $\sh{L}$ é muito amplo relativamente a $f$ \sref[II]{II.4.6.11}.
Existe então uma imersão aberta dominante sobre $Y$
\[
  i:X\longrightarrow P,
  \qquad P=\Proj(\sh{S}),
\]
onde $\sh{S}$ é uma $\sh{O}_Y$-álgebra quase-coerente graduada em graus positivos, $\sh{S}_1$ é de tipo finito e gera $\sh{S}$, e $\sh{L}$ é isomorfo a $i^*(\sh{O}_P(1))$ \sref[II]{II.4.4.7}.
Como $f$ é próprio, $i$ também é próprio \sref[II]{II.5.4.4} e, portanto, é um isomorfismo $X\simto P$.
Assim, ficamos reduzidos ao caso $X=\Proj(\sh{S})$ e $\sh{L}=\sh{O}_X(1)$.
O teorema é então consequência da seguinte proposição.
\end{proof}

\begin{proposition}[2.2.2]
\label{III.2.2.2}
Seja $A$ um anel noetheriano e seja $S$ uma $A$-álgebra graduada em graus positivos tal que o $A$-módulo $S_1$ possui um sistema de $r+1$ geradores e gera a álgebra $S$.
Ponhamos $X=\Proj(S)$.
Para todo $\sh{O}_X$-módulo coerente $\sh{F}$:
\begin{enumerate}
  \item[{\rm(i) }] os $A$-módulos $\HH^q(X,\sh{F})$ são do tipo finito;
  \item[{\rm(ii) }] $\HH^q(X,\sh{F})=0$ para $q>r$;
  \item[{\rm(iii) }] existe um inteiro $N$ tal que, para $n\geq N$,
  \[
    \HH^q(X,\sh{F}(n))=0
    \qquad\text{para todo }q>0;
  \]
  \item[{\rm(iv) }] existe um inteiro $N$ tal que, para $n\geq N$, $\sh{F}(n)$ é gerado por suas seções sobre $X$.
\end{enumerate}
\end{proposition}

\begin{proof}
Primeiro vamos mostrar como \sref{III.2.2.2} implica \sref{III.2.2.1}.
No caso particular $X=\Proj(\sh{S})$ reduzido anteriormente \sref{III.2.2.1}, $Y$ é quase compacto.
Pode, portanto, ser recoberto por um número finito de abertos afins $U_\alpha$, de anéis noetherianos, tais que a restrição de $\sh{S}_1$ a cada $U_\alpha$ seja gerada por um número finito de seções sobre $U_\alpha$.
Uma vez conhecida \sref{III.2.2.2}, basta tomar, nas partes (ii) e (iii) de \sref{III.2.2.1}, o maior dos inteiros correspondentes aos $U_\alpha$, tendo em conta \sref{III.1.4.11} e \sref[II]{II.3.4.7}.

Para provar \sref{III.2.2.2}, observemos que $X$ se identifica com um subesquema fechado de $P=\bb{P}_A^r$ \sref[II]{II.3.6.2}.
Além disso, se $j:X\to P$ é a injeção canônica, então $j_*(\sh{F})$ é um $\sh{O}_P$-módulo coerente e
\[
  j_*(\sh{F}(n))=(j_*\sh{F})(n)
\]
por \sref[II]{II.3.4.5} e \sref[II]{II.3.5.2}.
Por força de (G, II, Corolário do Teorema~4.9.1), ficamos, portanto, reduzidos a provar \sref{III.2.2.2} no caso particular
\[
  X=\bb{P}_A^r,
  \qquad S=A[T_0,\ldots,T_r].
\]

\oldpage[III]{101}

Como $X$ é coberto pelos $r+1$ abertos afins $D_+(T_i)$, a afirmação (ii) deduz-se de \sref{III.1.4.12}.
A afirmação (iv) já foi demonstrada em \sref[II]{II.2.7.9}.

Agora vamos demonstrar simultaneamente (i) e (iii).
Eles valem para $\sh{F}=\sh{O}_X(m)$ por \sref{III.2.1.13} e, portanto, também quando $\sh{F}$ é a soma direta de um número finito de $\sh{O}_X$-módulos da forma $\sh{O}_X(m_j)$.
São trivialmente verdadeiras para $q>r$ por (ii).
Procedemos por indução descendente sobre $q$.

O feixe $\sh{F}$ é isomorfo a um quociente de uma soma direta $\sh{E}$ de um número finito de feixes $\sh{O}_X(m_j)$ \sref[II]{II.2.7.10}.
Existe, portanto, uma sequência exata
\[
  0\longrightarrow\sh{R}\longrightarrow\sh{E}
   \longrightarrow\sh{F}\longrightarrow0,
\]
onde $\sh{R}$ é coerente \sref[0]{0.5.3.3} e $\sh{E}$ satisfaz (i) e (iii).
Como a mudança de grau é exata, para todo $n\in\bb{Z}$ temos também uma sequência exata
\[
  0\longrightarrow\sh{R}(n)\longrightarrow\sh{E}(n)
   \longrightarrow\sh{F}(n)\longrightarrow0,
\]
e, portanto, uma sequência exata de cohomologia
\[
  \HH^{q-1}(X,\sh{E}(n))
  \longrightarrow\HH^{q-1}(X,\sh{F}(n))
  \longrightarrow\HH^q(X,\sh{R}(n)).
\]
Como $\sh{E}(n)$ é uma soma direta de feixes $\sh{O}_X(n+m_j)$ \sref[II]{II.2.5.14}, o módulo $\HH^{q-1}(X,\sh{E}(n))$ é do tipo finito.
Analogamente, o módulo $\HH^q(X,\sh{R}(n))$ é do tipo finito pela hipótese de indução.
Como $A$ é noetheriano, deduz-se que $\HH^{q-1}(X,\sh{F}(n))$ é do tipo finito, o que demonstra (i).

Pela hipótese de indução podemos escolher $N$ de modo que $\HH^q(X,\sh{R}(n))=0$ para $n\geq N$.
Podemos também escolher $N$ para que $\HH^{q-1}(X,\sh{E}(n))=0$ para $n\geq N$, já que $\sh{E}$ satisfaz (iii).
A sequência exata de cohomologia dá então
\[
  \HH^{q-1}(X,\sh{F}(n))=0
  \qquad(n\geq N),
\]
o que conclui a demonstração.
\end{proof}

\begin{corollary}[2.2.3]
\label{III.2.2.3}
Sob as hipóteses de \sref{III.2.2.1}, seja
\[
  \sh{F}\longrightarrow\sh{G}\longrightarrow\sh{H}
\]
uma sequência exata de $\sh{O}_X$-módulos coerentes.
Existe um inteiro $N$ tal que, para $n\geq N$, a sequência
\[
  f_*(\sh{F}(n))
  \longrightarrow f_*(\sh{G}(n))
  \longrightarrow f_*(\sh{H}(n))
\]
é exata.
\end{corollary}

\begin{proof}
Sejam $\sh{F}'$, $\sh{G}'$ e $\sh{G}''$, respectivamente, o núcleo, a imagem e o conúcleo de $\sh{F}\to\sh{G}$.
Então $\sh{G}'$ é o núcleo e $\sh{G}''$ a imagem de $\sh{G}\to\sh{H}$; seja $\sh{H}''$ o conúcleo deste último homomorfismo.
Todos estes $\sh{O}_X$-módulos são coerentes \sref[0]{0.5.3.4}.
Como a mudança de grau é exata, basta provar que, para $n$ suficientemente grande, cada uma das sequências
\begin{align*}
  0&\longrightarrow f_*(\sh{F}'(n))
    \longrightarrow f_*(\sh{F}(n))
    \longrightarrow f_*(\sh{G}'(n))
    \longrightarrow0,\\
  0&\longrightarrow f_*(\sh{G}'(n))
    \longrightarrow f_*(\sh{G}(n))
    \longrightarrow f_*(\sh{G}''(n))
    \longrightarrow0,\\
  0&\longrightarrow f_*(\sh{G}''(n))
    \longrightarrow f_*(\sh{H}(n))
    \longrightarrow f_*(\sh{H}''(n))
    \longrightarrow0
\end{align*}
é exata.
Consequentemente, podemos supor que
\[
  0\longrightarrow\sh{F}\longrightarrow\sh{G}
   \longrightarrow\sh{H}\longrightarrow0
\]
é exata.
A sequência de cohomologia é então
\[
  0\longrightarrow f_*(\sh{F}(n))
  \longrightarrow f_*(\sh{G}(n))
  \longrightarrow f_*(\sh{H}(n))
  \longrightarrow \RR^1f_*(\sh{F}(n))
  \longrightarrow\cdots,
\]
e a conclusão deduz-se de \sref{III.2.2.1}, (ii).
\end{proof}

\oldpage[III]{102}

\begin{corollary}[2.2.4]
\label{III.2.2.4}
Seja $Y$ um pré-esquema noetheriano, seja $f:X\to Y$ um morfismo de tipo finito e seja $\sh{L}$ um $\sh{O}_X$-módulo invertível amplo em relação a $f$.
Para todo $\sh{O}_X$-módulo $\sh{F}$, ponhamos
\[
  \sh{F}(n)=\sh{F}\otimes_{\sh{O}_X}\sh{L}^{\otimes n}
  \qquad(n\in\bb{Z}).
\]
Seja
\[
  \sh{F}\longrightarrow\sh{G}\longrightarrow\sh{H}
\]
uma sequência exata de $\sh{O}_X$-módulos coerentes tal que os suportes de $\sh{F}$ e $\sh{H}$ sejam próprios sobre $Y$ \sref[II]{II.5.4.10}.
Então há um inteiro $N$ tal que, para $n\geq N$, a sequência
\[
  f_*(\sh{F}(n))
  \longrightarrow f_*(\sh{G}(n))
  \longrightarrow f_*(\sh{H}(n))
\]
é exata.
\end{corollary}

\begin{proof}
O raciocínio do início de \sref{III.2.2.1} mostra que, se o corolário for verdadeiro para $\sh{L}^{\otimes d}$, onde $d>0$, então será verdadeiro para $\sh{L}$.
Podemos, portanto, supor que $\sh{L}$ é muito amplo em relação a $f$ \sref[II]{II.4.6.11}.
Podemos então identificar $X$ com um subconjunto aberto de um $Y$-pré-esquema
\[
  Z=\Proj(\sh{S}),
\]
onde $\sh{S}$ é uma $\sh{O}_Y$-álgebra quase-coerente graduada em graus positivos, $\sh{S}_1$ é de tipo finito e gera $\sh{S}$, e
\[
  \sh{L}=i^*(\sh{O}_Z(1)),
\]
sendo $i:X\to Z$ a imersão canônica \sref[II]{II.4.4.7}.

Agora $\Supp(\sh{G})$ é fechado em $X$ e está contido em
\[
  \Supp(\sh{F})\cup\Supp(\sh{H});
\]
portanto, é próprio sobre $Y$.
Consequentemente, os suportes de $\sh{F}$, $\sh{G}$ e $\sh{H}$ são fechados em $Z$ \sref[II]{II.5.4.10}.
Assim,
\[
  \sh{F}'=i_*(\sh{F}),\qquad
  \sh{G}'=i_*(\sh{G}),\qquad
  \sh{H}'=i_*(\sh{H})
\]
são $\sh{O}_Z$-módulos coerentes, e a sequência
\[
  \sh{F}'\longrightarrow\sh{G}'\longrightarrow\sh{H}'
\]
é exata.
Se $g:Z\to Y$ é o morfismo estrutural, então $f=g\circ i$, e claramente
\[
  \sh{F}'(n)=i_*(\sh{F}(n)),
\]
com identidades análogas para $\sh{G}'$ e $\sh{H}'$.
A conclusão deduz-se aplicando \sref{III.2.2.3} a $\sh{F}'$, $\sh{G}'$ e $\sh{H}'$.
\end{proof}

\begin{env}[2.2.5]
\label{III.2.2.5}
\emph{Observações.}

(i) A afirmação (i) de \sref{III.2.2.1} continua sendo verdadeira se se supõe apenas que $Y$ é \emph{localmente noetheriano}.
Com efeito, a propriedade é manifestamente local sobre $Y$.
Além disso, as hipóteses de \sref{III.2.2.1} implicam que, para todo subconjunto aberto $U\subset Y$, a restrição de $f$ a $f^{-1}(U)$ é um morfismo projetivo
\[
  f^{-1}(U)\longrightarrow U
\]
\sref[II]{II.5.5.5}[iii], e $\sh{L}|_{f^{-1}(U)}$ é amplo para este morfismo \sref[II]{II.4.6.4}.

(ii) A afirmação (iii) de \sref{III.2.2.1} continua válida, como vimos, se se supõe apenas que $X$ é um esquema quase compacto, ou um pré-esquema cujo espaço subjacente é noetheriano, e que $f:X\to Y$ é quase compacto \sref[II]{II.4.6.8}.
É interessante observar, no entanto, que mesmo quando $Y$ é o \emph{espectro de um corpo} $K$ e $f$ é \emph{quase-projetivo}, a afirmação (ii) de \sref{III.2.2.1} pode deixar de se cumprir.

Por exemplo, seja
\[
  X'=\Spec(K[T_0,\ldots,T_r])
\]
e seja $X$ a união dos subconjuntos abertos afins $D(T_i)$ de $X'$, para $0\leq i\leq r$.
Como a imersão $X\to X'$ é quase compacta, o morfismo estrutural $f:X\to Y$ é quase afim \sref[II]{II.5.1.10} e, portanto, $\sh{O}_X$ é muito amplo em relação a $f$ \sref[II]{II.5.1.6}.
Mas
\[
  \Gamma(X,\sh{O}_X)
  =\bigcap_{i=0}^r(K[T_0,\ldots,T_r])_{T_i}
  =K[T_0,\ldots,T_r]
\]
por \sref[I]{I.8.2.1.1}.
Deduz-se de \sref{III.1.4.3.1} e \sref{III.1.1.3.5} que
\[
  \HH^r(X,\sh{O}_X^{\otimes n})
  =\HH^r(X,\sh{O}_X)
  =A\neq0
\]
para todo $n$.
\end{env}


\subsection{Aplicação a feixes graduados de álgebras e de módulos}
\label{subsection:III.2.3}

\begin{theorem}[2.3.1]
\label{III.2.3.1}
Seja $Y$ um pré-esquema noetheriano, seja $\sh{S}$ uma $\sh{O}_Y$-álgebra quase-coerente de tipo finito, graduada em graus positivos, seja
\[
  X=\Proj(\sh{S}),
\]
e seja $q:X\to Y$ o morfismo estrutural.
Seja $\sh{M}$ um $\sh{S}$-módulo graduado quase-coerente que satisfaça a condição \textbf{(TF) }.
Existe então um inteiro $N$ tal que, para $n\geq N$, o homomorfismo canônico \sref[II]{II.8.14.5.1}
\[
  \alpha_n:\sh{M}_n
  \longrightarrow q_*\bigl(\shProj_0(\sh{M}(n))\bigr)
  =q_*\bigl((\shProj(\sh{M}))_n\bigr)
\]
\oldpage[III]{103}
é bijetivo.
Em outras palavras, o homomorfismo canônico
\[
  \alpha:\sh{M}\longrightarrow\bbGamma_*(\shProj(\sh{M}))
\]
é um \textbf{(TN) }-isomorfismo.
\end{theorem}

\begin{proof}
Podemos restringir-nos ao caso em que $\sh{M}$ é um $\sh{S}$-módulo de tipo finito \sref[II]{II.3.4.2}.
Como $Y$ é quase compacto, existe um inteiro $d>0$ tal que $\sh{S}^{(d)}$ é gerada pelo $\sh{O}_Y$-módulo quase-coerente $\sh{S}_d$, que é do tipo finito \sref[II]{II.3.1.10} e, portanto, coerente porque $Y$ é noetheriano.
Agora, $\sh{M}$ é a soma direta dos $\sh{M}^{(d,k)}$, para $0\leq k<d$, e cada $\sh{M}^{(d,k)}$ é um $\sh{S}^{(d)}$-módulo quase-coerente de tipo finito, por \sref[II]{II.2.1.6}[iii].
Como a questão é local sobre $Y$, basta manifestamente demonstrar que cada homomorfismo canônico
\[
  \alpha:\sh{M}^{(d,k)}
  \longrightarrow
  \bbGamma_*\bigl((\shProj(\sh{M}))^{(d,k)}\bigr)
\]
é um \textbf{(TN) }-isomorfismo.
Tendo em conta \sref[II]{II.8.14.13}, e em particular o diagrama \sref[II]{II.8.14.13.4}, ficamos assim reduzidos a demonstrar o teorema quando $\sh{S}$ é gerada por $\sh{S}_1$ e $\sh{S}_1$ é um $\sh{O}_Y$-módulo coerente.

Como $Y$ é noetheriano, o raciocínio usado no início da demonstração de \sref{III.2.2.2} mostra que podemos reduzir-nos ao caso
\[
  Y=\Spec(A),\qquad
  \sh{S}=\widetilde{S},\qquad
  \sh{M}=\widetilde{M},
\]
onde $A$ é um anel noetheriano, $S_1$ é um $A$-módulo de tipo finito e $M$ é um $S$-módulo graduado de tipo finito.
Vamos mostrar que basta então provar o teorema para $M=S$.
Com efeito, no caso geral existe uma sequência exata
\[
  L'\longrightarrow L\longrightarrow M\longrightarrow0,
\]
onde $L$ e $L'$ são somas diretas de módulos graduados da forma $S(m)$.
Se o resultado vale para $M=S$, também vale para $M=S(m)$ e, portanto, para $L$ e $L'$.
Consideremos o diagrama comutativo
\[
  \xymatrix{
    \widetilde{L'_n}\ar[r]\ar[d]_{\alpha_n} &
    \widetilde{L_n}\ar[r]\ar[d]^{\alpha_n} &
    \widetilde{M_n}\ar[r]\ar[d]^{\alpha_n} &
    0\\
    q_*(\widetilde{L'}(n))\ar[r] &
    q_*(\widetilde{L}(n))\ar[r] &
    q_*(\widetilde{M}(n))\ar[r] &
    0 .
  }
\]
\oldpage[III]{104}
A segunda fila é exata por \sref{III.2.2.3} quando $n$ é suficientemente grande.
O mesmo acontece com a primeira fila e, como as duas flechas verticais à esquerda são isomorfismos, também é a terceira.

Resta, então, provar o teorema para $M=S$.
Suponhamos que, em primeiro lugar,
\[
  S=A[T_0,\ldots,T_r],
\]
onde os $T_i$ são indeterminadas.
Neste caso, a afirmação é precisamente \sref{III.2.1.11}[ii].
No caso geral, $S$ identifica-se com um quociente de um anel
\[
  S'=A[T_0,\ldots,T_r]
\]
por um ideal graduado, de modo que $X$ é um subesquema fechado de
\[
  X'=\bb{P}_A^r
\]
\sref[II]{II.2.9.2}.
Se $j:X\to X'$ é a injeção canônica, então $j_*(\widetilde{S}(n))$ é o $\sh{O}_{X'}$-módulo
\[
  (\shProj(\widetilde{S}))(n),
\]
onde $S$ é considerado como $S'$-módulo graduado; isso se deduz imediatamente de \sref[II]{II.2.8.7}.
Como $j_*(\widetilde{S}(n))$ é um $\sh{O}_{X'}$-módulo que satisfaz \textbf{(TF) }, o homomorfismo canônico
\[
  \alpha_n:S_n\longrightarrow
  \Gamma\bigl(X',j_*(\widetilde{S}(n))\bigr)
\]
é bijetivo para todo $n$ suficientemente grande, pelo caso polinomial que acabamos de demonstrar.
Isso conclui a demonstração, já que
\[
  \Gamma\bigl(X',j_*(\widetilde{S}(n))\bigr)
  =\Gamma(X,\widetilde{S}(n)).
\]
\end{proof}

\begin{corollary}[2.3.2]
\label{III.2.3.2}
Sob as hipóteses de \sref{III.2.3.1}, seja
\[
  \sh{S}_X=\bigoplus_{n\in\bb{Z}}\sh{O}_X(n),
\]
e seja $\sh{F}$ um $\sh{S}_X$-módulo graduado quase-coerente de tipo finito.
Então $\bbGamma_*(\sh{F})$ satisfaz a condição \textbf{(TF) }.
\end{corollary}

\begin{proof}
Na demonstração de \sref{III.2.3.1} vimos que $X$, que é isomorfo a $\Proj(\sh{S}^{(d)})$ \sref[II]{II.3.1.8}, é de tipo finito sobre $Y$ \sref[II]{II.3.4.1}.
Deduz-se de \sref[II]{II.8.14.9} que $\sh{F}$ é isomorfo a um $\sh{S}_X$-módulo graduado da forma $\shProj(\sh{M})$, onde $\sh{M}$ é um $\sh{S}$-módulo graduado quase-coerente do tipo finito.
Por \sref{III.2.3.1}, $\bbGamma_*(\sh{F})$ é \textbf{(TN) }-isomorfo a $\sh{M}$ e, portanto, satisfaz \textbf{(TF) }.
\end{proof}

\begin{env}[2.3.3]
\label{III.2.3.3}
\emph{Escolio.}
Seja $Y$ um pré-esquema noetheriano, seja $\sh{S}$ uma $\sh{O}_Y$-álgebra graduada que satisfaz as hipóteses de \sref{III.2.3.1}, e seja $X=\Proj(\sh{S})$.
Seja $\cat{K}_{\sh{S}}$ a categoria abeliana dos $\sh{S}$-módulos graduados quase-coerentes que satisfazem \textbf{(TF) }, e seja $\cat{K}'_{\sh{S}}$ a subcategoria de $\cat{K}_{\sh{S}}$ formada pelos $\sh{S}$-módulos que satisfazem \textbf{(TN) }.
Por último, seja $\cat{K}_X$ a categoria dos $\sh{S}_X$-módulos graduados quase-coerentes $\sh{F}$ de tipo finito; como $\sh{S}_X$ é periódica \sref[II]{II.8.14.4} e \sref[II]{II.8.14.12}, isto equivale a exigir que os $\sh{F}_i$ sejam $\sh{O}_X$-módulos coerentes.

Por \sref[II]{II.8.14.8}, \sref[II]{II.8.14.10} e \sref{III.2.3.2}, os funtores
\[
  \sh{M}\longmapsto\shProj(\sh{M}),
  \qquad
  \sh{F}\longmapsto\bbGamma_*(\sh{F})
\]
definem uma equivalência (T, I, 1.2) entre a categoria quociente
\[
  \cat{K}_{\sh{S}}/\cat{K}'_{\sh{S}}
\]
(T, I, 1.11) e a categoria $\cat{K}_X$.
Quando $\sh{S}$ é gerada por $\sh{S}_1$, a categoria $\cat{K}_X$ pode ser substituída pela categoria dos $\sh{O}_X$-módulos coerentes \sref[II]{II.8.14.12}.
\end{env}

\begin{proposition}[2.3.4]
\label{III.2.3.4}
Seja $Y$ um pré-esquema noetheriano.

(i) Seja $\sh{S}$ uma $\sh{O}_Y$-álgebra quase-coerente de tipo finito, graduada em graus positivos.
Ponhamos
\[
  X=\Proj(\sh{S}),\qquad
  \sh{S}_X=\shProj(\sh{S})
  =\bigoplus_{n\in\bb{Z}}\sh{O}_X(n).
\]
Então $\sh{S}_X$ é uma $\sh{O}_X$-álgebra graduada quase-coerente periódica \sref[II]{II.8.14.12}, suas componentes homogêneas
\[
  (\sh{S}_X)_n=\sh{O}_X(n)
\]
são $\sh{O}_X$-módulos coerentes e, se $d>0$ é um período de $\sh{S}_X$, então
\[
  (\sh{S}_X)_d=\sh{O}_X(d)
\]
é $Y$-amplo como $\sh{O}_X$-módulo invertível.
Além disso, o homomorfismo canônico
\[
  \alpha:\sh{S}\longrightarrow\bbGamma_*(\sh{S}_X)
\]
é um \textbf{(TN) }-isomorfismo.

(ii) Reciprocamente, seja $q:X\to Y$ um morfismo projetivo e seja $\sh{S}'$ uma $\sh{O}_X$-álgebra graduada cujas componentes homogêneas $\sh{S}'_n$, para $n\in\bb{Z}$, são $\sh{O}_X$-módulos coerentes.
Suponhamos que $\sh{S}'$ tenha um período $d>0$ tal que $\sh{S}'_d$ seja um $\sh{O}_X$-módulo invertível amplo para $q$.
Então
\[
  \sh{S}=\bigoplus_{n\geq0}q_*(\sh{S}'_n)
\]
é uma $\sh{O}_Y$-álgebra quase-coerente de tipo finito, graduada em graus positivos, e existe um $Y$-isomorfismo
\[
  r:X\simto\Proj(\sh{S})
\]
tal que $r^*(\shProj(\sh{S}))$ é isomorfo a $\sh{S}'$ como $\sh{O}_X$-álgebra graduada.
\end{proposition}

\begin{proof}
(i) Quase todas as afirmações já foram demonstradas; a última é um caso particular de \sref{III.2.3.2}.
A periodicidade de $\sh{S}_X$ foi demonstrada em \sref[II]{II.8.14.14}, e a afirmação de que existe um período $d>0$ para o qual $\sh{O}_X(d)$ é invertível e $Y$-amplo é precisamente \sref[II]{II.4.6.18}.
Finalmente, para $0\leq k<d$, $(\sh{S}_X)^{(d,k)}$ é um $(\sh{S}_X)^{(d)}$-módulo de tipo finito \sref[II]{II.8.14.14}.
Deduz-se que cada $(\sh{S}_X)_n$ é um $\sh{O}_X$-módulo quase-coerente de tipo finito \sref[II]{II.2.1.6}[ii]; como $\sh{O}_X$ é coerente, cada $(\sh{S}_X)_n$ também o é.

(ii) Substituindo o período $d$ por um de seus múltiplos, podemos supor que
\[
  \sh{L}=\sh{S}'_d
\]
é muito amplo em relação a $q$ \sref[II]{II.4.6.11}.
Por hipótese,
\[
  \sh{S}'^{(d)}
  =\bigoplus_{n\in\bb{Z}}\sh{L}^{\otimes n},
\]
e, portanto,
\[
  \sh{S}^{(d)}
  =\bigoplus_{n\geq0}q_*(\sh{L}^{\otimes n}).
\]
Por \sref[II]{II.3.1.8} e \sref[II]{II.3.2.9}, existe um $Y$-isomorfismo
\[
  s:X'=\Proj(\sh{S})\simto
  X''=\Proj(\sh{S}^{(d)})
\]
tal que
\oldpage[III]{105}
\[
  s^*(\sh{O}_{X''}(n))=\sh{O}_{X'}(nd)
\]
para todo $n$ \sref[II]{II.3.5.2}.
Assim, estabeleceremos a existência de um $Y$-isomorfismo $X\simto X'$ demonstrando a seguinte afirmação.

\begin{env}[2.3.4.1]
\label{III.2.3.4.1}
Seja $Y$ um pré-esquema noetheriano, seja $q:X\to Y$ um morfismo projetivo e seja $\sh{L}$ um $\sh{O}_X$-módulo invertível muito amplo para $q$.
Então
\[
  \sh{S}=\bigoplus_{n\geq0}q_*(\sh{L}^{\otimes n})
\]
é uma $\sh{O}_Y$-álgebra graduada quase-coerente de tipo finito tal que
\[
  \sh{S}_n=\sh{S}_1^n
\]
para todo $n$ suficientemente grande, e existe um $Y$-isomorfismo
\[
  r:X\simto P=\Proj(\sh{S})
\]
tal que
\[
  \sh{L}=r^*(\sh{O}_P(1)).
\]
\end{env}

Como $q$ é projetivo, deduz-se de \sref[II]{II.5.4.4} e \sref[II]{II.4.4.7} que existe um $Y$-isomorfismo
\[
  r':X\simto P'=\Proj(\sh{T}),
\]
onde $\sh{T}$ é uma $\sh{O}_Y$-álgebra quase-coerente, $\sh{T}_1$ é de tipo finito e gera $\sh{T}$, e
\[
  \sh{L}=r'^*(\sh{O}_{P'}(1)).
\]
Se $q':P'\to Y$ é o morfismo estrutural, então
\[
  \sh{S}=\bigoplus_{n\geq0}q'_*(\sh{O}_{P'}(n)).
\]
Por \sref{III.2.3.1}, o homomorfismo canônico
\[
  \alpha_n:\sh{T}_n\longrightarrow
  \sh{S}_n=q'_*(\sh{O}_{P'}(n))
\]
é bijetivo para todo $n$ suficientemente grande.
Como $\sh{T}_n=\sh{T}_1^n$, deduz-se \emph{a fortiori} que $\sh{S}_n=\sh{S}_1^n$ para todo $n$ suficientemente grande.

O homomorfismo canônico $\alpha:\sh{T}\to\sh{S}$ de $\sh{O}_Y$-álgebras graduadas é um \textbf{(TN) }-isomorfismo.
Consequentemente,
\phantomsection\label{III.2.3.5.1}
\[
  \Phi=\Proj(\alpha):\Proj(\sh{S})\longrightarrow\Proj(\sh{T})
\]
é um isomorfismo \sref[II]{II.3.6.1}.
Além disso, \sref[II]{II.3.5.2} e o fato de que os $\sh{T}$-módulos graduados obtidos de $\sh{S}(n)$ por restrição de escalares são \textbf{(TN) }-isomorfos a $\sh{T}(n)$ dão
\[
  \Phi_*(\sh{O}_P(n))=\sh{O}_{P'}(n)
\]
para todo $n$ \sref[II]{II.3.4.2}.
Só resta mostrar que $\sh{S}$ é de tipo finito sobre $\sh{O}_Y$.
Os módulos
\[
  \sh{S}_n=q'_*(\sh{O}_{P'}(n))
\]
são coerentes por \sref{III.2.2.1}; e, se $\sh{S}_n=\sh{S}_1^n$ para $n\geq n_0$, então $\sh{S}$ é gerada pelo módulo coerente
\[
  \bigoplus_{i\leq n_0}\sh{S}_i.
\]
Isto demonstra a afirmação \sref[I]{I.9.6.2}.

Voltemos à demonstração de \sref{III.2.3.4}, mantendo a notação.
Demonstramos que existe um $Y$-isomorfismo
\[
  r'':X\simto X''
\]
tal que
\[
  r''_*(\sh{L}^{\otimes n})=\sh{O}_{X''}(n)
\]
para todo $n\in\bb{Z}$.
Seja $q'':X''\to Y$ o morfismo estrutural.
A álgebra $\sh{S}'$ é a soma direta dos $\sh{S}'^{(d)}$-módulos graduados $\sh{S}'^{(d,k)}$.
Cada um deles é um $\sh{S}'^{(d)}$-módulo quase-coerente de tipo finito, pela periodicidade de $\sh{S}'$ e pela hipótese de que os $\sh{S}'_n$ são $\sh{O}_X$-módulos de tipo finito \sref[II]{II.8.14.12}.
Ponhamos
\[
  \sh{F}^{(k)}=r''_*(\sh{S}'^{(d,k)}),
\]
de modo que os $\sh{F}^{(k)}$ são $\sh{S}_{X''}$-módulos graduados quase-coerentes de tipo finito.
Por \sref[II]{II.8.14.8}, o homomorfismo canônico
\[
  \beta:\shProj\bigl(\bbGamma_*(\sh{F}^{(k)})\bigr)
  \longrightarrow\sh{F}^{(k)}
\]
é um isomorfismo de $\sh{S}_{X''}$-módulos.
Por outro lado,
\[
  q''_*\bigl((\sh{F}^{(k)})_n\bigr)
  =q_*\bigl((\sh{S}'^{(d,k)})_n\bigr),
\]
e para $n\geq0$ o último $\sh{O}_Y$-módulo é, por definição, igual a $(\sh{S}^{(d,k)})_n$.
Em outras palavras, a injeção canônica
\[
  \sh{S}^{(d,k)}
  \longrightarrow\bbGamma_*(\sh{F}^{(k)})
\]
é um \textbf{(TN) }-isomorfismo.
Portanto,
\[
  \shProj(\sh{S}^{(d,k)})
  =\shProj\bigl(\bbGamma_*(\sh{F}^{(k)})\bigr),
\]
e, consequentemente,
\[
  r''{}^*\bigl(\shProj(\sh{S}^{(d,k)})\bigr)
  \simeq\sh{S}'^{(d,k)}.
\]
Resta observar que, salvo isomorfismo canônico,
\[
  \shProj(\sh{S}^{(d,k)})
  =s_*\bigl((\shProj(\sh{S}))^{(d,k)}\bigr)
\]
\sref[II]{II.8.14.13.1}; isto demonstra que $r^*(\shProj(\sh{S}))$ é isomorfo a $\sh{S}'$.

Por fim, por \sref{III.2.3.2}, cada $\bbGamma_*(\sh{F}^{(k)})$ satisfaz a condição \textbf{(TF) } e, portanto, também a satisfaz cada $\sh{S}^{(d,k)}$.
Além disso, os $\sh{S}_n=q_*(\sh{S}'_n)$ são coerentes por \sref{III.2.2.1}, porque os $\sh{S}'_n$ são coerentes.
Deduz-se imediatamente que os $\sh{S}^{(d,k)}$ são $\sh{S}^{(d)}$-módulos de tipo finito.
Como \sref{III.2.3.4.1} mostra que $\sh{S}^{(d)}$ é uma $\sh{O}_Y$-álgebra de tipo finito, $\sh{S}$ é também uma $\sh{O}_Y$-álgebra de tipo finito.
\end{proof}

\oldpage[III]{106}
\begin{proposition}[2.3.5]
\label{III.2.3.5}
Seja $Y$ um pré-esquema noetheriano integral, seja $X$ um pré-esquema integral e seja
\[
  f:X\longrightarrow Y
\]
um morfismo projetivo birracional.
Existe então um ideal fracionário coerente
\[
  \sh{J}\subset\sh{R}(Y)
\]
\sref[II]{II.8.1.2} tal que $X$ é $Y$-isomorfo ao pré-esquema obtido por explosão de $\sh{J}$ \sref[II]{II.8.1.3}.
Além disso, existe um subconjunto aberto $U$ de $Y$ tal que a restrição de $f$ a $f^{-1}(U)$ é um isomorfismo de $f^{-1}(U)$ sobre $U$ (cf. \ \textbf{I}, 6.5.5), e tal que $\sh{J}|U$ é invertível.
\end{proposition}

\begin{proof}
Existe um $\sh{O}_X$-módulo invertível $\sh{L}$ muito amplo para $f$ \sref[II]{II.4.4.2} e \sref[II]{II.5.3.2}.
Aplicando \sref{III.2.3.4.1}, identificamos $X$ com $\Proj(\sh{S})$, onde
\[
  \sh{S}=\bigoplus_{n\geq0}f_*(\sh{L}^{\otimes n}).
\]
Os módulos $f_*(\sh{L}^{\otimes n})$ são sem torção sobre $\sh{O}_Y$ \sref[I]{I.7.4.5}.
Assim, $\sh{S}$ identifica-se canonicamente com um sub-$\sh{O}_Y$-módulo de
\[
  \sh{S}\otimes_{\sh{O}_Y}\sh{R}(Y)
\]
\sref[I]{I.7.4.1}.
O último é um feixe simples \sref[I]{I.7.3.6} e é, portanto, determinado por sua restrição a um subconjunto aberto não vazio.

Escolhamos um subconjunto aberto não vazio $U'\subset U$ sobre o qual $\sh{L}|f^{-1}(U')$ seja isomorfo a $\sh{O}_X|f^{-1}(U')$ e, portanto, a $\sh{O}_Y|U'$.
As restrições $f_*(\sh{L}^{\otimes n})|U'$ são então isomorfas a $\sh{O}_Y|U'$.
Segue-se que
\[
  \sh{S}\otimes_{\sh{O}_Y}\sh{R}(Y)
  \simeq\sh{R}(Y)[T],
\]
onde $T$ é uma indeterminada.
Por \sref{III.2.3.4.1}, $\sh{S}$ é \textbf{(TN) }-isomorfa à sub-$\sh{O}_Y$-álgebra gerada pela imagem canônica de $f_*(\sh{L})$ nesta álgebra de polinômios.
Sob a identificação mostrada, essa imagem é $\sh{J}\cdot T$, onde $\sh{J}$ é um sub-$\sh{O}_Y$-módulo coerente de $\sh{R}(Y)$ \sref{III.2.2.1}.
Sua restrição a $U'$ é isomorfa a $\sh{O}_Y|U'$ e, reduzindo $U$ se necessário, $\sh{J}|U$ é invertível.
Consequentemente,
\[
  \sh{S}\quad\text{é \textbf{(TN)}-isomorfa a}\quad
  \bigoplus_{n\geq0}\sh{J}^n,
\]
o que conclui a demonstração.
\end{proof}

\begin{corollary}[2.3.6]
\label{III.2.3.6}
Sob as hipóteses de \sref{III.2.3.5}, suponhamos ainda que, para todo sub-$\sh{O}_Y$-módulo coerente não nulo,
\[
  \sh{J}\subset\sh{R}(Y)
\]
existe um $\sh{O}_Y$-módulo invertível $\sh{L}$ tal que
\[
  \Gamma\bigl(Y,\sh{L}\otimes_{\sh{O}_Y}
  \shHom(\sh{J},\sh{O}_Y)\bigr)\neq0.
\]
Então, na declaração de \sref{III.2.3.5}, pode-se tomar $\sh{J}$ como um ideal de $\sh{O}_Y$.
Esta condição adicional é sempre satisfeita se houver um $\sh{O}_Y$-módulo amplo.
\end{corollary}

\begin{proof}
De fato, por \sref[0]{0.5.4.2},
\[
  \sh{L}\otimes\shHom(\sh{J},\sh{O}_Y)
  =\shHom\bigl(\sh{L}^{-1},\shHom(\sh{J},\sh{O}_Y)\bigr)
  =\shHom(\sh{J}\otimes\sh{L}^{-1},\sh{O}_Y).
\]
A hipótese diz, portanto, que existe um homomorfismo não nulo
\[
  u:\sh{J}\otimes\sh{L}^{-1}\longrightarrow\sh{O}_Y.
\]
Para todo $y\in Y$, o módulo $(\sh{J}\otimes\sh{L}^{-1})_y$ identifica-se com um sub-$\sh{O}_y$-módulo do corpo de frações $\sh{R}(Y)_y$ de $\sh{O}_y$ \sref[I]{I.7.1.5}.
Assim, $u_y$ é necessariamente injetivo, e $u$ identifica $\sh{J}\otimes\sh{L}^{-1}$ com um ideal $\sh{J}'$ de $\sh{O}_Y$.
Mas
\[
  \Proj\left(\bigoplus_{n\geq0}\sh{J}^n\right)
  \quad\text{e}\quad
  \Proj\left(\bigoplus_{n\geq0}
  (\sh{J}\otimes\sh{L}^{-1})^n\right)
\]
são $Y$-isomorfos \sref[II]{II.3.1.8}; isto demonstra a primeira afirmação.

Para a segunda, observe-se que
\[
  \sh{F}=\shHom_{\sh{O}_Y}(\sh{J},\sh{O}_Y)
\]
é coerente e não nulo, pois existe um subconjunto aberto $U$ de $Y$ sobre o qual $\sh{J}|U$ é invertível.
Se $\sh{L}$ é um $\sh{O}_Y$-módulo amplo, existe um inteiro $n$ tal que
\[
  \sh{F}(n)=\sh{F}\otimes\sh{L}^{\otimes n}
\]
é gerado por suas seções sobre $Y$ \sref[II]{II.4.5.5}.
\emph{A fortiori},
\[
  \Gamma(Y,\sh{F}(n))\neq0,
\]
o que demonstra a afirmação.
\end{proof}

\begin{corollary}[2.3.7]
\label{III.2.3.7}
Sejam $X$ e $Y$ esquemas integrais projetivos sobre um corpo $k$, e seja
\[
  f:X\longrightarrow Y
\]
um $k$-morfismo birracional.
Então $X$ é $k$-isomorfo a um $Y$-esquema obtido por explosão de um subesquema fechado $Y'$ de $Y$, não necessariamente reduzido.
\end{corollary}

\begin{proof}
\oldpage[III]{107}
Com efeito, $f$ é projetivo \sref[II]{II.5.5.5}[v] e, como $Y$ é projetivo sobre $k$, satisfaz-se a condição adicional de \sref{III.2.3.6}.
Basta, portanto, tomar o subesquema fechado $Y'$ de $Y$ definido pelo ideal coerente $\sh{J}$ de \sref{III.2.3.6}.
\end{proof}

\begin{env}[2.3.8]
\label{III.2.3.8}
\emph{Observação.}
No Capítulo~IV, ao estudar a noção de divisor, veremos que, se no enunciado de \sref{III.2.3.5} se supõe que os anéis locais $\sh{O}_y$, para $y\in Y$, são fatoriais---como acontece, por exemplo, quando $Y$ é não singular---, então $X$ pode ser obtido de $Y$ por explosão de um subpré-esquema fechado $Y'$ de $Y$ cujo espaço subjacente está contido em $Y-U$.
\end{env}


\subsection{Uma generalização do teorema fundamental}
\label{subsection:III.2.4}

\begin{theorem}[2.4.1]
\label{III.2.4.1}
Seja $Y$ um pré-esquema noetheriano, seja $\sh{S}$ uma $\sh{O}_Y$-álgebra quase-coerente de tipo finito, seja
\[
  f:X\longrightarrow Y
\]
um morfismo projetivo, ponhamos $\sh{S}'=f^*(\sh{S})$, e seja $\sh{M}$ um $\sh{S}'$-módulo quase-coerente de tipo finito.
Então:
\begin{enumerate}
  \item[(i) ] Para todo $p\in\bb Z$, $\RR^p f_*(\sh{M})$ é um $\sh{S}$-módulo de tipo finito.

  \item[(ii) ] Suponhamos ainda que $\sh{L}$ é um $\sh{O}_X$-módulo invertível amplo para $f$ e, para todo $n\in\bb Z$, ponhamos
  \[
    \sh{M}(n)=\sh{M}\otimes\sh{L}^{\otimes n}.
  \]
  Existe um inteiro $N$ tal que, para $n\geq N$,
  \[
  \label{III.2.4.1.1}
    \RR^p f_*(\sh{M}(n))=0
    \tag{2.4.1.1}
  \]
  para todo $p>0$, e o homomorfismo canônico
  \[
    f^*\bigl(f_*(\sh{M}(n))\bigr)\longrightarrow\sh{M}(n)
  \]
  \sref[0]{0.4.4.3} é sobrejetivo.
\end{enumerate}
\end{theorem}

\begin{proof}
Ponhamos
\[
  Y'=\Spec(\sh{S}),\qquad X'=\Spec(\sh{S}'),
\]
de modo que $X'=X\times_Y Y'$ \sref[II]{II.1.5.5}.
Sejam
\[
  g:Y'\longrightarrow Y,\qquad g':X'\longrightarrow X
\]
os morfismos estruturais, que são afins por definição, e seja
\[
  f'=f_{(Y')}:X'\longrightarrow Y'.
\]
Temos assim um diagrama comutativo
\[
  \xymatrix{
    X\ar[d]_f &
    X'\ar[l]_{g'}\ar[d]^{f'}\ar[dl]^{h}\\
    Y &
    Y'\ar[l]^{g},
  }
\]
onde $f'$ é projetivo \sref[II]{II.5.5.5}[iii] e
\[
  h=f\circ g'=g\circ f'.
\]

\emph{(i) } Seja $\widetilde{\sh{M}}$ o $\sh{O}_{X'}$-módulo associado ao $\sh{S}'$-módulo quase-coerente $\sh{M}$ quando $X'$ é considerado um $X$-esquema afim \sref[II]{II.1.4.3}; assim,
\[
  \sh{M}=g'_*(\widetilde{\sh{M}}).
\]
Como $\sh{M}$ é um $\sh{S}'$-módulo de tipo finito, $\widetilde{\sh{M}}$ é um $\sh{O}_{X'}$-módulo de tipo finito \sref[II]{II.1.4.5}.
Além disso, $h$ é de tipo finito, já que $g$ e $f'$ são de tipo finito \sref[II]{II.1.3.7} e \sref[I]{I.6.3.4}[ii].
Assim, $X'$ é noetheriano \sref[I]{I.6.3.7}, e $\widetilde{\sh{M}}$ é, portanto, coerente.

Como $g'$ é afim, o homomorfismo canônico
\[
  \RR^p f_*(\sh{M})
  \longrightarrow
  \RR^p h_*(\widetilde{\sh{M}})
\]
é bijetivo \sref{III.1.3.4}.
É, além disso, um homomorfismo de $\sh{S}$-módulos.
Com efeito, o homomorfismo canônico
\[
\label{III.2.4.1.2}
  g'_*(\sh{O}_{X'})
  \otimes_{\sh{O}_X}
  g'_*(\widetilde{\sh{M}})
  \longrightarrow
  g'_*(\widetilde{\sh{M}})
  \tag{2.4.1.2}
\]
define a estrutura de $\sh{S}'$-módulo sobre $\sh{M}$, já que
\[
  \sh{S}'=g'_*(\sh{O}_{X'}).
\]
Dá canonicamente lugar a um homomorfismo
\[
  f_*\bigl(g'_*(\sh{O}_{X'})\bigr)
  \otimes
  \RR^p f_*\bigl(g'_*(\widetilde{\sh{M}})\bigr)
  \longrightarrow
  \RR^p f_*\bigl(g'_*(\widetilde{\sh{M}})\bigr).
\]
\oldpage[III]{108}
Por \sref[0]{0.12.2.2}, e dado que \eqref{III.2.4.1.2} é obtido por sua vez, aplicando \sref[0]{0.4.2.2.1}, do homomorfismo
\[
  \sh{O}_{X'}\otimes\widetilde{\sh{M}}
  \longrightarrow\widetilde{\sh{M}}
\]
que define a estrutura de $\sh{O}_{X'}$-módulo sobre $\widetilde{\sh{M}}$, o diagrama
\[
  \xymatrix{
    f_*\bigl(g'_*(\sh{O}_{X'})\bigr)
    \otimes\RR^p f_*\bigl(g'_*(\widetilde{\sh{M}})\bigr)
    \ar[r]\ar[d] &
    \RR^p f_*\bigl(g'_*(\widetilde{\sh{M}})\bigr)\ar[d]\\
    h_*(\sh{O}_{X'})
    \otimes\RR^p h_*(\widetilde{\sh{M}})
    \ar[r] &
    \RR^p h_*(\widetilde{\sh{M}})
  }
\]
é comutativo \sref[0]{0.12.2.6}.
Compondo as flechas horizontais com o homomorfismo induzido pelo homomorfismo canônico
\[
  \sh{S}
  \longrightarrow f_*(f^*(\sh{S}))
  =f_*(\sh{S}')
  =f_*\bigl(g'_*(\sh{O}_{X'})\bigr)
  =h_*(\sh{O}_{X'}),
\]
obtemos a afirmação.

Por outro lado, como $g$ é afim e $f'$ é separado e quase compacto, o homomorfismo canônico
\[
  \RR^p h_*(\widetilde{\sh{M}})
  \longrightarrow
  g_*\bigl(\RR^p f'_*(\widetilde{\sh{M}})\bigr)
\]
é bijetivo \sref{III.1.4.14}.
Como antes, é um isomorfismo de $\sh{S}$-módulos, desta vez usando a comutatividade de \sref[0]{0.12.6.2}.
Agora, $f'$ é projetivo e $\widetilde{\sh{M}}$ é coerente, então $\RR^p f'_*(\widetilde{\sh{M}})$ é um $\sh{O}_{Y'}$-módulo coerente por \sref{III.2.2.1}.
Segue-se que
\[
  g_*\bigl(\RR^p f'_*(\widetilde{\sh{M}})\bigr)
\]
é um $\sh{S}$-módulo de tipo finito \sref[II]{II.1.4.5}.

\emph{(ii)} Ponhamos
\[
  \sh{L}'=g'^*(\sh{L}),
\]
um $\sh{O}_{X'}$-módulo invertível.
Para todo $n\in\bb Z$, salvo isomorfismo canônico,
\[
  g'_*\bigl(\widetilde{\sh{M}}\otimes\sh{L}'^{\otimes n}\bigr)
  =
  g'_*(\widetilde{\sh{M}})\otimes\sh{L}^{\otimes n}
  =
  \sh{M}(n)
\]
\sref[0]{0.5.4.10}.
Aplicando o raciocínio de \emph{(i) } a
\[
  \widetilde{\sh{M}}\otimes\sh{L}'^{\otimes n}
\]
mostra que
\[
  \RR^p f_*
  \bigl(g'_*(\widetilde{\sh{M}}\otimes\sh{L}'^{\otimes n})\bigr)
  \simeq
  g_*\bigl(\RR^p f'_*
  (\widetilde{\sh{M}}\otimes\sh{L}'^{\otimes n})\bigr).
\]
Agora, $\sh{L}'$ é amplo para $f'$ \sref[II]{II.4.6.13} [iii].
Deduz-se de \sref{III.2.2.1} que existe um inteiro $N$ tal que
\[
  \RR^p f'_*
  \bigl(\widetilde{\sh{M}}\otimes\sh{L}'^{\otimes n}\bigr)=0
\]
para todo $p>0$ e todo $n\geq N$.
Isso demonstra \eqref{III.2.4.1.1}.

Deduz-se mais uma vez de \sref{III.2.2.1} que pode ser escolhido $N$ de modo que, para $n\geq N$, o homomorfismo canônico
\[
  f'^*\bigl(f'_*(\widetilde{\sh{M}}\otimes\sh{L}'^{\otimes n})\bigr)
  \longrightarrow
  \widetilde{\sh{M}}\otimes\sh{L}'^{\otimes n}
\]
seja sobrejetivo.
Como $g'_*$ é exato \sref[II]{II.1.4.4}, o homomorfismo correspondente
\[
  g'_*\Bigl(
    f'^*\bigl(f'_*(\widetilde{\sh{M}}\otimes\sh{L}'^{\otimes n})\bigr)
  \Bigr)
  \longrightarrow
  g'_*(\widetilde{\sh{M}}\otimes\sh{L}'^{\otimes n})
  =\sh{M}(n)
\]
é sobrejetivo.
Mas \sref[II]{II.1.5.2} dá
\[
  g'_*\Bigl(
    f'^*\bigl(f'_*(\widetilde{\sh{M}}\otimes\sh{L}'^{\otimes n})\bigr)
  \Bigr)
  =
  f^*\Bigl(
    g_*\bigl(f'_*(\widetilde{\sh{M}}\otimes\sh{L}'^{\otimes n})\bigr)
  \Bigr),
\]
e $g_*\circ f'_*=f_*\circ g'_*$.
Consequentemente,
\[
  g'_*\Bigl(
    f'^*\bigl(f'_*(\widetilde{\sh{M}}\otimes\sh{L}'^{\otimes n})\bigr)
  \Bigr)
  =
  f^*\bigl(f_*(\sh{M}(n))\bigr),
\]
o que conclui a demonstração.
\end{proof}

\begin{env}[2.4.2]
\label{III.2.4.2}
Aplicaremos \sref{III.2.4.1}, em particular, quando $\sh{S}$ for uma $\sh{O}_Y$-álgebra graduada em graus positivos e
\[
  \sh{M}=\bigoplus_{k\in\bb Z}\sh{M}_k
\]
seja um $\sh{S}'$-módulo graduado.
\oldpage[III]{109}
Sob as mesmas hipóteses de finitude sobre $\sh{S}$ e $\sh{M}$, deduz-se de \sref{III.2.4.1} que
\[
  \bigoplus_{k\in\bb Z}\RR^p f_*(\sh{M}_k)
\]
é um $\sh{S}$-módulo de tipo finito para todo $p$.
Sob as hipóteses adicionais de \sref{III.2.4.1}[ii], existe também um inteiro $N$ tal que, para $n\geq N$,
\[
  \RR^p f_*(\sh{M}_k(n))=0
\]
para todo $p>0$ e todo $k\in\bb Z$, e o homomorfismo canônico
\[
  f^*\bigl(f_*(\sh{M}_k(n))\bigr)
  \longrightarrow
  \sh{M}_k(n)
\]
é sobrejetivo para todo $k\in\bb Z$.
\end{env}


\subsection{Característica de Euler--Poincar\'e e polinômio de Hilbert}
\label{subsection:III.2.5}

\begin{env}[2.5.1]
\label{III.2.5.1}
Seja $A$ um anel artiniano e seja $X$ um $A$-esquema projetivo sobre
\[
  Y=\Spec(A).
\]
Para todo $\sh{O}_X$-módulo coerente $\sh{F}$, os $A$-módulos
\[
  \HH^i(X,\sh{F})\qquad(i\geq0)
\]
são do tipo finito \sref{III.2.2.1} e, portanto, têm comprimento finito, já que $A$ é artiniano.
Além disso, \sref{III.2.2.1} mostra que $\HH^i(X,\sh{F})=0$ salvo para um número finito de $i\geq0$.
Assim, o inteiro
\[
\label{III.2.5.1.1}
  \chi_A(\sh{F})
  =
  \sum_{i=0}^{\infty}
  (-1)^i\operatorname{length}_A\bigl(\HH^i(X,\sh{F})\bigr)
  \tag{2.5.1.1}
\]
é definido para todo $\sh{O}_X$-módulo coerente $\sh{F}$.
Quando $A$ é um anel local artiniano, $\chi_A(\sh{F})$ chama-se a \emph{característica de Euler--Poincar\'e} de $\sh{F}$ em relação a $A$.
Para $\sh{F}=\sh{O}_X$, o inteiro $\chi_A(\sh{O}_X)$ chama-se o \emph{gênero aritmético} de $X$ em relação a $A$.
\end{env}

\begin{proposition}[2.5.2]
\label{III.2.5.2}
Seja
\[
  0\longrightarrow\sh{F}'
  \longrightarrow\sh{F}
  \longrightarrow\sh{F}''
  \longrightarrow0
\]
uma sequência exata de $\sh{O}_X$-módulos coerentes.
Então
\[
\label{III.2.5.2.1}
  \chi_A(\sh{F})
  =
  \chi_A(\sh{F}')+\chi_A(\sh{F}'').
  \tag{2.5.2.1}
\]
\end{proposition}

\begin{proof}
Os módulos de cohomologia de $\sh{F}$, $\sh{F}'$ e $\sh{F}''$ anulam-se em todos os graus suficientemente grandes.
Existe, portanto, um inteiro $r>0$ para o qual a longa sequência exata de cohomologia tem a forma
\[
  0\longrightarrow\HH^0(X,\sh{F}')
  \longrightarrow\HH^0(X,\sh{F})
  \longrightarrow\HH^0(X,\sh{F}'')
  \longrightarrow\HH^1(X,\sh{F}')
  \longrightarrow\cdots
\]
\[
  \cdots\longrightarrow\HH^r(X,\sh{F}')
  \longrightarrow\HH^r(X,\sh{F})
  \longrightarrow\HH^r(X,\sh{F}'')
  \longrightarrow0.
\]
Em uma sequência exata de $A$-módulos de comprimento finito com termos nulos em ambas as extremidades, a soma alternada dos comprimentos é nula \sref[0]{0.11.10.1}.
A aplicação deste fato dá imediatamente \eqref{III.2.5.2.1}.
\end{proof}

O resultado de \sref{III.2.5.2} aplica-se sempre que $X$ seja um $A$-esquema quase compacto e que os $A$-módulos $\HH^i(X,\sh{F})$ sejam de tipo finito para todo $\sh{O}_X$-módulo coerente $\sh{F}$ \sref{III.1.4.12}.

\begin{theorem}[2.5.3]
\label{III.2.5.3}
Seja $A$ um anel local artiniano, seja $X$ um esquema projetivo sobre
\[
  Y=\Spec(A),
\]
seja $\sh{L}$ um $\sh{O}_X$-módulo invertível muito amplo relativo a $Y$, e seja $\sh{F}$ um $\sh{O}_X$-módulo coerente.
Para todo $n\in\bb Z$, ponhamos
\[
  \sh{F}(n)=\sh{F}\otimes_{\sh{O}_X}\sh{L}^{\otimes n}.
\]
\begin{enumerate}
  \item[(i) ] Existe um único polinômio $P\in\bb Q[T]$ tal que
  \[
    \chi_A(\sh{F}(n))=P(n)
  \]
  para todo $n\in\bb Z$.
  O polinômio $P$ chama-se o \emph{polinômio de Hilbert} de $\sh{F}$ em relação a $A$.

  \item[(ii) ] Para todo $n$ suficientemente grande,
  \[
    \chi_A(\sh{F}(n))
    =
    \operatorname{length}_A\Gamma(X,\sh{F}(n)).
  \]

  \item[(iii)] O coeficiente do termo de maior grau de $\chi_A(\sh{F}(n))$ não é negativo.
\end{enumerate}
\end{theorem}

Além disso, demonstraremos no Capítulo~IV, na seção dedicada à dimensão, que \emph{o grau de $\chi_A(\sh{F}(n))$ é igual à dimensão do suporte de $\sh{F}$}.

\begin{proof}
Para todo $n$ suficientemente grande, tem-se
\[
  \HH^i(X,\sh{F}(n))=0
\]
para todo $i>0$ \sref{III.2.2.1}.
\oldpage[III]{110}
Consequentemente,
\[
  \chi_A(\sh{F}(n))
  =
  \operatorname{length}_A\HH^0(X,\sh{F}(n))
  =
  \operatorname{length}_A\Gamma(X,\sh{F}(n))
\]
para todo $n$ suficientemente grande, o que demonstra \emph{(ii) }.
Segue-se que $\chi_A(\sh{F}(n))\geq0$ para todo $n$ suficientemente grande e, portanto, \emph{(iii) } deduz-se de \emph{(i) }.
A afirmação de unicidade de \emph{(i) } é imediata, por isso resta provar a existência de $P$.

Primeiro, vamos mostrar que podemos supor
\[
  \mathfrak{m}\sh{F}=0,
\]
onde $\mathfrak{m}$ é o ideal maximal de $A$.
Com efeito, existe um inteiro $s>0$ tal que $\mathfrak{m}^s=0$ e, portanto, $\sh{F}(n)$ possui a filtração finita
\[
  \sh{F}(n)
  \supset\mathfrak{m}\sh{F}(n)
  \supset\cdots
  \supset\mathfrak{m}^{s-1}\sh{F}(n)
  \supset0.
\]
Uma indução que usa \eqref{III.2.5.2.1} dá
\[
  \chi_A(\sh{F}(n))
  =
  \sum_{k=1}^{s}
  \chi_A\left(
    \mathfrak{m}^{k-1}\sh{F}(n)/
    \mathfrak{m}^{k}\sh{F}(n)
  \right).
\]
Mas
\[
  \mathfrak{m}^{k-1}\sh{F}(n)/
  \mathfrak{m}^{k}\sh{F}(n)
  =
  \sh{F}'_k(n),
  \qquad
  \sh{F}'_k=
  \mathfrak{m}^{k-1}\sh{F}/\mathfrak{m}^{k}\sh{F},
\]
o que demonstra a redução.

Então, suponhamos que $\mathfrak{m}\sh{F}=0$.
Seja $X'$ o subesquema fechado de $X$ obtido como imagem inversa, pelo morfismo estrutural
\[
  X\longrightarrow\Spec(A),
\]
do único ponto fechado de $\Spec(A)$, e seja $j:X'\to X$ a injeção canônica.
Então
\[
  \sh{F}=j_*(\sh{F}')
\]
para um $\sh{O}_{X'}$-módulo coerente $\sh{F}'$, e $X'$ é projetivo sobre $\Spec(K)$, onde
\[
  K=A/\mathfrak{m}.
\]
Se $\sh{L}'=j^*(\sh{L})$, então $\sh{L}'$ é muito amplo relativo a $\Spec(K)$ \sref[II]{II.4.4.10}, e
\[
  \sh{F}(n)=j_*(\sh{F}'(n)),
  \qquad
  \sh{F}'(n)
  =
  \sh{F}'\otimes_{\sh{O}_{X'}}\sh{L}'^{\otimes n}
\]
\sref[0]{0.5.4.10}.
Segue-se que
\[
  \chi_A(\sh{F}(n))
  =
  \chi_K(\sh{F}'(n))
\]
(G, II, 4.9.1), e ficamos reduzidos ao caso em que $A$ é um corpo.

Agora, $X$ pode ser considerado um subesquema fechado de
\[
  P=\mathbf{P}_A^r
\]
para um $r$ conveniente \sref[II]{II.5.5.4}[ii].
Se $i:X\to P$ é a injeção canônica, o mesmo raciocínio dá
\[
  \chi_A(\sh{F}(n))
  =
  \chi_A\bigl(i_*(\sh{F})(n)\bigr).
\]
Podemos, portanto, limitar-nos ao caso
\[
  X=\mathbf{P}_A^r=\Proj(S),
  \qquad
  S=A[T_0,\ldots,T_r],
\]
onde $A$ é um corpo.

Neste caso, $\sh{F}=\widetilde{M}$ para um $S$-módulo graduado $M$ de tipo finito \sref[II]{II.2.7.8}.
Pelo teorema das sicigias de Hilbert (M, VIII, 6.5), $M$ possui uma resolução finita por $S$-módulos livres graduados de tipo finito:
\[
  0\longrightarrow L_q
  \longrightarrow L_{q-1}
  \longrightarrow\cdots
  \longrightarrow L_1
  \longrightarrow M
  \longrightarrow0.
\]
Como o funtor $M\mapsto\widetilde{M}$ é exato \sref[II]{II.2.5.4}, obtemos uma sequência exata
\[
  0\longrightarrow\widetilde{L}_q
  \longrightarrow\widetilde{L}_{q-1}
  \longrightarrow\cdots
  \longrightarrow\widetilde{L}_1
  \longrightarrow\widetilde{M}
  \longrightarrow0,
\]
e, portanto, para todo $n\in\bb Z$, uma sequência exata
\[
  0\longrightarrow\widetilde{L}_q(n)
  \longrightarrow\widetilde{L}_{q-1}(n)
  \longrightarrow\cdots
  \longrightarrow\widetilde{L}_1(n)
  \longrightarrow\widetilde{M}(n)
  \longrightarrow0.
\]
Uma indução sobre $q$, usando \sref{III.2.5.2}, dá
\[
  \chi_A(\widetilde{M}(n))
  =
  \sum_{j=1}^{q}
  (-1)^{j+1}\chi_A(\widetilde{L}_j(n)).
\]
Para provar \emph{(i) }, ficamos, portanto, reduzidos ao caso em que $M$ é graduado, livre e de tipo finito e, portanto, ao caso
\[
  M=S(h)
\]
para algum $h\in\bb Z$.
Então
\[
  \widetilde{M}(n)
  =
  \widetilde{M(n)}
  =
  \widetilde{S(n+h)}
\]
\sref[II]{II.2.5.15}, então o teorema deduz-se do seguinte lema.
\end{proof}

\oldpage[III]{111}
\begin{lemma}[2.5.3.1]
\label{III.2.5.3.1}
Seja $A$ um corpo, seja $r>0$ um inteiro e seja
\[
  X=\mathbf{P}_A^r.
\]
Então
\[
  \chi_A(\sh{O}_X(n))
  =
  \binom{n+r}{r}
\]
para todo $n\in\bb Z$.
\end{lemma}

\begin{proof}
Para $n\geq0$,
\[
  \chi_A(\sh{O}_X(n))
  =
  \operatorname{length}_A\HH^0(X,\sh{O}_X(n)),
\]
que é o número de monômios nos $T_i$ de grau total $n$, ou seja,
\[
  \binom{n+r}{r}
\]
\sref{III.2.1.12}.

Para $n\leq-r-1$, tem-se analogamente
\[
  \chi_A(\sh{O}_X(n))
  =
  (-1)^r\operatorname{length}_A\HH^r(X,\sh{O}_X(n)).
\]
Escrevendo $n=-r-h$, a dimensão de $\HH^r(X,\sh{O}_X(n))$ sobre $A$ é o número de sequências
\[
  (p_i)_{0\leq i\leq r}
\]
de inteiros positivos que satisfaçam
\[
  \sum_{i=0}^{r}p_i=r+h
\]
\sref{III.2.1.12}.
Equivalentemente, é o número de sequências de inteiros não negativos
\[
  (q_i)_{0\leq i\leq r}
\]
tais que
\[
  \sum_{i=0}^{r}q_i=h-1.
\]
Este número é
\[
  \binom{h+r-1}{r}
  =
  (-1)^r\binom{n+r}{r}.
\]
Finalmente, se $-r\leq n<0$, então
\[
  \binom{n+r}{r}=0,
\]
e $\HH^i(X,\sh{O}_X(n))=0$ para todo $i\geq0$ \sref{III.2.1.12}.
Isto demonstra o lema.
\end{proof}

\begin{corollary}[2.5.4]
\label{III.2.5.4}
Seja $A$ um anel local artiniano, seja $S$ uma $A$-álgebra graduada de tipo finito gerada por $S_1$, seja $M$ um $S$-módulo graduado de tipo finito e seja
\[
  X=\Proj(S).
\]
Então, para todo $n$ suficientemente grande,
\[
  \chi_A(\widetilde{M}(n))
  =
  \operatorname{length}_A(M_n).
\]
\end{corollary}

\begin{proof}
Para todo $n$ suficientemente grande, os módulos
\[
  M_n
  \quad\text{e}\quad
  \Gamma(X,\widetilde{M}(n))
\]
são isomorfos \sref{III.2.3.1}.
\end{proof}


\subsection{Aplicação: criterios de amplitud}
\label{subsection:III.2.6}

\begin{proposition}[2.6.1]
\label{III.2.6.1}
Seja $Y$ um pré-esquema noetheriano, seja $f:X\to Y$ um morfismo próprio,
e seja $\sh{L}$ um $\sh{O}_X$-módulo invertível.
As seguintes condições são equivalentes:
\begin{enumerate}
  \item[\emph{a) }] $\sh{L}$ é amplo para $f$.
  \item[\emph{b)}] Para todo $\sh{O}_X$-módulo coerente $\sh{F}$, existe
  um inteiro $N$ tal que, para $n\geq N$,
  \[
    \RR^qf_*(\sh{F}\otimes\sh{L}^{\otimes n})=0
  \]
  para todo $q>0$.
  \item[\emph{c)}] Para todo ideal coerente $\sh{J}$ de $\sh{O}_X$, existe
  um inteiro $N$ tal que, para $n\geq N$,
  \[
    \RR^1f_*(\sh{J}\otimes\sh{L}^{\otimes n})=0.
  \]
\end{enumerate}
\end{proposition}

\begin{proof}
Vimos que \emph{a) } implica \emph{b) } \sref{III.2.2.1},~(ii).
É imediato que \emph{b) } implica \emph{c) }, e resta demonstrar
que \emph{c) } implica \emph{a) }.
Podemos limitar-nos ao caso em que $Y$ é afim \sref[II]{II.4.6.4}.
Basta mostrar que, quando $h$ percorre as seções de
$\sh{L}^{\otimes n}$ sobre $X$, com $n>0$, os subconjuntos abertos afins $X_h$
formam um recobrimento de $X$ \sref[II]{II.4.5.2}.

Para isso, seja $x$ um ponto fechado de $X$ e seja $U$ uma vizinhança
aberta afim de $x$.
Encontraremos um inteiro $n$ e uma seção
\[
  h\in\Gamma(X,\sh{L}^{\otimes n})
\]
tal que $x\in X_h\subset U$.
O subconjunto aberto $X_h$ será então afim \sref[I]{I.1.3.6}.
A união de todos esses $X_h$ é um subconjunto aberto de $X$ que contém todo
ponto fechado e, portanto, é todo $X$, já que $X$ é noetheriano
\sref[I]{I.6.3.7} e \sref[0]{0.2.1.3}.

Sejam $\sh{J}$ e $\sh{J}'$, respectivamente, os ideais quase-coerentes de
$\sh{O}_X$ que definem os subesquemas fechados reduzidos cujos espaços
subjacentes são $X-U$ e $(X-U)\cup\{x\}$, respectivamente \sref[I]{I.5.2.1}.
Os ideais $\sh{J}$ e $\sh{J}'$ são coerentes \sref[I]{I.6.1.1},
$\sh{J}'\subset\sh{J}$, e
\[
  \sh{J}''=\sh{J}/\sh{J}'
\]
é um $\sh{O}_X$-módulo coerente \sref[0]{0.5.3.3}, com suporte
$\{x\}$, e $\sh{J}''_x=\kres(x)$.
Como $\sh{L}$ é localmente livre, a sequência
\[
  0\longrightarrow
  \sh{J}'\otimes\sh{L}^{\otimes n}
  \longrightarrow
  \sh{J}\otimes\sh{L}^{\otimes n}
  \longrightarrow
  \sh{J}''\otimes\sh{L}^{\otimes n}
  \longrightarrow0
\]
é exata para todo $n$.
Por hipótese, para todo $n$ suficientemente grande,
\[
  \HH^1(X,\sh{J}'\otimes\sh{L}^{\otimes n})=0.
\]

\oldpage[III]{112}

A sequência exata de cohomologia mostra, portanto, que
\[
  \Gamma(X,\sh{J}\otimes\sh{L}^{\otimes n})
  \longrightarrow
  \Gamma(X,\sh{J}''\otimes\sh{L}^{\otimes n})
\]
é sobrejetivo.
Uma seção $g$ de $\sh{J}''\otimes\sh{L}^{\otimes n}$ sobre $X$ tal que
$g(x)\neq0$ é, portanto, a imagem de uma seção
\[
  h\in\Gamma(X,\sh{J}\otimes\sh{L}^{\otimes n})
  \subset\Gamma(X,\sh{L}^{\otimes n});
\]
a inclusão deduz-se de \sref[0]{0.5.4.1}.
Por construção, $h(x)\neq0$, enquanto $h(z)=0$ para todo $z\notin U$.
Assim, $x\in X_h\subset U$, o que prova a afirmação.
\end{proof}

\begin{proposition}[2.6.2]
\label{III.2.6.2}
Seja $Y$ um pré-esquema noetheriano, seja $f:X\to Y$ um morfismo de
tipo finito, seja $g:X'\to X$ um morfismo finito sobrejetivo, seja
$\sh{L}$ um $\sh{O}_X$-módulo invertível, e ponhamos
\[
  \sh{L}'=g^*(\sh{L}).
\]
Suponhamos que existe um subconjunto $Z$ de $X$, próprio sobre $Y$
\sref[II]{II.5.4.10}, tal que, para todo $x\in X-Z$, ou $X$ é
normal em $x$, ou
\[
  \bigl(g_*(\sh{O}_{X'})\bigr)_x
\]
é um $\sh{O}_{X,x}$-módulo livre.
Então $\sh{L}$ é amplo para $f$ se e apenas se $\sh{L}'$ é amplo para
$f\circ g$.
\end{proposition}

\begin{proof}
\begin{env}[2.6.2.1]
\label{III.2.6.2.1}
Como $g$ é afim, a condição é necessária \sref[II]{II.5.1.12}.
Para provar que é suficiente, podemos supor que $Y$ é afim.
\sref[II]{II.4.6.4}.
Primeiro vamos mostrar que também podemos limitar-nos ao caso em que $X$ é
reducido.

Seja $j:X_{\mathrm{red}}\to X$ a imersão canônica, e ponhamos
\[
  X_1=X_{\mathrm{red}},
  \qquad
  X'_1=X'\times_X X_1.
\]
Então temos o diagrama comutativo
\[
  \xymatrix{
    X'\ar[d]_g &
    X'_1\ar[l]_{j'}\ar[d]^{g_1}\\
    X &
    X_1\ar[l]^{j}
  }
  \label{III.2.6.2.2}\tag{2.6.2.2}
\]
O morfismo $f\circ j$ é do tipo finito \sref[I]{I.6.3.4}, e
$g_1$ é finito \sref[II]{II.6.1.5},~(iii).
Se $\sh{L}'$ é amplo para $f\circ g$, então
$j'^*(\sh{L}')$ é amplo para $f\circ g\circ j'$, porque $j'$ é uma
imersão fechada \sref[II]{II.5.1.12} e \sref[I]{I.4.3.2}.
Se $Z_1=j^{-1}(Z)$, então $Z_1$ é próprio sobre $Y$
\sref[II]{II.5.4.10}.
Além disso, se $X$ é normal em um ponto $x$, também o é $X_{\mathrm{red}}$
nesse ponto; e, se
$\bigl(g_*(\sh{O}_{X'})\bigr)_x$ é um $\sh{O}_{X,x}$-módulo livre,
então \sref[II]{II.1.5.2} mostra que
\[
  \bigl((g_1)_*(\sh{O}_{X'_1})\bigr)_x
\]
é um $\sh{O}_{X_1,x}$-módulo livre.
Finalmente, como $X$ é noetheriano \sref[I]{I.6.3.7}, a amplitude de
$j^*(\sh{L})$ implica a de $\sh{L}$ \sref[II]{II.4.5.14}.
Como
\[
  j'^*(\sh{L}')=g_1^*(j^*(\sh{L})),
\]
Isto demonstra a redução anunciada.
Em seguida, assumimos que $Y$ é afim e que $X$ é reduzido.
\end{env}

As hipóteses de \sref[II]{II.6.6.11} são agora satisfeitas.
Existem um $Y$-pré-esquema reduzido $X_2$ e um
$Y$-morfismo finito birracional
\[
  h:X_2\to X
\]
cuja restrição a $h^{-1}(X-Z)$ é um isomorfismo sobre $X-Z$, e
tal que $h^*(\sh{L})$ é amplo.
Substituindo $X'$ por $X_2$, ficamos reduzidos a demonstrar a proposição sob
a hipótese adicional de que $g$ possui as propriedades que acabamos de listar para
$h$.
Seguimos denotando por $Z$ um subesquema fechado de $X$, próprio sobre
$Y$ \sref[II]{II.5.4.10}, cujo espaço subjacente é o subconjunto dado
$Z$.

\begin{env}[2.6.2.3]
\label{III.2.6.2.3}
Seja $X_1$ um subpreesquema fechado de $X$, seja $j:X_1\to X$ a
imersão canônica, e seja
\[
  X'_1=g^{-1}(X_1)=X'\times_X X_1,
\]
com imersão canônica $j':X'_1\to X'$.
O quadrado correspondente é novamente \eqref{III.2.6.2.2}.
Ponhamos
\[
  \sh{L}_1=j^*(\sh{L}),
  \qquad
  \sh{L}'_1=j'^*(\sh{L}')=g_1^*(\sh{L}_1).
\]
Então $\sh{L}'_1$ é amplo para $f\circ g\circ j'$
\sref[II]{II.5.1.12}.
Se $Z_1=j^{-1}(Z)$, o subpré-esquema fechado $Z_1$ de $X_1$ é próprio
sobre $Y$ \sref[II]{II.5.4.2},~(ii).
Em outras palavras, as hipóteses de \sref{III.2.6.2} verificam-se para
$X_1$, $\sh{L}_1$, $g_1$ e $Z_1$.
\end{env}

\oldpage[III]{113}

Isso nos permite demonstrar \sref{III.2.6.2} por indução noetheriana
\sref[0]{0.2.2.2} Quando a restrição de $g$ a $g^{-1}(X-Z)$ é um
isomorfismo sobre $X-Z$; como foi visto em \sref{III.2.6.2.1}, esse caso é
suficiente.
Basta mostrar que, se a conclusão de \sref{III.2.6.2} valer
para $\sh{L}_1$ sobre todo subpré-esquema fechado $X_1$ de $X$ cujo
o espaço subjacente não é $X$, então também vale para $\sh{L}$.

\begin{env}[2.6.2.4]
\label{III.2.6.2.4}
Ponhamos
\[
  \sh{A}=\sh{O}_X,
  \qquad
  \sh{B}=g_*(\sh{O}_{X'}).
\]
Como $g$ é birracional e $X$ é reduzido, $\sh{B}$ é uma
sub-$\sh{A}$-álgebra do feixe $\sh{R}(X)$ de anéis totais de frações,
e é um $\sh{A}$-módulo coerente.
Além disso,
\[
  \sh{B}|_{X-Z}=\sh{A}|_{X-Z}.
\]
Seja $\sh{K}$ o condutor de $\sh{B}$ sobre $\sh{A}$, ou seja, o maior
sub-$\sh{A}$-módulo quase-coerente de $\sh{A}$ tal que
\[
  \sh{B}\sh{K}\subset\sh{A}.
\]
Equivalentemente, $\sh{K}$ é o anulador do
$\sh{A}$-módulo $\sh{B}/\sh{A}$ \sref[0]{0.5.3.7}.
Segue-se que
\[
  \sh{B}\sh{K}=\sh{K}.
\]
É claro que $\sh{K}_x=\sh{A}_x$ em todo ponto que admita uma vizinhança
$W_x$ sobre o qual $g:g^{-1}(W_x)\to W_x$ é um isomorfismo.
Isto ocorre, em particular, em todo ponto de $X-Z$, e em uma
vizinhança de todo ponto genérico de uma componente irredutível de
$X$.

Seja
\[
  Z_1=\Spec(\sh{A}/\sh{K})
\]
o subpré-esquema fechado de $X$ definido por $\sh{K}$.
É próprio sobre $Y$, porque seu espaço subjacente é fechado em $Z$
\sref[II]{II.5.4.10}.
A definição de $\sh{K}$ também dá
\[
  \sh{B}|_{X-Z_1}=\sh{A}|_{X-Z_1}.
\]
Podemos, portanto, reduzir-nos ao caso
\[
  Z=\Spec(\sh{A}/\sh{K}).
\]
Como $X-Z_1$ é um subconjunto aberto não vazio de $X$, podemos também
supor que o espaço subjacente de $Z$ seja diferente de $X$.
\end{env}

\begin{env}[2.6.2.5]
\label{III.2.6.2.5}
Consideremos $X'$ como $\Spec(\sh{B})$, já que $g$ é afim, e seja
\[
  \sh{K}'=\widetilde{\sh{K}}.
\]
Este é um ideal coerente de $\sh{O}_{X'}$ que satisfaz
$g_*(\sh{K}')=\sh{K}$ \sref[II]{II.1.4.1}.
O subpré-esquema fechado
\[
  Z'=g^{-1}(Z)=Z\times_X X'
\]
de $X'$ é definido por $\sh{K}'$ e é igual a
\[
  \Spec(\sh{B}/\sh{K})
\]
\sref[II]{II.1.4.10}.
O morfismo $h:Z'\to Z$ é finito \sref[II]{II.6.1.5},~(iii);
portanto, $Z'$ é próprio sobre $Y$ \sref[II]{II.6.1.11} e
\sref[II]{II.5.4.2},~(ii).

Devemos demonstrar que, para todo $x\in X$ e toda vizinhança aberta $U$
de $x$, existem um inteiro $n>0$ e uma seção $s$ de
$\sh{L}^{\otimes n}$ sobre $X$ tais que
\[
  x\in X_s\subset U
\]
\sref[II]{II.4.5.2}.
Distinguimos dos casos.

\smallskip
\noindent
$1^\circ$ Vamos primeiro supor que $x\in X-Z$.
Podemos, evidentemente, supor que $U\subset X-Z$, de modo que
$U'=g^{-1}(U)$ não intersecta $Z'$.
Como $\sh{L}'$ é amplo, existem $n>0$ e uma seção $s'$ de
$\sh{L}'^{\otimes n}$ sobre $X'$ tais que
\[
  x'=g^{-1}(x)\in X'_{s'}\subset g^{-1}(U)
\]
\sref[II]{II.4.5.2}.
Também podemos supor que
$\sh{K}'\otimes\sh{L}'^{\otimes n}$ é gerado por suas seções globais
\sref[II]{II.4.5.5}.
Como $\sh{K}'_{x'}=\sh{O}_{X',x'}$, uma destas seções, digamos $s''$, é
não nula em $x'$.
Multiplicando por $s'$ e substituindo assim $n$ por $2n$, podemos obter
que
\[
  x'\in X'_{s''}\subset g^{-1}(U).
\]
Por \sref[0]{0.5.4.10}, existe um isomorfismo canônico
\[
  \Gamma(X,\sh{K}\otimes\sh{L}^{\otimes n})
  \xrightarrow{\sim}
  \Gamma(X',\sh{K}'\otimes\sh{L}'^{\otimes n}).
\]
A seção $s$ de $\sh{K}\otimes\sh{L}^{\otimes n}$ correspondente a
$s''$ possui manifestamente as propriedades exigidas.

\smallskip
\noindent
$2^\circ$ Agora vamos supor que $x\in Z$.
Seja $\sh{J}$ o ideal coerente de $\sh{O}_X$ que define o
subpreesquema fechado reduzido cujo espaço subjacente é $X-U$, e consideremos em
$\sh{B}$ os ideais coerentes

\oldpage[III]{114}

\[
  \sh{J}\sh{B}
  \qquad\text{e}\qquad
  \sh{J}\sh{K}
  =
  \sh{J}(\sh{K}\sh{B})
  =
  \sh{K}(\sh{J}\sh{B}).
\]
Temos o diagrama de inclusões
\[
  \xymatrix@R=1.5em{
    \sh{J}\sh{B}\ar@{^{(}->}[r] &
    \sh{B}\\
    \sh{J}\ar@{^{(}->}[r]\ar@{^{(}->}[u] &
    \sh{A}\ar@{^{(}->}[u]\\
    \sh{J}\sh{K}\sh{B}=\sh{J}\sh{K}
      \ar@{^{(}->}[r]\ar@{^{(}->}[u] &
    \sh{K}\ar@{^{(}->}[u]
  }
  \label{III.2.6.2.6}\tag{2.6.2.6}
\]
Seja $\sh{J}'=(\sh{J}\sh{B})^\sim$ o ideal coerente de
$\sh{O}_{X'}$ tal que
\[
  g_*(\sh{J}')=\sh{J}\sh{B}.
\]
Então
\[
  \sh{J}'\sh{K}'=(\sh{J}\sh{K}\sh{B})^\sim,
  \qquad
  \sh{J}'/\sh{J}'\sh{K}'
  =
  (\sh{J}\sh{B}/\sh{J}\sh{K}\sh{B})^\sim
\]
\sref[II]{II.1.4.4}.
Como $\sh{J}|_V=\sh{J}\sh{K}|_V$ sobre todo subconjunto aberto $V$ disjunto
de $Z$, o suporte de $\sh{J}'/\sh{J}'\sh{K}'$ está contido em
$Z'$.
Como $Z'$ é próprio sobre $Y$, \sref{III.2.2.4} mostra que, para todo
$n$ suficientemente grande, a aplicação canônica
\[
  \Gamma(X',\sh{J}'\otimes\sh{L}'^{\otimes n})
  \longrightarrow
  \Gamma\bigl(
    X',
    (\sh{J}'/\sh{J}'\sh{K}')\otimes\sh{L}'^{\otimes n}
  \bigr)
\]
é sobrejetiva.
Por \sref[0]{0.5.4.10}, segue-se que a aplicação canônica
\[
  \Gamma(X,\sh{J}\sh{B}\otimes\sh{L}^{\otimes n})
  \longrightarrow
  \Gamma\bigl(
    X,
    (\sh{J}\sh{B}/\sh{J}\sh{K}\sh{B})
      \otimes\sh{L}^{\otimes n}
  \bigr)
\]
é sobrejetiva.

Sejam $i:Z\to X$ e $i':Z'\to X'$ as imersões canônicas.
São parte do diagrama comutativo
\[
  \xymatrix{
    X'\ar[d]_g &
    Z'\ar[l]_{i'}\ar[d]^h\\
    X &
    Z\ar[l]^{i}
  }
\]
Ponhamos
\[
  \sh{M}=i^*(\sh{L}),
  \qquad
  \sh{M}'=i'^*(\sh{L}').
\]
Como $\sh{L}'$ é amplo, $\sh{M}'$ é amplo.
\sref[II]{II.5.1.12}, e
\[
  \sh{M}'=h^*(\sh{M}).
\]
A hipótese de indução noetheriana, aplicada pois $Z\neq X$, mostra agora
que $\sh{M}$ é amplo.
Por conseguinte, para todo $n$ suficientemente grande,
\[
  i^*(\sh{J}/\sh{J}\sh{K})\otimes\sh{M}^{\otimes n}
\]
é gerado por suas seções sobre $Z$ \sref[II]{II.4.5.5}.
Como
\[
  \sh{J}/\sh{J}\sh{K}
  =
  i_*\bigl(i^*(\sh{J}/\sh{J}\sh{K})\bigr),
\]
\sref[0]{0.5.4.10} dá uma seção $s$ de
\[
  (\sh{J}/\sh{J}\sh{K})\otimes\sh{L}^{\otimes n}
\]
sobre $X$, para algum $n$ suficientemente grande, tal que $s(x)\neq0$.
De fato, $\sh{J}_x=\sh{A}_x$ pela definição de $\sh{J}$, enquanto
$\sh{K}_x\neq\sh{A}_x$ por hipótese.

O diagrama \eqref{III.2.6.2.6} mostra que $s$ é também uma seção de
\[
  (\sh{J}\sh{B}/\sh{J}\sh{K}\sh{B})
  \otimes\sh{L}^{\otimes n}
\]
sobre $X$.
É, portanto, a imagem canônica de uma seção
\[
  t\in\Gamma(X,\sh{J}\sh{B}\otimes\sh{L}^{\otimes n}).
\]
Por definição, a imagem de $t$ módulo
$\sh{J}\sh{K}\sh{B}\otimes\sh{L}^{\otimes n}$ pertence a
\[
  \frac{\sh{J}\otimes\sh{L}^{\otimes n}}
       {\sh{J}\sh{K}\sh{B}\otimes\sh{L}^{\otimes n}}.
\]
O diagrama \eqref{III.2.6.2.6} implica, portanto, que $t$ é uma seção
de $\sh{J}\otimes\sh{L}^{\otimes n}$ e, em particular, uma seção de
$\sh{L}^{\otimes n}$.
Temos $t(x)\neq0$, logo $x\in X_t$.
Pela definição de $\sh{J}$, a seção $t$ anula-se sobre $X-U$, o
suporte de $\sh{O}_X/\sh{J}$.
Assim, $X_t\subset U$, e a demonstração está terminada.
\end{env}
\end{proof}

\oldpage[III]{115}

\begin{env}[2.6.3]
\label{III.2.6.3}
\emph{Observação.}
Quando $X$ é próprio sobre $Y$, a Proposição~\sref{III.2.6.2} pode ser demonstrada
mais simplesmente raciocinando como no teorema de Chevalley
\sref[II]{II.6.7.1}, com a ajuda da Proposição~\sref{III.2.6.1} e
do lema \sref[II]{II.6.7.1.1}.
\end{env}

\section{Teorema de finitude para os morfismos próprios}
\label{section:III.3}

\subsection{O lema de d\'evissage}
\label{subsection:III.3.1}

\begin{definition}[3.1.1]
\label{III.3.1.1}
Seja $\C$ uma categoria abeliana.
Diremos que um subconjunto $\C'$ do conjunto de objetos de $\C$ é \emph{exato} se $0\in\C'$ e se, para toda sucessão exata $0\to A'\to A\to A''\to 0$ em $\C$ tal que dois dos objetos $A$, $A'$, $A''$ pertencem a $\C'$, o terceiro pertence também a $\C'$.
\end{definition}

\begin{theorem}[3.1.2]
\label{III.3.1.2}
Seja $X$ um pré-esquema noetheriano; denotemos por $\C$ a categoria abeliana dos $\sh{O}_X$-módulos coerentes.
Seja $\C'$ um subconjunto exato de $\C$, e seja $X'$ um subconjunto fechado do espaço subjacente a $X$.
Suponhamos que, para todo subconjunto fechado irredutível $Y$ de $X'$, de ponto genérico $y$, existe um $\sh{O}_X$-módulo $\sh{G}\in\C'$ tal que $\sh{G}_y$ é um espaço vetorial sobre $\kres(y)$ de dimensão~$1$.
Então todo $\sh{O}_X$-módulo coerente cujo suporte esteja contido em $X'$ pertence a $\C'$ (e, em particular, se $X'=X$, tem-se $\C'=\C$).
\end{theorem}

\begin{proof}
Consideremos a seguinte propriedade $\textbf{P}(Y)$ de um subconjunto fechado $Y$ de $X'$: todo $\sh{O}_X$-módulo coerente cujo suporte esteja contido em $Y$ pertence a $\C'$.
Em virtude do princípio de indução noetheriana \sref[0]{0.2.2.2}, vemos que basta demonstrar que \emph{se $Y$ é um subconjunto fechado de $X'$ tal que a propriedade $\textbf{P}(Y')$ é verdadeira para todo subconjunto fechado $Y'$ de $Y$, distinto de $Y$, então $\textbf{P}(Y)$ é verdadeira}.

Seja, pois, $\sh{F}\in\C$ de suporte contido em $Y$, e demonstremos que $\sh{F}\in\C'$.
Denotemos também por $Y$ o subpré-esquema fechado reduzido de $X$ cujo espaço subjacente é $Y$ \sref[I]{I.5.2.1}; está definido por um feixe coerente de ideais $\sh{J}$ de $\sh{O}_X$.
Sabemos \sref[I]{I.9.3.4} que existe um inteiro $n>0$ tal que $\sh{J}^n\sh{F}=0$; para $1\leq k\leq n$ tem-se, portanto, uma sucessão exata
\[
  0\to\sh{J}^{k-1}\sh{F}/\sh{J}^k\sh{F}\to\sh{F}/\sh{J}^k\sh{F}\to\sh{F}/\sh{J}^{k-1}\sh{F}\to 0
\]
 de $\sh{O}_X$-módulos coerentes (\sref[I]{I.5.3.6} e \sref[I]{I.5.3.3}); como $\C'$ é exato, uma indução sobre $k$ mostra que basta provar que cada $\sh{F}_k=\sh{J}^{k-1}\sh{F}/\sh{J}^k\sh{F}$ pertence a $\C'$.
Reduzimo-nos assim a provar que $\sh{F}\in\C'$ sob a hipótese adicional $\sh{J}\sh{F}=0$; isto equivale a dizer que $\sh{F}=j_*(j^*(\sh{F}))$, onde $j$ é a imersão canônica $Y\to X$.
Consideremos agora dois casos:
\begin{enumerate}
  \item[(a)] $Y$ é \emph{redutível}.
    Escrevamos $Y=Y'\cup Y''$, onde $Y'$ e $Y''$ são subconjuntos fechados de $Y$, distintos de $Y$; denotemos também por $Y'$ e $Y''$ os subpré-esquemas fechados reduzidos de $X$ cujos espaços subjacentes são, respectivamente, $Y'$ e $Y''$, definidos respectivamente pelos feixes de ideais $\sh{J}'$ e $\sh{J}''$ de $\sh{O}_X$.
    Ponhamos $\sh{F}'=\sh{F}\otimes_{\sh{O}_X}(\sh{O}_X/\sh{J}')$ e $\sh{F}''=\sh{F}\otimes_{\sh{O}_X}(\sh{O}_X/\sh{J}'')$.
    Os homomorfismos canônicos $\sh{F}\to\sh{F}'$ e $\sh{F}\to\sh{F}''$ definem assim um homomorfismo $u:\sh{F}\to\sh{F}'\oplus\sh{F}''$.
    Demonstremos que, para todo $z\not\in Y'\cap Y''$, o homomorfismo $u_z:\sh{F}_z\to\sh{F}_z'\oplus\sh{F}_z''$ é \emph{bijetivo}.
    Com efeito, tem-se $\sh{J}'\cap\sh{J}''=\sh{J}$, pois a questão é local e
\oldpage[III]{116}
    a igualdade anterior deduz-se de (\sref[I]{I.5.2.1} e \sref[I]{I.1.1.5}); se $z\not\in Y''$, tem-se $\sh{J}_z'=\sh{J}_z$, de onde $\sh{F}_z'=\sh{F}_z$ e $\sh{F}_z''=0$, o que prova nossa afirmação neste caso; raciocina-se do mesmo modo se $z\not\in Y'$.
    Por conseguinte, o núcleo e o conúcleo de $u$, que pertencem a $\C$ \sref[0]{0.5.3.4}, têm seu suporte em $Y'\cap Y''$ e, por hipótese, pertencem a $\C'$; pela mesma razão, $\sh{F}'$ e $\sh{F}''$ pertencem a $\C'$, e portanto também $\sh{F}'\oplus\sh{F}''$, pois $\C'$ é exato.
    A conclusão obtém-se então considerando as duas sucessões exatas
    \[
      0\to\Im u\to\sh{F}'\oplus\sh{F}''\to\Coker u\to 0,
    \]
    \[
      0\to\Ker u\to\sh{F}\to\Im u\to 0,
    \]
    e a hipótese de que $\C'$ é exato.
  \item[(b)] $Y$ é irredutível e, por conseguinte, o subpré-esquema $Y$ de $X$ é \emph{íntegro}.
    Se $y$ é seu ponto genérico, tem-se $(\sh{O}_Y)_y=\kres(y)$ e, como $j^*(\sh{F})$ é um $\sh{O}_Y$-módulo coerente, $\sh{F}_y=(j^*(\sh{F}))_y$ é um espaço vetorial sobre $\kres(y)$ de dimensão finita~$m$.
    Por hipótese existe um $\sh{O}_X$-módulo coerente $\sh{G}\in\C'$ (necessariamente de suporte $Y$) tal que $\sh{G}_y$ é um espaço vetorial sobre $\kres(y)$ de dimensão~$1$.
    Existe, portanto, um $\kres(y)$-isomorfismo $(\sh{G}_y)^m\isoto\sh{F}_y$, que é também um $\sh{O}_Y$-isomorfismo; como $\sh{G}^m$ e $\sh{F}$ são coerentes, existe uma vizinhança aberta $W$ de $y$ em $X$ e um isomorfismo $\sh{G}^m|W\isoto\sh{F}|W$ \sref[0]{0.5.2.7}.
    Seja $\sh{H}$ o grafo deste isomorfismo, que é um $(\sh{O}_X|W)$-submódulo coerente de $(\sh{G}^m\oplus\sh{F})|W$, canonicamente isomorfo a $\sh{G}^m|W$ e a $\sh{F}|W$; existe então um $\sh{O}_X$-submódulo coerente $\sh{H}_0$ de $\sh{G}^m\oplus\sh{F}$ que induz $\sh{H}$ sobre $W$ e $0$ sobre $X\setmin Y$, dado que $\sh{G}^m$ e $\sh{F}$ têm suporte $Y$~\sref[I]{I.9.4.7}.
    As restrições $v:\sh{H}_0\to\sh{G}^m$ e $w:\sh{H}_0\to\sh{F}$ das projeções canônicas de $\sh{G}^m\oplus\sh{F}$ são então homomorfismos de $\sh{O}_X$-módulos coerentes que, sobre $W$ e sobre $X\setmin Y$, se reduzem a isomorfismos; em outras palavras, os núcleos e conúcleos de $v$ e $w$ têm seu suporte no fechado $Y\setmin(Y\cap W)$, distinto de $Y$.
    Pertencem a $\C'$; por outro lado, $\sh{G}^m\in\C'$, pois $\sh{G}\in\C'$ e $\C'$ é exato.
    Da exatidão de $\C'$ concluímos sucessivamente que $\sh{H}_0\in\C'$ e logo que $\sh{F}\in\C'$.
C.Q.D.
\end{enumerate}
\end{proof}

\begin{corollary}[3.1.3]
\label{III.3.1.3}
Suponhamos além disso que o subconjunto exato $\C'$ de $\C$ tem a propriedade de que todo fator direto coerente de um $\sh{O}_X$-módulo coerente $\sh{M}\in\C'$ pertence também a $\C'$.
Neste caso, a conclusão do Teorema~\sref{III.3.1.2} continua válida quando a condição «$\sh{G}_y$ é um espaço vetorial sobre $\kres(y)$ de dimensão~$1$» é substituída por $\sh{G}_y\neq 0$ (o que equivale a $\Supp(\sh{G})=Y$).
\end{corollary}

\begin{proof}
O raciocínio do Teorema~\sref{III.3.1.2} somente deve modificar-se no caso~(b); agora $\sh{G}_y$ é um espaço vetorial sobre $\kres(y)$ de dimensão $q>0$ e, portanto, tem-se um $\sh{O}_Y$-isomorfismo $(\sh{G}_y)^m\isoto(\sh{F}_y)^q$; o final do raciocínio do Teorema~\sref{III.3.1.2} prova então que $\sh{F}^q\in\C'$, e a hipótese adicional sobre $\C'$ implica que $\sh{F}\in\C'$.
\end{proof}

\subsection{O teorema de finitude: caso dos esquemas ordinários}
\label{subsection:III.3.2}
\label{III.3.2}

\begin{theorem}[3.2.1]
\label{III.3.2.1}
Seja $Y$ um pré-esquema localmente noetheriano e seja $f:X\to Y$ um morfismo próprio.
Para todo $\sh{O}_X$-módulo coerente $\sh{F}$, os $\sh{O}_Y$-módulos $\RR^q f_*(\sh{F})$ são coerentes para $q\geq 0$.
\end{theorem}

\begin{proof}
Como a questão é local sobre $Y$, podemos supor que $Y$ é noetheriano, e portanto que $X$ é noetheriano~\sref[I]{I.6.3.7}.
Os $\sh{O}_X$-módulos coerentes $\sh{F}$ para os quais é válida a conclusão do Teorema~\sref{III.3.2.1} formam um subconjunto \emph{exato} $\C'$ da categoria $\C$ dos $\sh{O}_X$-módulos coerentes.
\oldpage[III]{117}
Com efeito, seja $0\to\sh{F}'\to\sh{F}\to\sh{F}''\to 0$ uma sucessão exata de $\sh{O}_X$-módulos coerentes; suponhamos, por exemplo, que $\sh{F}'$ e $\sh{F}''$ pertencem a $\C'$; tem-se a sucessão exata longa de cohomologia
\[
  \RR^{q-1}f_*(\sh{F}'')\xrightarrow{\partial}\RR^q f_*(\sh{F}')\to\RR^q f_*(\sh{F})\to\RR^q f_*(\sh{F}'')\xrightarrow{\partial}\RR^{q+1}f_*(\sh{F}'),
\]
na qual, por hipótese, os quatro termos exteriores são coerentes; o mesmo ocorre com o termo central $\RR^q f_*(\sh{F})$ por (\sref[0]{0.5.3.4}~e~\sref[0]{0.5.3.3}).
Demonstra-se do mesmo modo que, quando $\sh{F}$ e $\sh{F}'$ (resp.~$\sh{F}$ e $\sh{F}''$) pertencem a $\C'$, também pertence $\sh{F}''$ (resp.~$\sh{F}'$).
Além disso, todo \emph{fator direto} coerente $\sh{F}'$ de um $\sh{O}_X$-módulo $\sh{F}\in\C'$ pertence a $\C'$: com efeito, $\RR^q f_*(\sh{F}')$ é então um fator direto de $\RR^q f_*(\sh{F})$~(G,~II,~4.4.4), pelo que é de tipo finito; como é quase-coerente~\sref{III.1.4.10}, é coerente, pois $Y$ é noetheriano.
Em virtude do Corolário~\sref{III.3.1.3}, reduzimo-nos a provar que, quando $X$ é \emph{irredutível} de ponto genérico $x$, existe \emph{um} $\sh{O}_X$-módulo coerente $\sh{F}$ pertencente a $\C'$ tal que $\sh{F}_x\neq 0$: com efeito, uma vez estabelecido este ponto, pode aplicar-se a qualquer subpré-esquema fechado irredutível $Y$ de $X$, pois se $j:Y\to X$ é a imersão canônica, $f\circ j$ é próprio~\sref[II]{II.5.4.2}, e se $\sh{G}$ é um $\sh{O}_Y$-módulo coerente de suporte $Y$, então $j_*(\sh{G})$ é um $\sh{O}_X$-módulo coerente tal que $\RR^q(f\circ j)_*(\sh{G})=\RR^q f_*(j_*(\sh{G}))$~(G,~II,~4.9.1); podemos, portanto, aplicar o Corolário~\sref{III.3.1.3}.

Em virtude do lema de Chow~\sref[II]{II.5.6.2}, existem um pré-esquema irredutível $X'$ e um morfismo \emph{projetivo} e sobrejetivo $g:X'\to X$ tais que $f\circ g:X'\to Y$ é \emph{projetivo}.
Existe um $\sh{O}_{X'}$-módulo $\sh{L}$ amplo para $g$~\sref[II]{II.5.3.1}; aplicamos o teorema fundamental dos morfismos projetivos~\sref{III.2.2.1} a $g:X'\to X$ e a $\sh{L}$: existe então um inteiro $n$ tal que $\sh{F}=g_*(\sh{O}_{X'}(n))$ é um $\sh{O}_X$-módulo coerente e $\RR^q g_*(\sh{O}_{X'}(n))=0$ para todo $q>0$; além disso, como $g^*(g_*(\sh{O}_{X'}(n)))\to\sh{O}_{X'}(n)$ é sobrejetivo para $n$ suficientemente grande~\sref{III.2.2.1}, vemos que podemos supor, no ponto genérico $x$ de $X$, que $\sh{F}_x\neq 0$~\sref[II]{II.3.4.7}.
Por outro lado, como $f\circ g$ é projetivo e $Y$ é noetheriano, os $\RR^q(f\circ g)_*(\sh{O}_{X'}(n))$ são \emph{coerentes}~\sref{III.2.2.1}.
Ora, $\RR^\bullet(f\circ g)_*(\sh{O}_{X'}(n))$ é o termo de convergência de uma sucessão espectral de Leray cujo termo $E_2$ é dado por $E_2^{pq}=\RR^p f_*(\RR^q g_*(\sh{O}_{X'}(n)))$; o anterior mostra que esta sucessão espectral degenera, e sabemos então~\sref[0]{0.11.1.6} que $E_2^{p0}=\RR^p f_*(\sh{F})$ é isomorfo a $\RR^p(f\circ g)_*(\sh{O}_{X'}(n))$, o que termina a demonstração.
\end{proof}

\begin{corollary}[3.2.2]
\label{III.3.2.2}
Seja $Y$ um pré-esquema localmente noetheriano.
Para todo morfismo próprio $f:X\to Y$, a imagem direta por $f$ de qualquer $\sh{O}_X$-módulo coerente é um $\sh{O}_Y$-módulo coerente.
\end{corollary}

\begin{corollary}[3.2.3]
\label{III.3.2.3}
Seja $A$ um anel noetheriano e seja $X$ um esquema próprio sobre $A$; para todo $\sh{O}_X$-módulo coerente $\sh{F}$, os $\HH^p(X,\sh{F})$ são $A$-módulos de tipo finito, e existe um inteiro $r>0$ tal que, para todo $\sh{O}_X$-módulo coerente $\sh{F}$ e todo $p>r$, $\HH^p(X,\sh{F})=0$.
\end{corollary}

\begin{proof}
A segunda afirmação já se demonstrou em~\sref{III.1.4.12}; a primeira deduz-se do teorema de finitude~\sref{III.3.2.1}, tendo em conta o Corolário~\sref{III.1.4.11}.
\end{proof}

Em particular, se $X$ é um \emph{esquema algébrico próprio} sobre um corpo $k$, então, para todo $\sh{O}_X$-módulo coerente $\sh{F}$, os $\HH^p(X,\sh{F})$ são espaços vetoriais sobre $k$ de \emph{dimensão finita}.

\begin{corollary}[3.2.4]
\label{III.3.2.4}
Seja $Y$ um pré-esquema localmente noetheriano e seja $f:X\to Y$ um morfismo de tipo finito.
Para todo $\sh{O}_X$-módulo coerente $\sh{F}$ cujo suporte seja próprio sobre $Y$~\sref[II]{II.5.4.10}, os $\sh{O}_Y$-módulos $\RR^q f_*(\sh{F})$ são coerentes.
\end{corollary}

\oldpage[III]{118}
\begin{proof}
Como a questão é local sobre $Y$, podemos supor que $Y$ é noetheriano, e o mesmo ocorre então com $X$~\sref[I]{I.6.3.7}.
Por hipótese, todo subpré-esquema fechado $Z$ de $X$ cujo espaço subjacente seja $\Supp(\sh{F})$ é próprio sobre $Y$; em outras palavras, se $j:Z\to X$ é a imersão canônica, então $f\circ j:Z\to Y$ é próprio.
Podemos supor que $Z$ é tal que $\sh{F}=j_*(\sh{G})$, onde $\sh{G}=j^*(\sh{F})$ é um $\sh{O}_Z$-módulo coerente~\sref[I]{I.9.3.5}; como $\RR^q f_*(\sh{F})=\RR^q(f\circ j)_*(\sh{G})$ pelo Corolário~\sref{III.1.3.4}, a conclusão deduz-se imediatamente do Teorema~\sref{III.3.2.1}.
\end{proof}

\subsection{Generalização do teorema de finitude (esquemas ordinários)}
\label{subsection:III.3.3}

\begin{proposition}[3.3.1]
\label{III.3.3.1}
Seja $Y$ um pré-esquema noetheriano, seja $\sh{S}$ uma $\sh{O}_Y$-álgebra quase-coerente de tipo finito, graduada em graus positivos, seja $Y'=\Proj(\sh{S})$, e seja $g:Y'\to Y$ o morfismo estrutural.
Seja $f:X\to Y$ um morfismo próprio, seja $\sh{S}'=f^*(\sh{S})$, e seja $\sh{M}=\bigoplus_{k\in\bb{Z}}\sh{M}_k$ um $\sh{S}'$-módulo graduado quase-coerente de tipo finito.
Então os $\RR^p f_*(\sh{M})=\bigoplus_{k\in\bb{Z}}\RR^p f_*(\sh{M}_k)$ são $\sh{S}$-módulos graduados de tipo finito para todo $p$.
Suponhamos além disso que $\sh{S}$ está gerada por $\sh{S}_1$; então, para todo $p\in\bb{Z}$, existe um inteiro $k_p$ tal que, para todo $k\geq k_p$ e todo $r\geq 0$, tem-se
\[
\label{III.3.3.1.1}
  \RR^p f_*(\sh{M}_{k+r})=\sh{S}_r\RR^p f_*(\sh{M}_k).
  \tag{3.3.1.1}
\] 
\end{proposition}

\begin{proof}
A primeira afirmação é idêntica ao enunciado do Teorema~\sref{III.2.4.1}[i], no qual simplesmente se substitui «morfismo projetivo» por «morfismo próprio».
Na demonstração do Teorema~\sref{III.2.4.1}[i], a hipótese sobre $f$ foi utilizada somente para demonstrar (com a notação daquela demonstração) que $\RR^p f_*'(\widetilde{\sh{M}})$ é um $\sh{O}_{Y'}$-módulo coerente.
Sob as hipóteses da Proposição~\sref{III.3.3.1}, $f'$ é próprio~\sref[II]{II.5.4.2}[iii], de modo que podemos retomar sem modificação a demonstração do Teorema~\sref{III.2.4.1}[i], graças ao teorema de finitude~\sref{III.3.2.1}.

Para a segunda afirmação, basta observar que existe um recobrimento aberto afim finito $(U_i)$ de $Y$ tal que as restrições aos $U_i$ dos dois membros de~(\hyperref[III.3.3.1.1]{3.3.1.1}) são iguais para todo $k\geq k_{p,i}$~\sref[II]{II.2.1.6}[ii]; basta tomar como $k_p$ o maior dos $k_{p,i}$.
\end{proof}

\begin{corollary}[3.3.2]
\label{III.3.3.2}
Seja $A$ um anel noetheriano, seja $\mathfrak{m}$ um ideal de $A$, seja $X$ um $A$-esquema próprio e seja $\sh{F}$ um $\sh{O}_X$-módulo coerente.
Então, para todo $p\geq 0$, a soma direta $\bigoplus_{k\geq 0}\HH^p(X,\mathfrak{m}^k\sh{F})$ é um módulo de tipo finito sobre o anel $S=\bigoplus_{k\geq 0}\mathfrak{m}^k$; em particular, existe um inteiro $k_p\geq 0$ tal que, para todo $k\geq k_p$ e todo $r\geq 0$, tem-se
\[
\label{III.3.3.2.1}
  \HH^p(X,\mathfrak{m}^{k+r}\sh{F})=\mathfrak{m}^r\HH^p(X,\mathfrak{m}^k\sh{F}).
  \tag{3.3.2.1}
\]
\end{corollary}

\begin{proof}
Basta aplicar a Proposição~\sref{III.3.3.1} com $Y=\Spec(A)$, $\sh{S}=\widetilde{S}$ e $\sh{M}_k=\mathfrak{m}^k\sh{F}$, tendo em conta o Corolário~\sref{III.1.4.11}.
\end{proof}

Recordemos que a estrutura de $S$-módulo sobre $\bigoplus_{k\geq 0}\HH^p(X,\mathfrak{m}^k\sh{F})$ obtém-se considerando, para todo $a\in\mathfrak{m}^r$, a aplicação $\HH^p(X,\mathfrak{m}^k\sh{F})\to\HH^p(X,\mathfrak{m}^{k+r}\sh{F})$ que procede da passagem à cohomologia da aplicação de multiplicação $\mathfrak{m}^k\sh{F}\to\mathfrak{m}^{k+r}\sh{F}$ definida por $a$~\sref{III.2.4.1}.

\subsection{Teorema de finitude: caso dos esquemas formais}
\label{subsection:III.3.4}

\oldpage[III]{119}
Os resultados desta seção (salvo a definição~\sref{III.3.4.1}) não serão utilizados no restante deste capítulo.

\begin{env}[3.4.1]
\label{III.3.4.1}
Sejam $\mathfrak{X}$ e $\mathfrak{S}$ dois pré-esquemas formais localmente noetherianos~\sref[I]{I.10.4.2}, e seja $f:\mathfrak{X}\to\mathfrak{S}$ um morfismo de pré-esquemas formais.
Diremos que $f$ é um morfismo \emph{próprio} se satisfaz as condições seguintes:
\begin{enumerate}
  \item[1.º] \emph{$f$ é um morfismo de tipo finito~\sref[I]{I.10.13.3}}.
  \item[2.º] \emph{Se $\sh{K}$ é um feixe de ideais de definição de $\mathfrak{S}$ e se põe $\sh{J}=f^*(\sh{K})\sh{O}_\mathfrak{X}$, $X_0=(\mathfrak{X},\sh{O}_\mathfrak{X}/\sh{J})$, $S_0=(\mathfrak{S},\sh{O}_\mathfrak{S}/\sh{K})$, então o morfismo $f_0:X_0\to S_0$ induzido por $f$~\sref[I]{I.10.5.6} é próprio}.
\end{enumerate}
É imediato que esta definição não depende do feixe de ideais de definição $\sh{K}$ de $\sh{S}$ considerado; com efeito, se $\sh{K}'$ é um segundo feixe de ideais de definição tal que $\sh{K}'\subset\sh{K}$, e se se põe $\sh{J}'=f^*(\sh{K}')\sh{O}_\mathfrak{X}$, $X_0'=(\mathfrak{X},\sh{O}_\mathfrak{X}/\sh{J}')$, $S_0'=(\sh{S},\sh{O}_\mathfrak{S}/\sh{K}')$, então o morfismo $f_0':X_0'\to S_0'$ induzido por $f$ é tal que o diagrama
\[
  \xymatrix{
    X_0\ar[r]^{f_0}\ar[d]_i &
    S_0\ar[d]^j\\
    X_0'\ar[r]_{f_0'} &
    S_0'
  }
\]
é comutativo, sendo $i$ e $j$ imersões sobrejetivas; em virtude de~\sref[II]{II.5.4.5}, é equivalente dizer que $f_0$ ou $f_0'$ é próprio.

Observemos que, para todo $n\geq 0$, se se põe $X_n=(\mathfrak{X},\sh{O}_\mathfrak{X}/\sh{J}^{n+1})$, $S_n=(\mathfrak{S},\sh{O}_\mathfrak{S}/\sh{K}^{n+1})$, então o morfismo $f_n:X_n\to S_n$ induzido por $f$~\sref[I]{I.10.5.6} é próprio para todo $n$ na medida em que o é para $n=0$~\sref[II]{II.5.4.6}.

Se $g:Y\to Z$ é um morfismo próprio de pré-esquemas ordinários localmente noetherianos, $Z'$ é um subconjunto fechado de $Z$, e $Y'$ é um subconjunto fechado de $Y$ tal que $g(Y')\subset Z'$, então a extensão $\widehat{g}:Y_{/Y'}\to Z_{/Z'}$ de $g$ aos completamentos~\sref[I]{I.10.9.1} é um morfismo próprio de pré-esquemas formais, como se deduz da definição e de~\sref[II]{II.5.4.5}.

Sejam $\mathfrak{X}$ e $\mathfrak{S}$ dois pré-esquemas formais localmente noetherianos e seja $f:\mathfrak{X}\to\mathfrak{S}$ um morfismo \emph{de tipo finito}~\sref[I]{I.10.13.3}; conservando a notação anterior, diremos que um subconjunto $Z$ do espaço subjacente a $\mathfrak{X}$ é \emph{próprio} sobre $\mathfrak{S}$ (ou próprio para $f$) se, considerado como subconjunto de $X_0$, $Z$ é \emph{próprio sobre $S_0$}~\sref[II]{II.5.4.10}.
Todas as propriedades dos subconjuntos próprios de pré-esquemas ordinários enunciadas em~\sref[II]{II.5.4.10} continuam válidas para os subconjuntos próprios de pré-esquemas formais, como se deduz imediatamente das definições.
\end{env}

\begin{theorem}[3.4.2]
\label{III.3.4.2}
Sejam $\mathfrak{X}$ e $\mathfrak{Y}$ pré-esquemas formais localmente noetherianos, e seja $f:\mathfrak{X}\to\mathfrak{Y}$ um morfismo próprio.
Para todo $\sh{O}_\mathfrak{X}$-módulo coerente $\sh{F}$, os $\sh{O}_\mathfrak{Y}$-módulos $\RR^q f_*(\sh{F})$ são coerentes para todo $q\geq 0$.
\end{theorem}

Seja $\sh{J}$ um feixe de ideais de definição de $\mathfrak{Y}$, seja $\sh{K}=f^*(\sh{J})\sh{O}_\mathfrak{X}$, e consideremos os $\sh{O}_\mathfrak{X}$-módulos
\[
\label{III.3.4.2.1}
  \sh{F}_k=\sh{F}\otimes_{\sh{O}_\mathfrak{Y}}(\sh{O}_\mathfrak{Y}/\sh{J}^{k+1})=\sh{F}/\sh{K}^{k+1}\sh{F}\quad(k\geq 0)
  \tag{3.4.2.1}
\]
que formam evidentemente um \emph{sistema projetivo} de $\sh{O}_\mathfrak{X}$-módulos topológicos, tais que
\oldpage[III]{120}
\[
  \sh{F}=\varprojlim_k\sh{F}_k
\]
por \sref[I]{I.10.11.3}.
Por outro lado, do Teorema~\sref{III.3.4.2} deduz-se que cada $\RR^q f_*(\sh{F})$, por ser coerente, está dotado naturalmente de uma estrutura de $\sh{O}_\mathfrak{Y}$-módulo topológico~\sref[I]{I.10.11.6}, e o mesmo ocorre com os $\RR^q f_*(\sh{F}_k)$.
Aos homomorfismos canônicos $\sh{F}\to\sh{F}_k=\sh{F}/\sh{K}^{k+1}\sh{F}$ correspondem canonicamente homomorfismos
\[
  \RR^q f_*(\sh{F})\to\RR^q f_*(\sh{F}_k),
\]
que são necessariamente contínuos para as estruturas de $\sh{O}_\mathfrak{Y}$-módulo topológico anteriores~\sref[I]{I.10.11.6} e formam um sistema projetivo, fornecendo, no limite, um homomorfismo funtorial canônico
\[
\label{III.3.4.2.2}
  \RR^q f_*(\sh{F})\to\varprojlim_k\RR^q f_*(\sh{F}_k),
  \tag{3.4.2.2}
\]
que é um homomorfismo contínuo de $\sh{O}_\mathfrak{Y}$-módulos topológicos.
Junto com o Teorema~\sref{III.3.4.2}, demonstraremos o
\begin{corollary}[3.4.3]
\label{III.3.4.3}
Cada um dos homomorfismos~(\hyperref[III.3.4.2.2]{3.4.2.2}) é um isomorfismo topológico.
Além disso, se $\mathfrak{Y}$ é noetheriano, o sistema projetivo $(\RR^q f_*(\sh{F}/\sh{K}^{k+1}\sh{F}))_{k\geq 0}$ satisfaz a condição \emph{(ML)}~\sref[0]{0.13.1.1}.
\end{corollary}
Começaremos estabelecendo o Teorema~\sref{III.3.4.2} e o Corolário~\sref{III.3.4.3} quando $Y$ é um esquema formal afim noetheriano~\sref[I]{I.10.4.1}:
\begin{corollary}[3.4.4]
\label{III.3.4.4}
Sob as hipóteses do Teorema~\sref{III.3.4.2}, suponhamos além disso que $\mathfrak{Y}=\Spf(A)$, onde $A$ é um anel ádico noetheriano.
Seja $\mathfrak{J}$ um ideal de definição de $A$, e ponhamos $\sh{F}_k=\sh{F}/\mathfrak{J}^{k+1}\sh{F}$ para $k\geq 0$.
Então os $\HH^n(\mathfrak{X},\sh{F})$ são $A$-módulos de tipo finito; o sistema projetivo $(\HH^n(\mathfrak{X},\sh{F}_k))_{k\geq 0}$ satisfaz a condição \emph{(ML)} para todo $n$; se se põe
\[
\label{III.3.4.4.1}
  N_{n,k}=\Ker\big(\HH^n(\mathfrak{X},\sh{F})\to\HH^n(\mathfrak{X},\sh{F}_k)\big)
  \tag{3.4.4.1}
\]
(igual também a $\Im(\HH^n(\mathfrak{X},\mathfrak{J}^{k+1}\sh{F})\to\HH^n(\mathfrak{X},\sh{F}))$ pela sucessão exata de cohomologia), então os $N_{n,k}$ definem sobre $\HH^n(\mathfrak{X},\sh{F})$ uma filtração $\mathfrak{J}$-boa~\sref[0]{0.13.7.7}; finalmente, o homomorfismo canônico
\[
\label{III.3.4.4.2}
  \HH^n(\mathfrak{X},\sh{F})\to\varprojlim_k\HH^n(\mathfrak{X},\sh{F}_k)
  \tag{3.4.4.2}
\]
é um isomorfismo topológico para todo $n$ (o membro esquerdo está dotado da topologia $\mathfrak{J}$-ádica e os $\HH^n(\mathfrak{X},\sh{F}_k)$ da topologia discreta).
\end{corollary}

Ponhamos
\[
\label{III.3.4.4.3}
  S=\gr(A)=\bigoplus_{k\geq 0}\mathfrak{J}^k/\mathfrak{J}^{k+1},\ \sh{M}=\gr(\sh{F})=\bigoplus_{k\geq 0}\mathfrak{J}^k\sh{F}/\mathfrak{J}^{k+1}\sh{F}.
\]
Sabemos que $\mathfrak{J}^\Delta$ é um feixe de ideais de definição de $\mathfrak{Y}$~\sref[I]{I.10.3.1}; sejam $\sh{K}=f^*(\mathfrak{J}^\Delta)\sh{O}_\mathfrak{X}$, $X_0=(\mathfrak{X},\sh{O}_\mathfrak{X}/\sh{K})$, $Y_0=(\mathfrak{Y},\sh{O}_\mathfrak{Y}/\mathfrak{J}^\Delta)=\Spec(A_0)$, com $A_0=A/\mathfrak{J}$.
Está claro que os $\sh{M}_k=\mathfrak{J}^k\sh{F}/\mathfrak{J}^{k+1}\sh{F}$ são $\sh{O}_{X_0}$-módulos coerentes~\sref[I]{I.10.11.3}.
Consideremos, por outro lado, a $\sh{O}_{X_0}$-álgebra graduada quase-coerente
\[
\label{III.3.4.4.4}
  \sh{S}=\sh{O}_{X_0}\otimes_{A_0}S=\gr(\sh{O}_\mathfrak{X})=\bigoplus_{k\geq 0}\sh{K}^k/\sh{K}^{k+1}.
  \tag{3.4.4.4}
\]

A hipótese de que $\sh{F}$ é um $\sh{O}_\mathfrak{X}$-módulo de tipo finito implica primeiro que $\sh{M}$ é
\oldpage[III]{121}
um $\sh{S}$-módulo graduado \emph{de tipo finito}.
Com efeito, a questão é local sobre $\mathfrak{X}$, e podemos supor que $\mathfrak{X}=\Spf(B)$, onde $B$ é um anel ádico noetheriano, e que $\sh{F}=N^\Delta$, onde $N$ é um $B$-módulo de tipo finito~\sref[I]{I.10.10.5}; tem-se além disso $X_0=\Spec(B_0)$, onde $B_0=B/\mathfrak{J}B$, e os $\sh{O}_{X_0}$-módulos quase-coerentes $\sh{S}$ e $\sh{M}$ são respectivamente iguais a $\widetilde{S'}$ e $\widetilde{M'}$, onde $S'=\bigoplus_{k\geq 0}((\mathfrak{J}^k/\mathfrak{J}^{k+1})\otimes_{A_0}B_0)$ e $M'=\bigoplus_{k\geq 0}((\mathfrak{J}^k/\mathfrak{J}^{k+1})\otimes_{A_0}N_0)$, com $N_0=N/\mathfrak{J}N$; tem-se evidentemente $M'=S'\otimes_{B_0}N_0$, e como $N_0$ é um $B_0$-módulo de tipo finito, $M'$ é um $S'$-módulo de tipo finito, o que prova nossa afirmação~\sref[I]{I.1.3.13}.

Como o morfismo $f_0:X_0\to Y_0$ é \emph{próprio} por hipótese, podemos aplicar o Corolário~\sref{III.3.3.2} a $\sh{S}$, $\sh{M}$ e ao morfismo $f_0$: tendo em conta o Corolário~\sref{III.1.4.11}, concluímos que, para \emph{todo $n\geq 0$}, $\bigoplus_{k\geq 0}\HH^n(X_0,\sh{M}_k)$ é um $S$-módulo graduado \emph{de tipo finito}.
Isto prova que a condição ($\text{F}_n$) de~\sref[0]{0.13.7.7} é satisfeita para \emph{todo $n\geq 0$}, quando se considera o sistema projetivo estrito $(\sh{F}/\mathfrak{J}^k\sh{F})_{k\geq 0}$ de feixes de grupos abelianos sobre $X_0$, cada um dotado de sua estrutura natural de «$A$-módulo filtrado».
Podemos, pois, aplicar~\sref[0]{0.13.7.7}, que demonstra que:
\begin{enumerate}
  \item[1.º] O sistema projetivo $(\HH^n(\mathfrak{X},\sh{F}_k))_{k\geq 0}$ satisfaz a condição (ML).
  \item[2.º] Se $\HH^{\prime n}=\varprojlim_k\HH^n(\mathfrak{X},\sh{F}_k)$, então $\HH^{\prime n}$ é um $A$-módulo de tipo finito.
  \item[3.º] A filtração definida sobre $\HH^{\prime n}$ pelos núcleos dos homomorfismos canônicos $\HH^{\prime n}\to\HH^n(\mathfrak{X},\sh{F}_k)$ é $\mathfrak{J}$-boa.
\end{enumerate}

Observemos, por outro lado, que se se põe $X_k=(\mathfrak{X},\sh{O}_\mathfrak{X}/\sh{K}^{k+1})$, então $\sh{F}_k$ é um $\sh{O}_{X_k}$-módulo coerente~\sref[I]{I.10.11.3}; e se $U$ é um aberto afim de $X_0$, então $U$ é também um aberto afim de cada $X_k$~\sref[I]{I.5.1.9}, de modo que $\HH^n(U,\sh{F}_k)=0$ para todo $n>0$ e todo $k$~\sref{III.1.3.1}, e $\HH^0(U,\sh{F}_k)\to\HH^0(U,\sh{F}_h)$ é sobrejetivo para $h\leq k$~\sref[I]{I.1.3.9}.
Estamos, pois, nas condições de~\sref[0]{0.13.3.2}, e a aplicação de~\sref[0]{0.13.3.1} demonstra que $\HH^{\prime n}$ identifica-se canonicamente com $\HH^n(\mathfrak{X},\varprojlim_k\sh{F}_k)=\HH^n(\mathfrak{X},\sh{F})$; isto termina a demonstração do Corolário~\sref{III.3.4.4}.

\begin{env}[3.4.5]
\label{III.3.4.5}
Voltemos à demonstração de~\sref{III.3.4.2} e~\sref{III.3.4.3}.
Demonstremos primeiro as proposições no caso $\mathfrak{Y}=\Spf(A)$ considerado em~\sref{III.3.4.4}; para isso, para todo $g\in A$, apliquemos~\sref{III.3.4.4} ao esquema formal afim noetheriano induzido sobre o aberto $\mathfrak{Y}_g=\mathfrak{D}(g)$ de $\mathfrak{Y}$, que é igual a $\Spf(A_{\{g\}})$, e ao pré-esquema formal induzido por $\mathfrak{X}$ sobre $f^{-1}(\mathfrak{Y}_g)$; observemos que $\mathfrak{Y}_g$ é também um aberto afim do pré-esquema $Y_k=(\mathfrak{Y},\sh{O}_\mathfrak{Y}/(\mathfrak{J}^\Delta)^{k+1})$, e como $\sh{F}_k$ é um $\sh{O}_{X_k}$-módulo coerente, tem-se
\[
  \HH^n(f^{-1}(\mathfrak{Y}_g),\sh{F}_k)=\Gamma(\mathfrak{Y}_g,\RR^n f_*(\sh{F}_k))
\]
para todo $k\geq 0$, em virtude do Corolário~\sref{III.1.4.11}.
O homomorfismo canônico
\[
  \HH^n(f^{-1}(\mathfrak{Y}_g),\sh{F})\to\varprojlim_k\Gamma(\mathfrak{Y}_g,\RR^n f_*(\sh{F}_k))
\]
é um isomorfismo; mas tem-se~\sref[0]{0.3.2.6}
\[
  \varprojlim_k\Gamma(\mathfrak{Y}_g,\RR^n f_*(\sh{F}_k))=\Gamma(\mathfrak{Y}_g,\varprojlim_k\RR^n f_*(\sh{F}_k)),
\]
\oldpage[III]{122}
e como o feixe $\RR^n f_*(\sh{F})$ é o feixe associado ao pré-feixe $\mathfrak{Y}_g\mapsto\HH^n(f^{-1}(\mathfrak{Y}_g),\sh{F})$ sobre os $\mathfrak{Y}_g$~\sref[0]{0.3.2.1}, demonstramos que o homomorfismo~(\hyperref[III.3.4.2.2]{3.4.2.2}) é \emph{bijetivo}.
Demonstremos agora que $\RR^n f_*(\sh{F})$ é um $\sh{O}_\mathfrak{Y}$-módulo coerente e, mais precisamente, que se tem
\[
\label{III.3.4.5.1}
  \RR^n f_*(\sh{F})=\big(\HH^n(\mathfrak{X},\sh{F})\big)^\Delta.
  \tag{3.4.5.1}
\]

Com a notação anterior, como $\sh{F}_k$ é um $\sh{O}_{X_k}$-módulo coerente~\sref{III.1.4.13}, tem-se
\[
  \Gamma(\mathfrak{Y}_g,\RR^n f_*(\sh{F}_k))=(\Gamma(\mathfrak{Y},\RR^n f_*(\sh{F}_k)))_g=(\HH^n(\mathfrak{X},\sh{F}_k))_g.
\]

Ora, os $\HH^n(\mathfrak{X},\sh{F}_k)$ formam um sistema projetivo que satisfaz (ML), e seu limite projetivo $\HH^n(\mathfrak{X},\sh{F})$ é um $A$-módulo de tipo finito.
Concluímos~\sref[0]{0.13.7.8} que se tem
\[
  \varprojlim_k\big((\HH^n(\mathfrak{X},\sh{F}_k))_g\big)=\HH^n(\mathfrak{X},\sh{F})\otimes_A A_{\{g\}}=\Gamma(\mathfrak{Y}_g,(\HH^n(\mathfrak{X},\sh{F}))^\Delta),
\]
tendo em conta~\sref[I]{I.10.10.8} aplicado a $A$ e $A_{\{g\}}$; isto demonstra~(\hyperref[III.3.4.5.1]{3.4.5.1}), pois $\Gamma(\mathfrak{Y}_g,\RR^n f_*(\sh{F}))=\varprojlim_k\Gamma(\mathfrak{Y}_g,\RR^n f_*(\sh{F}_k))$.

Como (\hyperref[III.3.4.2.2]{3.4.2.2}) é então um isomorfismo de $\sh{O}_\mathfrak{Y}$-módulos coerentes, é necessariamente um isomorfismo \emph{topológico}~\sref[I]{I.10.11.6}.
Finalmente, das relações $\RR^n f_*(\sh{F}_k)=(\HH^n(\mathfrak{X},\sh{F}_k))^\Delta$ deduz-se que o sistema projetivo $(\RR^n f_*(\sh{F}_k))_{k\geq 0}$ satisfaz (ML)~\sref[I]{I.10.10.2}.

Uma vez demonstrados \sref{III.3.4.2} e \sref{III.3.4.3} no caso em que o pré-esquema formal $\mathfrak{Y}$ é afim noetheriano, a passagem ao caso geral é imediata para~\sref{III.3.4.2} e para a primeira afirmação de~\sref{III.3.4.3}, que são locais sobre $\mathfrak{Y}$.
Quanto à segunda afirmação de~\sref{III.3.4.3}, como $\mathfrak{Y}$ é noetheriano, basta recobri-lo por um número finito de abertos afins noetherianos $U_i$ e observar que as restrições do sistema projetivo $(\RR^q f_*(\sh{F}_k))$ a cada um dos $U_i$ satisfazem (ML).
\end{env}

No curso da demonstração, provamos, além disso:
\begin{corollary}[3.4.6]
\label{III.3.4.6}
Sob as hipóteses do Corolário~\sref{III.3.4.4}, o homomorfismo canônico
\[
\label{III.3.4.6.1}
  \HH^q(\mathfrak{X},\sh{F})\to\Gamma(\mathfrak{Y},\RR^q f_*(\sh{F}))
  \tag{3.4.6.1}
\]
é bijetivo.
\end{corollary}

\section{O teorema fundamental dos morfismos próprios. Aplicações}
\label{section:III.4}


\subsection{O teorema fundamental}
\label{subsection:III.4.1}

\begin{env}[4.1.1]
\label{III.4.1.1}
Sejam $X$ e $Y$ esquemas ordinários noetherianos, seja $f:X\to Y$ um
morfismo próprio, seja $Y'$ um subconjunto fechado de $Y$, e ponhamos
\[
  X'=f^{-1}(Y').
\]
Denotamos por $\widehat{X}$ e $\widehat{Y}$ os esquemas formais
$X_{/X'}$ e $Y_{/Y'}$, obtidos completando $X$ e $Y$ ao longo de
$X'$ e $Y'$, respectivamente \sref[I]{I.10.8.5}.
Denotamos por $\widehat{f}:\widehat{X}\to\widehat{Y}$ a extensão de
$f$ a estes completamentos \sref[I]{I.10.9.1}.
Para todo $\sh{O}_X$-módulo coerente $\sh{F}$, denotamos por

\oldpage[III]{123}

\[
  \widehat{\sh{F}}=\sh{F}_{/X'}
\]
seu completamento ao longo de $X'$ \sref[I]{I.10.8.4}; este é um
$\sh{O}_{\widehat{X}}$-módulo coerente \sref[I]{I.10.8.8}.
\end{env}

\begin{env}[4.1.2]
\label{III.4.1.2}
Seja $\sh{J}$ um ideal coerente de $\sh{O}_Y$ tal que
\[
  \Supp(\sh{O}_Y/\sh{J})=Y'
\]
\sref[I]{I.5.2.1}.
Sabemos \sref[I]{I.4.4.5} que
\[
  \sh{K}=f^*(\sh{J})\sh{O}_X
\]
é um ideal coerente de $\sh{O}_X$ que satisfaz
\[
  \Supp(\sh{O}_X/\sh{K})=X'.
\]
Para todo $k\geq0$, consideremos o $\sh{O}_X$-módulo coerente
\[
  \sh{F}_k
  =
  \sh{F}\otimes_{\sh{O}_Y}(\sh{O}_Y/\sh{J}^{k+1})
  =
  \sh{F}/\sh{K}^{k+1}\sh{F}.
\]
Os $\sh{O}_Y$-módulos
\[
  \RR^n f_*(\sh{F})
  \quad\text{e}\quad
  \RR^n f_*(\sh{F}_k)
\]
são coerentes para todo $n$ \sref{III.3.2.1}.
Para todo $k\geq0$ e todo $n$, o homomorfismo canônico
$\sh{F}\to\sh{F}_k$ induz, por funtorialidade, um homomorfismo
\[
  \RR^n f_*(\sh{F})
  \longrightarrow
  \RR^n f_*(\sh{F}_k).
  \tag{4.1.2.1}
  \label{III.4.1.2.1}
\]
Como $\sh{F}_k$ é um $\sh{O}_X/\sh{K}^{k+1}$-módulo,
$\RR^n f_*(\sh{F}_k)$ é um
$\sh{O}_Y/\sh{J}^{k+1}$-módulo \sref[0]{0.12.2.1}.
Deduzimos, pois, de \eqref{III.4.1.2.1} um homomorfismo
\[
  \RR^n f_*(\sh{F})
  \otimes_{\sh{O}_Y}(\sh{O}_Y/\sh{J}^{k+1})
  \longrightarrow
  \RR^n f_*(\sh{F}_k).
  \tag{4.1.2.2}
  \label{III.4.1.2.2}
\]
Os dois membros de \eqref{III.4.1.2.2} formam sistemas projetivos.
O limite projetivo do membro esquerdo é o completamento
\[
  \bigl(\RR^n f_*(\sh{F})\bigr)_{/Y'},
\]
que denotaremos por
\[
  \widehat{\RR^n f_*(\sh{F})}.
\]
Os homomorfismos \eqref{III.4.1.2.2} são compatíveis com os sistemas
projetivos e definem, portanto, ao passar ao limite, um homomorfismo
canônico
\[
  \vphi_n:
  \widehat{\RR^n f_*(\sh{F})}
  \longrightarrow
  \varprojlim_k\RR^n f_*(\sh{F}_k).
  \tag{4.1.2.3}
  \label{III.4.1.2.3}
\]
Além disso, \eqref{III.4.1.2.2} é um homomorfismo de
$\sh{O}_Y/\sh{J}^{k+1}$-módulos.
Por \sref[I]{I.10.8.3}, pode, portanto, ser considerado um
homomorfismo contínuo de $\sh{O}_{\widehat{Y}}$-módulos topológicos
munidos de suas topologias pseudodiscretas \sref[0]{0.3.8.1}.
Consequentemente, $\vphi_n$ é um homomorfismo contínuo de
$\sh{O}_{\widehat{Y}}$-módulos topológicos.
\end{env}

\begin{env}[4.1.3]
\label{III.4.1.3}
Seja $i:\widehat{X}\to X$ o morfismo canônico de espaços anelados
definido em \sref[I]{I.10.8.7}.
Temos o diagrama comutativo
\[
  \xymatrix{
    X_k\ar[r]^{h_k}\ar[dr]_{i_k} &
    \widehat{X}\ar[d]^{i}\\
    & X
  }
  \tag{4.1.3.1}
  \label{III.4.1.3.1}
\]
onde $X_k$ é o subesquema fechado de $X$ definido por
$\sh{K}^{k+1}$, $i_k$ é a imersão canônica, e $h_k$ é o
morfismo de espaços anelados que corresponde à identidade sobre
os espaços subjacentes e ao homomorfismo canônico
\[
  \sh{O}_{\widehat{X}}
  \longrightarrow
  \sh{O}_X/\sh{K}^{k+1}
\]
\sref[I]{I.10.5.2}.
Além disso,
\[
  \widehat{\sh{F}}=i^*(\sh{F})
\]
salvo isomorfismo canônico \sref[I]{I.10.8.8}.
Temos
\[
  \HH^n(X_k,i_k^*(\sh{F}_k))
  =
  \HH^n(X,\sh{F}_k)
  \tag{4.1.3.2}
  \label{III.4.1.3.2}
\]
salvo isomorfismo canônico.

\oldpage[III]{124}

Com efeito,
\[
  \sh{F}_k=(i_k)_*(i_k^*(\sh{F}_k))
\]
(G,~II,~4.9.1).
O homomorfismo canônico
\[
  \HH^n(\widehat{X},\widehat{\sh{F}})
  \longrightarrow
  \HH^n(X_k,h_k^*(\widehat{\sh{F}}))
\]
\sref[0]{0.12.1.3.5} pode escrever-se também
\[
  \HH^n(\widehat{X},\widehat{\sh{F}})
  \longrightarrow
  \HH^n(X,\sh{F}_k).
  \tag{4.1.3.3}
  \label{III.4.1.3.3}
\]
Estes homomorfismos formam um sistema projetivo e definem, portanto,
um homomorfismo canônico
\[
  \psi_{n,X}:
  \HH^n(\widehat{X},\widehat{\sh{F}})
  \longrightarrow
  \varprojlim_k\HH^n(X,\sh{F}_k).
  \tag{4.1.3.4}
  \label{III.4.1.3.4}
\]
Substituindo $X$ por um subconjunto aberto da forma $f^{-1}(V)$,
onde $V$ é um aberto afim de $Y$, e tendo em conta
\sref{III.1.4.11}, obtemos homomorfismos
\[
  \psi_{n,V}:
  \HH^n(\widehat{X}\cap f^{-1}(V),\widehat{\sh{F}})
  \longrightarrow
  \varprojlim_k\Gamma(V,\RR^n f_*(\sh{F}_k)).
  \tag{4.1.3.5}
  \label{III.4.1.3.5}
\]
Estes homomorfismos comutam manifestamente com a restrição de $V$
a um subconjunto aberto afim menor.
Definem, portanto, um homomorfismo canônico de feixes
\[
  \psi_n:
  \RR^n\widehat{f}_*(\widehat{\sh{F}})
  \longrightarrow
  \varprojlim_k\RR^n f_*(\sh{F}_k).
  \tag{4.1.3.6}
  \label{III.4.1.3.6}
\]
\end{env}

\begin{env}[4.1.4]
\label{III.4.1.4}
Por fim, seja $j:\widehat{Y}\to Y$ o morfismo canônico de espaços
anelados \sref[I]{I.10.8.7}.
Como $\RR^n f_*(\sh{F})$ é um $\sh{O}_Y$-módulo coerente
\sref{III.3.2.1}, temos, salvo isomorfismo canônico,
\[
  j^*(\RR^n f_*(\sh{F}))
  =
  \widehat{\RR^n f_*(\sh{F})}
\]
\sref[I]{I.10.8.8}.
Existe, consequentemente, um homomorfismo canônico
\[
  \rho_n:
  \widehat{\RR^n f_*(\sh{F})}
  =
  j^*(\RR^n f_*(\sh{F}))
  \longrightarrow
  \RR^n\widehat{f}_*(i^*(\sh{F}))
  =
  \RR^n\widehat{f}_*(\widehat{\sh{F}})
  \tag{4.1.4.1}
  \label{III.4.1.4.1}
\]
definido para os espaços anelados em geral; veja-se a demonstração
de \sref{III.1.4.15}.
Demonstraremos que o diagrama
\[
  \xymatrix{
    \widehat{\RR^n f_*(\sh{F})}
      \ar[rr]^{\rho_n}\ar[dr]_{\vphi_n}
    &&
    \RR^n\widehat{f}_*(\widehat{\sh{F}})
      \ar[dl]^{\psi_n}\\
    &
    \displaystyle\varprojlim_k\RR^n f_*(\sh{F}_k)
    &
  }
  \tag{4.1.4.2}
  \label{III.4.1.4.2}
\]
é comutativo.
Basta demonstrar a comutatividade do diagrama formado pelos
homomorfismos correspondentes de pré-feixes.
Podemos, portanto, restringir-nos ao caso em que $Y$ é afim, e basta
demonstrar que
\[
  \xymatrix{
    \widehat{\HH^n(X,\sh{F})}
      \ar[rr]^{\rho_n}\ar[dr]_{\vphi_n}
    &&
    \HH^n(\widehat{X},\widehat{\sh{F}})
      \ar[dl]^{\psi_{n,X}}\\
    &
    \displaystyle\varprojlim_k\HH^n(X,\sh{F}_k)
    &
  }
  \tag{4.1.4.3}
  \label{III.4.1.4.3}
\]
é comutativo.
A comutatividade de \eqref{III.4.1.3.1}, junto com as relações
entre os grupos de cohomologia estabelecidas em \sref{III.4.1.3}, dá
o diagrama comutativo

\oldpage[III]{125}

\[
  \xymatrix{
    \HH^n(X,\sh{F})
      \ar[r]\ar[dr]
    &
    \HH^n(\widehat{X},\widehat{\sh{F}})
      =
    \HH^n(\widehat{X},i^*(\sh{F}))
      \ar[d]\\
    &
    \HH^n(X,\sh{F}_k)
      =
    \HH^n(X_k,i_k^*(\sh{F})).
  }
\]
Segue imediatamente a comutatividade de \eqref{III.4.1.4.3}.
\end{env}

\begin{theorem}[4.1.5]
\label{III.4.1.5}
Seja $f:X\to Y$ um morfismo próprio de esquemas noetherianos, seja
$Y'$ um subconjunto fechado de $Y$, e ponhamos $X'=f^{-1}(Y')$.
Então, para todo $\sh{O}_X$-módulo coerente $\sh{F}$,
\[
  \RR^n\widehat{f}_*(\widehat{\sh{F}})
\]
é um $\sh{O}_{\widehat{Y}}$-módulo coerente, e os homomorfismos
$\vphi_n$, $\psi_n$ e $\rho_n$ do diagrama
\eqref{III.4.1.4.2} são isomorfismos topológicos.
\end{theorem}

\begin{proof}
Basta demonstrar que $\vphi_n$ e $\psi_n$ são isomorfismos.
Com efeito, $\RR^n f_*(\sh{F})$ é coerente \sref{III.3.2.1}, logo
seu completamento $\widehat{\RR^n f_*(\sh{F})}$ é coerente
\sref[I]{I.10.8.8}; a bicontinuidade de $\vphi_n$, $\psi_n$ e
$\rho_n$ é então automática \sref[I]{I.10.11.6}.
\end{proof}

\begin{env}[4.1.6]
\label{III.4.1.6}
\emph{Observações.}
\begin{enumerate}
  \item[(i)] Ponhamos
  \[
    \widehat{\sh{F}}_k
    =
    \widehat{\sh{F}}/
    \widehat{\sh{K}}^{\,k+1}\widehat{\sh{F}}.
  \]
  Então, imediatamente,
  \[
    \sh{F}_k=i_*(\widehat{\sh{F}}_k),
  \]
  e o homomorfismo canônico \eqref{III.4.1.3.6} é precisamente o
  homomorfismo já definido em (\hyperref[III.3.4.2.2]{3.4.2.2}):
  \[
    \RR^n\widehat{f}_*(\widehat{\sh{F}})
    \longrightarrow
    \varprojlim_k
    \RR^n\widehat{f}_*(\widehat{\sh{F}}_k).
    \tag{4.1.6.1}
    \label{III.4.1.6.1}
  \]
  Assim, o fato de que $\psi_n$ seja um isomorfismo é um caso particular
  de \sref{III.3.4.3}.
  Daremos, não obstante, mais adiante uma demonstração direta, evitando
  as delicadas considerações relativas aos limites projetivos de
  sucessões espectrais \sref[0]{0.13.7} nas quais se apoia o teorema
  geral \sref{III.3.4.3}.

  \item[(ii)] Como os $\psi_n$ são isomorfismos, dizer que os
  $\vphi_n$ são isomorfismos equivale a dizer que
  \[
    \rho_n=\psi_n^{-1}\circ\vphi_n
  \]
  são isomorfismos.
  O Teorema~\sref{III.4.1.5} afirma, em particular, que a formação de
  $\RR^n f_*$ comuta com o completamento; pode chamar-se, portanto, o
  primeiro teorema de comparação entre as teorias ``algébrica'' e
  ``formal''.
\end{enumerate}
Começaremos demonstrando a forma afim de \sref{III.4.1.5}.
\end{env}

\begin{corollary}[4.1.7]
\label{III.4.1.7}
Sob as hipóteses de \sref{III.4.1.5}, suponhamos além disso que
$Y=\Spec(A)$, onde $A$ é noetheriano, e que
$\sh{J}=\widetilde{\mathfrak{J}}$, onde $\mathfrak{J}$ é um ideal
de $A$, de modo que
\[
  \sh{F}_k=\sh{F}/\mathfrak{J}^{k+1}\sh{F}.
\]
O homomorfismo canônico
\[
  \vphi_n:
  \widehat{\HH^n(X,\sh{F})}
  \longrightarrow
  \varprojlim_k\HH^n(X,\sh{F}_k)
  \tag{4.1.7.1}
  \label{III.4.1.7.1}
\]
onde o primeiro termo é o completamento separado de
$\HH^n(X,\sh{F})$ para a topologia $\mathfrak{J}$-pré-ádica, é um
isomorfismo.
Para todo $n$, o sistema projetivo
\[
  \bigl(\HH^n(X,\sh{F}_k)\bigr)_{k\geq0}
\]
satisfaz a condição \emph{(ML)}, e o homomorfismo canônico
\[
  \psi_n:
  \HH^n(\widehat{X},\widehat{\sh{F}})
  \longrightarrow
  \varprojlim_k\HH^n(X,\sh{F}_k)
  \tag{4.1.7.2}
  \label{III.4.1.7.2}
\]

\oldpage[III]{126}

é um isomorfismo.
Por fim, a filtração sobre $\HH^n(X,\sh{F})$ definida pelos
núcleos dos homomorfismos canônicos
\[
  \HH^n(X,\sh{F})\longrightarrow\HH^n(X,\sh{F}_k)
\]
é $\mathfrak{J}$-boa \sref[0]{0.13.7.7}, e $\vphi_n$ é um
isomorfismo topológico.\footnote{A demonstração seguinte, mais simples
que a demonstração original, e o complemento relativo à filtração
de $\HH^n(X,\sh{F})$, nos foram comunicados por J.-P. Serre.}
\end{corollary}

\begin{proof}
Fixemos $n\geq0$ durante toda esta demonstração e, para simplificar,
ponhamos
\[
  \mathrm{H}=\HH^n(X,\sh{F}),
  \qquad
  \mathrm{H}_k=\HH^n(X,\sh{F}_k).
  \tag{4.1.7.3}
  \label{III.4.1.7.3}
\]
Ponhamos também
\[
  \mathrm{R}_k
  =
  \Ker(\mathrm{H}\longrightarrow\mathrm{H}_k),
  \qquad
  \text{um sub-$A$-módulo de $\mathrm{H}$.}
  \tag{4.1.7.4}
  \label{III.4.1.7.4}
\]
A sucessão exata de cohomologia
\[
  \HH^n(X,\mathfrak{J}^{k+1}\sh{F})
  \longrightarrow
  \HH^n(X,\sh{F})
  \longrightarrow
  \HH^n(X,\sh{F}_k)
  \xrightarrow{\partial}
  \HH^{n+1}(X,\mathfrak{J}^{k+1}\sh{F})
  \longrightarrow
  \HH^{n+1}(X,\sh{F})
\]
mostra que
\[
  \mathrm{R}_k
  =
  \Im\bigl(
    \HH^n(X,\mathfrak{J}^{k+1}\sh{F})
    \longrightarrow
    \HH^n(X,\sh{F})
  \bigr).
\]
Ponhamos
\[
  \begin{split}
  \mathrm{Q}_k
  &=
  \Ker\bigl(
    \HH^{n+1}(X,\mathfrak{J}^{k+1}\sh{F})
    \longrightarrow
    \HH^{n+1}(X,\sh{F})
  \bigr)
  \\
  &=
  \Im\bigl(
    \HH^n(X,\sh{F}_k)
    \longrightarrow
    \HH^{n+1}(X,\mathfrak{J}^{k+1}\sh{F})
  \bigr).
  \end{split}
  \tag{4.1.7.5}
  \label{III.4.1.7.5}
\]
Temos, portanto, a sucessão exata
\[
  0\longrightarrow\mathrm{R}_k
  \longrightarrow\mathrm{H}
  \longrightarrow\mathrm{H}_k
  \longrightarrow\mathrm{Q}_k
  \longrightarrow0.
  \tag{4.1.7.6}
  \label{III.4.1.7.6}
\]

\begin{env}[4.1.7.7]
\label{III.4.1.7.7}
Seja $x\in\mathfrak{J}^m$, onde $m\geq0$.
A multiplicação por $x$ sobre $\mathfrak{J}^k\sh{F}$ é um
homomorfismo
\[
  \mathfrak{J}^k\sh{F}
  \longrightarrow
  \mathfrak{J}^{k+m}\sh{F},
\]
e dá, portanto, um homomorfismo
\[
  \mu_{x,m}:
  \HH^n(X,\mathfrak{J}^k\sh{F})
  \longrightarrow
  \HH^n(X,\mathfrak{J}^{k+m}\sh{F}).
  \tag{4.1.7.8}
  \label{III.4.1.7.8}
\]
Seja
\[
  S=\bigoplus_{k\geq0}\mathfrak{J}^k
\]
a $A$-álgebra graduada.
Sabemos que as multiplicações $\mu_{x,m}$ dotam a
\[
  \mathrm{E}
  =
  \bigoplus_{k\geq0}
  \HH^n(X,\mathfrak{J}^k\sh{F})
\]
de uma estrutura de $S$-módulo graduado de tipo finito
\sref{III.3.3.2}; o anel graduado $S$ é noetheriano
\sref[II]{II.2.1.5}.
\end{env}

\begin{lemma}[4.1.7.9]
\label{III.4.1.7.9}
Os submódulos $\mathrm{R}_k$ de $\mathrm{H}$ definem uma
filtração $\mathfrak{J}$-boa sobre $\mathrm{H}$.
\end{lemma}

\begin{proof}
Demonstremos primeiro que
\[
  \mathfrak{J}^m\mathrm{R}_k
  \subset
  \mathrm{R}_{k+m},
  \tag{4.1.7.10}
  \label{III.4.1.7.10}
\]
onde a multiplicação sobre
$\mathrm{H}=\HH^n(X,\sh{F})$ por um elemento
$x\in\mathfrak{J}^m$ é a aplicação $\mu_{x,0}$.
Para todo $x\in\mathfrak{J}^m$, o diagrama
\[
  \xymatrix{
    \mathfrak{J}^{k+1}\sh{F}
      \ar[r]^{x}\ar[d]
    &
    \mathfrak{J}^{k+m+1}\sh{F}
      \ar[d]
    \\
    \sh{F}
      \ar[r]_{x}
    &
    \sh{F}
  }
  \tag{4.1.7.11}
  \label{III.4.1.7.11}
\]

\oldpage[III]{127}

é comutativo; suas flechas horizontais são as multiplicações por
$x$, e suas flechas verticais são as injeções canônicas.
Consequentemente, o diagrama correspondente
\[
  \xymatrix{
    \HH^n(X,\mathfrak{J}^{k+1}\sh{F})
      \ar[r]^{\mu_{x,m}}\ar[d]
    &
    \HH^n(X,\mathfrak{J}^{k+m+1}\sh{F})
      \ar[d]
    \\
    \HH^n(X,\sh{F})
      \ar[r]_{\mu_{x,0}}
    &
    \HH^n(X,\sh{F})
  }
\]
é comutativo.
A interpretação de $\mathrm{R}_k$ como a imagem de
\[
  \HH^n(X,\mathfrak{J}^{k+1}\sh{F})
  \longrightarrow
  \HH^n(X,\sh{F})
\]
demonstra, portanto, \eqref{III.4.1.7.10}.
Mostra também que o $S$-módulo graduado
\[
  \mathrm{R}=\bigoplus_{k\geq0}\mathrm{R}_k
\]
é um quociente do sub-$S$-módulo
\[
  \mathrm{M}
  =
  \bigoplus_{k\geq0}
  \HH^n(X,\mathfrak{J}^{k+1}\sh{F})
\]
de $\mathrm{E}$.
A observação anterior mostra, pois, que $\mathrm{R}$ é um
$S$-módulo de tipo finito.
Isto equivale à afirmação do lema
(Bourbaki, \emph{Alg.\ comm.}, cap.~III, \S3, n.º~1, teorema~1).
\end{proof}

\begin{env}[4.1.7.12]
\label{III.4.1.7.12}
Consideremos agora o $S$-módulo graduado
\[
  \mathrm{N}
  =
  \bigoplus_{k\geq0}
  \HH^{n+1}(X,\mathfrak{J}^{k+1}\sh{F}),
\]
definido como em \eqref{III.4.1.7.8}.
Também é um $S$-módulo de tipo finito por \sref{III.3.3.2}.
Por \eqref{III.4.1.7.5}, temos
$\mathrm{Q}_k\subset\mathrm{N}_k$ para todo $k$, e o diagrama
\eqref{III.4.1.7.11}, substituindo $n$ por $n+1$, mostra que
\[
  S_m\mathrm{Q}_k
  =
  \mathfrak{J}^m\mathrm{Q}_k
  \subset
  \mathrm{Q}_{k+m}.
\]
Em outros termos, $\mathrm{Q}$ é um sub-$S$-módulo graduado de
$\mathrm{N}$ e é, portanto, de tipo finito.
\end{env}

\begin{env}[4.1.7.13]
\label{III.4.1.7.13}
Denotemos por $\alpha_m$ a injeção canônica
\[
  \mathfrak{J}^m\longrightarrow A,
\]
que também pode escrever-se $S_m\to S_0$.
Como $\mathfrak{J}^{k+1}\sh{F}_k=0$, o $A$-módulo
$\HH^n(X,\sh{F}_k)$ é anulado por $\mathfrak{J}^{k+1}$.
Como $\mathrm{Q}_k$ é a imagem do $A$-homomorfismo
\[
  \HH^n(X,\sh{F}_k)
  \longrightarrow
  \HH^{n+1}(X,\mathfrak{J}^{k+1}\sh{F}),
\]
$\mathrm{Q}_k$ é igualmente anulado por $\mathfrak{J}^{k+1}$.
Equivalentemente, no $S$-módulo $\mathrm{Q}$,
\[
  \alpha_{k+1}(S_{k+1})\mathrm{Q}_k=0.
  \tag{4.1.7.14}
  \label{III.4.1.7.14}
\]
Como $\mathrm{Q}$ é um $S$-módulo de tipo finito, existem inteiros
$k_0$ e $h$ tais que
\[
  \mathrm{Q}_{k+h}=S_h\mathrm{Q}_k
  \qquad(k\geq k_0)
\]
\sref[II]{II.2.1.6}[ii].
Esta relação e \eqref{III.4.1.7.14} implicam que existe um inteiro
$r>0$ tal que
\[
  \alpha_r(S_r)\mathrm{Q}=0.
  \tag{4.1.7.15}
  \label{III.4.1.7.15}
\]
\end{env}

\begin{env}[4.1.7.16]
\label{III.4.1.7.16}
A injeção canônica
\[
  \mathfrak{J}^{k+m}\sh{F}
  \longrightarrow
  \mathfrak{J}^k\sh{F}
\]
induz em cohomologia um $A$-homomorfismo
\[
  \nu_m:
  \HH^{n+1}(X,\mathfrak{J}^{k+m}\sh{F})
  \longrightarrow
  \HH^{n+1}(X,\mathfrak{J}^k\sh{F}).
  \tag{4.1.7.17}
  \label{III.4.1.7.17}
\]
Para todo $x\in\mathfrak{J}^m$, temos manifestamente a
fatoração
\[
  \xymatrix@C=4.5em{
    \HH^{n+1}(X,\mathfrak{J}^k\sh{F})
      \ar[r]^{\mu_{x,m}}
    &
    \HH^{n+1}(X,\mathfrak{J}^{k+m}\sh{F})
      \ar[r]^{\nu_m}
    &
    \HH^{n+1}(X,\mathfrak{J}^k\sh{F}),
  }
  \tag{4.1.7.18}
  \label{III.4.1.7.18}
\]
cuja composta é $\mu_{x,0}$.
\end{env}

\oldpage[III]{128}

Segue-se que, para todo sub-$A$-módulo $\mathrm{P}$ de
$\HH^{n+1}(X,\mathfrak{J}^k\sh{F})$, temos no $S$-módulo
$\mathrm{N}$
\[
  \nu_m(S_m\mathrm{P})
  =
  \alpha_m(S_m)\mathrm{P}.
  \tag{4.1.7.19}
  \label{III.4.1.7.19}
\]

\begin{lemma}[4.1.7.20]
\label{III.4.1.7.20}
Existe um inteiro $m>0$ tal que
\[
  \nu_m(\mathrm{Q}_{k+m})=0
\]
para todo $k\geq k_0$.
\end{lemma}

\begin{proof}
Tomemos $m$ múltiplo de $h$ com $m\geq r$.
Como
\[
  \mathrm{Q}_{k+m}=S_m\mathrm{Q}_k
  \qquad(k\geq k_0),
\]
as relações \eqref{III.4.1.7.19} e \eqref{III.4.1.7.15} dão
\[
  \nu_m(\mathrm{Q}_{k+m})
  =
  \alpha_m(S_m)\mathrm{Q}_k
  \subset
  \alpha_r(S_r)\mathrm{Q}_k
  =
  0.
\]
\end{proof}

\begin{env}[4.1.7.21]
\label{III.4.1.7.21}
O diagrama comutativo
\[
  \xymatrix@C=2.5em{
    \HH^n(X,\sh{F})
      \ar[r]
    &
    \HH^n(X,\sh{F}_k)
      \ar[r]^-{\partial}
    &
    \HH^{n+1}(X,\mathfrak{J}^{k+1}\sh{F})
      \ar[r]
    &
    \HH^{n+1}(X,\sh{F})
    \\
    \HH^n(X,\sh{F})
      \ar[r]\ar[u]
    &
    \HH^n(X,\sh{F}_{k+m})
      \ar[r]^-{\partial}\ar[u]
    &
    \HH^{n+1}(X,\mathfrak{J}^{k+m+1}\sh{F})
      \ar[r]\ar[u]
    &
    \HH^{n+1}(X,\sh{F})
      \ar[u]
  }
\]
procede, por sua vez, do diagrama comutativo
\[
  \xymatrix{
    0\ar[r]
    &
    \mathfrak{J}^{k+1}\sh{F}\ar[r]
    &
    \sh{F}\ar[r]
    &
    \sh{F}_k\ar[r]
    &
    0
    \\
    0\ar[r]
    &
    \mathfrak{J}^{k+m+1}\sh{F}\ar[r]\ar[u]
    &
    \sh{F}\ar[r]\ar[u]
    &
    \sh{F}_{k+m}\ar[r]\ar[u]
    &
    0.
  }
\]
Aqui as flechas verticais são as aplicações canônicas.
Obtemos, portanto, o diagrama comutativo
\[
  \xymatrix@C=2.2em{
    0\ar[r]
    &
    \mathrm{R}_k\ar[r]
    &
    \mathrm{H}\ar[r]
    &
    \mathrm{H}_k\ar[r]
    &
    \mathrm{Q}_k\ar[r]
    &
    0
    \\
    0\ar[r]
    &
    \mathrm{R}_{k+m}\ar[r]\ar[u]
    &
    \mathrm{H}\ar[r]\ar@{=}[u]
    &
    \mathrm{H}_{k+m}\ar[r]\ar[u]
    &
    \mathrm{Q}_{k+m}\ar[r]\ar[u]^{\nu_m}
    &
    0,
  }
\]
de linhas exatas.
Para $k\geq k_0$, a última flecha vertical é nula por
\sref{III.4.1.7.20}.
Assim, a imagem de $\mathrm{H}_{k+m}$ em $\mathrm{H}_k$ está contida
em
\[
  \Ker(\mathrm{H}_k\longrightarrow\mathrm{Q}_k)
  =
  \Im(\mathrm{H}\longrightarrow\mathrm{H}_k).
\]
Reciprocamente, a comutatividade mostra que contém
$\Im(\mathrm{H}\to\mathrm{H}_k)$ e, portanto, que é igual a esta
imagem.
O mesmo ocorre, pois, com as imagens em $\mathrm{H}_k$ de
$\mathrm{H}_{k'}$ para $k'\geq k+m$.
Isto demonstra a condição \emph{(ML)} para o sistema projetivo
$(\mathrm{H}_k)_{k\geq0}$.

Além disso, para todo subconjunto aberto afim $U$ de $X$,
\[
  \HH^i(U,\sh{F}_k)=0
  \qquad(i>0)
\]
por \sref{III.1.3.1}, e, para $m>0$, a aplicação
\[
  \HH^0(U,\sh{F}_{k+m})
  \longrightarrow
  \HH^0(U,\sh{F}_k)
\]
é sobrejetiva \sref[I]{I.1.3.9}.
Podemos, portanto,

\oldpage[III]{129}

aplicar \sref[0]{0.13.3.1}; o homomorfismo canônico
\[
  \HH^n(\widehat{X},\widehat{\sh{F}})
  \longrightarrow
  \varprojlim_k\HH^n(X,\sh{F}_k)
\]
é bijetivo para todo $n\geq0$.

Como o sistema projetivo
$(\mathrm{H}/\mathrm{R}_k)_{k\geq0}$ é estrito, podemos passar ao
limite projetivo nas sucessões exatas
\[
  0
  \longrightarrow
  \mathrm{H}/\mathrm{R}_k
  \longrightarrow
  \mathrm{H}_k
  \longrightarrow
  \mathrm{Q}_k
  \longrightarrow
  0
  \tag{4.1.7.22}
  \label{III.4.1.7.22}
\]
por \sref[0]{0.13.2.2}.
Como $\nu_m(\mathrm{Q}_{k+m})=0$, temos
\[
  \varprojlim_k\mathrm{Q}_k=0,
\]
e, portanto, um isomorfismo topológico
\[
  \varprojlim_k(\mathrm{H}/\mathrm{R}_k)
  \isoto
  \varprojlim_k\mathrm{H}_k.
\]
Mas a filtração $(\mathrm{R}_k)$ de $\mathrm{H}$ é
$\mathfrak{J}$-boa e define, portanto, sobre $\mathrm{H}$ a
topologia $\mathfrak{J}$-pré-ádica.
Assim, $\varprojlim_k(\mathrm{H}/\mathrm{R}_k)$ é o completamento
separado de $\mathrm{H}$ para esta topologia.
Isto termina a demonstração de \sref{III.4.1.7}.
\end{env}
\end{proof}

\begin{env}[4.1.8]
\label{III.4.1.8}
Podemos terminar agora a demonstração de \sref{III.4.1.5}.
Para todo subconjunto aberto afim $V$ de $Y$,
\[
  \Gamma\bigl(
    V,\widehat{\RR^n f_*(\sh{F})}
  \bigr)
\]
é o completamento separado de
$\Gamma(V,\RR^n f_*(\sh{F}))$ para a topologia
$\mathfrak{J}$-pré-ádica, se
$\sh{J}|V=\widetilde{\mathfrak{J}}$, porque
$\RR^n f_*(\sh{F})$ é um $\sh{O}_Y$-módulo coerente
\sref[I]{I.10.8.4}.
Além disso,
\[
  \Gamma\left(
    V,\varprojlim_k\RR^n f_*(\sh{F}_k)
  \right)
  =
  \varprojlim_k
  \Gamma\bigl(V,\RR^n f_*(\sh{F}_k)\bigr)
\]
por \sref[0]{0.3.2.6}.
De \sref{III.4.1.7} e \sref{III.1.4.11} segue-se que
$\vphi_n$ é um isomorfismo topológico.
Do mesmo modo, novamente por \sref{III.1.4.11}, \sref{III.4.1.7}
mostra que o homomorfismo $\psi_{n,V}$ de \eqref{III.4.1.3.5} é um
isomorfismo.
Pela definição de
$\RR^n\widehat{f}_*(\widehat{\sh{F}})$, segue-se que $\psi_n$ é um
isomorfismo.
\end{env}

\begin{corollary}[4.1.9]
\label{III.4.1.9}
Sob as hipóteses de \sref{III.4.1.4}, para todo subconjunto aberto
afim $V$ de $Y$, o homomorfismo canônico
\[
  \HH^n\bigl(
    \widehat{X}\cap f^{-1}(V),
    \widehat{\sh{F}}
  \bigr)
  \longrightarrow
  \Gamma\bigl(
    \widehat{Y}\cap V,
    \RR^n\widehat{f}_*(\widehat{\sh{F}})
  \bigr)
\]
é bijetivo.
\end{corollary}

\begin{env}[4.1.10]
\label{III.4.1.10}
\emph{Observação.}
Seja $f:X\to Y$ um morfismo de tipo finito entre esquemas ordinários
noetherianos, e seja $\sh{F}$ um $\sh{O}_X$-módulo coerente cujo
suporte é próprio sobre $Y$ \sref[II]{II.5.4.10}.
Sabemos então \sref{III.3.2.4} que $\RR^n f_*(\sh{F})$ é um
$\sh{O}_Y$-módulo coerente para todo $n\geq0$.
Além disso, sempre podemos supor que
\[
  \sh{F}=u_*(\sh{G}),
  \qquad
  \sh{G}=u^*(\sh{F}),
\]
onde $\sh{G}$ é um $\sh{O}_Z$-módulo coerente, $Z$ é um
subesquema fechado adequado de $X$ cujo espaço subjacente é
$\Supp(\sh{F})$, e $u:Z\to X$ é a injeção canônica
\sref[I]{I.9.3.5}.
Se
\[
  \sh{G}_k
  =
  \sh{G}\otimes_{\sh{O}_Y}
  (\sh{O}_Y/\sh{J}^{k+1}),
\]
então
\[
  \sh{G}_k=u^*(\sh{F}_k),
  \qquad
  \RR^n f_*(\sh{F}_k)
  =
  \RR^n(f\circ u)_*(\sh{G}_k)
\]
e
\[
  \RR^n f_*(\sh{F})
  =
  \RR^n(f\circ u)_*(\sh{G})
\]
por \sref{III.1.3.4}; e, tendo em conta \sref[I]{I.10.9.5},
\[
  \RR^n\widehat{f}_*(\widehat{\sh{F}})
  =
  \RR^n\widehat{(f\circ u)}_*(\widehat{\sh{G}}).
\]
Podemos, portanto, aplicar \sref{III.4.1.5} a $\sh{G}$ e ao morfismo
próprio $f\circ u$.
Segue-se que, sob estas hipóteses, os resultados de
\sref{III.4.1.5} continuam válidos para $\sh{F}$ e $f$.
\end{env}


\subsection{Casos particulares e variantes}
\label{subsection:III.4.2}

A forma mais útil do teorema de comparação \sref{III.4.1.5} é a
seguinte.

\begin{proposition}[4.2.1]
\label{III.4.2.1}
Seja $Y$ um esquema localmente noetheriano, seja $f:X\to Y$ um
morfismo próprio, e seja $\sh{F}$ um $\sh{O}_X$-módulo coerente.
Então, para todo $y\in Y$ e todo $p$,
\[
  \bigl(\RR^p f_*(\sh{F})\bigr)_y
\]

\oldpage[III]{130}

é um $\sh{O}_y$-módulo de tipo finito e está, portanto, separado
para a topologia $\mathfrak{m}_y$-pré-ádica.
Além disso, existe um isomorfismo topológico canônico
\[
  \widehat{\bigl(\RR^p f_*(\sh{F})\bigr)_y}
  \isoto
  \varprojlim_k
  \HH^p\left(
    f^{-1}(y),
    \sh{F}\otimes_{\sh{O}_Y}
    (\sh{O}_y/\mathfrak{m}_y^k)
  \right).
  \tag{4.2.1.1}
  \label{III.4.2.1.1}
\]
Aqui o primeiro termo é o completamento de
$\bigl(\RR^p f_*(\sh{F})\bigr)_y$ para a topologia
$\mathfrak{m}_y$-pré-ádica.
No segundo termo, para todo $k\geq0$, $f^{-1}(y)$ se considera
como o espaço subjacente do esquema
\[
  X\times_Y\Spec(\sh{O}_y/\mathfrak{m}_y^k)
\]
\sref[I]{I.3.6.1}.
\end{proposition}

\begin{proof}
Como $\sh{O}_y$ é um anel local noetheriano e
$\bigl(\RR^p f_*(\sh{F})\bigr)_y$ é um $\sh{O}_y$-módulo de tipo
finito \sref{III.3.2.1}, a topologia $\mathfrak{m}_y$-pré-ádica sobre
este módulo é separada \sref[0]{0.7.3.5}.

Quando $Y$ é noetheriano e $y$ é fechado, as demais afirmações se
seguem de \sref{III.4.1.7}, substituindo $Y$ por uma vizinhança afim de
$y$ e tomando $Y'=\{y\}$; aqui utilizamos (G,~II,~4.9.1).

Em geral, ponhamos
\[
  Y_1=\Spec(\sh{O}_y),
  \qquad
  X_1=X\times_Y Y_1,
  \qquad
  \sh{F}_1
  =
  \sh{F}\otimes_{\sh{O}_Y}\sh{O}_{Y_1},
\]
e seja
\[
  f_1=f\times 1_{Y_1}:X_1\longrightarrow Y_1.
\]
O esquema $Y_1$ é noetheriano, o morfismo $f_1$ é próprio
\sref[II]{II.5.4.2}[iii], e $\sh{F}_1$ é coerente
\sref[0]{0.5.3.11}.
Seja $y_1$ o único ponto fechado de $Y_1$.
A proposição já demonstrada se aplica a $f_1$, $\sh{F}_1$ e
$y_1$.
Temos
\[
  \sh{O}_{y_1}=\sh{O}_y,
  \qquad
  f_1^{-1}(y_1)=f^{-1}(y)
\]
\sref[I]{I.3.6.5}.
Além disso, os esquemas
\[
  X\times_Y\Spec(\sh{O}_y/\mathfrak{m}_y^k)
  \quad\text{e}\quad
  X_1\times_{Y_1}\Spec(\sh{O}_y/\mathfrak{m}_y^k)
\]
identificam-se canonicamente \sref[I]{I.3.3.9}, assim como
\[
  \sh{F}_1
  \otimes_{\sh{O}_{Y_1}}
  (\sh{O}_y/\mathfrak{m}_y^k)
  \quad\text{e}\quad
  \sh{F}
  \otimes_{\sh{O}_Y}
  (\sh{O}_y/\mathfrak{m}_y^k)
\]
\sref[I]{I.9.1.6}.
Resta observar que
\[
  \RR^p f_{1*}(\sh{F}_1)
  \simeq
  \RR^p f_*(\sh{F})
  \otimes_{\sh{O}_Y}
  \sh{O}_{Y_1}.
\]
Isto se segue de \sref{III.1.4.15}, pois o morfismo local
$\Spec(\sh{O}_y)\to Y$ é plano
\sref[0]{0.6.7.1}, \sref[I]{I.2.4.2}.
\end{proof}

O corolário seguinte utiliza a terminologia da teoria da
dimensão (capítulo~IV) e não se aplicará antes desse capítulo.

\begin{corollary}[4.2.2]
\label{III.4.2.2}
Seja $Y$ um esquema localmente noetheriano, seja $f:X\to Y$ um
morfismo próprio, seja $y\in Y$, e seja $r$ a dimensão de
$f^{-1}(y)$.
Então, para todo $\sh{O}_X$-módulo coerente $\sh{F}$, os feixes
$\RR^p f_*(\sh{F})$ se anulam em uma vizinhança de $y$ para todo
$p>r$.
\end{corollary}

\begin{proof}
Para todo $k$, temos
\[
  \HH^p\left(
    f^{-1}(y),
    \sh{F}\otimes_{\sh{O}_Y}
    (\sh{O}_y/\mathfrak{m}_y^k)
  \right)
  =
  0
\]
por (G,~II,~4.15.2).
De \sref{III.4.2.1} segue-se que o completamento separado de
$\bigl(\RR^p f_*(\sh{F})\bigr)_y$ para a topologia
$\mathfrak{m}_y$-pré-ádica é nulo.
Como esta topologia está separada, temos também
\[
  \bigl(\RR^p f_*(\sh{F})\bigr)_y=0.
\]
A conclusão se segue porque $\RR^p f_*(\sh{F})$ é coerente
\sref[0]{0.5.2.2}.
\end{proof}

\begin{env}[4.2.3]
\label{III.4.2.3}
O resultado \sref{III.4.2.1} é utilizado especialmente quando $p=0$.
Obtemos o corolário seguinte.
\end{env}

\begin{corollary}[4.2.4]
\label{III.4.2.4}
Sob as hipóteses de \sref{III.4.2.1}, existe um isomorfismo
topológico canônico
\[
  \widehat{\bigl(f_*(\sh{F})\bigr)_y}
  \isoto
  \varprojlim_k
  \Gamma\left(
    f^{-1}(y),
    \sh{F}_y/\mathfrak{m}_y^k\sh{F}_y
  \right).
\]
\end{corollary}


\subsection{Teorema de conexidade de Zariski}
\label{subsection:III.4.3}
\label{III.4.3}

Os resultados desta seção e da seguinte generalizam teoremas
bem conhecidos de Zariski, e todos podem ser deduzidos de
\sref{III.4.2.4}. São consequências do teorema seguinte.

\begin{theorem}[4.3.1]
\label{III.4.3.1}
\textnormal{(Teorema de conexidade).}
Seja $Y$ um esquema localmente noetheriano e seja $f:X\to Y$ um
morfismo próprio. Então
\[
  \sh{A}(X)=f_*(\sh{O}_X)
\]
é uma $\sh{O}_Y$-álgebra coerente.

\oldpage[III]{131}

Seja $Y'$ o $Y$-esquema finito tal que
$\sh{A}(Y')=\sh{A}(X)$; está determinado salvo $Y$-isomorfismo
\sref[II]{II.1.3.1}, \sref[II]{II.6.1.3}.
Se $f'=\sh{A}(e)$ é o $Y$-morfismo $X\to Y'$ deduzido do
isomorfismo identidade
\[
  e:\sh{A}(Y')\isoto\sh{A}(X)
\]
\sref[II]{II.1.2.7}, então $f'$ é próprio,
$f'_*(\sh{O}_X)$ é isomorfo a $\sh{O}_{Y'}$, e toda fibra
$f'^{-1}(y')$, $y'\in Y'$, é conexa e não vazia.
\end{theorem}

\begin{proof}
Seja $g:Y'\to Y$ o morfismo estrutural. Para demonstrar que o
homomorfismo
\[
  \theta:\sh{O}_{Y'}\longrightarrow f'_*(\sh{O}_X)
\]
que intervém na definição de $f'$ é bijetivo, basta, como
$Y'$ é afim sobre $Y$, demonstrar que
\[
  g_*(\theta):
  g_*(\sh{O}_{Y'})
  \longrightarrow
  g_*\bigl(f'_*(\sh{O}_X)\bigr)
  =
  f_*(\sh{O}_X)
\]
é a identidade \sref[II]{II.1.4.2}. Isto se segue das
definições, pois
$g_*(\sh{O}_{Y'})=\sh{A}(Y')$ e
$f_*(\sh{O}_X)=\sh{A}(X)$.
A coerência de $\sh{A}(X)$ é um caso particular do teorema de
finitude \sref{III.3.2.1}. Como $f$ é próprio e $g$ é separado,
$f'$ é próprio \sref[II]{II.5.4.3}[i]. Resta demonstrar
\sref{III.4.3.2}.
\end{proof}

\begin{corollary}[4.3.2]
\label{III.4.3.2}
Sob as hipóteses de \sref{III.4.3.1}, suponhamos além disso que
$f_*(\sh{O}_X)$ é isomorfo a $\sh{O}_Y$. Então toda fibra
$f^{-1}(y)$ de $f$ é conexa e não vazia.
\end{corollary}

\begin{proof}
A hipótese de que $f_*(\sh{O}_X)$ é isomorfo a $\sh{O}_Y$ implica
que $f$ já é dominante, logo sobrejetivo porque $f$ é uma
aplicação fechada. Como em \sref{III.4.2.1}, podemos reduzir-nos ao
caso em que $y$ é fechado em $Y$. Como $f^{-1}(y)$ é um espaço
noetheriano, tem um número finito de componentes conexas; estas são
também as componentes conexas do espaço subjacente do
completamento $\widehat{X}$ ao longo de $f^{-1}(y)$.

Se $Z_i$, $1\leq i\leq n$, são estas componentes conexas, então
\[
  \Gamma(\widehat{X},\sh{O}_{\widehat{X}})
  =
  \prod_{i=1}^{n}\Gamma(Z_i,\sh{O}_{\widehat{X}}),
\]
e nenhum dos anéis do membro direito é nulo, pois a seção
unidade não é nula em nenhum ponto de $\widehat{X}$. Aplicando
\sref{III.4.1.5} a $\sh{F}=\sh{O}_X$, cujo completamento ao longo de
$f^{-1}(y)$ é $\sh{O}_{\widehat{X}}$, vê-se que
$\Gamma(\widehat{X},\sh{O}_{\widehat{X}})$ é isomorfo ao
completamento $\mathfrak{m}_y$-ádico separado
$\widehat{\sh{O}}_y$ do anel local $\sh{O}_y$. É, portanto, um
anel local e não pode ser um produto direto de vários anéis não
nulos. Assim, $n=1$.
\end{proof}

\begin{corollary}[4.3.3]
\label{III.4.3.3}
Sob as hipóteses de \sref{III.4.3.1}, para todo $y\in Y$ o conjunto
de componentes conexas da fibra $f^{-1}(y)$ está em correspondência
bijetiva com o conjunto finito de pontos da fibra
$g^{-1}(y)$, onde $g:Y'\to Y$ é o morfismo estrutural; em outros
termos, corresponde ao conjunto de ideais maximais de
$\bigl(f_*(\sh{O}_X)\bigr)_y$.
\end{corollary}

\begin{proof}
Como $Y'$ é finito sobre $Y$, o espaço $g^{-1}(y)$ é finito e
discreto \sref[II]{II.6.1.7}. Como
\[
  f^{-1}(y)=f'^{-1}\bigl(g^{-1}(y)\bigr),
\]
a afirmação se segue desta observação e de
\sref{III.4.3.1}.
\end{proof}

Isto dá uma interpretação notável do $Y$-esquema $Y'$ definido em
\sref{III.4.3.1}. A fatoração
$f=g\circ f'$ do morfismo próprio $f$ é análoga à fatoração
obtida por K.~Stein para aplicações holomorfas de espaços
analíticos, e será chamada daqui em diante de \emph{fatoração de Stein}
de $f$.

\begin{remark}[4.3.4]
\label{III.4.3.4}
Seja $k$ uma extensão do corpo $\kres(y)$. Se
\[
  f^{-1}(y)\otimes_{\kres(y)}k
  =
  X\times_Y\Spec(k)
\]
é conexo, então $f^{-1}(y)$ também é conexo, pois é a imagem
do primeiro por um morfismo de projeção
\sref[I]{I.3.4.7}. Para um morfismo $f:X\to Y$ de esquemas e um ponto
$y\in Y$, diremos que a fibra $f^{-1}(y)$ é
\emph{geometricamente conexa} se, para toda extensão $k$ de
$\kres(y)$, o esquema
$f^{-1}(y)\otimes_{\kres(y)}k=X\times_Y\Spec(k)$ é conexo.

\oldpage[III]{132}

Sob as hipóteses de \sref{III.4.3.2}, a conclusão pode
ser reforçada: as fibras $f^{-1}(y)$ são geometricamente conexas.
Com efeito, para toda extensão $k$ de $\kres(y)$ existem um anel
local noetheriano $A$ e um homomorfismo local
\[
  \varphi:\sh{O}_y\longrightarrow A
\]
tal que $A$ é um $\sh{O}_y$-módulo plano e o corpo residual de
$A$ é $\kres(y)$-isomorfo a $k$ \sref[0]{0.10.3.1}.
Seja $Y_1=\Spec(A)$, e seja $h:Y_1\to Y$ o morfismo local
correspondente a $\varphi$, que leva o único ponto fechado $y_1$ de
$Y_1$ a $y$ \sref[I]{I.2.4.1}. Ponhamos
\[
  X_1=X\times_Y Y_1,
  \qquad
  f_1=f\times 1_{Y_1}.
\]
O esquema $Y_1$ é noetheriano, $f_1$ é próprio
\sref[II]{II.5.4.2}, e $f_1^{-1}(y_1)$ é um
$\kres(y_1)$-esquema isomorfo a $X\times_Y\Spec(k)$. Basta, portanto,
demonstrar que
\[
  f_{1*}(\sh{O}_{X_1})=\sh{O}_{Y_1},
\]
de modo que \sref{III.4.3.2} possa ser aplicado a $f_1$. Agora $h$ é
plano por \sref[I]{I.2.4.2}; portanto, por \sref{III.1.4.15} com
$q=0$,
\[
  f_{1*}(\sh{O}_{X_1})
  =
  h^*\bigl(f_*(\sh{O}_X)\bigr)
  =
  h^*(\sh{O}_Y)
  =
  \sh{O}_{Y_1}.
\]

No caso geral de \sref{III.4.3.1}, o mesmo argumento mostra,
com a notação desse teorema, que
\[
  f_{1*}(\sh{O}_{X_1})
  =
  h^*\bigl(g_*(\sh{O}_{Y'})\bigr),
\]
e que a fatoração de Stein $f_1=g_1\circ f'_1$ de $f_1$
satisfaz
\[
  g_1=g\times 1_{Y_1}
\]
\sref[II]{II.1.5.2}, sendo o $Y_1$-esquema finito correspondente
$Y'_1=Y'\times_Y Y_1$. Tendo em conta a transitividade das
fibras \sref[I]{I.3.6.4}, segue-se de
\sref{III.4.3.3} que o número de componentes conexas de
$f_1^{-1}(y_1)$ é o número de elementos de
\[
  g_1^{-1}(y_1)
  =
  g^{-1}(y)\otimes_{\kres(y)}k.
\]
Se $k$ for escolhido algebricamente fechado, este número será independente
da extensão algebricamente fechada escolhida e igual ao
\emph{número geométrico de pontos} de $g^{-1}(y)$
\sref[I]{I.6.4.7}, ou, equivalentemente, a
\[
  \sum_i[\kres(y'_i):\kres(y)]_s,
\]
onde os $y'_i$ percorrem o conjunto finito $g^{-1}(y)$. Este número
chama-se \emph{número geométrico de componentes conexas} de
$f^{-1}(y)$. Os corpos $\kres(y'_i)$ são precisamente os corpos
residuais do anel semilocal $\bigl(f_*(\sh{O}_X)\bigr)_y$.
\end{remark}

\begin{proposition}[4.3.5]
\label{III.4.3.5}
Sejam $X$ e $Y$ esquemas íntegros localmente noetherianos, e seja
$f:X\to Y$ um morfismo próprio dominante. Para todo $y\in Y$, o
número de componentes conexas de $f^{-1}(y)$ é no máximo o número
de ideais maximais do fecho integral $\sh{O}'_y$ de
$\sh{O}_y$ no corpo de funções racionais $\mathrm{R}(X)$.
\end{proposition}

\begin{proof}
Para todo aberto $U$ de $Y$,
\[
  \Gamma\bigl(U,f_*(\sh{O}_X)\bigr)
  =
  \Gamma\bigl(f^{-1}(U),\sh{O}_X\bigr)
\]
é a interseção dos anéis locais $\sh{O}_x$ para
$x\in f^{-1}(U)$ \sref[I]{I.8.2.1.1}. Segue-se imediatamente que
$\bigl(f_*(\sh{O}_X)\bigr)_y$ é um subanel de $\mathrm{R}(X)$ que
contém $\sh{O}_y$. Como $f_*(\sh{O}_X)$ é um
$\sh{O}_Y$-módulo coerente, $\bigl(f_*(\sh{O}_X)\bigr)_y$ é um
$\sh{O}_y$-módulo finito e, portanto, está contido em $\sh{O}'_y$.
Todo ideal maximal de tal anel é a interseção deste anel com
um ideal maximal de $\sh{O}'_y$
\cite[vol.~I, pp.~257 and 259]{I-13}, o que demonstra a proposição.
\end{proof}

\begin{definition}[4.3.6]
\label{III.4.3.6}
Um anel local íntegro se diz \emph{unirrama} se seu fecho integral
é novamente um anel local. Um ponto $y$ de um esquema íntegro
$Y$ se diz \emph{unirrama} se o anel local $\sh{O}_y$ é
unirrama; este é, em particular, o caso quando $Y$ é normal em
$y$.
\end{definition}

Seja $A$ um anel local íntegro e seja $K$ seu corpo de frações. O
anel $A$ é unirrama se e somente se todo subanel $A_1$ de $K$ que
contenha $A$ e seja uma $A$-álgebra finita é local. Com efeito, seja
$A'$ o fecho integral de $A$. Pelo primeiro teorema de
Cohen--Seidenberg (Bourbaki, \emph{Alg\`ebre commutative}, capítulo~V,
\S2, n.º~1, teorema~1), todo ideal maximal de $A_1$ é a contração
a $A_1$ de um ideal maximal de $A'$; assim, se $A'$ é local, também
o é $A_1$.

\oldpage[III]{133}

Reciprocamente, $A'$ é o limite indutivo da família dirigida de
$A$-subálgebras finitas $A_\alpha$ de $A'$. Se cada $A_\alpha$ é
local, então, para $A_\alpha\subset A_\beta$, o ideal maximal de
$A_\alpha$ é a contração do ideal maximal de $A_\beta$, pelo
mesmo argumento; portanto, $A'$ é local
\sref[0]{0.10.3.1.3}.

Observemos também que, se o completamento de um anel local
noetheriano $A$ é íntegro ---propriedade que se exprime dizendo que
$A$ é \emph{analiticamente íntegro}---, então $A$ é unirrama.
Seja $\mathfrak{m}$ o ideal maximal de $A$, seja $K$ seu corpo de
frações, e seja $K'$ o corpo de frações de $\widehat{A}$;
então
\[
  K'=K\otimes_A\widehat{A}.
\]
Seja $A'_F$ uma $A$-subálgebra finita de $K$. O subanel $B_F$ de
$K'$ gerado por $\widehat{A}$ e $A'_F$ é isomorfo a
\[
  A'_F\otimes_A\widehat{A};
\]
é um $\widehat{A}$-módulo finito e é o completamento de $A'_F$
para a topologia $\mathfrak{m}$-ádica
\sref[0]{0.7.3.3}, \sref[0]{0.7.3.6}. Como $A'_F$ é semilocal
(Bourbaki, \emph{Alg\`ebre commutative}, capítulo~IV, \S2, n.º~5,
corolário~3 da proposição~9), enquanto seu completamento é
íntegro, $A'_F$ só pode ter um ideal maximal $\mathfrak{m}'_F$, e
$\mathfrak{m}'_F\cap A=\mathfrak{m}$. Isto demonstra a afirmação.

\begin{corollary}[4.3.7]
\label{III.4.3.7}
Sob as hipóteses de \sref{III.4.3.5}, suponhamos que o fecho
algébrica de $\mathrm{R}(Y)$ em $\mathrm{R}(X)$ tem grau separável
$n$ e que $y\in Y$ é unirrama. Então a fibra $f^{-1}(y)$ tem no
máximo $n$ componentes conexas. Em particular, se o fecho
algébrica de $\mathrm{R}(Y)$ em $\mathrm{R}(X)$ é radicial sobre
$\mathrm{R}(Y)$, então $f^{-1}(y)$ é conexa.
\end{corollary}

\begin{proof}
Seja $\sh{O}'_y$ o fecho integral de $\sh{O}_y$, e seja
$\sh{O}''_y$ o fecho integral de $\sh{O}_y$ em $\mathrm{R}(X)$.
Esta última é também o fecho integral de $\sh{O}'_y$ em
$\mathrm{R}(X)$. Se $\sh{O}'_y$ é local, então $\sh{O}''_y$ é
um anel semilocal que tem no máximo $n$ ideais maximais
\cite[vol.~I, p.~289, Theorem~22]{I-13}.

Este corolário é essencialmente a forma na qual Zariski enuncia seu
``teorema de conexidade'' para esquemas algébricos.
\end{proof}

\begin{remark}[4.3.8]
\label{III.4.3.8}
Se, além das hipóteses de \sref{III.4.3.7}, $Y$ é normal em
$y$, então a fibra $f^{-1}(y)$ é geometricamente conexa: com a
notação de \sref{III.4.3.4}, $g^{-1}(y)$ é um único ponto $y'$, e
$\kres(y')$ é radicial sobre $\kres(y)$.
\end{remark}

\begin{definition}[4.3.9]
\label{III.4.3.9}
Seja $Y$ um esquema localmente noetheriano. Um morfismo de tipo finito
$f:X\to Y$ se diz \emph{universalmente aberto} se, para todo esquema
irredutível localmente noetheriano $Y'$ e todo morfismo dominante
$g:Y'\to Y$, toda componente irredutível de
$X'=X\times_Y Y'$ domina $Y'$.
\end{definition}

Se $Y$ é irredutível, isto significa que, se $\eta$ e $\eta'$ são
os pontos genéricos de $Y$ e $Y'$, respectivamente, e se
$f'=f\times 1_{Y'}$, toda componente irredutível de $X'$ tem seu
ponto genérico em $f'^{-1}(\eta')$. Consequentemente, para todo
aberto $U$ de $Y$, o morfismo restrito $f^{-1}(U)\to U$ é
universalmente aberto.

\begin{corollary}[4.3.10]
\label{III.4.3.10}
Sejam $X$ e $Y$ esquemas íntegros localmente noetherianos, e seja
$f:X\to Y$ próprio, dominante e universalmente aberto. Se o fecho
algébrica de $\mathrm{R}(Y)$ em $\mathrm{R}(X)$ é radicial sobre
$\mathrm{R}(Y)$, então toda fibra $f^{-1}(y)$, $y\in Y$, é
geometricamente conexa.
\end{corollary}

\begin{proof}
Podemos reduzir-nos ao caso $Y=\Spec(B)$, onde $B$ é um anel
íntegro noetheriano. Por \sref[II]{II.7.1.7}, existe um anel local
noetheriano integralmente fechado $A$ que domina $\sh{O}_y$ e tem
$\mathrm{R}(Y)$ como corpo de frações. Ponhamos $Y'=\Spec(A)$, e
seja $h:Y'\to Y$ o morfismo correspondente à injeção canônica
$B\to A$. É birracional, logo dominante; além disso, se $y'$ é o
único ponto fechado de $Y'$, então $h(y')=y$.

Ponhamos
\[
  X'=X\times_Y Y',
  \qquad
  f'=f\times 1_{Y'}.
\]
Sejam $\eta$, $\eta'$ e $\xi$ os pontos genéricos de $Y$, $Y'$ e
$X$, respectivamente. Então $f(\xi)=\eta$ e
$h^{-1}(\eta)=\{\eta'\}$; além disso,
$\kres(\eta)=\kres(\eta')=\mathrm{R}(Y)$. Assim, $f'^{-1}(\eta')$ é
isomorfo a $f^{-1}(\eta)$ \sref[I]{I.3.6.4}, e tem um único ponto
genérico, pois $\xi$ é o ponto genérico de
$f^{-1}(\eta)$ \sref[0]{0.2.1.8}. Pela abertura universal, toda
componente irredutível de $X'$ tem seu ponto genérico em
$f'^{-1}(\eta')$. Portanto, $X'$ é irredutível; seu ponto genérico
$\xi'$ é o de $f'^{-1}(\eta')$, e
$\kres(\xi')=\kres(\xi)$.

\oldpage[III]{134}

Ponhamos $X''=X'_{\mathrm{red}}$ e $f''=f'_{\mathrm{red}}$. Então
$X''$ é íntegro e noetheriano, $f''$ é próprio
\sref[II]{II.5.4.6}, e os espaços subjacentes das fibras
$f'^{-1}(y')$ e $f''^{-1}(y')$ são iguais. Além disso,
\[
  \mathrm{R}(X'')=\kres(\xi')=\mathrm{R}(X).
\]
Assim, $f''$ satisfaz as hipóteses de \sref{III.4.3.8}, e
$f''^{-1}(y')$ é geometricamente conexa.

Seja agora $k$ uma extensão qualquer de $\kres(y)$. Existe uma
extensão $k_1$ de $\kres(y)$ na qual tanto $\kres(y')$ como $k$
podem ser vistos como subextensões (Bourbaki, \emph{Alg\`ebre},
capítulo~V, \S4, proposição~2). Pelo parágrafo anterior,
\[
  f''^{-1}(y')\times_{Y'}\Spec(k_1)
\]
é conexo. Tem o mesmo esquema reduzido associado que
$f'^{-1}(y')\times_{Y'}\Spec(k_1)$
\sref[I]{I.5.1.8}, logo este último é conexo; é isomorfo a
$f^{-1}(y)\times_Y\Spec(k_1)$ \sref[I]{I.3.6.4}. Portanto,
$f^{-1}(y)\times_Y\Spec(k)$ é conexo, e a observação do início
de \sref{III.4.3.4} termina a demonstração.
\end{proof}

\begin{remark}[4.3.11]
\label{III.4.3.11}
\emph{(i)}
O argumento anterior se deve, no essencial, a Zariski
\cite{I-20}, salvo que, em seu contexto, pode-se tomar como $A$ o
fecho integral de $\sh{O}_y$, pois este é um anel noetheriano para
os anéis locais da geometria algébrica clássica. Zariski
demonstra também que, se $Y$ é a variedade de Chow de um espaço
projetivo $\mathbf{P}_k^r$ sobre um corpo $k$, e $X$ é o
subconjunto fechado de $\mathbf{P}_k^r\times_kY$ que define a
correspondência de Chow entre $\mathbf{P}_k^r$ e $Y$, então a
projeção $X\to Y$ é universalmente aberta (\emph{loc.\ cit.},
lema da p.~82).
Parece que esta é a única propriedade formal das ``coordenadas de
Chow'' utilizada em certas aplicações. Em tal situação é, por
isso, útil substituir a linguagem da especialização de ciclos no
espaço projetivo pela linguagem das fibras de um morfismo próprio,
eventualmente suposto universalmente aberto ou submetido a condições
análogas de regularidade local.

\emph{(ii)}
No capítulo~IV veremos que um morfismo universalmente aberto
$f:X\to Y$ pode definir-se equivalentemente exigindo que, para todo
morfismo $Z\to Y$, o morfismo obtido por mudança de base
$f_{(Z)}:X_{(Z)}\to Z$ seja aberto. Também pode demonstrar-se que, se
$f$ satisfaz as hipóteses de \sref{III.4.3.10}, e se $y$ é uma
especialização de $y'$, então o número geométrico de componentes
conexas de $f^{-1}(y)$ é no máximo o de $f^{-1}(y')$.
\end{remark}

\begin{corollary}[4.3.12]
\label{III.4.3.12}
Sob as hipóteses de \sref{III.4.3.5}, suponhamos além disso que
$\mathrm{R}(Y)$ é algebricamente fechado em $\mathrm{R}(X)$, e seja
$y$ um ponto normal de $Y$. Então $f^{-1}(y)$ é geometricamente
conexa, e existe uma vizinhança aberta $U$ de $y$ em $Y$ tal que
\[
  f_*\bigl(\sh{O}_X|_{f^{-1}(U)}\bigr)
  \simeq
  \sh{O}_Y|_U.
\]
Em particular, se $Y$ é normal e $\mathrm{R}(Y)$ é algebricamente
fechado em $\mathrm{R}(X)$, então
$f_*(\sh{O}_X)\simeq\sh{O}_Y$.
\end{corollary}

\begin{proof}
A afirmação relativa a $f^{-1}(y)$ é um caso particular de
\sref{III.4.3.8}.

\oldpage[III]{135}

Se
\[
  X\xrightarrow{f'}Y'\xrightarrow{g}Y
\]
é a fatoração de Stein de $f$ \sref{III.4.3.3}, segue-se que
$g^{-1}(y)$ consiste em um único ponto $y'$. Além disso,
\[
  \sh{O}_y
  \subset
  \sh{O}_{y'}
  =
  \bigl(f_*(\sh{O}_X)\bigr)_y
  \subset
  \mathrm{R}(X).
\]
Como $\sh{O}_{y'}$ é finito sobre $\sh{O}_y$, e portanto
\emph{a fortiori} algébrico sobre $\mathrm{R}(Y)$ dentro do corpo
de frações pertinente, a hipótese obriga que esteja contido em
$\mathrm{R}(Y)$. Como $y$ é normal, necessariamente
$\sh{O}_{y'}=\sh{O}_y$. Assim, $g$ é um isomorfismo local em $y'$
\sref[I]{I.6.5.4}, o que demonstra a primeira parte. A segunda
afirmação se segue da primeira: a hipótese adicional faz com que $g$
seja bijetivo e um isomorfismo local em uma vizinhança de cada ponto de
$Y'$, logo um isomorfismo.
\end{proof}

O fato de que \sref{III.4.3.7} esteja disponível para esquemas fornece,
por exemplo, a aplicação seguinte.

\begin{proposition}[4.3.13]
\label{III.4.3.13}
Seja $A$ um anel local noetheriano unirrama, seja $\mathfrak{a}$ um
ideal de definição de $A$, e sejam
\[
  A_0=A/\mathfrak{a},
  \qquad
  S=\operatorname{gr}_{\mathfrak{a}}(A)
\]
o anel graduado associado à filtração $\mathfrak{a}$-pré-ádica.
Suponhamos que $S$ é uma $A_0$-álgebra graduada gerada por $S_1$,
onde $S_1$ é um $A_0$-módulo finito. Então $\Proj(S)$ é um
$A_0$-esquema conexo.
\end{proposition}

\begin{proof}
Seja $\mathfrak{m}$ o ideal maximal de $A$. O esquema
$Y=\Spec(A)$ é íntegro, e o ponto $y$ correspondente a
$\mathfrak{m}$ é seu único ponto fechado. Para certo inteiro $p$,
\[
  \mathfrak{m}^p\subset\mathfrak{a}\subset\mathfrak{m},
\]
de modo que $V(\mathfrak{a})=\{\mathfrak{m}\}$. Ponhamos
\[
  S'=\bigoplus_{n\geq0}\mathfrak{a}^n
  \qquad\text{e}\qquad
  X=\Proj(S').
\]
Este é o $Y$-esquema obtido pela explosão do ideal
$\mathfrak{a}$. É íntegro, e o morfismo estrutural
$f:X\to Y$ é birracional \sref[II]{II.8.1.4} e projetivo.
Portanto, \sref{III.4.3.7} se aplica e mostra que $f^{-1}(y)$ é
conexo. Mas o espaço subjacente de $f^{-1}(y)$ é o de
\[
  \Proj(S'\otimes_A A_0)
\]
\sref[I]{I.3.6.1}, \sref[II]{II.2.8.10}; como
$S'\otimes_AA_0=S$ por definição, segue-se a proposição.
\end{proof}


\subsection{O ``teorema principal'' de Zariski}
\label{subsection:III.4.4}

\begin{proposition}[4.4.1]
\label{III.4.4.1}
Seja $Y$ um esquema localmente noetheriano e seja $f:X\to Y$ um
morfismo próprio. Seja $X'$ o conjunto de pontos $x\in X$ que estão
isolados em sua fibra $f^{-1}(f(x))$. Então $X'$ é aberto em $X$.
Se
\[
  f=g\circ f'
\]
é a fatoração de Stein de $f$ \sref{III.4.3.3}, a restrição
de $f'$ a $X'$ é um isomorfismo de $X'$ sobre o subesquema induzido
em um subconjunto aberto $U$ de $Y'$, e
$X'=f'^{-1}(U)$.
\end{proposition}

\begin{proof}
Como $g^{-1}(f(x))$ é finito e discreto
\sref{III.4.3.3}, \sref[II]{II.6.1.7}, o ponto $x$ está isolado em
$f^{-1}(f(x))$ se e somente se está isolado em
$f'^{-1}(f'(x))$. Podemos, portanto, reduzir-nos ao caso $f'=f$, de
modo que $f_*(\sh{O}_X)=\sh{O}_Y$. Se $x\in X'$, a fibra
$f^{-1}(f(x))$, que é conexa por \sref{III.4.3.2}, deve constar
então do único ponto $x$.

Como $f$ é fechado, para toda vizinhança aberta $V$ de $x$ em $X$, o
conjunto fechado $f(X-V)$ não contém $y=f(x)$. Se $U$ é seu
complementar em $Y$, então $f^{-1}(U)\subset V$. Assim, as
imagens inversas por $f$ de um sistema fundamental de vizinhanças
abertas de $y$ formam um sistema fundamental de vizinhanças abertas de
$x$. A igualdade $f_*(\sh{O}_X)=\sh{O}_Y$ e a definição da
imagem direta de um feixe \sref[0]{0.3.4.1}, \sref{III.4.2.1} implicam
agora que, se $f=(\psi,\theta)$, o homomorfismo
\[
  \theta_y^\#:\sh{O}_y\longrightarrow\sh{O}_x
\]
é um isomorfismo. Portanto, existem vizinhanças abertas $V$ de $x$ e
$U$ de $y$ tais que a restrição de $f$ é um isomorfismo
$V\isoto U$ \sref[I]{I.6.5.4}.

\oldpage[III]{136}

Pelo parágrafo anterior podemos dispor além disso que
$f^{-1}(U)=V$. Segue-se imediatamente que $V\subset X'$, o que
termina a demonstração.
\end{proof}

A proposição seguinte foi demonstrada por Chevalley para esquemas
algébricos.

\begin{proposition}[4.4.2]
\label{III.4.4.2}
Seja $Y$ um esquema localmente noetheriano e seja $f:X\to Y$ um
morfismo. As condições seguintes são equivalentes:
\begin{enumerate}
  \item[{\rm a)}] $f$ é finito;
  \item[{\rm b)}] $f$ é afim e próprio;
  \item[{\rm c)}] $f$ é próprio e, para todo $y\in Y$, o conjunto
    $f^{-1}(y)$ é finito.
\end{enumerate}
\end{proposition}

\begin{proof}
A condição~a) implica b) por \sref[II]{II.6.1.2} e
\sref[II]{II.6.1.11}. Se $f$ é próprio e afim, também o é
\[
  f^{-1}(y)\longrightarrow\Spec(\kres(y))
\]
\sref[II]{II.1.6.2}, \sref[II]{II.5.4.2}. O teorema de finitude
\sref{III.3.2.1}, aplicado ao feixe estrutural de $f^{-1}(y)$, mostra
que $f^{-1}(y)=\Spec(A)$ para uma $\kres(y)$-álgebra finita $A$. Assim,
$f^{-1}(y)$ é um conjunto finito \sref[II]{II.6.1.7}, e b) implica
c).

Por fim, se c) for satisfeita, a fibra $f^{-1}(y)$ é um esquema
algébrico sobre $\kres(y)$ cujo conjunto subjacente é finito e é,
portanto, discreto \sref[I]{I.6.4.4}. Com a notação de
\sref{III.4.4.1}, temos $X'=X$, logo
$f':X\to Y'$ é um isomorfismo. Como $g$ é finito, c) implica a).
\end{proof}

\begin{theorem}[4.4.3]
\label{III.4.4.3}
\textnormal{(``Teorema principal'' de Zariski).}
Seja $Y$ um esquema noetheriano, seja $f:X\to Y$ um morfismo
quase-projetivo, e seja $X'$ o conjunto de pontos $x\in X$ que estão
isolados em sua fibra $f^{-1}(f(x))$. Então $X'$ é aberto em $X$,
e o subesquema induzido por $X$ sobre $X'$ é isomorfo a um
subesquema induzido sobre um subconjunto aberto de um $Y$-esquema
$Y'$ finito sobre $Y$.
\end{theorem}

\begin{proof}
Existe um $Y$-esquema projetivo $Z$ tal que $X$ é
$Y$-isomorfo ao subesquema induzido sobre um subconjunto aberto de
$Z$ \sref[II]{II.5.3.2}, \sref[II]{II.5.5.1}. Reduzimo-nos assim ao
caso em que $f$ é projetivo, logo próprio \sref[II]{II.5.5.3}, e o
teorema se segue imediatamente de
\sref{III.4.4.1}.
\end{proof}

\begin{remark}[4.4.4]
\label{III.4.4.4}
Se $X$ é reduzido (respectivamente, se $X$ é irredutível e $X'$ não
é vazio), pode exigir-se em \sref{III.4.4.3} que $Y'$ seja reduzido
(respectivamente, irredutível). Com efeito, $Y'$ pode ser substituído pelo
fecho esquemático $\overline{X'}$ de $X'$ em $Y'$
\sref[I]{I.9.5.11}, \sref[II]{II.6.1.5}. Se $X'$ é reduzido, também
o é $\overline{X'}$ \sref[I]{I.9.5.9}; se $X'$ não é vazio e é
irredutível, então $\overline{X'}$ também é irredutível.
\end{remark}

\begin{corollary}[4.4.5]
\label{III.4.4.5}
Seja $Y$ um esquema localmente noetheriano, seja $f:X\to Y$ um
morfismo de tipo finito, e seja $x$ um ponto de $X$ isolado em sua fibra
$f^{-1}(f(x))$. Então existe uma vizinhança aberta de $x$ em $X$
isomorfa a um subconjunto aberto de um $Y$-esquema finito sobre $Y$.
\end{corollary}

\begin{proof}
Ponhamos $y=f(x)$, seja $U$ uma vizinhança aberta afim de $y$ em $Y$, e
seja $V$ uma vizinhança aberta afim de $x$ em $X$, contida em
$f^{-1}(U)$. Como $Y$ é separado, a inclusão $U\to Y$ é afim
\sref[II]{II.1.6.3}; e como $V$ é afim sobre $U$ (\emph{ibid.}), a
restrição de $f$ a $V$ é um morfismo afim $V\to Y$
\sref[II]{II.1.6.2}. A fortiori, esta restrição é
quase-projetiva, pois é de tipo finito
\sref[I]{I.6.3.5}, \sref[II]{II.5.3.4}. Resta aplicar
\sref{III.4.4.3} a esta restrição.
\end{proof}

O Corolário \sref{III.4.4.5} pode ser enunciado na linguagem da
álgebra comutativa como segue.

\begin{corollary}[4.4.6]
\label{III.4.4.6}
Seja $A$ um anel noetheriano, seja $B$ uma $A$-álgebra de tipo
finito, seja $\mathfrak q$ um ideal primo de $B$, e seja $\mathfrak p$
sua imagem inversa em $A$. Suponhamos que $\mathfrak q$ é ao mesmo tempo
maximal e minimal no conjunto de ideais primos de $B$ cuja imagem
inversa é $\mathfrak p$. Então existem um elemento
$g\in B-\mathfrak q$, uma $A$-álgebra finita $A'$, e um elemento
$f'\in A'$ tais que as $A$-álgebras $B_g$ e $A'_{f'}$ são
isomorfas.
\end{corollary}

\begin{proof}
Apliquemos \sref{III.4.4.5} a $Y=\Spec(A)$ e $X=\Spec(B)$. A
hipótese sobre $\mathfrak q$ diz precisamente que $\mathfrak q$
está isolado em sua fibra \sref[I]{I.1.1.7}.
\end{proof}

A consequência seguinte parece menos geral.

\oldpage[III]{137}

\begin{corollary}[4.4.7]
\label{III.4.4.7}
Seja $A$ um anel local noetheriano, seja $B$ uma $A$-álgebra de tipo
finito, e seja $\mathfrak n$ um ideal primo de $B$ cuja imagem inversa
em $A$ é o ideal maximal $\mathfrak m$. Suponhamos que $\mathfrak n$
é maximal em $B$ e minimal no conjunto de ideais primos de $B$
cuja imagem inversa é $\mathfrak m$ (o que significa também que
$\mathfrak mB$ é $\mathfrak n$-primário). Então existem uma
$A$-álgebra finita $A'$ e um ideal maximal $\mathfrak m'$ de $A'$
(cuja imagem inversa em $A$ é $\mathfrak m$) tais que $B_{\mathfrak
n}$ é isomorfo, como $A$-álgebra, a $A'_{\mathfrak m'}$.
\end{corollary}

O caso particular seguinte de \sref{III.4.4.7} também é chamado, às
vezes, de ``teorema principal''.

\begin{corollary}[4.4.8]
\label{III.4.4.8}
Sob as hipóteses de \sref{III.4.4.7}, suponhamos além disso que $A$ e
$B$ são domínios íntegros com o mesmo corpo de frações $K$.
Se $A$ é integralmente fechado, então $A=B$.
\end{corollary}

\begin{proof}
A Observação \sref{III.4.4.4} mostra que, ao aplicar
\sref{III.4.4.7}, podemos supor que $A'$ é um domínio íntegro com
corpo de frações $K$. A hipótese sobre $A$ dá então
$A'=A$, logo $B_{\mathfrak n}=A$. Como
$A\subset B\subset B_{\mathfrak n}$, segue-se que $A=B$.
\end{proof}

O enunciado \sref{III.4.4.8} é a forma na qual Zariski enunciou seu
``teorema principal'' (estendido aqui a anéis locais íntegros
noetherianos arbitrários).

Os corolários anteriores eram variantes locais de
\sref{III.4.4.3}, que é um resultado global. Eis outra
consequência de natureza global.

\begin{corollary}[4.4.9]
\label{III.4.4.9}
Seja $Y$ um esquema íntegro localmente noetheriano, e seja
$f:X\to Y$ um morfismo separado de tipo finito que é birracional.
Suponhamos além disso que $Y$ é normal e que toda fibra $f^{-1}(y)$,
para $y\in Y$, é finita. Então $f$ é uma imersão aberta. Se,
além disso, $f$ é fechado (em particular, se $f$ é próprio), então
$f$ é um isomorfismo.
\end{corollary}

\begin{proof}
Seja $x\in X$ e ponhamos $y=f(x)$. Como $f^{-1}(y)$ é um esquema
algébrico sobre $\kres(y)$, a hipótese de que é finito implica que
é discreto \sref[I]{I.6.4.4}. Além disso, $\sh{O}_y$ é integralmente
fechado, e $\sh{O}_x$ e $\sh{O}_y$ têm o mesmo corpo de frações
\sref[I]{I.7.1.5}. Podemos, portanto, aplicar
\sref{III.4.4.8}; se $f=(\psi,\theta)$, o homomorfismo
\[
  \theta_y^\#:\sh{O}_y\longrightarrow\sh{O}_x
\]
é bijetivo. De \sref[I]{I.6.5.4} segue-se que $f$ é um isomorfismo
local. Como $f$ é separado e $X$ é íntegro, $f$ é uma imersão
aberta \sref[I]{I.8.2.8}. A última afirmação se segue do fato de
que $f$ é dominante.
\end{proof}

\begin{proposition}[4.4.10]
\label{III.4.4.10}
Seja $Y$ um esquema localmente noetheriano, e seja $f:X\to Y$ um
morfismo localmente de tipo finito. O conjunto $X'$ de pontos
$x\in X$ que estão isolados em sua fibra $f^{-1}(f(x))$ é aberto em
$X$.
\end{proposition}

\begin{proof}
A questão é local sobre $X$ e $Y$, pelo que podemos supor que
$X$ e $Y$ são afins noetherianos e que $f$ é de tipo finito. O
morfismo $f$ é então afim e de tipo finito, logo
quase-projetivo \sref[II]{II.5.3.4}; basta aplicar
\sref{III.4.4.3}.
\end{proof}

\oldpage[III]{138}

\begin{corollary}[4.4.11]
\label{III.4.4.11}
Seja $Y$ um esquema localmente noetheriano, e seja $f:X\to Y$ um
morfismo próprio. O conjunto $U$ de pontos $y\in Y$ para os quais
$f^{-1}(y)$ é discreto é aberto em $Y$, e o morfismo restrito
$f^{-1}(U)\to U$ é finito. Em particular, todo morfismo próprio
quase-finito $X\to Y$ é finito.
\end{corollary}

\begin{proof}
O complementar de $U$ em $Y$ é a imagem por $f$ de $X-X'$, que
é fechado em $X$ por \sref{III.4.4.10}; como $f$ é uma aplicação
fechada, $U$ é aberto. Além disso, \sref[II]{II.6.2.2} mostra que
$f^{-1}(y)$ é finito para todo $y\in U$. Como o morfismo restrito
$f^{-1}(U)\to U$ é próprio \sref[II]{II.5.4.1}, é finito por
\sref{III.4.4.2}.
\end{proof}

\begin{remark}[4.4.12]
\label{III.4.4.12}
\emph{(i)}
Como se anunciou em \sref[II]{II.6.2.7}, demonstraremos no capítulo~V
que, se $Y$ é localmente noetheriano, todo morfismo separado
quase-finito $f:X\to Y$ é quase-afim, logo quase-projetivo. Seguir-se-á
que a conclusão do teorema principal \sref{III.4.4.3} continua
válida quando se supõe apenas que $f$ é separado e de tipo finito.
Com efeito, \sref{III.4.4.10} mostra que $X'$ é aberto em $X$; e
como $X$ é localmente noetheriano, a restrição de $f$ a $X'$ continua
sendo de tipo finito \sref[I]{I.6.3.5}, logo quase-finita pela
definição de $X'$, e claramente separada. Podemos, portanto, aplicar
\sref{III.4.4.3} a esta restrição.

\emph{(ii)}
No capítulo~IV daremos uma demonstração mais elementar de
\sref{III.4.4.10}, utilizando a teoria da dimensão.
\end{remark}


\subsection{Completamentos de módulos de homomorfismos}
\label{subsection:III.4.5}

\begin{proposition}[4.5.1]
\label{III.4.5.1}
Seja $A$ um anel noetheriano, seja $\mathfrak J$ um ideal de $A$, seja
$X$ um $A$-esquema de tipo finito, e sejam $\sh{F}$ e $\sh{G}$ dois
$\sh{O}_X$-módulos coerentes cujos suportes têm interseção
própria sobre $Y=\Spec(A)$ \sref[II]{II.5.4.10}. Então, para todo
inteiro $n\geq0$,
\[
  \Ext^n_{\sh{O}_X}(X;\sh{F},\sh{G})
\]
é um $A$-módulo finito, e seu completamento separado para a topologia
$\mathfrak J$-pré-ádica identifica-se canonicamente, com a notação de
\sref{III.4.1.7}, com
\[
  \Ext^n_{\sh{O}_{\hat{X}}}
    (\hat{X};\hat{\sh{F}},\hat{\sh{G}}).
\]
\end{proposition}

\begin{proof}
Sabemos por (T, 4.2) que existe uma sucessão espectral birregular
$\mathrm{E}(\sh{F},\sh{G})$ cujo termo limite é
\[
  \Ext_{\sh{O}_X}(X;\sh{F},\sh{G})
\]
e cujos termos $\mathrm{E}_2$ são
\[
  \mathrm{E}_2^{pq}
  =\HH^p\bigl(X,\shExt^q_{\sh{O}_X}(\sh{F},\sh{G})\bigr).
\]
O $\sh{O}_X$-módulo
$\shExt^q_{\sh{O}_X}(\sh{F},\sh{G})$ é coerente
\sref[0]{0.12.3.3}; seu suporte está contido na interseção dos
suportes de $\sh{F}$ e $\sh{G}$ (T, 4.2.2) e, consequentemente, é
próprio sobre $Y$ \sref[II]{II.5.4.10}. De \sref{III.3.2.4} se segue
que os $\mathrm{E}_2^{pq}$ são $A$-módulos finitos, e depois, de
\sref[0]{0.11.1.8}, que o mesmo ocorre com todo termo
$\mathrm{E}_r^{pq}$ da sucessão espectral e com seu termo limite.

Por outro lado, se $i:\hat{X}\to X$ é o morfismo canônico, então
$\hat{\sh{F}}$ e $\hat{\sh{G}}$ identificam-se canonicamente com
$i^*(\sh{F})$ e $i^*(\sh{G})$, e $i$ é plano
\sref[I]{I.10.8.8}, \sref[I]{I.10.8.9}. Para todo $q\geq0$ existe,
portanto, um $i$-morfismo canônico
\[
  u_q:\shExt^q_{\sh{O}_X}(\sh{F},\sh{G})
  \longrightarrow
  \shExt^q_{\sh{O}_{\hat{X}}}
    (\hat{\sh{F}},\hat{\sh{G}})
\]
\sref[0]{0.12.3.4}, e o $\sh{O}_{\hat{X}}$-homomorfismo
correspondente
\[
  u_q^\#:
  i^*\bigl(\shExt^q_{\sh{O}_X}(\sh{F},\sh{G})\bigr)
  \longrightarrow
  \shExt^q_{\sh{O}_{\hat{X}}}
    (\hat{\sh{F}},\hat{\sh{G}})
\]
é um isomorfismo \sref[0]{0.12.3.5}. Em outros termos,
\sref[I]{I.10.8.8} identifica
\[
  \shExt^q_{\sh{O}_{\hat{X}}}
    (\hat{\sh{F}},\hat{\sh{G}})
\]
canonicamente com o completamento de
$\shExt^q_{\sh{O}_X}(\sh{F},\sh{G})$, com a notação de
\sref{III.4.1.7}.

O teorema de comparação \sref{III.4.1.10} mostra agora, para todo
$p\geq0$, que
\[
  \HH^p\bigl(\hat{X},
    \shExt^q_{\sh{O}_{\hat{X}}}
      (\hat{\sh{F}},\hat{\sh{G}})\bigr)
\]
identifica-se canonicamente com o completamento separado, para a
topologia $\mathfrak J$-pré-ádica, de
\[
  \HH^p\bigl(X,\shExt^q_{\sh{O}_X}(\sh{F},\sh{G})\bigr).
\]

\oldpage[III]{139}

Denotemos por $\mathrm{E}(\hat{\sh{F}},\hat{\sh{G}})$ a sucessão
espectral birregular de (T, 4.2) associada a $\hat{\sh{F}}$ e
$\hat{\sh{G}}$. Se $\hat{A}$ é o completamento separado de $A$ para
a topologia $\mathfrak J$-pré-ádica, então, salvo isomorfismo
canônico,
\[
  \mathrm{E}_2^{pq}(\hat{\sh{F}},\hat{\sh{G}})
  =
  \mathrm{E}_2^{pq}(\sh{F},\sh{G})\otimes_A\hat{A}
\]
\sref[0]{0.7.3.3}.

O morfismo plano $i$ define um homomorfismo canônico de sucessões
espectrais
\[
  \varphi:\mathrm{E}(\sh{F},\sh{G})
  \longrightarrow
  \mathrm{E}(\hat{\sh{F}},\hat{\sh{G}})
  =
  \mathrm{E}(i^*(\sh{F}),i^*(\sh{G})).
\]
Sobre os termos $\mathrm{E}_2$ e sobre o termo limite,
respectivamente, este é o homomorfismo
\[
  \varphi_2^{pq}:
  \HH^p\bigl(X,\shExt^q_{\sh{O}_X}(\sh{F},\sh{G})\bigr)
  \longrightarrow
  \HH^p\bigl(\hat{X},
    \shExt^q_{\sh{O}_{\hat{X}}}
      (\hat{\sh{F}},\hat{\sh{G}})\bigr)
\]
e o homomorfismo
\[
  \varphi^n:
  \Ext^n_{\sh{O}_X}(X;\sh{F},\sh{G})
  \longrightarrow
  \Ext^n_{\sh{O}_{\hat{X}}}
    (\hat{X};\hat{\sh{F}},\hat{\sh{G}}),
\]
deduzidos de $u_q$ e $u_0$, respectivamente, por funtorialidade
\sref[0]{0.12.3.4}.

Após tensorizar com $\hat{A}$, os $\varphi_r^{pq}$ e
$\varphi^n$ dão homomorfismos de $\hat{A}$-módulos
\[
  \psi_r^{pq}:
  \mathrm{E}_r^{pq}(\sh{F},\sh{G})\otimes_A\hat{A}
  \longrightarrow
  \mathrm{E}_r^{pq}(\hat{\sh{F}},\hat{\sh{G}})
\]
e
\[
  \psi^n:
  \Ext^n_{\sh{O}_X}(X;\sh{F},\sh{G})\otimes_A\hat{A}
  \longrightarrow
  \Ext^n_{\sh{O}_{\hat{X}}}
    (\hat{X};\hat{\sh{F}},\hat{\sh{G}}).
\]
Como $\hat{A}$ é um $A$-módulo plano \sref[0]{0.7.3.3}, os
$\hat{A}$-módulos
\[
  \mathrm{E}_r^{pq}(\sh{F},\sh{G})\otimes_A\hat{A}
\]
formam uma sucessão espectral birregular cujo termo limite é
\[
  \Ext^n_{\sh{O}_X}(X;\sh{F},\sh{G})\otimes_A\hat{A},
\]
e os $\psi_r^{pq}$ e $\psi^n$ formam um morfismo de sucessões
espectrais. Como os $\psi_2^{pq}$ são isomorfismos, o mesmo ocorre
com os $\psi^n$ \sref[0]{0.11.1.5}.
\end{proof}

\begin{corollary}[4.5.2]
\label{III.4.5.2}
Sob as hipóteses de \sref{III.4.5.1}, suponhamos além disso que $A$ é
um anel noetheriano $\mathfrak J$-ádico. Então, para todo inteiro
$n\geq0$,
\[
  \Ext^n_{\sh{O}_X}(X;\sh{F},\sh{G})
  \isoto
  \Ext^n_{\sh{O}_{\hat{X}}}
    (\hat{X};\hat{\sh{F}},\hat{\sh{G}})
\]
canonicamente.
\end{corollary}

\begin{proof}
O $A$-módulo finito
$\Ext^n_{\sh{O}_X}(X;\sh{F},\sh{G})$ está separado e completo para a
topologia $\mathfrak J$-pré-ádica \sref[0]{0.7.3.6}.
\end{proof}

O caso particular $n=0$ de \sref{III.4.5.1} pode ser enunciado como
segue.

\begin{corollary}[4.5.3]
\label{III.4.5.3}
Sob as hipóteses de \sref{III.4.5.1}, para todo homomorfismo
$u:\sh{F}\to\sh{G}$, denotemos por
$\hat{u}:\hat{\sh{F}}\to\hat{\sh{G}}$ o homomorfismo completado
\sref[I]{I.10.8.4}. Existe um isomorfismo canônico
\[
\label{III.4.5.3.1}
  \left(\Hom_{\sh{O}_X}(\sh{F},\sh{G})\right)^\wedge
  \isoto
  \Hom_{\sh{O}_{\hat{X}}}(\hat{\sh{F}},\hat{\sh{G}}),
  \tag{4.5.3.1}
\]
onde o membro esquerdo é o completamento separado, para a
topologia $\mathfrak J$-pré-ádica, do $A$-módulo
$\Hom_{\sh{O}_X}(\sh{F},\sh{G})$. Este isomorfismo obtém-se passando
aos completamentos separados no homomorfismo
$u\mapsto\hat{u}$.
\end{corollary}


\subsection{Relações entre morfismos formais e morfismos ordinários}
\label{subsection:III.4.6}

\begin{proposition}[4.6.1]
\label{III.4.6.1}
Sejam $Y$ um esquema localmente noetheriano, $f:X\to Y$ um
morfismo próprio, $\sh{F}$ um $\sh{O}_X$-módulo coerente que é
plano sobre $Y$, e $y\in Y$. Suponhamos que, para algum inteiro
$n$,
\[
  \HH^n\bigl(f^{-1}(y),
    \sh{F}\otimes_{\sh{O}_Y}\kres(y)\bigr)=0.
\]
Então existe uma vizinhança $U$ de $y$ em $Y$ tal que
\[
  \RR^n f_*(\sh{F})|U=0,
\]
e, para todo inteiro $p\geq0$, o homomorfismo canônico
\[
  \bigl(\RR^{n-1}f_*(\sh{F})\bigr)_y
  \longrightarrow
  \HH^{n-1}\bigl(f^{-1}(y),
    \sh{F}\otimes_{\sh{O}_Y}
      (\sh{O}_y/\mathfrak m_y^{p+1})\bigr)
\]
\eqref{III.4.2.1.1} é sobrejetivo.
\end{proposition}

\begin{proof}
O $\sh{O}_Y$-módulo $\RR^n f_*(\sh{F})$ é coerente
\sref{III.3.2.1}; portanto, a primeira afirmação se seguirá uma vez que provemos
\[
  \bigl(\RR^n f_*(\sh{F})\bigr)_y=0
\]
\sref[0]{0.5.2.2}. Por \sref{III.4.2.1}, basta provar que
\[
  \HH^n\bigl(f^{-1}(y),
    \sh{F}\otimes_{\sh{O}_Y}
      (\sh{O}_y/\mathfrak m_y^{p+1})\bigr)=0
\]
para todo $p$. Isto é certo por hipótese para $p=0$; procedemos por
indução sobre $p$.

Ponhamos
\[
  X_p=X\times_Y
    \Spec(\sh{O}_y/\mathfrak m_y^{p+1}).
\]
Então $X_{p-1}$ é um subesquema fechado de $X_p$ com o mesmo
espaço subjacente \sref[I]{I.3.6.1}. A hipótese de indução
\[
  \HH^n\bigl(X_{p-1},
    \sh{F}\otimes_{\sh{O}_Y}
      (\sh{O}_y/\mathfrak m_y^p)\bigr)=0
\]
dá, portanto,
\[
  \HH^n\bigl(X_p,
    \sh{F}\otimes_{\sh{O}_Y}
      (\sh{O}_y/\mathfrak m_y^p)\bigr)=0.
\]
Por outro lado, a sucessão exata de $\sh{O}_{X_p}$-módulos
\[
  0\longrightarrow
  \mathfrak m_y^p\sh{F}/\mathfrak m_y^{p+1}\sh{F}
  \longrightarrow
  \sh{F}/\mathfrak m_y^{p+1}\sh{F}
  \longrightarrow
  \sh{F}/\mathfrak m_y^p\sh{F}
  \longrightarrow0
\]
dá a sucessão exata de cohomologia
\[
  \HH^n\bigl(X_p,
    \mathfrak m_y^p\sh{F}/\mathfrak m_y^{p+1}\sh{F}\bigr)
  \longrightarrow
  \HH^n\bigl(X_p,\sh{F}/\mathfrak m_y^{p+1}\sh{F}\bigr)
  \longrightarrow
  \HH^n\bigl(X_p,\sh{F}/\mathfrak m_y^p\sh{F}\bigr).
\]
Basta, pois, demonstrar que
\[
\label{III.4.6.1.1}
  \HH^n\bigl(X_p,
    \mathfrak m_y^p\sh{F}/\mathfrak m_y^{p+1}\sh{F}\bigr)=0.
  \tag{4.6.1.1}
\]
Com efeito, o termo central será então um submódulo do último termo,
que é zero pela hipótese de indução.

\oldpage[III]{140}

A fibra
\[
  Z=f^{-1}(y)=X\times_Y\Spec(\kres(y))
\]
é um subesquema fechado de $X_p$, e
$\mathfrak m_y^p\sh{F}/\mathfrak m_y^{p+1}\sh{F}$ é anulado por
$\mathfrak m_y$. Portanto, pode ser considerado um
$\sh{O}_Z$-módulo, e
\[
  \HH^n\bigl(Z,
    \mathfrak m_y^p\sh{F}/\mathfrak m_y^{p+1}\sh{F}\bigr)
  =
  \HH^n\bigl(X_p,
    \mathfrak m_y^p\sh{F}/\mathfrak m_y^{p+1}\sh{F}\bigr).
\]
Provaremos agora que o $\sh{O}_Z$-homomorfismo canônico
\[
\label{III.4.6.1.2}
  (\sh{F}/\mathfrak m_y\sh{F})
    \otimes_{\kres(y)}
      (\mathfrak m_y^p/\mathfrak m_y^{p+1})
  \longrightarrow
  \mathfrak m_y^p\sh{F}/\mathfrak m_y^{p+1}\sh{F}
  \tag{4.6.1.2}
\]
é bijetivo. Uma vez estabelecido isto, o fato de que
$\mathfrak m_y^p/\mathfrak m_y^{p+1}$ seja um
$\kres(y)$-módulo livre dá
\[
\begin{split}
  \HH^n\bigl(Z,
    \mathfrak m_y^p\sh{F}/\mathfrak m_y^{p+1}\sh{F}\bigr)
  &=
  \HH^n\bigl(Z,\sh{F}/\mathfrak m_y\sh{F}\bigr)
    \otimes_{\kres(y)}
      (\mathfrak m_y^p/\mathfrak m_y^{p+1})\\
  &=0
\end{split}
\]
\sref[0]{0.12.2.3}, dado que
$\HH^n(Z,\sh{F}/\mathfrak m_y\sh{F})=0$ por hipótese. Isto prova
\eqref{III.4.6.1.1}.

Para terminar a demonstração da primeira afirmação, resta demonstrar
que \eqref{III.4.6.1.2} é bijetivo. A questão pode ser verificada ponto a ponto em
$X$, e $\sh{F}_x$ é um $\sh{O}_y$-módulo plano por hipótese para
todo $x\in f^{-1}(y)$. Basta, pois, aplicar
\sref[0]{0.10.2.1}, c), dado que
$\sh{F}_x/\mathfrak m_y\sh{F}_x$ é plano sobre o corpo
$\kres(y)=\sh{O}_y/\mathfrak m_y$.

Para a segunda afirmação de \sref{III.4.6.1}, reduzimo-nos imediatamente, como
em \sref{III.4.2.1}, ao caso em que $Y$ é afim e $y$ é
fechado. O mesmo argumento demonstra, para todo $k>0$, que
\[
\label{III.4.6.1.3}
  \HH^n\bigl(X_{k+1},
    \mathfrak m_y^k\sh{F}/
      \mathfrak m_y^{k+p+1}\sh{F}\bigr)=0.
  \tag{4.6.1.3}
\]

\oldpage[III]{141}

De \sref{III.4.2.1} segue-se que
\[
\label{III.4.6.1.4}
  \bigl(\RR^n f_*(\mathfrak m_y^k\sh{F})\bigr)_y=0.
  \tag{4.6.1.4}
\]
A sucessão exata de cohomologia dá agora a exatidão de
\[
  \bigl(\RR^{n-1}f_*(\sh{F})\bigr)_y
  \longrightarrow
  \bigl(\RR^{n-1}f_*
    (\sh{F}/\mathfrak m_y^p\sh{F})\bigr)_y
  \longrightarrow
  \bigl(\RR^n f_*(\mathfrak m_y^p\sh{F})\bigr)_y=0.
\]
Como $y$ é fechado e $Y$ é afim, temos
\sref{III.1.4.11}
\[
  \RR^{n-1}f_*(\sh{F}/\mathfrak m_y^p\sh{F})
  =
  \bigl(\HH^{n-1}
    (X,\sh{F}/\mathfrak m_y^p\sh{F})\bigr)^{\sim}
  =
  \bigl(\HH^{n-1}
    (f^{-1}(y),\sh{F}/\mathfrak m_y^p\sh{F})\bigr)^{\sim}
\]
(G, II, 4.9.1). O módulo
\[
  \HH^{n-1}
    (f^{-1}(y),\sh{F}/\mathfrak m_y^p\sh{F})
\]
é um $\sh{O}_y/\mathfrak m_y^p$-módulo, e portanto
\[
  \bigl(\RR^{n-1}f_*
    (\sh{F}/\mathfrak m_y^p\sh{F})\bigr)_y
  =
  \HH^{n-1}
    (f^{-1}(y),\sh{F}/\mathfrak m_y^p\sh{F}).
\]
Isto termina a demonstração.
\end{proof}

\begin{corollary}[4.6.2]
\label{III.4.6.2}
Sejam $Y$ um esquema localmente noetheriano, $f:X\to Y$ um
morfismo próprio e plano, $\sh{F}$ e $\sh{G}$ dois
$\sh{O}_X$-módulos localmente livres, e $y$ um ponto de $Y$. Ponhamos
\[
  X_y=f^{-1}(y)=X\times_Y\Spec(\kres(y)),\qquad
  \sh{F}_y=\sh{F}\otimes_{\sh{O}_Y}\kres(y),\qquad
  \sh{G}_y=\sh{G}\otimes_{\sh{O}_Y}\kres(y),
\]
e suponhamos que
\[
\label{III.4.6.2.1}
  \HH^1\bigl(
    X_y,\shHom_{\sh{O}_{X_y}}(\sh{F}_y,\sh{G}_y)
  \bigr)=0.
  \tag{4.6.2.1}
\]
Então, para todo homomorfismo $u_0:\sh{F}_y\to\sh{G}_y$, existem
uma vizinhança aberta $U$ de $y$ em $Y$ e um homomorfismo
\[
  u:\sh{F}|_{f^{-1}(U)}\longrightarrow
    \sh{G}|_{f^{-1}(U)}
\]
tal que $u_0=u\otimes1$.
\end{corollary}

\begin{proof}
A hipótese nos permite aplicar \sref{III.4.6.1}, com $n=1$
e $p=0$, ao $\sh{O}_X$-módulo coerente
\[
  \sh{H}=\shHom_{\sh{O}_X}(\sh{F},\sh{G}).
\]
Com efeito, $\sh{H}$ é localmente livre e, portanto, com maior razão, plano sobre
$Y$, enquanto que o $\sh{O}_{X_y}$-módulo
$\sh{H}\otimes_{\sh{O}_Y}\kres(y)$ identifica-se com
$\shHom_{\sh{O}_{X_y}}(\sh{F}_y,\sh{G}_y)$
\sref[0]{0.6.2.2}. Podemos supor que $Y=\Spec(A)$ é afim.
Então \sref{III.1.4.11} dá
\[
  \RR^0f_*(\sh{H})
  =
  \bigl(\Hom_{\sh{O}_X}(\sh{F},\sh{G})\bigr)^{\sim},
\]
e, consequentemente,
\[
  \bigl(\RR^0f_*(\sh{H})\bigr)_y
  =
  \Hom_{\sh{O}_X}(\sh{F},\sh{G})
    \otimes_A\sh{O}_y.
\]
Por \sref{III.4.6.1}, o homomorfismo canônico
\[
  \Hom_{\sh{O}_X}(\sh{F},\sh{G})
    \otimes_A\sh{O}_y
  \longrightarrow
  \Hom_{\sh{O}_{X_y}}(\sh{F}_y,\sh{G}_y)
\]
é sobrejetivo. Isto prova o corolário, dado que todo elemento do
módulo da esquerda pode escrever-se como $u\otimes(1/s)$, onde
$s\in A\setminus\mathfrak m_y$.
\end{proof}

\begin{corollary}[4.6.3]
\label{III.4.6.3}
Sob as hipóteses de \sref{III.4.6.2}, se $u_0$ é injetivo
(respectivamente sobrejetivo, bijetivo), pode-se escolher $u$
injetivo (respectivamente sobrejetivo, bijetivo).
\end{corollary}

\begin{proof}
Podemos restringir-nos ao caso $U=Y$. Basta provar
que, se $u_0$ é injetivo (respectivamente sobrejetivo), então
\[
  \Ker u_x=0
  \quad\text{(respectivamente }\Coker u_x=0\text{)}
\]
para todo $x\in f^{-1}(y)$. Com efeito, $\Ker u$ e $\Coker u$ são
$\sh{O}_X$-módulos coerentes \sref[0]{0.5.3.4}; portanto, existe uma
vizinhança $V$ de $f^{-1}(y)$ em $X$ sobre a qual a restrição de
$\Ker u$ (respectivamente $\Coker u$) é zero \sref[0]{0.5.2.2}.
Como $f$ é fechado, existe uma vizinhança aberta
$U'\subset U$ de $y$ tal que $f^{-1}(U')\subset V$, o que prova
a afirmação depois de substituir $U$ por $U'$.

\oldpage[III]{142}

Por hipótese,
\[
  u_x\otimes1:
  \sh{F}_x\otimes_{\sh{O}_y}\kres(y)
  \longrightarrow
  \sh{G}_x\otimes_{\sh{O}_y}\kres(y)
\]
é injetivo (respectivamente sobrejetivo), $\sh{F}_x$ e $\sh{G}_x$
são $\sh{O}_x$-módulos livres finitos, e $\sh{O}_x$ é plano sobre
$\sh{O}_y$. Se $u_x\otimes1$ é injetivo, a injetividade de $u_x$
se segue de \sref[0]{0.10.2.4}. Se $u_x\otimes1$ é sobrejetivo,
então, com maior razão, o homomorfismo induzido
\[
  \sh{F}_x/\mathfrak m_x\sh{F}_x
  \longrightarrow
  \sh{G}_x/\mathfrak m_x\sh{G}_x
\]
é sobrejetivo. Como $\sh{G}_x$ é um $\sh{O}_x$-módulo finito e
$\sh{O}_x$ é um anel local de ideal maximal $\mathfrak m_x$, a
conclusão se segue do lema de Nakayama (Bourbaki,
\emph{Alg\`ebre}, cap.~VIII, \S\,6, n.º~3, cor.~4 da prop.~6).
\end{proof}

\begin{corollary}[4.6.4]
\label{III.4.6.4}
Sejam $Y$ um esquema localmente noetheriano, $f:X\to Y$ um
morfismo próprio e plano, e $y$ um ponto de $Y$, e ponhamos
$X_y=X\times_Y\Spec(\kres(y))$. Seja $\sh{E}_0$ um
$\sh{O}_{X_y}$-módulo localmente livre tal que
\[
\label{III.4.6.4.1}
  \HH^1\bigl(
    X_y,\shHom_{\sh{O}_{X_y}}(\sh{E}_0,\sh{E}_0)
  \bigr)=0.
  \tag{4.6.4.1}
\]
Sejam $\sh{F}$ e $\sh{G}$ dois $\sh{O}_X$-módulos localmente livres
tais que $\sh{F}_y$ e $\sh{G}_y$, com a notação de
\sref{III.4.6.2}, são isomorfos a $\sh{E}_0$. Então existe uma
vizinhança aberta $U$ de $y$ tal que
$\sh{F}|_{f^{-1}(U)}$ e $\sh{G}|_{f^{-1}(U)}$ são isomorfos.
\end{corollary}

\begin{corollary}[4.6.5]
\label{III.4.6.5}
Sob as hipóteses de \sref{III.4.6.4} sobre $f$, $X$ e $Y$,
suponhamos que
\[
  \HH^1(X_y,\sh{O}_{X_y})=0.
\]
Se $\sh{F}$ e $\sh{G}$ são dois $\sh{O}_X$-módulos invertíveis tais
que $\sh{F}_y$ e $\sh{G}_y$ são isomorfos, então existe uma vizinhança
aberta $U$ de $y$ tal que $\sh{F}|_{f^{-1}(U)}$ e
$\sh{G}|_{f^{-1}(U)}$ são isomorfos.
\end{corollary}

\begin{proof}
Basta aplicar \sref{III.4.6.4} aos módulos
$\sh{F}^{-1}\otimes\sh{G}$ e $\sh{O}_X$.
\end{proof}

\begin{remarks}[4.6.6]
\label{III.4.6.6}
\begin{enumerate}
\item[(i)] Utilizando \sref{III.4.6.5}, estabeleceremos no capítulo~V
  a classificação dos feixes invertíveis sobre um fibrado projetivo
  anunciada em \sref[II]{II.4.2.7}.
\item[(ii)] O resultado de \sref{III.4.6.1} aparecerá em \S\,7 como
  consequência de proposições mais gerais.
\end{enumerate}
\end{remarks}

\begin{proposition}[4.6.7]
\label{III.4.6.7}
Seja $Z$ um esquema localmente noetheriano, e sejam $X$ e $Y$
dois $Z$-esquemas cujos morfismos estruturais
\[
  g:X\longrightarrow Z,\qquad h:Y\longrightarrow Z
\]
são próprios. Seja $f:X\to Y$ um $Z$-morfismo, seja $z$ um ponto
de $Z$, e ponhamos
\[
  f_z=f\times_Z1:
  X\times_Z\Spec(\kres(z))
  \longrightarrow
  Y\times_Z\Spec(\kres(z)).
\]
\begin{enumerate}
\item[(i)] Se $f_z$ é finito (respectivamente uma imersão fechada),
  existe uma vizinhança aberta $U$ de $z$ tal que o morfismo
  \[
    g^{-1}(U)\longrightarrow h^{-1}(U)
  \]
  induzido por $f$ é finito (respectivamente uma imersão fechada).
\item[(ii)] Suponhamos além disso que $g$ é plano. Se $f_z$ é um
  isomorfismo, existe uma vizinhança aberta $U$ de $z$ tal que
  o morfismo $g^{-1}(U)\to h^{-1}(U)$ induzido por $f$ é um
  isomorfismo.
\end{enumerate}
\end{proposition}

\begin{proof}
Em ambos os casos basta provar que, para todo
$y\in h^{-1}(z)$, existe uma vizinhança $V_y$ de $y$ tal que
a restrição
\[
  f^{-1}(V_y)\longrightarrow V_y
\]
de $f$ é finita (respectivamente uma imersão fechada, um isomorfismo).
Com efeito, se $V$ é a união dos $V_y$, então
$f^{-1}(V)\to V$ é finito (respectivamente uma imersão fechada, um
isomorfismo) \sref[II]{II.6.1.1}, \sref[I]{I.4.2.4}. Como $h$ é
fechado, existe uma vizinhança aberta $U$ de $z$ tal que
$h^{-1}(U)\subset V$, e a proposição se segue.

\noindent\textit{(i)} Observe-se primeiro que $f$ é próprio
\sref[II]{II.5.4.3}. Se
$f_z$ é finito, a vizinhança requerida $V_y$, para cada
$y\in h^{-1}(z)$, existe por \sref{III.4.4.11}.

\oldpage[III]{143}

Para tratar o caso em que $f_z$ é uma imersão fechada, podemos
supor já, portanto, que $f$ é finito. Assim,
$X=\Spec(\sh{B})$, onde $\sh{B}$ é uma
$\sh{O}_Y$-álgebra coerente, e $f$ corresponde \sref[II]{II.1.2.7} ao
homomorfismo canônico
\[
  u:\sh{O}_Y\longrightarrow\sh{B}.
\]
Se, para todo $y\in h^{-1}(z)$, o homomorfismo
$u_y:\sh{O}_y\to\sh{B}_y$ é sobrejetivo, então $u$ é sobrejetivo em
alguma vizinhança $V_y$ de $y$, dado que $\Coker u$ é coerente
\sref[0]{0.5.3.4}, \sref[0]{0.5.2.2}. Agora o morfismo finito
$f_z$ corresponde a
\[
  v=u\otimes1:
  \sh{O}_Y\otimes_{\sh{O}_Z}\kres(z)
  \longrightarrow
  \sh{B}\otimes_{\sh{O}_Z}\kres(z),
\]
e a hipótese de que $f_z$ é uma imersão fechada implica que
\[
  u_y\otimes1:
  \sh{O}_y\otimes_{\sh{O}_z}\kres(z)
  \longrightarrow
  \sh{B}_y\otimes_{\sh{O}_z}\kres(z)
\]
é sobrejetivo. Como $\sh{B}_y$ é um $\sh{O}_y$-módulo finito e
$\sh{O}_y$ é um anel local noetheriano, a conclusão se segue, como
em \sref{III.4.6.3}, do lema de Nakayama.

\noindent\textit{(ii)} O mesmo raciocínio mostra que agora
basta provar
que $u_y$ é bijetivo, sabendo que $u_y\otimes1$ é bijetivo.
A injetividade requerida se segue do lema seguinte; a sobrejetividade
se segue do argumento precedente.
\end{proof}

\begin{lemma}[4.6.7.1]
\label{III.4.6.7.1}
Sejam $A$ e $B$ anéis locais noetherianos, seja
$\rho:A\to B$ um homomorfismo local, e seja $u:N\to M$ um
homomorfismo de $B$-módulos. Suponhamos que $M$ é plano sobre $A$,
que $N$ é um $B$-módulo finito, e que
\[
  u\otimes1:
  N\otimes_A k\longrightarrow M\otimes_A k,
\]
onde $k$ é o corpo residual de $A$, é injetivo. Então $N$ é
plano sobre $A$ e $u$ é injetivo.
\end{lemma}

\begin{proof}
Para provar a primeira afirmação, devemos demonstrar que, para todo par de
$A$-módulos finitos $P,Q$ e todo $A$-homomorfismo injetivo
$v:P\to Q$, o homomorfismo
\[
  1_N\otimes v:N\otimes_A P\longrightarrow N\otimes_A Q
\]
é injetivo. Consideremos o diagrama comutativo
\[
  \xymatrix{
    N\otimes_A P
      \ar[r]^-{1_N\otimes v}
      \ar[d]_{u\otimes1_P}
      &
    N\otimes_A Q
      \ar[d]^{u\otimes1_Q}
    \\
    M\otimes_A P
      \ar[r]_{1_M\otimes v}
      &
    M\otimes_A Q .
  }
\]
Como $1_M\otimes v$ é injetivo pela planicidade de $M$, basta
provar que $u\otimes1_P$ é injetivo.

Seja $\mathfrak m$ o ideal maximal de $A$. A
filtração $\mathfrak m$-ádica sobre o $A$-módulo $N\otimes_A P$ é
também sua filtração $\mathfrak mB$-ádica como $B$-módulo. A
topologia definida por esta filtração é separada: $B$ é
noetheriano, $\mathfrak mB$ está contido no radical de $B$, e
$N\otimes_A P$ é um $B$-módulo finito, dado que $N$ é finito sobre $B$
e $P$ é finito sobre $A$ \sref[0]{0.7.3.5}. Basta, portanto,
provar que
\[
  \operatorname{gr}(u\otimes1_P):
  \operatorname{gr}_\bullet(N\otimes_A P)
  \longrightarrow
  \operatorname{gr}_\bullet(M\otimes_A P)
\]
é injetivo, onde os módulos graduados associados se tomam com
respeito às filtrações $\mathfrak m$-ádicas (Bourbaki,
\emph{Alg\`ebre commutative}, cap.~III, \S\,2, n.º~8, cor.~1 do
teor.~1).

Como $M$ é plano sobre $A$, os homomorfismos
\[
  M\otimes_A(\mathfrak m^nP)
  \longrightarrow
  \mathfrak m^n(M\otimes_A P)
\]
são bijetivos. Consequentemente, o homomorfismo canônico
\[
  \vphi_M:
  \operatorname{gr}_0(M)\otimes_A\operatorname{gr}_\bullet(P)
  \longrightarrow
  \operatorname{gr}_\bullet(M\otimes_A P)
\]
é bijetivo.

\oldpage[III]{144}

Temos o diagrama comutativo
\[
  \xymatrix{
    \operatorname{gr}_0(N)\otimes_A\operatorname{gr}_\bullet(P)
      \ar[r]^-{\operatorname{gr}_0(u)\otimes1}
      \ar[d]_{\vphi_N}
      &
    \operatorname{gr}_0(M)\otimes_A\operatorname{gr}_\bullet(P)
      \ar[d]^{\vphi_M}
    \\
    \operatorname{gr}_\bullet(N\otimes_A P)
      \ar[r]_{\operatorname{gr}(u\otimes1)}
      &
    \operatorname{gr}_\bullet(M\otimes_A P).
  }
\]
Aqui $\vphi_M$ é bijetivo e $\vphi_N$ é sobrejetivo.
Além disso, $\operatorname{gr}_0(u)$ é injetivo por hipótese, e
\[
\begin{split}
  \operatorname{gr}_0(N)\otimes_A\operatorname{gr}_\bullet(P)
    &=
  \operatorname{gr}_0(N)\otimes_k\operatorname{gr}_\bullet(P),\\
  \operatorname{gr}_0(M)\otimes_A\operatorname{gr}_\bullet(P)
    &=
  \operatorname{gr}_0(M)\otimes_k\operatorname{gr}_\bullet(P).
\end{split}
\]
Assim, $\operatorname{gr}_0(u)\otimes1$ também é injetivo. Do
diagrama segue-se que $\operatorname{gr}(u\otimes1_P)$ é
injetivo, o que prova a primeira afirmação. A segunda se segue do
mesmo argumento tomando $P=A$.
\end{proof}

\begin{proposition}[4.6.8]
\label{III.4.6.8}
Seja $Z$ um esquema localmente noetheriano, e sejam $X$ e $Y$
dois $Z$-esquemas cujos morfismos estruturais
$g:X\to Z$ e $h:Y\to Z$ são próprios. Seja $Z'$ um subconjunto fechado
de $Z$, ponhamos
\[
  X'=g^{-1}(Z'),\qquad Y'=h^{-1}(Z'),
\]
e denotemos por
\[
  \widehat X=X_{/X'},\qquad
  \widehat Y=Y_{/Y'},\qquad
  \widehat Z=Z_{/Z'}
\]
os completamentos formais ao longo destes subconjuntos fechados. Seja $f:X\to Y$
um $Z$-morfismo e seja
$\widehat f:\widehat X\to\widehat Y$ sua extensão aos completamentos
formais. Para que $\widehat f$ seja um isomorfismo
(respectivamente uma imersão fechada), é necessário e suficiente
que exista uma vizinhança aberta $U$ de $Z'$ tal que o
morfismo
\[
  g^{-1}(U)\longrightarrow h^{-1}(U)
\]
induzido por $f$ seja um isomorfismo (respectivamente uma imersão fechada).
\end{proposition}

\begin{proof}
A suficiência da condição é imediata
\sref[I]{I.10.14.7}. Para provar sua necessidade, basta, pelo
mesmo raciocínio de \sref{III.4.6.7}, demonstrar que para todo
$y\in Y'$ existe uma vizinhança aberta $V_y$ de $y$ tal que
$f^{-1}(V_y)\to V_y$ seja um isomorfismo (respectivamente uma
imersão fechada). Ficamos assim reduzidos ao caso
\[
  Y=Z=\Spec(A),
\]
onde $A$ é noetheriano.

Por hipótese \sref[I]{I.10.9.1}, \sref[I]{I.10.14.2}, a fibra
$f^{-1}(y)$ consiste de um único ponto para todo $y\in Y'$. Como
$f$ é próprio \sref[II]{II.5.4.3}, \sref{III.4.4.11} dá uma vizinhança
aberta $U$ de $y$ tal que $f^{-1}(U)\to U$ é finito. Podemos,
portanto, supor que $f$ é finito, de modo que
\[
  X=\Spec(B),
\]
onde $B$ é uma $A$-álgebra finita.

Se $Y'=V(\mathfrak J)$, então
\[
  \widehat Y=\Spf(\widehat A),
  \qquad
  \widehat X=\Spf(\widehat B),
\]
onde $\widehat A$ é o completamento separado de $A$ para a
topologia $\mathfrak J$-ádica, e $\widehat B$ é o completamento
separado de $B$ para a topologia $\mathfrak JB$-ádica, ou,
equivalentemente, o completamento separado do $A$-módulo $B$ para
a topologia $\mathfrak J$-ádica. Além disso, $\widehat f$
corresponde à extensão contínua
\[
  \widehat{\vphi}:\widehat A\longrightarrow\widehat B
\]
do homomorfismo canônico de anéis $\vphi:A\to B$, e a
hipótese é que $\widehat{\vphi}$ é sobrejetivo (respectivamente
bijetivo) \sref[I]{I.10.14.2}. Mas $\widehat{\vphi}$ é também a
extensão contínua de $\vphi$ considerado como um homomorfismo de
$A$-módulos. De \sref[I]{I.10.8.14} segue-se que existe uma
vizinhança aberta $U$ de $Y'$ tal que a restrição a $U$ do
homomorfismo de $\sh{O}_Y$-módulos
\[
  \widetilde{\vphi}:\widetilde A\longrightarrow\widetilde B
\]
é sobrejetiva (respectivamente bijetiva), o que termina a demonstração.
\end{proof}


\subsection{Um critério de amplitude}
\label{subsection:III.4.7}

\begin{theorem}[4.7.1]
\label{III.4.7.1}
Sejam $Y$ um esquema localmente noetheriano, $f:X\to Y$ um
morfismo próprio, $\sh{L}$ um $\sh{O}_X$-módulo invertível,
e $y$ um ponto de $Y$. Ponhamos
\[
  X_y=X\times_Y\Spec(\kres(y))=f^{-1}(y),
\]
e seja $g:X_y\to X$ o morfismo canônico. Se
\[
  \sh{L}_y=g^*(\sh{L})
  =\sh{L}\otimes_{\sh{O}_Y}\kres(y)
\]
é amplo sobre $X_y$, então existe uma vizinhança aberta $U$ de $y$
em $Y$ tal que $\sh{L}|_{f^{-1}(U)}$ é amplo para a restrição
de $f$ a $f^{-1}(U)$.
\end{theorem}

\begin{proof}
\oldpage[III]{145}

\noindent\textit{I)} Ponhamos
\[
  Y'=\Spec(\sh{O}_y),\qquad
  X'=X\times_Y Y',\qquad
  \sh{L}'=\sh{L}\otimes_{\sh{O}_Y}\sh{O}_y.
\]
Provamos primeiro que $\sh{L}'$ é amplo para $f'=f_{(Y')}$. Temos
o diagrama comutativo
\[
  \xymatrix{
    X\ar[d]_f
      &
    X'\ar[l]\ar[d]_{f'}
      &
    X_y\ar[l]\ar[d]_{f_y}
    \\
    Y
      &
    Y'\ar[l]
      &
    \Spec(\kres(y))\ar[l].
  }
\]
Como $f'$ é próprio \sref[II]{II.5.4.2}, (iii), e $\sh{O}_y$ é
noetheriano, \sref[II]{II.4.6.6} mostra que podemos reduzir-nos ao caso
\[
  Y=Y'=\Spec(\sh{O}_y),\qquad X=X',
\]
supor que $\sh{L}_y$ é amplo para $f_y$, e provar que $\sh{L}$
é amplo para $f$. Aplicamos o critério
\sref{III.2.6.1}, c), e provamos que para todo
$\sh{O}_X$-módulo coerente $\sh{F}$ existe um inteiro $N$ tal que
\[
  \HH^1(X,\sh{F}(n))=0
  \quad\text{para todo }n\geq N,
  \qquad
  \sh{F}(n)=\sh{F}\otimes\sh{L}^{\otimes n}.
\]

O ponto $y$ é fechado em $Y$ e corresponde ao ideal maximal
$\mathfrak m$ de $\sh{O}_y$. Portanto, $X_y$ é o
subesquema fechado de $X$ definido pelo ideal coerente
\[
  \sh{J}
  =
  f^*(\widetilde{\mathfrak m})\sh{O}_X
  =
  \mathfrak m\sh{O}_X
\]
\sref[I]{I.4.4.5}, e $g$ é sua injeção canônica. Consideremos
a $\kres(y)$-álgebra graduada
\[
  S=\operatorname{gr}(\sh{O}_y)
   =\bigoplus_{j\geq0}
      \mathfrak m^j/\mathfrak m^{j+1}.
\]
É de tipo finito porque $\sh{O}_y$ é noetheriano. A
$\sh{O}_X$-álgebra
\[
  \sh{S}=f^*(\widetilde S)
\]
é quase-coerente e de tipo finito, e é anulada por
$\sh{J}$. Assim, se $\sh{S}_y=g^*(\sh{S})$, então $\sh{S}_y$ é uma
$\sh{O}_{X_y}$-álgebra quase-coerente de tipo finito e
\[
  \sh{S}=g_*(\sh{S}_y).
\]

Ponhamos
\[
  \sh{M}_j
  =
  \mathfrak m^j\sh{F}/\mathfrak m^{j+1}\sh{F},
  \qquad
  \sh{M}
  =
  \bigoplus_{j\geq0}\sh{M}_j
  =
  \operatorname{gr}(\sh{F}).
\]
Como $\sh{F}$ é coerente, $\sh{M}$ é um
$\sh{S}$-módulo quase-coerente de tipo finito \sref[0]{0.10.1.1}; também é
anulado por $\sh{J}$. Consequentemente, se
\[
  \sh{M}'_j=g^*(\sh{M}_j),
  \qquad
  \sh{M}'=g^*(\sh{M})
    =\bigoplus_{j\geq0}\sh{M}'_j,
\]
então $\sh{M}'$ é um $\sh{S}_y$-módulo graduado quase-coerente de
tipo finito e $\sh{M}=g_*(\sh{M}')$. Além disso, pondo
\[
  \sh{M}'_j(n)
  =
  \sh{M}'_j\otimes\sh{L}_y^{\otimes n},
\]
temos $\sh{M}'_j(n)=g^*(\sh{M}_j(n))$.

O morfismo $f_y$ é próprio \sref[II]{II.5.4.2}, (iii), e
$\sh{L}_y$ é amplo; portanto, $f_y$ é projetivo
\sref[II]{II.5.5.4}, \sref[II]{II.4.6.11}. Podemos, pois, aplicar
\sref{III.2.4.1}, (ii), a
$\Spec(\kres(y))$, $f_y$, $\sh{S}_y$, $\sh{L}_y$ e $\sh{M}'$.
Existe um inteiro $N$ tal que, para $n\geq N$,
\[
  \HH^q(X_y,\sh{M}'_j(n))=0
\]
para todo $q>0$ e todo $j$. Segue-se que
\[
  \HH^q(X,\sh{M}_j(n))=0
\]
para todo $q>0$ e todo $j$ (G, II, 4.9.1).

Ponhamos agora
\[
  \sh{F}(n)_j
  =
  \sh{F}(n)/\mathfrak m^{j+1}\sh{F}(n).
\]
Para $j\geq1$ temos
\[
  \sh{F}(n)_{j-1}
  =
  \sh{F}(n)_j/\sh{M}_j(n),
  \qquad
  \sh{F}(n)_0=\sh{M}_0(n).
\]
Assim, $\HH^1(X,\sh{F}(n)_0)=0$, e a sucessão exata de cohomologia
dá, para todo $j\geq1$,
\[
  \HH^1(X,\sh{F}(n)_j)
  =
  \HH^1(X,\sh{F}(n)_{j-1}).
\]
Portanto, $\HH^1(X,\sh{F}(n)_j)=0$ para todo $j\geq0$. Por
\sref{III.4.2.1}, segue-se que
\[
  \HH^1(X,\sh{F}(n))=0,
\]
o que prova a afirmação.

\oldpage[III]{146}

\noindent\textit{II)} Voltemos à notação do início da
demonstração. Podemos supor que $Y=\Spec(A)$ é afim. Como
$f'$ é de tipo finito e $\sh{L}'$ é amplo para $f'$, existe um
inteiro $m>0$ tal que $\sh{L}'^{\otimes m}$ é muito amplo para
$f'$ \sref[II]{II.4.6.11}. Substituindo $\sh{L}$ por
$\sh{L}^{\otimes m}$, basta considerar o caso em que
$\sh{L}'$ é muito amplo para $f'$ e provar que
$\sh{L}|_{f^{-1}(U)}$ é muito amplo para a restrição de $f$.

Como $f'$ é próprio, existe, para um inteiro adequado $r>0$, uma
$Y'$-imersão fechada
\[
  j:X'\longrightarrow P=\bb{P}_{Y'}^r
\]
tal que $\sh{L}'\simeq j^*(\sh{O}_P(1))$
\sref[II]{II.5.5.4}, (ii). Essa imersão fechada corresponde
canonicamente a um $\sh{O}_{X'}$-homomorfismo sobrejetivo
\[
  u:\sh{O}_{X'}^{r+1}\longrightarrow\sh{L}'
\]
\sref[II]{II.4.2.3}, ou, equivalentemente, \sref[0]{0.5.1.1} a
$r+1$ seções $s'_i$ $(0\leq i\leq r)$ de $\sh{L}'$ sobre $X'$
que geram este $\sh{O}_{X'}$-módulo.

Por definição, estas são também seções de $f'_*(\sh{L}')$ sobre
$Y'$, e
\[
  f'_*(\sh{L}')
  =
  f_*(\sh{L})\otimes_{\sh{O}_Y}\sh{O}_{Y'}
\]
\sref[0]{0.5.4.10}. Como $Y$ é afim e $\sh{O}_y$ é o
anel local de $A$ no ideal primo $\mathfrak j_y$, cada $s'_i$
tem a forma
\[
  s'_i=s''_i/t_i,
\]
onde $s''_i$ é uma seção de $f_*(\sh{L})$ sobre $Y$ e
$t_i\in A\setminus\mathfrak j_y$. Segue-se que existem uma
vizinhança aberta $V$ de $y$ em $Y$ e seções $s_i$ de
$f_*(\sh{L})|_V$ tais que $s'_i=s_i/1$ (o espaço $Y'$ está
contido em $V$; veja-se \sref[I]{I.2.4.2}).

As seções $s_i$ são seções de $\sh{L}$ sobre $f^{-1}(V)$ e
definem, portanto, um homomorfismo
\[
  v:
  \bigl(\sh{O}_X|_{f^{-1}(V)}\bigr)^{r+1}
  \longrightarrow
  \sh{L}|_{f^{-1}(V)}.
\]
Por construção, é sobrejetivo em todo ponto de $f^{-1}(y)$.
Como $\Coker(v)$ é coerente \sref[0]{0.5.3.4}, seu suporte é
fechado \sref[0]{0.5.2.2}; portanto, existe uma vizinhança aberta
$W\subset f^{-1}(V)$ de $f^{-1}(y)$ sobre a qual $v$ é sobrejetivo.
Como $f$ é fechado, após restringir em torno de $y$, podemos escolher uma
vizinhança aberta $U\subset V$ com $f^{-1}(U)\subset W$.
A conclusão se segue agora de \sref[II]{II.4.2.3}.
\end{proof}


\subsection{Morfismos finitos de esquemas formais}
\label{subsection:III.4.8}

\begin{proposition}[4.8.1]
\label{III.4.8.1}
Seja $\mathfrak{Y}$ um esquema formal localmente noetheriano,
$\sh{K}$ um ideal de definição de $\mathfrak{Y}$, e
$f:\mathfrak{X}\to\mathfrak{Y}$ um morfismo de esquemas formais.
As condições seguintes são equivalentes:
\begin{enumerate}
  \item[\rm{a)}] $\mathfrak{X}$ é localmente noetheriano, $f$ é um
    morfismo ádico \sref[I]{I.10.12.1}, e, se
    $\sh{J}=f^*(\sh{K})\sh{O}_{\mathfrak{X}}$, o morfismo
    \[
      f_0:
      (\mathfrak{X},\sh{O}_{\mathfrak{X}}/\sh{J})
      \longrightarrow
      (\mathfrak{Y},\sh{O}_{\mathfrak{Y}}/\sh{K})
    \]
    induzido por $f$ é finito.
  \item[\rm{b)}] $\mathfrak{X}$ é localmente noetheriano e é o
    limite indutivo de um sistema indutivo ádico de $(Y_n)$, $(X_n)$,
    tal que o morfismo $X_0\to Y_0$ é finito.
  \item[\rm{c)}] Todo ponto de $\mathfrak{Y}$ possui uma vizinhança aberta
    formal afim noetheriana $V$ tal que $f^{-1}(V)$ é
    um aberto formal afim e
    $\Gamma(f^{-1}(V),\sh{O}_{\mathfrak{X}})$ é um
    $\Gamma(V,\sh{O}_{\mathfrak{Y}})$-módulo de tipo finito.
\end{enumerate}
\end{proposition}

\begin{proof}
Segue imediatamente de \sref[I]{I.10.12.3} que \emph{a)} implica
\emph{b)}. Para ver que \emph{b)} implica \emph{c)}, podemos supor
que
\[
  \mathfrak{Y}=\Spf(B),\qquad \mathfrak{X}=\Spf(A),
\]
onde $B$ é um anel noetheriano ádico e $\mathfrak{R}$ é um ideal
de definição de $B$. Por hipótese, $X_0$ é um esquema afim
cujo anel $A_0$ é um $B/\mathfrak{R}$-módulo de tipo finito
\sref[II]{II.6.1.3}. Por \sref[I]{I.5.1.9}, cada $X_n$ é um esquema
afim; se $A_n$ denota seu anel, a condição \emph{b)} implica que,
para $m\leq n$,
\[
  A_m\simeq A_n/\mathfrak{R}^{m+1}A_n.
\]
Segue-se que $\mathfrak{X}\simeq\Spf(A)$, onde
$A=\varprojlim A_n$, e a conclusão se segue de
\sref[0]{0.7.2.9}.

Finalmente, para provar que \emph{c)} implica \emph{a)}, podemos restringir-nos
novamente ao caso
\[
  \mathfrak{Y}=\Spf(B),\qquad \mathfrak{X}=\Spf(A),
\]
onde $A$ é uma $B$-álgebra finita. Como
$A/\mathfrak{R}A$ é então uma
$B/\mathfrak{R}$-álgebra finita, \sref[I]{I.10.10.9} mostra que se
satisfazem as condições de \emph{a)}.
\end{proof}

\begin{definition}[4.8.2]
\label{III.4.8.2}
Quando se satisfazem as condições equivalentes \emph{a)}, \emph{b)} e \emph{c)}
de \sref{III.4.8.1}, diz-se que o morfismo $f$ é
\emph{finito}; também se diz que $\mathfrak{X}$ é um
\emph{$\mathfrak{Y}$-esquema formal finito}, ou um esquema
formal \emph{finito sobre $\mathfrak{Y}$}.
\end{definition}

\begin{proposition}[4.8.3]
\label{III.4.8.3}
\begin{enumerate}
  \item[\rm{(i)}] Uma imersão fechada de esquemas formais localmente
    noetherianos é um morfismo finito.
  \item[\rm{(ii)}] O composto de dois morfismos finitos de esquemas formais
    localmente noetherianos é um morfismo finito.
  \item[\rm{(iii)}] Sejam $\mathfrak{X}$, $\mathfrak{Y}$ e
    $\mathfrak{S}$ esquemas formais localmente noetherianos,
    $f:\mathfrak{X}\to\mathfrak{S}$ um morfismo finito, e
    $g:\mathfrak{Y}\to\mathfrak{S}$ um morfismo. Então a projeção
    \[
      \mathfrak{X}\times_{\mathfrak{S}}\mathfrak{Y}
      \longrightarrow\mathfrak{Y}
    \]
    é finita.
  \item[\rm{(iv)}] Seja $\mathfrak{S}$ um esquema formal localmente
    noetheriano, e sejam $\mathfrak{X}'$ e $\mathfrak{Y}'$
    esquemas formais localmente noetherianos tais que
    $\mathfrak{X}'\times_{\mathfrak{S}}\mathfrak{Y}'$ é localmente
    noetheriano. Se $\mathfrak{X}$ e $\mathfrak{Y}$ são
    $\mathfrak{S}$-esquemas formais localmente noetherianos e
    \[
      f:\mathfrak{X}\to\mathfrak{X}',\qquad
      g:\mathfrak{Y}\to\mathfrak{Y}'
    \]
    são $\mathfrak{S}$-morfismos finitos, então
    $f\times_{\mathfrak{S}}g$ é finito.
  \item[\rm{(v)}] Sejam
    \[
      f:\mathfrak{X}\to\mathfrak{Y},\qquad
      g:\mathfrak{Y}\to\mathfrak{Z}
    \]
    morfismos de esquemas formais localmente noetherianos tais que
    $g$ é de tipo finito e separado. Se $g\circ f$ é finito,
    então $f$ é finito.
\end{enumerate}
\end{proposition}

\begin{proof}
A afirmação \emph{(i)} é imediata. Utilizando o critério \emph{a)} de
\sref{III.4.8.1}, as demais afirmações se reduzem imediatamente às
proposições correspondentes para morfismos de esquemas ordinários
\sref[II]{II.6.1.5}. Deixamos os detalhes ao leitor, seguindo
o modelo de \sref[I]{I.10.13.5}.
\end{proof}

\begin{corollary}[4.8.4]
\label{III.4.8.4}
Sob as hipóteses de \sref[I]{I.10.9.9}, se $f$ é um morfismo
finito, também o é sua extensão $\widehat{f}$ às
completamentos.
\end{corollary}

\begin{corollary}[4.8.5]
\label{III.4.8.5}
Seja $\mathfrak{X}$ um esquema formal finito sobre
$\mathfrak{Y}$, e seja
$f:\mathfrak{X}\to\mathfrak{Y}$ o morfismo estrutural. Para
todo subconjunto aberto $U\subset\mathfrak{Y}$, o esquema formal
$f^{-1}(U)$ é finito sobre $U$.
\end{corollary}

\begin{proposition}[4.8.6]
\label{III.4.8.6}
Se $f:\mathfrak{X}\to\mathfrak{Y}$ é um morfismo finito de esquemas
formais localmente noetherianos, então
$f_*(\sh{O}_{\mathfrak{X}})$ é uma
$\sh{O}_{\mathfrak{Y}}$-álgebra coerente.
\end{proposition}

\begin{proof}
Podemos considerar $f$ como o limite indutivo de um sistema indutivo
$(f_n)$ de morfismos $f_n:X_n\to Y_n$. Demonstremos que os
$f_n$ são finitos e que $f_*(\sh{O}_{\mathfrak{X}})$ é
isomorfo ao limite projetivo dos
$(f_n)_*(\sh{O}_{X_n})$; isto provará a afirmação
\sref[I]{I.10.10.5}. Basta restringir-se ao caso
\[
  \mathfrak{Y}=\Spf(B),\qquad \mathfrak{X}=\Spf(A),
\]
e observar que, se $\mathfrak{R}$ é um ideal de definição
de $B$ e $A$ é um $B$-módulo de tipo finito, então
$A/\mathfrak{R}^{n+1}A$ é um módulo de tipo finito sobre
$B/\mathfrak{R}^{n+1}B$, e $A$ é o limite projetivo dos
$A/\mathfrak{R}^{n+1}A$.
\end{proof}

Reciprocamente:

\begin{proposition}[4.8.7]
\label{III.4.8.7}
Seja $\mathfrak{Y}$ um esquema formal localmente noetheriano e
$\sh{A}$ uma $\sh{O}_{\mathfrak{Y}}$-álgebra coerente. Existe
um esquema formal $\mathfrak{X}$ finito sobre $\mathfrak{Y}$,
único salvo $\mathfrak{Y}$-isomorfismo único, tal que
\[
  f_*(\sh{O}_{\mathfrak{X}})=\sh{A},
\]
onde $f:\mathfrak{X}\to\mathfrak{Y}$ é o morfismo estrutural.
\end{proposition}

\begin{proof}
\oldpage[III]{148}
Seja $\sh{K}$ um ideal de definição de $\mathfrak{Y}$, e ponhamos
\[
  Y_n=(\mathfrak{Y},
       \sh{O}_{\mathfrak{Y}}/\sh{K}^{n+1}),
  \qquad
  \sh{A}_n=\sh{A}/\sh{K}^{n+1}\sh{A}.
\]
Claramente, $\sh{A}_n$ é uma $\sh{O}_{Y_n}$-álgebra finita e, portanto,
define um $Y_n$-esquema finito
\[
  X_n=\Spec(\sh{A}_n)
\]
\sref[II]{II.6.1.3}. Para $m\leq n$, o homomorfismo canônico
sobrejetivo
\[
  h_{mn}:\sh{A}_n\longrightarrow\sh{A}_m
\]
define um morfismo $u_{mn}:X_m\to X_n$ para o qual o diagrama
\[
  \xymatrix{
    X_n\ar[d]_{f_n} &
    X_m\ar[l]_{u_{mn}}\ar[d]^{f_m}\\
    Y_n &
    Y_m\ar[l]
  }
\]
é comutativo, onde $f_n$ denota o morfismo estrutural; além disso,
identifica $X_m$ com
$X_n\times_{Y_n}Y_m$, como se segue imediatamente de
\sref[II]{II.1.4.6}. O esquema formal $\mathfrak{X}$ obtido
como limite indutivo do sistema indutivo $(X_n)$ é, portanto,
localmente noetheriano, e seu morfismo estrutural
\[
  f:\mathfrak{X}\longrightarrow\mathfrak{Y},
\]
o limite indutivo do sistema $(f_n)$, é finito
\sref{III.4.8.1} \sref[I]{I.10.12.3.1}. Além disso, a demonstração de
\sref{III.4.8.6} mostra que
$f_*(\sh{O}_{\mathfrak{X}})$ é o limite projetivo das
$\sh{A}_n$, e é, portanto, igual a $\sh{A}$
\sref[I]{I.10.10.6}. A afirmação de unicidade se segue do
resultado mais geral seguinte.
\end{proof}

\begin{proposition}[4.8.8]
\label{III.4.8.8}
Seja $\mathfrak{Y}$ um esquema formal localmente noetheriano, e
sejam $\mathfrak{X}$ e $\mathfrak{X}'$
$\mathfrak{Y}$-esquemas formais finitos sobre $\mathfrak{Y}$, com
morfismos estruturais
\[
  f:\mathfrak{X}\to\mathfrak{Y},\qquad
  f':\mathfrak{X}'\to\mathfrak{Y}.
\]
Existe uma bijeção canônica
\[
  \Hom_{\mathfrak{Y}}(\mathfrak{X},\mathfrak{X}')
  \xrightarrow{\ \sim\ }
  \Hom_{\sh{O}_{\mathfrak{Y}}}
  \bigl(f'_*(\sh{O}_{\mathfrak{X}'}),
        f_*(\sh{O}_{\mathfrak{X}})\bigr)\footnotemark.
\]
\footnotetext{A última expressão denota o conjunto de homomorfismos de
$\sh{O}_{\mathfrak{Y}}$-álgebras
$f'_*(\sh{O}_{\mathfrak{X}'})\to f_*(\sh{O}_{\mathfrak{X}})$.}
\end{proposition}

\begin{proof}
A aplicação $h\mapsto\sh{A}(h)$ se define como em
\sref[II]{II.1.1.2}. Para provar que é bijetiva, reduzimo-nos de
imediato ao caso em que
$\mathfrak{Y}=\Spf(B)$ é um esquema formal afim noetheriano.
Então
\[
  \mathfrak{X}=\Spf(A),\qquad
  \mathfrak{X}'=\Spf(A'),
\]
onde $A$ e $A'$ são $B$-álgebras finitas, e
\[
  f_*(\sh{O}_{\mathfrak{X}})=A^\Delta,\qquad
  f'_*(\sh{O}_{\mathfrak{X}'})=A'^\Delta.
\]
A conclusão se segue da correspondência bijetiva entre
$\mathfrak{Y}$-morfismos
$\mathfrak{X}\to\mathfrak{X}'$ e homomorfismos contínuos de
$B$-álgebras $A'\to A$ \sref[I]{I.10.2.2}, junto com a
correspondência bijetiva entre homomorfismos de $B$-módulos
$A'\to A$ e homomorfismos de
$\sh{O}_{\mathfrak{Y}}$-módulos
$A'^\Delta\to A^\Delta$ \sref[I]{I.10.10.2.3}.
\end{proof}

\begin{corollary}[4.8.9]
\label{III.4.8.9}
Sob a bijeção canônica de \sref{III.4.8.8}, as imersões
fechadas $\mathfrak{X}\to\mathfrak{X}'$ correspondem a
homomorfismos sobrejetivos de $\sh{O}_{\mathfrak{Y}}$-álgebras
\[
  f'_*(\sh{O}_{\mathfrak{X}'})
  \longrightarrow
  f_*(\sh{O}_{\mathfrak{X}}).
\]
\end{corollary}

\begin{proof}
A questão é local sobre $\mathfrak{Y}$ e, portanto, reduz-se à
definição das imersões fechadas de esquemas formais localmente
noetherianos \sref[I]{I.10.14.2}.
\end{proof}

\begin{corollary}[4.8.10]
\label{III.4.8.10}
Com a notação e as hipóteses de \sref{III.4.8.1}, um morfismo
ádico $f$ é uma imersão fechada se e somente se $f_0$ é uma
imersão fechada de esquemas ordinários.
\end{corollary}

\begin{proof}
Isto se segue imediatamente de \sref{III.4.8.9} e do critério de
sobrejetividade para um homomorfismo de
$\sh{O}_{\mathfrak{Y}}$-módulos coerentes \sref[I]{I.10.11.5}.
\end{proof}

\begin{proposition}[4.8.11]
\label{III.4.8.11}
Um morfismo $f:\mathfrak{X}\to\mathfrak{Y}$ de esquemas formais
localmente noetherianos
\oldpage[III]{149}
é finito se e somente se é próprio e todas suas fibras
$f^{-1}(y)$, para $y\in\mathfrak{Y}$, são finitas.
\end{proposition}

\begin{proof}
Pelas definições \sref{III.3.4.1} e \sref{III.4.8.2}, a
afirmação reduz-se imediatamente à mesma afirmação para $f_0$
(com a notação de \sref{III.4.8.1}), que é precisamente
\sref{III.4.4.2}.
\end{proof}

\section{Um teorema de existência para os feixes algébricos coerentes}
\label{section:III.5}


\subsection{Enunciado do teorema}
\label{subsection:III.5.1}

\begin{env}[5.1.1]
\label{III.5.1.1}
Seja $A$ um anel noetheriano ádico e seja $\mathfrak{J}$ um ideal de
definição de $A$, de modo que $A$ seja separado e completo para a
topologia $\mathfrak{J}$-ádica. Se $Y=\Spec(A)$, o pré-esquema formal
afim $\Spf(A)$ identifica-se com o completamento
$\widehat{Y}$ de $Y$ ao longo do subconjunto fechado
$Y'=\mathrm{V}(\mathfrak{J})$ \sref[I]{I.10.10.1}. Seja $X$ um
$Y$-pré-esquema ordinário de tipo finito, de morfismo estrutural
$f:X\to Y$. Denotemos por $\widehat{X}$ o completamento de $X$ ao longo
do subconjunto fechado $X'=f^{-1}(Y')$, ou, equivalentemente, o
$\widehat{Y}$-pré-esquema formal $X\times_Y\widehat{Y}$, e por
$\widehat{f}:\widehat{X}\to\widehat{Y}$ a extensão de $f$ às
completamentos. Finalmente, para todo $\sh{O}_X$-módulo coerente
$\sh{F}$, denotemos por $\widehat{\sh{F}}$ seu completamento
$\sh{F}_{/X'}$, que é um
$\sh{O}_{\widehat{X}}$-módulo coerente.
\end{env}

\begin{proposition}[5.1.2]
\label{III.5.1.2}
Com as hipóteses e a notação de \sref{III.5.1.1}, seja $\sh{F}$
um $\sh{O}_X$-módulo coerente cujo suporte é próprio sobre $Y$
\sref[II]{II.5.4.10}. Os homomorfismos canônicos
\sref{III.4.1.4}
\[
  \rho_i:
  \HH^i(X,\sh{F})
  \longrightarrow
  \HH^i(\widehat{X},\widehat{\sh{F}})
\]
são isomorfismos.
\end{proposition}

\begin{proof}
O $A$-módulo $\HH^i(X,\sh{F})$ é de tipo finito
\sref{III.3.2.4} e, portanto, identifica-se com seu completamento
separado para a topologia $\mathfrak{J}$-pré-ádica
\sref[0]{0.7.3.6}. A proposição é, por conseguinte, um caso particular
de \sref{III.4.1.10}.
\end{proof}

Recordemos que os isomorfismos canônicos $\rho_i$ comutam com os
homomorfismos de conexão associados a toda sucessão exata
\[
  0\longrightarrow\sh{F}'
  \longrightarrow\sh{F}
  \longrightarrow\sh{F}''
  \longrightarrow0
\]
de $\sh{O}_X$-módulos coerentes
\sref[0]{0.12.1.6} \sref[I]{I.10.8.9}.

\begin{corollary}[5.1.3]
\label{III.5.1.3}
Sejam $\sh{F}$ e $\sh{G}$ $\sh{O}_X$-módulos coerentes tais que
a interseção de seus suportes seja própria sobre $Y$. O
homomorfismo canônico
\[
  \Hom_{\sh{O}_X}(\sh{F},\sh{G})
  \longrightarrow
  \Hom_{\sh{O}_{\widehat{X}}}
  (\widehat{\sh{F}},\widehat{\sh{G}})
  \tag{5.1.3.1}
  \label{III.5.1.3.1.eq}
\]
que envia um homomorfismo $u:\sh{F}\to\sh{G}$ a seu completamento
$\widehat{u}:\widehat{\sh{F}}\to\widehat{\sh{G}}$, é um
isomorfismo. Além disso, quando $f$ é fechado,
$\widehat{u}$ é injetivo (respectivamente, sobrejetivo) se e somente
se $u$ é injetivo (respectivamente, sobrejetivo).
\end{corollary}

\begin{proof}
A primeira afirmação é um caso particular de \sref{III.4.5.3}; também
decorre do fato de que a fonte de \eqref{III.5.1.3.1.eq} é
um $A$-módulo de tipo finito e, portanto, identifica-se com seu
completamento separado. Para a segunda afirmação, observemos que, por
\sref[I]{I.10.8.14}, $\widehat{u}$ é injetivo (respectivamente,
sobrejetivo) se e somente se existe uma vizinhança de $X'$ em
$X$ sobre a qual $u$ é injetivo (respectivamente, sobrejetivo). A
conclusão deduz-se então do lema seguinte.
\end{proof}

\oldpage[III]{150}

\begin{lemma}[5.1.3.1]
\label{III.5.1.3.1}
Sob as hipóteses de \sref{III.5.1.1}, suponhamos além disso que
$f$ é fechado. Toda vizinhança de $X'$ em $X$ é então igual
a $X$.
\end{lemma}

\begin{proof}
Reduzimo-nos primeiro ao caso $f(X)=Y$. Com efeito, por hipótese,
$f(X)$ é um subconjunto fechado $Z$ de $Y$. Podemos substituir $f$ por
$f_{\mathrm{red}}$ \sref[I]{I.6.3.4} e supor assim que $X$
e $Y$ são reduzidos; podemos então substituir $Y$ pelo subpré-esquema
fechado reduzido cujo espaço subjacente é $Z$
\sref[I]{I.5.2.2}, pois todo ideal de $A$ é fechado e todo
anel quociente de $A$ é, portanto, ádico e noetheriano.

Tem-se então $f(X')=Y'$. Se $V$ é uma vizinhança aberta de $X'$
em $X$, o conjunto $f(X-V)$ é fechado em $Y$ e não intersecta $Y'$.
Isto é impossível, a menos que $X-V$ seja vazio, pois $\mathfrak{J}$
está contido no radical de $A$
\sref[I]{I.1.1.15} \sref[0]{0.7.1.10}. Portanto, $V=X$.
\end{proof}

Se nos restringimos aos $\sh{O}_X$-módulos coerentes cujo
suporte é próprio sobre $Y$, \sref{III.5.1.3} afirma, em linguagem
categórica, que o funtor
\[
  \sh{F}\longmapsto\widehat{\sh{F}}
\]
é plenamente fiel da categoria de tais $\sh{O}_X$-módulos para
a categoria dos $\sh{O}_{\widehat{X}}$-módulos coerentes. Fornece, assim,
uma equivalência da primeira categoria com uma subcategoria plena
da segunda \sref[0]{0.8.1.6}. O teorema de existência
mostrará que, quando $X$ é próprio sobre $Y$, esta subcategoria é em
realidade, a categoria de todos os
$\sh{O}_{\widehat{X}}$-módulos coerentes. Mais precisamente:

\begin{theorem}[5.1.4]
\label{III.5.1.4}
Seja $A$ um anel noetheriano ádico, seja $Y=\Spec(A)$, seja
$\mathfrak{J}$ um ideal de definição de $A$, e ponhamos
$Y'=\mathrm{V}(\mathfrak{J})$. Seja $f:X\to Y$ um morfismo separado
de tipo finito e ponhamos $X'=f^{-1}(Y')$. Sejam
\[
  \widehat{Y}=Y_{/Y'}=\Spf(A),\qquad
  \widehat{X}=X_{/X'}
\]
os completamentos de $Y$ e $X$ ao longo de $Y'$ e $X'$, e seja
$\widehat{f}:\widehat{X}\to\widehat{Y}$ a extensão de $f$ aos
completamentos. Então o funtor
\[
  \sh{F}\longmapsto\sh{F}_{/X'}=\widehat{\sh{F}}
\]
é uma equivalência da categoria dos
$\sh{O}_X$-módulos coerentes cujo suporte é próprio sobre $\Spec(A)$ com a
categoria dos $\sh{O}_{\widehat{X}}$-módulos coerentes cujo suporte
é próprio sobre $\Spf(A)$.
\end{theorem}

Em outras palavras, tendo em conta \sref{III.5.1.3}:

\begin{corollary}[5.1.5]
\label{III.5.1.5}
Um $\sh{O}_{\widehat{X}}$-módulo é isomorfo ao completamento de
um $\sh{O}_X$-módulo coerente cujo suporte é próprio sobre
$\Spec(A)$ se e somente se é coerente e seu suporte é próprio
sobre $\Spf(A)$.
\end{corollary}

O caso mais importante é o seguinte:

\begin{corollary}[5.1.6]
\label{III.5.1.6}
Suponhamos que $X$ é próprio sobre $Y=\Spec(A)$. Então o funtor
$\sh{F}\mapsto\widehat{\sh{F}}$ é uma equivalência entre a
categoria dos $\sh{O}_X$-módulos coerentes e a categoria dos
$\sh{O}_{\widehat{X}}$-módulos coerentes.
\end{corollary}

\begin{env}[5.1.7]
\label{III.5.1.7}
\emph{Escólio.}
Tendo em conta a caracterização dos feixes coerentes sobre
pré-esquemas formais \sref[I]{I.10.11.3}, vemos que, sob as
condições de \sref{III.5.1.1}, dar um
$\sh{O}_X$-módulo coerente cujo suporte é próprio sobre $\Spec(A)$
equivale---se
\[
  Y_n=\Spec(A/\mathfrak{J}^{n+1}),\qquad
  X_n=X\times_Y Y_n
\]
---a dar um sistema projetivo $(\sh{F}_n)$ de
$\sh{O}_{X_n}$-módulos coerentes tal que, para $m\leq n$,
\[
  \sh{F}_m
  =
  \sh{F}_n\otimes_{\sh{O}_{Y_n}}\sh{O}_{Y_m},
  \qquad\text{ou, equivalentemente}\qquad
  \sh{F}_m
  =
  \sh{F}_n/\mathfrak{J}^{m+1}\sh{F}_n,
\]
e tal que o suporte de $\sh{F}_0$ seja um subconjunto de $X_0$
próprio sobre $Y_0$. Por \sref[I]{I.10.11.4}, os homomorfismos de
$\sh{O}_X$-módulos coerentes são interpretados de modo análogo como
homomorfismos de sistemas projetivos de
$\sh{O}_{X_n}$-módulos coerentes.

\oldpage[III]{151}

Em todas as aplicações conhecidas, $A$ é de fato um anel
noetheriano ádico local. Assim, os $Y_n$ são espectros de anéis
artinianos locais, e os resultados desta subseção e das
anteriores reduzem, em grande medida, a geometria algébrica sobre um
anel noetheriano ádico local à geometria algébrica sobre anéis
artinianos locais.
\end{env}

\begin{corollary}[5.1.8]
\label{III.5.1.8}
Sob as hipóteses de \sref{III.5.1.4}, a aplicação
\[
  Z\longmapsto\widehat{Z}=Z_{/(Z\cap X')}
\]
é uma bijeção do conjunto dos subpré-esquemas fechados $Z$ de $X$
que são próprios sobre $Y$ com o conjunto dos subpré-esquemas
formais fechados de $\widehat{X}$ que são próprios sobre
$\widehat{Y}$.
\end{corollary}

\begin{proof}
Um subpré-esquema formal fechado de $\widehat{X}$ é da forma
\[
  \bigl(T,
    (\sh{O}_{\widehat{X}}/\sh{A})|_T\bigr),
\]
onde $\sh{A}$ é um ideal coerente de
$\sh{O}_{\widehat{X}}$ \sref[I]{I.10.14.2}. Se $T$ é próprio sobre
$\widehat{Y}$, \sref{III.5.1.4} mostra que
$\sh{O}_{\widehat{X}}/\sh{A}$ é isomorfo a um
$\sh{O}_{\widehat{X}}$-módulo da forma
$\widehat{\sh{F}}$, onde $\sh{F}$ é um
$\sh{O}_X$-módulo coerente cujo suporte é próprio sobre $Y$. Além disso,
\sref{III.5.1.3} mostra que o homomorfismo canônico
\[
  \sh{O}_{\widehat{X}}
  \longrightarrow
  \sh{O}_{\widehat{X}}/\sh{A}
\]
é da forma $\widehat{u}$, onde
$u:\sh{O}_X\to\sh{F}$ é um homomorfismo sobrejetivo de
$\sh{O}_X$-módulos. Por conseguinte,
$\sh{F}=\sh{O}_X/\sh{N}$ para um ideal coerente $\sh{N}$ de
$\sh{O}_X$, e
$\sh{A}=\widehat{\sh{N}}$ \sref[I]{I.10.8.8}. A conclusão se
deduz agora de \sref[I]{I.10.14.7}.
\end{proof}


\subsection{Demonstração do teorema de existência: caso projetivo e quase-projetivo}
\label{subsection:III.5.2}

\begin{env}[5.2.1]
\label{III.5.2.1}
Sob as condições de \sref{III.5.1.4}, diremos
provisoriamente que um $\sh{O}_{\widehat{X}}$-módulo coerente é
\emph{algebrizável} se é isomorfo ao completamento
$\widehat{\sh{F}}$ de um $\sh{O}_X$-módulo coerente $\sh{F}$ cujo
suporte é próprio sobre $Y$.
\end{env}

\begin{lemma}[5.2.2]
\label{III.5.2.2}
Sejam $\sh{F}'$ e $\sh{G}'$ $\sh{O}_{\widehat{X}}$-módulos
algebrizáveis. Para todo homomorfismo
$u:\sh{F}'\to\sh{G}'$, os módulos
$\Ker(u)$, $\Im(u)$ e $\Coker(u)$ são algebrizáveis.
\end{lemma}

\begin{proof}
Escrevamos
\[
  \sh{F}'=\widehat{\sh{F}},\qquad
  \sh{G}'=\widehat{\sh{G}},
\]
onde $\sh{F}$ e $\sh{G}$ são $\sh{O}_X$-módulos coerentes cujos
suportes são próprios sobre $Y$. Por \sref{III.5.1.3},
$u=\widehat{v}$ para um homomorfismo
$v:\sh{F}\to\sh{G}$. Como o completamento é exato,
$\Ker(\widehat{v})\simeq\widehat{\Ker(v)}$; e como o suporte
de $\Ker(v)$ está contido no de $\sh{F}$, deduz-se que
$\Ker(u)$ é algebrizável. Os argumentos para $\Im(u)$ e
$\Coker(u)$ são os mesmos.
\end{proof}

\begin{proposition}[5.2.3]
\label{III.5.2.3}
Seja $A$ um anel noetheriano ádico, seja $\mathfrak{J}$ um ideal de
definição de $A$, seja $\mathfrak{Y}=\Spf(A)$, e seja
$f:\mathfrak{X}\to\mathfrak{Y}$ um morfismo próprio de pré-esquemas
formais. Ponhamos
\[
  Y_k=\Spec(A/\mathfrak{J}^{k+1}),\qquad
  X_k=\mathfrak{X}\times_{\mathfrak{Y}}Y_k,
\]
e, para todo $\sh{O}_{\mathfrak{X}}$-módulo $\sh{F}$, ponhamos
\[
  \sh{F}_k
  =
  \sh{F}\otimes_{\sh{O}_{\mathfrak{X}}}\sh{O}_{X_k}
  =
  \sh{F}/\mathfrak{J}^{k+1}\sh{F}.
\]
Seja $\sh{L}$ um $\sh{O}_{\mathfrak{X}}$-módulo invertível, e
suponhamos que
$\sh{L}_0=\sh{L}/\mathfrak{J}\sh{L}$ é um
$\sh{O}_{X_0}$-módulo amplo. Para todo
$\sh{O}_{\mathfrak{X}}$-módulo $\sh{F}$ e todo inteiro $n$, ponhamos
\[
  \sh{F}(n)=\sh{F}\otimes\sh{L}^{\otimes n}.
\]
Então, para todo $\sh{O}_{\mathfrak{X}}$-módulo coerente $\sh{F}$,
existe um inteiro $n_0$ tal que, para todo $n\geq n_0$:
\begin{enumerate}
  \item[\rm{(i)}] o homomorfismo canônico
    \[
      \HH^0(\mathfrak{X},\sh{F}(n))
      \longrightarrow
      \HH^0(\mathfrak{X},\sh{F}_k(n))
    \]
    é sobrejetivo para todo $k\geq0$;
  \item[\rm{(ii)}]
    $\HH^q(\mathfrak{X},\sh{F}(n))=0$ para todo $q>0$.
\end{enumerate}
\end{proposition}

\begin{proof}
Os espaços subjacentes a $\mathfrak{X}$ e $X_0$ são os mesmos.
Os feixes
\[
  \sh{M}_k
  =
  \mathfrak{J}^k\sh{F}/\mathfrak{J}^{k+1}\sh{F},
\]
por serem anulados por $\mathfrak{J}$, podem ser considerados
$\sh{O}_{X_0}$-módulos coerentes \sref[0]{0.5.3.10}. Além disso, se
\[
  \sh{M}_k(n)
  =
  \sh{M}_k\otimes_{\sh{O}_{X_0}}\sh{L}_0^{\otimes n},
\]
então
\[
  \sh{M}_k(n)
  =
  \mathfrak{J}^k\sh{F}(n)/
  \mathfrak{J}^{k+1}\sh{F}(n).
\]
Como $\sh{L}_0$ é amplo para $f_0:X_0\to Y_0$ e $f_0$ é
próprio,
\oldpage[III]{152}
$f_0$ é projetivo \sref[II]{II.5.5.4}. Seja
\[
  S=\bigoplus_{k\geq0}
  \mathfrak{J}^k/\mathfrak{J}^{k+1}
\]
a $A$-álgebra graduada associada à
filtração $\mathfrak{J}$-ádica de $A$. É de tipo finito
porque $A$ é noetheriano. Ponhamos
\[
  \sh{S}'=f_0^*(\widetilde{S}),\qquad
  \sh{M}=\bigoplus_{k\geq0}\sh{M}_k.
\]
Então $\sh{S}'$ é uma $\sh{O}_{X_0}$-álgebra quase-coerente de
tipo finito, e $\sh{M}$ é um $\sh{S}'$-módulo graduado
quase-coerente de tipo finito, pois $\sh{F}_0$ é coerente e
gera $\sh{M}$ como $\sh{S}'$-módulo. O Teorema
\sref{III.2.4.1},~\emph{(ii)}, fornece, portanto, um inteiro $n_0$
tal que, para $n\geq n_0$ e todo $k$,
\[
  \HH^q(X_0,\sh{M}_k(n))=0
  \qquad(q>0).
  \tag{5.2.3.1}
  \label{III.5.2.3.1}
\]
Considerado agora como $\sh{O}_{\mathfrak{X}}$-módulo,
$\sh{M}_k(n)$ satisfaz, por conseguinte, também
\[
  \HH^q(\mathfrak{X},\sh{M}_k(n))=0
  \qquad(q>0,\ n\geq n_0).
\]

Apliquemos a sucessão exata de cohomologia a
\[
  0\longrightarrow
  \frac{\mathfrak{J}^k\sh{F}(n)}
       {\mathfrak{J}^{k+1}\sh{F}(n)}
  \longrightarrow
  \frac{\mathfrak{J}^h\sh{F}(n)}
       {\mathfrak{J}^{k+1}\sh{F}(n)}
  \longrightarrow
  \frac{\mathfrak{J}^h\sh{F}(n)}
       {\mathfrak{J}^{k}\sh{F}(n)}
  \longrightarrow0.
\]
Para $0\leq h<k$, $n\geq n_0$ e $q>0$, uma indução sobre $k-h$
dá
\[
  \HH^q\left(
    \mathfrak{X},
    \frac{\mathfrak{J}^h\sh{F}(n)}
         {\mathfrak{J}^{k}\sh{F}(n)}
  \right)=0,
  \tag{5.2.3.2}
  \label{III.5.2.3.2}
\]
e, em particular, para $h=0$,
\[
  \HH^q(\mathfrak{X},\sh{F}_k(n))=0
  \qquad(n\geq n_0,\ k\geq0,\ q>0).
  \tag{5.2.3.3}
  \label{III.5.2.3.3}
\]
Outra parte da sucessão exata de cohomologia, com $h=0$, dá
\[
  \HH^0(\mathfrak{X},\sh{F}_{k+1}(n))
  \longrightarrow
  \HH^0(\mathfrak{X},\sh{F}_k(n))
  \longrightarrow
  \HH^1\left(
    \mathfrak{X},
    \frac{\mathfrak{J}^k\sh{F}(n)}
         {\mathfrak{J}^{k+1}\sh{F}(n)}
  \right)
  =0.
  \tag{5.2.3.4}
  \label{III.5.2.3.4}
\]
Deduz-se que, para $h\leq k$, a aplicação canônica
\[
  \HH^0(\mathfrak{X},\sh{F}_k(n))
  \longrightarrow
  \HH^0(\mathfrak{X},\sh{F}_h(n))
  \tag{5.2.3.5}
  \label{III.5.2.3.5}
\]
é sobrejetiva. Assim, para todo $q$, o sistema projetivo
\[
  \bigl(\HH^q(\mathfrak{X},\sh{F}_k(n))\bigr)_{k\geq0}
\]
satisfaz a condição (ML) quando $n\geq n_0$.

Além disso, todo aberto formal afim $U$ de $\mathfrak{X}$ é um
aberto afim de cada $X_k$ \sref[I]{I.10.5.2}. Portanto,
$\HH^q(U,\sh{F}_k(n))=0$ para $q>0$ \sref[I]{I.1.3.1}, e
\[
  \HH^0(U,\sh{F}_k(n))
  \longrightarrow
  \HH^0(U,\sh{F}_h(n))
\]
é sobrejetivo para $h\leq k$ \sref[I]{I.1.3.9}. Portanto, são satisfeitas
as hipóteses de \sref[0]{0.13.3.1}. Para $n\geq n_0$
obtemos:
\begin{enumerate}
  \item[$1^\circ$] para todo $q>0$, a aplicação
    \[
      \HH^q(\mathfrak{X},\sh{F}(n))
      \longrightarrow
      \varprojlim_k
      \HH^q(\mathfrak{X},\sh{F}_k(n))
    \]
    é bijetiva; portanto, por \eqref{III.5.2.3.3},
    $\HH^q(\mathfrak{X},\sh{F}(n))=0$;
  \item[$2^\circ$] o homomorfismo
    \[
      \HH^0(\mathfrak{X},\sh{F}(n))
      \longrightarrow
      \varprojlim_k
      \HH^0(\mathfrak{X},\sh{F}_k(n))
    \]
    é bijetivo. Como os homomorfismos
    \eqref{III.5.2.3.5} são sobrejetivos, também o é cada homomorfismo
    \[
      \varprojlim_k
      \HH^0(\mathfrak{X},\sh{F}_k(n))
      \longrightarrow
      \HH^0(\mathfrak{X},\sh{F}_h(n)).
    \]
\end{enumerate}
Isto demonstra as duas afirmações.
\end{proof}

\begin{corollary}[5.2.4]
\label{III.5.2.4}
Sob as hipóteses de \sref{III.5.2.3}, para todo
$\sh{O}_{\mathfrak{X}}$-módulo coerente $\sh{F}$ existe um inteiro $N$
tal que, para $n\geq N$, $\sh{F}(n)$ está gerado por suas seções
sobre
\oldpage[III]{153}
$\mathfrak{X}$. Equivalentemente, $\sh{F}$ é isomorfo a um quociente
de um $\sh{O}_{\mathfrak{X}}$-módulo da forma
$(\sh{O}_{\mathfrak{X}}(-n))^k$.
\end{corollary}

\begin{proof}
Como $X_0$ é noetheriano, a hipótese sobre $\sh{L}_0$ e
\sref[II]{II.4.5.5} fornecem um inteiro $n_0$ tal que
$\sh{F}_0(n)$ está gerado por suas seções sobre $\mathfrak{X}$ para
$n\geq n_0$. Podemos tomar $n_0$ suficientemente grande para que
\[
  \Gamma(\mathfrak{X},\sh{F}(n))
  \longrightarrow
  \Gamma(\mathfrak{X},\sh{F}_0(n))
\]
seja também sobrejetivo para $n\geq n_0$ \sref{III.5.2.3}. Existem,
portanto, um número finito de seções
$s_i\in\Gamma(\mathfrak{X},\sh{F}(n))$ cujas imagens geram
$\sh{F}_0(n)$ \sref[0]{0.5.2.3}. Como $\mathfrak{J}$ está
contido no ideal maximal do anel local em todo ponto
de $\mathfrak{X}$, o lema de Nakayama aplicado a esses anéis locais
mostra que os $s_i$ geram $\sh{F}(n)$
\sref[0]{0.5.1.1}.
\end{proof}

\begin{env}[5.2.5]
\label{III.5.2.5}
\emph{Demonstração do teorema de existência: caso projetivo.}
Com a notação de \sref{III.5.1.4}, suponhamos que $f$ é
projetivo, de modo que existe um $\sh{O}_X$-módulo amplo $\sh{L}$
\sref[I]{I.5.5.4}. Por definição,
\[
  X_n=\widehat{X}\times_{\widehat{Y}}Y_n
\]
é o subpré-esquema fechado
$X\times_Y Y_n=f^{-1}(Y_n)$ de $X$. Se
\[
  \sh{L}'=\sh{L}\otimes_{\sh{O}_X}\sh{O}_{\widehat{X}}
\]
é o completamento de $\sh{L}$, então
$\sh{L}'_0=\sh{L}/\mathfrak{J}\sh{L}$, considerado como
$\sh{O}_{X_0}$-módulo; este módulo é amplo
\sref[II]{II.4.6.13}[(i~\emph{bis})].

Podemos, portanto, aplicar \sref{III.5.2.4} a $\sh{L}'$ e a todo
$\sh{O}_{\widehat{X}}$-módulo coerente $\sh{F}$. Deduz-se que
$\sh{F}$ é isomorfo a um quociente de
\[
  \sh{G}=
  \bigl(\sh{L}'^{\otimes(-n)}\bigr)^k
\]
para inteiros adequados $n>0$ e $k>0$. O módulo $\sh{G}$ é o
completamento de $(\sh{L}^{\otimes(-n)})^k$
\sref[I]{I.10.8.10} e, portanto, é algebrizável. Seja
$u:\sh{G}\to\sh{F}$ a sobrejeção canônica, e ponhamos
$\sh{H}=\Ker(u)$; este é um
$\sh{O}_{\widehat{X}}$-módulo coerente \sref[0]{0.5.3.4}. Aplicando o mesmo
argumento a $\sh{H}$, obtemos um
$\sh{O}_{\widehat{X}}$-módulo algebrizável $\sh{K}$ e um homomorfismo
$v:\sh{K}\to\sh{G}$ tal que $\sh{H}=\Im(v)$. Assim,
$\sh{F}=\Coker(v)$, e $\sh{F}$ é algebrizável por
\sref{III.5.2.2}.
\end{env}

\begin{env}[5.2.6]
\label{III.5.2.6}
\emph{Demonstração do teorema de existência: caso quase-projetivo.}
Conservemos a notação de \sref{III.5.1.4} e suponhamos agora que $f$
é quase-projetivo. Existe um morfismo projetivo
$g:Z\to Y$ tal que $X$ identifica-se com o $Y$-pré-esquema
induzido sobre um subconjunto aberto de $Z$ \sref[II]{II.5.3.2}. Se
$Z'=g^{-1}(Y')$, então $X'=X\cap Z'$. Por conseguinte, o completamento
$\widehat{X}=X_{/X'}$ identifica-se com o pré-esquema formal
induzido por $\widehat{Z}=Z_{/Z'}$ sobre o subconjunto aberto $X\cap Z'$ de
$\widehat{Z}$ \sref[I]{I.10.8.5}.

Seja $\sh{F}'$ um $\sh{O}_{\widehat{X}}$-módulo coerente cujo
suporte $T'$ seja próprio sobre $\widehat{Y}$. Por definição, existe
um subpré-esquema fechado de $X'$ de espaço subjacente
$T'\subset X'$ tal que a restrição $T'\to Y'$ de $f$ é
própria. Portanto, $T'$ é próprio sobre $Y$ e, consequentemente, fechado em
$Z'$ \sref[II]{II.5.4.10}. Deduz-se que $\sh{F}'$ é induzido
sobre $\widehat{X}$ pelo $\sh{O}_{\widehat{Z}}$-módulo $\sh{G}'$
obtido pela colagem de $\sh{F}'$, sobre o aberto $\widehat{X}$, com o
feixe zero sobre o aberto $\widehat{Z}-T'$; ambos os feixes coincidem sobre
sua interseção $\widehat{X}-T'$.

O módulo $\sh{G}'$ é coerente. Por \sref{III.5.2.5}, existe um
$\sh{O}_Z$-módulo coerente $\sh{G}$ tal que
$\sh{G}'=\widehat{\sh{G}}$. Seja $T$ o suporte de $\sh{G}$;
então $T'=T\cap Z'$ \sref[I]{I.10.8.12}. Se $h$ é a restrição
de $g$ ao subpré-esquema fechado reduzido de $Z$ cujo espaço subjacente
é $T$, então
\[
  T'=h^{-1}(Y')=T\cap g^{-1}(Y').
\]
Assim, $X\cap T$ é um subconjunto aberto de $T$ que contém $T'$.

\oldpage[III]{154}
Como $h$ é próprio \sref[II]{II.5.4.2}, e portanto fechado,
\sref{III.5.1.3.1} dá $X\cap T=T$. Assim, $T\subset X$; e
como $T$ é fechado em $Z$, é próprio sobre $Y$. Se $\sh{F}$ é
o feixe induzido sobre $X$ por $\sh{G}$, seu completamento
$\widehat{\sh{F}}$ é induzido sobre $\widehat{X}$ por
$\widehat{\sh{G}}$ \sref[I]{I.10.8.4} e, portanto, é igual a
$\sh{F}'$. Isto termina a demonstração.
\end{env}


\subsection{Demonstração do teorema de existência: caso geral}
\label{subsection:III.5.3}

\begin{lemma}[5.3.1]
\label{III.5.3.1}
Sob as condições de \sref{III.5.1.4}, seja
\[
  0\longrightarrow\sh{H}\longrightarrow\sh{F}
  \longrightarrow\sh{G}\longrightarrow0
\]
uma sucessão exata de $\sh{O}_{\widehat{X}}$-módulos coerentes.
Se $\sh{G}$ e $\sh{H}$ são algebrizáveis, então $\sh{F}$ é
algebrizável.
\end{lemma}

\begin{proof}
Suponhamos que
\[
  \sh{G}=\widehat{\sh{B}},
  \qquad
  \sh{H}=\widehat{\sh{C}},
\]
onde $\sh{B}$ e $\sh{C}$ são $\sh{O}_X$-módulos coerentes cujos
suportes são próprios sobre $Y$. O módulo $\sh{F}$ define canonicamente
um elemento do $A$-módulo
\[
  \Ext^1_{\sh{O}_{\widehat{X}}}
  (\widehat{X};\widehat{\sh{B}},\widehat{\sh{C}})
\]
\sref[0]{0.12.3.2}. As hipóteses implicam que este $A$-módulo é
canonicamente isomorfo a
\[
  \Ext^1_{\sh{O}_X}(X;\sh{B},\sh{C})
\]
\sref{III.4.5.2}. Existe, portanto, uma sucessão exata
\[
  0\longrightarrow\sh{C}\longrightarrow\sh{A}
  \longrightarrow\sh{B}\longrightarrow0
\]
de $\sh{O}_X$-módulos coerentes cuja classe canônica é enviada para a
classe correspondente a $\sh{F}$. Por definição, tendo em conta
\sref[I]{I.10.8.8}, (ii), isto significa que $\sh{F}$ é isomorfo a
$\widehat{\sh{A}}$. O lema segue, pois,
$\Supp(\sh{A})$ está contido em
$\Supp(\sh{B})\cup\Supp(\sh{C})$ e é, portanto, próprio sobre $Y$.
\end{proof}

\begin{corollary}[5.3.2]
\label{III.5.3.2}
Sob as condições de \sref{III.5.1.1}, seja
$u:\sh{F}\to\sh{G}$ um homomorfismo de
$\sh{O}_{\widehat{X}}$-módulos coerentes. Se $\sh{G}$, $\Ker(u)$ e
$\Coker(u)$ são algebrizáveis, então $\sh{F}$ é algebrizável.
\end{corollary}

\begin{proof}
O Lema \sref{III.5.2.2}, aplicado a
$\sh{G}\to\Coker(u)$, mostra que $\Im(u)$ é algebrizável. Basta então
aplicar o Lema \sref{III.5.3.1} à sucessão exata
\[
  0\longrightarrow\Ker(u)\longrightarrow\sh{F}
  \longrightarrow\Im(u)\longrightarrow0.
\]
\end{proof}

\begin{lemma}[5.3.3]
\label{III.5.3.3}
Sob as condições de \sref{III.5.1.1}, seja $h:Z\to Y$ um
morfismo de tipo finito, seja $\widehat{Z}$ o completamento de $Z$
ao longo de $Z'=h^{-1}(Y')$, seja $g:Z\to X$ um $Y$-morfismo próprio,
e seja $\widehat{g}:\widehat{Z}\to\widehat{X}$ sua extensão aos
completamentos. Para todo $\sh{O}_{\widehat{Z}}$-módulo
algebrizável $\sh{F}'$, o módulo
$\widehat{g}_*(\sh{F}')$ é um
$\sh{O}_{\widehat{X}}$-módulo algebrizável.
\end{lemma}

\begin{proof}
Se $\sh{F}$ é um $\sh{O}_Z$-módulo coerente tal que
$\sh{F}'=\widehat{\sh{F}}$, o primeiro teorema de comparação
\sref{III.4.1.5} identifica
$\widehat{g}_*(\widehat{\sh{F}})$ com o completamento de
$g_*(\sh{F})$.
\end{proof}

\begin{lemma}[5.3.4]
\label{III.5.3.4}
Seja $X$ um esquema ordinário noetheriano, seja $X'$ um subconjunto
fechado de $X$, seja $f:Z\to X$ um morfismo próprio, e ponhamos
\[
  Z'=f^{-1}(X'),\qquad
  \widehat{X}=X_{/X'},\qquad
  \widehat{Z}=Z_{/Z'}.
\]
Seja $\widehat{f}:\widehat{Z}\to\widehat{X}$ a extensão de $f$
aos completamentos. Seja $\sh{M}$ um ideal coerente de
$\sh{O}_X$ tal que, se
\[
  U=X-\Supp(\sh{O}_X/\sh{M}),
\]
a restrição $f^{-1}(U)\to U$ de $f$ é um isomorfismo. Então,
para todo $\sh{O}_{\widehat{X}}$-módulo coerente $\sh{F}$, existe
um inteiro $n>0$ tal que o núcleo e o conúcleo do homomorfismo
canônico
\[
  \rho_{\sh{F}}:
  \sh{F}\longrightarrow
  \widehat{f}_*\bigl(\widehat{f}^{\,*}(\sh{F})\bigr)
\]
\sref[0]{0.4.4.3} são anulados por
$\widehat{\sh{M}}^{\,n}$.
\end{lemma}

\begin{proof}
Podemos restringir-nos ao caso $X=\Spec(B)$, onde $B$ é um
anel noetheriano; então $X'=\mathrm{V}(\mathfrak{K})$ para um ideal
$\mathfrak{K}$ de $B$. Reduzimo-nos primeiro ao caso em que $B$ é
noetheriano ádico e $\mathfrak{K}$ é um ideal de definição.

Seja $B_1$ o completamento separado de $B$ para a
topologia $\mathfrak{K}$-pré-ádica.
\oldpage[III]{155}
Se $\mathfrak{K}_1=\mathfrak{K}B_1$, então $B_1$ é um anel
noetheriano ádico e $\mathfrak{K}_1$ é um ideal de definição.
Ponhamos $X_1=\Spec(B_1)$, e seja $h:X_1\to X$ o morfismo
correspondente ao homomorfismo canônico $B\to B_1$. Se
$X'_1=h^{-1}(X')$, então $X'_1=\mathrm{V}(\mathfrak{K}_1)$. Finalmente,
ponhamos
\[
  Z_1=Z\times_X X_1,\qquad f_1:Z_1\longrightarrow X_1.
\]
O morfismo $f_1$ é próprio \sref[II]{II.5.4.2}. Denotemos por
$\widehat{X}_1$ o completamento de $X_1$ ao longo de $X'_1$, por
\[
  \widehat{Z}_1=
  Z_1\times_{X_1}\widehat{X}_1
\]
o completamento de $Z_1$ ao longo de
$Z'_1=f_1^{-1}(X'_1)$, e por
$\widehat{f}_1$ a extensão de $f_1$ aos completamentos.

A extensão
$\widehat{h}:\widehat{X}_1\to\widehat{X}$ corresponde à
aplicação identidade de $B_1$ e é, portanto, um isomorfismo
\sref[I]{I.10.9.1}. O morfismo correspondente
$\widehat{Z}_1\to\widehat{Z}$ é também, por conseguinte, um isomorfismo,
e estes isomorfismos identificam $\widehat{f}_1$ com
$\widehat{f}$. Além disso,
$\sh{M}_1=h^*(\sh{M})$ é um ideal coerente de $\sh{O}_{X_1}$ e
\[
  \Supp(\sh{O}_{X_1}/\sh{M}_1)
  =h^{-1}\bigl(\Supp(\sh{O}_X/\sh{M})\bigr)
\]
\sref[I]{I.9.1.13}. Assim, se
$U_1=X_1-\Supp(\sh{O}_{X_1}/\sh{M}_1)$, então
$U_1=h^{-1}(U)$, e a restrição
$f_1^{-1}(U_1)\to U_1$ é um isomorfismo
\sref[I]{I.3.2.7}. Finalmente, $\widehat{\sh{M}}$ e
$\widehat{\sh{M}}_1$ são identificados por $\widehat{h}$
\sref[I]{I.10.9.5}. Todas as hipóteses do lema são, pois,
satisfeitas por $X_1$, $X'_1$, $f_1$ e $\sh{M}_1$; doravante
podemos supor que $B$ é noetheriano ádico e que $\mathfrak{K}$ é
um ideal de definição.

Tem-se então $\widehat{X}=\Spf(B)$ e
$\sh{F}=N^\Delta$, onde $N$ é um $B$-módulo de tipo finito. Assim,
$\sh{F}=\widehat{\sh{G}}$, onde $\sh{G}$ é o
$\sh{O}_X$-módulo coerente $\widetilde{N}$ \sref[I]{I.10.10.5}, e
\[
  \widehat{f}^{\,*}(\sh{F})
  =\widehat{f^*(\sh{G})}
\]
\sref[I]{I.10.9.5}. Pelo primeiro teorema de comparação
\sref{III.4.1.5},
\[
  \widehat{f}_*
  \bigl(\widehat{f^*(\sh{G})}\bigr)
  \simeq
  \widehat{f_*\bigl(f^*(\sh{G})\bigr)},
\]
e, por \sref{III.5.1.3}, o homomorfismo canônico
$\rho_{\sh{F}}$ é precisamente $\widehat{\rho}_{\sh{G}}$.

O núcleo $\sh{P}$ e o conúcleo $\sh{R}$ de
\[
  \rho_{\sh{G}}:
  \sh{G}\longrightarrow f_*\bigl(f^*(\sh{G})\bigr)
\]
são $\sh{O}_X$-módulos coerentes, e suas restrições a $U$ são
nulas. Existe, portanto, um inteiro $n>0$ tal que
\[
  \sh{M}^n\sh{P}=\sh{M}^n\sh{R}=0
\]
\sref[I]{I.9.3.4}. Deduz-se que
\[
  \widehat{\sh{M}}^{\,n}\widehat{\sh{P}}
  =
  \widehat{\sh{M}}^{\,n}\widehat{\sh{R}}
  =0
\]
\sref[I]{I.10.8.8} \sref[I]{I.10.8.10}, o que demonstra o lema.
\end{proof}

\begin{env}[5.3.5]
\label{III.5.3.5}
\emph{Fim da demonstração do teorema de existência.}
Conservemos as hipóteses de \sref{III.5.1.4}. Utilizemos a indução
noetheriana \sref[0]{0.2.2.2}, supondo conhecido o teorema para todo
subpré-esquema fechado $T$ de $X$ cujo espaço subjacente é distinto
de $X$; aqui $\widehat{T}$ denota o completamento de $T$ ao longo de
$T\cap X'$. Podemos supor que $X$ não é vazio.

Como $f$ é separado e de tipo finito, o lema de Chow
\sref[II]{II.5.6.1} fornece um $Y$-esquema $Z$ e um $Y$-morfismo
$g:Z\to X$ tais que o morfismo estrutural $h:Z\to Y$ é
quase-projetivo, $g$ é projetivo e sobrejetivo, e existe um
subconjunto aberto não vazio $U$ de $X$ no qual
$g^{-1}(U)\to U$ é um isomorfismo. Seja $\sh{M}$ um ideal coerente
de $\sh{O}_X$ que define um subpré-esquema fechado de espaço subjacente
$X-U$ \sref[I]{I.5.2.2}, e seja $\sh{F}$ um
$\sh{O}_{\widehat{X}}$-módulo coerente cujo suporte $E$ é próprio sobre $Y$.
Denotemos por $\widehat{Z}$ o completamento de $Z$ ao longo de
$h^{-1}(Y')$, e por
$\widehat{g}:\widehat{Z}\to\widehat{X}$ a extensão de $g$ aos
completamentos.

O módulo $\widehat{g}^{\,*}(\sh{F})$ é um
$\sh{O}_{\widehat{Z}}$-módulo coerente cujo suporte está contido em
$g^{-1}(E)$ e é, portanto, próprio sobre $Y$, pois $g$ é
projetivo e, portanto, próprio \sref[II]{II.5.4.6}. Como $h$ é
quase-projetivo, $\widehat{g}^{\,*}(\sh{F})$ é algebrizável por
\sref{III.5.2.6}.
\oldpage[III]{156}
Deduz-se de \sref{III.5.3.3}, dado que $g$ é próprio, que
\[
  \widehat{g}_*
  \bigl(\widehat{g}^{\,*}(\sh{F})\bigr)
\]
é um $\sh{O}_{\widehat{X}}$-módulo algebrizável.

Apliquemos \sref{III.5.3.4} a $\sh{F}$ e $g$. O núcleo $\sh{P}$
e o conúcleo $\sh{R}$ de
\[
  \rho_{\sh{F}}:
  \sh{F}\longrightarrow
  \widehat{g}_*\bigl(\widehat{g}^{\,*}(\sh{F})\bigr)
\]
são anulados por uma potência $\widehat{\sh{M}}^{\,n}$. Seja $T$
o subpré-esquema fechado de $X$ definido por $\sh{M}^n$, cujo
espaço subjacente é $X-U$, e seja $j:T\to X$ a imersão
canônica. Sua extensão
$\widehat{j}:\widehat{T}\to\widehat{X}$ é a imersão canônica
\sref[I]{I.10.14.7}, e podemos escrever
\[
  \sh{P}=
  \widehat{j}_*\bigl(\widehat{j}^{\,*}(\sh{P})\bigr),
  \qquad
  \sh{R}=
  \widehat{j}_*\bigl(\widehat{j}^{\,*}(\sh{R})\bigr).
\]
Como $U$ não é vazio, a hipótese de indução mostra que
$\widehat{j}^{\,*}(\sh{P})$ e
$\widehat{j}^{\,*}(\sh{R})$ são $\sh{O}_{\widehat{T}}$-módulos
algebrizáveis. O Lema \sref{III.5.3.3} mostra então
que $\sh{P}$ e $\sh{R}$ são algebrizáveis, e
\sref{III.5.3.2} demonstra finalmente que $\sh{F}$ é algebrizável.
\end{env}


\subsection{Aplicação: comparação entre morfismos de pré-esquemas ordinários e morfismos de pré-esquemas formais. Pré-esquemas formais algebrizáveis}
\label{subsection:III.5.4}

\begin{theorem}[5.4.1]
\label{III.5.4.1}
Seja $A$ um anel noetheriano ádico, seja $\mathfrak{J}$ um ideal
de definição de $A$, e ponhamos
\[
  S=\Spec(A),\qquad S'=\mathrm{V}(\mathfrak{J}).
\]
Seja $u:X\to S$ um morfismo próprio e seja $v:Y\to S$ um
morfismo separado de tipo finito. Sejam
$\widehat{S}$, $\widehat{X}$ e $\widehat{Y}$ os completamentos
de $S$, $X$ e $Y$ ao longo de $S'$, $u^{-1}(S')$ e $v^{-1}(S')$,
respectivamente. Para todo $S$-morfismo $f:X\to Y$, seja
$\widehat{f}:\widehat{X}\to\widehat{Y}$ sua extensão às
completamentos. Então a aplicação $f\mapsto\widehat{f}$ é uma bijeção
\[
  \Hom_S(X,Y)
  \xrightarrow{\ \sim\ }
  \Hom_{\widehat{S}}(\widehat{X},\widehat{Y}).
\]
\end{theorem}

\begin{proof}
Demonstremos primeiro a injetividade. Suponhamos que $f,g:X\to Y$ são
$S$-morfismos tais que $\widehat{f}=\widehat{g}$. Existe então
uma vizinhança aberta $V$ de $X'=u^{-1}(S')$ sobre a qual $f$ e $g$
coincidem \sref[I]{I.10.9.4}. Como $u$ é fechado,
\sref{III.5.1.3.1} dá $V=X$, e portanto $f=g$.

Demonstremos agora a sobrejetividade. Seja
$h:\widehat{X}\to\widehat{Y}$ um
$\widehat{S}$-morfismo. Ponhamos
\[
  Z=X\times_S Y,
\]
e denotemos as projeções canônicas por $p:Z\to X$ e
$q:Z\to Y$. O pré-esquema $Z$ é de tipo finito sobre $S$
\sref[I]{I.6.3.4} e, portanto, noetheriano. Seja
$\widehat{Z}$ seu completamento ao longo de
\[
  Z'=p^{-1}\bigl(u^{-1}(S')\bigr).
\]
O pré-esquema formal $\widehat{Z}$ identifica-se canonicamente com
$\widehat{X}\times_{\widehat{S}}\widehat{Y}$, e suas projeções
sobre $\widehat{X}$ e $\widehat{Y}$ identificam-se com
$\widehat{p}$ e $\widehat{q}$ \sref[I]{I.10.9.7}.

Como $Y$ é separado sobre $S$, $\widehat{Y}$ é separado sobre
$\widehat{S}$ \sref[I]{I.10.15.7}. O morfismo grafo
\[
  \Gamma_h=(1_{\widehat{X}},h):
  \widehat{X}\longrightarrow\widehat{Z}
\]
é, por conseguinte, uma imersão fechada \sref[I]{I.10.15.4}. Seja
$\mathfrak{T}$ o subpré-esquema formal fechado de
$\widehat{Z}$ associado a esta imersão, e seja
$j:\mathfrak{T}\to\widehat{Z}$ a imersão canônica. Assim,
\[
  \Gamma_h=j\circ w,
\]
onde $w:\widehat{X}\to\mathfrak{T}$ é um isomorfismo
\sref[I]{I.10.14.3} cujo inverso é
$\widehat{p}\circ j$. Além disso, $\mathfrak{T}$ é próprio sobre
$\widehat{S}$ porque $\widehat{X}$ o é. Deduz-se de
\sref{III.5.1.8} que existe um subpré-esquema fechado $T$ de $Z$
tal que
\[
  \mathfrak{T}
  =\widehat{T}
  =T_{/(T\cap Z')},
\]
e que $j=\widehat{i}$, onde $i:T\to Z$ é a imersão
canônica \sref[I]{I.10.14.7}.

O morfismo $p\circ i:T\to X$ é um isomorfismo: sua extensão aos
completamentos é o isomorfismo
$\widehat{p}\circ\widehat{i}$, e basta aplicar
\sref{III.4.6.8}, observando como antes que $S$ é a única
vizinhança de $S'$ em $S$.
\oldpage[III]{157}
Seja $g:X\to T$ o inverso de $p\circ i$, e ponhamos
\[
  f=q\circ i\circ g.
\]
Então $f:X\to Y$ tem por grafo $\Gamma_f=i\circ g$. Como
$\widehat{g}$ é o inverso de
$\widehat{p\circ i}$, tem-se
\[
  \widehat{\Gamma_f}
  =\widehat{i}\circ\widehat{g}
  =j\circ w
  =\Gamma_h.
\]
Mas $\widehat{\Gamma_f}=\Gamma_{\widehat{f}}$
\sref[I]{I.10.9.8}; portanto, $h=\widehat{f}$, o que termina a
demonstração.
\end{proof}

Em linguagem categórica, o funtor
$X\mapsto\widehat{X}$ é, portanto, plenamente fiel
\sref[0]{0.8.1.6} da categoria dos pré-esquemas próprios sobre
$\Spec(A)$ para a categoria dos pré-esquemas formais próprios sobre
$\Spf(A)$, para todo anel noetheriano ádico $A$. Fornece, por
conseguinte, uma equivalência entre a primeira categoria e uma
subcategoria plena da segunda. Os objetos desta subcategoria se
denominarão doravante \emph{pré-esquemas formais algebrizáveis}.
Para tal pré-esquema formal $\mathfrak{X}$, existe um pré-esquema
ordinário $X$, próprio sobre $\Spec(A)$ e determinado salvo isomorfismo
único, tal que $\mathfrak{X}\simeq\widehat{X}$.

\begin{env}[5.4.2]
\label{III.5.4.2}
\emph{Escólio.}
Com a notação de \sref{III.5.4.1}, ponhamos
\[
  S_n=\Spec(A/\mathfrak{J}^{n+1}),\qquad
  X_n=X\times_S S_n,\qquad
  Y_n=Y\times_S S_n.
\]
Deduz-se de \sref{III.5.4.1} e \sref[I]{I.10.12.3} que
dar um $S$-morfismo $f:X\to Y$ equivale a dar um sistema indutivo
ádico $(S_n)$ \sref[I]{I.10.12.2} de
$S_n$-morfismos $f_n:X_n\to Y_n$.
\end{env}

\begin{remark}[5.4.3]
\label{III.5.4.3}
Contrariamente ao que poderia sugerir o teorema de existência
\sref{III.5.1.6}, existem pré-esquemas formais próprios sobre $\Spf(A)$
que \emph{não são algebrizáveis}, assim como existem espaços analíticos
compactos que não procedem de variedades algébricas complexas. Mais
adiante encontraremos tais pré-esquemas na «teoria de moduli»,
que, quando o corpo base é $\mathbb{C}$, trata precisamente das
variações infinitesimais da estrutura complexa de uma variedade
algébrica completa; é bem sabido que tais variações podem dar
lugar a variedades analíticas que não são algébricas.
\end{remark}

\begin{proposition}[5.4.4]
\label{III.5.4.4}
Seja $A$ um anel noetheriano ádico, ponhamos
$\mathfrak{S}=\Spf(A)$, e sejam
\[
  g:\mathfrak{X}\to\mathfrak{S},
  \qquad
  h:\mathfrak{Y}\to\mathfrak{S}
\]
dois morfismos próprios de pré-esquemas formais. Seja
$f:\mathfrak{X}\to\mathfrak{Y}$ um
$\mathfrak{S}$-morfismo. Se $f$ é finito e
$\mathfrak{Y}$ é algebrizável, então $\mathfrak{X}$ é
algebrizável.
\end{proposition}

\begin{proof}
As hipóteses sobre $g$ e $h$ já implicam que $f$ é próprio
\sref{III.3.4.1}; para que $f$ seja finito, basta que a fibra
$f^{-1}(y)$ seja finita para todo $y\in\mathfrak{Y}$
\sref[0]{0.4.8.11}. Por hipótese,
\[
  \sh{B}=f_*(\sh{O}_{\mathfrak{X}})
\]
é uma $\sh{O}_{\mathfrak{Y}}$-álgebra coerente
\sref[0]{0.4.8.6}. Escrevamos
$\mathfrak{Y}=\widehat{Y}$ e $h=\widehat{w}$, onde
$w:Y\to\Spec(A)$ é um morfismo próprio de pré-esquemas ordinários. O
teorema de existência fornece uma $\sh{O}_Y$-álgebra coerente $\sh{C}$
tal que $\sh{B}=\widehat{\sh{C}}$. Ponhamos
\[
  X=\Spec(\sh{C}),
\]
e seja $u:X\to Y$ o morfismo estrutural. A definição de
$\mathfrak{X}$ a partir de $\sh{B}$ \sref[0]{0.4.8.7} mostra então que
$\mathfrak{X}$ é canonicamente isomorfo a $\widehat{X}$ e
que $f$ se identifica com $\widehat{u}$; basta verificá-lo
quando $Y$ é afim.
\end{proof}

A Proposição \sref{III.5.1.8} é um caso particular de
\sref{III.5.4.4}.

\begin{theorem}[5.4.5]
\label{III.5.4.5}
Seja $A$ um anel noetheriano ádico, seja $\mathfrak{J}$ um ideal
de definição de $A$, ponhamos
\[
  S=\Spec(A),\qquad
  \mathfrak{S}=\widehat{S}=\Spf(A),
\]
e seja $f:\mathfrak{X}\to\mathfrak{S}$ um morfismo próprio de
pré-esquemas formais. Ponhamos
\[
  S_k=\Spec(A/\mathfrak{J}^{k+1}),\qquad
  X_k=\mathfrak{X}\times_{\mathfrak{S}}S_k,
\]
e, para todo $\sh{O}_{\mathfrak{X}}$-módulo $\sh{F}$, ponhamos
\[
  \sh{F}_k
  =\sh{F}\otimes_{\sh{O}_{\mathfrak{X}}}\sh{O}_{X_k}
  =\sh{F}/\mathfrak{J}^{k+1}\sh{F}.
\]
\oldpage[III]{158}
Seja $\sh{L}$ um
$\sh{O}_{\mathfrak{X}}$-módulo invertível, e suponhamos que
\[
  \sh{L}_0=\sh{L}/\mathfrak{J}\sh{L}
\]
é um $\sh{O}_{X_0}$-módulo amplo. Então $\mathfrak{X}$ é
algebrizável. Além disso, se $X$ é um $S$-pré-esquema próprio tal
que $\mathfrak{X}\simeq\widehat{X}$, existe um
$\sh{O}_X$-módulo amplo $\sh{M}$ tal que
$\sh{L}\simeq\widehat{\sh{M}}$; em particular, $X$ é projetivo
sobre $S$.
\end{theorem}

\begin{proof}
Apliquemos \sref{III.5.2.3} com
$\sh{F}=\sh{O}_{\mathfrak{X}}$. Existe um inteiro $n_0$ tal
que, para $n\geq n_0$, o homomorfismo canônico
\[
  \Gamma(\mathfrak{X},\sh{L}^{\otimes n})
  \longrightarrow
  \Gamma(X_0,\sh{L}_0^{\otimes n})
\]
é sobrejetivo. Escolhamos $n\geq n_0$ suficientemente grande para que
$\sh{L}_0^{\otimes n}$ seja muito amplo para $S_0$
\sref[II]{II.4.5.10}. Como
$f_0:X_0\to S_0$ é próprio,
$\Gamma(X_0,\sh{L}_0^{\otimes n})$ é um $A$-módulo de tipo finito
\sref{III.3.2.1}. Existe, por conseguinte, um
$A$-submódulo de tipo finito
\[
  E\subset\Gamma(\mathfrak{X},\sh{L}^{\otimes n})
\]
cuja imagem em $\Gamma(X_0,\sh{L}_0^{\otimes n})$ é todo o
módulo.

Para todo $k\geq0$, consideremos o homomorfismo
\[
  u_k:
  E/\mathfrak{J}^{k+1}E
  \longrightarrow
  \Gamma(X_k,\sh{L}_k^{\otimes n})
\]
induzido pela injeção canônica
$E\to\Gamma(\mathfrak{X},\sh{L}^{\otimes n})$. O módulo
$(f_k)_*(\sh{L}_k^{\otimes n})$ é quase-coerente, e
\[
  \Gamma\bigl(S_k,(f_k)_*(\sh{L}_k^{\otimes n})\bigr)
  =
  \Gamma(X_k,\sh{L}_k^{\otimes n}).
\]
Assim, $u_k$ define um homomorfismo
\[
  \widetilde{u}_k:
  \widetilde{E/\mathfrak{J}^{k+1}E}
  \longrightarrow
  (f_k)_*(\sh{L}_k^{\otimes n}),
\]
e, portanto, por adjunção, um homomorfismo
\[
  \widetilde{u}_k^{\,\sharp}:
  f_k^*\bigl(\widetilde{E/\mathfrak{J}^{k+1}E}\bigr)
  \longrightarrow
  \sh{L}_k^{\otimes n}.
\]
Ponhamos
\[
  \sh{G}_k=
  f_k^*\bigl(\widetilde{E/\mathfrak{J}^{k+1}E}\bigr).
\]
Tem-se
$\sh{G}_0=\sh{G}_k/\mathfrak{J}\sh{G}_k$
\sref[I]{I.9.1.5}, e
\[
  \widetilde{u}_0^{\,\sharp}:
  \sh{G}_0/\mathfrak{J}\sh{G}_0
  \longrightarrow
  \sh{L}_0^{\otimes n}/
  \mathfrak{J}\sh{L}_0^{\otimes n}
\]
obtém-se de $\widetilde{u}_k^{\,\sharp}$ passando aos
quocientes.

Pela definição de $E$, $\widetilde{u}_0^{\,\sharp}$ é o
homomorfismo canônico
\[
  \sigma:
  f_0^*\bigl((f_0)_*(\sh{L}_0^{\otimes n})\bigr)
  \longrightarrow
  \sh{L}_0^{\otimes n}.
\]
Como $\sh{L}_0^{\otimes n}$ é muito amplo,
$\widetilde{u}_0^{\,\sharp}$ é sobrejetivo
\sref[II]{II.4.4.3}. O lema de Nakayama mostra então que todo
$\widetilde{u}_k^{\,\sharp}$ é sobrejetivo.

Cada $\widetilde{u}_k^{\,\sharp}$ define, portanto,
\sref[II]{II.4.2.2} um $S_k$-morfismo
\[
  g_k:X_k\longrightarrow
  P_k=\mathbf{P}(E/\mathfrak{J}^{k+1}E).
\]
Para $h\leq k$,
\[
  P_h=P_k\times_{S_k}S_h
\]
por \sref[II]{II.4.1.3}. As relações entre os
$\widetilde{u}_k^{\,\sharp}$ e \sref[II]{II.4.2.10} mostram que
$(g_k)$ é um sistema indutivo ádico $(S_k)$
\sref[I]{I.10.12.2}. Os $g_k$ definem, por conseguinte, um
$\mathfrak{S}$-morfismo de pré-esquemas formais
\[
  g:\mathfrak{X}\longrightarrow\mathfrak{P},
\]
onde $\mathfrak{P}$ é o limite indutivo do sistema $(P_k)$,
ou, equivalentemente, o completamento $\widehat{P}$ de
$P=\mathbf{P}(E)$. O fato de que
$\sh{L}_0^{\otimes n}$ seja muito amplo implica que $g_0$ é uma imersão fechada
\sref[II]{II.4.4.3}. Portanto, $g$ é uma imersão fechada de
pré-esquemas formais \sref[0]{0.4.8.10}, e $\mathfrak{X}$ é algebrizável
por \sref{III.5.1.8}.

O teorema de existência \sref{III.5.1.6} mostra agora que $\sh{L}$ é
o completamento $\widehat{\sh{M}}$ de um
$\sh{O}_X$-módulo invertível $\sh{M}$. Além disso,
$\sh{L}^{\otimes n}$ é o completamento de
$\sh{M}^{\otimes n}$ \sref[I]{I.10.8.10}, e os homomorfismos
$\widetilde{u}_k^{\,\sharp}$ definem um único homomorfismo
\[
  v:f^*(\widetilde{E})
  \longrightarrow
  \sh{M}^{\otimes n}
\]
\sref{III.5.1.7}. Como os
$\widetilde{u}_k^{\,\sharp}$ são sobrejetivos,
$\widehat{v}$ é sobrejetivo \sref[I]{I.10.11.5}, e portanto também
o é $v$ \sref{III.5.1.3}. Seja $r:X\to P$ o morfismo definido
por $v$ \sref[II]{II.4.2.2}. Sua extensão aos completamentos é
$g$. Como $g$ é uma imersão fechada, também o é $r$, por
\sref{III.5.1.8} e \sref{III.5.4.1}. Deduz-se que
$\sh{M}^{\otimes n}$ é muito amplo \sref[II]{II.4.4.6} e, portanto,
que $\sh{M}$ é amplo \sref[II]{II.4.5.10}.
\end{proof}

\begin{remark}[5.4.6]
\label{III.5.4.6}
Seja $A$ um anel noetheriano ádico, ponhamos
$\mathfrak{S}=\Spf(A)$, e seja
$f:\mathfrak{X}\to\mathfrak{S}$ um morfismo próprio de pré-esquemas
formais. Seja $\sh{N}$ o ideal coerente de
$\sh{O}_{\mathfrak{X}}$ tal que, para todo subconjunto aberto
formal afim $U$ de $\mathfrak{X}$,
\[
  \Gamma(U,\sh{N})
\]
é o nilradical de
$\Gamma(U,\sh{O}_{\mathfrak{X}})$. A existência deste ideal
deduz-se facilmente de \sref[I]{I.10.10.2} e do fato de que
todo homomorfismo de anéis $B\to C$ envia o nilradical de $B$ em
o de $C$. Seja $\mathfrak{X}'$ o subpré-esquema formal fechado de
$\mathfrak{X}$ definido por $\sh{N}$
\sref[I]{I.10.14.2}. Seria interessante saber se a
algebraizabilidade de $\mathfrak{X}'$ implica a algebraizabilidade do
próprio $\mathfrak{X}$.

Sem dúvida, chegar-se-ia a uma solução deste problema se se soubesse
classificar, por exemplo mediante invariantes
\oldpage[III]{159}
de natureza cohomológica, as \emph{extensões} de um feixe estrutural
$\sh{O}_{\mathfrak{X}}$---para um pré-esquema ordinário ou
formal---por um ideal de quadrado nulo. Equivalentemente, estas são as
$\sh{O}_{\mathfrak{X}}$-álgebras $\sh{A}$ tais que
$\sh{O}_{\mathfrak{X}}\simeq\sh{A}/\sh{J}$, onde $\sh{J}$ é um
ideal de quadrado nulo de $\sh{A}$.
\end{remark}


\subsection{Uma decomposição de certos pré-esquemas}
\label{subsection:III.5.5}

\begin{proposition}[5.5.1]
\label{III.5.5.1}
Seja $A$ um anel noetheriano ádico, seja $\mathfrak{J}$ um ideal
de definição de $A$, e ponhamos $Y=\Spec(A)$. Seja $f:X\to Y$ um
morfismo separado de tipo finito, e ponhamos
\[
  Y_0=\Spec(A/\mathfrak{J}),\qquad
  X_0=X\times_Y Y_0=f^{-1}(Y_0).
\]
Seja $Z_0$ um subconjunto aberto de $X_0$ que é próprio sobre $Y_0$.
Então existe um subconjunto aberto e fechado $Z$ de $X$, próprio sobre
$Y$, tal que
\[
  Z\cap X_0=Z_0.
\]
\end{proposition}

\begin{proof}
Por hipótese, existe um subconjunto aberto $T$ de $X$ tal que
$T\cap X_0=Z_0$. Seja $\widehat{T}$ o completamento ao longo de $Z_0$
do pré-esquema induzido por $X$ sobre $T$. O suporte de
$\sh{O}_{\widehat{T}}$ é $Z_0$, que é próprio sobre $Y_0$;
portanto, $\widehat{T}$ é próprio sobre
$\widehat{Y}=\Spf(A)$ \sref{III.3.4.1}.

Deduz-se de \sref{III.5.1.8} que existe um subpré-esquema fechado
$Z$ de $T$, próprio sobre $Y$, tal que, se
$i:Z\to T$ é a imersão canônica, então
\[
  \widehat{i}:\widehat{Z}\longrightarrow\widehat{T}
\]
é um isomorfismo; aqui $\widehat{Z}$ é o completamento de $Z$
ao longo de $Z_0$. Por \sref{III.4.6.8}, existe uma vizinhança aberta
$V$ de $Z_0$ em $T$ tal que
$i^{-1}(V)\to V$ é um isomorfismo. Mas $i^{-1}(V)$ é uma
vizinhança de $Z_0$ em $Z$ e, portanto, é igual a $Z$ por
\sref{III.5.1.3.1}. Assim, $Z$ é aberto em $T$, e portanto em $X$.
Também é fechado em $X$, pois $Z$ é próprio sobre $Y$ e $X$ é
separado sobre $Y$. Isto demonstra a proposição.
\end{proof}

\begin{corollary}[5.5.2]
\label{III.5.5.2}
Se $X_0$ é próprio sobre $Y_0$, então $X$ é a união de dois
subconjuntos abertos disjuntos $Z$ e $Z'$ tais que $Z$ é próprio sobre
$Y$ e contém $X_0$. Além disso, todo subconjunto fechado $P$ de $X$
que é próprio sobre $Y$ está contido em $Z$.
\end{corollary}

\begin{proof}
Apliquemos \sref{III.5.5.1} com $Z_0=X_0$, e ponhamos $Z'=X-Z$. Para a
última afirmação, $P\cap Z'$ é fechado em $P$ e, portanto, próprio
sobre $Y$. Se não fosse vazio, $f(P\cap Z')$ seria um subconjunto fechado
não vazio de $Y$ e, portanto, intersectaria $Y_0$
\sref{III.5.1.3.1}. Isto contradiz
$Z'\cap X_0=\varnothing$ e demonstra que $P\subset Z$.
\end{proof}

\bigskip
\noindent\emph{(Continuará.)}

% La página impresa 504 está en blanco.
\clearpage
\thispagestyle{empty}
\null
\clearpage
\thispagestyle{plain}

\oldpage[III]{161}
\section*{BIBLIOGRAFIA (continuação)}

\begin{flushleft}
[27] P. Gabriel, \emph{Des catégories abéliennes}, tese, Paris (1961).

\medskip
[28] J. E. Roos, \emph{Sur les foncteurs dérivés de $\varprojlim$.
Applications}, C. R. Acad. Sci. Paris, vol. CCLII, 1961, pp. 3702--3704.

\medskip
[29] H. Cartan, \emph{Séminaire de l'École Normale Supérieure},
13.er anho (1960--1961), exposição n.º 11.
\end{flushleft}

\clearpage
\oldpage[III]{162}
\section*{ÍNDICE De Notações}

\begin{flushleft}
$\Set$: 0, 8.1.1.

$h_X(Y)$, $h_X(u)$ ($X$, $Y$ objetos de uma categoria $\C$, $u$ um morfismo
de $\C$): 0, 8.1.1.

$h_w(Y)$ ($Y$ um objeto e $w$ um morfismo de $\C$): 0, 8.1.2.

$Z_k(E_r^{pq})$, $B_k(E_r^{pq})$ ($k\geq r$): 0, 11.1.1.

$Z_\infty(E_2^{pq})$, $B_\infty(E_2^{pq})$: 0, 11.1.1.

$K^\bullet$, $K_\bullet$ (complexos): 0, 11.2.1.

$K^{\bullet\bullet}$, $K_{\bullet\bullet}$ (bicomplexos): 0, 11.2.1.

$B^i(K^\bullet)$, $Z^i(K^\bullet)$, $H^i(K^\bullet)$,
$H^\bullet(K^\bullet)$: 0, 11.2.1.

$B_i(K_\bullet)$, $Z_i(K_\bullet)$, $H_i(K_\bullet)$,
$H_\bullet(K_\bullet)$: 0, 11.2.1.

$T(K^\bullet)$, $T(K_\bullet)$ ($T$ um funtor): 0, 11.2.1.

$E(K^\bullet)$ ($K^\bullet$ um complexo filtrado): 0, 11.2.2.

$K^{i,\bullet}$, $K^{\bullet,j}$, $Z_{\mathrm{II}}^p(K^{i,\bullet})$,
$B_{\mathrm{II}}^p(K^{i,\bullet})$, $H_{\mathrm{II}}^p(K^{i,\bullet})$,
$Z_{\mathrm I}^p(K^{\bullet,j})$, $B_{\mathrm I}^p(K^{\bullet,j})$,
$H_{\mathrm I}^p(K^{\bullet,j})$: 0, 11.3.1.

$H_{\mathrm{II}}^p(K^{\bullet\bullet})$,
$Z_{\mathrm I}^q(H_{\mathrm{II}}^p(K^{\bullet\bullet}))$,
$B_{\mathrm I}^q(H_{\mathrm{II}}^p(K^{\bullet\bullet}))$,
$H_{\mathrm I}^q(H_{\mathrm{II}}^p(K^{\bullet\bullet}))$,
$H^n(K^{\bullet\bullet})$: 0, 11.3.1.

$F_{\mathrm I}^p(K^{\bullet\bullet})$,
$F_{\mathrm{II}}^p(K^{\bullet\bullet})$, ${}'E(K^{\bullet\bullet})$,
${}''E(K^{\bullet\bullet})$: 0, 11.3.2.

$K_{i,\bullet}$, $K_{\bullet,j}$, $H_p^{\mathrm{II}}(K_{i,\bullet})$,
$H_p^{\mathrm I}(K_{\bullet,j})$, $H_p^{\mathrm{II}}(K_{\bullet\bullet})$,
$H_p^{\mathrm I}(K_{\bullet\bullet})$,
$H_q^{\mathrm I}(H_p^{\mathrm{II}}(K_{\bullet\bullet}))$,
$H_q^{\mathrm{II}}(H_p^{\mathrm I}(K_{\bullet\bullet}))$,
$H_n(K_{\bullet\bullet})$: 0, 11.3.5.

$\RR^\bullet T(K^\bullet)$ ($T$ um funtor covariante aditivo): 0, 11.4.3.

$\RR^\bullet T(K^\bullet,K'^{\bullet})$ ($T$ um bifuntor covariante
aditivo): 0, 11.4.6.

$\LL_\bullet T(K_\bullet)$ ($T$ um funtor covariante aditivo): 0, 11.6.2.

$\LL_\bullet T(K_\bullet,K'_\bullet)$ ($T$ um bifuntor covariante
aditivo): 0, 11.6.6.

$\LL_\bullet T(K_{\bullet\bullet})$ ($T$ um funtor covariante aditivo):
0, 11.7.2.

$|\sigma|$, $\Sigma(A)$, $\mathrm C_\bullet(A)$, $\mathrm L_\bullet(A)$
($A$ um conjunto finito, $\sigma$ um simplexo de $A$): 0, 11.8.1.

$\mathrm C^\bullet(A,B;\sh{S})$, $\mathrm L^\bullet(A,B;\sh{S})$
($A$, $B$ conjuntos finitos, $\sh{S}$ um sistema de coeficientes): 0, 11.8.4.

$\mathrm P^\bullet(A,B;\sh{S})$ ($A$, $B$ conjuntos finitos, $\sh{S}$ um
sistema de coeficientes): 0, 11.8.6.

$\mathrm C^\bullet(A,B;\sh{S}^\bullet)$,
$\mathrm L^\bullet(A,B;\sh{S}^\bullet)$ ($A$, $B$ conjuntos finitos,
$\sh{S}^\bullet$ um complexo de sistemas de coeficientes): 0, 11.8.8.

$\mathrm P^\bullet(A,B;\sh{S}^\bullet)$ ($A$, $B$ conjuntos finitos,
$\sh{S}^\bullet$ um complexo de sistemas de coeficientes): 0, 11.8.10.

$\chi(M^\bullet)$ ($M^\bullet$ um complexo finito de módulos de comprimento finito):
0, 11.10.1.

$\check C^\bullet(\mathfrak U,\sh{F})$ ($\sh{F}$ um feixe de grupos abelianos
sobre $X$, $\mathfrak U$ um recobrimento aberto de $X$): 0, 12.1.2.

$\RR^p f_*(\sh{F})$ ($f:X\to Y$ um morfismo de espaços anelados,
$\sh{F}$ um feixe de $\sh{O}_X$-módulos): 0, 12.2.1.

$\sh{H}^p(f,\sh{K}^\bullet)$,
$\sh{H}^p_f(\sh{K}^\bullet)$, ${}'E(f,\sh{K}^\bullet)$,
${}''E(f,\sh{K}^\bullet)$ ($f:X\to Y$ um morfismo de espaços anelados,
$\sh{K}^\bullet$ um complexo de $\sh{O}_X$-módulos): 0, 12.4.1.

$\mathbf H^\bullet(X,\sh{K}^\bullet)$,
$\mathbf H^\bullet(V,\sh{K}^\bullet)$ ($X$ um espaço anelado,
$\sh{K}^\bullet$ um complexo de $\sh{O}_X$-módulos, $V$ aberto em $X$):
0, 12.4.2.

$\operatorname{gr}^\bullet(A)$ ($A$ um sistema projetivo que satisfaz (ML)):
0, 13.4.2.

$E(A)$ ($A$ um objeto filtrado): 0, 13.6.4.

$K_\bullet(\mathbf f)$, $i_{\mathbf f}$ ($\mathbf f$ uma família finita de
elementos de um anel): III, 1.1.1.

$K^\bullet(\mathbf f,M)$, $K_\bullet(\mathbf f,M)$ ($\mathbf f$ uma família
finita de elementos de um anel $A$, $M$ um $A$-módulo): III, 1.1.2.

$H^\bullet(\mathbf f,M)$, $H_\bullet(\mathbf f,M)$ ($\mathbf f$ uma família
finita de elementos de um anel $A$, $M$ um $A$-módulo): III, 1.1.3.

$C^\bullet((\mathbf f),M)$, $H^\bullet((\mathbf f),M)$ ($\mathbf f$ uma
família finita de elementos de um anel $A$, $M$ um $A$-módulo): III, 1.1.6.

$H^\bullet(U,\sh{F}^{(*)})$,
$H^\bullet(\mathfrak U,\sh{F}^{(*)})$ ($\sh{F}$ um
$\sh{O}_X$-módulo quase-coerente, $U$ aberto em $X$, $\mathfrak U$ um recobrimento aberto de $X$):
III, 2.1.1.
\end{flushleft}

\clearpage
\oldpage[III]{163}
\section*{ÍNDICE TERMINOLÓGICO}

\begin{flushleft}
Termo de convergência de uma sucessão espectral: 0, 11.1.1.

Algebrizável ($\sh{O}_X$-módulo): III, 5.2.1.

Algebrizável (esquema formal): III, 5.4.2.

Analiticamente íntegro (anel): III, 4.3.6.

Aplicação quase-compacta: 0, 9.1.1.

Aumento de uma resolução: 0, 11.4.1.

Característica de Euler--Poincaré de um complexo: 0, 11.10.1.

Característica de Euler--Poincaré de um $\sh{O}_X$-módulo coerente
($X$ projetivo sobre $\Spec(A)$, $A$ um anel artiniano): III, 2.5.1.

Cocadeia bialternada: 0, 11.8.4.

Complexo definido por um bicomplexo: 0, 11.3.1.

Complexo de álgebra exterior: III, 1.1.1.

Condição (ML): 0, 13.1.2.

Construtível (subconjunto, conjunto): 0, 9.1.2.

Construtível (função): 0, 9.3.1.

Produto cup: 0, 12.1.2.

Dihomomorfismo para estruturas algébricas sobre categorias: 0, 8.2.1.

Exato (subconjunto) em uma categoria abeliana: III, 3.1.1.

Essencialmente constante (sistema projetivo): 0, 13.4.2.

Fatoração de Stein: III, 4.3.3.

Codiscreta, cosseparada, discreta, exaustiva, finita, separada
(filtração): 0, 11.1.3.

Final (objeto) de uma categoria: 0, 8.1.10.

Finito (morfismo) de pré-esquemas formais: III, 4.8.2.

Finito (pré-esquema formal) sobre $\mathfrak Y$,
$\mathfrak Y$-finito (pré-esquema formal): III, 4.8.2.

Funtor covariante canônico $\C\to\CHom(\C\op,\Set)$: 0, 8.1.2.

Género aritmético: III, 2.5.1.

Geometricamente conexa (fibra): III, 4.3.4.

Homomorfismo de estruturas algébricas sobre categorias: 0, 8.2.1.

Hipercohomologia de um funtor com respeito a um complexo $K^\bullet$:
0, 11.4.3.

Hipercohomologia de um bifuntor com respeito a dois complexos
$K^\bullet$, $K'^{\bullet}$: 0, 11.4.6.

Hiper-homologia de um funtor com respeito a um complexo $K_\bullet$:
0, 11.6.2.

Hiper-homologia de um bifuntor com respeito a dois complexos $K_\bullet$,
$K'_\bullet$: 0, 11.6.6.

Hiper-homologia de um funtor com respeito a um bicomplexo: 0, 11.7.2.

Localmente construtível (conjunto): 0, 9.1.11.

Lei de composição externa (interna) sobre os objetos de uma categoria: 0, 8.2.1.

\oldpage[III]{164}

$S$-$\C$-módulo filtrado, $\operatorname{gr}^\bullet(S)$-$\C$-módulo graduado,
$\operatorname{gr}^{\bullet\bullet}(S)$-$\C$-módulo bigraduado: 0, 13.6.6.

Morfismo de sucessões espectrais: 0, 11.1.2.

Número geométrico de componentes conexas de uma fibra: III, 4.3.4.

Objeto grupo em $\C$, $\C$-grupo, $\C$-anel, $\C$-módulo: 0, 8.2.3.

Objeto de bordos, objeto de cobordes, objeto de cociclos, objeto de
ciclos: 0, 11.2.1.

Objeto de imagens universais: 0, 13.1.1.

Objeto graduado associado a um sistema projetivo que satisfaz (ML): 0, 13.4.2.

Objeto graduado limitado inferiormente (superiormente): 0, 11.2.1.

Subconjunto próprio (sobre $\mathfrak Y$) de um pré-esquema formal: III, 3.4.1.

Plena (subcategoria): 0, 8.1.5.

Plenamente fiel (funtor): 0, 8.1.5.

Polinômio de Hilbert: III, 2.5.3.

Próprio (morfismo) de pré-esquemas formais: III, 3.4.1.

Representável (funtor): 0, 8.1.8.

Cohomológica, direita, esquerda, homológica, injetiva, livre, plana, projetiva
(resolução): 0, 11.4.1.

Resolução de Cartan--Eilenberg direita, injetiva: 0, 11.4.2.

Resolução de Cartan--Eilenberg esquerda, projetiva: 0, 11.6.1.

Retrocompacto (conjunto): 0, 9.1.1.

Estrito (sistema projetivo): 0, 13.4.2.

Sucessão espectral em uma categoria abeliana: 0, 11.1.1.

Fracamente convergente, regular, coregular, birregular (sucessão espectral):
0, 11.1.3.

Sucessão espectral degenerada: 0, 11.1.6.

Sucessão espectral de um funtor relativa a um objeto filtrado: 0, 13.6.4.

Sucessões espectrais de um bicomplexo: 0, 11.3.2 e 11.3.5.

Sucessões espectrais de hipercohomologia de um funtor com respeito a um complexo:
0, 11.4.3.

Sucessões espectrais de hiper-homologia de um funtor com respeito a um complexo:
0, 11.6.2.

Sucessões espectrais de hiper-homologia de um funtor com respeito a um bicomplexo:
0, 11.7.2.

Sistema de coeficientes: 0, 11.8.4.

Unirrama (anel, ponto): III, 4.3.6.

Universalmente aberto (morfismo): III, 4.3.9.
\end{flushleft}

\clearpage
\oldpage[III]{165}
\section*{ÍNDICE GERAL}

\begin{flushleft}
\textsc{Capítulo 0.} --- \textbf{Preliminares (continuação)} \dotfill 5

\medskip
\S~8. Funtores representáveis \dotfill 5

8.1. Funtores representáveis \dotfill 5

8.2. Estruturas algébricas em categorias \dotfill 9

\medskip
\S~9. Conjuntos construtíveis \dotfill 12

9.1. Conjuntos construtíveis \dotfill 12

9.2. Subconjuntos construtíveis de espaços noetherianos \dotfill 14

9.3. Funções construtíveis \dotfill 16

\medskip
\S~10. Suplemento sobre os módulos planos \dotfill 17

10.1. Relações entre módulos planos e módulos livres \dotfill 17

10.2. Critérios locais de planicidade \dotfill 18

10.3. Existência de extensões planas de anéis locais \dotfill 20

\medskip
\S~11. Suplemento sobre álgebra homológica \dotfill 23

11.1. Lembrete sobre as sucessões espectrais \dotfill 23

11.2. A sucessão espectral de um complexo filtrado \dotfill 27

11.3. As sucessões espectrais de um bicomplexo \dotfill 29

11.4. Hipercohomologia de um funtor com respeito a um complexo $K^\bullet$
\dotfill 32

11.5. Passagem ao limite indutivo em hipercohomologia \dotfill 35

11.6. Hiper-homologia de um funtor com respeito a um complexo $K_\bullet$
\dotfill 39

11.7. Hiper-homologia de um funtor com respeito a um bicomplexo
$K_{\bullet\bullet}$ \dotfill 41

11.8. Suplemento sobre a cohomologia dos complexos simpliciais \dotfill 43

11.9. Um lema sobre complexos de tipo finito \dotfill 46

11.10. Característica de Euler--Poincaré de um complexo de módulos de comprimento finito
\dotfill 48

\medskip
\S~12. Suplemento sobre a cohomologia de feixes \dotfill 49

12.1. Cohomologia de feixes de módulos sobre espaços anelados \dotfill 49

12.2. Imagens diretas superiores \dotfill 57

12.3. Suplementos sobre os funtores Ext de feixes \dotfill 60

12.4. Hipercohomologia do funtor de imagem direta \dotfill 62

\oldpage[III]{166}

\S~13. Limites projetivos em álgebra homológica \dotfill 64

13.1. A condição de Mittag--Leffler \dotfill 64

13.2. A condição de Mittag--Leffler para grupos abelianos \dotfill 65

13.3. Aplicação: cohomologia de um limite projetivo de feixes \dotfill 68

13.4. Condição de Mittag--Leffler e objetos graduados associados a sistemas
projetivos \dotfill 69

13.5. Limites projetivos de sucessões espectrais de complexos filtrados
\dotfill 71

13.6. Sucessão espectral de um funtor relativa a um objeto dotado de uma
filtração finita \dotfill 73

13.7. Funtores derivados de um limite projetivo de argumentos \dotfill 75

\medskip
\textsc{Capítulo III.} --- \textbf{Estudo cohomológico dos feixes coerentes}
\dotfill 81

\medskip
\S~1. Cohomologia dos esquemas afins \dotfill 82

1.1. Lembrete sobre o complexo de álgebra exterior \dotfill 82

1.2. Cohomologia de \v Cech de um recobrimento aberto \dotfill 85

1.3. Cohomologia de um esquema afim \dotfill 88

1.4. Aplicação à cohomologia de pré-esquemas quaisquer \dotfill 89

\medskip
\S~2. Estudo cohomológico dos morfismos projetivos \dotfill 95

2.1. Cálculos explícitos de certos grupos de cohomologia \dotfill 95

2.2. O teorema fundamental dos morfismos projetivos \dotfill 100

2.3. Aplicação a feixes graduados de álgebras e módulos \dotfill 102

2.4. Uma generalização do teorema fundamental \dotfill 107

2.5. Característica de Euler--Poincaré e polinômio de Hilbert \dotfill 109

2.6. Aplicação: critérios de amplitude \dotfill 111

\medskip
\S~3. Teorema de finitude para os morfismos próprios \dotfill 115

3.1. O lema de dévissage \dotfill 115

3.2. O teorema de finitude: caso dos esquemas ordinários \dotfill 116

3.3. Generalização do teorema de finitude (esquemas ordinários) \dotfill 118

3.4. Teorema de finitude: caso dos esquemas formais \dotfill 119

\medskip
\S~4. O teorema fundamental dos morfismos próprios. Aplicações \dotfill 122

4.1. O teorema fundamental \dotfill 122

4.2. Casos particulares e variantes \dotfill 129

4.3. Teorema de conexão de Zariski \dotfill 130

4.4. «Teorema principal» de Zariski \dotfill 135

4.5. Completamentos de módulos de homomorfismos \dotfill 138

4.6. Relações entre morfismos formais e morfismos ordinários \dotfill 139

4.7. Um critério de amplitude \dotfill 145

4.8. Morfismos finitos de pré-esquemas formais \dotfill 146

\oldpage[III]{167}

\S~5. Um teorema de existência para os feixes algébricos coerentes \dotfill 149

5.1. Enunciado do teorema \dotfill 149

5.2. Demonstração do teorema de existência: caso projetivo e quase-projetivo
\dotfill 151

5.3. Demonstração do teorema de existência: caso geral \dotfill 154

5.4. Aplicação: comparação entre morfismos de pré-esquemas ordinários e
morfismos de pré-esquemas formais. Pré-esquemas formais algebrizáveis \dotfill 156

5.5. Uma decomposição de certos pré-esquemas \dotfill 159

\medskip
\textsc{Bibliografia (continuação)} \dotfill 161

\textsc{Índice de notações} \dotfill 162

\textsc{Índice terminológico} \dotfill 163

\begin{flushright}
\emph{Recebido em 28 de junho de 1961.}
\end{flushright}
\end{flushleft}


\clearpage
\phantomsection
\addcontentsline{toc}{section}{Referências consolidadas da edição}
\section*{Referências consolidadas da edição}
\nocite{*}
\bibliography{the}
\bibliographystyle{amsalpha}

\end{document}
